\documentclass[11pt,a4paper,leqno]{article}

\usepackage[T1]{fontenc}
\usepackage[utf8]{inputenc}
\usepackage[english,french]{babel}
\usepackage{amsmath,amssymb,amsfonts,mathrsfs}
\usepackage{mathtools}
\usepackage{tikz}
\usetikzlibrary{cd,arrows.meta}
\usepackage{geometry}
\geometry{margin=2.4cm}
\usepackage{enumitem}
\usepackage[normalem]{ulem}
\usepackage{hyperref}
\hypersetup{
  colorlinks=true,
  linkcolor=blue,
  citecolor=blue,
  urlcolor=blue,
  hypertexnames=false,
  pdftitle={SGA 7-I: Monodromy Groups in Algebraic Geometry - English Reader},
  pdfauthor={Interlanguage Mathematical Translation Project},
  pdfsubject={Source-first English translation from the complete French working transcription of LNM 288},
  pdfkeywords={SGA 7; monodromy; algebraic geometry; English translation}
}
\let\mathunderline\underline
\renewcommand{\underline}[1]{\ifmmode\mathunderline{#1}\else\uline{#1}\fi}
\emergencystretch=3em

\DeclareUnicodeCharacter{2192}{\ensuremath{\rightarrow}}
\DeclareUnicodeCharacter{2297}{\ensuremath{\otimes}}
\DeclareUnicodeCharacter{2295}{\ensuremath{\oplus}}
\DeclareUnicodeCharacter{2228}{\ensuremath{\vee}}
\DeclareUnicodeCharacter{223C}{\ensuremath{\sim}}
\DeclareUnicodeCharacter{2248}{\ensuremath{\approx}}
\DeclareUnicodeCharacter{00D7}{\ensuremath{\times}}
\DeclareUnicodeCharacter{03B1}{\ensuremath{\alpha}}
\DeclareUnicodeCharacter{03B2}{\ensuremath{\beta}}
\DeclareUnicodeCharacter{00A7}{\S}
\DeclareUnicodeCharacter{010C}{\v{C}}
\DeclareUnicodeCharacter{00C3}{\~{A}}
\DeclareUnicodeCharacter{02D9}{\.{}}

\providecommand{\Zl}{\mathbb{Z}_{\ell}}
\providecommand{\Ql}{\mathbb{Q}_{\ell}}
\providecommand{\Gal}{\operatorname{Gal}}
\providecommand{\Spec}{\operatorname{Spec}}
\providecommand{\Isom}{\operatorname{Isom}}
\providecommand{\Aut}{\operatorname{Aut}}
\providecommand{\Hom}{\operatorname{Hom}}
\providecommand{\Ext}{\operatorname{Ext}}
\providecommand{\Biext}{\operatorname{Biext}}
\providecommand{\Pic}{\operatorname{Pic}}
\providecommand{\Coker}{\operatorname{Coker}}
\providecommand{\Ker}{\operatorname{Ker}}
\providecommand{\Im}{\operatorname{Im}}
\providecommand{\Tor}{\operatorname{Tor}}
\providecommand{\prof}{\operatorname{prof}}
\providecommand{\codim}{\operatorname{codim}}
\providecommand{\pr}{\operatorname{pr}}
\providecommand{\id}{\mathrm{id}}
\providecommand{\sym}{\mathrm{sym}}
\providecommand{\Corr}{\operatorname{Corr}}
\providecommand{\Proj}{\operatorname{Proj}}
\providecommand{\Comp}{\operatorname{Comp}}
\providecommand{\cC}{\mathcal{C}}
\providecommand{\cY}{\mathcal{Y}}
\providecommand{\cT}{\mathcal{T}}
\providecommand{\cPsi}{\Psi}
\providecommand{\Dcat}{D}
\providecommand{\La}{\Lambda}
\providecommand{\Ghat}{\widehat{G}}
\providecommand{\uHom}{\underline{\operatorname{Hom}}}
\providecommand{\uExt}{\underline{\operatorname{Ext}}}
\providecommand{\uTor}{\underline{\operatorname{Tor}}}
\providecommand{\cl}{\operatorname{c\ell}}
\providecommand{\e}{\mathrm{e}}

\begin{document}
\selectlanguage{english}
\begin{titlepage}
\centering
\vspace*{2.2in}
{\LARGE\textbf{SGA 7-I}}\\[0.6em]
{\Large Monodromy Groups in Algebraic Geometry}\\[1.5em]
{\large Complete English reader}\\[1.2em]
{\normalsize Exposés I, II, VI, VII, VIII, and IX}
\vfill
\end{titlepage}
\tableofcontents
\clearpage
\phantomsection\addcontentsline{toc}{section}{Exposé I}
% BEGIN INCLUDED SGA7I_English_current_buildable_source_20260823/components/01_expose_I_opening_through_0_0.tex
% SOURCE AUTHORITY:
% expose_I_body.tex lines 1--54; scan indices 12--13; source folios 1--2.
% Hard stop before source line 57, paragraph 0.1.

%%%%%%%%%%  scan idx 12  |  folio 1  %%%%%%%%%%

\SGARefAliasTarget{sga7.sga7i.ei.nroot}SGA 7

\underline{Expos\'e I}

\underline{Summary of the First Lectures by A. Grothendieck}

\underline{written up by P. Deligne}

\underline{Contents}

0. Preliminaries

1. Arithmetic proof of the monodromy theorem

2. Vanishing cycles

3. Geometric proof of the monodromy theorem

4. Vanishing criteria for the sheaves of vanishing cycles

5. Action of monodromy on the $\pi_1$

6. Appendix by P. Deligne: arithmetic proof of the stable reduction theorem

0. \underline{Preliminaries}

\SGARefTarget{sga7.sga7i.ei.n0-0}{0.0}. \underline{Terminology}. Recall the following definitions:

(\SGARefTarget{sga7.sga7i.ei.n0-0-1}{0.0.1}) \underline{trait}: the spectrum of a discrete valuation ring. The generic and closed points of a trait are often denoted by $\eta$ and $s$.

(\SGARefTarget{sga7.sga7i.ei.n0-0-2}{0.0.2}) \underline{strictly local}: for a ring, local henselian with separably closed residue field; for a scheme, the spectrum of such a ring.

(\SGARefTarget{sga7.sga7i.ei.n0-0-3}{0.0.3}) \underline{geometric point of $S$} (in this expos\'e): a morphism $\bar{x}:\operatorname{Spec}(\bar{k})\longrightarrow S$ with image $x\in S$, where $k(\bar{x})\underset{\mathrm{dfn}}{=}\bar{k}$ is a separable closure of $k(x)$. For a henselian trait $S=\operatorname{Spec}(V)$, a generic geometric point of $S$ (that is, one with image $\eta$) is often denoted by $\bar{\eta}$. The residue field of the valuation ring $\bar{V}$, the integral closure of $V$ in $k(\bar{\eta})$, is a separably closed algebraic extension of $k(s)$ and defines a geometric point $\bar{s}$ of $S$ centered at $s$.

%%%%%%%%%%  scan idx 13  |  folio 2  %%%%%%%%%%

(\SGARefTarget{sga7.sga7i.ei.n0-0-4}{0.0.4}) \underline{strict localization of $X$ at the geometric point $\bar{x}$}: SGA 4 VIII 4.4.

(\SGARefTarget{sga7.sga7i.ei.n0-0-5}{0.0.5}) \underline{essentially of finite type over $S$}: for a local (respectively, henselian local; respectively, strictly local) $S$-scheme, the spectrum of a local ring (respectively, the henselization of a local ring; respectively, a strict localization) of a scheme of finite type over $S$.

(\SGARefTarget{sga7.sga7i.ei.n0-0-6}{0.0.6}) $F\longmapsto F(n)$: Tate twist.

(\SGARefTarget{sga7.sga7i.ei.n0-0-7}{0.0.7}) An endomorphism $T$ of a finite-dimensional vector space is \underline{quasi-unipotent} if its eigenvalues are roots of unity, that is, if there exist $N\geqslant 1$ and $M\geqslant 1$ such that $(T^N-1)^M=0$.

% END INCLUDED SGA7I_English_current_buildable_source_20260823/components/01_expose_I_opening_through_0_0.tex

% BEGIN INCLUDED SGA7I_English_current_buildable_source_20260823/components/02_expose_I_0_1_through_1_4.tex
% SOURCE AUTHORITY:
% expose_I_body.tex lines 57--246; scan indices 13--18; source folios 2--7.
% Hard stop before source line 249, section 2.

\SGARefTarget{sga7.sga7i.ei.n0-1}{0.1}. \underline{Review of $\ell$-adic cohomology}

Let $X$ be a scheme of finite type over an algebraically closed field $k$ of characteristic exponent $p$, and let $\ell$ be a prime different from $p$.

The $\ell$-adic cohomology groups of $X$ were defined in SGA 4 and SGA 5.

a) $H^i(X,\mathbb{Z}/\ell^n)$ is the $i$-th cohomology group of the \`etale site $X_{\mathrm{et}}$ of $X$ with coefficients in $\mathbb{Z}/\ell^n$;

b) by definition,
\[
H^i(X,\mathbb{Z}_\ell)=\varprojlim_n H^i(X,\mathbb{Z}/\ell^n);
\]

c) by definition,
\[
H^i(X,\mathbb{Q}_\ell)=H^i(X,\mathbb{Z}_\ell)\otimes_{\mathbb{Z}_\ell}\mathbb{Q}_\ell.
\]

The $\ell$-adic cohomology groups of $X$ with proper supports are defined in the same way from the groups $H^i_c(X,\mathbb{Z}/\ell^n)$ of SGA 4 XVII.

In each of the following cases, one knows how to prove that these groups $H^i(X)$ have reasonable finiteness properties:

(\SGARefTarget{sga7.sga7i.ei.n0-1-1}{0.1.1}) for cohomology with proper supports;

(\SGARefTarget{sga7.sga7i.ei.n0-1-2}{0.1.2}) when $X$ is proper over $k$;

(\SGARefTarget{sga7.sga7i.ei.n0-1-3}{0.1.3}) when $X$ is smooth over $k$ (by (\SGARefLink{sga7.sga7i.ei.n0-1-1}{0.1.1}) and Poincar\'e duality, SGA 4 XVIII);

(\SGARefTarget{sga7.sga7i.ei.n0-1-4}{0.1.4}) when Hironaka's resolution of singularities is available for schemes of finite type over $k$ of dimension at most $\dim(X)$;

(\SGARefTarget{sga7.sga7i.ei.n0-1-5}{0.1.5}) when $i=0,1$.

%%%%%%%%%%  scan idx 14  |  folio 3  %%%%%%%%%%

In all these cases, the groups $H^i(X,\mathbb{Z}/\ell^n)$ are finite, the groups $H^i(X,\mathbb{Z}_\ell)$ are finitely generated, the groups $H^i(X,\mathbb{Q}_\ell)$ are finite-dimensional, and (except perhaps in case (\SGARefLink{sga7.sga7i.ei.n0-1-5}{0.1.5})) there are split short exact sequences

\begin{equation*}
0 \longrightarrow H^i(X,\mathbb{Z}_\ell)\otimes \mathbb{Z}/\ell^n
\longrightarrow H^i(X,\mathbb{Z}/\ell^n)
\longrightarrow \Tor_1(H^{i+1}(X,\mathbb{Z}_\ell),\mathbb{Z}/\ell^n)
\longrightarrow 0.
\end{equation*}
Moreover, in each case the proof of the finiteness theorem also proves a generic constructibility theorem: if $k$ is the field of fractions of an integral noetherian scheme $S$, and $X$ is the generic fibre of a morphism $f_1:X_1\longrightarrow S$ of finite type, then there exists a nonempty open subset $U\subset S$ such that over $U$ the sheaves $R^i f_*\mathbb{Z}/\ell^n$ are locally constant and their formation commutes with every base change.

\SGARefTarget{sga7.sga7i.ei.n0-2}{0.2}. We no longer suppose that $k$ is algebraically closed. Let $\bar{k}$ be an algebraic closure of $k$ (or a separable closure, which amounts to the same thing here). We consider the ``geometric'' cohomology groups $H^i(X_{\bar{k}})$, where $X_{\bar{k}}$ is obtained from $X$ by extension of scalars from $k$ to $\bar{k}$. By transport of structure, the Galois group $\Gal(\bar{k}/k)$ acts on these groups. In $\mathbb{Q}_\ell$-cohomology, in any of the cases (\SGARefLink{sga7.sga7i.ei.n0-1-1}{0.1.1})--(\SGARefLink{sga7.sga7i.ei.n0-1-5}{0.1.5}), this gives continuous representations of $\Gal(\bar{k}/k)$ in a linear group $\mathrm{GL}(n,\mathbb{Q}_\ell)$.

\SGARefTarget{sga7.sga7i.ei.n0-3}{0.3}. \underline{Review of traits} (\SGARefLink{sga7.sga7i.ei.n0-0-1}{0.0.1})

For the proofs, we refer to [4]. Let $S$ be a henselian trait, and let $\eta$, $s$, $V$, $\bar V$, $\bar\eta$, and $\bar s$ be as in (\SGARefLink{sga7.sga7i.ei.n0-0-1}{0.0.1}) and (\SGARefLink{sga7.sga7i.ei.n0-0-3}{0.0.3}). Let $p$ denote the characteristic exponent of $k(s)$, and let $\Gal(\bar\eta/\eta)$ (respectively, $\Gal(\bar s/s)$) denote the Galois group $\Gal(k(\bar\eta)/k(\eta))$ (respectively, the corresponding group for $\bar s/s$). By transport of structure, $\Gal(\bar\eta/\eta)$ acts on $k(\bar s)$; hence there is a homomorphism from $\Gal(\bar\eta/\eta)$ onto $\Gal(\bar s/s)$ whose kernel is the \underline{inertia group} $I$. The subfield of $k(\bar\eta)$ fixed by $I$ is a maximal unramified extension of $k(\eta)$.

The profinite group $I$ has a largest pro-$p$ subgroup $P$. The quotient $I/P$ is the \underline{tame inertia group}. Canonically,

\begin{equation*}
I/P\simeq \varprojlim_n\mu_n(k(\bar s))
\underset{\mathrm{dfn}}{=}\widehat{\mathbb{Z}}(1)(k(\bar s)).
\end{equation*}

%%%%%%%%%%  scan idx 15  |  folio 4  %%%%%%%%%%

If $\phi_n$ is the canonical map from $I/P$ to $\mu_n(\bar k)$ and $u$ is an element of $\bar V$ of valuation $1/n$, then, for $\sigma\in I$,

\begin{equation*}
\sigma(u)\equiv\phi_n(\sigma)\,u
\qquad (\text{modulo elements of valuation greater than }1/n).
\end{equation*}
% Printed source: \phi_n(u)u. The argument must be \sigma because \phi_n is
% a character of I/P; the English text transparently corrects this source typo.
In summary,

\begin{equation}\tag{\SGARefTarget{sga7.sga7i.ei.n0-3-1}{0.3.1}}
\begin{aligned}
\Gal(\bar\eta/\eta)&\supset I\supset P,\\
\Gal(\bar\eta/\eta)/I&\simeq\Gal(\bar s/s),\\
I/P&\simeq\widehat{\mathbb{Z}}(1)(k(\bar s)),\\
P&:\ \text{a pro-$p$ group}\quad(\{e\}\text{ if }p=1),
\end{aligned}
\end{equation}
and the action (by inner automorphisms in $\Gal(\bar\eta/\eta)$) of $\Gal(\bar\eta/\eta)/I$ on $I/P$ is identified with the natural action of $\Gal(\bar s/s)$ on $\widehat{\mathbb{Z}}(1)(k(\bar s))$.

\SGARefTarget{sga7.sga7i.ei.n0-4}{0.4}. More generally, let $S$ be strictly local and regular, let $D$ be a regular divisor in $S$, and let $p'$ be the characteristic exponent of the generic point of $D$. Let $\bar\eta$ be a geometric generic point of $S$, and let $\bar s$ be the resulting geometric point centered at the closed point $s$. By Abhyankar's lemma, the fundamental group $\pi_1(S-D,\bar\eta)$ has a filtration analogous to (\SGARefLink{sga7.sga7i.ei.n0-3-1}{0.3.1}):

\begin{equation}\tag{\SGARefTarget{sga7.sga7i.ei.n0-4-1}{0.4.1}}
\begin{aligned}
\pi_1(S-D,\bar\eta)&\supset I\supset P,\\
\pi_1(S-D,\bar\eta)/I&\simeq\Gal(\bar s/s),\\
I/P&\simeq\widehat{\mathbb{Z}}(1)(k(\bar s)),\\
P&:\ \text{with no quotient of order prime to }p'.
\end{aligned}
\end{equation}
The action by inner automorphisms of $\pi_1(S-D,\bar\eta)/I$ on $I/P$ is identified with the natural action of $\Gal(\bar s/s)$ on $\widehat{\mathbb{Z}}(1)(k(\bar s))$.

\SGARefTarget{sga7.sga7i.ei.n0-5}{0.5}. \underline{N\'eron desingularization}

The method of N\'eron desingularization gives the following result.

\underline{Lemma} \SGARefTarget{sga7.sga7i.ei.n0-5-1}{0.5.1}. \underline{Let $S\longrightarrow S_o$ be a morphism of henselian traits, with $S=\Spec(V)$ and $S_o=\Spec(V_o)$, and let $\pi$ be a uniformizer of $V_o$. Suppose that $\pi$ is a uniformizer of $V$, and that the extensions $k(s)/k(s_o)$ and $k(\eta)/k(\eta_o)$ are separable. Then $V$ is a filtered direct limit of smooth henselian $V_o$-algebras $B_i$ essentially of finite type over $V_o$.}

%%%%%%%%%%  scan idx 16  |  folio 5  %%%%%%%%%%

To construct the desired $V_o$-algebras, one takes $V_o\subset B\subset V$, with $B$ of finite type, and desingularizes $\Spec(B)$ by N\'eron's method ([1], \S~4).

Let $S=\Spec(V)$ be a henselian trait.

(\SGARefTarget{sga7.sga7i.ei.n0-5-2}{0.5.2}) If $S$ is of equal characteristic, with uniformizer $\pi$, one may apply (\SGARefLink{sga7.sga7i.ei.n0-5-1}{0.5.1}) with $S_o=\mathbf{F}_p[\pi]_{(\pi)}$. The extension $s/s_o$ is separable because $\mathbf{F}_p$ is perfect; the extension $\eta/\eta_o$ is separable because $\pi$, being a uniformizer, is not a $p$-th power (if $p\neq1$), and hence belongs to a $p$-basis. It follows that $S$ is a limit of spectra of smooth henselian $\mathbf{F}_p[\pi]_{(\pi)}$-algebras essentially of finite type.

(\SGARefTarget{sga7.sga7i.ei.n0-5-3}{0.5.3}) If $S$ is complete and of mixed characteristic, with perfect residue field $k$, then $V$ is an Eisenstein extension of the ring of Witt vectors (that is, the Cohen ring) $W(k)$:

\begin{equation*}
V\simeq W(k)[\pi]/\left(\pi^n+\sum_{i=0}^{n-1}a_i\pi^i\right).
\end{equation*}
If the $a_i$ are perturbed slightly, one obtains an isomorphic extension. We may therefore suppose that the $a_i$ belong to $W(k')$, where $k'$ is of finite type over $\mathbf{F}_p$. Let $k_o$ be the perfect closure of $k'$ in $k$, and set

\begin{equation*}
V_o=W(k_o)[\pi]/\left(\pi^n+\sum_{i=0}^{n-1}a_i\pi^i\right).
\end{equation*}
Lemma \SGARefLink{sga7.sga7i.ei.n0-5-1}{0.5.1} applies to $V/V_o$, and the residue field of $V_o$ is a radicial extension of a field of finite type over $\mathbf{F}_p$.

1. \underline{Arithmetic proof of the monodromy theorem}

The reader will find in [5], Appendix, a proof of the following result of A. Grothendieck.

\underline{Proposition} \SGARefTarget{sga7.sga7i.ei.n1-1}{1.1}. \underline{Let $S$ be a henselian trait, let $\ell$ be a prime different from the characteristic exponent $p$ of $k(s)$, and let $\rho$ be a continuous representation of $\Gal(\bar\eta/\eta)$ in $\mathrm{GL}(n,\mathbb{Q}_\ell)$. Suppose that the following condition holds:}

$(*_\ell)$ \quad \underline{No finite extension of $k(s)$ contains all roots of unity of $\ell$-power order.}

%%%%%%%%%%  scan idx 17  |  folio 6  %%%%%%%%%%

\underline{Then there exists an open subgroup $I_1$ of the inertia group $I$ such that $\rho(\sigma)$ is unipotent for every $\sigma\in I_1$.}

Condition $(*_\ell)$ is automatically satisfied when $k(s)$ is of finite type over the prime field (or is radicial over such a field).

The following theorem follows immediately from \SGARefLink{sga7.sga7i.ei.n1-1}{1.1}.

\underline{Monodromy theorem} \SGARefTarget{sga7.sga7i.ei.n1-2}{1.2}. \underline{With the preceding notation, let $X$ be a scheme of finite type over $\eta$, and let $H$ be a cohomology space of $X_{\bar\eta}$ with coefficients in $\mathbb{Q}_\ell$, of one of the types (\SGARefLink{sga7.sga7i.ei.n0-1-1}{0.1.1})--(\SGARefLink{sga7.sga7i.ei.n0-1-5}{0.1.5}). If condition $(*_\ell)$ is satisfied, the action on $H$ of every element of $I$ is quasi-unipotent.}
% Printed source refers to (1.1.1)--(1.1.5); these do not exist. The intended
% reference is the list (0.1.1)--(0.1.5) immediately above.

The limit arguments of \SGARefLink{sga7.sga7i.ei.n0-5}{0.5} allow us to dispense with hypothesis $(*_\ell)$.

\underline{Variant} \SGARefTarget{sga7.sga7i.ei.n1-3}{1.3}. \underline{Condition $(*_\ell)$ in \SGARefLink{sga7.sga7i.ei.n1-2}{1.2} is superfluous.}

Let $H'_o$ be a $\mathbb{Z}_\ell$-cohomology space of $X_{\bar\eta}$ giving rise to $H$, and put $H_o=H'_o/\text{torsion}$. By hypothesis,
$H=H_o\otimes_{\mathbb{Z}_\ell}\mathbb{Q}_\ell$, and the action of $\Gal(\bar\eta/\eta)$ is induced by

\begin{equation*}
\rho:\Gal(\bar\eta/\eta)\longrightarrow\mathrm{GL}(H_o).
\end{equation*}
Let $I_{\ell'}$ be the largest prime-to-$\ell$ subgroup of $I$. Its image in $\mathrm{GL}(H_o)$ is an $\ell$-adic Lie group of order prime to $\ell$, and is therefore finite.

Make an extension of traits $S'\longrightarrow S$. This does not change $H_o$ (SGA 4 XVI 6.6), and $\rho(I')$ has finite index in $\rho(I)$, because $I'/P'$ has finite index in $I/P$, by what precedes. It therefore suffices to prove the theorem over $S'$; we may suppose that

a) $\Gal(\bar\eta/\eta)$ acts trivially on $H_o/\ell H_o$;

b) if $S$ is of mixed characteristic, then $S$ is complete with perfect residue field (EGA $\mathrm{O}_{III}$ 10.3.1).

Put $S=\Spec(V)$, define $V_o$ as in (\SGARefLink{sga7.sga7i.ei.n0-5-2}{0.5.2}) and (\SGARefLink{sga7.sga7i.ei.n0-5-3}{0.5.3}), and apply (\SGARefLink{sga7.sga7i.ei.n0-5-1}{0.5.1}). We have

\begin{equation*}
V=\varinjlim_i B_i.
\end{equation*}

%%%%%%%%%%  scan idx 18  |  folio 7  %%%%%%%%%%

Here $S_i=\Spec(B_i)$ is essentially of finite type and smooth over $V_o$. Let $D_i\subset S_i$ be defined by the equation $\pi=0$. For sufficiently large $i$, $H_o$ comes from a twisted constant $\mathbb{Z}_\ell$-sheaf $\mathcal H_i$ on $S_i-D_i$ (by ``generic constructibility,'' (\SGARefLink{sga7.sga7i.ei.n0-1}{0.1})), and $\mathcal H_i/\ell\mathcal H_i$ is constant.

The filtrations (\SGARefLink{sga7.sga7i.ei.n0-3-1}{0.3.1}) and (\SGARefLink{sga7.sga7i.ei.n0-4-1}{0.4.1}) give rise to commutative diagrams:

\begin{equation}\tag{\SGARefTarget{sga7.sga7i.ei.n1-3-1}{1.3.1}}
\begin{tikzcd}[column sep=small,row sep=large]
\Gal(\bar\eta/\eta) \arrow[d] & \supset & I \arrow[d] & \supset & P \arrow[d] \\
\pi_1(S_i-D_i,\bar\eta) & \supset & I_i & \supset & P_i
\end{tikzcd}
\qquad\text{and}\qquad
\begin{tikzcd}[row sep=large]
I/P \arrow[r,dash,"\sim"] \arrow[d] & \widehat{\mathbb Z}(1) \arrow[d,equal] \\
I_i/P_i \arrow[r,dash,"\sim"] & \widehat{\mathbb Z}(1)
\end{tikzcd}
\end{equation}

% The authority has two distinct commutative diagrams separated by the word
% ``and.'' The horizontal isomorphisms in the right-hand diagram are rules
% labelled \sim, not arrows; its right vertical connector is equality.
The representation on $H_o$ of $\Gal(\bar\eta/\eta)$ is induced by an action of $\pi_1(S_i-D_i,\bar\eta)$ on $H_o$ that is trivial on $H_o/\ell H_o$. The group
\[
\Ker\bigl(\mathrm{GL}(H_o)\longrightarrow\mathrm{GL}(H_o/\ell H_o)\bigr)
\]
is a pro-$\ell$ group. Thus $P_i$ acts trivially, and we apply the following variant of \SGARefLink{sga7.sga7i.ei.n1-1}{1.1}.
% Printed source writes GL(H) -> GL(H/ell H), but the lattice being acted on
% and reduced here is H_o. The English text makes this evident subscript repair.

\underline{Variant} \SGARefTarget{sga7.sga7i.ei.n1-4}{1.4}. \underline{Retain the notation of \SGARefLink{sga7.sga7i.ei.n0-4}{0.4}, and let $\rho$ be a continuous representation of $\Gal(\bar\eta/\eta)$ in $\mathrm{GL}(n,\mathbb{Q}_\ell)$. Suppose that $k(s)$ satisfies $(*_\ell)$ and that $\rho(P)$ is finite. Then the action of $I$ is quasi-unipotent.}
% Printed source has Gal(n,Q_l), an unambiguous typographical error for GL.

% END INCLUDED SGA7I_English_current_buildable_source_20260823/components/02_expose_I_0_1_through_1_4.tex

% BEGIN INCLUDED SGA7I_English_current_buildable_source_20260823/components/03_expose_I_section_2.tex
% SOURCE AUTHORITY:
% expose_I_body.tex lines 249--366; scan indices 18--21; source folios 7--10.
% Hard stop before source line 368, section 3.

2. \underline{Vanishing cycles}

\SGARefTarget{sga7.sga7i.ei.n2-1}{2.1}. Let $S$ be a strictly local trait, and retain the notation of \SGARefLink{sga7.sga7i.ei.n0-3}{0.3}. Thus, by hypothesis, $s=\bar s$. Let $f:X\longrightarrow S$ be a scheme of finite type over $S$. The formalism of vanishing cycles is a tool for studying the relation between the cohomology of $X_s$ and that of $X_{\bar\eta}$, as well as the action of the inertia group on the latter, using information about $f$ that is \underline{local} on $X$.

A key result of this kind is the specialization theorem (SGA 4 XVI 2.2), according to which $X_s$ and $X_{\bar\eta}$ have the same cohomology when $f$ is proper and smooth.

In the transcendental theory, if $f:X\longrightarrow D$ is a proper morphism from an analytic space $X$ to the disk $D$, vanishing cycles arise as follows:

a) After shrinking $D$ if necessary, one may suppose that $X$ is a topological fibre bundle over the punctured disk $D^*$ (Thom's isotopy theorem). In particular, the fibres

%%%%%%%%%%  scan idx 19  |  folio 8  %%%%%%%%%%

$X_t$ $(t\in D^*)$ are all isomorphic.

b) It is not unreasonable to picture the special fibre $X_s$ as obtained from \lq\lq the\rq\rq\ general fibre $X_t$ by contracting polyhedra in $X_t$; one would then have a map $\pi:X_t\longrightarrow X_s$. The method of vanishing cycles consists in studying the difference between the cohomology of $X_s$ and that of $X_t$ by means of the Leray spectral sequence for $\pi$.

The geometric theory explained below is very close to this transcendental theory: the general point $t\in D$ must be replaced by $\bar\eta$, the geometric generic point of $S$. Conversely, the transcendental theory can be modeled on the geometric theory by replacing $t$ or $\bar\eta$ with the universal cover $\widetilde D^*$ of $D^*$, and $X_t$ or $X_{\bar\eta}$ with $X\times_D\widetilde D^*$. This eliminates many $\varepsilon$'s and $\eta$'s, but sometimes loses information.

\SGARefTarget{sga7.sga7i.ei.n2-2}{2.2}. Fix a torsion ring $\Lambda$ in which $p$ is invertible (for example, $\Lambda=\mathbb Z/\ell^n\mathbb Z$, with $\ell$ prime to $p$). Consider the morphisms

\begin{center}
\begin{tikzcd}[column sep=large,row sep=huge]
X_s \arrow[r,"i"] \arrow[d,"f_s"] & X \arrow[d,"f"] & & X_{\bar\eta} \arrow[ll,"\bar j"] \arrow[d,"f_{\bar\eta}"] \\
s \arrow[r] & S & \eta \arrow[l] & \bar\eta \arrow[l]
\end{tikzcd}
\end{center}
% Authority geometry: seven nodes and eight plain arrows. The top row is
% X_s -> X <- X_{bar eta}; the bottom row is s -> S <- eta <- bar eta;
% the three vertical arrows are f_s, f, and f_{bar eta}.
The sheaves of vanishing cycles $\psi^i$ on $X_s$ are

\begin{equation*}\tag{\SGARefTarget{sga7.sga7i.ei.n2-2-1}{2.2.1}}
\psi^i=i^*R^i\bar j_*\Lambda.
\end{equation*}
If $f$ is \underline{proper}, then (SGA 4 XII 5.5)

\begin{equation*}\tag{\SGARefTarget{sga7.sga7i.ei.n2-2-2}{2.2.2}}
i^*:H^p(X,R^q\bar j_*\Lambda)\xrightarrow{\ \sim\ }
H^p(X_s,i^*R^q\bar j_*\Lambda),
\end{equation*}
% The fresh French TeX drops the bar from j on the right of (2.2.2); the
% printed authority has bar-j on both sides.
and the Leray spectral sequence for $\bar j$ becomes

\begin{equation*}\tag{\SGARefTarget{sga7.sga7i.ei.n2-2-3}{2.2.3}}
H^p(X_s,\psi^q)\Longrightarrow H^{p+q}(X_{\bar\eta},\Lambda).
\end{equation*}
When nothing happens at infinity, this spectral sequence can sometimes be obtained without a properness hypothesis. That some hypothesis at infinity is necessary is already clear from the trivial case $X_s=\varnothing$.

%%%%%%%%%%  scan idx 20  |  folio 9  %%%%%%%%%%

The definitions immediately give:

\underline{Proposition} \SGARefTarget{sga7.sga7i.ei.n2-3}{2.3}. \underline{Let $\bar x$ be a geometric point of $X_s$ (for example, a closed point), let $X_{(\bar x)}$ be the strict henselization of $X$ at $\bar x$, and let $X_{(\bar x)\bar\eta}$ be the geometric generic fibre of the projection $X_{(\bar x)}\longrightarrow S$. Then}

\begin{equation*}
\psi^i_{\bar x}=H^i(X_{(\bar x)\bar\eta},\Lambda).
\end{equation*}
% The subscript bar-x is visible in the printed formula but absent from the
% fresh French TeX; it is restored here.

\underline{Corollary} \SGARefTarget{sga7.sga7i.ei.n2-4}{2.4}. \underline{Wherever $f$ is smooth, $\psi^i=0$ for $i>0$ and $\psi^0=\Lambda$.}

This is SGA 4 XV 2.1.

The group $\Gal(\bar\eta/\eta)$ (which, by hypothesis, equals the inertia group) acts on the $\psi^i$ by transport of structure, and

\underline{Scholium} \SGARefTarget{sga7.sga7i.ei.n2-5}{2.5}. \underline{The spectral sequence (\SGARefLink{sga7.sga7i.ei.n2-2-3}{2.2.3}) is $\Gal(\bar\eta/\eta)$-equivariant.}

\underline{Variant} \SGARefTarget{sga7.sga7i.ei.n2-6}{2.6}. Put $\bar S=\Spec(\bar V)$ and $\bar X=X\times_S\bar S$. The scheme $X_{\bar\eta}$ is the complement of $X_s=X_{\bar s}$ in $\bar X$. Put

\begin{equation*}
\phi^i=\underline{H}^{\,i+1}_{X_s}(\bar X,\Lambda)
\end{equation*}
(the local-cohomology sheaves). There is an exact sequence of sheaves

\begin{equation*}
0\longrightarrow\phi^{-1}\longrightarrow\Lambda\longrightarrow\psi^0
\longrightarrow\phi^0\longrightarrow0,
\end{equation*}
and, for $i>0$, isomorphisms $\psi^i\xrightarrow{\ \sim\ }\phi^i$. If $f$ is smooth, the $\phi^i$ vanish.

Put $\phi^i_{\mathrm{gl}}=H^i_{X_s}(\bar X,\Lambda)$, the global analogue of $\phi^i$. There is a spectral sequence

\begin{equation*}\tag{\SGARefTarget{sga7.sga7i.ei.n2-6-1}{2.6.1}}
H^p(X_s,\phi^q)\Longrightarrow\phi^{p+q}_{\mathrm{gl}},
\end{equation*}
and a long exact cohomology sequence which, when $f$ is proper, becomes

\begin{equation*}\tag{\SGARefTarget{sga7.sga7i.ei.n2-6-2}{2.6.2}}
\cdots\longrightarrow H^i(X_s,\Lambda)\longrightarrow H^i(X_{\bar\eta},\Lambda)
\longrightarrow\phi^i_{\mathrm{gl}}\longrightarrow\cdots.
\end{equation*}

\underline{Variant} \SGARefTarget{sga7.sga7i.ei.n2-7}{2.7}. Let $K_t$ be the maximal tamely ramified extension of $K$ in $\bar K$. Put $X_t=X\times_S\Spec(K_t)$, and let $\bar j_t$ be the projection from $X_t$ to $X$. Define tame variants of the $\psi^i$ by

%%%%%%%%%%  scan idx 21  |  folio 10  %%%%%%%%%%

\begin{equation*}
\psi^i_t=R^i\bar j_{t*}\Lambda.
\end{equation*}
Since $P$ is a pro-$p$ group and $p$ is invertible in $\Lambda$, one obtains

\begin{equation}\tag{\SGARefTarget{sga7.sga7i.ei.n2-7-1}{2.7.1}}
H^*(X_t,\Lambda)=H^*(X_{\bar\eta},\Lambda)^P,
\end{equation}
\begin{equation}\tag{\SGARefTarget{sga7.sga7i.ei.n2-7-2}{2.7.2}}
\psi^i_t=(\psi^i)^P.
\end{equation}
When $f$ is proper, the Leray spectral sequence for $\bar j_t$ is obtained from (\SGARefLink{sga7.sga7i.ei.n2-2-3}{2.2.3}) by passing to $P$-invariants:

\begin{equation}\tag{\SGARefTarget{sga7.sga7i.ei.n2-7-3}{2.7.3}}
H^p(X_s,\psi^q_t)\Longrightarrow H^{p+q}(X_t,\Lambda)
=H^{p+q}(X_{\bar\eta},\Lambda)^P.
\end{equation}

% END INCLUDED SGA7I_English_current_buildable_source_20260823/components/03_expose_I_section_2.tex

% BEGIN INCLUDED SGA7I_English_current_buildable_source_20260823/components/04_expose_I_section_3.tex
% SOURCE AUTHORITY:
% expose_I_body.tex lines 368--457; scan indices 21--24; source folios 10--13.
% Hard stop before source line 459, section 4.

3. \underline{Geometric proof of the monodromy theorem}

Let $\Lambda$ be as in \SGARefLink{sga7.sga7i.ei.n2-2}{2.2}.

\underline{Conjecture} \SGARefTarget{sga7.sga7i.ei.n3-1}{3.1} (\underline{purity}). \underline{Let $A$ be an
excellent regular strictly local ring} (\SGARefLink{sga7.sga7i.ei.n0-0-2}{0.0.2}) \underline{and let $D$ be a
regular divisor in $\Spec(A)$ defined by a parameter $x$. Put
$U=\Spec(A)-D=\Spec(A[1/x])$. Then}
% The French TeX/print has the unrelated cross-reference (0.2) and A[1/X].
% The definition of "strictly local" is (0.0.2), and the named parameter is x.

\begin{equation}\tag{\SGARefTarget{sga7.sga7i.ei.n3-1-1}{3.1.1}}
H^{q}(U,\Lambda)=
\begin{cases}
\Lambda & \text{if }q=0,\\[4pt]
\Lambda(-1) & \text{if }q=1\quad\text{(notation (\SGARefLink{sga7.sga7i.ei.n0-0-6}{0.0.6})),}\\[4pt]
0 & \text{if }q>1.
\end{cases}
\end{equation}

Here are some cases in which this conjecture is known (SGA 4 XIX): for
$q=0$ or $1$; when $A$ has characteristic $0$; when $A$ has equal
characteristic $p$ and resolution of singularities is available in
characteristic $p$ in dimensions $\leqslant\dim(A)$; or when
$(\Spec(A),D)$ is a smooth pair (the relative purity theorem, SGA 4 XVI
3.7).

When purity is available and
$D=\sum_{i\in B}D_i$ is a normal-crossings divisor in $\Spec(A)$ defined
by part of a regular system of parameters, the cohomology of
$U=\Spec(A)-D$ is given by

%%%%%%%%%%  scan idx 22  |  folio 11  %%%%%%%%%%

\begin{equation}\tag{\SGARefTarget{sga7.sga7i.ei.n3-1-2}{3.1.2}}
\begin{aligned}
H^0(U,\Lambda)&=\Lambda,\\[4pt]
H^1(U,\Lambda)&=\Lambda(-1)^B,\\[4pt]
H^q(U,\Lambda)&=\bigwedge\nolimits^q H^1(U,\Lambda)
  \qquad\text{(via the cup product).}
\end{aligned}
\end{equation}
These formulas are identical to those which, over $\mathbf C$, give the
cohomology of $D^{*B}\times D^k$.

\SGARefTarget{sga7.sga7i.ei.n3-2}{3.2}. Let $S$ be a strictly local trait, and retain the notation of
(\SGARefLink{sga7.sga7i.ei.n0-0-1}{0.0.1}) and (\SGARefLink{sga7.sga7i.ei.n0-0-3}{0.0.3}). Let $f:X\longrightarrow S$ be a morphism of finite
type. Suppose that $X$ is regular, that $X_\eta$ is smooth, and that
$X_s$ is a normal-crossings divisor in $X$. We shall suppose that this
divisor is, globally and not merely locally for the \'{e}tale topology, a
sum of regular divisors:
\[
(X_s)_{\mathrm{red}}=\sum_{i\in B}D_i
\]
(the general case can be recovered by localization). Denote their
multiplicities by $m_i$, so that
$X_s=\sum_{i\in B}m_iD_i$.

Let $x$ be a point of $X_s$, let $\bar x$ be a geometric point of $X$
lying over $x$ (for example, $x$ may be closed and equal to $\bar x$),
put $C=\{i\mid x\in D_i\}$, and let $K$ be the complex concentrated in
degrees $0$ and $-1$
\begin{equation*}
K:\mathbf Z^C\xrightarrow{\ d\ }\mathbf Z,
\qquad d((n_j))=\sum_jm_jn_j.
\end{equation*}

\underline{Theorem} \SGARefTarget{sga7.sga7i.ei.n3-3}{3.3} (\underline{uses purity}). \underline{The groups
of tame vanishing cycles $\psi^i_{t\bar x}$ are as follows:}

(i) \underline{There is a canonical isomorphism of graded
$\Lambda$-algebras}
\begin{equation*}
\psi^*_{t\bar x}\simeq
\bigwedge\nolimits^*\bigl(H_1(K)\otimes\Lambda(-1)\bigr)
\otimes_\Lambda\psi^0_{t\bar x}.
\end{equation*}

(ii) $\psi^0_{t\bar x}$ \underline{is the} $\Lambda$-\underline{algebra}
$\Lambda^E$, \underline{where $E$ is a set on which the tame inertia
group} $I_t=I/P=\widehat{\mathbf Z}(1)(k(s))$
\underline{acts transitively, and whose cardinality is the largest
prime-to-$p$ divisor of $\#H_0(K)$.}

Here is the method of proof.

a) Let $X_{(\bar x)}$ be the strict localization of $X$ at $\bar x$,
let $U=X_{(\bar x)}-X_{(\bar x)s}$ (the complement of a
normal-crossings divisor), let $t_j$ be a local equation for $D_j$, put
\[
X_n=X_{(\bar x)}\bigl[(t_j^{1/n})\bigr]
\qquad ((n,p)=1),
\]
which is a regular scheme, and let $U_n$ be the inverse image of $U$ in
$X_n$. By (\SGARefLink{sga7.sga7i.ei.n3-1-2}{3.1.2}),
\[
H^q(U_n,\Lambda)=\bigwedge\nolimits^q\bigl(\Lambda(-1)^C\bigr).
\]
For $m=nd$, the pullback map from $H^q(U_n,\Lambda)$ to
$H^q(U_m,\Lambda)$ is multiplication
% The French TeX/print has m=xd. In this inverse system x is a point while
% m,n,d are the indices; the pullback relation is necessarily m=nd.

%%%%%%%%%%  scan idx 23  |  folio 12  %%%%%%%%%%

by $d^q$. Putting
$\widetilde U=\varprojlim_{(n,p)=1}U_n$, one therefore has
\begin{equation}\tag{\SGARefTarget{sga7.sga7i.ei.n3-3-1}{3.3.1}}
\begin{cases}
H^0(\widetilde U,\Lambda)=\Lambda,\\[6pt]
H^q(\widetilde U,\Lambda)=\displaystyle\varinjlim_n
  H^q(U_n,\Lambda)=0 & \text{if }q\ne0.
\end{cases}
\end{equation}

b) $\widetilde U$ is a pro-covering of $U$ with group
$\widehat{\mathbf Z}(1)(k(s))^C$. By definition,
\[
\psi^*_{t\bar x}=H^*(U_{\eta_t},\Lambda).
\]
% The French working TeX misreads the printed Greek eta in U_{eta_t} as n.
$\widetilde U$ is a pro-covering, with group
$H_1(K)\otimes\widehat{\mathbf Z}(1)(k(s))$, of each connected component
of $U_{\eta_t}$, and these components are permuted transitively by the
tame inertia group. Formula (\SGARefLink{sga7.sga7i.ei.n3-3-1}{3.3.1}) and the Hochschild--Serre spectral
sequence now give the formulas in \SGARefLink{sga7.sga7i.ei.n3-3}{3.3}.

Pictorially, one may say that $U_{\bar\eta}$ is the fibre of a morphism
$\delta$ from the cohomological
$K(\widehat{\mathbf Z}(1)(k(s))^C,1)$ represented by $U$ to the
cohomological $K(\widehat{\mathbf Z}(1)(k(s)),1)$ represented by $\eta$;
the morphism $\delta$ induces
\[
d:(n_j)\longmapsto\sum_jm_jn_j:
\widehat{\mathbf Z}(1)(k(s))^C
\longrightarrow\widehat{\mathbf Z}(1)(k(s)).
\]

\underline{Corollary} \SGARefTarget{sga7.sga7i.ei.n3-4}{3.4} (\underline{uses purity}). \underline{Suppose
that $f$ is proper. Let $N$ be the least common multiple of the
multiplicities $m_j$ $(j\in B)$. Let $q\geqslant0$, and let $q'$ be the
minimum of $q$ and the greatest number of divisors $D_j$ passing through
one point. For $T$ in the tame inertia group $I_t=I/P$, one has}
\begin{equation*}
(T^N-1)^{q'+1}=0
\quad\underline{\text{on}}\quad
H^q(X_{\bar\eta},\Lambda)^P.
\end{equation*}
% The French print has exponent q'. For q=0 that asserts 1=0; the
% filtration argument and the bound restated in 3.8 require q'+1.
This follows at once from \SGARefLink{sga7.sga7i.ei.n3-3}{3.3}, Scholium \SGARefLink{sga7.sga7i.ei.n2-5}{2.5}, and Variant \SGARefLink{sga7.sga7i.ei.n2-7}{2.7}.

The required purity results are available for $H^1$, in characteristic
$0$, or when $S$ is essentially of finite type over a field.

\underline{Theorem} \SGARefTarget{sga7.sga7i.ei.n3-5}{3.5}. \underline{Let $C$ be a smooth projective curve
over the fraction field $k(\eta)$ of a discrete valuation ring $V$. The
inertia group $I$ acts on $H^1(C_{\bar\eta},\mathbf Q_\ell)$ by
quasi-unipotent automorphisms of level $2$: there exists $N$ such that,
for $T\in I$,}
\begin{equation*}
(T^N-1)^2=0
\quad\text{on}\quad H^1(C_{\bar\eta},\mathbf Q_\ell).
\end{equation*}
We reduce to the case in which $V$ is strictly local. The assertion is
that there exists an open subgroup $I_1$ of $I$ such that
$(T-1)^2=0$ on $H^1$ for $T\in I_1$.

%%%%%%%%%%  scan idx 24  |  folio 13  %%%%%%%%%%

We may extend the scalars of $k(\eta)$ to a finite extension, replacing
$V$ by its normalization in that extension. Since the image of the
pro-$p$ group $P$ in an $\ell$-adic Lie group
$\mathrm{GL}(n,\mathbf Q_\ell)$ is finite, we may in particular suppose
that $P$ acts trivially.

Let $C_o$ be a minimal model of $C$ over $V$ (for the existence of
$C_o$, which is based on Abhyankar's theorem, see the discussion in [3],
\S~2, p.~87; one could instead reduce to the case where $V$ is excellent
by passing to its completion and invoking SGA 4 XVI 1.6).

Let $C_1$ be a regular model of $C$, obtained from $C_o$ by blowups, whose
special fibre is a normal-crossings divisor. The arguments of \SGARefLink{sga7.sga7i.ei.n3-3}{3.3} and
\SGARefLink{sga7.sga7i.ei.n3-4}{3.4} apply to $C_1$, $\psi_t^0$, $\psi_t^1$, and
\[
H^1(C_{\bar\eta},\mathbb Z/\ell^n)^P
=H^1(C_{\bar\eta},\mathbb Z/\ell^n),
\]
with $N$ independent of $n$; Theorem \SGARefLink{sga7.sga7i.ei.n3-5}{3.5} follows.
% The print uses a letter o in psi_t^o, but this is cohomological degree zero.

\underline{Remark} \SGARefTarget{sga7.sga7i.ei.n3-6}{3.6}. Let $A$ be an abelian variety over $k(\eta)$.
Since $A$ is a quotient of a Jacobian, it follows again that the action
of $I$ on $H^1(A_{\bar\eta},\mathbb Q_\ell)$, or on $T_\ell(A)$, is
quasi-unipotent of level $2$.
% The print omits the geometric bar from the fibre in this sentence.

\underline{Remark} \SGARefTarget{sga7.sga7i.ei.n3-7}{3.7}. The action of $I$ on
$H^1(X_{\bar\eta},\mathbb Q_\ell)$ is in fact quasi-unipotent of level
$2$ for every scheme $X$ of finite type over $\eta$.
% The print has Q here; the surrounding l-adic statement requires Q_ell.

Theorem \SGARefLink{sga7.sga7i.ei.n3-5}{3.5} (or Remark \SGARefLink{sga7.sga7i.ei.n3-6}{3.6}) is the basis for the proof of the stable
reduction theorem for abelian varieties that will be given in Expos\'{e}
IX, \S~3.

\SGARefTarget{sga7.sga7i.ei.n3-8}{3.8}. It is clear that whenever purity and resolution of singularities
are available (for example, in characteristic $0$), the preceding method
proves the monodromy theorem and gives bounds for the nilpotence
exponent. Thus, in equal characteristic $0$, if $X$ is proper and smooth
over $\eta$ and $T$ belongs to an open subgroup $I_1$ of $I$, one has
\begin{equation*}
(T-1)^{q+1}=0
\quad\text{on}\quad H^q(X_{\bar\eta},\mathbb Z_\ell)
\qquad(T\in I_1\subset I).
\end{equation*}

% END INCLUDED SGA7I_English_current_buildable_source_20260823/components/04_expose_I_section_3.tex

% BEGIN INCLUDED SGA7I_English_current_buildable_source_20260823/components/05_expose_I_section_4.tex
% SOURCE AUTHORITY:
% expose_I_body.tex lines 459--589; scan indices 24--28; source folios 13--17.
% Hard stop before source line 591, section 5.

4. \underline{Vanishing criteria for the sheaves of vanishing cycles}

\SGARefTarget{sga7.sga7i.ei.n4-1}{4.1}. Let $S$ be a strictly local trait and let
$f:X\longrightarrow S$ be a scheme of finite type over $S$. Put
$n=\dim(X_\eta)$.

If $X'$ is the scheme-theoretic closure of $X_\eta$ in $X$, then:

a) On $X_s-X'_s$, one has $\phi^{-1}=\Lambda$, and the other
$\phi^i$ vanish;

%%%%%%%%%%  scan idx 25  |  folio 14  %%%%%%%%%%

b) on $X'_s$, the $\phi^i$ and $\psi^i$ relative to $X/S$ or to
$X'/S$ coincide, and $\phi^{-1}=0$.

Since $X'$ is flat over $S$, this often permits a reduction to the case
in which $X$ is flat over $S$.

\underline{Theorem} \SGARefTarget{sga7.sga7i.ei.n4-2}{4.2}.

(i) \underline{One has $\psi^i=0$ for $i>n$.}

(ii) \underline{More precisely, let $x$ be a point of $X_s$ whose
closure has dimension $k$}
($k=\operatorname{trdeg}_{k(s)}k(x)$),
\underline{and let $\bar x$ be a geometric point lying over $x$. Then}
\begin{equation*}
\psi^i_{\bar x}=0
\quad\underline{\text{for}}\quad i>n-k,
\qquad\text{i.e.}\qquad
d(\psi^i)\leqslant n-i
\quad(\underline{\text{notation of SGA 4 XIV 2.1}}).
\end{equation*}
Indeed,
\[
\psi^i_{\bar x}=H^i(X_{(\bar x)\bar\eta},\Lambda),
\]
and $X_{(\bar x)\bar\eta}$ is an inverse limit of affine schemes of
dimension at most $n$ over $k(\bar\eta)$. Assertion (i) therefore follows
from the affine Lefschetz theorem (SGA 4 XIV 3.2) by passage to the
limit.

To prove (ii), we may suppose that $X$ is flat over $S$ (\SGARefLink{sga7.sga7i.ei.n4-1}{4.1}) and that
there is a quasi-finite flat $S$-morphism
\[
g:X\longrightarrow\mathbb A^k_S
\]
such that $g(x)$ is the generic point of $\mathbb A^k_s$. Let $S'$ be
the strict localization of $\mathbb A^k_S$ at $\bar x$. Then:

a) $S'_{\bar\eta}$ is the spectrum of a field, and
$\Gal(\bar\eta'/s'_{\bar\eta})$ is a $p$-group $Q$;

b) by (i), applied to $g$,
\[
H^i(X_{(\bar x)\bar\eta'},\Lambda)=0
\qquad\text{for }i>n-k;
\]

c) one has
\[
H^i(X_{(\bar x)\bar\eta},\Lambda)
=H^i(X_{(\bar x)\bar\eta'},\Lambda)^Q,
\]
which proves the theorem.

\underline{Corollary} \SGARefTarget{sga7.sga7i.ei.n4-3}{4.3}. \underline{If $f$ is proper and flat and the
non-smooth locus of $f$ in $X_s$ has dimension at most $d$, the
specialization morphism}
\[
H^i(X_s)\longrightarrow H^i(X_{\bar\eta})
\]
\underline{is an isomorphism for $i>n+d+1$ and an epimorphism for
$i=n+d+1$.}

\SGARefTarget{sga7.sga7i.ei.n4-4}{4.4}. To prove, in certain cases, that the $\phi^i$ vanish for small
$i$, we shall use local duality, either directly or through the local
Lefschetz theorems of SGA 2 XIV. This tool is available only in
characteristic $0$, or in equal characteristic $p$ when resolution of
singularities \`a la Hironaka is available in the dimensions under
consideration.

%%%%%%%%%%  scan idx 26  |  folio 15  %%%%%%%%%%

\underline{Theorem} \SGARefTarget{sga7.sga7i.ei.n4-5}{4.5}. \underline{Suppose that $S$ has characteristic
$0$. Let $f:X\longrightarrow S$ be a morphism of finite type and let
$x$ be a closed point of $X_s$. Suppose that:}

(a) \underline{$x$ is an isolated non-smooth point in its fibre;}

(b) \underline{$X$ has relative geometric depth at least $n$ at $x$:
in a neighbourhood of $x$, the scheme $X$ is identified with a
subscheme defined by $k$ equations in a smooth scheme of relative
dimension $N$ at $x$, with $n\leqslant N-k$.}

\underline{Then $\phi^i_x=0$ for $i<n$.}

Let $X_x$ be the strict localization of $X$ at $x$, and put
$V=X_x-\{x\}$. For every trait $S'$ finite over $S$, put similarly
\[
X'_x=X_x\times_S S',\qquad V'=V\times_S S'.
\]
Let $\bar S$ be the normalization of $S$ in $\bar\eta$, and put
\[
\bar X_x=X_x\times_S\bar S,
\qquad \bar V=V\times_S\bar S.
\]
We have:

(1) If $\phi^i=0$ on $(X_x)_s$ (i.e. at every generization of $x$ in
$X_s$) for $i\leqslant m$, then
\begin{equation*}
H^{i-1}_{\{x\}}(\bar X_x)
=\widetilde H^{,i}(\bar V)
\xrightarrow{\ \sim\ }
\widetilde H^{,i}(\bar V_{\bar\eta})
=\phi^i
\qquad\text{for }i\leqslant m.
\end{equation*}
The $\phi^i$ are indeed the cohomology sheaves of $\bar V$ with support
in $V_s$; under the hypotheses of \SGARefLink{sga7.sga7i.ei.n4-5}{4.5}, condition (a) makes (1)
applicable with $m=\infty$ (SGA 4 XV 2.1).

(2) Hypothesis (b) of \SGARefLink{sga7.sga7i.ei.n4-5}{4.5} is stable under base change. By SGA 2 XIV
5.6, it implies
\begin{equation*}
\widetilde H^{,i}(V')
=H^{i-1}_{\{x\}}(X'_x)=0
\qquad\text{for }i<n.
\end{equation*}
We conclude by observing that $\widetilde H^{,i}(\bar V)$ is the
filtered colimit of the $\widetilde H^{,i}(V')$.

Combining \SGARefLink{sga7.sga7i.ei.n4-2}{4.2} and \SGARefLink{sga7.sga7i.ei.n4-5}{4.5} gives:

\underline{Corollary} \SGARefTarget{sga7.sga7i.ei.n4-6}{4.6}. \underline{Let $S$ be a henselian trait of
characteristic $0$, let $X/S$ be a relative complete intersection of
relative dimension $n$, and let $x\in X_s$ be an isolated non-smooth
point in its fibre. Then $\phi^i_x=0$ for $i\ne n$.}

\underline{Corollary} \SGARefTarget{sga7.sga7i.ei.n4-7}{4.7}. \underline{Under the hypotheses of \SGARefLink{sga7.sga7i.ei.n4-6}{4.6},
$\phi^n_x$ is a flat $\Lambda$-module.}

Indeed, if $\Lambda'$ is a finite $\Lambda$-algebra, a general nonsense
(SGA 4 XVII 5.2.11) gives, for the $\phi^i_{\Lambda'}$,

%%%%%%%%%%  scan idx 27  |  folio 16  %%%%%%%%%%

\begin{equation*}
\phi^{n-i}_{\Lambda'}
=\Tor_i^\Lambda(\Lambda',\phi^n_\Lambda),
\end{equation*}
and we apply \SGARefLink{sga7.sga7i.ei.n4-6}{4.6} to $\Lambda'$.

\underline{Variant} \SGARefTarget{sga7.sga7i.ei.n4-8}{4.8}. In \SGARefLink{sga7.sga7i.ei.n4-5}{4.5}, omit condition (a), and let $Z$ be the
non-smooth locus of $f$. One may still conclude that
\begin{equation*}
\phi^i_x=0\qquad\text{for }i<n-\dim Z_s.
\end{equation*}
We proceed by descending induction on $\dim\overline{\{x\}}$. The
induction hypothesis permits us to apply \SGARefLink{sga7.sga7i.ei.n4-5}{4.5} (1) with
$m=n-\dim Z_s$, and the proof is completed as before.

\SGARefTarget{sga7.sga7i.ei.n4-9}{4.9}. It is possible to obtain a vanishing criterion using only
information about the geometric fibres of $f$. For complete
intersections, however, the result is one degree weaker than \SGARefLink{sga7.sga7i.ei.n4-5}{4.5}.

\underline{Proposition} \SGARefTarget{sga7.sga7i.ei.n4-10}{4.10}. \underline{Under the general hypotheses
of \SGARefLink{sga7.sga7i.ei.n4-5}{4.5}, suppose that:}

(a) \underline{$x$ is an isolated non-smooth point in its fibre;}

(b) \underline{for every point $y$ of $X_{\bar\eta}$ whose closure in
$X_{\bar\eta}$ has dimension $d(y)$ and which is a generization of $x$,
the \'{e}tale depth satisfies}
\begin{equation*}
\operatorname{depth}_{\mathrm{\acute{e}t},y}(X_{\bar\eta})
\geqslant n-d(y),
\quad\text{i.e.}\quad
H^i_{\{y\}}(X_{\bar\eta})=0
\quad\underline{\text{for }}i\leqslant n-d(y);
\end{equation*}

(c) \underline{$\operatorname{depth}_{\mathrm{\acute{e}t},x}(X_s)
\geqslant n$.}

\underline{Then $\phi^i_x=0$ for $i<n-1$.}

Retain the notation of \SGARefLink{sga7.sga7i.ei.n4-5}{4.5}. We may apply \SGARefLink{sga7.sga7i.ei.n4-5}{4.5} (1) with $m=\infty$.
By the local Lefschetz theorem (SGA 2 XIV), hypothesis (b) gives
\begin{center}
\begin{tikzcd}[row sep=large, column sep=huge]
H^i(V') \arrow[r, "\sim"] & H^i(V_s) & {(\ i<n\text{-}1\ )} \\
H^i(V') \arrow[r, hook] & H^i(V_s) & {(\ i=n\text{-}1\ )}
\end{tikzcd}
\end{center}
% Direct authority comparison: the first arrow is an isomorphism arrow
% with the tilde above; the second is a rightward hook arrow, not a
% leftward arrowhead. There are no vertical or diagonal arrows.
Condition (c) then gives
\[
H^i(\bar V)=\varinjlim H^i(V')=0
\qquad\text{for }i\leqslant n-1,
\]
which proves \SGARefLink{sga7.sga7i.ei.n4-10}{4.10}.

\SGARefTarget{sga7.sga7i.ei.n4-11}{4.11}. As in \SGARefLink{sga7.sga7i.ei.n4-8}{4.8}, one can give a variant of \SGARefLink{sga7.sga7i.ei.n4-10}{4.10} in which (a) is not
assumed.

%%%%%%%%%%  scan idx 28  |  folio 17  %%%%%%%%%%

% END INCLUDED SGA7I_English_current_buildable_source_20260823/components/05_expose_I_section_4.tex

% BEGIN INCLUDED SGA7I_English_current_buildable_source_20260823/components/06_expose_I_section_5.tex
% SOURCE AUTHORITY:
% expose_I_body.tex lines 591--688; scan indices 28--31; source folios 17--20.
% Hard stop before source line 689, section 6.

5. \underline{Action of monodromy on $\pi_1$}

In this section (Grothendieck's expos\'{e} of 19 November 1968), we
shall use the semistable reduction theorem for curves. Artin and Winters
[2] gave a direct proof of this theorem; it can also be deduced [3] from
the analogous theorem for abelian varieties (IX \S~3, using \SGARefLink{sga7.sga7i.ei.n3-6}{3.6}). We
shall also use Schlessinger's theory (Rim's Expos\'{e} VI).

\SGARefTarget{sga7.sga7i.ei.n5-1}{5.1}. Let $k$ be an algebraically closed field, let $C$ be a projective
curve over $k$, and let $x_1,\ldots,x_n$ be a finite set of closed points
of $C$. The pointed curve
\[
X=(C;x_1,\ldots,x_n)
\]
is called \underline{stable} if:

(a) $C$ is reduced, and its only singularities are double points with
distinct tangents;

(b) the $x_i$ are distinct smooth points of $C$;

(c) if $C_1$ is an irreducible component of $C$, and $E$ is the set of
points of the normalization $C'_1$ of $C_1$ lying above a singular point
of $C$ or above one of the $x_i$, then, if $C'_1$ has genus $0$
(respectively $1$), the set $E$ has at least $3$ (respectively $1$)
elements.
% The print says "a singular point of C_1". Stability requires all points
% above the singular locus of the whole curve C, including intersections
% with other irreducible components.

Condition (c) is equivalent to each of the following:

(c') $X$ has no nontrivial infinitesimal automorphisms.

(c'') If $\omega$ is the invertible dualizing sheaf, the invertible
sheaf
\[
\omega'=\omega\Bigl(\sum_i x_i\Bigr)
\]
is ample.

An $S$-scheme $C$ equipped with sections $x_1,\ldots,x_n$ is a
\underline{stable curve} over $S$ if its geometric fibres are stable
curves. In [3], \S\S~1--2, only the case $n=0$ is considered. The
extension of the results of loc. cit. to the general case is immediate:

(1) If $X/k$ is stable, then
$H^1(C,\omega'^{\otimes n})=0$ for $n\geqslant2$, and
$\omega'^{\otimes n}$ is very ample for $n\geqslant3$ (cf. [3] 1.2).

(2) The formal moduli scheme $M$ of $X$ is smooth, and $C$ remains
singular over a normal-crossings divisor of $M$.

%%%%%%%%%%  scan idx 29  |  folio 18  %%%%%%%%%%

(3) If $X$ and $Y$ are stable over $S$, then $\Isom_S(X,Y)$ is finite
and unramified over $S$ (isomorphisms are required to respect the marked
points).

(4) If $X/k$ is stable, the map
\[
\Aut_k(X)\longrightarrow\Aut_k\Pic^\circ(C)
\]
is injective.

Finally, we have:

\underline{Proposition} \SGARefTarget{sga7.sga7i.ei.n5-1-1}{5.1.1}. \underline{Let
$X_\eta=(C_\eta;x_1,\ldots,x_n)$ be stable over the generic point
$\eta$ of a trait $S$, with $C_\eta$ smooth. There exist a finite
separable extension $\eta'$ of $\eta$ and a stable curve $X'$ over the
normalization $S'$ of $S$ in $\eta'$ such that}
\begin{equation*}
X'\otimes_{S'}\eta'=X_\eta\otimes_\eta\eta'.
\end{equation*}
% The print omits the prime from the second occurrence of eta in the
% phrase defining S'. The extension and the displayed base change make
% eta' mathematically forced.

This follows from the usual semistable reduction theorem, which permits
one to choose $\eta'$ so that the geometric special fibre of the minimal
model $C''$ of $C_{\eta'}$ over $S'$ satisfies (a). One then blows up,
iteratively if necessary, \underline{smooth} points of $C''$ so that
(b) is satisfied. Next one contracts the chains of smooth rational
components of self-intersection $-2$ which contain none of the
$\bar x_i$ in the special fibre; the result is $C'$ (cf. [3], proof of
2.3, p.~88).
% The print says "self intersection two"; the omitted minus sign is the
% standard and mathematically necessary value for the contracted chains.

As in loc. cit., one can check that $C'$ exists as soon as the Jacobian
of $C_\eta$ has semistable reduction. The smoothness hypothesis on
$C_\eta$ is moreover unnecessary.

\SGARefTarget{sga7.sga7i.ei.n5-2}{5.2}. Let $X_\eta$ be a geometrically connected scheme of finite type
over the fraction field of a strictly local trait $S$. For a geometric
point $\xi$ of $X_{\bar\eta}$, there is an exact sequence
\begin{equation}\tag{\SGARefTarget{sga7.sga7i.ei.n5-2-1}{5.2.1}}
\begin{tikzcd}[column sep=large, row sep=tiny]
0 \arrow[r] & \pi_1(X_{\bar\eta},\xi) \arrow[r]
  & \pi_1(X,\xi) \arrow[r]
  & \Gal(\bar\eta/\eta) \arrow[r] & 0 .
\end{tikzcd}
\end{equation}
% The source typescript omits the arrowhead on the first exact-sequence
% map. The English reader supplies the mathematically required arrow.
It induces a homomorphism from $\Gal(\bar\eta/\eta)$ to the group of
outer automorphisms of $\pi_1(X_{\bar\eta})$. Passing to the largest
prime-to-$p$ quotient $\pi_1^{(p)}(X_{\bar\eta})$ of
$\pi_1(X_{\bar\eta})$, one obtains
\begin{equation}\tag{\SGARefTarget{sga7.sga7i.ei.n5-2-2}{5.2.2}}
\Gal(\bar\eta/\eta)
\longrightarrow
\Aut_{\mathrm{ext}}\bigl(\pi_1^{(p)}(X_{\bar\eta})\bigr).
\end{equation}

Suppose that $X_\eta$ has a rational point $x$, and let $\bar x$ be the
geometric point
\[
\bar\eta\longrightarrow\eta\xrightarrow{\ x\ }X_\eta.
\]
For $\xi=\bar x$, the exact sequence (\SGARefLink{sga7.sga7i.ei.n5-2-1}{5.2.1}) splits canonically, giving

%%%%%%%%%%  scan idx 30  |  folio 19  %%%%%%%%%%

\begin{equation}\tag{\SGARefTarget{sga7.sga7i.ei.n5-2-3}{5.2.3}}
[x]:\Gal(\bar\eta/\eta)
\longrightarrow\Aut\bigl(\pi_1(X_{\bar\eta},\bar x)\bigr)
\longrightarrow\Aut\bigl(\pi_1^{(p)}(X_{\bar\eta},\bar x)\bigr).
\end{equation}
% The print refers here to the "exact sequence (5.2.2)". The only exact
% sequence is (5.2.1); (5.2.2) is the induced outer action.

\underline{Theorem} \SGARefTarget{sga7.sga7i.ei.n5-3}{5.3}. \underline{The image of $P$ under (\SGARefLink{sga7.sga7i.ei.n5-2-2}{5.2.2})
(or under (\SGARefLink{sga7.sga7i.ei.n5-2-3}{5.2.3}), which amounts to the same thing) is finite.}

\SGARefTarget{sga7.sga7i.ei.n5-3-1}{5.3.1}. \underline{Reduction to the case in which $X_\eta$ is
geometrically normal and integral.}

Let $X'_\eta$ be the normalization of $X_\eta$, let
$(X'_{i\eta})_i$ be the family of connected components of $X'_\eta$,
put
\[
X''_\eta=X'_\eta\times_{X_\eta}X'_\eta,
\]
and let $(X''_{j\eta})_j$ be the family of connected components of
$X''_\eta$. After a finite extension of $k(\eta)$, we may suppose that
the $X'_{i\eta}$ are geometrically normal and integral and have rational
points $x_i$, and that the $X''_{j\eta}$ are geometrically connected and
have rational points $y_j$.

For $k=1$ or $2$, and whenever $p_k(y_j)\in X'_{i\eta}$, choose a class
of paths $\ell_{kij}$ from $\operatorname{pr}_k(\bar y_j)$ to $\bar x_i$.
It is an element of a torsor under
$\pi_1(X'_{i\bar\eta},\bar x_i)$. Let $\bar\ell_{kij}$ be the
corresponding prime-to-$p$ path (the same path modulo
$\Ker(\pi\longrightarrow\pi^{(p)})$). We may choose
$\bar\ell_{kij}$ to be invariant under $P$: the set of prime-to-$p$
paths is an inverse limit of finite sets of cardinality prime to $p$ on
which the pro-$p$ group $P$ acts.

Consequently, if a finite-index subgroup $P'$ of $P$ acts trivially on
the groups $\pi_1^{(p)}(X'_{i\bar\eta},\bar x_i)$, the Van Kampen
formulas---equivalently, the fact that the normalization map
$X'_{\bar\eta}\longrightarrow X_{\bar\eta}$ is effective descent for
\'{e}tale coverings---show that $P'$ acts trivially on
$\pi_1^{(p)}(X_{\bar\eta},\bar x)$.
% The print carries a spurious component index on the target of the
% normalization map; effective descent concerns X'_{bar eta} -> X_{bar eta}.

\SGARefTarget{sga7.sga7i.ei.n5-3-2}{5.3.2}. If $X_{\bar\eta}$ is normal and $U$ is an open subset, then
$\pi_1(U)$ maps onto $\pi_1(X_{\bar\eta})$; we may therefore suppose
that $X_{\bar\eta}$ is smooth and affine. By cutting with
\underline{generic} hyperplane sections (which replaces $S$ by its
strict localization at the generic point of the special fibre of an
$\mathbb A^n_S$) and applying Bertini, we reduce to the case in which
$X_\eta$ is \underline{smooth, geometrically connected, and of
dimension one}.

After a finite extension of $k(\eta)$, we may suppose that $X_\eta$ is
the complement in a smooth projective curve $\bar X_\eta$ of a finite
set of rational points $x_i$, whose number we are free to increase. By
the semistable reduction theorem (\SGARefLink{sga7.sga7i.ei.n5-1-1}{5.1.1}), we reduce to the case in which
there is a stable curve $X=(C;x_1,\ldots,x_n)$ over $S$ such that
\[
X_\eta=C_\eta-\{x_1,\ldots,x_n\}.
\]
We then apply the following result.

%%%%%%%%%%  scan idx 31  |  folio 20  %%%%%%%%%%

\underline{Theorem} \SGARefTarget{sga7.sga7i.ei.n5-4}{5.4}. \underline{Let $S$ be a strictly local scheme,
let $f:\bar X\longrightarrow S$ be a proper flat morphism whose special
fibre is a curve with ordinary double points as its only non-smooth
points, and let $Y$ be a subscheme of $\bar X$ which is the union of a
finite number of disjoint sections contained in the smooth locus. Let
$x$ be a rational point of $X=\bar X-Y$. Let $U$ be the open subset of
$S$ over which $X$ is smooth, and let $\xi$ be a geometric point of $U$.
Then the image of $\pi_1(U,\xi)$ in
$\Aut(\pi_1^{(p)}(X_\xi,x\xi))$ is abelian and prime to $p$.}

Observe that $\pi_1(U,\xi)$ can act on
$\pi_1^{(p)}(X_\xi,x\xi)$ only because the groups $\pi_1^{(p)}$
satisfy the specialization theorem.

We reduce first to the case in which $S$ is complete, and then to the
universal case in which $S$ is a formal moduli scheme for the special
fibre equipped with its sections. Then $S$ is the spectrum of a formal
power-series ring over a Cohen ring:
\[
S=\Spec\bigl(W[[t_1,\ldots,t_n]]\bigr),
\]
and $U$ is the complement of a normal-crossings divisor with equation
$t_1\cdots t_m=0$. In this universal case, Abhyankar's lemma shows that
$\pi_1(U,\xi)$ is abelian and prime to $p$, which proves the theorem.

\underline{Remark} \SGARefTarget{sga7.sga7i.ei.n5-5}{5.5}. One can prove a result analogous to \SGARefLink{sga7.sga7i.ei.n5-4}{5.4} in
arbitrary relative dimension by taking a generic plane section of
dimension one and applying Bertini.


% END INCLUDED SGA7I_English_current_buildable_source_20260823/components/06_expose_I_section_5.tex

% BEGIN INCLUDED SGA7I_English_current_buildable_source_20260823/components/07_expose_I_section_6.tex
6. \underline{Appendix}: \underline{arithmetic proof of the stable reduction theorem}

(by P. Deligne)

Let $S$ be a trait and let $A_\eta$ be an abelian variety over $k(\eta)$,
of dimension $d\neq 0$. For a local morphism of traits
$S'\longrightarrow S$, denote by $A_{\eta'}$ the abelian variety over
$k(\eta')$ obtained from $A_\eta$ by extension of scalars. Let
$\mathcal A_{S'}$ be the N\'eron model of $A_{\eta'}$ over $S'$, let
$\mathcal A_{S',s'}$ be its special fibre, and let
$\mathcal A^\circ_{S',s'}$ be the identity component of
$\mathcal A_{S',s'}$. The stable reduction theorem is as follows.

%%%%%%%%%%  scan idx 32  |  folio 21  %%%%%%%%%%

\underline{Theorem \SGARefTarget{sga7.sga7i.ei.n6-1}{6.1}. For a suitable $S'$, $A_{\eta'}$ has stable
reduction, i.e. $\mathcal A^\circ_{S',s'}$ is an extension of an abelian
variety by a torus.}

It follows fairly readily from \SGARefLink{sga7.sga7i.ei.n6-1}{6.1} that one may choose $S'$ such that
$k(\eta')$ is a finite separable extension of $k(\eta)$. Suppose first
that the residue field of $S$ is of finite type over the prime field (the
same argument would work if $k(s)$ were radicial over such a field).

Let $\ell$ be a prime different from the residual characteristic exponent
$p$, let $T_\ell(A)$ be the Tate module, let
$V_\ell(A)=T_\ell(A)\otimes\mathbf Q_\ell$, and let
\[
 \rho:\Gal(\bar\eta/\eta)\longrightarrow
 \mathrm{GL}(T_\ell(A))
\]
be the natural representation. After a preliminary extension of scalars,
we are reduced to the case in which the inertia group $I$ acts
unipotently (\SGARefLink{sga7.sga7i.ei.n1-2}{1.2}), and hence through its largest pro-$\ell$ quotient
$\mathbf Z_\ell(1)$ (\SGARefLink{sga7.sga7i.ei.n0-3}{0.3}). Let $T$ be the image in
$\mathrm{GL}(T_\ell(A))$ of a generator of $\mathbf Z_\ell(1)$, and let
$r$ be the largest integer $\geq 0$ such that $(T-1)^r\neq 0$. We have:

\underline{Lemma \SGARefTarget{sga7.sga7i.ei.n6-2}{6.2}. The map}
\[
 (\sigma,v)\in I\times V_\ell(A)
 \longmapsto (\rho(\sigma)-1)^r(v)
\]
\underline{induces the morphism and the isomorphism in the commutative
diagram}
\begin{equation}\tag{\SGARefTarget{sga7.sga7i.ei.n6-2-1}{6.2.1}}
\begin{tikzcd}[column sep=huge, row sep=huge]
V_\ell(A)_I\otimes\mathbf Q_\ell(r)
  \arrow[r,"M"] \arrow[d]
& V_\ell(A)^I \\
\bigl(V_\ell(A)/\Ker(T-1)^r\bigr)\otimes\mathbf Q_\ell(r)
  \arrow[r,"\sim"]
& \operatorname{Im}(T-1)^r \arrow[u]
\end{tikzcd}
\end{equation}
$V_\ell(A)_I$ \underline{and} $V_\ell(A)^I$ \underline{are
$\Gal(\bar s/s)$-modules, and $M$ is Galois-equivariant.}

\SGARefTarget{sga7.sga7i.ei.n6-3}{6.3}. An $\ell$-adic representation
\[
 \tau:\Gal(\bar s/s)\longrightarrow\mathrm{GL}(V)
\]
is said to be of weight $n$ ($n\in\mathbf Z$) if the following condition
is satisfied:

(*) For a suitable model $X$ of $k(s)$ (that is, an irreducible variety
of finite type over the prime field such that $k(s)=k(X)$), $\tau$
factors through $\pi_1(X,\bar s)$ and, for every closed point $x$ of
$X$, the eigenvalues of $\tau(F_x)$ are algebraic numbers all of whose
complex conjugates have absolute value $|k(x)|^{n/2}$ (here $F_x$
denotes the geometric Frobenius $\phi_x^{-1}$).

%%%%%%%%%%  scan idx 33  |  folio 22  %%%%%%%%%%

If $V$ (respectively $W$) has weight $n$ (respectively $m$) and
$n\neq m$, then $\Hom_{\Gal}(V,W)=0$.

\SGARefTarget{sga7.sga7i.ei.n6-4}{6.4}. It follows from the universal property of the N\'eron model that
$V_\ell(A)^I=V_\ell(\mathcal A^\circ_{S,s})$. If $a$ is the dimension
of the largest abelian quotient of $\mathcal A^\circ_{S,s}$ and $m$ is
the dimension of its multiplicative part, then, by Weil,
$V_\ell(A)^I$ is an extension of a representation of dimension $2a$ and
weight $-1$ by a representation of dimension $m$ and weight $-2$.

\SGARefTarget{sga7.sga7i.ei.n6-5}{6.5}. Let $A^*$ be the abelian variety dual to $A$. Recall that
$V_\ell(A^*)$ and $V_\ell(A)$ (and hence $V_\ell(A^*)^I$ and
$V_\ell(A)_I$) are in duality with values in $\mathbf Q_\ell(1)$. If
$a^*$ and $m^*$ are defined for $A^*$ as $a$ and $m$ are for $A$, it
follows from \SGARefLink{sga7.sga7i.ei.n6-4}{6.4} that $V_\ell(A)_I$ is an extension of a representation
of dimension $m^*$ and weight $0$ by a representation of dimension
$2a^*$ and weight $-1$.

\SGARefTarget{sga7.sga7i.ei.n6-6}{6.6}. The Galois module $\mathbf Q_\ell(r)$ has weight $-2r$. Since
$\operatorname{Im}(T-1)^r\neq 0$, it follows from \SGARefLink{sga7.sga7i.ei.n6-4}{6.4}, \SGARefLink{sga7.sga7i.ei.n6-5}{6.5}, and the isomorphism in \SGARefLink{sga7.sga7i.ei.n6-2}{6.2}
that $r\leq 1$. If $r=0$, then
\[
\left\{
\begin{array}{l}
a+m\leq d,\\[1ex]
2a+m=2d,
\end{array}
\right.
\]
so $m=0$, $a=d$, and $A_\eta$ has good reduction. Suppose that $r=1$.
It then follows from \SGARefLink{sga7.sga7i.ei.n6-2}{6.2} that $\operatorname{Im}(T-1)$ has weight $-2$ and that
$V_\ell(A)/V_\ell(A)^I$ has weight $0$; these spaces have the same
dimension $m_1\leq m$.

We then have
\[
 \dim V_\ell(A)^I=2d-m_1=2a+m,
\]
and
\[
 2\dim\mathcal A_{S,s}\geq 2(a+m)\geq
 2a+m+m_1=2d=2\dim\mathcal A_{S,s},
\]
so $A_\eta$ has stable reduction
($\dim\mathcal A_{S,s}=a+m$). At the same time, we have obtained
$m_1=m$; that is,
\[
 \operatorname{Im}(T-1)\subset V_\ell(A)^I\simeq
 V_\ell(\mathcal A^\circ_{S,s})
\]
is the $V_\ell$ of the multiplicative part of
$\mathcal A^\circ_{S,s}$.

\SGARefTarget{sga7.sga7i.ei.n6-7}{6.7}. To treat the general case, we shall use (\SGARefLink{sga7.sga7i.ei.n0-5}{0.5}) (cf. \SGARefLink{sga7.sga7i.ei.n1-3}{1.3}). We reduce
to the case in which:

%%%%%%%%%%  scan idx 34  |  folio 23  %%%%%%%%%%

(a) the conditions of (\SGARefLink{sga7.sga7i.ei.n0-5-2}{0.5.2}) or (\SGARefLink{sga7.sga7i.ei.n0-5-3}{0.5.3}) are satisfied;

(b) $\Gal(\bar\eta/\eta)$ acts trivially on $A_\ell$;

(c) the action of $I$ on $T_\ell(A)$ is unipotent.

Under these hypotheses, we shall prove that $A_\eta$ has stable
reduction. Hypothesis (a) allows us to apply (\SGARefLink{sga7.sga7i.ei.n0-5-1}{0.5.1}) and write
$S=\varprojlim S_i$ as in \SGARefLink{sga7.sga7i.ei.n1-3}{1.3}, whose notation we retain. Let $s_i$ be
the closed point of $S_i$.

For sufficiently large $i$, the identity component $\mathcal A^\circ$
of the N\'eron model $\mathcal A_S$ descends to a smooth group scheme
$\mathcal A_i$ with connected fibres over $S_i$, and
$(\mathcal A_i)_\ell$ is constant over $S_i-D_i$. The group
$\pi_1(S_i-D_i,\bar\eta)$ acts on $T_\ell(A)$ and acts trivially on
$T_\ell(A)/\ell T_\ell(A)\simeq A_\ell$. As in \SGARefLink{sga7.sga7i.ei.n1-3}{1.3}, one sees that
$P_i$ acts trivially. We then have, by (\SGARefLink{sga7.sga7i.ei.n1-3-1}{1.3.1}),
\[
 T_\ell(A)_I=T_\ell(A)_{I_i},\qquad
 T_\ell(\mathcal A_{i,s_i})=T_\ell(\mathcal A^\circ_s)
 =T_\ell(A)^I=T_\ell(A)^{I_i},
\]
and (\SGARefLink{sga7.sga7i.ei.n6-2-1}{6.2.1}) is a diagram of $\Gal(\bar s_i/s_i)$-modules to which the
arguments of \SGARefLink{sga7.sga7i.ei.n6-4}{6.4}--\SGARefLink{sga7.sga7i.ei.n6-6}{6.6} may be applied.

%%%%%%%%%%  scan idx 35  |  folio 24  %%%%%%%%%%

\begin{center}
\underline{BIBLIOGRAPHY}
\end{center}

\begin{enumerate}[label={[\arabic*]},leftmargin=3em]
\item M. Artin, ``Algebraic approximation of structures over complete
local rings,'' \emph{Publ. Math. I.H.E.S.} \underline{36} (1969),
pp. 23--58.

\item M. Artin and G. Winters, ``Degenerate fibres and reduction of
curves.''

\item P. Deligne and D. Mumford, ``The irreducibility of the space of
curves of a given genus,'' \emph{Publ. Math. I.H.E.S.} \underline{36}
(1969), pp. 75--110.

\item J.-P. Serre, \emph{Corps locaux}, Publ. Math. Univ.
Nancago--Hermann, 1962.

\item J.-P. Serre and J. Tate, ``Good reduction of abelian varieties,''
\emph{Ann. of Math.} \underline{88}, no. 3 (1968), pp. 492--517.
\end{enumerate}

% END INCLUDED SGA7I_English_current_buildable_source_20260823/components/07_expose_I_section_6.tex

\clearpage
\phantomsection\addcontentsline{toc}{section}{Exposé II}
% BEGIN INCLUDED SGA7I_English_current_buildable_source_20260823/components/08_expose_II_complete.tex
%%%%%%%%%%  scan idx 36  |  folio 25  %%%%%%%%%%

\begin{center}
\underline{\SGARefAliasTarget{sga7.sga7i.eii.nroot}Expos\'e II}

\vspace{1em}
\underline{\underline{Finiteness properties of the fundamental group}}

\vspace{1em}
\underline{by Mich\`ele Raynaud}
\end{center}

Let $k$ be a separably closed field of characteristic $p\geq 0$, let
$X$ be a connected scheme of finite type over $k$, and let
$\Pi_1^{(p)}(X)$ be the largest ``prime-to-$p$'' quotient of $\Pi_1(X)$.
In this expos\'e we prove that, when $X$ is a surface,
$\Pi_1^{(p)}(X)$ is topologically finitely presented, and that the same
holds for arbitrary $X$ if resolution of singularities is assumed.

For a surface, the proof method, due to J.-P. Murre, consists in reducing
to a fibration over $P^1$ having properties (i), (ii), and (iii) of the
following proposition.

\underline{Proposition \SGARefTarget{sga7.sga7i.eii.n2-1}{2.1}. Let $k$ be an algebraically closed field and
let $S$ be a projective, irreducible, smooth surface over $k$. Then there
exists a blow-up}
\[
 g:S'\longrightarrow S,
\]
\underline{where $S'$ is a regular surface equipped with a fibration}
\[
 f:S'\longrightarrow P^1,
\]
\underline{such that $f$ has a section $\sigma$ and the following
conditions hold:}

(i) \underline{$f$ is smooth at the points of $\sigma(P^1)$.}

(ii) \underline{The fibres of $f$ are geometrically integral curves
having only ordinary singularities.}

(iii) \underline{There exists a nonempty open subset $U$ of $P^1$ such
that the fibres of $f$ over the points of $U$ are smooth.}

The proposition follows from the existence of a projective embedding
$S\longrightarrow P^r$ for which a Lefschetz pencil exists (XVII 2.5).
Indeed, let $D=\{H_t\}_{t\in P^1}$ be such a pencil of hyperplanes. Recall
that it has the following properties (loc. cit. \SGARefLink{sga7.sga7i.eii.n2-2}{2.2}):

%%%%%%%%%%  scan idx 37  |  folio 26  %%%%%%%%%%

(a) The axis $\Delta$ of $D$ meets $S$ transversely; the closed
subscheme $\Delta\cap S$ is therefore reduced and consists of finitely
many closed points $x_1,\ldots,x_n$ of $S$.

(b) There exists a nonempty open subset $U$ of $P^1$ such that, for
every $t\in U$, $H_t$ meets $S$ transversely.

(c) For every $t\in P^1-U$, $H_t$ meets $S$ transversely except at one
point, which is an ordinary quadratic singularity.

Through every point $x$ of $S$ not belonging to $\Delta$ there passes a
unique hyperplane $H_t$. The map assigning $t\in P^1$ to $x$ is a
rational map $a:S\dashrightarrow P^1$ with domain of definition
$S-(\Delta\cap S)$.

Consider the blow-up $g:S'\longrightarrow S$ of the closed subscheme
$S\cap\Delta$. The scheme $S'$ is regular, and the fibre $S'_{x_i}$ over
$x_i$ is the projective bundle defined by the tangent space of $S$ at
$x_i$, hence is isomorphic to $P^1$ (EGA IV 19.4.2, 19.4.4). Moreover,
the rational map $f=ag$ is everywhere defined (XVIII 3.1). Finally, to
each $x_i$ one can associate the section of $f$ that sends
$t\in P^1$ to the tangent vector at $x_i$ of $H_t\cap S$.

Let us verify conditions (i), (ii), and (iii). For each $t\in P^1$, the
curves $S'_t$ and $S_t=H_t\cap S$ are isomorphic. Indeed, the morphism
$S'_t\longrightarrow S_t$ is an isomorphism except possibly at the
points $x_i$, where $S_t$ is regular; moreover, since $S'_t$ is a
complete intersection in $S'$, it is integral, and is therefore
isomorphic to $S_t$.

The asserted properties of the fibres follow from (a)--(c). Finally,
$f$ is smooth at the points of $\sigma$, since $f$ is flat and the
fibres $S'_t$ are geometrically regular at $\sigma(t)$.

%%%%%%%%%%  scan idx 38  |  folio 27  %%%%%%%%%%

\underline{Corollary \SGARefTarget{sga7.sga7i.eii.n2-2}{2.2}. With the notation of \SGARefLink{sga7.sga7i.eii.n2-1}{2.1}, let $Y$ be a
normal-crossings divisor in $S$, and let $Y'=Y\times_S S'$. Then $f$,
$g$, and $U$ may be chosen so that $Y'\times_{P^1}U$ is a divisor with
normal crossings relative to $U$.}

It is known that, if a Lefschetz pencil exists for an embedding
$S\longrightarrow P^r$, the Lefschetz pencils form an open subset of the
variety of lines in $P^r$ (XVII 3.2.1). The same is true of the set of
Lefschetz pencils whose axis does not meet $Y$ and for which there is a
nonempty open subset $U_1$ of $P^1$ such that $Y\cap H_t$ is smooth for
$t\in U_1$. It is then enough to construct $f$ and $g$ from a pencil
satisfying these two conditions and replace $U$ by $U\cap U_1$.

\SGARefTarget{sga7.sga7i.eii.n2-3-0}{2.3.0}. Let $G$ be a profinite group. Recall that $G$ is said to be
topologically finitely generated if there is an integer $n$ such that
$G$ is isomorphic to a quotient of the free pro-group $F(n)$ on $n$
generators. Put $K=\Ker(F(n)\longrightarrow G)$. The group $G$ is said
to be topologically finitely presented if, in addition, there are
finitely many elements $k_1,\ldots,k_r$ of $K$ such that the smallest
closed normal subgroup of $K$ containing $k_1,\ldots,k_r$ is $K$.

\underline{Theorem \SGARefTarget{sga7.sga7i.eii.n2-3-1}{2.3.1}. Let $k$ be a separably closed field of
characteristic $p\geq 0$, and let $X$ be a connected scheme of finite
type over $k$. Suppose that schemes of finite type and dimension at most
$\dim X$ over an algebraic closure of $k$ admit strong resolution of
singularities (SGA 5 I 3.1.5). This condition is satisfied when
$\dim X=2$, by [1], or when $\operatorname{char}k=0$, by [2]. Then the
fundamental group $\Pi_1^{(p)}(X)$ is topologically finitely presented.}

By SGA 1 IX 4.10, we may suppose that $k$ is algebraically closed.

1) \underline{Reduction to the case in which $X$ is a regular surface.}

If $f:X'\longrightarrow X$ is a morphism of finite type, put
$X''=X'\times_X X'$, and let $(X'_a)_{a\in A}$ and
$(X''_b)_{b\in B}$ be the families of connected components of $X'$ and
$X''$, respectively. When $f$ is an effective descent morphism for the
category of finite \'etale coverings, the fundamental group $\Pi_1(X)$
can be computed explicitly

%%%%%%%%%%  scan idx 39  |  folio 28  %%%%%%%%%%

in terms of the fundamental groups $\Pi_1(X'_a)$ and $\Pi_1(X''_b)$.
It follows (SGA 1 IX 5.2, \SGARefLink{sga7.sga7i.eix.n5-3}{5.3}) that if the $\Pi_1(X'_a)$ are
topologically finitely generated, then so is $\Pi_1(X)$; and if the
$\Pi_1(X'_a)$ are topologically finitely presented while the
$\Pi_1(X''_b)$ are topologically finitely generated, then $\Pi_1(X)$ is
topologically finitely presented.

Since $X$ admits resolution of singularities, there is a regular scheme
$X'$ and a proper birational morphism $f:X'\longrightarrow X$. Since
$f$ is an effective descent morphism for finite \'etale coverings, it
suffices to prove the theorem under the additional assumption that $X$
is regular. One first deduces that $\Pi_1(X)$ is topologically finitely
generated; knowing that $\Pi_1(X')$ is topologically finitely presented
and that the $\Pi_1(X''_b)$ are topologically finitely generated then
shows that $\Pi_1(X)$ is topologically finitely presented.

Now let $(U_i)_{i\in I}$ be a finite affine open covering of $X$, and
let $f:\bigsqcup U_i\longrightarrow X$ be the canonical morphism. Since
$f$ is an effective descent morphism for finite \'etale coverings, it is
enough to prove the theorem for the $U_i$. We are thus reduced to the
case in which $X$ is affine and smooth.

Assume henceforth that $X$ is affine and smooth over $k$, with
$\dim X>2$. The Lefschetz theorem for quasi-projective schemes
([3] V 7.4) applies: if $Y$ is a sufficiently general hyperplane section
of $X$, then
\[
 \Pi_1(Y)\simeq\Pi_1(X).
\]
It therefore suffices to prove the theorem for $Y$, reducing us to the
case $\dim X=2$.

2) \underline{The case in which $X$ is a smooth surface over $k$.}

By the strong resolution theorem for surfaces, there exist a smooth,
integral, projective surface $S$ and an open immersion
$i:X\longrightarrow S$ such that $Y=S-X$ is a normal-crossings divisor.
We apply Corollary \SGARefLink{sga7.sga7i.eii.n2-2}{2.2} and retain its notation.

%%%%%%%%%%  scan idx 40  |  folio 29  %%%%%%%%%%

We show that it is enough to prove that the fundamental group of
$X'=X\times_S S'$ is topologically finitely presented. Since $X$ and
$X'$ are regular and
\[
 \operatorname{codim}\!\left(\bigcup_{1\leq i\leq n}\{x_i\},X\right)
 \geq 2,
\]
every \'etale covering of $X'$ induces on
$X-\bigcup_{1\leq i\leq n}\{x_i\}$ an \'etale covering which extends
uniquely to an \'etale covering of $X$ (SGA 2 1.11). This gives an
isomorphism
\[
 \Pi_1(X')\simeq\Pi_1(X).
\]

Consider the cartesian diagram
\[
\begin{tikzcd}[column sep=large,row sep=large]
X' \arrow[d,"h"'] & X'_U \arrow[l] \arrow[d,"h_U"] \\
P^1 & U \arrow[l]
\end{tikzcd}
\]
where $h=f|_{X'}$. Let $L$ be the set of primes different from $p$, and
let $\bar\eta$ be a geometric point over the generic point $\eta$ of
$P^1$. We verify the hypotheses of SGA 1 XIII 4.8. The morphism $h$ has
a section $\sigma$, because $f$ does. To prove condition (a) of
SGA 1 XIII 4.7, note first that, since the fibres of $h$ are
geometrically integral, $h$ is $0$-acyclic and locally $0$-acyclic
(SGA 4 XV 4.1, 1.16). It remains to show that, for every locally constant
sheaf of sets $F$ on $X'$, $(F,h)$ is cohomologically proper in
dimensions $\leq 0$. If $\xi$ is a closed point of $P^1$ and $Z$ denotes
the integral closure of $\Spec\mathcal O_{P^1,\xi}$ in $\bar\eta$, then
$X'_Z=X'\times_{P^1}Z$ satisfies condition $(S_2)$ at its closed points
and is regular at every other point, hence is normal. Consequently, an
\'etale covering of $X'_Z$ whose restriction to $X'_{\bar\eta}$ is
connected is necessarily connected. This proves the assertion, and
hence condition (a).

Since $X'_U$ is the complement in $S'_U$ of a divisor with normal
crossings relative to $U$, $h_U$ is locally $1$-aspherical for $L$
(SGA 4 XV 2.1); and, for every finite constant sheaf of $L$-groups $F$
on $X'_U$, $(F,h_U)$ is cohomologically proper in dimensions $\leq 1$
(SGA 1 XIII 2.4), which gives condition (b). Since $P^1$ is normal and
$T=P^1-U$ has codimension $1$, we have

%%%%%%%%%%  scan idx 41  |  folio 30  %%%%%%%%%%

\[
 \operatorname{prof\,\acute et}_T(P^1)\geq 2,
\]
which gives condition (c). The condition
\[
 \operatorname{prof\,hop}_x(X')\geq 3
\]
holds at every closed point of $X'$ by the purity theorem, which gives
condition (e). Finally, for every $t\in P^1$, $h$ is smooth at
$\sigma(t)$, and hence is locally $1$-aspherical for $L$ at $\sigma(t)$,
which gives condition (d).

By SGA 1 XIII 4.7, if $a$ is a geometric point of
$X'_{\bar\eta}$, there is an exact sequence
\[
 1\longrightarrow\Pi_1^{(p)}(X'_{\bar\eta},a)
 \longrightarrow\Pi'_1(X'_U,a)
 \longrightarrow\Pi_1(U,a)\longrightarrow 1,
\]
and the choice of the section $\sigma$ defines a splitting of this exact
sequence. Let $K$ be the image of $\Pi_1(U,a)$ in
$\Aut(\Pi_1^{(p)}(X'_{\bar\eta},a))$. The group of coinvariants under
$K$,
\[
 \Pi_1^{(p)}(X'_{\bar\eta},a)_K,
\]
is the quotient of $\Pi_1^{(p)}(X'_{\bar\eta},a)$ by the closed normal
subgroup generated by the elements $x^{-1}q(x)$, where
$x\in\Pi_1^{(p)}(X'_{\bar\eta},a)$ and $q\in K$. By loc. cit. (4.7.4),
there is an isomorphism
\[
 \Pi_1^{(p)}(X',a)\simeq
 \Pi_1^{(p)}(X'_{\bar\eta},a)_K.
\]

Write $P^1-U=\coprod_{1\leq i\leq n}\{t_i\}$. For each $i$, let
$\bar T_i$ be the strict localization of $P^1$ at a geometric point
$\bar t_i$ over $t_i$, and let $s_i$ be the generic point of $\bar T_i$.
An \'etale covering of $U$ whose restriction to every $s_i$ is trivial
extends to an \'etale covering of $P^1$, and is therefore trivial. For
each $i$, let $\overline{k(s_i)}$ be an algebraic closure of $k(s_i)$.
The preceding assertion says that $\Pi_1(U,a)$ is the smallest closed
normal subgroup containing the images of the Galois groups
$\Gal(\overline{k(s_i)}/k(s_i))$ in $\Pi_1(U,a)$. By I \SGARefLink{sga7.sga7i.ei.n5-3}{5.3}, the image
$K_i$ of $\Gal(\overline{k(s_i)}/k(s_i))$ in
$\Aut(\Pi_1^{(p)}(X'_{\bar\eta},a))$ is finite. Since $K$ is the
smallest closed normal subgroup containing all the $K_i$, and since
$\Pi_1^{(p)}(X'_{\bar\eta},a)$ is topologically finitely presented
(SGA 1 XIII 3.2), this proves that
$\Pi_1^{(p)}(X'_{\bar\eta},a)_K$, and hence also the isomorphic group
$\Pi_1^{(p)}(X)$, is topologically finitely presented.

%%%%%%%%%%  scan idx 42  |  folio 31  %%%%%%%%%%

\underline{Remark \SGARefTarget{sga7.sga7i.eii.n2-3-2}{2.3.2}.} Resolution of singularities is not needed to
prove Theorem \SGARefLink{sga7.sga7i.eii.n2-3-1}{2.3.1} when $X$ is the open complement of a
normal-crossings divisor in a projective scheme smooth over $k$.

\begin{center}\rule{1.5cm}{0.4pt}\end{center}

\begin{center}\underline{BIBLIOGRAPHY}\end{center}

\begin{enumerate}[label={[\arabic*]},leftmargin=3em]
\item S. Abhyankar, ``Local uniformization on algebraic surfaces over
ground fields of characteristic $p\neq0$,'' \emph{Ann. of Math.},
vol. 63, no. 3, 1956, pp. 491--526.

\item H. Hironaka, ``Resolution of singularities of an algebraic variety
over a field of characteristic zero,'' \emph{Ann. of Math.}, vol. 79,
1964.

\item M. Raynaud, \emph{Th\'eor\`eme de Lefschetz en cohomologie des
faisceaux coh\'erents et en cohomologie \'etale}, thesis, forthcoming.
\end{enumerate}

% END INCLUDED SGA7I_English_current_buildable_source_20260823/components/08_expose_II_complete.tex

\clearpage
\phantomsection\addcontentsline{toc}{section}{Exposé VI}
% BEGIN INCLUDED SGA7I_English_current_buildable_source_20260823/components/09_expose_VI_opening_through_1_1.tex
%%%%%%%%%%  scan idx 43  |  folio 32  %%%%%%%%%%

\begin{center}
\underline{\SGARefAliasTarget{sga7.sga7i.evi.nroot}Expos\'e VI}

\vspace{1em}
\underline{\underline{FORMAL DEFORMATION THEORY}}

\vspace{1em}
\underline{by D. S. Rim}
\end{center}

1. \underline{Formal existence theorem}

We recall the notion of a cofibered category (S.G.A. 1 IV). Let
$\underline C$ be a fixed category. A cofibered category $\underline A$
over $\underline C$ is, by definition, a category $\underline A$
together with a functor $P:\underline A\longrightarrow\underline C$
such that:

\underline{Ax. 1.} For every map $f$ in $\underline C$ and every object
$a$ in $\underline A$ for which $P(a)$ is the source of $f$, there is a
cocartesian $f$-morphism $\alpha:a\longrightarrow b$. Thus, for every
$f$-morphism $\alpha':a\longrightarrow b'$, there is a unique
$T$-morphism $\tau:b\longrightarrow b'$ in $\underline A$ such that
$\alpha'=\tau\alpha$, where $T$ is the target of $f$.

\underline{Ax. 2.} A composite of cocartesian morphisms in
$\underline A$ is again cocartesian.

The following properties of a cofibered category
$P:\underline A\longrightarrow\underline C$ follow immediately from
the definitions.

(Cancellation) Let
$a\overset{\alpha}{\underset{\beta}{\rightrightarrows}}a'$ be two
cocartesian maps such that $P(\alpha)=P(\beta)$. If there is a
cocartesian map $b\xrightarrow{\gamma}a$ such that
$\alpha\gamma=\beta\gamma$, then $\alpha=\beta$.

(Factorization) Given cocartesian maps
$a\xrightarrow{\alpha}a'$ and $a\xrightarrow{\beta}a''$, there is a
cocartesian map $a'\xrightarrow{\gamma}a''$ such that
$\gamma\alpha=\beta$ if and only if there is a map
$f:P(a')\longrightarrow P(a'')$ in $\underline C$ such that
$fP(\alpha)=P(\beta)$. Moreover, $\gamma$ is unique.

Let $\underline A$ be a cofibered category over $\underline C$. For each
object $S$ of $\underline C$, denote by $\underline A(S)$ the
subcategory of $\underline A$ whose objects are the objects $a$ for
which $P(a)=S$, and in which
\[
 \Hom_{\underline A(S)}(a,a')=
 \{\alpha\in\Hom_{\underline A}(a,a')\mid P(\alpha)=\id_S\}.
\]
A cofibered groupoid over $\underline C$ is a cofibered category
$\underline A$ over $\underline C$ such that $\underline A(S)$

%%%%%%%%%%  scan idx 44  |  folio 33  %%%%%%%%%%

is a groupoid for every object $S$ of $\underline C$. Recall that a
groupoid is a category in which every morphism is an isomorphism. Thus a
cofibered groupoid over $\underline C$ is a cofibered category
$P:\underline A\longrightarrow\underline C$ in which every map of
$\underline A$ is cocartesian. Furthermore, for every
$\alpha\in\operatorname{Fl}(\underline A)$, $\alpha$ is an isomorphism
in $\underline A$ if and only if $P(\alpha)$ is an isomorphism.

Throughout the rest of this section, fix a complete noetherian local ring
$\Lambda$ with residue field $k$, and let $\underline C_\Lambda$ be the
category of artinian local $\Lambda$-algebras with residue field $k$,
whose maps are local homomorphisms over $\Lambda$. From \SGARefLink{sga7.sga7i.evi.n1-2}{1.2} onward, a
cofibered groupoid $\underline A$ over $\underline C_\Lambda$ is always
assumed to satisfy $\underline A(k)=\{e\}$, where $\{e\}$ denotes the
category in which $\Hom_{\{e\}}(a,b)$ consists of exactly one element
for every two objects $a,b$ of $\{e\}$. A map
$a'\xrightarrow{\alpha}a$ in $\underline A$ will simply be called a
$\underline C_\Lambda$-map if $P(a')=P(a)$ and
$P(\alpha)=1_{P(a)}$. By abuse of language, a map
$a'\xrightarrow{\alpha}a$ in $\underline A$ will be called surjective
if $P(\alpha):P(a')\longrightarrow P(a)$ is surjective in
$\underline C_\Lambda$.

\underline{Ex. \SGARefTarget{sga7.sga7i.evi.n1-1.item-a}{\SGARefAliasTarget{sga7.sga7i.evi.n1-1}1.1}(a).} Let
$F:\underline C\longrightarrow(\text{categories})$ be a functor.
Define the category $\underline F$ as follows. An object of
$\underline F$ is a pair $(S,a)$ with
$S\in\mathrm{ob}(\underline C)$ and $a\in\mathrm{ob}F(S)$; a map
$(S,a)\longrightarrow(S',a')$ is a pair $(f,u')$, where
$f:S\longrightarrow S'$ is a map in $\underline C$ and
$u':F(f)a\longrightarrow a'$ is a map in $F(S')$. The composite of
\[
 (S,a)\xrightarrow{(f,u')}(S',a')
 \xrightarrow{(g,u'')}(S'',a'')
\]
is
\[
 (gf,\,u''\mathbin{\cdot}F(g)u'):(S,a)\longrightarrow(S'',a'').
\]
With respect to the functor
$P:\underline F\longrightarrow\underline C$ given by
$(S,a)\longmapsto S$, $\underline F$ becomes a cofibered category over
$\underline C$. If $F(S)$ is a groupoid for every
$S\in\mathrm{ob}(\underline C)$, then $\underline F$ is a cofibered
groupoid over $\underline C$. In particular, a functor
$F:\underline C\longrightarrow(\text{groups})$ yields a cofibered
group, and a functor $F:\underline C\longrightarrow(\mathrm{Ens})$
yields a cofibered discrete groupoid. A natural transformation
$\varphi:F\longrightarrow G$, where $F,G:\underline C\longrightarrow
(\text{categories})$, induces a cocartesian functor
$\varphi:\underline F\longrightarrow\underline G$ over $\underline C$.

\underline{Ex. \SGARefTarget{sga7.sga7i.evi.n1-1.item-b}{1.1}(b).} Let $X$ and $Y$ be schemes over $\Lambda$
together with a fixed

%%%%%%%%%%  scan idx 45  |  folio 34  %%%%%%%%%%

isomorphism
\[
 X\otimes_\Lambda k\xleftrightarrow{\ \tau\ }Y\otimes_\Lambda k.
\]
Then the functor
$\Isom_{X,Y}:\underline C_\Lambda\longrightarrow(\text{groupoids})$
given by
\[
 \Isom_{X,Y}(S)=
 \left\{X\otimes_\Lambda S\xleftrightarrow{\ \sigma\ }
 Y\otimes_\Lambda S\;\middle|\;
 \sigma\otimes_\Lambda k=\tau\right\}
\]
yields a cofibered groupoid
$P:\underline{\Isom}_{X,Y}\longrightarrow\underline C_\Lambda$.

\underline{Ex. \SGARefTarget{sga7.sga7i.evi.n1-1.item-c}{1.1}(c).} Let $X_0$ be a fixed scheme over $k$, and let
$\underline M_{X_0}$ be the category of deformations of $X_0$ over
$\underline C_\Lambda$ (see \S4). Thus an object of
$\underline M_{X_0}$ is a pair $(R,X)$, where
$R\in\underline C_\Lambda$ and $X$ is a deformation of $X_0$ to $R$.
Then the functor
$P:\underline M_{X_0}\longrightarrow\underline C_\Lambda$ given by
$(R,X)\longmapsto R$ defines a cofibered groupoid over
$\underline C_\Lambda$.

\underline{Ex. \SGARefTarget{sga7.sga7i.evi.n1-1.item-d}{1.1}(d).} For any groupoid $X$, denote by $[X]$ the
discrete groupoid consisting of the isomorphism classes of objects of
$X$. If $P:\underline F\longrightarrow\underline C_\Lambda$ is a
cofibered groupoid, we obtain a functor
$[\underline F]:\underline C_\Lambda\longrightarrow(\mathrm{Ens})$
given by $[\underline F](S)=[\underline F(S)]$, and hence a cofibered
discrete groupoid which, by abuse of notation, is again denoted by
$[\underline F]$.

% END INCLUDED SGA7I_English_current_buildable_source_20260823/components/09_expose_VI_opening_through_1_1.tex

% BEGIN INCLUDED SGA7I_English_current_buildable_source_20260823/components/10_expose_VI_1_2_through_1_3.tex
\underline{Definition \SGARefTarget{sga7.sga7i.evi.n1-2}{1.2}. A cofibered groupoid $\underline A$ over
$\underline C_\Lambda$ is called semi-homogeneous if the following
properties hold:}

(1) \underline{Every diagram}
\[
\begin{tikzcd}[row sep=large,column sep=large]
a'' \arrow[r,dashed] \arrow[d,dotted] & a' \arrow[d] \\
a_1 \arrow[r] & a
\end{tikzcd}
\]
\underline{of solid arrows in $\underline A$, where
$a'\longrightarrow a$ is surjective, can be completed to a commutative
diagram by including the dotted arrows.}

(2) \underline{Let}
\[
\begin{tikzcd}[column sep=large]
& R \\
R\underset{k}{\times}k[\varepsilon]
 \arrow[ur,"\pi_1"] \arrow[dr,"\pi_2"'] & \\
& k[\varepsilon]
\end{tikzcd}
\]
\underline{be the projection maps, where $k[\varepsilon]$ is the ring
of dual numbers over $k$, and let}
\[
\begin{tikzcd}[column sep=huge]
a' \arrow[dr,"\alpha"] & \\
& a \\
b' \arrow[ur,"\beta"'] &
\end{tikzcd}
\]
\underline{be $\pi_1$-maps in $\underline A$. Then there is a
$\underline C_\Lambda$-map $\gamma:a'\longrightarrow b'$ such that
$\alpha=\beta\gamma$ if and only if
$(\pi_2)_*(a')\simeq(\pi_2)_*(b')$.}

\underline{Remark \SGARefTarget{sga7.sga7i.evi.n1-3.item-a}{\SGARefAliasTarget{sga7.sga7i.evi.n1-3}1.3}(a).} A reformulation of the notion in terms of
$2$-categories will be given later in \SGARefLink{sga7.sga7i.evi.n2-6.item-b}{2.6}(b). If $\underline A$ is
semi-homogeneous, then so is $[\underline A]$. In particular, \SGARefLink{sga7.sga7i.evi.n1-2}{1.2}(2)
implies that the canonical map

%%%%%%%%%%  scan idx 46  |  folio 35  %%%%%%%%%%

\[
 [\underline A(R\times_k k[\varepsilon])]
 \longrightarrow
 [\underline A(R)]\times[\underline A(k[\varepsilon])]
\]
is bijective.

\underline{Remark \SGARefTarget{sga7.sga7i.evi.n1-3.item-b}{1.3}(b).} Consider a commutative diagram
\[
\begin{tikzcd}[row sep=large,column sep=large]
& a' \arrow[dl,"v_1"'] \arrow[dr,"v_2"] & \\
a_1 \arrow[dr,"u_1"'] & & a_2 \arrow[dl,"u_2"] \\
& a &
\end{tikzcd}
\]
in a cofibered category $\underline A$ over $\underline C_\Lambda$.
Put $R=P(a)$, $R'=P(a')$, and $R_i=P(a_i)$ for $i=1,2$. If $a''$ is
the direct image under the map
\[
 (P(v_1),P(v_2)):R'\longrightarrow R_1\times_RR_2,
\]
we obtain a commutative diagram
\[
\begin{tikzcd}[row sep=large,column sep=large]
& a' \arrow[d] \arrow[ddl,"v_1"'] \arrow[ddr,"v_2"] & \\
& a'' \arrow[dl] \arrow[dr] & \\
a_1 \arrow[dr,"u_1"'] & & a_2 \arrow[dl,"u_2"] \\
& a &
\end{tikzcd}
\]
Thus \SGARefLink{sga7.sga7i.evi.n1-2}{1.2}(1) may be replaced by:

(1)$'$ Every diagram
\[
\begin{tikzcd}[cramped,column sep=small,row sep=small]
a_1 \arrow[dr] & & a_2 \arrow[dl] \\
& a &
\end{tikzcd}
\qquad\text{in }\underline A\text{ above}\qquad
\begin{tikzcd}[cramped,column sep=small,row sep=small]
R_1 \arrow[dr] & & R_2 \arrow[dl] \\
& R &
\end{tikzcd}
\quad\text{in }\underline C_\Lambda,
\]
where $R_2\longrightarrow R$ is surjective, can be completed to a
commutative diagram
\[
\begin{tikzcd}[row sep=large,column sep=large]
& a' \arrow[dl] \arrow[dr] & \\
a_1 \arrow[dr] & & a_2 \arrow[dl] \\
& a &
\end{tikzcd}
\qquad\text{above}\qquad
\begin{tikzcd}[row sep=large,column sep=large]
& R_1\times_RR_2 \arrow[dl,"\pi_1"'] \arrow[dr,"\pi_2"] & \\
R_1 \arrow[dr] & & R_2 \arrow[dl] \\
& R &
\end{tikzcd}
\]
where $\pi_i$ ($i=1,2$) are the projections.

\underline{Remark \SGARefTarget{sga7.sga7i.evi.n1-3.item-c}{1.3}(c).} If $\underline A$ is a semi-homogeneous
cofibered groupoid over $\underline C_\Lambda$, then
$[\underline A(k[\varepsilon])]$ is canonically endowed with a vector
space structure over $k$, where $k[\varepsilon]$ is the ring of dual
numbers over $k$. Indeed, let $\underline{\underline V}$ be the category
of

%%%%%%%%%%  scan idx 47  |  folio 36  %%%%%%%%%%

finite-dimensional vector spaces over $k$. Every product-preserving
functor $H:\underline{\underline V}\longrightarrow(\mathrm{Ens})$
factors canonically as
\[
\begin{tikzcd}[row sep=large,column sep=large]
\underline{\underline V} \arrow[rr,"H"]
 \arrow[dr,"\overline H"'] & & (\mathrm{Ens}) \\
& \underline{\underline V} \arrow[ur] &
\end{tikzcd}
\]
where $\underline{\underline V}\longrightarrow(\mathrm{Ens})$ is the
forgetful functor. On the other hand, the functor
\[
 G:\underline{\underline V}\longrightarrow\underline C_\Lambda,
 \qquad W\longmapsto k[W]=k\oplus W,\qquad W^2=0,
\]
takes products in $\underline{\underline V}$ to fibre products over $k$
in $\underline C_\Lambda$. If $\underline A$ is a semi-homogeneous
cofibered groupoid over $\underline C_\Lambda$, then
\[
 [\underline A(k[V]\times_k k[W])]
 \longrightarrow
 [\underline A(k[V])]\times[\underline A(k[W])]
\]
is bijective by \SGARefLink{sga7.sga7i.evi.n1-3.item-a}{1.3}(a). Hence the functor
$V\longmapsto[\underline A(k[V])]$ commutes with products, proving the
assertion.

% END INCLUDED SGA7I_English_current_buildable_source_20260823/components/10_expose_VI_1_2_through_1_3.tex

% BEGIN INCLUDED SGA7I_English_current_buildable_source_20260823/components/11_expose_VI_1_4_through_1_6.tex
%%%%%%%%%%  scan idx 47  |  folio 36 (continued)  %%%%%%%%%%

\underline{Definition \SGARefTarget{sga7.sga7i.evi.n1-4}{1.4}.} (1) \underline{Let $R$ be an object of
$\underline C_\Lambda$. A small extension of $R$ in
$\underline C_\Lambda$ is a surjective map
$R'\xrightarrow{f}R$ in $\underline C_\Lambda$ such that
$(\ker f)^2=0$ and $\ker f\simeq k$. In other words,
$R'\xrightarrow{f}R$ is a small extension of $R$ if}
\[
0\longrightarrow k\xrightarrow{t}R'\xrightarrow{f}R
 \longrightarrow0
\]
\underline{is exact for some $t\in R'$ with $t^2=0$.}

(2) \underline{Let $\underline A$ be a cofibered groupoid over
$\underline C_\Lambda$, and let $a$ be an object of $\underline A$.
A small prolongation of $a$ is an arrow
$a'\xrightarrow{\alpha}a$ in $\underline A$ such that $P(\alpha)$ is
a small extension in $\underline C_\Lambda$. An arrow
$a'\xrightarrow{\alpha}a$ will be called versal for small
prolongations of $a$ if, for every small prolongation
$a'_1\longrightarrow a$, there exists a commutative diagram}
\[
\begin{tikzcd}[row sep=large,column sep=large]
a' \arrow[rr] \arrow[dr,"\alpha"'] & & a'_1 \arrow[dl] \\
& a &
\end{tikzcd}
\]

\underline{Proposition \SGARefTarget{sga7.sga7i.evi.n1-5}{1.5}. Let $\underline A$ be a semi-homogeneous
cofibered groupoid over $\underline C_\Lambda$.}

(1) \underline{Let}
\[
\begin{tikzcd}[column sep=large]
& R \\
R\times_k k[\varepsilon]
 \arrow[ur,"\pi_1"] \arrow[dr,"\pi_2"'] & \\
& k[\varepsilon]
\end{tikzcd}
\]
\underline{be the projection maps in $\underline C_\Lambda$, and let
$a'\xrightarrow{\alpha}a$ be a $\pi_1$-map. Then there exists a map
$a\xrightarrow{\beta}a'$ such that $\alpha\beta=1_a$ if and only if
there exists a map $a\longrightarrow(\pi_2)_*a'$.}

%%%%%%%%%%  scan idx 48  |  folio 37  %%%%%%%%%%

(2) \underline{Let $R'\xrightarrow{f}R$ be a small extension in
$\underline C_\Lambda$, and consider a commutative diagram}
\[
\begin{tikzcd}[row sep=large,column sep=small]
& a'' \arrow[dl,"\alpha'"'] \arrow[dr,"\alpha'_1"]
  & & & & & R'\times_RR' \arrow[dl] \arrow[dr] & \\
a' \arrow[dr,"\alpha"'] & & a' \arrow[dl,"\alpha_1"]
  & & \underline{\text{above}} & R' \arrow[dr] & & R' \arrow[dl] \\
& a & & & & & R &
\end{tikzcd}
\]
\underline{Define the canonical map
$\widetilde f:R'\times_RR'\longrightarrow k[\ker f]$ by}
\[
 \widetilde f(r'_1,r'_2)=r'_1(0)+(r'_1-r'_2).
\]
\underline{If there exists a map
$a'\longrightarrow(\widetilde f)_*(a'')$, then there exists a map
$\gamma:a'\longrightarrow a'$ such that
$\alpha=\alpha_1\gamma$.}

\underline{Proof.} (1) The ``only if'' part is clear. Assume conversely
that there exists a map
$\varphi:a\longrightarrow(\pi_2)_*a'$. Set
\[
 h:R\longrightarrow R\times_k k[\varepsilon],
 \qquad h=1\times P(\varphi).
\]
The composite $\pi_1h$ is the identity. We therefore obtain maps
\[
 a\mathrel{\mathop{\rightleftarrows}^{\gamma}_{\alpha_1}}h_*a,
 \qquad \alpha_1\gamma=1_a,
 \qquad P(\gamma)=h,
 \qquad P(\alpha_1)=\pi_1.
\]
On the other hand, $\pi_2h=P(\varphi)$, so there exists a
$\pi_2$-map
\[
 h_*a\longrightarrow(\pi_2)_*a'.
\]
Thus $(\pi_2)_*(h_*a)=(\pi_2)_*a'$, and \SGARefLink{sga7.sga7i.evi.n1-2}{1.2}(2) gives a
$\underline C_\Lambda$-map $u:h_*a\longrightarrow a'$ such that
$\alpha u=\alpha_1$. Define $\beta=u\gamma$. Then
\[
 \alpha\beta=\alpha u\gamma=\alpha_1\gamma=1_a.
\]

(2) The map
\[
 1\times\widetilde f:R'\times_RR'
 \longrightarrow R'\times_k k[\ker f]
\]
is an isomorphism compatible with the projections onto the first
factor, and $\pi_2(1\times\widetilde f)=\widetilde f$, where
$\pi_2:R'\times_k k[\ker f]\longrightarrow k[\ker f]$ is the second
projection. Since $R'\xrightarrow fR$ is a small extension,
$k[\ker f]\simeq k[\varepsilon]$. Hence, if there exists a map
\[
 a'\longrightarrow(\widetilde f)_*a''
   =(\pi_2)_*\bigl((1\times\widetilde f)_*a''\bigr),
\]
part (1) yields a map $\gamma':a'\longrightarrow a''$ such that
$\alpha'\gamma'=1_{a'}$. Define
$\gamma=\alpha'_1\gamma':a'\longrightarrow a'$. Then
$\alpha_1\gamma=\alpha$.

%%%%%%%%%%  scan idx 49  |  folio 38  %%%%%%%%%%

\underline{Corollary \SGARefTarget{sga7.sga7i.evi.n1-6}{1.6}. Let $\underline A$ be a semi-homogeneous
cofibered groupoid over $\underline C_\Lambda$, and consider a diagram}
\[
\begin{tikzcd}[column sep=large]
a' \arrow[dr,"\alpha"'] & & a_1 \arrow[dl,"\alpha_1"] \\
& a &
\end{tikzcd}
\]
\underline{in $\underline A$, where $\alpha_1$ is a small
prolongation. Assume that $a$ is versal for small prolongations of
$e$. Then there exists a map $\rho:a'\longrightarrow a_1$ such that
$\alpha=\alpha_1\rho$ if and only if there exists a map
$f:P(a')\longrightarrow P(a_1)$ such that
$P(\alpha)=P(\alpha_1)f$ in $\underline C_\Lambda$.}

\underline{Proof.} Assume that $P(\alpha)=P(\alpha_1)f$, where
$f:P(a')\longrightarrow P(a_1)$. Replacing $a'$ by its direct image
under $f$, we may and shall assume that
$P(\alpha)=P(\alpha_1)$. Put $R'=P(a_1)$ and $R=P(a)$. Since
$\underline A$ is semi-homogeneous, we obtain a commutative diagram
\[
\begin{tikzcd}[row sep=large,column sep=small]
& a'' \arrow[dl,"\beta"'] \arrow[dr,"\beta_1"]
  & & & & & R'\times_RR' \arrow[dl] \arrow[dr] & \\
a' \arrow[dr,"\alpha"'] & & a_1 \arrow[dl,"\alpha_1"]
  & & \text{above} & R' \arrow[dr] & & R' \arrow[dl] \\
& a & & & & & R &
\end{tikzcd}
\]
Since $a$ is versal for small prolongations of $e$, so is $a'$;
the assertion therefore follows from \SGARefLink{sga7.sga7i.evi.n1-5}{1.5}(2).


% END INCLUDED SGA7I_English_current_buildable_source_20260823/components/11_expose_VI_1_4_through_1_6.tex

% BEGIN INCLUDED SGA7I_English_current_buildable_source_20260823/components/12_expose_VI_1_7_through_1_8.tex
%%%%%%%%%%  scan idx 49  |  folio 38 (continued)  %%%%%%%%%%

\underline{Definition \SGARefTarget{sga7.sga7i.evi.n1-7}{1.7}. For each object $R$ of
$\underline C_\Lambda$, put}
\[
 \overline\Omega_R=M/(m_\Lambda R+M^2),
\]
\underline{where $m_\Lambda$ and $M$ are the maximal ideals of
$\Lambda$ and $R$, respectively. This is a finite-dimensional vector
space over $k$, and $R\longmapsto\overline\Omega_R$ defines a functor}
\[
 \overline\Omega:\underline C_\Lambda
   \longrightarrow\underline{\underline V},
\]
\underline{where $\underline{\underline V}$ is the category of
finite-dimensional vector spaces over $k$. If $\underline A$ is a
cofibered category over $\underline C_\Lambda$, the composite functor
$\underline A\xrightarrow P\underline C_\Lambda
\longrightarrow\underline{\underline V}$ is again denoted
$\overline\Omega$, by abuse of notation.}

\underline{Proposition \SGARefTarget{sga7.sga7i.evi.n1-8}{1.8}. Let $\underline A$ be a semi-homogeneous
cofibered groupoid over $\underline C_\Lambda$. Given an object $a$ of
$\underline A$ that is versal for small prolongations of $e$, there
exists an arrow $a'\longrightarrow a$ such that:}

%%%%%%%%%%  scan idx 50  |  folio 39  %%%%%%%%%%

(i) \underline{$a'$ is versal for small prolongations of $a$;}

(ii) \underline{$P(a')\longrightarrow P(a)$ is surjective, and its
kernel is annihilated by the maximal ideal of $P(a)$;}

(iii) \underline{$\overline\Omega_{a'}\longrightarrow
\overline\Omega_a$ is an isomorphism.}

\underline{Proof.} Put $R=P(a)$, and consider an exact sequence
\[
 0\longrightarrow J\longrightarrow S\longrightarrow R
 \longrightarrow0,
\]
where $S$ is a formal power-series ring over $\Lambda$ in finitely
many variables and $J\subset M^2$, with $M$ the maximal ideal of $S$.
Let $\Sigma$ be the set of ideals $J_1$ of $S$ such that

(i) $J\supset J_1\supset MJ$, and

(ii) there exists a $\pi_1$-morphism $a_1\longrightarrow a$ in
$\underline A$, where $\pi_1:S/J_1\longrightarrow R$ is the canonical
map.

Property \SGARefLink{sga7.sga7i.evi.n1-2}{1.2}(1) of $\underline A$, together with the fact that $J/MJ$
is a finite-dimensional vector space over $k$, implies that $\Sigma$
has a smallest ideal. Denote it by $J'$, and choose a
$\pi'$-morphism $a'\xrightarrow g a$, where
$\pi':S/J'\longrightarrow R$ is the canonical map. Since
$J'\subset J\subset M^2$, the map
$\overline\Omega_{a'}\longrightarrow\overline\Omega_a$ is an
isomorphism.

We claim that $a'\xrightarrow g a$ is versal for small prolongations
of $a$. Indeed, let $a_1\longrightarrow a$ be a small prolongation,
and put $R_1=P(a_1)$. We obtain an exact commutative diagram
\[
\begin{tikzcd}[column sep=small]
0 \arrow[r] & J/MJ \arrow[r] \arrow[d,"h"]
 & S/MJ \arrow[r] \arrow[d,"h_1"]
 & R \arrow[r] \arrow[d,equal] & 0 \\
0 \arrow[r] & k \arrow[r] & R_1 \arrow[r] & R \arrow[r] & 0.
\end{tikzcd}
\]
If $h=0$, then $\pi_1:R_1\longrightarrow R$ admits a cross-section,
so $h_1$ induces a map $f:S/J'\longrightarrow R_1$. If $h\ne0$,
then $h_1:S/MJ\longrightarrow R_1$ is surjective and hence, by the
definition of $J'$, again induces $f:S/J'\longrightarrow R_1$.
In either case we obtain a commutative diagram
\[
\begin{tikzcd}[column sep=large]
S/J' \arrow[rr,"f"] \arrow[dr,"\pi'"']
 & & R_1 \arrow[dl,"\pi_1"] \\
& R &
\end{tikzcd}
\]
Since $a$ is versal for small prolongations of $e$, Corollary \SGARefLink{sga7.sga7i.evi.n1-6}{1.6}
gives a commutative diagram

%%%%%%%%%%  scan idx 51  |  folio 40  %%%%%%%%%%

\[
\begin{tikzcd}[column sep=large]
a' \arrow[rr,"\rho"] \arrow[dr] & & a_1 \arrow[dl] \\
& a &
\end{tikzcd}
\]
as required.


% END INCLUDED SGA7I_English_current_buildable_source_20260823/components/12_expose_VI_1_7_through_1_8.tex

% BEGIN INCLUDED SGA7I_English_current_buildable_source_20260823/components/13_expose_VI_pro_categories_through_1_10.tex
%%%%%%%%%%  scan idx 51  |  folio 40 (continued)  %%%%%%%%%%

Let $P:\underline A\longrightarrow\underline C_\Lambda$ be a
cofibered category in groupoids. It then gives a cofibered category in
groupoids
\[
 \operatorname{pro}(P):\operatorname{pro}(\underline A)
 \longrightarrow\operatorname{pro}(\underline C_\Lambda)
\]
(for the definition of the category of pro-objects
$\operatorname{pro}(\underline A)$, see A. Grothendieck,
S\'eminaire Bourbaki 195), fitting into the commutative diagram
\[
\begin{tikzcd}[column sep=large,row sep=large]
\underline A \arrow[r] \arrow[d,"P"']
 & \operatorname{pro}(\underline A)
   \arrow[d,"\operatorname{pro}(P)"] \\
\underline C_\Lambda \arrow[r]
 & \operatorname{pro}(\underline C_\Lambda).
\end{tikzcd}
\]

Let $\widehat{\underline C}_\Lambda$ be the full subcategory of
$\operatorname{pro}(\underline C_\Lambda)$ whose objects are systems
$(R_i)_{i\in\mathbb N}$ with $R_i\in\underline C_\Lambda$ such that
each map $\varphi_i:R_i\longrightarrow R_{i-1}$ is surjective and,
for all sufficiently large $i$, induces an isomorphism
\[
 \frac{m_i}{m_\Lambda R_i+m_i^2}
 \xrightarrow{\ \sim\ }
 \frac{m_{i-1}}{m_\Lambda R_{i-1}+m_{i-1}^2},
\]
where $\mathbb N$ denotes the ordered category of natural numbers and
$m_i$ is the maximal ideal of $R_i$.

Let $\widehat{\underline A}$ be the full subcategory of
$\operatorname{pro}(\underline A)$ defined by
\[
 \operatorname{ob}(\widehat{\underline A})=
 \left\{a\in\operatorname{ob}(\operatorname{pro}(\underline A))
 \mathrel{\Big|}
 \operatorname{pro}(P)(a)\in
 \operatorname{ob}(\widehat{\underline C}_\Lambda)\right\}.
\]
We obtain a cofibered category in groupoids
$\widehat P:\widehat{\underline A}\longrightarrow
\widehat{\underline C}_\Lambda$ and a commutative diagram
\[
\begin{tikzcd}[column sep=large,row sep=large]
\underline A \arrow[r,hook] \arrow[d,"P"']
 & \widehat{\underline A} \arrow[d,"\widehat P"] \\
\underline C_\Lambda \arrow[r,hook]
 & \widehat{\underline C}_\Lambda.
\end{tikzcd}
\]
The category $\widehat{\underline C}_\Lambda$ is canonically
equivalent to the category of complete local $\Lambda$-algebras of
formally finite type over $\Lambda$ with fixed residue field $k$; we
henceforth make this identification.

\underline{Definition \SGARefTarget{sga7.sga7i.evi.n1-9}{1.9}. Let $\underline A$ be a cofibered groupoid
over $\underline C_\Lambda$. An object $\xi$ of
$\widehat{\underline A}$ is called quasi-initial if, for every object
$a$ of $\widehat{\underline A}$, there exists a map $\xi\to a$ in
$\widehat{\underline A}$. A quasi-initial object $\xi$ is called
minimal if the morphism
$\xi:h_R\longrightarrow[\widehat{\underline A}]$ induces a bijection}

%%%%%%%%%%  scan idx 52  |  folio 41  %%%%%%%%%%

\[
 \xi(k[\varepsilon]):
 \operatorname{Hom}_{\Lambda\text{-alg}}(R,k[\varepsilon])
 \xrightarrow{\ \sim\ }
 [\underline A(k[\varepsilon])],
 \qquad R=P(\xi).
\]
\underline{An object $\xi$ of $\widehat{\underline A}$ is called a
versal formal object of $\underline A$ if every surjective map
$\alpha:a'\to a$ in $\underline A$ induces a surjective map}
\[
 \operatorname{Hom}_{\widehat{\underline A}}(\xi,a')
 \longrightarrow
 \operatorname{Hom}_{\widehat{\underline A}}(\xi,a).
\]
\underline{A versal formal object $\xi$ is called minimal if
$\xi:h_R\longrightarrow[\widehat{\underline A}]$ induces the above
bijection on $k[\varepsilon]$, where $R=P(\xi)$.}

\underline{Proposition \SGARefTarget{sga7.sga7i.evi.n1-10}{1.10}. Let $\underline A$ be a cofibered
groupoid over $\underline C_\Lambda$.}

(i) \underline{Every (minimally) versal formal object of
$\underline A$ is a quasi-initial (minimal) object of
$\widehat{\underline A}$.}

(ii) \underline{Suppose that $\xi$ is an object of
$\widehat{\underline A}$ for which}
\[
 \widetilde\xi(k[\varepsilon]):
 \operatorname{Hom}_{\Lambda\text{-alg}}(R,k[\varepsilon])
 \longrightarrow[\underline A(k[\varepsilon])]
\]
\underline{is injective, where $P(\xi)=R$. Then every endomorphism of
$\xi$ in $\widehat{\underline A}$ is an automorphism. In particular,
a quasi-initial minimal object of $\widehat{\underline A}$, if it
exists, is unique up to noncanonical isomorphism.}

\underline{Proof.} (i) Let $\xi$ be a versal formal object of
$\underline A$. For an object $a$ of $\widehat{\underline A}$, choose
a presentation by surjective maps in $\underline A$,
\[
 \cdots\longrightarrow a_n\xrightarrow{\alpha_n}a_{n-1}
 \longrightarrow\cdots\longrightarrow
 a_1\xrightarrow{\alpha_1}a_0,
 \qquad a=\varprojlim_n a_n.
\]
Suppose maps $\varphi_i:\xi\longrightarrow a_i$ have been chosen with
$\alpha_i\varphi_i=\varphi_{i-1}$ for all $i<n$. Since $\xi$ is
versal and $\alpha_n$ is surjective, there is a map
$\varphi_n:\xi\longrightarrow a_n$ with
$\alpha_n\varphi_n=\varphi_{n-1}$. Thus
$\varphi=\varprojlim_n\varphi_n$ defines a map $\xi\longrightarrow a$.

(ii) Let $\alpha:\xi\longrightarrow\xi$ be an endomorphism in
$\widehat{\underline A}$. For every
$\lambda\in\operatorname{Hom}_{\Lambda\text{-alg}}
(R,k[\varepsilon])$, we have
\[
 \lambda_*\xi\simeq
 (\lambda\circ P(\alpha))_*\xi.
\]
Injectivity therefore gives
$\lambda=\lambda\circ P(\alpha)$ for every $\lambda$. It follows that
$P(\alpha):R\longrightarrow R$ is surjective. Every surjective
endomorphism of a noetherian ring is an automorphism, so $P(\alpha)$,
and hence $\alpha$, is an automorphism.


% END INCLUDED SGA7I_English_current_buildable_source_20260823/components/13_expose_VI_pro_categories_through_1_10.tex

% BEGIN INCLUDED SGA7I_English_current_buildable_source_20260823/components/14_expose_VI_1_11_through_1_14.tex
%%%%%%%%%%  scan idx 52  |  folio 41 (continued)  %%%%%%%%%%

\underline{Theorem \SGARefTarget{sga7.sga7i.evi.n1-11}{1.11} (Schlessinger). Let $\underline A$ be a
cofibered groupoid over $\underline C_\Lambda$. Then $\underline A$
admits a minimally versal formal object if and only if:}

%%%%%%%%%%  scan idx 53  |  folio 42  %%%%%%%%%%

(i) \underline{$\underline A$ is semi-homogeneous;}

(ii) \underline{$[\underline A(k[\varepsilon])]$ is a
finite-dimensional vector space over $k$ (cf. \SGARefLink{sga7.sga7i.evi.n1-3.item-c}{1.3}(c)).}

\underline{Proof. Necessity.} Let $\xi\in\widehat{\underline A}$ be a
minimally versal formal object of $\underline A$, and consider a
diagram
\[
\begin{tikzcd}[column sep=large,row sep=large]
a_1 \arrow[dr] & & a_2 \arrow[dl,"\alpha"] \\
& a &
\end{tikzcd}
\]
in $\underline A$, where $\alpha$ is surjective. Since $\xi$ is
quasi-initial, there is a map $\xi\longrightarrow a_1$. Since
$\alpha$ is surjective, the composite
$\xi\longrightarrow a_1\longrightarrow a$ lifts to $a_2$; hence we
obtain a commutative diagram
\[
\begin{tikzcd}[column sep=large,row sep=large]
& \xi \arrow[dl] \arrow[dr] & \\
a_1 \arrow[dr] & & a_2 \arrow[dl,"\alpha"] \\
& a &
\end{tikzcd}
\]
which proves condition \SGARefLink{sga7.sga7i.evi.n1-2}{1.2}(1).

We now prove condition \SGARefLink{sga7.sga7i.evi.n1-2}{1.2}(2). Let
\[
\begin{tikzcd}[column sep=large]
& S \\
S\times_k k[\varepsilon]
 \arrow[ur,"\pi_1"] \arrow[dr,"\pi_2"'] & \\
& k[\varepsilon]
\end{tikzcd}
\]
be the projections, and let
\[
\begin{tikzcd}[column sep=large]
a' \arrow[dr,"\alpha"'] & & b' \arrow[dl,"\beta"] \\
& a &
\end{tikzcd}
\]
be $\pi_1$-maps in $\underline A$. By the definition of $\xi$, there
is a commutative diagram
\[
\begin{tikzcd}[column sep=large,row sep=large]
& \xi \arrow[dl,"u"'] \arrow[dr,"v"] & \\
a' \arrow[dr,"\alpha"'] & & b' \arrow[dl,"\beta"] \\
& a &
\end{tikzcd}
\]
and therefore $\pi_1P(u)=\pi_1P(v)$. Suppose now that
$(\pi_2)_*a'\simeq(\pi_2)_*b'$. Minimality of $\xi$ gives
$\pi_2P(u)=\pi_2P(v)$, whence

%%%%%%%%%%  scan idx 54  |  folio 43  %%%%%%%%%%

\[
 P(u)=\pi_1P(u)\times\pi_2P(u)
     =\pi_1P(v)\times\pi_2P(v)=P(v).
\]
Since $\underline A$ is a cofibered groupoid, there is a unique
isomorphism $\sigma:a'\xrightarrow{\sim}b'$ in
$\underline A(S\times_k k[\varepsilon])$ such that $\sigma u=v$.
The equality $\beta\sigma u=\alpha u$ then implies
$\beta\sigma=\alpha$.

\underline{Sufficiency.} Let $c_1,\ldots,c_n$ be objects of
$\underline A(k[\varepsilon])$ whose classes form a $k$-basis of
$[\underline A(k[\varepsilon])]$. Choose
\[
 a_1\in\underline A(
 k[\varepsilon]\times_k\cdots\times_k k[\varepsilon])
\]
such that $(\pi_j)_*a_1\simeq c_j$. Then
\[
 \widetilde a_1(k[\varepsilon]):
 \operatorname{Hom}_{\Lambda\text{-alg}}(R_1,k[\varepsilon])
 \longrightarrow[\underline A(k[\varepsilon])]
\]
is bijective, where $R_1=P(a_1)$. In particular,
$a_1\longrightarrow e$ is versal for small prolongations of $e$.
Choose $a_2\longrightarrow a_1$ versal for small prolongations of
$a_1$ such that
$\overline\Omega_{a_2}\longrightarrow\overline\Omega_{a_1}$ is
bijective. Continue inductively, choosing
$a_{n+1}\longrightarrow a_n$ versal for small prolongations and
inducing a bijection on $\overline\Omega$, and set
\[
 \xi=\varprojlim_n a_n.
\]
Then $\overline\Omega_\xi\longrightarrow\overline\Omega_{a_1}$ is
bijective, so
\[
 \xi(k[\varepsilon]):
 \operatorname{Hom}_{\Lambda\text{-alg}}(R,k[\varepsilon])
 \xrightarrow{\ \sim\ }
 [\underline A(k[\varepsilon])]
\]
is bijective.

It remains to prove versality. Consider a diagram
\[
\begin{tikzcd}
& a' \arrow[d,"\alpha"] \\
\xi \arrow[r,"h"] & a
\end{tikzcd}
\]
where $\alpha$ is surjective. We must lift $h$ to $a'$. By induction
we may assume that $a'$ is a small prolongation of $a$. Since $P(a)$
is artinian, $h$ factors through some $a_q$,
\[
 \xi\longrightarrow a_q\longrightarrow a.
\]
Semi-homogeneity gives a commutative diagram
\[
\begin{tikzcd}
& a'' \arrow[r] \arrow[d,"\alpha'"] & a' \arrow[d,"\alpha"] \\
\xi \arrow[r] & a_q \arrow[r] & a,
\end{tikzcd}
\]
where $\alpha'$ is also a small prolongation. Since
$a_{q+1}\longrightarrow a_q$ is versal for small prolongations of
$a_q$, there is a commutative diagram
\[
\begin{tikzcd}
& & a'' \arrow[d,"\alpha'"] \\
\xi \arrow[r] & a_{q+1} \arrow[r] \arrow[ur] & a_q,
\end{tikzcd}
\]
which supplies the required lift.

\underline{Remark \SGARefTarget{sga7.sga7i.evi.n1-12}{1.12}.} Let

%%%%%%%%%%  scan idx 55  |  folio 44  %%%%%%%%%%

\[
 \cdots\longrightarrow a_n\longrightarrow a_{n-1}
 \longrightarrow\cdots\longrightarrow a_2\longrightarrow a_1
 \longrightarrow a_0=e
\]
be an infinite sequence of arrows in $\underline A$ such that
\[
\left\{
\begin{aligned}
&a_i\text{ is versal for small prolongations of }a_{i-1},\\
&P(a_i)\longrightarrow P(a_{i-1})\text{ is surjective},\\
&\overline\Omega_{a_i}\longrightarrow
 \overline\Omega_{a_{i-1}}\text{ is bijective for }i\gg0.
\end{aligned}
\right.
\]
The preceding proof shows that $\varprojlim_i a_i$ is a versal formal
object of $\underline A$.

\underline{Remark \SGARefTarget{sga7.sga7i.evi.n1-13}{1.13}.} Let $\underline A$ be a cofibered groupoid
over $\underline C_\Lambda$. If $\underline A$ satisfies \SGARefLink{sga7.sga7i.evi.n1-11}{1.11}(i) and
(ii), then so does $[\underline A]$; moreover,
$\xi\in\operatorname{ob}(\widehat{\underline A})$ is a minimally
versal formal object of $\underline A$ if and only if it is so for
$[\underline A]$, by 1.10(ii).

\underline{Remark \SGARefTarget{sga7.sga7i.evi.n1-14}{1.14}.} Let $\underline A$ be a cofibered groupoid
over $\underline C_\Lambda$ admitting a minimally versal formal
object $\xi$. Any invariant of $R=P(\xi)$, a complete local
$\Lambda$-algebra of formally finite type over $\Lambda$, is
automatically an invariant of $\underline A$ over
$\underline C_\Lambda$. Thus, for example, we say that $\underline A$
is (formally) smooth over $\underline C_\Lambda$ if $R$ is
(formally) smooth over $\Lambda$, that is, a formal power-series
algebra over $\Lambda$; and we say that $\underline A$ is irreducible over
$\underline C_\Lambda$ if $\operatorname{Spec}(R)$ is irreducible.

In the next section we introduce smoothness for a functor
$\underline A\longrightarrow\underline B$ over
$\underline C_\Lambda$, where $\underline A$ and $\underline B$ are
cofibered categories in groupoids over $\underline C_\Lambda$.
It follows from \SGARefLink{sga7.sga7i.evi.n2-4}{2.4} that $\underline A$ being (formally) smooth over
$\underline C_\Lambda$ is equivalent to the functor
$P:\underline A\longrightarrow\underline C_\Lambda$ being smooth,
so no ambiguity arises. We define
\[
 \dim_{\underline C_\Lambda}\underline A
   =\dim(k\otimes_\Lambda R).
\]

For this, let $\underline C_k$ be the category of artinian local
$k$-algebras with residue field $k$. The functor
$\underline C_k\longrightarrow\underline C_\Lambda$ is fully
faithful, and its left adjoint is
\[
 \underline C_\Lambda\longrightarrow\underline C_k,
 \qquad R\longmapsto k\otimes_\Lambda R.
\]
Let $\underline A|_{\underline C_k}$ be the pullback of
$P:\underline A\longrightarrow\underline C_\Lambda$ along
$\underline C_k\longrightarrow\underline C_\Lambda$. If $\xi$ is a
minimally versal formal object for $\underline A$ over
$\underline C_\Lambda$, then $k\otimes_\Lambda\xi$ is minimally
versal for $\underline A|_{\underline C_k}$ over $\underline C_k$;
the arrow $\xi\longrightarrow k\otimes_\Lambda\xi$ is cocartesian
over $R\longrightarrow k\otimes_\Lambda R$. Consequently,
\[
 \dim_{\underline C_\Lambda}\underline A
 =\dim_{\underline C_k}(\underline A|_{\underline C_k}).
\]

% END INCLUDED SGA7I_English_current_buildable_source_20260823/components/14_expose_VI_1_11_through_1_14.tex

% BEGIN INCLUDED SGA7I_English_current_buildable_source_20260823/components/15_expose_VI_1_15_through_1_17.tex
%%%%%%%%%%  scan idx 56  |  folio 45  %%%%%%%%%%

\textbf{\SGARefTarget{sga7.sga7i.evi.n1-15}{1.15}.} In the remainder of this paragraph, we briefly explain
how to modify the preceding theory to allow residue-field extensions.

Let $\Lambda$ and $K$ be complete noetherian local rings, with $K$ a
$\Lambda$-algebra. Assume that:

(a) the natural map $\Lambda\longrightarrow K$ is a local
homomorphism;

(b) the residue field $K'$ of $K$ is an extension of finite type of
the residue field $k$ of $\Lambda$.

The most interesting case is $K=K'$ with $K$ a finite inseparable
extension of $k$.

Denote by $\widehat C_{\Lambda,K}$ the category whose objects are
complete noetherian local rings $R$ equipped with a
$\Lambda$-algebra structure and a $\Lambda$-epimorphism
$s:R\longrightarrow K$. Let $C_{\Lambda,K}$ be its full subcategory
consisting of pairs $(R,s)$ for which $\ker(s)$ has finite length as
an $R$-module. The category $\widehat C_{\Lambda,K}$ identifies with
a full subcategory of $\operatorname{Pro}(C_{\Lambda,K})$. If $K$ is
artinian, then every object of $C_{\Lambda,K}$ is artinian.

For a finite-dimensional vector space $V$ over $K'$, write $D_V(K)$
for the algebra of dual numbers over $K$ defined by $V$. Its elements
have the form $x+v\varepsilon$, where $x\in K$, $v\in V$, and
$\varepsilon^2=0$, with multiplication
\[
 (x+v\varepsilon)(x'+v'\varepsilon)
 =xx'+(xv'+x'v)\varepsilon.
\]
The map
\[
 s:D_V(K)\longrightarrow K,
 \qquad s(x+v\varepsilon)=x,
\]
makes $D_V(K)$ an object of $C_{\Lambda,K}$. The role previously
played by $k[\varepsilon]$ is now played by $D_{K'}(K)$.

As in \SGARefLink{sga7.sga7i.evi.n1-8}{1.8}, every cofibered groupoid $\underline A$ over
$C_{\Lambda,K}$ extends naturally to a cofibered category
$\widehat{\underline A}$ over $\widehat C_{\Lambda,K}$. A map in
$\underline A$ or $\widehat{\underline A}$ is called
\underline{surjective} when the corresponding algebra homomorphism
in $C_{\Lambda,K}$ or $\widehat C_{\Lambda,K}$ is surjective.

%%%%%%%%%%  scan idx 57  |  folio 46  %%%%%%%%%%

\underline{We assume that $\underline A(K)=\{e\}$.}

\underline{Definition \SGARefTarget{sga7.sga7i.evi.n1-16}{1.16}. A cofibered groupoid $\underline A$ over
$C_{\Lambda,K}$ is called semi-homogeneous if the following
properties hold:}

(I) \underline{Every diagram of solid arrows}
\[
\begin{tikzcd}[row sep=large,column sep=large]
a'' \arrow[r,dashed] \arrow[d,dashed]
 & a' \arrow[d] \\
a_1 \arrow[r] & a
\end{tikzcd}
\]
\underline{in $\underline A$, where $a'\longrightarrow a$ is
surjective, can be completed to a commutative diagram by including the
dotted arrows.}

(II) \underline{For every $R$, the natural map}
\[
 [\underline A(R\times_KD_{K'}(K))]
 \longrightarrow
 [\underline A(R)]\times[\underline A(D_{K'}(K))]
\]
\underline{is bijective.}

If $\underline A$ is semi-homogeneous, then so is $[\underline A]$.

\textbf{\SGARefTarget{sga7.sga7i.evi.n1-17}{1.17}.} First restrict $\underline A$ to the full subcategory
of $C_{\Lambda,K}$ whose objects are the $D_V(K)$. Put
\[
 \Omega=\Omega^1_{K/\Lambda}\otimes_KK',
 \qquad d:K\longrightarrow\Omega.
\]
A map $f:D_V(K)\longrightarrow D_W(K)$ can be written
\[
 f(x+v\varepsilon)
 =x+\bigl(u(v)+D(dx)\bigr)\varepsilon,
\]
with $u\in\operatorname{Hom}_{K'}(V,W)$ and
$D\in\operatorname{Hom}_{K'}(\Omega,W)$. Identifying $f$ with
$(u,D)$, composition is given by
\[
 (u,D)\circ(u',D')=(uu',D+uD').
\]
Condition (II) implies:

(II$'$) the functor from $K'$-vector spaces to sets
\[
 V\longmapsto[\underline A(D_V(K))]
\]
commutes with products.


% END INCLUDED SGA7I_English_current_buildable_source_20260823/components/15_expose_VI_1_15_through_1_17.tex

% BEGIN INCLUDED SGA7I_English_current_buildable_source_20260823/components/16_expose_VI_lemma_1_18.tex
%%%%%%%%%%  scan idx 58  |  folio 47  %%%%%%%%%%

\underline{Lemma \SGARefTarget{sga7.sga7i.evi.n1-18}{1.18}. Assume that condition (II$'$) is satisfied by
$\underline A$. Then:}

(1) $[\underline A(D_{K'}(K))]$ carries a unique vector-space
structure such that, functorially in $V$,
\[
 [\underline A(D_V(K))]
 \simeq
 V\otimes_{K'}[\underline A(D_{K'}(K))].
\]

(2) There is a linear map
\[
 \gamma:\operatorname{Hom}(\Omega,K')
 \longrightarrow [\underline A(D_{K'}(K))]
\]
such that, for every $D\in\operatorname{Hom}(\Omega,K')$, the map
\[
 [\underline A((1,D))]:
 [\underline A(D_{K'}(K))]
 \longrightarrow [\underline A(D_{K'}(K))]
\]
is $x\longmapsto x+\gamma(D)$.

(3) For every
$(u,D):D_V(K)\longrightarrow D_W(K)$, the induced map
\[
 (u,D):[\underline A(D_V(K))]
 \longrightarrow[\underline A(D_W(K))],
\]
that is,
\[
 (u,D):
 V\otimes[\underline A(D_{K'}(K))]
 \longrightarrow
 W\otimes[\underline A(D_{K'}(K))],
\]
is given by
\[
 x\longmapsto u(x)+\gamma(D).
\]

Introduce the notation
\[
\begin{aligned}
 F(V)&=[\underline A(D_V(K))],\\
 A&=F(K'),\\
 F(u,D):F(V)&\longrightarrow F(W)
   &&\text{for the map induced by }
     (u,D):D_V(K)\longrightarrow D_W(K),\\
 F(D):F(V)&\longrightarrow F(V),
   &F(D)&=F(1,D).
\end{aligned}
\]
By (II$'$), and because $K'$ is a $K'$-vector-space object in the
category of $K'$-vector spaces, $F(K')$ is a $K'$-vector-space object
in the category of sets, hence a $K'$-vector space. This proves (1).
We henceforth identify $F(V)$ with $V\otimes A$.

The identity
\[
 (u,D)=(1,D)\circ(u,0)
\]
gives
\[
 F(u,D)=F(D)\circ(u\otimes1_A).
\]
The functoriality of $F$ amounts to the identities
\[
\begin{aligned}
 F(0)&=\operatorname{Id},\\
 F(D)\circ F(D')&=F(D+D'),\\
 (u\otimes1)\circ F(D)
   &=F(uD)\circ(u\otimes1_A).
\end{aligned}
\]

For any $D:\Omega\longrightarrow V$ and any $n\in\mathbb N$, let
\[
 i:V\lhook\joinrel\longrightarrow V\oplus K'^{\,n},
 \qquad
 p:V\oplus K'^{\,n}\longrightarrow K'^{\,n}
\]
be the canonical inclusion and projection. Then $piD=0$, and the
following diagram is commutative:

%%%%%%%%%%  scan idx 59  |  folio 48  %%%%%%%%%%

\[
\begin{tikzcd}[column sep=huge,row sep=large]
 V\otimes A
   \arrow[r]
   \arrow[d,"{F(D)}"]
 & (V\oplus K'^{\,n})\otimes A
   \arrow[r]
   \arrow[d,"{F(i,D)}"]
 & K'^{\,n}\otimes A
   \arrow[d,equal]\\
 V\otimes A
   \arrow[r]
 & (V\oplus K'^{\,n})\otimes A
   \arrow[r]
 & K'^{\,n}\otimes A.
\end{tikzcd}
\]

Let $x:K'^{\,n}\longrightarrow V$ be a map, put
$x_i=x(e_i)$, and let $(a_i)$ be a family of elements of $A$. We
have
\[
 (1_V+x)\circ i\circ D=D.
\]
Consequently, if
\[
 F(i,D)\Bigl(\sum_i e_i a_i\Bigr)
 =\sum_i e_i a_i+a,
 \qquad a\in V\otimes A,
\]
then
\[
 F(D)\Bigl(\sum_i x_i a_i\Bigr)
 =\sum_i x_i a_i+a,
\]
and $a$ is independent of the $x_i$. Hence
\[
\begin{aligned}
 F(D)(y)&=y+G(D),\\
 G(D+D')&=G(D)+G(D'),\\
 uG(D)&=G(uD)
   \qquad\text{for every }u:V\longrightarrow W.
\end{aligned}
\]
Every such $G$ is given, for some $\gamma\in\Omega\otimes A$, by
\[
 G(D)=(D\otimes1_A)(\gamma).
\]
This proves (2) and (3).

% END INCLUDED SGA7I_English_current_buildable_source_20260823/components/16_expose_VI_lemma_1_18.tex

% BEGIN INCLUDED SGA7I_English_current_buildable_source_20260823/components/17_expose_VI_1_19_through_1_20.tex
%%%%%%%%%%  scan idx 59  |  folio 48 (continued)  %%%%%%%%%%

\underline{Definition \SGARefTarget{sga7.sga7i.evi.n1-19}{1.19}.} (1) An object
$\xi\in\widehat{\underline A}$ is called a \underline{versal formal
object} if every surjective map $\alpha:a'\longrightarrow a$ in
$\underline A$ induces a surjective map
\[
 \operatorname{Hom}_{\widehat{\underline A}}(\xi,a')
 \longrightarrow
 \operatorname{Hom}_{\widehat{\underline A}}(\xi,a).
\]

(2) A versal formal object $\xi\in\widehat{\underline A}$, with
image $(R,s)$ in $\widehat C_{\Lambda,K}$, is called
\underline{minimal} if the following condition holds. For any two
maps
\[
 \varphi_1,\varphi_2:R\rightrightarrows D_{K'}(K),
\]
if $\varphi_1(\xi)$ is isomorphic to $\varphi_2(\xi)$, then there is
a derivation $D:K\longrightarrow K'$ such that
\[
 (1,D)(\varphi_1)=\varphi_2,
 \qquad\text{that is,}\qquad
 \varphi_2-\varphi_1=Ds.
\]

\underline{Theorem \SGARefTarget{sga7.sga7i.evi.n1-20}{1.20}. If $\underline A$ is semi-homogeneous and}
\[
 \dim_{K'}[\underline A(D_{K'}(K))]<\infty,
\]
\underline{then $\underline A$ admits a minimally versal formal
object, and any two such objects are isomorphic.}

%%%%%%%%%%  scan idx 60  |  folio 49  %%%%%%%%%%

With the notation of \SGARefLink{sga7.sga7i.evi.n1-18}{1.18}, let $H^*$ be a subspace of
$[\underline A(D_{K'}(K))]$ complementary to the image of $\gamma$,
and let $H$ be the dual of $H^*$. Set
\[
 R_0=D_H(K).
\]
Choose $\xi_0\in\underline A(R_0)$ whose class, under the
identifications
\[
 [\underline A(R_0)]
 \simeq H\otimes[\underline A(D_{K'}(K))]
 \simeq
 \operatorname{Hom}\bigl(H^*,[\underline A(D_{K'}(K))]\bigr),
\]
is the inclusion
$H^*\hookrightarrow[\underline A(D_{K'}(K))]$. Then:

(1) for every $V$, the maps
\[
 \operatorname{Hom}_{C_{\Lambda,K}}
   \bigl(D_H(K),D_V(K)\bigr)
 \longrightarrow[\underline A(D_V(K))]
\]
are surjective;

(2) the minimality condition of \SGARefLink{sga7.sga7i.evi.n1-19}{1.19}(2) holds for
$(R_0,\xi_0)$.

Choose $S\in\widehat C_{\Lambda,K}$ formally smooth over $\Lambda$,
together with an epimorphism $t:S\longrightarrow R_0$, such that
\[
 \Omega_{R_0/\Lambda}\otimes K'
 \xrightarrow{\ \sim\ }
 \Omega_{S/\Lambda}\otimes K'.
\]
If $M$ is the maximal ideal of $S$, define inductively, for
$n\geq1$,
\[
 R_n=S/I_n,
 \qquad
 \xi_n\in\operatorname{Ob}(\underline A(R_n)),
\]
by the following conditions:

(a) $I_{n+1}$ is the intersection of the ideals $I$ such that
\[
 I_nM\subset I\subset I_n
\]
and $\xi_n$ extends to an object $\xi\in\underline A(S/I)$;

(b) $\xi_{n+1}$ is an extension of $\xi_n$.

If
\[
 R=\varprojlim_nR_n,
 \qquad
 \xi=\varprojlim_n\xi_n,
\]
then $(R,\xi)$ is a minimally versal formal object, and it is unique
up to isomorphism.


% END INCLUDED SGA7I_English_current_buildable_source_20260823/components/17_expose_VI_1_19_through_1_20.tex

% BEGIN INCLUDED SGA7I_English_current_buildable_source_20260823/components/18_expose_VI_section_2_opening_through_2_2.tex
\section*{\SGARefTarget{sga7.sga7i.evi.n2}{2}. Prorepresentable cofibered groupoids}

Let $\underline A$ be a cofibered groupoid over
$\underline C_\Lambda$ admitting a minimally versal formal object
$\xi$. In general, $\xi$ alone does not determine $\underline A$ up
to $\underline C_\Lambda$-equivalence. One must also know when two
maps
\[
 R\mathrel{\substack{\xrightarrow{\ f\ }\\[-.6ex]
                      \xrightarrow[\ g\ ]{}}}S
\]
give rise to a $\underline C_\Lambda$-isomorphism
$\xi(f)\longrightarrow\xi(g)$, where $R=P(\xi)$. As an application
of \SGARefLink{sga7.sga7i.evi.n1-11}{1.11} and \SGARefLink{sga7.sga7i.evi.n1-12}{1.12}, we now consider the problem of recovering
$\underline A$, up to $\underline C_\Lambda$-equivalence. We obtain
a satisfactory answer for a broad class of cofibered groupoids.

%%%%%%%%%%  scan idx 61  |  folio 50  %%%%%%%%%%

Many cofibered groupoids that arise in practice satisfy a property
slightly stronger than semi-homogeneity; it will be defined
explicitly in \SGARefLink{sga7.sga7i.evi.n2-5}{2.5}. To investigate them, it is convenient to recast
the preceding discussion in terms of ``smooth functors.'' All the
notation and terminology of the preceding section remain in force
throughout this section.

\underline{Definition \SGARefTarget{sga7.sga7i.evi.n2-1}{2.1}.} Let
$\varphi:\underline A\longrightarrow\underline B$ be a functor of
cofibered groupoids over $\underline C_\Lambda$. We say that
$\varphi$ is \underline{smooth} if every surjective map
$\beta:b'\longrightarrow\varphi(a)$ in $\underline B$ can be
completed to a commutative diagram
\begin{equation}\tag{*}
\begin{tikzcd}[row sep=large,column sep=huge]
 \varphi(a') \arrow[r] \arrow[dr,"{\varphi(\alpha)}"']
 & b' \arrow[d,"\beta"]\\
 & \varphi(a)
\end{tikzcd}
\end{equation}
for some $\alpha:a'\longrightarrow a$ in $\underline A$. We say
that $\varphi$ is \underline{minimally smooth} if it is smooth and
\[
 \varphi(k[\varepsilon]):
 [\underline A(k[\varepsilon])]
 \longrightarrow
 [\underline B(k[\varepsilon])]
\]
is bijective.

\underline{Remark \SGARefTarget{sga7.sga7i.evi.n2-2.item-a}{\SGARefAliasTarget{sga7.sga7i.evi.n2-2}2.2}(a).} Let
$\varphi:\underline A\longrightarrow\underline B$ be a functor of
cofibered groupoids over $\underline C_\Lambda$. The following
conditions are equivalent:

(i) $\varphi$ is smooth;

(ii) every surjective map
$\beta:b'\longrightarrow\varphi(a)$ in $\underline B$ can be
completed to a diagram $(*)$ in which
$\varphi(a')\longrightarrow b'$ is a
$\underline C_\Lambda$-isomorphism;

(iii) every small prolongation
$\beta:b'\longrightarrow\varphi(a)$ in $\underline B$ can be
completed to a commutative diagram $(*)$.

\underline{Remark \SGARefTarget{sga7.sga7i.evi.n2-2.item-b}{2.2}(b).} For every cofibered groupoid
$\underline A$ over $\underline C_\Lambda$, the natural functor
\[
 \underline A\longrightarrow[\underline A]
\]
is minimally smooth.

\underline{Remark \SGARefTarget{sga7.sga7i.evi.n2-2.item-c}{2.2}(c).} For each
$R\in\widehat{\underline C}_\Lambda$, let
\[
 h_R:\underline C_\Lambda\longrightarrow(\mathrm{Ens})
\]
be the functor defined by
\[
 h_R(S)=\operatorname{Hom}_{\Lambda\text{-alg}}(R,S)
 \qquad(S\in\underline C_\Lambda).
\]
It determines a cofibered discrete groupoid $\underline h_R$ over
$\underline C_\Lambda$. A map $R\longrightarrow S$ in
$\widehat{\underline C}_\Lambda$ induces a functor
\[
 \underline h_S\longrightarrow\underline h_R.
\]
This functor is smooth if and only if $S$ is a formal power-series
ring over $R$.

%%%%%%%%%%  scan idx 62  |  folio 51  %%%%%%%%%%

\underline{Remark \SGARefTarget{sga7.sga7i.evi.n2-2.item-d}{2.2}(d).} Given
$R\in\widehat{\underline C}_\Lambda$, denote by
$(R,\underline C_\Lambda)$ the category of arrows in
$\widehat{\underline C}_\Lambda$ with source $R$ and target in
$\underline C_\Lambda$. Let $\underline A$ be a cofibered groupoid
over $\underline C_\Lambda$. Given
$\xi\in\operatorname{Ob}\widehat{\underline A}$ with $P(\xi)=R$,
let
\[
 \underline A'\longrightarrow(R,\underline C_\Lambda)
\]
be the pullback of
$\underline A\longrightarrow\underline C_\Lambda$ along the
canonical functor
$(R,\underline C_\Lambda)\longrightarrow\underline C_\Lambda$.
Choose a section $\xi^\#$ such that
$\xi^\#(1_R)=\xi$. Thus, for every map
$f:R\longrightarrow S$, we choose an $f$-map
$\xi\longrightarrow\xi(f)$.

We obtain a functor
\[
 \xi^\#:\underline h_R\longrightarrow\underline A
\]
over $\underline C_\Lambda$ by sending
$(R\xrightarrow{f}S)$ to $\xi(f)$. Up to natural equivalence, this
functor does not depend on the choice of $\xi^\#$.

The assertion that
$\xi^\#:\underline h_R\longrightarrow\underline A$ is
(minimally) smooth means that every surjective map
$a'\xrightarrow{\alpha}\xi(f)$ in $\underline A$ can be completed
to a commutative diagram
\[
\begin{tikzcd}[row sep=large,column sep=huge]
 \xi(f') \arrow[r] \arrow[dr,"{\xi(\,)}"']
 & a' \arrow[d,"\alpha"]\\
 & \xi(f)
\end{tikzcd}
\]
(and, in the minimal case, that
$\underline h_R(k[\varepsilon])\longrightarrow
\underline A(k[\varepsilon])$ is bijective). This means precisely
that $\xi$ is a (minimally) versal formal object of $\underline A$
over $\underline C_\Lambda$.

Thus the existence of a (minimally) versal formal object of
$\underline A$ is equivalent to the existence of a (minimally)
smooth functor
\[
 \underline h_R\longrightarrow\underline A
\]
over $\underline C_\Lambda$ for some
$R\in\operatorname{Ob}(\widehat{\underline C}_\Lambda)$.


% END INCLUDED SGA7I_English_current_buildable_source_20260823/components/18_expose_VI_section_2_opening_through_2_2.tex

% BEGIN INCLUDED SGA7I_English_current_buildable_source_20260823/components/19_expose_VI_2_3_through_2_7.tex
The following general properties of smoothness are immediate.

\underline{Proposition \SGARefTarget{sga7.sga7i.evi.n2-3}{2.3}.} Let
\[
 \underline A\xrightarrow{\ \varphi\ }\underline B
 \xrightarrow{\ \psi\ }\underline C
\]
be functors of cofibered groupoids over $\underline C_\Lambda$.

(a) If $\varphi$ and $\psi$ are both smooth, then so is
$\psi\circ\varphi$.

(b) If $\psi\circ\varphi$ is smooth and $\varphi$ is surjective on
isomorphism classes of objects, then $\psi$ is smooth.

%%%%%%%%%%  scan idx 63  |  folio 52  %%%%%%%%%%

\underline{Remark \SGARefTarget{sga7.sga7i.evi.n2-4}{2.4}.} Let $\underline A$ be a cofibered groupoid
over $\underline C_\Lambda$ admitting a versal formal object $\xi$.
By \SGARefLink{sga7.sga7i.evi.n2-2.item-d}{2.2}(d), this means that
\[
 \xi:\underline h_R\longrightarrow\underline A
\]
is smooth, where $R=P(\xi)$. Therefore the following assertions are
equivalent:

\begin{itemize}
\item $\underline A$ is formally smooth over $\underline C_\Lambda$
      (cf. \SGARefLink{sga7.sga7i.evi.n1-13}{1.13}), that is, $R$ is a formal power-series algebra over
      $\Lambda$;
\item the composite functor
      \[
       \underline h_R\xrightarrow{\ \xi\ }\underline A
       \xrightarrow{\ P\ }\underline C_\Lambda
       \;(=\underline h_\Lambda)
      \]
      is smooth;
\item $P:\underline A\longrightarrow\underline C_\Lambda$ is
      smooth.
\end{itemize}
The last equivalence follows from Proposition \SGARefLink{sga7.sga7i.evi.n2-3}{2.3}(b).

We now consider a stronger property than semi-homogeneity.

\underline{Definition \SGARefTarget{sga7.sga7i.evi.n2-5}{2.5}.} A cofibered groupoid $\underline A$ over
$\underline C_\Lambda$ is called \underline{homogeneous} if every
diagram
\[
\begin{tikzcd}[row sep=large,column sep=huge]
 & a' \arrow[d]\\
 b' \arrow[r] & a
\end{tikzcd}
\]
in which $a'\longrightarrow a$ is surjective admits a fibre product
$a'\times_a b'$ in $\underline A$ such that
\[
 P(a'\times_a b')=
 P(a')\times_{P(a)}P(b').
\]

\underline{Remark \SGARefTarget{sga7.sga7i.evi.n2-6.item-a}{\SGARefAliasTarget{sga7.sga7i.evi.n2-6}2.6}(a).} Let $\underline A$ be a cofibered
groupoid over $\underline C_\Lambda$, and consider a diagram
\[
(*)\qquad
\begin{tikzcd}[row sep=small,column sep=small]
 & a' \arrow[dl] \arrow[dr] & \\
 a_1 \arrow[dr] & & a_2 \arrow[dl]\\
 & a &
\end{tikzcd}
\qquad\text{in $\underline A$ above the diagram}\qquad
\begin{tikzcd}[row sep=small,column sep=small]
 & R_1\times_R R_2 \arrow[dl,"{\pi_1}"'] \arrow[dr,"{\pi_2}"] & \\
 R_1 \arrow[dr] & & R_2 \arrow[dl]\\
 & R &
\end{tikzcd}
\]
in $\underline C_\Lambda$, where the $\pi_i$ are the projections.
If a fibre product $a_1\times_a a_2$ exists in $\underline A$,
there is a unique map
\[
 \varphi:a'\longrightarrow a_1\times_a a_2
\]
compatible with the projections. If
\[
 P(a_1\times_a a_2)=R_1\times_R R_2,
\]
then $P(\varphi)$ is the identity map of $R_1\times_R R_2$;
hence $\varphi$ is a $\underline C_\Lambda$-isomorphism, that is,
$a'=a_1\times_a a_2$. Thus, if $\underline A$ is homogeneous,
every diagram $(*)$ defines a fibre product $a_1\times_a a_2$
whenever $R_1\longrightarrow R$ is surjective. It follows at once

%%%%%%%%%%  scan idx 64  |  folio 53  %%%%%%%%%%

that every homogeneous cofibered groupoid is, a fortiori,
semi-homogeneous.

\underline{Remark \SGARefTarget{sga7.sga7i.evi.n2-6.item-b}{2.6}(b).} Consider a diagram of categories and
functors
\[
\begin{tikzcd}[row sep=large,column sep=large]
 X \arrow[dr,"{\varphi}"'] & & Y \arrow[dl,"{\psi}"]\\
 & Z &
\end{tikzcd}
\]
The ordinary fibre product $X\times_ZY$ is the category whose
objects are pairs $(x,y)$ satisfying $\varphi(x)=\psi(y)$ and whose
maps $(x,y)\longrightarrow(x',y')$ are pairs
\[
 \alpha:x\longrightarrow x',\qquad
 \beta:y\longrightarrow y'
\]
such that $\varphi(\alpha)=\psi(\beta)$. The fibre product
$X\times_Z^2Y$ in the sense of $2$-categories is the category whose
objects are triples $(x,y;\gamma)$, where
\[
 \varphi(x)\xrightarrow{\ \gamma\ }\psi(y)
\]
is an isomorphism. A map
$(x,y;\gamma)\longrightarrow(x',y';\gamma')$ is a pair
$\alpha:x\longrightarrow x'$ and $\beta:y\longrightarrow y'$ for
which the diagram
\[
\begin{tikzcd}[row sep=huge,column sep=huge]
 \varphi(x) \arrow[r,"{\varphi(\alpha)}"] \arrow[d,"{\gamma}"']
 & \varphi(x') \arrow[d,"{\gamma'}"]\\
 \psi(y) \arrow[r,"{\psi(\beta)}"'] & \psi(y')
\end{tikzcd}
\]
commutes.

The notions of a (semi-)homogeneous cofibered groupoid and of a
smooth functor can be rephrased in terms of fibre products of
categories in this $2$-categorical sense. Let
$P:\underline A\longrightarrow\underline C_\Lambda$ be a
cofibered category, and consider $R_1\times_RR_2$ in
$\underline C_\Lambda$. Choose direct-image functors
\[
 (\pi_i)_*:\underline A(R_1\times_RR_2)
 \longrightarrow\underline A(R_i),\qquad i=1,2.
\]
They induce a functor
\[
 \underline A(R_1\times_RR_2)\longrightarrow
 \underline A(R_1)\times^2_{\underline A(R)}\underline A(R_2).
\]
Up to canonical natural equivalence, it is independent of the
chosen direct-image functors. The groupoid $\underline A$ is
homogeneous if and only if this functor is an equivalence whenever
$R_1\longrightarrow R$ is surjective; it is semi-homogeneous if and
only if
\[
 [\underline A(R_1\times_RR_2)]\longrightarrow
 [\underline A(R_1)]
 \times_{[\underline A(R)]}[\underline A(R_2)]
\]

%%%%%%%%%%  scan idx 65  |  folio 54  %%%%%%%%%%

is surjective whenever $R_1\longrightarrow R$ is surjective, and is
bijective whenever $R=k$ and $R_1=k[\varepsilon]$.

Now let $\varphi:\underline A\longrightarrow\underline B$ be a
functor of cofibered groupoids over $\underline C_\Lambda$. For
every map $R'\longrightarrow R$ in $\underline C_\Lambda$,
$\varphi$ induces, up to canonical natural equivalence, a functor
\[
 \underline A(R')\longrightarrow
 \underline A(R)\times^2_{\underline B(R)}\underline B(R').
\]
The functor $\varphi$ is smooth if and only if this functor is
surjective on isomorphism classes of objects whenever
$R'\longrightarrow R$ is surjective.

\underline{Remark \SGARefTarget{sga7.sga7i.evi.n2-6.item-c}{2.6}(c).} If $\underline A$ is homogeneous, then
for any finite-dimensional $k$-vector spaces $V$ and $W$ the functor
\[
 \underline A(k[V\times W])\longrightarrow
 \underline A(k[V])\times\underline A(k[W])
\]
is an equivalence of categories. Recall our convention
$\underline A(k)=\{e\}$. In particular,
\[
 \operatorname{Aut}_{\underline A(k[V\times W])}(1_{V\times W})
 \xrightarrow{\ \sim\ }
 \operatorname{Aut}_{\underline A(k[V])}(1_V)
 \times
 \operatorname{Aut}_{\underline A(k[W])}(1_W),
\]
where $1_V$ is the direct image of $e$ under the unique map
$k\longrightarrow k[V]$ in $\underline C_\Lambda$. Consequently,
if $\underline A$ is homogeneous, both
$\operatorname{Aut}(1_\varepsilon)$ and
$[\underline A(k[\varepsilon])]$ carry canonical $k$-vector-space
structures.

\underline{Remark \SGARefTarget{sga7.sga7i.evi.n2-6.item-d}{2.6}(d).} For every
$R\in\widehat{\underline C}_\Lambda$, the cofibered discrete
groupoid $\underline h_R$ over $\underline C_\Lambda$ is
homogeneous.

If $\underline A$ is semi-homogeneous, then so is $[\underline A]$.
Homogeneity of $\underline A$, however, does not in general imply
homogeneity of $[\underline A]$.

\underline{Proposition \SGARefTarget{sga7.sga7i.evi.n2-7}{2.7}.} Let $\underline A$ be a homogeneous
cofibered groupoid over $\underline C_\Lambda$. The following
conditions are equivalent.

(i) $[\underline A]$ is homogeneous; that is,
\[
 [\underline A(R'\times_RS)]\longrightarrow
 [\underline A(R')]
 \times_{[\underline A(R)]}[\underline A(S)]
\]
is bijective for every small extension $R'\longrightarrow R$.

(ii) For every small prolongation $a'\longrightarrow a$ in
$\underline A$, the map
\[
 \operatorname{Aut}_{\underline A(R')}(a')
 \longrightarrow
 \operatorname{Aut}_{\underline A(R)}(a)
\]

%%%%%%%%%%  scan idx 66  |  folio 55  %%%%%%%%%%

is surjective, where $R'=P(a')$ and $R=P(a)$.

\underline{Proof.} (i)$\Rightarrow$(ii). Let
$f:R'\longrightarrow R$ be a small extension in
$\underline C_\Lambda$, and let $a'\longrightarrow a$ be an
$f$-map in $\underline A$. Given
$\tau\in\operatorname{Aut}_{\underline A(R)}(a)$, consider
$a'\times_a a'$ and $a'\times_a a'_\tau$, where
$a'_\tau\longrightarrow a$ is the composite
\[
 a'\longrightarrow a\xrightarrow{\ \tau\ }a.
\]
Since
\[
 [\underline A(R'\times_RR')]
 \longrightarrow
 [\underline A(R')]
 \times_{[\underline A(R)]}[\underline A(R')]
\]
is bijective, there is an isomorphism
\[
 a'\times_a a'\xrightarrow{\ u\ }a'\times_a a'_\tau.
\]
Projection onto the second factor then gives an isomorphism
$\tau':a'\longrightarrow a'$ lifting $\tau$.

(ii)$\Rightarrow$(i). Let $a''_i$ $(i=1,2)$ be objects of
$\underline A(R'\times_RS)$ with the same image under
\[
 [\underline A(R'\times_RS)]\longrightarrow
 [\underline A(R')]
 \times_{[\underline A(R)]}[\underline A(S)].
\]
We obtain a diagram of solid arrows
\[
\begin{tikzcd}[column sep=small,row sep=small]
 && a''_2 \arrow[ddll] \arrow[ddrr] &&\\
 && a''_1 \arrow[u,dashed] \arrow[dl] \arrow[dr] &&\\
 a'_2 \arrow[ddrr] & a'_1 \arrow[l,"{\sim}"'] \arrow[dr]
 && b_1 \arrow[r,"{\tau}"] \arrow[dl] & b_2 \arrow[ddll]\\
 && a_1 \arrow[d,"{\bar\tau}"] &&\\
 && a_2 &&
\end{tikzcd}
\qquad\text{above}\qquad
\begin{tikzcd}[column sep=small,row sep=small]
 & R'\times_RS \arrow[dl] \arrow[dr] &\\
 R' \arrow[dr] && S \arrow[dl]\\
 & R &
\end{tikzcd}
\]
where $\bar\tau:a_1\longrightarrow a_2$ is the isomorphism induced
by $\tau$, so the lower-right square commutes. Under condition (ii),
the given isomorphism $a'_1\xrightarrow{\sim}a'_2$ can be replaced
so that the left square also commutes. Since
\[
 a''_i=a'_i\times_{a_i}b_i\qquad(i=1,2)
\]
by Remark \SGARefLink{sga7.sga7i.evi.n2-6.item-a}{2.6}(a), we obtain an isomorphism
$a''_1\xrightarrow{\sim}a''_2$. Hence
\[
 [\underline A(R'\times_RS)]\longrightarrow
 [\underline A(R')]
 \times_{[\underline A(R)]}[\underline A(S)]
\]
is injective, and therefore bijective because $\underline A$ is
homogeneous by hypothesis. \hfill$\square$


% END INCLUDED SGA7I_English_current_buildable_source_20260823/components/19_expose_VI_2_3_through_2_7.tex

% BEGIN INCLUDED SGA7I_English_current_buildable_source_20260823/components/20_expose_VI_2_8_through_2_9.tex
Let $\underline A$ be a homogeneous cofibered groupoid over
$\underline C_\Lambda$, and let $f:R'\longrightarrow R$ be a small
extension in $\underline C_\Lambda$. There is a canonical map

%%%%%%%%%%  scan idx 67  |  folio 56  %%%%%%%%%%

\[
 \widetilde f:R'\times_RR'\longrightarrow k[\operatorname{Ker}f],
 \qquad
 (r'_1,r'_2)\longmapsto r'_1(0)+(r'_1-r'_2),
\]
and
\[
 1\times\widetilde f:R'\times_RR'
 \longrightarrow R'\times_k k[\operatorname{Ker}f]
\]
is an isomorphism. Given two $f$-maps
$a_i\longrightarrow a$ $(i=1,2)$ in $\underline A$, the fibre
product $a_1\times_a a_2$ is an object over $R'\times_RR'$. We may
therefore form
\[
 \tau(a_1,a_2)=
 (\widetilde f)_*(a_1\times_a a_2)
 \in\underline A(k[\operatorname{Ker}f]).
\]
Since $1\times\widetilde f$ is an isomorphism, there is an
isomorphism
\[
 a_1\times_a a_2\simeq a_1\times\tau(a_1,a_2)
\]
compatible with projection onto the first factor. Statement \SGARefLink{sga7.sga7i.evi.n1-5}{1.5}(2)
now takes the following more definite form.

\underline{Lemma \SGARefTarget{sga7.sga7i.evi.n2-8}{2.8}.} Let $\underline A$ be a homogeneous
cofibered groupoid over $\underline C_\Lambda$, and let
\[
 a_i\xrightarrow{\ \alpha_i\ }a\qquad(i=1,2)
\]
be $f$-maps in $\underline A$.

(1) There exists $\rho:a_1\longrightarrow a_2$ in $\underline A$
with $\alpha_2\circ\rho=\alpha_1$ if and only if there is an arrow
$a_1\longrightarrow\tau(a_1,a_2)$ in $\underline A$.

(2) There exists such a $\rho$ with $P(\rho)=\operatorname{id}_{R'}$
if and only if
\[
 \tau(a_1,a_2)=0
 \quad\text{in}\quad
 [\underline A(k[\operatorname{Ker}f])].
\]

\underline{Proof.} We prove both assertions simultaneously. The
existence of $\rho:a_1\longrightarrow a_2$ satisfying
$\alpha_2\circ\rho=\alpha_1$ (and $P(\rho)=\operatorname{id}_{R'}$)
is equivalent successively to the existence of
\[
 h:a_1\longrightarrow a_1\times_a a_2,
 \qquad \operatorname{pr}_1\circ h=\operatorname{id}_{a_1},
\]
with $P(h)$ the diagonal map; to the existence of
\[
 g:a_1\longrightarrow a_1\times\tau(a_1,a_2),
 \qquad \operatorname{pr}_1\circ g=\operatorname{id}_{a_1},
\]
with
\[
 P(g):R'\longrightarrow R'\times_k k[\operatorname{Ker}f],
 \qquad r'\longmapsto(r',r'(0));
\]
and to the existence of an arrow
$\gamma:a_1\longrightarrow\tau(a_1,a_2)$ in $\underline A$, with
\[
 P(\gamma):R'\longrightarrow k[\operatorname{Ker}f],
 \qquad r'\longmapsto r'(0).
\]
Such an arrow factors through
$R'\longrightarrow k\longrightarrow k[\operatorname{Ker}f]$ if and
only if $\tau(a_1,a_2)=0$. \hfill$\square$

\underline{Proposition \SGARefTarget{sga7.sga7i.evi.n2-9}{2.9}.} (a) Let
\[
 \underline A\xrightarrow{\ \varphi\ }\underline B
 \xrightarrow{\ \psi\ }\underline C
\]
be functors of cofibered groupoids over $\underline C_\Lambda$.
Assume that

(i) $\underline B$ is semi-homogeneous and $\underline C$ is
homogeneous;

(ii) the map
\[
 \psi(k[\varepsilon]):[\underline B(k[\varepsilon])]
 \longrightarrow[\underline C(k[\varepsilon])]
\]
is injective.

Then $\psi\circ\varphi$ is smooth if and only if $\varphi$ and
$\psi$ are both

%%%%%%%%%%  scan idx 68  |  folio 57  %%%%%%%%%%

smooth.

(b) Let $\varphi:\underline A\longrightarrow\underline B$ be a
smooth functor of cofibered groupoids over $\underline C_\Lambda$.
If $\underline A$ is semi-homogeneous and $\underline B$ is
homogeneous, then
\[
 \widehat\varphi:\widehat{\underline A}
 \longrightarrow\widehat{\underline B}
\]
is smooth. Explicitly, for every surjective map
$b'\xrightarrow{\beta}\widehat\varphi(a)$ in
$\widehat{\underline B}$, there are a map
$\alpha:a'\longrightarrow a$ in $\widehat{\underline A}$ and a
$\widehat{\underline C}_\Lambda$-isomorphism
\[
 \rho:\widehat\varphi(a')\xrightarrow{\ \sim\ }b'
\]
such that $\beta\circ\rho=\varphi(\alpha)$.

\underline{Proof.} (a) Suppose $\psi\circ\varphi$ is smooth. If
$\varphi$ is smooth, then $\psi$ is smooth by Proposition \SGARefLink{sga7.sga7i.evi.n2-3}{2.3}(b),
so it remains to prove that $\varphi$ is smooth. Let
$b'\xrightarrow{\beta}\varphi(a)$ be a small prolongation in
$\underline B$. Since $\psi\circ\varphi$ is smooth, there are an
arrow $a'\xrightarrow{\alpha}a$ in $\underline A$ and a
$\underline C_\Lambda$-isomorphism
\[
 \rho:\psi(b')\longrightarrow\psi\varphi(a')
\]
for which
\[
\begin{tikzcd}[column sep=huge,row sep=huge]
 \psi(b') \arrow[rr,"{\rho}"] \arrow[dr,"{\psi(\beta)}"']
 && \psi\varphi(a') \arrow[dl,"{\psi\varphi(\alpha)}"]\\
 & \psi\varphi(a) &
\end{tikzcd}
\]
commutes. Since $\underline B$ is semi-homogeneous, we can choose a
commutative diagram
\[
\begin{tikzcd}[column sep=small,row sep=small]
 & b'' \arrow[dl] \arrow[dr] &
 &&& R'\times_RR' \arrow[dl] \arrow[dr] &\\
 b' \arrow[dr,"{\beta}"'] &&
 \varphi(a') \arrow[dl,"{\varphi(\alpha)}"]
 & \text{above} &
 R' \arrow[dr] && R' \arrow[dl]\\
 & \varphi(a) &&&& R &
\end{tikzcd}
\]
where $R=P(a)$ and $R'=P(b')=P(a')$. Homogeneity of
$\underline C$ gives
\[
 \psi(b'')=
 \psi(b')\times_{\psi\varphi(a)}\psi\varphi(a')
\]
by Remark \SGARefLink{sga7.sga7i.evi.n2-6.item-a}{2.6}(a). Lemma \SGARefLink{sga7.sga7i.evi.n2-8}{2.8}(2) then gives
\[
 0=\tau(\psi\varphi(a'),\psi(b'))
   =(\widetilde f)_*\psi(b'').
\]
By hypothesis (ii), it follows that $(\widetilde f)_*(b'')=0$.
Statement \SGARefLink{sga7.sga7i.evi.n1-5}{1.5}(2) now provides a map
$\gamma:b'\longrightarrow\varphi(a')$ satisfying
$\beta=\varphi(\alpha)\circ\gamma$. Thus $\varphi$ is smooth, and
then so is $\psi$. The converse follows from Proposition \SGARefLink{sga7.sga7i.evi.n2-3}{2.3}(a).

(b) Let a surjective map
$\beta:b'\longrightarrow\widehat\varphi(a)$ in
$\widehat{\underline B}$ be given. It is represented by a sequence
of commutative diagrams

%%%%%%%%%%  scan idx 69  |  folio 58  %%%%%%%%%%

\[
\begin{tikzcd}[column sep=large]
 {} \arrow[r,dashed] & b'_3 \arrow[r,"{\beta'_2}"] \arrow[d,"{\beta_3}"]
 & b'_2 \arrow[r,"{\beta'_1}"] \arrow[d,"{\beta_2}"]
 & b'_1 \arrow[d,"{\beta_1}"]\\
 {} \arrow[r,dashed] & \varphi(a_3) \arrow[r,"{\varphi(\alpha_2)}"']
 & \varphi(a_2) \arrow[r,"{\varphi(\alpha_1)}"']
 & \varphi(a_1)
\end{tikzcd}
\]
in $\underline B$, where every $\alpha_v$, $\beta_v$, and
$\beta'_v$ is surjective and
$\beta=\varprojlim_v\beta_v$. It is enough to show inductively that
every diagram of solid arrows below, with
$\alpha_n$, $\beta_n$, $\beta_{n+1}$, and $\beta'_n$ surjective,
can be completed by the dotted arrows:
\[
\begin{tikzcd}[column sep=large,row sep=large]
 \varphi(a'_{n+1})
   \arrow[rrrr,dashed,"{\varphi(\,)}"]
   \arrow[dr,dashed,leftrightarrow]
   \arrow[ddr,dashed,"{\varphi(\,)}"']
 &&&& \varphi(a'_n)
   \arrow[dl,leftrightarrow]
   \arrow[ddl,"{\varphi(\alpha'_n)}"]\\
 & b'_{n+1} \arrow[rr,"{\beta'_n}"] \arrow[d,"{\beta_{n+1}}"']
 && b'_n \arrow[d,"{\beta_n}"'] &\\
 & \varphi(a_{n+1}) \arrow[rr,"{\varphi(\alpha_n)}"']
 && \varphi(a_n) &
\end{tikzcd}
\]
By hypothesis $\underline A$ is semi-homogeneous and
$\alpha_n:a'_{n}\longrightarrow a_n$ is surjective. Hence there is
a commutative diagram
\[
\begin{tikzcd}[column sep=small,row sep=small]
 & a' \arrow[dl,"{\gamma}"'] \arrow[dr,"{\gamma'}"] &
 &&& R_{n+1}\times_{R_n}R'_n \arrow[dl] \arrow[dr] &\\
 a_{n+1} \arrow[dr,"{\alpha_n}"'] &&
 a'_n \arrow[dl,"{\alpha'_n}"]
 & \text{in $\underline A$ over} &
 R_{n+1} \arrow[dr] && R'_n \arrow[dl]
 & \text{in $\underline C_\Lambda$}\\
 & a_n &&&& R_n &&
\end{tikzcd}
\]
Homogeneity of $\underline B$ implies that
\[
\begin{tikzcd}[column sep=huge,row sep=huge]
 & \varphi(a') \arrow[dl,"{\varphi(\gamma)}"']
                  \arrow[dr,"{\varphi(\gamma')}"] &\\
 \varphi(a_{n+1}) \arrow[dr,"{\varphi(\alpha_n)}"']
 && \varphi(a'_n) \arrow[dl,"{\varphi(\alpha'_n)}"]\\
 & \varphi(a_n) &
\end{tikzcd}
\]

%%%%%%%%%%  scan idx 70  |  folio 59  %%%%%%%%%%

is a fibre-product diagram in $\underline B$. Thus there is a unique
map $\delta:b'_{n+1}\longrightarrow\varphi(a')$ with the required
commutativity. Since $\beta_{n+1}$ and $\beta'_n$ are surjective, so
is $\delta$. Smoothness of $\varphi$ therefore gives a map
$a'_{n+1}\xrightarrow{\alpha'}a'$ in $\underline A$ and a
$\underline C_\Lambda$-isomorphism
\[
 \rho:\varphi(a'_{n+1})\xrightarrow{\ \sim\ }b'_{n+1}
\]
such that $\delta\circ\rho=\varphi(\alpha')$. We obtain the
commutative diagram
\[
\begin{tikzcd}[row sep=huge,column sep=huge]
 \varphi(a'_{n+1})
   \arrow[rrrr,"{\varphi(\gamma'\circ\alpha')}"]
   \arrow[dr,"{\rho}"]
   \arrow[ddr,"{\varphi(\gamma\circ\alpha')}"']
 &&&& \varphi(a'_n)
   \arrow[dl,"{\sim}"']
   \arrow[ddl,"{\varphi(\alpha'_n)}"]\\
 & b'_{n+1} \arrow[rr] \arrow[d]
 && b'_n \arrow[d] &\\
 & \varphi(a_{n+1}) \arrow[rr]
 && \varphi(a_n) &
\end{tikzcd}
\]
which completes the induction and proves (b). \hfill$\square$


% END INCLUDED SGA7I_English_current_buildable_source_20260823/components/20_expose_VI_2_8_through_2_9.tex

% BEGIN INCLUDED SGA7I_English_current_buildable_source_20260823/components/21_expose_VI_corollary_2_10.tex
\underline{Corollary \SGARefTarget{sga7.sga7i.evi.n2-10}{2.10}.} (a) Let
$\varphi:\underline A\longrightarrow\underline B$ be a smooth functor
of cofibered groupoids over $\underline C_\Lambda$. Let $\xi$ be a
versal formal object of $\underline A$, and let $\eta$ be a minimally
versal formal object of $\underline B$. If $\underline B$ is
homogeneous, then for every map
$\beta:\eta\longrightarrow\varphi(\xi)$ in
$\widehat{\underline B}$, the ring $P(\xi)$ is a formal power-series
ring over $P(\eta)$ via
\[
 P(\beta):P(\eta)\longrightarrow P(\varphi(\xi))=P(\xi).
\]

(b) Let $\underline A$ be a homogeneous cofibered groupoid over
$\underline C_\Lambda$, and let $\xi$ and $\eta$ be versal formal
objects of $\underline A$. If $\eta$ is minimal, then for every map
$\alpha:\eta\longrightarrow\xi$ in $\widehat{\underline A}$, the
ring $P(\xi)$ is a formal power-series ring over $P(\eta)$ via
\[
 P(\alpha):P(\eta)\longrightarrow P(\xi).
\]

\underline{Proof.} (a) A map
$\beta:\eta\longrightarrow\varphi(\xi)$ induces a commutative
diagram
\[
\begin{tikzcd}[row sep=huge,column sep=huge]
 \underline h_{P(\xi)} \arrow[r,"{\widetilde\beta}"]
                           \arrow[d,"{\widetilde\xi}"']
 & \underline h_{P(\eta)} \arrow[d,"{\widetilde\eta}"]\\
 \underline A \arrow[r,"{\varphi}"'] & \underline B.
\end{tikzcd}
\]
Since $\widetilde\xi$ and $\varphi$ are smooth, the composite
$\widetilde\eta\circ\widetilde\beta
=\varphi\circ\widetilde\xi$ is smooth. On the other hand,

%%%%%%%%%%  scan idx 71  |  folio 60  %%%%%%%%%%

\[
 \widetilde\eta(k[\varepsilon]):
 \underline h_{P(\eta)}(k[\varepsilon])
 \longrightarrow[\underline B(k[\varepsilon])]
\]
is an isomorphism. Proposition \SGARefLink{sga7.sga7i.evi.n2-9}{2.9}(a) therefore shows that
$\widetilde\beta$ is smooth, which means precisely that $P(\xi)$ is
a formal power-series ring over $P(\eta)$.

(b) Apply (a) with $\underline A=\underline B$. \hfill$\square$

% END INCLUDED SGA7I_English_current_buildable_source_20260823/components/21_expose_VI_corollary_2_10.tex

% BEGIN INCLUDED SGA7I_English_current_buildable_source_20260823/components/22_expose_VI_representable_cogroupoids_intro.tex
We conclude by considering the ``representability'' of a cofibered
groupoid. Let $\underline C$ be a fixed category. Every object $R$ of
$\underline C$ defines a functor
\[
 h_R:\underline C\longrightarrow(\mathrm{Sets}),
 \qquad h_R(S)=\operatorname{Hom}_{\underline C}(R,S),
\]
and hence a fully faithful functor
\[
 h:\underline C^{\mathrm{op}}
 \longrightarrow
 \underline{\operatorname{Hom}}(\underline C,(\mathrm{Sets})).
\]
A \emph{cocategory in $\underline C$} is a diagram
\[
\begin{tikzcd}[column sep=huge]
 h_{R_1}\times_{h_{R_0}}h_{R_1} \arrow[r]
 & h_{R_1}
   \arrow[r,shift left=3,"{s_1}"]
   \arrow[r,shift right=3,"{s_2}"']
 & h_{R_0} \arrow[l,"{\epsilon}" description]
\end{tikzcd}
\qquad
\text{with }s_i\circ\epsilon=\operatorname{id}\quad(i=1,2),
\]
where the first arrow is denoted by $c$ and
\[
 h_{R_1}\times_{h_{R_0}}h_{R_1}
 = (h_{R_1})_{s_1}\times_{h_{R_0}}{}_{s_2}(h_{R_1}),
\]
such that

(i) $c$ is associative and compatible with $s_i$ for $i=1,2$;

(ii) both composites
\begin{align*}
 h_{R_1}&\xrightarrow{\ \sim\ }
 h_{R_0}\times_{h_{R_0}}h_{R_1}
 \xrightarrow{\ \epsilon\times1\ }
 h_{R_1}\times_{h_{R_0}}h_{R_1}
 \xrightarrow{\ c\ }h_{R_1},\\[1ex]
 h_{R_1}&\xrightarrow{\ \sim\ }
 h_{R_1}\times_{h_{R_0}}h_{R_0}
 \xrightarrow{\ 1\times\epsilon\ }
 h_{R_1}\times_{h_{R_0}}h_{R_1}
 \xrightarrow{\ c\ }h_{R_1}
\end{align*}
are identities.

If $(R_0,R_1,s_1,s_2,\epsilon,c)$ is a cocategory in
$\underline C$, then for every $S\in\underline C$ we define a
category $h_{(R_0,R_1,\ldots)}(S)$ by
\[
 \operatorname{ob}h_{(R_0,R_1,\ldots)}(S)
   =\operatorname{Hom}_{\underline C}(R_0,S),
 \qquad
 \operatorname{Fl}h_{(R_0,R_1,\ldots)}(S)
   =\operatorname{Hom}_{\underline C}(R_1,S),
\]
with composition supplied by
\[
 c:h_{R_1}\times_{h_{R_0}}h_{R_1}\longrightarrow h_{R_1}.
\]
The categories so obtained depend functorially on $S$ and define a
cofibered category
$\underline h_{(R_0,R_1,\ldots)}$ over $\underline C$. For example,
every object $R$ of $\underline C$ defines a discrete cocategory
\[
 \underline h_R
 =\underline h_{(R,R,s_1,s_2,\epsilon,c)}
\]
by taking $s_1,s_2,\epsilon$, and $c$ all to be identity maps. A
homomorphism
\[
 (R_0,R_1,\ldots)\longrightarrow(R'_0,R'_1,\ldots)
\]

%%%%%%%%%%  scan idx 72  |  folio 61  %%%%%%%%%%

of cocategories is a pair of maps of represented functors
$h_{R_0}\longrightarrow h_{R'_0}$ and
$h_{R_1}\longrightarrow h_{R'_1}$ whose induced map
\[
 h_{(R_0,R_1,\ldots)}(S)
 \longrightarrow h_{(R'_0,R'_1,\ldots)}(S)
\]
is a functor for every $S\in\underline C$. Every cocategory
$(R_0,R_1,\ldots)$ has, in particular, a homomorphism
\[
 \epsilon:\underline h_{R_0}
 \longrightarrow\underline h_{(R_0,R_1,\ldots)}.
\]
A \emph{cogroupoid in $\underline C$} is a cocategory
$(R_0,R_1,\ldots)$ for which $h_{(R_0,R_1,\ldots)}(S)$ is a groupoid
for every $S\in\underline C$; equivalently,
$\underline h_{(R_0,R_1,\ldots)}$ is a category cofibered in
groupoids over $\underline C$.

In practice, cocategories and cogroupoids arise as follows. Let
$P:\underline F\longrightarrow\underline C$ be a cofibered category,
and fix an object $\xi$ of $\underline F$. Put $R_0=P(\xi)$, and let
$(R_0,\underline C)$ denote the category of arrows in $\underline C$
with source $R_0$. Let
$P:\underline F'\longrightarrow(R_0,\underline C)$ be the pullback of
$\underline F$ along the canonical functor
$(R_0,\underline C)\longrightarrow\underline C$. Choose a
cocartesian section $\xi^{\#}$ of $\underline F'$ over
$(R_0,\underline C)$ such that
$\xi^{\#}(\operatorname{id}_{R_0})=\xi$. Thus, for every arrow
$R_0\xrightarrow fS$ in $\underline C$, we choose a cocartesian
$f$-map $\xi\longrightarrow\xi^{\#}(f)$.

Define
\[
 \operatorname{Hom}_{\xi^{\#}}:
 (R_0\amalg R_0,\underline C)\longrightarrow(\mathrm{Sets})
\]
by
\[
 X=(f_1,f_2)\longmapsto
 \operatorname{Hom}_{\underline F(S)}
   (\xi^{\#}(f_1),\xi^{\#}(f_2)),
\]
where $S$ is the common target of $f_1$ and $f_2$. Up to canonical
natural equivalence, this functor is independent of the choice of
$\xi^{\#}$; henceforth we write simply $\xi$ for $\xi^{\#}$.

Suppose the functor is represented by an object
\[
 X_0=
 \left(R_0\mathrel{\substack{\xrightarrow{\ s_1\ }\\[-.6ex]
                              \xrightarrow[\ s_2\ ]{}}}R_1\right)
\]
of $(R_0\amalg R_0,\underline C)$. Thus there is an $R_0$-map
$\varphi:\xi(s_1)\longrightarrow\xi(s_2)$ in $\underline F$ such
that
\[
 \widetilde\varphi:
 \operatorname{Hom}_{(R_0\amalg R_0,\underline C)}(X_0,X)
 \longrightarrow\operatorname{Hom}_\xi(X)
\]
is bijective for every $X$ in $(R_0\amalg R_0,\underline C)$. Then
$X_0$ naturally defines a cocategory $(R_0,R_1,\ldots)$ in
$\underline C$. Indeed, a triple of arrows

%%%%%%%%%%  scan idx 73  |  folio 62  %%%%%%%%%%

\[
\begin{tikzcd}[column sep=huge]
 R_0 \arrow[r,shift left=3,"u"]
     \arrow[r,"v"]
     \arrow[r,shift right=3,"w"']
 & S
\end{tikzcd}
\]
in $\underline C$ defines a composition
\[
 \operatorname{Hom}_{\underline F(S)}(\xi(v),\xi(w))
 \times
 \operatorname{Hom}_{\underline F(S)}(\xi(u),\xi(v))
 \longrightarrow
 \operatorname{Hom}_{\underline F(S)}(\xi(u),\xi(w)),
\]
and hence a composition rule
\[
 h_{R_1}(S)\times_{h_{R_0}(S)}h_{R_1}(S)
 \longrightarrow h_{R_1}(S),
\]
functorial in $S$. Consequently, $\operatorname{Hom}_\xi$ defines the
cocategory $(R_0,R_1,\ldots)$ in $\underline C$ and the cofibered
category $\underline h_{(R_0,R_1,\ldots)}$ over $\underline C$.
Explicitly,
\[
 \operatorname{ob}\underline h_{(R_0,R_1,\ldots)}
 =\coprod_{S\in\operatorname{ob}(\underline C)}
   \operatorname{Hom}_{\underline C}(R_0,S).
\]
An arrow
\[
 (R_0\xrightarrow fS)\longrightarrow(R_0\xrightarrow{f'}S')
\]
is a pair of arrows $S\xrightarrow hS'$ and $R_1\xrightarrow uS'$ in
$\underline C$ such that
\[
 u\circ s_1=h\circ f,
 \qquad u\circ s_2=f'.
\]
The structural functor
$P:\underline h_{(R_0,R_1,\ldots)}\longrightarrow\underline C$ sends
$(R_0\xrightarrow fS)$ to $S$.

The pair $(\xi,\varphi)$ now defines a functor
\[
 (\xi,\varphi):\underline h_{(R_0,R_1,\ldots)}
 \longrightarrow\underline F
 \quad\text{over }\underline C
\]
by
\[
 (R_0\xrightarrow fS)\longmapsto\xi(f)
\]
on objects and by
\[
 \left[(R_0\xrightarrow fS)
 \xrightarrow{(h,u)}(R_0\xrightarrow{f'}S')\right]
 \longmapsto
 \left[\xi(f)\xrightarrow{h_*}\xi(hf)
 \xrightarrow{u(\varphi)}\xi(f')\right]
\]
on arrows. For every $S\in\underline C$, the functor
\[
 (\xi,\varphi)(S):
 \underline h_{(R_0,R_1,\ldots)}(S)\longrightarrow\underline F(S)
\]
is fully faithful by the representability of
$\operatorname{Hom}_\xi$. It follows that
\[
 (\xi,\varphi):\underline h_{(R_0,R_1,\ldots)}
 \longrightarrow\underline F
\]
is a $\underline C$-equivalence if and only if $\xi$ is quasi-initial
in $\underline F$ over $\underline C$, that is, if every object $a$
of $\underline F$ admits a cocartesian arrow $\xi\longrightarrow a$.

One may instead consider the functor
\[
 \operatorname{Isom}_\xi:
 (R_0\amalg R_0,\underline C)\longrightarrow(\mathrm{Sets}),
\]

%%%%%%%%%%  scan idx 74  |  folio 63  %%%%%%%%%%

given by
\[
 X=(f_1,f_2)\longmapsto
 \operatorname{Isom}_{\underline F(S)}(\xi(f_1),\xi(f_2)),
\]
where $S$ is the common target of $f_1$ and $f_2$. If this functor is
represented by an object $(R_0\rightrightarrows R_1)$ of
$(R_0\amalg R_0,\underline C)$, then it defines a cogroupoid
$(R_0,R_1,\ldots)$ in $\underline C$, together with a functor
\[
 (\xi,\varphi):\underline h_{(R_0,R_1,\ldots)}
 \longrightarrow\underline F.
\]
When $\underline F$ is a cofibered groupoid, this functor is a
$\underline C$-equivalence if and only if $\xi$ is quasi-initial in
$\underline F$ over $\underline C$.

% END INCLUDED SGA7I_English_current_buildable_source_20260823/components/22_expose_VI_representable_cogroupoids_intro.tex

% BEGIN INCLUDED SGA7I_English_current_buildable_source_20260823/components/23_expose_VI_2_11_and_2_12_ab.tex
\underline{Definition and Proposition \SGARefTarget{sga7.sga7i.evi.n2-11}{2.11}.} A cogroupoid
$(R_0,R_1,\ldots)$ in $\widehat{\underline C}_\Lambda$ is called
\underline{smooth} if either, and hence both, of the following
equivalent conditions holds:

(i) $R_1$ is a formal power-series ring over $R_0$ via $s_1$ (and
via $s_2$);

(ii) the functor
\[
 \epsilon:\underline h_{R_0}
 \longrightarrow\underline h_{(R_0,R_1,\ldots)}
\]
is smooth over $\underline C_\Lambda$.

A \underline{minimally smooth cogroupoid}, also called a
\underline{normalized smooth cogroupoid}, in
$\widehat{\underline C}_\Lambda$ is a smooth cogroupoid
$(R_0,R_1,\ldots)$ for which
$\underline h_{(R_0,R_1,\ldots)}(k[\varepsilon])$ is totally
disconnected; equivalently,
\[
 s_1(k[\varepsilon])=s_2(k[\varepsilon]).
\]

\underline{Proof.} Smoothness of
$\epsilon:\underline h_{R_0}\to
\underline h_{(R_0,R_1,\ldots)}$ means that, for every surjective
map
\[
 (S',f')\xrightarrow{(h,u)}(S,f)
\]
in $\underline h_{(R_0,R_1,\ldots)}$, there are a lift $f''$ of $f$
to $S'$ and an isomorphism
\[
 (S',f'')\xrightarrow{(1,u')}(S',f').
\]
Equivalently, given $u\in h_{R_1}(S)$ and a lift $f'$ of $s_1(u)$
to $S'$, there is a lift $u'$ of $u$ to $S'$ such that
$s_1(u')=f'$, for every surjection $S'\to S$ in
$\underline C_\Lambda$. This is precisely the smoothness of
\[
 s_1:\underline h_{R_1}\longrightarrow\underline h_{R_0}
\]
over $\underline C_\Lambda$, and that in turn says that $R_1$ is a
formal power-series ring over $R_0$ via $s_1$. \hfill$\square$

\underline{Proposition \SGARefTarget{sga7.sga7i.evi.n2-12}{2.12}.} Let $(R_0,R_1,\ldots)$ be a cogroupoid
in $\widehat{\underline C}_\Lambda$.

(a) If $(R_0,R_1,\ldots)$ is smooth, then
$\underline h_{(R_0,R_1,\ldots)}$ is homogeneous.

(b) If $s_1(k[\varepsilon])=s_2(k[\varepsilon])$, then the following
conditions are equivalent:

(i) $(R_0,R_1,\ldots)$ is smooth;

(ii) $\underline h_{(R_0,R_1,\ldots)}$ is homogeneous;

(iii) $\underline h_{(R_0,R_1,\ldots)}$ is semi-homogeneous.

%%%%%%%%%%  scan idx 75  |  folio 64  %%%%%%%%%%

(c) If $\underline h_{(R_0,R_1,\ldots)}$ is semi-homogeneous, then
it is equivalent to $\underline h_{(S_0,S_1,\ldots)}$ for some
minimally smooth cogroupoid $S_*$. 

\underline{Proof.} (a) Consider a diagram
\[
\begin{tikzcd}[column sep=huge,row sep=large]
 (S_1,f_1) \arrow[dr,"{(h_1,u_1)}"']
 && (S_2,f_2) \arrow[dl,"{(h_2,u_2)}"]\\
 & (S,f) &
\end{tikzcd}
\]
in $\underline h_{(R_0,R_1,\ldots)}$, where
$h_2:S_2\longrightarrow S$ is surjective in
$\underline C_\Lambda$. Since
$s_1:\underline h_{R_1}\longrightarrow\underline h_{R_0}$ is
smooth, $u_1^{-1}u_2$ lifts to an element $v\in h_{R_1}(S_2)$ with
$s_1(v)=f_2$. We have
\[
 h_2\circ s_2(v)=h_1\circ f_1,
\]
so $f_1\times s_2(v)$ is an element of
$h_{R_0}(S_1\times_SS_2)$ and gives a commutative diagram
\[
\begin{tikzcd}[column sep=huge,row sep=large]
 & (S_1\times_SS_2,,f_1\times s_2(v))
   \arrow[dl,"{\pi_1}"']
   \arrow[dr,"{(\pi_2,v^{-1})}"] &\\
 (S_1,f_1) \arrow[dr,"{(h_1,u_1)}"']
 && (S_2,f_2) \arrow[dl,"{(h_2,u_2)}"]\\
 & (S,f) &
\end{tikzcd}
\]
We claim that its top object is the fibre product
\[
 (S_1,f_1)\times_{(S,f)}(S_2,f_2)
\]
in $\underline h_{(R_0,R_1,\ldots)}$. Indeed, for every object
$(T,g)$, a map to the top object is equivalently a tuple
$(h'_1\times h'_2,w_1,w_2)$ satisfying
\begin{align*}
 h_1w_1&=h_2w_2,&
 s_1(w_1)&=h'_1g,&s_2(w_1)&=f_1,\\
 s_1(w_2)&=h'_2g,&
 s_2(w_2)&=s_2(v).
\end{align*}
Putting $w'_2=v^{-1}w_2$ gives a bijection with the tuples
$(h'_1\times h'_2,w_1,w'_2)$ satisfying
\begin{align*}
 u_2h_2w'_2&=u_1h_1w_1,&
 s_1(w_1)&=h'_1g,&s_2(w_1)&=f_1,\\
 s_1(w'_2)&=h'_2g,&
 s_2(w'_2)&=f_2.
\end{align*}
These are exactly the pairs of arrows from $(T,g)$ to $(S_1,f_1)$
and $(S_2,f_2)$ whose composites to $(S,f)$ coincide. Thus
\begin{multline*}
 \operatorname{Hom}((T,g),(S_1\times_SS_2,
                 f_1\times s_2(v)))\\
 \simeq
 \operatorname{Hom}((T,g),(S_1,f_1))
 \times_{\operatorname{Hom}((T,g),(S,f))}
 \operatorname{Hom}((T,g),(S_2,f_2)),
\end{multline*}
which proves (a).

%%%%%%%%%%  scan idx 76  |  folio 65  %%%%%%%%%%

(b) It remains only to prove (iii)$\Rightarrow$(i). If
$\underline h_{(R_0,R_1,\ldots)}$ is semi-homogeneous, then
\[
 \dim[\underline h_{(R_0,R_1,\ldots)}(k[\varepsilon])]
 =\dim h_{R_0}(k[\varepsilon])<\infty,
\]
so it admits a minimally versal formal object. On the other hand,
$(R_0,1_{R_0})$ is a quasi-initial minimal object of
$\widehat{\underline h}_{(R_0,R_1,\ldots)}$, because
\[
 h_{R_0}(k[\varepsilon])
 \longrightarrow
 [\underline h_{(R_0,R_1,\ldots)}(k[\varepsilon])]
\]
is bijective. Proposition \SGARefLink{sga7.sga7i.evi.n1-10}{1.10}(ii) therefore makes
$(R_0,1_{R_0})$ a minimally versal formal object of
$\underline h_{(R_0,R_1,\ldots)}$. Equivalently,
\[
 \epsilon:\underline h_{R_0}
 \longrightarrow\underline h_{(R_0,R_1,\ldots)}
\]
is minimally smooth. This proves (i).

% END INCLUDED SGA7I_English_current_buildable_source_20260823/components/23_expose_VI_2_11_and_2_12_ab.tex

% BEGIN INCLUDED SGA7I_English_current_buildable_source_20260823/components/24_expose_VI_2_12_c_through_2_14.tex
(c) Let $(S_0,\xi_0)$ be a minimally versal formal object of
$\underline h_{(R_0,R_1,\ldots)}$ (Theorem \SGARefLink{sga7.sga7i.evi.n1-11}{1.11}). For any object
\[
 \xi:R_0\longrightarrow T
 \qquad\text{of }\underline h_{(R_0,R_1,\ldots)}(T),
\]
consider the functor
\[
 \operatorname{Isom}(\xi,\xi):
 \underline C_\Lambda/(T\widehat\otimes_\Lambda T)
 \longrightarrow(\mathrm{Sets}).
\]
It assigns to $R$ the set of isomorphisms between the two images of
$\xi$ on $R$ obtained from
\[
 R_0\xrightarrow{\ \xi\ }T
 \mathrel{\substack{\xrightarrow{}\\[-.6ex]\xrightarrow[]{}}}
 T\widehat\otimes_\Lambda T\longrightarrow R.
\]
This functor is represented by
\[
 R_1\otimes_{R_0\widehat\otimes_\Lambda R_0}
       (T\widehat\otimes_\Lambda T).
\]
Passing to the inverse limit shows that
$\operatorname{Isom}(\xi_0,\xi_0)$ is prorepresented by an object
$S_1$, and that
\[
 \underline h_{(R_0,R_1,\ldots)}
 \simeq\underline h_{(S_0,S_1,\ldots)}.
\]
By (b), the cogroupoid $S_*$ is minimally smooth. \hfill$\square$

\underline{Definition \SGARefTarget{sga7.sga7i.evi.n2-13}{2.13}.} A cofibered groupoid $\underline A$
over $\underline C_\Lambda$ is called \underline{(minimally)
prorepresentable} if it is $\underline C_\Lambda$-equivalent to
$\underline h_{(R_0,R_1,\ldots)}$ for some (normalized) cogroupoid
$(R_0,R_1,\ldots)$ in $\widehat{\underline C}_\Lambda$.

\underline{Proposition \SGARefTarget{sga7.sga7i.evi.n2-14}{2.14}.} If $(R_0,R_1,\ldots)$ is a normalized
cogroupoid in $\widehat{\underline C}_\Lambda$, then every
homomorphism
\[
 (\alpha_0,\alpha_1):(R_0,R_1,\ldots)
 \longrightarrow(R_0,R_1,\ldots)
\]
which induces a $\underline C_\Lambda$-equivalence
\[
 \underline h_{(\alpha_0,\alpha_1)}:
 \underline h_{(R_0,R_1,\ldots)}
 \longrightarrow\underline h_{(R_0,R_1,\ldots)}
\]
is an isomorphism. In

%%%%%%%%%%  scan idx 77  |  folio 66  %%%%%%%%%%

particular, a normalized prorepresentation of a cofibered groupoid
over $\underline C_\Lambda$, if it exists, is unique up to
isomorphism.

\underline{Proof.} Consider the commutative diagram
\[
\begin{tikzcd}[column sep=huge,row sep=huge]
 \underline h_{(R_0,R_1,\ldots)}
   \arrow[r,"{\underline h_{(\alpha_0,\alpha_1)}}"]
 & \underline h_{(R_0,R_1,\ldots)}\\
 \underline h_{R_0} \arrow[u,"\epsilon"']
   \arrow[r,"{\underline h_{\alpha_0}}"']
 & \underline h_{R_0} \arrow[u,"\epsilon"]
\end{tikzcd}
\]
in which $\underline h_{(\alpha_0,\alpha_1)}$ is a
$\underline C_\Lambda$-equivalence. Since
$\underline h_{(R_0,R_1,\ldots)}(k[\varepsilon])$ is totally
disconnected, both
\[
 [\underline h_{(\alpha_0,\alpha_1)}(k[\varepsilon])]
 \quad\text{and}\quad
 [\epsilon(k[\varepsilon])]
\]
are bijections, and hence so is
$\underline h_{\alpha_0}(k[\varepsilon])$. Thus
$\alpha_0:R_0\longrightarrow R_0$ is a surjective endomorphism and
therefore an automorphism. It follows that $\alpha_1$ is also an
automorphism of $R_1$. \hfill$\square$

% END INCLUDED SGA7I_English_current_buildable_source_20260823/components/24_expose_VI_2_12_c_through_2_14.tex

% BEGIN INCLUDED SGA7I_English_current_buildable_source_20260823/components/25_expose_VI_2_15_through_2_17.tex
\underline{Lemma \SGARefTarget{sga7.sga7i.evi.n2-15}{2.15}.} Let $A$ and $B$ be discrete groupoids over
$\underline C_\Lambda$, and let $\varphi:A\longrightarrow B$ be a
minimally smooth functor. If $A$ and $B$ are homogeneous, then
$\varphi$ is a $\underline C_\Lambda$-equivalence.

\underline{Proof.} We must show that
$\varphi(S):A(S)\longrightarrow B(S)$ is bijective for every
$S\in\operatorname{ob}(\underline C_\Lambda)$. It is surjective for
every $S$ and bijective for $S=k[\varepsilon]$. We argue by induction
on the length of $S$. Let $S'\xrightarrow fS$ be a small extension
in $\underline C_\Lambda$. Since $A$ and $B$ are homogeneous
discrete groupoids, the canonical isomorphism
\[
 1\times\widetilde f:S'\times_SS'
 \xrightarrow{\ \sim\ }S'\times_k k[\operatorname{Ker}f]
\]
induces a commutative diagram

%%%%%%%%%%  scan idx 78  |  folio 67  %%%%%%%%%%

\[
\begin{tikzcd}[column sep=huge,row sep=huge]
 A(S')\times_{A(S)}A(S') \arrow[r,"\sim"] \arrow[d]
 & A(S')\times A(k[\operatorname{Ker}f]) \arrow[d]\\
 B(S')\times_{B(S)}B(S') \arrow[r,"\sim"']
 & B(S')\times B(k[\operatorname{Ker}f]).
\end{tikzcd}
\]
Thus $A(S')$ and $B(S')$ are principal homogeneous sets over $A(S)$
and $B(S)$, with structural groups identified by
\[
 A(k[\operatorname{Ker}f])\xrightarrow{\ \sim\ }
 B(k[\operatorname{Ker}f]).
\]
If $\varphi(S)$ is bijective, it follows that $\varphi(S')$ is
injective, and hence bijective. This completes the induction.
\hfill$\square$

\underline{Corollary \SGARefTarget{sga7.sga7i.evi.n2-16}{2.16}.} Let $\underline A$ be a homogeneous
groupoid over $\underline C_\Lambda$ such that
\[
 \operatorname{Aut}_{\underline A(k[\varepsilon])}(1_\varepsilon)
\]
is finite-dimensional. For every
$\xi\in\operatorname{ob}(\widehat{\underline A})$, the functor
\[
 \operatorname{Isom}_\xi:
 \widehat{\underline C}_{R\widehat\otimes_\Lambda R}
 \longrightarrow(\mathrm{Sets})
\]
is prorepresentable, where $R=P(\xi)$.

\underline{Proof.} If $\operatorname{Isom}_\xi$ is homogeneous,
Theorem \SGARefLink{sga7.sga7i.evi.n1-11}{1.11} supplies a minimally smooth functor
\[
 \Phi:\underline h_{R_1}\longrightarrow
 \underline{\operatorname{Isom}}_\xi.
\]
Lemma \SGARefLink{sga7.sga7i.evi.n2-15}{2.15} makes $\Phi$ a
$\widehat{\underline C}_{R\widehat\otimes_\Lambda R}$-equivalence,
and hence $\operatorname{Isom}_\xi$ is prorepresentable. It remains
to prove homogeneity. Consider a diagram
\[
\begin{tikzcd}[column sep=large,row sep=large]
 T_1 \arrow[dr] && T_2 \arrow[dl]\\
 & T &
\end{tikzcd}
\]
in $\widehat{\underline C}_{R\widehat\otimes_\Lambda R}$, with
$T_2\longrightarrow T$ surjective. Since $\underline A$ is
homogeneous, the fibre products
\[
 \xi_i(T_1)\times_{\xi_i(T)}\xi_i(T_2)
\]
exist in $\underline A$, and there are canonical isomorphisms
\[
 \xi_i(T_1\times_TT_2)\xrightarrow{\ \sim\ }
 \xi_i(T_1)\times_{\xi_i(T)}\xi_i(T_2)
 \qquad(i=1,2)
\]
(cf. \SGARefLink{sga7.sga7i.evi.n2-6.item-a}{2.6}(a)). Therefore
\[
 \operatorname{Isom}_\xi(T_1\times_TT_2)
 \longrightarrow
 \operatorname{Isom}_\xi(T_1)
 \times_{\operatorname{Isom}_\xi(T)}
 \operatorname{Isom}_\xi(T_2)
\]
is an isomorphism. \hfill$\square$

\underline{Theorem \SGARefTarget{sga7.sga7i.evi.n2-17}{2.17}.} Let $\underline A$ be a cofibered groupoid
over $\underline C_\Lambda$. Then it is prorepresentable by a
(normalized) smooth cogroupoid if and only if

%%%%%%%%%%  scan idx 79  |  folio 68  %%%%%%%%%%

(1) $\underline A$ is homogeneous;

(2) $[\underline A(k[\varepsilon])]$ and
$\operatorname{Aut}_{\underline A(k[\varepsilon])}(1_\varepsilon)$
are finite-dimensional vector spaces over $k$.

\underline{Proof.} Suppose first that $\underline A$ is
prorepresented by a smooth cogroupoid. Thus $\underline A$ is
$\underline C_\Lambda$-equivalent to
\[
 \underline h=\underline h_{(R_0,R_1,\ldots)}
\]
for a smooth cogroupoid $(R_0,R_1,\ldots)$ in
$\widehat{\underline C}_\Lambda$. Proposition \SGARefLink{sga7.sga7i.evi.n2-12}{2.12}(a) shows that
$\underline h$ is homogeneous. Moreover,
\begin{align*}
 [\underline h(k[\varepsilon])]
 &=\operatorname{Coker}\!\left(
   \operatorname{Hom}_{\Lambda\text{-alg}}(R_1,k[\varepsilon])
   \mathrel{\substack{\xrightarrow{\ s_1\ }\\[-.6ex]
                       \xrightarrow[\ s_2\ ]{}}}
   \operatorname{Hom}_{\Lambda\text{-alg}}(R_0,k[\varepsilon])
   \right),\\[1ex]
 \operatorname{Aut}_{\underline h}(1_\varepsilon)
 &=\operatorname{Ker}\!\left(
   \operatorname{Hom}_{\Lambda\text{-alg}}(R_1,k[\varepsilon])
   \xrightarrow{\ s_1\ }
   \operatorname{Hom}_{\Lambda\text{-alg}}(R_0,k[\varepsilon])
   \right)\\
 &=\operatorname{Ker}\!\left(
   \operatorname{Hom}_{\Lambda\text{-alg}}(R_1,k[\varepsilon])
   \xrightarrow{\ s_2\ }
   \operatorname{Hom}_{\Lambda\text{-alg}}(R_0,k[\varepsilon])
   \right).
\end{align*}
Both are finite-dimensional vector spaces over $k$, so the same is
true for $\underline A$.

Conversely, assume (1) and (2). Since $\underline A$ is homogeneous
and $[\underline A(k[\varepsilon])]$ is finite-dimensional, it has a
minimally versal formal object $\xi$. Put $P(\xi)=R_0$ and consider
\[
 \operatorname{Isom}_\xi:
 \widehat{\underline C}_{R_0\widehat\otimes_\Lambda R_0}
 \longrightarrow(\mathrm{Sets}),
\]
given by
\[
 T\longmapsto
 \operatorname{Isom}_{\underline A(T)}(\xi_1(T),\xi_2(T)),
\]
where $\xi_i(T)=\xi(f_i)$ and $f_i:R_0\longrightarrow T$ are the
composite maps in
\[
\begin{tikzcd}[column sep=large]
 R_0 \arrow[r,shift left=2,"s_1"]
     \arrow[r,shift right=2,"s_2"']
 & R_0\widehat\otimes_\Lambda R_0 \arrow[r] & T.
\end{tikzcd}
\]
Corollary \SGARefLink{sga7.sga7i.evi.n2-16}{2.16} prorepresents this functor by an object
\[
 R_1\in
 \operatorname{ob}(\widehat{\underline C}_{R_0\widehat\otimes_\Lambda R_0})
\]
and an isomorphism
$\xi_1(R_1)\xrightarrow{\varphi}\xi_2(R_1)$ for which
\[
 \Phi:\underline h_{R_1}\longrightarrow
 \underline{\operatorname{Isom}}_\xi
\]
is an isomorphism over
$\widehat{\underline C}_{R_0\widehat\otimes_\Lambda R_0}$. The pair
$(R_0,R_1)$ therefore defines a cogroupoid, and
\[
 (\xi,\varphi):\underline h_{(R_0,R_1,\ldots)}
 \longrightarrow\underline A
\]
is a $\underline C_\Lambda$-equivalence. Because $\xi$ is minimally
versal, the composite
\[
 \underline h_{(R_0,R_0,\ldots)}
 \xrightarrow{\ \epsilon\ }
 \underline h_{(R_0,R_1,\ldots)}
 \xrightarrow{\ (\xi,\varphi)\ }
 \underline A
\]
is minimally smooth. Since $(\xi,\varphi)$ is an equivalence,
$\epsilon$ is minimally smooth as well. Thus
$(R_0,R_1,\ldots)$ is a normalized smooth cogroupoid in
$\widehat{\underline C}_\Lambda$. \hfill$\square$

% END INCLUDED SGA7I_English_current_buildable_source_20260823/components/25_expose_VI_2_15_through_2_17.tex

% BEGIN INCLUDED SGA7I_English_current_buildable_source_20260823/components/26_expose_VI_2_18_through_2_22.tex
%%%%%%%%%%  scan idx 80  |  folio 69  %%%%%%%%%%

\underline{Corollary \SGARefTarget{sga7.sga7i.evi.n2-18}{2.18}.} A functor
$F:\underline C_\Lambda\longrightarrow(\mathrm{Sets})$ is
prorepresentable if and only if it is homogeneous and
$\dim F(k[\varepsilon])<\infty$.

\underline{Corollary \SGARefTarget{sga7.sga7i.evi.n2-19}{2.19}.} Every smooth cogroupoid in
$\widehat{\underline C}_\Lambda$ can be normalized up to isomorphism.

Given a cogroupoid $(R_0,R_1,\ldots)$ in
$\widehat{\underline C}_\Lambda$, set $\bar R_1=R_1/J$, where $J$ is
the ideal of $R_1$ generated by
\[
 \{s_1(r)-s_2(r)\mid r\in R_0\}.
\]
Then $(R_0,\bar R_1,\ldots)$ defines a formal group scheme over $R_0$
with the induced composition law.

Now let $\underline A$ be a cofibered groupoid over
$\underline C_\Lambda$, prorepresented by a normalized smooth
cogroupoid $(R_0,R_1,\ldots)$ via $(\xi,\varphi)$. Let
\[
 \underline C_k\lhook\joinrel\longrightarrow
 \underline C_{R_0}\lhook\joinrel\longrightarrow
 \underline C_{R_0\widehat\otimes_\Lambda R_0}
\]
be the embeddings corresponding to the surjections
\[
 R_0\widehat\otimes_\Lambda R_0\longrightarrow R_0\longrightarrow k,
\]
and let $\operatorname{Aut}_\xi$ and $\operatorname{Aut}^{o}$ be the
restrictions of $\operatorname{Isom}_\xi$ to
$\underline C_{R_0}$ and $\underline C_k$, respectively. Thus
$\operatorname{Aut}_\xi(S)$ is the automorphism group of $\xi(S)$ for
$S\in\operatorname{ob}(\underline C_{R_0})$, while
$\operatorname{Aut}^{o}(S)$ is the automorphism group of $1_S$ for
$S\in\operatorname{ob}(\underline C_k)$, where $1_S$ is the direct
image of $1\in\underline A(k)$ under $k\longrightarrow S$.

Since $R_1$ prorepresents $\operatorname{Isom}_\xi$ on
$\underline C_{R_0\widehat\otimes_\Lambda R_0}$, it follows that
$\bar R_1$ prorepresents $\operatorname{Aut}_\xi$ on
$\underline C_{R_0}$, and that
$k\otimes_{R_0}\bar R_1$ prorepresents $\operatorname{Aut}^{o}$ on
$\underline C_k$. Although $\underline A$ is prorepresented by the
smooth cogroupoid $(R_0,R_1,\ldots)$, neither
$(R_0,\bar R_1,\ldots)$ nor
$(k,k\otimes_{R_0}\bar R_1,\ldots)$ is smooth in general. More
precisely:

\underline{Corollary \SGARefTarget{sga7.sga7i.evi.n2-20}{2.20}.} Let $\underline A$ be a homogeneous
cofibered groupoid over $\underline C_\Lambda$, prorepresented by a
normalized smooth cogroupoid $(R_0,R_1,\ldots)$ via
$(\xi,\varphi)$.

\begin{enumerate}[label=(\alph*)]
\item $(R_0,\bar R_1,\ldots)$ and
$(k,k\otimes_{R_0}\bar R_1,\ldots)$ prorepresent
$\operatorname{Aut}_\xi$ and $\operatorname{Aut}^{o}$, respectively.

\item The following statements are equivalent:
\begin{enumerate}[label=(\roman*)]
\item the formal group scheme
$\underline h_{(R_0,\bar R_1,\ldots)}$ is formally smooth over $R_0$;
\item for every surjective arrow $a'\longrightarrow a$ in
$\underline A$, the map
\[
 \operatorname{Aut}_{P(a')}(a')\longrightarrow
 \operatorname{Aut}_{P(a)}(a)
\]
is surjective;
\item $R_0$ prorepresents the functor $[\underline A]$ over
$\underline C_\Lambda$.
\end{enumerate}
\end{enumerate}

%%%%%%%%%%  scan idx 81  |  folio 70  %%%%%%%%%%

\underline{Proof.} We have already noted (a). For (b), the
equivalence (i)$\Leftrightarrow$(ii) is clear, whereas
(ii)$\Leftrightarrow$(iii) follows from \SGARefLink{sga7.sga7i.evi.n2-7}{2.7} and \SGARefLink{sga7.sga7i.evi.n2-15}{2.15}, since
\[
 \widetilde\xi:\underline h_{R_0}\longrightarrow[\underline A]
\]
is minimally smooth. \hfill$\square$

\underline{Remark \SGARefTarget{sga7.sga7i.evi.n2-21}{2.21}.} Let $\underline A$ be a cofibered groupoid
over $\underline C_\Lambda$, and let
$\underline A|\underline C_k$ be the pullback of
$p:\underline A\longrightarrow\underline C_\Lambda$ along
$\underline C_k\lhook\joinrel\longrightarrow\underline C_\Lambda$.
If $\underline A$ is prorepresented by $(R_0,R_1,\ldots)$ via
$(\xi,\varphi)$, then $\underline A|\underline C_k$ is
prorepresented by
\[
 (k\otimes_\Lambda R_0,k\otimes_\Lambda R_1,\ldots)
\]
via $(k\otimes_\Lambda\xi,k\otimes_\Lambda\varphi)$. Here
$k\otimes_\Lambda\xi$ is the direct image of $\xi$ under
$R_0\longrightarrow k\otimes_\Lambda R_0$, and
$k\otimes_\Lambda\varphi$ is the image of $\varphi$ under
\[
 \operatorname{Isom}_{\underline A(R_0)}
   \bigl(\xi(s_1),\xi(s_2)\bigr)
 \longrightarrow
 \operatorname{Isom}_{\underline A(k\otimes_\Lambda R_1)}
   \bigl(k\otimes_\Lambda\xi(s_1),
         k\otimes_\Lambda\xi(s_2)\bigr).
\]
If $(R_0,R_1,\ldots)$ is normalized, then so is
$(k\otimes_\Lambda R_0,k\otimes_\Lambda R_1,\ldots)$; moreover, the
relative dimension of $R_1$ over $R_0$ equals that of
$k\otimes_\Lambda R_1$ over $k\otimes_\Lambda R_0$, since $R_1$ is
smooth over $R_0$.

\underline{Remark \SGARefTarget{sga7.sga7i.evi.n2-22}{2.22}.} Suppose that $\underline A$ is
prorepresented by smooth cogroupoids $(R_0,R_1,\ldots)$ and
$(R'_0,R'_1,\ldots)$, with $(R_0,R_1,\ldots)$ normalized. We obtain a
homomorphism
\[
 (i_0,i_1):(R_0,R_1,\ldots)\longrightarrow(R'_0,R'_1,\ldots).
\]
By \SGARefLink{sga7.sga7i.evi.n2-10}{2.10}(b), $R'_0$ is a formal power-series ring over $R_0$, and
$(i_0,i_1)$ induces an isomorphism
\[
 R'_0\otimes_{R_0}(R_0,R_1,\ldots)
 \xrightarrow{\ \sim\ }(R'_0,R'_1,\ldots).
\]

% END INCLUDED SGA7I_English_current_buildable_source_20260823/components/26_expose_VI_2_18_through_2_22.tex

% BEGIN INCLUDED SGA7I_English_current_buildable_source_20260823/components/27_expose_VI_section_3_opening_through_3_1_b.tex
%%%%%%%%%%  scan idx 81  |  folio 70  %%%%%%%%%%

\section*{\SGARefTarget{sga7.sga7i.evi.n3}{3}. The coherent sheaves $\underline D_X^i$}

One important functor which is usually not prorepresentable, but
nevertheless admits a formal versal object, is the deformation
functor. With this in view, we recall some basic facts about
extensions of commutative rings and their obstruction theory. This
section contains nothing new, except perhaps for its systematic use
of Grothendieck cohomology on a site. Its purpose is to make this
expos\'e as self-contained as possible; in fact, we prove somewhat
more than will be needed later. The basic results are stated in
Theorem \SGARefLink{sga7.sga7i.evi.n3-12}{3.12}. Throughout, all rings are commutative and have a unit.

%%%%%%%%%%  scan idx 82  |  folio 71  %%%%%%%%%%

Given a ring $R$, let $\underline C_R$ denote the category of
$R$-algebras. For every $A\in\operatorname{Ob}(\underline C_R)$, put
\[
 \operatorname{Cov}(A)=
 \{\text{surjective arrows }\varphi:A'\longrightarrow A
   \text{ in }\underline C_R\}.
\]
This defines a pretopology, and hence a site, on $\underline C_R$.
For every $A\in\operatorname{Ob}(\underline C_R)$, let
$\underline C_{A|R}$ be the category of arrows in $\underline C_R$
with target $A$. We endow $\underline C_{A|R}$ with the topology
induced by the forgetful functor
\[
 j_A:\underline C_{A|R}\longrightarrow\underline C_R.
\]
The category $\underline C_{A|R}$ has a final object
$(1_A:A\longrightarrow A)$ and an initial object
$(i_A:R\longrightarrow A)$, and it admits both fibre products and
coproducts. When there is no risk of confusion, we write simply $B$
for an object $(B\longrightarrow A)$ of $\underline C_{A|R}$.

Let $\mathscr C$ be an abelian category. A $\mathscr C$-valued sheaf
on $\underline C_{A|R}$ is, by definition, a functor
\[
 F:\underline C_{A|R}^{\mathrm{op}}\longrightarrow\mathscr C
\]
such that, for every surjective arrow $B'\longrightarrow B$ in
$\underline C_{A|R}$, the sequence
\[
 0\longrightarrow F(B)\longrightarrow F(B')
 \mathrel{\substack{\longrightarrow\\[-.6ex]\longrightarrow}}
 F(B'\mathbin{\times_B}B')
\]
is exact. The \underline{additive} presheaves are those which
preserve coproducts; thus
\[
 F(B_1\otimes_R B_2)=F(B_1)\oplus F(B_2).
\]
For example, an $A$-module $M$ defines the functor
\[
 \operatorname{Der}_R(-,M):
 \underline C_{A|R}^{\mathrm{op}}\longrightarrow(A\text{-mod}),
\]
which sends $B$ to the $R$-linear derivations
$\operatorname{Der}_R(B,M)$, with $M$ regarded as a $B$-module via
the structural map $B\longrightarrow A$. This is plainly an
additive sheaf on $\underline C_{A|R}$. Here $(A\text{-mod})$ denotes
the category of $A$-modules.

A presheaf $F$ on $\underline C_{A|R}$ with values in
$(A\text{-mod})$ is said to be of \underline{finite type} if $F(B)$
is an $A$-module of finite type whenever $B$ is an $R$-algebra of
finite type. For example, if $A$ is noetherian and $M$ is an
$A$-module of finite type, then $\operatorname{Der}_R(-,M)$ is an
additive sheaf of finite type on $\underline C_{A|R}$.

We record a little more notation and a few basic facts from
homological algebra.

\underline{\SGARefTarget{sga7.sga7i.evi.n3-1.item-a}{\SGARefAliasTarget{sga7.sga7i.evi.n3-1}3.1}(a).} Throughout the remainder of this section,
$\widetilde{\underline C}_{A|R}$ and
$\overset{\vee}{\underline C}_{A|R}$ denote, respectively, the
categories of sheaves and presheaves with values in
$(A\text{-mod})$. There is a left-exact inclusion functor
\[
 i_{A|R}:\widetilde{\underline C}_{A|R}
 \longrightarrow\overset{\vee}{\underline C}_{A|R}
\]
and an exact associated-sheaf functor
\[
 G_{A|R}:\overset{\vee}{\underline C}_{A|R}
 \longrightarrow\widetilde{\underline C}_{A|R},
\]
left adjoint to $i_{A|R}$.

%%%%%%%%%%  scan idx 83  |  folio 72  %%%%%%%%%%

An object $B\in\operatorname{Ob}(\underline C_{A|R})$ induces a
continuous functor
$j:\underline C_{B|R}\longrightarrow\underline C_{A|R}$ and hence
a commutative diagram
\begin{equation}\tag{*}
\begin{tikzcd}[column sep=huge,row sep=huge]
 \widetilde{\underline C}_{A|R}
   \arrow[r,"\widetilde j"] \arrow[d,"i_{A|R}"']
 & \widetilde{\underline C}_{B|R} \arrow[d] \\
 \overset{\vee}{\underline C}_{A|R}
   \arrow[r,"j^{\vee}"']
 & \overset{\vee}{\underline C}_{B|R}.
\end{tikzcd}
\end{equation}
The restriction functor $\widetilde j$, as well as $j^{\vee}$, is
exact, because the site $\underline C_{B|R}$ is simply the category
of arrows in $\underline C_{A|R}$ with target $B$, endowed with the
induced topology.

A ring homomorphism $S\longrightarrow R$ induces a continuous
functor
$\beta:\underline C_{A|R}\longrightarrow\underline C_{A|S}$ and
hence a commutative diagram
\begin{equation}\tag{**}
\begin{tikzcd}[column sep=huge,row sep=huge]
 \widetilde{\underline C}_{A|S}
   \arrow[r,"\widetilde\beta"] \arrow[d,"i_{A|S}"']
 & \widetilde{\underline C}_{A|R} \arrow[d] \\
 \overset{\vee}{\underline C}_{A|S}
   \arrow[r,"\beta^{\vee}"']
 & \overset{\vee}{\underline C}_{A|R}.
\end{tikzcd}
\end{equation}
Notice that $\widetilde\beta$ need not be exact.

\underline{\SGARefTarget{sga7.sga7i.evi.n3-1.item-b}{3.1}(b).} The right-derived functor of the inclusion
\[
 i_{A|R}:\widetilde{\underline C}_{A|R}
 \longrightarrow\overset{\vee}{\underline C}_{A|R}
\]
will be denoted by $\underline H^{\bullet}_{A|R}$. For every sheaf
$F$ on $\underline C_{A|R}$, put
\[
 H^{\bullet}(A|R,F)=
 [\underline H^{\bullet}_{A|R}(F)](A).
\]

For $B'\longrightarrow B$ in $\operatorname{Cov}(B)$, define
\[
 H^0(\{B'\longrightarrow B\},F)=
 \ker\!\left(
 F(B')\mathrel{\substack{\longrightarrow\\[-.6ex]\longrightarrow}}
 F(B'\mathbin{\times_B}B')
 \right).
\]
We denote the right-derived functors of
$H^0(\{B'\longrightarrow B\},-)$ by
$H^{\bullet}(\{B'\longrightarrow B\},-)$. These are the cohomology
groups of the \v{C}ech semisimplicial module

%%%%%%%%%%  scan idx 84  |  folio 73  %%%%%%%%%%

\[
\begin{tikzcd}[column sep=large]
 F(B')
   \arrow[r,shift left=2] \arrow[r,shift right=2]
 & F(B'\times_BB')
   \arrow[r,shift left=3] \arrow[r] \arrow[r,shift right=3]
 & F(B'\times_BB'\times_BB')
   \arrow[r,shift left=4.5] \arrow[r,shift left=1.5]
   \arrow[r,shift right=1.5] \arrow[r,shift right=4.5]
 & \cdots .
\end{tikzcd}
\]

The right-derived functor of
\[
 \check{\underline H}^{,0}_{A|R}:
 \overset{\vee}{\underline C}_{A|R}
 \longrightarrow\overset{\vee}{\underline C}_{A|R}
\]
will be denoted by $\check{\underline H}^{\,\bullet}_{A|R}$, where
\[
 [\check{\underline H}^{,0}_{A|R}(F)](B)=
 \varinjlim_{\{B'\to B\}\in\operatorname{Cov}(B)}
 H^0(\{B'\longrightarrow B\},F).
\]
For every presheaf $F$ on $\underline C_{A|R}$, put
\[
 \check H^{\bullet}(A|R,F)=
 [\check{\underline H}^{\,\bullet}_{A|R}(F)](A)=
 \varinjlim_{\{A'\to A\}\in\operatorname{Cov}(A)}
 H^{\bullet}(\{A'\longrightarrow A\},F).
\]
For every sheaf $F$ on $\underline C_{A|R}$ there is the usual
spectral sequence
\[
 \check H^p\bigl(A|R,\underline H^q_{A|R}(F)\bigr)
 \Longrightarrow H^{p+q}(A|R,F).
\]
Here
$\check H^0(A|R,\underline H^q_{A|R}(F))=0$ for every $q>0$.
Consequently,
\[
 \check H^1(A|R,F)=H^1(A|R,F),
\]
and the sequence
\[
 0\longrightarrow\check H^2(A|R,F)
 \longrightarrow H^2(A|R,F)
 \longrightarrow
 \check H^1\bigl(A|R,\underline H^1_{A|R}(F)\bigr)
 \longrightarrow\check H^3(A|R,F)
\]
is exact.

% END INCLUDED SGA7I_English_current_buildable_source_20260823/components/27_expose_VI_section_3_opening_through_3_1_b.tex

% BEGIN INCLUDED SGA7I_English_current_buildable_source_20260823/components/28_expose_VI_3_1_c_through_3_4.tex
\underline{\SGARefTarget{sga7.sga7i.evi.n3-1.item-c}{3.1}(c).} If $B$ is an object of
$\underline C_{A|R}$, diagram $(*)$ in \SGARefLink{sga7.sga7i.evi.n3-1.item-a}{3.1}(a), together with the
exactness of $\widetilde j$ and $j^{\vee}$, gives a canonical
isomorphism
\[
 H^{\bullet}(B|R,\widetilde jF)=
 [\underline H^{\bullet}_{A|R}(F)](B).
\]
On the other hand, a ring homomorphism $S\longrightarrow R$ induces
diagram $(**)$ in \SGARefLink{sga7.sga7i.evi.n3-1.item-a}{3.1}(a). The functor $\widetilde\beta$ need not be
exact; however, $\beta$ preserves fibre products, and we obtain a
spectral sequence
\[
 \underline H^p_{A|R}(R^q\widetilde\beta F)
 \Longrightarrow
 \beta^{\vee}\underline H^{p+q}_{A|S}(F),
\]
or, on evaluating at $A$,
\[
 H^p(A|R,R^q\widetilde\beta F)
 \Longrightarrow H^{p+q}(A|S,F).
\]
Since there is a canonical isomorphism
\[
 \widetilde\beta=
 G_{A|R}\circ\beta^{\vee}\circ i_{A|S}
\]
and the functor
$\beta^{\#}=G_{A|R}\circ\beta^{\vee}$ is exact, we have

%%%%%%%%%%  scan idx 85  |  folio 74  %%%%%%%%%%

\[
 R^{\bullet}\widetilde\beta F=
 \beta^{\#}\underline H^{\bullet}_{A|S}(F).
\]
Thus the preceding spectral sequence may be written
\[
 H^p\bigl(A|R,\beta^{\#}\underline H^q_{A|S}(F)\bigr)
 \Longrightarrow H^{p+q}(A|S,F).
\]

\underline{Lemma \SGARefTarget{sga7.sga7i.evi.n3-2}{3.2}.} Let
$B\in\operatorname{Ob}(\underline C_{A|R})$, and choose a
surjective arrow $P\longrightarrow B$ in $\underline C_{B|R}$,
where $P$ is a polynomial algebra over $R$. Then, for every
$F\in\operatorname{Ob}(\overset{\vee}{\underline C}_{A|R})$, the
natural map
\[
 H^{\bullet}(\{P\longrightarrow B\},F)
 \longrightarrow[\check{\underline H}^{\,\bullet}(F)](B)
\]
is an isomorphism.

\underline{Proof.} It suffices to prove that the natural map
\[
 H^0(\{P\longrightarrow B\},F)
 \longrightarrow[\check{\underline H}^{,0}(F)](B)
 =\varinjlim_{\{B'\to B\}\in\operatorname{Cov}(B)}
   H^0(\{B'\longrightarrow B\},F)
\]
is an isomorphism. This is clear because, for every
$\{B'\longrightarrow B\}\in\operatorname{Cov}(B)$, there is a map
of covers
\[
 \{P\longrightarrow B\}\longrightarrow
 \{B'\longrightarrow B\}.
\]
\hfill$\square$

\underline{Corollary \SGARefTarget{sga7.sga7i.evi.n3-3}{3.3}.}

\begin{enumerate}[label=(\alph*)]
\item Suppose that $A$ is noetherian. If
$F\in\operatorname{Ob}(\widetilde{\underline C}_{A|R})$ is of
finite type, then $\underline H^n(F)$ is of finite type for every
$n$.

\item If $P$ is a polynomial algebra over $R$, then
\[
 H^n(P|R,F)=0
\]
for every $n>0$ and every
$F\in\operatorname{Ob}(\widetilde{\underline C}_{P|R})$.

\item A sequence
\[
 0\longrightarrow F'\longrightarrow F\longrightarrow F''
 \longrightarrow0
\]
in $\widetilde{\underline C}_{A|R}$ is exact if and only if
\[
 0\longrightarrow F'(P)\longrightarrow F(P)
 \longrightarrow F''(P)\longrightarrow0
\]
is exact for every polynomial $R$-algebra
$P\in\operatorname{Ob}(\underline C_{A|R})$.
\end{enumerate}

\underline{Proof.}

(a) In view of the spectral sequence
\[
 \check H^p\bigl(B|R,\underline H^q(F)\bigr)
 \Longrightarrow H^{p+q}(B|R,F)
\]
and the equality
$\check H^0(\underline H^q(F))=0$ for $q>0$, it is enough to show
that $\check{\underline H}^{,n}(G)$ is of finite type for every
$n$ and every presheaf $G$ of finite type. By Lemma \SGARefLink{sga7.sga7i.evi.n3-2}{3.2},
\[
 \check H^{\bullet}(B|R,G)=
 H^{\bullet}(\{P\longrightarrow B\},G),
\]
where $P$ is a polynomial algebra over $R$ and
$\{P\longrightarrow B\}\in\operatorname{Cov}(B)$. The latter is
the cohomology of the complex

%%%%%%%%%%  scan idx 86  |  folio 75  %%%%%%%%%%

\[
 \cdots
 G(P\times_BP\times_BP)
 \mathrel{\substack{\longleftarrow\\[-.6ex]\longleftarrow\\[-.6ex]\longleftarrow}}
 G(P\times_BP)
 \mathrel{\substack{\longleftarrow\\[-.6ex]\longleftarrow}}
 G(P).
\]
If $B$ is an $R$-algebra of finite type, $P$ may be chosen of
finite type over $R$. All the iterated fibre products
$P\times_B\cdots\times_BP$ are then $R$-algebras of finite type.
If $G$ is of finite type, their images under $G$ are $A$-modules of
finite type. Hence
$H^{\bullet}(\{P\longrightarrow B\},G)$ is an $A$-module of finite
type because $A$ is noetherian.

(b) Let $P$ be a polynomial algebra over $R$. Every surjective
$R$-algebra map $B\longrightarrow P$ admits a section, so
\[
 H^i(\{B\longrightarrow P\},G)=0\qquad(i>0).
\]
It follows that $\check H^i(P|R,G)=0$ for every $i>0$ and every
presheaf $G$ on $\underline C_{P|R}$. Consequently the spectral
sequence
\[
 \check H^p\bigl(P|R,\underline H^q_{P|R}(F)\bigr)
 \Longrightarrow H^{p+q}(P|R,F)
\]
degenerates, and
\[
 \check H^0\bigl(P|R,\underline H^n(F)\bigr)
 \simeq H^n(P|R,F)
\]
for every $n$. For a sheaf $F$, however,
$\check H^0(P|R,\underline H^n(F))=0$ whenever $n>0$. Thus
$H^n(P|R,F)=0$ for every $n>0$.

(c) If
$0\longrightarrow F'\longrightarrow F\longrightarrow F''
\longrightarrow0$ is exact in
$\widetilde{\underline C}_{A|R}$, then the sequence
\[
 0\longrightarrow F'(B)\longrightarrow F(B)
 \longrightarrow F''(B)
 \longrightarrow H^1(B|R,F')\longrightarrow H^1(B|R,F)
 \longrightarrow\cdots
\]
is exact. If $P\in\operatorname{Ob}(\underline C_{A|R})$ is a
polynomial algebra over $R$, then $H^1(P|R,F')=0$ by (b), and hence
\[
 0\longrightarrow F'(P)\longrightarrow F(P)
 \longrightarrow F''(P)\longrightarrow0
\]
is exact.

Conversely, let $\underline P_{A|R}$ be the full subcategory of
$\underline C_{A|R}$ with
\[
 \operatorname{Ob}(\underline P_{A|R})=
 \{P\in\operatorname{Ob}(\underline C_{A|R})\mid
   P\text{ is a polynomial algebra over }R\},
\]
endowed with the induced topology, and let
$p:\underline P_{A|R}\longrightarrow\underline C_{A|R}$ be the
inclusion. Every object of $\underline C_{A|R}$ is covered by an
object of $\underline P_{A|R}$; the comparison lemma therefore
shows that
\[
 \widetilde p:\widetilde{\underline C}_{A|R}
 \longrightarrow\widetilde{\underline P}_{A|R}
\]
is an equivalence of categories. If the sequence obtained by
evaluating $F'\longrightarrow F\longrightarrow F''$ at every
$P\in\operatorname{Ob}(\underline P_{A|R})$ is exact, then
\[
 0\longrightarrow\widetilde pF'\longrightarrow\widetilde pF
 \longrightarrow\widetilde pF''\longrightarrow0
\]
is exact in $\widetilde{\underline P}_{A|R}$, and therefore the
original sequence is exact in
$\widetilde{\underline C}_{A|R}$. \hfill$\square$

%%%%%%%%%%  scan idx 87  |  folio 76  %%%%%%%%%%

\underline{Remark \SGARefTarget{sga7.sga7i.evi.n3-4}{3.4}.} Every surjective arrow in
$\underline P_{A|R}$ admits a section. It follows at once that every
presheaf on $\underline P_{A|R}$ is a sheaf; that is,
\[
 \widetilde{\underline P}_{A|R}=
 \overset{\vee}{\underline P}_{A|R}.
\]
Consequently, the restriction functor
\[
 \widetilde{\underline C}_{A|R}
 \longrightarrow\overset{\vee}{\underline P}_{A|R}
\]
is an equivalence of categories.

% END INCLUDED SGA7I_English_current_buildable_source_20260823/components/28_expose_VI_3_1_c_through_3_4.tex

% BEGIN INCLUDED SGA7I_English_current_buildable_source_20260823/components/29_expose_VI_3_5_through_3_6.tex
\underline{Lemma \SGARefTarget{sga7.sga7i.evi.n3-5}{3.5}.} Consider a diagram of $R$-algebras
\[
\begin{tikzcd}[column sep=large,row sep=large]
 & C \arrow[d]\\
 B' \arrow[r,"\text{surjective}"'] & B
\end{tikzcd}
\]
and let $S$ be an $R$-algebra.

\begin{enumerate}[label=(\roman*)]
\item If $\operatorname{Tor}_1^R(B,S)=0$, then the canonical map
\[
 (B'\times_BC)\otimes_RS
 \longrightarrow
 (B'\otimes_RS)\times_{B\otimes_RS}(C\otimes_RS)
\]
is an isomorphism.

\item Suppose that $B'$ is flat over $R$. If
\[
 \operatorname{Tor}_i^R(C,S)=0\quad(0<i<n)
 \qquad\text{and}\qquad
 \operatorname{Tor}_i^R(B,S)=0\quad(0<i\leq n),
\]
then
\[
 \operatorname{Tor}_i^R(B'\times_BC,S)=0\qquad(0<i<n).
\]
\end{enumerate}

\underline{Proof.} Let
$0\longrightarrow J\longrightarrow B'\longrightarrow B
\longrightarrow0$ be exact.

(i) If $\operatorname{Tor}_1^R(B,S)=0$, then
\[
 0\longrightarrow J\otimes_RS\longrightarrow B'\otimes_RS
 \longrightarrow B\otimes_RS\longrightarrow0
\]
is exact, and hence so is
\[
 0\longrightarrow J\otimes_RS
 \longrightarrow
 (B'\otimes_RS)\times_{B\otimes_RS}(C\otimes_RS)
 \longrightarrow C\otimes_RS\longrightarrow0.
\]
We therefore obtain a commutative diagram with exact rows:
\begin{center}
\scriptsize
\begin{tikzcd}[column sep=small,row sep=large]
 & J\otimes_RS \arrow[r] \arrow[d,equal]
 & (B'\times_BC)\otimes_RS \arrow[r] \arrow[d]
 & C\otimes_RS \arrow[r] \arrow[d,equal]
 & 0\\
 0 \arrow[r] & J\otimes_RS \arrow[r]
 & (B'\otimes_RS)\times_{B\otimes_RS}(C\otimes_RS) \arrow[r]
 & C\otimes_RS \arrow[r] & 0.
\end{tikzcd}
\end{center}

%%%%%%%%%%  scan idx 88  |  folio 77  %%%%%%%%%%

It follows that
\[
 (B'\times_BC)\otimes_RS
 \longrightarrow
 (B'\otimes_RS)\times_{B\otimes_RS}(C\otimes_RS)
\]
is an isomorphism.

(ii) Since $B'$ is flat over $R$ and
$\operatorname{Tor}_i^R(B,S)=0$ for $0<i\leq n$, the long exact
Tor sequence gives
\[
 \operatorname{Tor}_i^R(J,S)=0\qquad(0<i<n).
\]
Applying the same sequence to
\[
 0\longrightarrow J\longrightarrow B'\times_BC
 \longrightarrow C\longrightarrow0
\]
and using the hypothesis on $C$ gives
\[
 \operatorname{Tor}_i^R(B'\times_BC,S)=0\qquad(0<i<n).
\]
\hfill$\square$

\underline{Corollary \SGARefTarget{sga7.sga7i.evi.n3-6}{3.6}.}

\begin{enumerate}[label=(\alph*)]
\item Let $A$ and $S$ be $R$-algebras. If
$\operatorname{Tor}_1^R(A,S)=0$, then for every
$G\in\operatorname{Ob}
(\overset{\vee}{\underline C}_{A\otimes_RS|S})$ the canonical map
\[
 \check H^{\bullet}(A\otimes_RS|S,G)
 \longrightarrow\check H^{\bullet}(A|R,\check SG)
\]
is an isomorphism. Here
\[
 S:\underline C_{A|R}\longrightarrow
 \underline C_{A\otimes_RS|S},
 \qquad X\longmapsto X\otimes_RS.
\]
In particular, if $S$ is flat over $R$, then
\[
 \check S\,\check{\underline H}^{\,\bullet}(G)
 \longrightarrow
 \check{\underline H}^{\,\bullet}(\check SG)
\]
is an isomorphism for every such $G$.

\item Let $P\longrightarrow B$ be a surjective map of $R$-algebras,
where $P$ is a polynomial algebra over $R$. If
\[
 \operatorname{Tor}_i^R(B,S)=0\qquad(0<i\leq n),
\]
then, for every iterated fibre product of copies of $P$ over $B$,
\[
 \operatorname{Tor}_i^R
 (P\times_B P\times_B\cdots\times_B P,S)=0
 \qquad(0<i<n).
\]
\end{enumerate}

\underline{Proof.}

(a) Given $B\in\operatorname{Ob}(\underline C_{A|R})$, choose a
cover $\{P\longrightarrow B\}\in\operatorname{Cov}(B)$ with $P$ a
polynomial algebra over $R$. If
$\operatorname{Tor}_1^R(B,S)=0$, Lemma \SGARefLink{sga7.sga7i.evi.n3-5}{3.5}(i) shows that the natural
map
\[
 \left(\overbrace{P\times_B\cdots\times_BP}^{n}\right)\otimes_RS
 \longrightarrow
 \overbrace{(P\otimes_RS)\times_{B\otimes_RS}\cdots
 \times_{B\otimes_RS}(P\otimes_RS)}^{n}
\]
is an isomorphism for every $n$. Consequently,
\begin{align*}
 \check H^{\bullet}(B|R,\check SG)
 &=H^{\bullet}(\{P\longrightarrow B\},\check SG)\\
 &=H^{\bullet}(\{P\otimes_RS\longrightarrow B\otimes_RS\},G)\\
 &=\check H^{\bullet}(B\otimes_RS|S,G).
\end{align*}
If $S$ is flat over $R$, then
$\operatorname{Tor}_1^R(X,S)=0$ for every
$X\in\operatorname{Ob}(\underline C_{A|R})$. The preceding
isomorphism is therefore natural in $X$, and gives
\[
 \check S\,\check{\underline H}^{\,\bullet}(G)
 \xrightarrow{\ \sim\ }
 \check{\underline H}^{\,\bullet}(\check SG).
\]

(b) Put
\[
 P^{(m)}=\overbrace{P\times_B\cdots\times_BP}^{m+1}.
\]
For $m=0$ there is nothing to prove, since $P$ is flat over $R$.
Suppose inductively that
$\operatorname{Tor}_i^R(P^{(m-1)},S)=0$ for $0<i<n$. Since
$P\longrightarrow B$ is surjective and
$P^{(m)}=P\times_BP^{(m-1)}$, Lemma \SGARefLink{sga7.sga7i.evi.n3-5}{3.5}(ii) gives
\[
 \operatorname{Tor}_i^R(P^{(m)},S)=0\qquad(0<i<n).
\]
\hfill$\square$

% END INCLUDED SGA7I_English_current_buildable_source_20260823/components/29_expose_VI_3_5_through_3_6.tex

% BEGIN INCLUDED SGA7I_English_current_buildable_source_20260823/components/30_expose_VI_proposition_3_7.tex
%%%%%%%%%%  scan idx 89  |  folio 78  %%%%%%%%%%

\underline{Proposition \SGARefTarget{sga7.sga7i.evi.n3-7}{3.7}.}

\begin{enumerate}[label=(\alph*)]
\item Let $A$ and $S$ be $R$-algebras. If
\[
 \operatorname{Tor}_i^R(A,S)=0\qquad(0<i\leq n),
\]
then, for every sheaf $F$ on
$\underline C_{A\otimes_RS|S}$, the canonical map
\[
 H^i(A\otimes_RS|S,F)
 \longrightarrow H^i(A|R,\widetilde SF)
\]
is an isomorphism for every $i\leq n$. Here
\[
 S:\underline C_{A|R}\longrightarrow
 \underline C_{A\otimes_RS|S},
 \qquad X\longmapsto X\otimes_RS.
\]

\item Let $P$ be a polynomial algebra over $R$. Then, for every
additive sheaf $F$ on $\underline C_{A\otimes_RP|R}$, the
canonical map
\[
 H^n(A\otimes_RP|R,F)\longrightarrow H^n(A|R,F)
\]
is an isomorphism for every $n>0$.

\item Let $S\longrightarrow R\longrightarrow A$ be ring
homomorphisms. Then, for every additive sheaf $F$ on
$\underline C_{A|S}$, there is an exact sequence
\begin{multline*}
0\longrightarrow H^0(A|R,F_R)\longrightarrow H^0(A|S,F)
\longrightarrow H^0(R|S,F)\longrightarrow H^1(A|R,F_R)\\
\longrightarrow H^1(A|S,F)\longrightarrow H^1(R|S,F)
\longrightarrow\cdots\longrightarrow H^n(A|R,F_R)\\
\longrightarrow H^n(A|S,F)\longrightarrow H^n(R|S,F)
\longrightarrow H^{n+1}(A|R,F_R)\longrightarrow\cdots,
\end{multline*}
where $F_R$ is the sheaf on $\underline C_{A|R}$ defined by
\[
 F_R(X)=\ker\bigl(F(X)\longrightarrow F(R)\bigr)
\]
for every $X\in\operatorname{Ob}(\underline C_{A|R})$, via the
forgetful functor
$\underline C_{A|R}\longrightarrow\underline C_{A|S}$.
\end{enumerate}

\underline{Proof.}

(a) Consider the continuous functor
\[
 S:\underline C_{A|R}\longrightarrow
 \underline C_{A\otimes_RS|S},
 \qquad S(X)=X\otimes_RS.
\]
Since $S$ sends polynomial algebras over $R$ to polynomial
algebras over $S$, the functor
\[
 \widetilde S:\widetilde{\underline C}_{A\otimes_RS|S}
 \longrightarrow\widetilde{\underline C}_{A|R}
\]
is exact by \SGARefLink{sga7.sga7i.evi.n3-3}{3.3}(c). We therefore obtain the canonical map
\[
 \check S\,\underline H^{\bullet}_{A\otimes_RS|S}(F)
 \longrightarrow
 \underline H^{\bullet}_{A|R}(\widetilde SF)
\]
and the following commutative diagram of spectral sequences:
\begin{center}
\small
\begin{tikzcd}[column sep=huge,row sep=large]
E_S^{p,q}(S(X))=
 \check H^p\bigl(X\otimes_RS|S,
 \underline H^q_{A\otimes_RS|S}(F)\bigr)
 \arrow[r,Rightarrow] \arrow[d]
& H^{p+q}(X\otimes_RS|S,F) \arrow[dd]\\
\check H^p\bigl(X|R,
 \check S\,\underline H^q_{A\otimes_RS|S}(F)\bigr)
 \arrow[d] & \\
E_R^{p,q}(X)=
 \check H^p\bigl(X|R,
 \underline H^q_{A|R}(\widetilde SF)\bigr)
 \arrow[r,Rightarrow]
& H^{p+q}(X|R,\widetilde SF).
\end{tikzcd}
\end{center}

%%%%%%%%%%  scan idx 90  |  folio 79  %%%%%%%%%%

By \SGARefLink{sga7.sga7i.evi.n3-6}{3.6}(a), the map
\[
 E_S^{p,q}(S(X))\longrightarrow
 \check H^p\bigl(X|R,
 \check S\,\underline H^q_{A\otimes_RS|S}(F)\bigr)
\]
is an isomorphism whenever
$\operatorname{Tor}_1^R(X,S)=0$. Suppose inductively that
\[
 [\check S\,\underline H^q_{A\otimes_RS|S}(F)](X)
 \longrightarrow
 [\underline H^q_{A|R}(\widetilde SF)](X)
\]
is an isomorphism for every $q<n$ and every
$X\in\operatorname{Ob}(\underline C_{A|R})$ satisfying
\[
 \operatorname{Tor}_i^R(X,S)=0\qquad(0<i\leq n).
\]
It follows from \SGARefLink{sga7.sga7i.evi.n3-6}{3.6}(b) that
$E_S^{p,q}(S(X))\longrightarrow E_R^{p,q}(X)$ is an isomorphism
for every $q<n$ and every $p$, under the same Tor hypothesis.
Moreover,
$E_S^{0,n}(S(X))\longrightarrow E_R^{0,n}(X)$ is an isomorphism,
since both sides are zero. Hence
\[
 H^n(X\otimes_RS|S,F)
 \longrightarrow H^n(X|R,\widetilde SF)
\]
is an isomorphism whenever
$\operatorname{Tor}_i^R(X,S)=0$ for $0<i\leq n$. Equivalently,
\[
 [\check S\,\underline H^i_{A\otimes_RS|S}(F)](X)
 \longrightarrow
 [\underline H^i_{A|R}(\widetilde SF)](X)
\]
is an isomorphism for every $i\leq n$ under that hypothesis.

(b) Let
\[
 f:\underline C_{A|R}\longrightarrow
 \underline C_{A\otimes_RP|R},
 \qquad f(X)=X\otimes_RP.
\]
Choose a surjective arrow $P'\longrightarrow X$ in
$\underline C_{A|R}$, with $P'$ polynomial over $R$. Then
\[
 P'\otimes_RP\longrightarrow X\otimes_RP
\]
is a surjective arrow in $\underline C_{A\otimes_RP|R}$, and
$P'\otimes_RP$ is again polynomial over $R$. Since $P$ is flat
over $R$, \SGARefLink{sga7.sga7i.evi.n3-5}{3.5}(i) gives
\begin{align*}
 \check H^{\bullet}(X|R,\check fG)
 &=H^{\bullet}(\{P'\longrightarrow X\},\check fG)\\
 &=H^{\bullet}(\{P'\otimes_RP\longrightarrow X\otimes_RP\},G)\\
 &=\check H^{\bullet}(X\otimes_RP|R,G)
\end{align*}
for every
$G\in\operatorname{Ob}(\overset{\vee}{\underline C}_{A\otimes_RP|R})$.
Thus $\widetilde f$ sends acyclic sheaves to acyclic sheaves. On
the other hand, $f$ sends polynomial algebras over $R$ to
polynomial algebras over $R$, so
\[
 \widetilde f:\widetilde{\underline C}_{A\otimes_RP|R}
 \longrightarrow\widetilde{\underline C}_{A|R}
\]
is exact by \SGARefLink{sga7.sga7i.evi.n3-3}{3.3}(c). It follows that
\[
 H^{\bullet}(A\otimes_RP|R,F)
 \longrightarrow H^{\bullet}(A|R,\widetilde fF)
\]
is an isomorphism for every
$F\in\operatorname{Ob}(\widetilde{\underline C}_{A\otimes_RP|R})$.

For each $X\in\operatorname{Ob}(\underline C_{A|R})$ there is a
canonical map
\[
 (\widetilde fF)(X)=F(X\otimes_RP)
 \longrightarrow F(X)\oplus F(P),
\]
and hence a canonical map
$\widetilde fF\longrightarrow F\oplus F(P)$, where $F(P)$ denotes
the constant sheaf and $F$ on the right denotes the restriction of
$F$ to $\underline C_{A|R}$. If $F$ is additive, this map is an
isomorphism. Therefore, for every $n>0$,
\[
 H^n(A|R,\widetilde fF)
 \xrightarrow{\ \sim\ }
 H^n(A|R,F\oplus F(P))=H^n(A|R,F),
\]
because the constant sheaf $F(P)$ is flasque. This proves (b).

%%%%%%%%%%  scan idx 91  |  folio 80  %%%%%%%%%%

(c) This argument was communicated by Rinehart. Consider the
forgetful functor
$\beta:\underline C_{A|R}\longrightarrow\underline C_{A|S}$.
There is a spectral sequence
\[
 E^{p,q}=H^p\bigl(A|R,
 \beta^{\#}\underline H^q_{A|S}(F)\bigr)
 \Longrightarrow H^{p+q}(A|S,F),
\]
where
\[
 \beta^{\#}=G_{A|R}\circ\beta^{\vee},
 \qquad
 G_{A|R}:\overset{\vee}{\underline C}_{A|R}
 \longrightarrow\widetilde{\underline C}_{A|R}
\]
is the associated-sheaf functor.

Suppose now that $F$ is additive. We claim that
$\beta^{\#}\underline H^q_{A|S}(F)$ is a constant sheaf for every
$q>0$. Indeed, if $P$ is a polynomial algebra over $R$, then
$P=R\otimes_SP'$ for a polynomial algebra $P'$ over $S$, and
\[
 [\beta^{\#}\underline H^q_{A|S}(F)](P)
 =H^q(P|S,F)
 =H^q(R\otimes_SP'|S,F)
 \longrightarrow H^q(R|S,F)
\]
is an isomorphism for $q>0$ by (b), and is surjective for $q=0$.
By \SGARefLink{sga7.sga7i.evi.n3-3}{3.3}(c), the canonical map
\[
 \beta^{\#}\underline H^q_{A|S}(F)
 \longrightarrow H^q(R|S,F)
\]
to the corresponding constant sheaf is therefore an isomorphism
for $q>0$ and a surjection for $q=0$. Hence $E^{p,q}=0$ whenever
$p,q>0$, and the spectral sequence yields the exact sequence
\[
 (*)\quad
 0\longrightarrow E^{1,0}\longrightarrow H^1
 \longrightarrow E^{0,1}\longrightarrow E^{2,0}
 \longrightarrow H^2\longrightarrow E^{0,2}
 \longrightarrow E^{3,0}\longrightarrow H^3\longrightarrow\cdots.
\]
Here
\[
 E^{0,q}=[\beta^{\#}\underline H^q_{A|S}(F)](A)
 \xrightarrow{\ \sim\ }H^q(R|S,F)\quad(q>0),
 \qquad
 E^{p,0}=H^p(A|R,\widetilde\beta F).
\]
On the other hand, the exact sequence
\[
 0\longrightarrow F_R\longrightarrow\widetilde\beta F
 \longrightarrow F(R)\longrightarrow0
\]
gives the exact diagram
\begin{center}
\scriptsize
\begin{tikzcd}[column sep=small,row sep=large]
0 \arrow[r]
& H^0(A|R,F_R) \arrow[r]
& H^0(A|R,\widetilde\beta F) \arrow[r] \arrow[d,equal]
& H^0(A|R,F(R)) \arrow[r] \arrow[d,equal]
& H^1(A|R,F_R) \arrow[r]
& H^1(A|R,\widetilde\beta F) \arrow[r]
& 0\\
&& H^0(A|S,F) & H^0(R|S,F) &&&
\end{tikzcd}
\end{center}
labelled $(**)$ in the source. Combining $(*)$ and $(**)$ gives
the desired exact sequence. \hfill$\square$

% END INCLUDED SGA7I_English_current_buildable_source_20260823/components/30_expose_VI_proposition_3_7.tex

% BEGIN INCLUDED SGA7I_English_current_buildable_source_20260823/components/31_expose_VI_corollaries_3_8_and_3_9.tex
%%%%%%%%%%  scan idx 92  |  folio 81  %%%%%%%%%%

\underline{Corollary \SGARefTarget{sga7.sga7i.evi.n3-8}{3.8}.}

\begin{enumerate}[label=(\alph*)]
\item Let $A$ and $B$ be $R$-algebras such that
\[
 \operatorname{Tor}_i^R(A,B)=0\qquad(0<i\leq n).
\]
Then, for every additive sheaf
$F\in\operatorname{Ob}(\widetilde{\underline C}_{A\otimes_RB|R})$,
the canonical map
\[
 H^i(A\otimes_RB|R,F)
 \longrightarrow H^i(A|R,F)\oplus H^i(B|R,F)
\]
is an isomorphism for every $i\leq n$.

\item Let $A$ be an $R$-algebra such that
$\operatorname{Tor}_i^R(A,A)=0$ for every $i>0$. If the map
$A\otimes_RA\longrightarrow A$ is an isomorphism, then, for every
additive sheaf
$F\in\operatorname{Ob}(\widetilde{\underline C}_{A|R})$, one has
\[
 H^i(A|R,F)=0
\]
for every $i$.
\end{enumerate}

\underline{Proof.}

(a) Let
$F\in\operatorname{Ob}(\widetilde{\underline C}_{A\otimes_RB|R})$
be additive. Applied to
$R\longrightarrow A\longrightarrow A\otimes_RB$, Proposition
\SGARefLink{sga7.sga7i.evi.n3-7}{3.7}(c) gives the exact sequence
\[
 (*)\quad\cdots\longrightarrow
 H^i(A\otimes_RB|A,F_A)
 \longrightarrow H^i(A\otimes_RB|R,F)
 \longrightarrow H^i(A|R,F)\longrightarrow\cdots.
\]
Consider the commutative diagram
\[
\begin{tikzcd}[column sep=huge,row sep=huge]
H^i(A\otimes_RB|A,F_A) \arrow[r] \arrow[d]
& H^i(A\otimes_RB|R,F) \arrow[d]\\
H^i(B|R,\widetilde A F_A) \arrow[r,"\approx"]
& H^i(B|R,F).
\end{tikzcd}
\]
Here
\[
 A:\underline C_{B|R}\longrightarrow
 \underline C_{A\otimes_RB|A},
 \qquad X\longmapsto A\otimes_RX,
\]
and the bottom map is induced by
$\widetilde A F_A\longrightarrow F$, which is an isomorphism
because $F$ is additive. By functoriality, the commutativity of the
diagram reduces to the case $i=0$.

If $\operatorname{Tor}_i^R(A,B)=0$ for $0<i\leq n$, then
Proposition \SGARefLink{sga7.sga7i.evi.n3-7}{3.7}(a) shows that the left vertical map is an
isomorphism for $i\leq n$. Consequently, both composites
\[
 H^i(A\otimes_RB|A,F_A)
 \longrightarrow H^i(A\otimes_RB|R,F)
 \longrightarrow H^i(B|R,F)
\]
and, by symmetry,

%%%%%%%%%%  scan idx 93  |  folio 82  %%%%%%%%%%

\[
 H^i(A\otimes_RB|B,F_B)
 \longrightarrow H^i(A\otimes_RB|R,F)
 \longrightarrow H^i(A|R,F)
\]
are isomorphisms for $i\leq n$. It follows that $(*)$ breaks into
split short exact sequences
\[
 0\longrightarrow H^i(A\otimes_RB|A,F_A)
 \longrightarrow H^i(A\otimes_RB|R,F)
 \longrightarrow H^i(A|R,F)\longrightarrow0
\]
for $i\leq n$, and hence that
\[
 H^i(A\otimes_RB|R,F)
 \longrightarrow H^i(A|R,F)\oplus H^i(B|R,F)
\]
is an isomorphism in that range.

(b) By (a), the map
\[
 H^i(A\otimes_RA|R,F)
 \longrightarrow H^i(A|R,F)\oplus H^i(A|R,F)
\]
is an isomorphism for every $i$. On the other hand, the assumed
isomorphism $A\otimes_RA\longrightarrow A$ induces an isomorphism
\[
 H^i(A|R,F)\longrightarrow H^i(A\otimes_RA|R,F).
\]
Their composite is the diagonal map
\[
 H^i(A|R,F)\longrightarrow
 H^i(A|R,F)\oplus H^i(A|R,F).
\]
Since it is an isomorphism, $H^i(A|R,F)=0$ for every $i$.
\hfill$\square$

\underline{Corollary \SGARefTarget{sga7.sga7i.evi.n3-9}{3.9}.} Let $S$ be a multiplicative subset of
an $R$-algebra $A$. Suppose that
\[
 F\in\operatorname{Ob}
 (\widetilde{\underline C}_{S^{-1}A|R})
\]
is additive and that $F_A$ is also additive. Then the canonical map
\[
 H^i(S^{-1}A|R,F)\longrightarrow H^i(A|R,F)
\]
is an isomorphism for every $i$. Here
$F_A\in\operatorname{Ob}
(\widetilde{\underline C}_{S^{-1}A|A})$ is defined by
\[
 F_A(X)=\ker\bigl(F(X)\longrightarrow F(A)\bigr)
\]
for every
$X\in\operatorname{Ob}(\underline C_{S^{-1}A|A})$.

\underline{Proof.} Applied to
$R\longrightarrow A\longrightarrow S^{-1}A$, Proposition \SGARefLink{sga7.sga7i.evi.n3-7}{3.7}(c)
gives the exact sequence
\[
 \cdots\longrightarrow H^n(S^{-1}A|A,F_A)
 \longrightarrow H^n(S^{-1}A|R,F)
 \longrightarrow H^n(A|R,F)\longrightarrow\cdots.
\]
The map
\[
 S^{-1}A\otimes_AS^{-1}A\longrightarrow S^{-1}A
\]
is an isomorphism, and $F_A$ is additive by hypothesis. Corollary
\SGARefLink{sga7.sga7i.evi.n3-8}{3.8}(b) therefore gives
$H^n(S^{-1}A|A,F_A)=0$ for every $n$. The middle map in the exact
sequence is consequently an isomorphism for every $n$.
\hfill$\square$

% END INCLUDED SGA7I_English_current_buildable_source_20260823/components/31_expose_VI_corollaries_3_8_and_3_9.tex

% BEGIN INCLUDED SGA7I_English_current_buildable_source_20260823/components/32_expose_VI_definition_and_lemma_3_10.tex
%%%%%%%%%%  scan idx 94  |  folio 83  %%%%%%%%%%

Let $A$ be an $R$-algebra. An $A$-module $M$ defines an additive
sheaf
\[
 \operatorname{Der}_R(-,M)
 \in\operatorname{Ob}(\widetilde{\underline C}_{A|R}),
 \qquad X\longmapsto\operatorname{Der}_R(X,M).
\]
We make the following definition.

\underline{Definition \SGARefTarget{sga7.sga7i.evi.n3-10}{3.10}.} Let $A$ be an $R$-algebra. For every
$A$-module $M$, set
\[
 D^i(A|R,M)=H^i(A|R,\operatorname{Der}_R(-,M)).
\]

We now compute the groups of low degree, namely $D^i(A|R,M)$ for
$i=1,2$. Choose a surjective arrow
$P\longrightarrow A$ in $\underline C_{A|R}$ with $P$ a
polynomial algebra over $R$, and write
\[
 0\longrightarrow I\longrightarrow P\longrightarrow A
 \longrightarrow0.
\]

\underline{Lemma \SGARefTarget{sga7.sga7i.evi.n3-10.occ-2}{3.10}.}
\begin{align*}
 D^1(A|R,M)
 &=\operatorname{Coker}\bigl(
   \operatorname{Der}_R(P,M)\longrightarrow
   \operatorname{Hom}_P(I,M)\bigr),\\
 D^2(A|R,M)
 &=H^1\bigl(A|R,
   \underline H^1_{A|R}(\operatorname{Der}_R(-,M))\bigr).
\end{align*}

\underline{Proof.} Consider the simplicial algebra
\[
 (*)\qquad
 \cdots\ P_A^{(2)}
 \mathrel{\substack{\longrightarrow\\[-.6ex]\longrightarrow\\[-.6ex]\longrightarrow}}
 P_A^{(1)}
 \mathrel{\substack{\longrightarrow\\[-.6ex]\longrightarrow}}
 P_A^{(0)},
\]
where
\[
 P_A^{(n)}=
 \overbrace{P\times_A\cdots\times_AP}^{n+1}.
\]
For every $n$ there is an exact sequence
\[
 0\longrightarrow I^{(n)}\longrightarrow P_A^{(n)}
 \longrightarrow A\longrightarrow0,
 \qquad
 I^{(n)}=\overbrace{I\times\cdots\times I}^{n+1}.
\]
Put $\overline I=I/I^2$, and let $\Delta$ denote the diagonal
map. It is straightforward to verify that, for every $n$, the
sequence
\[
 0\longrightarrow
 \operatorname{Hom}_A
 \bigl(\overline I^{(n)}/\Delta(\overline I),M\bigr)
 \longrightarrow\operatorname{Der}_R(P_A^{(n)},M)
 \xrightarrow{\operatorname{Der}(\Delta,M)}
 \operatorname{Der}_R(P,M)\longrightarrow0
 \tag{a}
\]
is exact. We also have the exact sequence
\[
 0\longrightarrow
 \operatorname{Hom}_A
 \bigl(\overline I^{(n)}/\Delta(\overline I),M\bigr)
 \longrightarrow\operatorname{Hom}_A(\overline I^{(n)},M)
 \xrightarrow{\operatorname{Hom}(\Delta,M)}
 \operatorname{Hom}_A(\overline I,M)\longrightarrow0.
 \tag{b}
\]
The sequences (a) and (b) yield the corresponding exact sequences
of simplicial $A$-modules. Since $\{I\longrightarrow0\}$ and
$\{I\longrightarrow I\}$ admit sections,
\[
 H^i(\{I\longrightarrow0\},\operatorname{Hom}_A(-,M))
 =H^i(\{I\longrightarrow I\},\operatorname{Hom}_A(-,M))=0
 \qquad(i>0).
\]

%%%%%%%%%%  scan idx 95  |  folio 84  %%%%%%%%%%

It follows that
\[
\left\{
\begin{aligned}
0\longrightarrow\operatorname{Der}_R(A,M)
&\longrightarrow\operatorname{Der}_R(P,M)
\longrightarrow\operatorname{Hom}_A(\overline I,M)\\
&\longrightarrow
H^1(\{P\longrightarrow A\},\operatorname{Der}_R(-,M))
\longrightarrow0
&&\text{is exact},\\
H^i(\{P\longrightarrow A\},\operatorname{Der}_R(-,M))
&=0 &&(i>1).
\end{aligned}
\right.
\]
By \SGARefLink{sga7.sga7i.evi.n3-1.item-b}{3.1}(b), therefore,
\begin{align*}
D^1(A|R,M)
 &=\check H^1(A|R,\operatorname{Der}_R(-,M))\\
 &=H^1(\{P\longrightarrow A\},\operatorname{Der}_R(-,M))\\
 &=\operatorname{Coker}\bigl(
   \operatorname{Der}_R(P,M)\longrightarrow
   \operatorname{Hom}_A(\overline I,M)\bigr)\\
 &=\operatorname{Coker}\bigl(
   \operatorname{Der}_R(P,M)\longrightarrow
   \operatorname{Hom}_P(I,M)\bigr),\\
D^2(A|R,M)
 &=H^1\bigl(A|R,
   \underline H^1_{A|R}(\operatorname{Der}_R(-,M))\bigr).
\end{align*}
\hfill$\square$

Choose a free $P$-module $F$ and a $P$-linear map
$\gamma:F\longrightarrow P$ such that $\gamma(F)=I$. Associated
with $\gamma$ is a Koszul complex, denoted by $K_I$ by abuse of
notation.

% END INCLUDED SGA7I_English_current_buildable_source_20260823/components/32_expose_VI_definition_and_lemma_3_10.tex

% BEGIN INCLUDED SGA7I_English_current_buildable_source_20260823/components/33_expose_VI_lemma_3_11.tex
\underline{Lemma \SGARefTarget{sga7.sga7i.evi.n3-11}{3.11}.} There is an isomorphism
\[
 D^2(A|R,M)\simeq
 \operatorname{Coker}\bigl(
 H^1(K_I,M)\longrightarrow
 \operatorname{Hom}_A(H_1(K_I),M)\bigr).
\]

\underline{Proof.} We have
\begin{align*}
D^2(A|R,M)
 &=H^1\bigl(A|R,
   \underline H^1_{A|R}(\operatorname{Der}_R(-,M))\bigr)\\
 &=\ker\left(
   H^1(P_A^{(1)}|R,M)
   \mathrel{\substack{\longrightarrow\\[-.6ex]\longrightarrow\\[-.6ex]\longrightarrow}}
   H^1(P_A^{(2)}|R,M)\right),
\end{align*}
because
\[
 [\underline H^1_{A|R}(\operatorname{Der}_R(-,M))](P_A^{(0)})
 =[\underline H^1_{A|R}(\operatorname{Der}_R(-,M))](P)=0.
\]

Let
\[
 0\longrightarrow K\longrightarrow F\xrightarrow{\gamma}I
 \longrightarrow0
\]
be exact. Write $P[F]$ for the symmetric algebra over $P$
generated by $F$. Then $P[F]$ and
$P[F]\otimes_PP[F]$ are polynomial algebras over $R$, and there
are surjective maps
\begin{align*}
P[F]&\longrightarrow P_A^{(1)}=P\times_AP,
&X&\longmapsto(0,\gamma(X)),\\
P[F]\otimes_PP[F]&\longrightarrow
P_A^{(2)}=P\times_AP\times_AP,
&X\otimes1&\longmapsto(0,\gamma(X),0),\\[-.3ex]
&&1\otimes X&\longmapsto(0,0,\gamma(X)),
\end{align*}
for $X\in F$. Put
\[
 J=\ker(P[F]\longrightarrow P_A^{(1)}),
 \qquad
 J'=\ker(P[F]\otimes_PP[F]\longrightarrow P_A^{(2)}).
\]
If $F_\gamma$ denotes the ideal generated by
$\{X-\gamma(X)\mid X\in F\}$, then
\begin{align*}
J&=(F)\cap(F_\gamma)=(K)+(F\cdot F_\gamma),\\
J'&=(F\otimes1+1\otimes F)
 \cap(F_\gamma\otimes1+1\otimes F)
 \cap(F\otimes1+1\otimes F_\gamma)\\
&=(K\otimes1+1\otimes K)
 +(F\cdot F_\gamma\otimes1+1\otimes F\cdot F_\gamma)
 +(F\otimes F).
\end{align*}

%%%%%%%%%%  scan idx 96  |  folio 85  %%%%%%%%%%

The three maps
$\pi_i:P_A^{(2)}\longrightarrow P_A^{(1)}$ admit lifts
$\widehat\pi_i:P[F]\otimes_PP[F]\longrightarrow P[F]$ given, for
$X\in F$, by
\[
\begin{array}{c@{\qquad}c@{\qquad}c}
\widehat\pi_0:
\begin{cases}
 X\otimes1\longmapsto\gamma(X)-X,\\
 1\otimes X\longmapsto X,
\end{cases}
&
\widehat\pi_1:
\begin{cases}
 X\otimes1\longmapsto0,\\
 1\otimes X\longmapsto X,
\end{cases}
&
\widehat\pi_2:
\begin{cases}
 X\otimes1\longmapsto X,\\
 1\otimes X\longmapsto0.
\end{cases}
\end{array}
\]
A direct computation shows that
$\widehat\pi_0-\widehat\pi_1+\widehat\pi_2$ annihilates
\[
 K\otimes1+1\otimes K+1\otimes F\cdot F_\gamma
\]
and maps
$F\cdot F_\gamma\otimes1+F\otimes F$ onto
$F\cdot F_\gamma$. For every
\[
 D\in\operatorname{Der}_R(P[F]\otimes_PP[F],M)
\]
one has
$D(F\cdot F_\gamma\otimes1+F\otimes F)=0$, because $I\cdot M=0$.
Consequently,
\begin{align*}
D^2(A|R,M)
 &=\ker\left(
   D^1(P_A^{(1)}|R,M)
   \mathrel{\substack{\longrightarrow\\[-.6ex]\longrightarrow\\[-.6ex]\longrightarrow}}
   D^1(P_A^{(2)}|R,M)\right)\\
 &=\{[f]\in D^1(P_A^{(1)}|R,M)
       \mid f(F\cdot F_\gamma)=0\}\\
 &=\operatorname{Coker}\left(
   \operatorname{Der}_R(P[F],M)\longrightarrow
   \operatorname{Hom}_{P[F]}
   \left(\frac{(F)\cap(F_\gamma)}{F\cdot F_\gamma},M\right)
   \right)\\
 &=\operatorname{Coker}\left(
   \operatorname{Der}_P(P[F],M)\longrightarrow
   \operatorname{Hom}_{P[F]}
   \left(\frac{(F)\cap(F_\gamma)}{(F)(F_\gamma)},M\right)
   \right).
\end{align*}

The sequences
\[
 0\longrightarrow(F)\longrightarrow P[F]
 \xrightarrow{P[0]}P\longrightarrow0,
 \qquad
 0\longrightarrow(F_\gamma)\longrightarrow P[F]
 \xrightarrow{P[\gamma]}P_\gamma\longrightarrow0
\]
are exact. Hence
\[
 \frac{(F)\cap(F_\gamma)}{(F)(F_\gamma)}
 =\operatorname{Tor}_1^{P[F]}(P,P_\gamma)
 =H_1(K_I),
\]
where $K_I$ is the Koszul complex associated with
$F\xrightarrow{\gamma}I\subset P$. We therefore obtain
\begin{align*}
D^2(A|R,M)
 &=\operatorname{Coker}\bigl(
   \operatorname{Hom}_P(F,M)\longrightarrow
   \operatorname{Hom}_P(H_1(K_I),M)\bigr)\\
 &=\operatorname{Coker}\bigl(
   H^1(K_I,M)\longrightarrow
   \operatorname{Hom}_A(H_1(K_I),M)\bigr).
\end{align*}
\hfill$\square$

We can now state the basic facts needed below, summarizing the
results already established.

% END INCLUDED SGA7I_English_current_buildable_source_20260823/components/33_expose_VI_lemma_3_11.tex

% BEGIN INCLUDED SGA7I_English_current_buildable_source_20260823/components/34_expose_VI_theorem_3_12_statement.tex
%%%%%%%%%%  scan idx 97  |  folio 86  %%%%%%%%%%

\underline{Theorem \SGARefTarget{sga7.sga7i.evi.n3-12}{3.12}.}

\begin{enumerate}[label=(\arabic*)]
\item If
\[
 0\longrightarrow M'\longrightarrow M\longrightarrow M''
 \longrightarrow0
\]
is an exact sequence of $A$-modules, then the sequence
\begin{multline*}
0\longrightarrow D^0(A|R,M')\longrightarrow D^0(A|R,M)
\longrightarrow D^0(A|R,M'')\longrightarrow D^1(A|R,M')\\
\longrightarrow D^1(A|R,M)\longrightarrow D^1(A|R,M'')
\longrightarrow\cdots\longrightarrow D^n(A|R,M')\\
\longrightarrow D^n(A|R,M)\longrightarrow D^n(A|R,M'')
\longrightarrow\cdots
\end{multline*}
is exact.

\item Let $S\longrightarrow R\longrightarrow A$ be ring
homomorphisms. Then, for every $A$-module $M$, the sequence
\begin{multline*}
0\longrightarrow D^0(A|R,M)\longrightarrow D^0(A|S,M)
\longrightarrow D^0(R|S,M)\longrightarrow D^1(A|R,M)\\
\longrightarrow D^1(A|S,M)\longrightarrow D^1(R|S,M)
\longrightarrow\cdots\longrightarrow D^n(A|R,M)\\
\longrightarrow D^n(A|S,M)\longrightarrow D^n(R|S,M)
\longrightarrow\cdots
\end{multline*}
is exact.

\item Let $A$ and $S$ be $R$-algebras such that
$\operatorname{Tor}_i^R(A,S)=0$ for every $i>0$. Then, for every
$A\otimes_RS$-module $M$, the canonical map
\[
 D^i(A\otimes_RS|S,M)\longrightarrow D^i(A|R,M)
\]
is an isomorphism for every $i$. In particular, if $A$ is flat
over $R$, then
\[
 D^i(A\otimes_RS|S,M)\xrightarrow{\ \sim\ }D^i(A|R,M)
\]
for every $i$, every $R$-algebra $S$, and every
$A\otimes_RS$-module $M$.

\item Let $S$ be a multiplicative subset of $A$. For every
$S^{-1}A$-module $N$, the canonical map
\[
 D^i(S^{-1}A|R,N)\longrightarrow D^i(A|R,N)
\]
is an isomorphism for every $i$. Consequently, for every
$A$-module $M$, there is a canonical isomorphism
\[
 D^i(S^{-1}A|R,S^{-1}M)
 \xrightarrow{\ \sim\ }S^{-1}D^i(A|R,M)
\]
for every $i$.

\item Let $A$ and $B$ be $R$-algebras such that
$\operatorname{Tor}_i^R(A,B)=0$ for $0<i\leq n$. Then, for every
$A\otimes_RB$-module $M$, the canonical map
\[
 D^i(A\otimes_RB|R,M)
 \longrightarrow D^i(A|R,M)\oplus D^i(B|R,M)
\]
is an isomorphism for $i\leq n$. In particular, suppose that
$\operatorname{Tor}_i^R(A,A)=0$ for every $i>0$ and that the
diagonal
\[
 \Delta:\operatorname{Spec}(A)\longrightarrow
 \operatorname{Spec}(A)\times_{\operatorname{Spec}(R)}
 \operatorname{Spec}(A)
\]
is an open immersion---for example, suppose that $A$ is \emph{étale}
over $R$. Then
\[
 D^i(A|R,M)=0
\]
for every $i>0$ and every $A$-module $M$.

%%%%%%%%%%  scan idx 98  |  folio 87  %%%%%%%%%%

\item Let $R$ be noetherian, let $A$ be an $R$-algebra essentially
of finite type, and let $E$ be a flat $A$-module. Then the canonical
map
\[
 E\otimes_AD^i(A|R,M)
 \xrightarrow{\ \sim\ }
 D^i(A|R,E\otimes_AM)
\]
is an isomorphism for every $i\geq0$ and every $A$-module $M$.

\item If $R$ is noetherian, $A$ is an $R$-algebra essentially of
finite type, and $M$ is an $A$-module of finite type, then
$D^i(A|R,M)$ is an $A$-module of finite type for every $i$.

\item Let $P$ be a polynomial algebra over $R$, and let
\[
 0\longrightarrow I\longrightarrow P\longrightarrow A
 \longrightarrow0
\]
be exact. Then
\[
 D^0(A|R,M)=\operatorname{Der}_R(A,M),
\]
\begin{align*}
D^1(A|R,M)
 &=\operatorname{Coker}\bigl(
   \operatorname{Der}_R(P,M)\longrightarrow
   \operatorname{Hom}_A(I/I^2,M)\bigr)\\
 &=\left\{
   \begin{array}{c}
   \text{isomorphism classes of commutative}\
   R\text{-algebra extensions of }A\text{ by }M
   \end{array}\right\},\\
D^2(A|R,M)
 &=\operatorname{Coker}\bigl(
   H^1(K_I,M)\longrightarrow
   \operatorname{Hom}_A(H_1(K_I),M)\bigr).
\end{align*}
\end{enumerate}

% END INCLUDED SGA7I_English_current_buildable_source_20260823/components/34_expose_VI_theorem_3_12_statement.tex

% BEGIN INCLUDED SGA7I_English_current_buildable_source_20260823/components/35_expose_VI_theorem_3_12_proof.tex
\underline{Proof.}

\begin{enumerate}[label=(\arabic*)]
\item If
\[
 0\longrightarrow M'\longrightarrow M\longrightarrow M''
 \longrightarrow0
\]
is an exact sequence of $A$-modules, then, for every polynomial
algebra $P$ over $R$, the sequence
\[
 0\longrightarrow\operatorname{Der}_R(P,M')
 \longrightarrow\operatorname{Der}_R(P,M)
 \longrightarrow\operatorname{Der}_R(P,M'')
 \longrightarrow0
\]
is exact. Hence, by \SGARefLink{sga7.sga7i.evi.n3-3}{3.3}(c),
\[
 0\longrightarrow\operatorname{Der}_R(-,M')
 \longrightarrow\operatorname{Der}_R(-,M)
 \longrightarrow\operatorname{Der}_R(-,M'')
 \longrightarrow0
\]
is an exact sequence of sheaves.

\item The functor $\operatorname{Der}_S(-,M)$ is an additive sheaf
on $\underline C_{A|S}$, and
\[
 \operatorname{Der}_R(B,M)
 =\ker\bigl(\operatorname{Der}_S(B,M)
 \longrightarrow\operatorname{Der}_S(R,M)\bigr)
\]
for every $B\in\operatorname{Ob}(\underline C_{A|R})$.
The assertion therefore follows from \SGARefLink{sga7.sga7i.evi.n3-7}{3.7}(c).

\item Let
\[
 S:\underline C_{A|R}\longrightarrow
 \underline C_{A\otimes_RS|S},
 \qquad X\longmapsto X\otimes_RS.
\]
Then
$\widetilde S\operatorname{Der}_S(-,M)
=\operatorname{Der}_R(-,M)$, so the assertion follows from \SGARefLink{sga7.sga7i.evi.n3-7}{3.7}(a).

%%%%%%%%%%  scan idx 99  |  folio 88  %%%%%%%%%%

\item This follows from \SGARefLink{sga7.sga7i.evi.n3-9}{3.9}. Indeed,
$F=\operatorname{Der}_R(-,N)$ is an additive sheaf on
$\underline C_{S^{-1}A|R}$, and
$F_A\in\operatorname{Ob}
(\widetilde{\underline C}_{S^{-1}A|A})$ is also additive, since
\begin{align*}
 F_A(X)
 &=\ker\bigl(F(X)\longrightarrow F(A)\bigr)\\
 &=\ker\bigl(\operatorname{Der}_R(X,N)
   \longrightarrow\operatorname{Der}_R(A,N)\bigr)\\
 &=\operatorname{Der}_A(X,N).
\end{align*}
Thus $F_A=\operatorname{Der}_A(-,N)$.

\item The first assertion follows from \SGARefLink{sga7.sga7i.evi.n3-8}{3.8}(a). Now suppose that
$\operatorname{Tor}_i^R(A,A)=0$ for every $i>0$ and that
\[
 \Delta:\operatorname{Spec}(A)\longrightarrow
 \operatorname{Spec}(A)\times_{\operatorname{Spec}(R)}
 \operatorname{Spec}(A)
\]
is an open immersion. Let
$\varphi:A\otimes_RA\longrightarrow A$ be the multiplication map.
For every maximal ideal $\mathfrak m$ of $A$, the induced map
\[
 (A\otimes_RA)_{\varphi^{-1}(\mathfrak m)}
 \longrightarrow A_{\mathfrak m}
\]
is an isomorphism. The first assertion, together with \SGARefLink{sga7.sga7i.evi.n3-12}{3.12}(4),
therefore implies that the diagonal map
\[
 \Delta_{\mathfrak m}:D^i(A|R,M)_{\mathfrak m}
 \longrightarrow D^i(A|R,M)_{\mathfrak m}
 \oplus D^i(A|R,M)_{\mathfrak m}
\]
is an isomorphism for every maximal ideal $\mathfrak m$ of $A$.
Consequently,
\[
 \Delta:D^i(A|R,M)\longrightarrow
 D^i(A|R,M)\oplus D^i(A|R,M)
\]
is an isomorphism, and hence $D^i(A|R,M)=0$ for every $i>0$.

\item In view of \SGARefLink{sga7.sga7i.evi.n3-12}{3.12}(4), we may assume that $A$ is an
$R$-algebra of finite type. Consider the morphism of sheaves
\[
 E\otimes_A\operatorname{Der}_R(-,M)
 \longrightarrow
 \operatorname{Der}_R(-,E\otimes_AM)
\]
on $\widetilde{\underline C}_{A|R}$. Since $E$ is flat over $A$,
the map
\[
 E\otimes_A\operatorname{Der}_R(B,M)
 \xrightarrow{\ \sim\ }
 \operatorname{Der}_R(B,E\otimes_AM)
\]
is an isomorphism whenever
$B\in\operatorname{Ob}(\underline C_{A|R})$ is of finite type
over $R$. By \SGARefLink{sga7.sga7i.evi.n3-3}{3.3}(b), it follows that
\[
 D^i\bigl(A|R,E\otimes_A\operatorname{Der}_R(-,M)\bigr)
 \longrightarrow
 D^i\bigl(A|R,\operatorname{Der}_R(-,E\otimes_AM)\bigr)
\]
is an isomorphism. On the other hand,
$E\otimes_A-$ is exact, and therefore
\[
 E\otimes_AD^i(A|R,M)
 =D^i\bigl(A|R,E\otimes_A\operatorname{Der}_R(-,M)\bigr).
\]
This proves the assertion.

\item This follows from \SGARefLink{sga7.sga7i.evi.n3-3}{3.3}(a).

%%%%%%%%%%  scan idx 100  |  folio 89  %%%%%%%%%%

\item In view of \SGARefLink{sga7.sga7i.evi.n3-10.occ-2}{3.10} and \SGARefLink{sga7.sga7i.evi.n3-11}{3.11}, it suffices to show that
\[
 \operatorname{Coker}\bigl(
 \operatorname{Der}_R(P,M)\longrightarrow
 \operatorname{Hom}_A(I/I^2,M)\bigr)
 =\operatorname{Exalcomm}(A|R,M),
\]
where $\operatorname{Exalcomm}(A|R,M)$ denotes the set of
isomorphism classes of $R$-algebra extensions of $A$ by $M$.
Let
\[
 0\longrightarrow M\longrightarrow A'\longrightarrow A
 \longrightarrow0
\]
be such an extension. Since $P$ is a polynomial algebra over $R$,
we obtain a commutative diagram
\[
\begin{tikzcd}[column sep=large]
0 \arrow[r] & I \arrow[r] \arrow[d]
& P \arrow[r] \arrow[d]
& A \arrow[r] \arrow[d,equal] & 0 \\
0 \arrow[r] & M \arrow[r]
& A' \arrow[r]
& A \arrow[r] & 0.
\end{tikzcd}
\]
The map $I\longrightarrow M$ is determined by the extension $A'$,
up to the action of $\operatorname{Der}_R(P,M)$. We thus obtain a
map
\[
 \Phi:\operatorname{Exalcomm}(A|R,M)
 \longrightarrow
 \operatorname{Coker}\bigl(
 \operatorname{Der}_R(P,M)\longrightarrow
 \operatorname{Hom}_A(I/I^2,M)\bigr).
\]
Conversely, a $P$-linear map $I\longrightarrow M$ yields, by
pushout, an extension fitting into a diagram
\[
\begin{tikzcd}[column sep=large]
0 \arrow[r] & I \arrow[r] \arrow[d]
& P \arrow[r] \arrow[d]
& A \arrow[r] \arrow[d,equal] & 0 \\
0 \arrow[r] & M \arrow[r]
& A' \arrow[r]
& A \arrow[r] & 0.
\end{tikzcd}
\]
Two $P$-linear maps $I\rightrightarrows M$ that differ by an
element of $\operatorname{Der}_R(P,M)$ yield isomorphic
extensions. Hence there is a map
\[
 \Psi:\operatorname{Coker}\bigl(
 \operatorname{Der}_R(P,M)\longrightarrow
 \operatorname{Hom}_A(I/I^2,M)\bigr)
 \longrightarrow\operatorname{Exalcomm}(A|R,M).
\]
It is straightforward to verify that $\Phi$ and $\Psi$ are inverse
to each other.
\end{enumerate}
\hfill$\square$

% END INCLUDED SGA7I_English_current_buildable_source_20260823/components/35_expose_VI_theorem_3_12_proof.tex

% BEGIN INCLUDED SGA7I_English_current_buildable_source_20260823/components/36_expose_VI_corollary_3_13.tex
%%%%%%%%%%  scan idx 100  |  folio 89  %%%%%%%%%%

\underline{Corollary \SGARefTarget{sga7.sga7i.evi.n3-13}{3.13}.}

\begin{enumerate}[label=(\alph*)]
\item The $R$-algebra $A$ is formally smooth over $R$ if and only if
$D^1(A|R,M)=0$ for every $A$-module $M$.

\item Let $R$ be noetherian and let $A$ be an $R$-algebra of finite
type. Then $A$ is a relative complete intersection over $R$ if and
only if $D^2(A|R,M)=0$ for every $A$-module $M$.

\item Let $R$ be noetherian and let $A$ be an $R$-algebra of finite
type. For every multiplicative subset $S$ of $A$, there is a canonical
isomorphism
\[
 D^i(S^{-1}A|R,S^{-1}M)
 \xrightarrow{\ \sim\ }S^{-1}D^i(A|R,M)
 \qquad(i\geq0).
\]

\item Let $R$ be noetherian and let $A$ be a local $R$-algebra
essentially of finite type. Write $\widehat A$ for the
$\mathfrak m$-adic completion of $A$, where $\mathfrak m$ is the
maximal ideal of $A$. Then, for every $A$-module $E$ of finite type,
there is a canonical isomorphism
\[
 D^1(\widehat A|R,\widehat A\otimes_AE)
 \xrightarrow{\ \sim\ }
 \widehat A\otimes_AD^1(A|R,E).
\]
\end{enumerate}

\underline{Proof.}

Part (a) follows immediately from
$D^1(A|R,M)=\operatorname{Exalcomm}(A|R,M)$.

For (b), let $P$ be a polynomial algebra over $R$ in finitely many
variables, and let
\[
 0\longrightarrow I\longrightarrow P\longrightarrow A
 \longrightarrow0
\]
be exact. After localization, we may suppose that $I$ is generated by
a $P$-regular sequence. Then $H_1(K_I)=0$, and hence
\[
 D^2(A|R,M)
 =\operatorname{Coker}\bigl(
 H^1(K_I,M)\longrightarrow
 \operatorname{Hom}_A(H_1(K_I),M)\bigr)=0.
\]
Conversely, suppose that $D^2(A|R,M)=0$ for every $A$-module $M$.
Choose an exact sequence
\[
 0\longrightarrow K\longrightarrow F\longrightarrow I
 \longrightarrow0,
\]
where $F$ is a free $P$-module of finite type. If
\[
 K_0=\operatorname{Im}\bigl(\textstyle\bigwedge^2F
 \longrightarrow\bigwedge^1F\bigr),
\]
then $H_1(K_I)=K/K_0$. Since
$D^2(A|R,H_1(K_I))=0$, the map
\[
 H^1(K_I,H_1(K_I))
 \longrightarrow
 \operatorname{Hom}_A(H_1(K_I),H_1(K_I))
\]
is surjective. Consequently, the exact sequence
\[
 0\longrightarrow K/K_0\longrightarrow F/IF
 \longrightarrow I/I^2\longrightarrow0
\]
splits as a sequence of $A$-modules. In particular, $I/I^2$ is
projective. After localizing $A$ once more, we may suppose that
$I/I^2$ is free and that $F/IF\longrightarrow I/I^2$ is an
isomorphism. The split exact sequence then gives
$H_1(K_I)=K/K_0=0$. Thus $I$ is generated by a $P$-regular sequence,
so $A$ is a relative complete intersection over $R$.

Part (c) follows immediately from \SGARefLink{sga7.sga7i.evi.n3-12}{3.12}(4) and (6).

%%%%%%%%%%  scan idx 101  |  folio 90  %%%%%%%%%%

For (d), \SGARefLink{sga7.sga7i.evi.n3-12}{3.12}(6) first gives canonical isomorphisms
\[
 D^i(A|R,\widehat A\otimes_AE)
 \xrightarrow{\ \sim\ }
 \widehat A\otimes_AD^i(A|R,E)
\]
for every $i$. It therefore remains to prove that
\[
 D^1(\widehat A|R,\widehat A\otimes_AE)
 \longrightarrow D^1(A|R,\widehat A\otimes_AE)
\]
is an isomorphism. We use the description of $D^1$ in terms of
algebra extensions.

Let
\[
 0\longrightarrow\widehat A\otimes_AE\longrightarrow A'
 \xrightarrow{\ \varphi\ }\widehat A\longrightarrow0
\]
be an $R$-algebra extension whose restriction to $A$ splits. Since
$\widehat A\otimes_AE$ is an $\widehat A$-module of finite type,
$A'$ is a local $R$-algebra complete for its $\mathfrak m'$-adic
topology, where $\mathfrak m'$ is the maximal ideal of $A'$. If
$\psi:A\longrightarrow A'$ is an $R$-algebra map lifting the
canonical map $A\longrightarrow\widehat A$, then $\psi$ extends
uniquely to a map
$\widehat\psi:\widehat A\longrightarrow A'$ satisfying
$\varphi\circ\widehat\psi=\operatorname{id}_{\widehat A}$.
This proves that
\[
 D^1(\widehat A|R,\widehat A\otimes_AE)
 \longrightarrow D^1(A|R,\widehat A\otimes_AE)
\]
is injective.

%%%%%%%%%%  scan idx 102  |  folio 91  %%%%%%%%%%

Now let
\[
 0\longrightarrow\widehat A\otimes_AE\longrightarrow B
 \xrightarrow{\ \varphi\ }A\longrightarrow0
\]
be an $R$-algebra extension, and put
$\widehat E=\widehat A\otimes_AE$. We first claim that the topology
on $\widehat E$ induced by the $\mathfrak m_B$-adic topology on $B$
coincides with its topology as an $\widehat A$-module. Indeed, since
$A$ is essentially of finite type over $R$, there is a subalgebra
$B_0\subset B$ such that $B_0$ is a local $R$-algebra essentially
of finite type and $\varphi(B_0)=A$. Put $E_0=\widehat E\cap B_0$,
and denote the maximal ideals of $B$, $B_0$, and $A$ by
$\mathfrak m_B$, $\mathfrak m_0$, and $\mathfrak m$, respectively.
Then
\[
 \mathfrak m_B=\mathfrak m_0+\widehat E,
 \qquad
 \mathfrak m_B^{n+1}=\mathfrak m_0^{n+1}
   +\mathfrak m^n\widehat E
\]
for every $n$. Since $B_0$ is noetherian, the Artin--Rees theorem
gives an integer $r>0$ such that, for every $n\geq r$,
\begin{align*}
 \widehat E\cap\mathfrak m_B^{n+1}
 &=\mathfrak m^n\widehat E
   +\widehat E\cap\mathfrak m_0^{n+1}\\
 &=\mathfrak m^n\widehat E
   +E_0\cap\mathfrak m_0^{n+1}\\
 &=\mathfrak m^n\widehat E
   +\mathfrak m^{n+1-r}(E_0\cap\mathfrak m_0^r)
 \subset\mathfrak m^{n+1-r}\widehat E.
\end{align*}
This proves the assertion about the two topologies. In particular,
$\widehat E$ is complete for the topology induced by the
$\mathfrak m_B$-adic topology on $B$, and hence completion gives an
exact commutative diagram
\[
\begin{tikzcd}[column sep=large]
0 \arrow[r] & \widehat E \arrow[r] \arrow[d,equal]
& B \arrow[r] \arrow[d]
& A \arrow[r] \arrow[d] & 0 \\
0 \arrow[r] & \widehat E \arrow[r]
& \widehat B \arrow[r]
& \widehat A \arrow[r] & 0.
\end{tikzcd}
\]
Here $\widehat B$ denotes the $\mathfrak m_B$-adic completion of
$B$. Thus
\[
 D^1(\widehat A|R,\widehat E)
 \longrightarrow D^1(A|R,\widehat E)
\]
is surjective, and the proof is complete.
\hfill$\square$

% END INCLUDED SGA7I_English_current_buildable_source_20260823/components/36_expose_VI_corollary_3_13.tex

% BEGIN INCLUDED SGA7I_English_current_buildable_source_20260823/components/37_expose_VI_3_14_through_3_16.tex
%%%%%%%%%%  scan idx 103  |  folio 92  %%%%%%%%%%

\underline{Corollary \SGARefTarget{sga7.sga7i.evi.n3-14}{3.14}.}
Let $R$ be noetherian, and let $A$ be an $R$-algebra of finite type
that is smooth everywhere except at one closed point
$x\in\operatorname{Spec}(A)$. Then there are isomorphisms
\[
 D^1(A|R,A)\simeq D^1(A_x|R,A_x)
 \simeq D^1(\widehat A_x|R,\widehat A_x).
\]

\underline{Proof.}
Since $A$ is smooth over $R$ away from $x$,
$D^1(A|R,A)$ is an $A$-module of finite type supported on
$\{x\}$. Hence
\[
 D^1(A|R,A)=D^1(A|R,A)_x
 =\widehat{D^1(A|R,A)_x}.
\]
The assertion now follows from \SGARefLink{sga7.sga7i.evi.n3-13}{3.13}.
\hfill$\square$

\underline{Remark \SGARefTarget{sga7.sga7i.evi.n3-15}{3.15}.}
Let $R$ be a topological ring and $A$ a topological $R$-algebra.
For every $A$-module $E$, define
\[
 D^i_{\mathrm{top}}(A|R,E)
 =\varinjlim_\alpha
 D^i(A_\alpha|R_\alpha,A_\alpha\otimes_AE),
\]
where $A_\alpha=A/J_\alpha$, $R_\alpha=R/I_\alpha$, and
$J_\alpha$ and $I_\alpha$ are open in $A$ and $R$, respectively.
Thus $A$ is formally smooth as a topological $R$-algebra precisely
when $D^1_{\mathrm{top}}(A|R,E)=0$ for every $A$-module $E$.

Now let $X$ be a scheme over a ring $R$, and let $\underline E$ be a
quasi-coherent $\mathcal O_X$-module. We define sheaves
\[
 \underline D^i_{X|R}(\underline E)
 =\underline D^i(X|R,\underline E)
\]
on $X$ from the presheaves on affine open subsets $U\subset X$
given by
\[
 \bigl[\underline D^i_{X|R}(\underline E)\bigr](U)
 =D^i\bigl(\Gamma(U,\mathcal O_X)|R,
            \Gamma(U,\underline E)\bigr).
\]
If $R$ is noetherian and $X$ is of finite type over $R$, the
localization property \SGARefLink{sga7.sga7i.evi.n3-13}{3.13}(c) shows that these presheaves canonically
form quasi-coherent $\mathcal O_X$-modules.

When no confusion is possible, we abbreviate
\[
 \underline D^i_{X|R}
 =\underline D^i_{X|R}(\mathcal O_X).
\]
All the assertions of \SGARefLink{sga7.sga7i.evi.n3-13}{3.13} have corresponding formulations for the
quasi-coherent sheaves $\underline D^i_{X|R}(\underline E)$. We leave
those formulations to the reader and record below the properties
most relevant to our purpose.

\underline{Corollary \SGARefTarget{sga7.sga7i.evi.n3-15.occ-2}{3.15}.}
Let $R$ be noetherian and let $X$ be an $R$-scheme of finite type.

%%%%%%%%%%  scan idx 104  |  folio 93  %%%%%%%%%%

\begin{enumerate}[label=(\alph*)]
\item For every coherent $\mathcal O_X$-module $\underline E$,
$\underline D^i_{X|R}(\underline E)$ is again coherent for every
$i$. Moreover,
\[
 \operatorname{Supp}\underline D^i_{X|R}(\underline E)
 \subset X^{\mathrm{Sing}}
 \qquad(i>0),
\]
where
\[
 X^{\mathrm{Sing}}
 =\{x\in X\mid X\text{ is not smooth over }R\text{ at }x\}.
\]

\item If $X$ is flat over $R$, then for every $R$-algebra $S$ there
is a canonical isomorphism
\[
 \underline D^i_{X\otimes_RS\,|\,S}(\underline E)
 \xrightarrow{\ \sim\ }
 \underline D^i_{X|R}(\underline E)
\]
for every $i$.

\item Suppose that $X$ is a relative complete intersection over
$R$. For every surjective ring homomorphism $S\longrightarrow R$,
there is an exact sequence
\[
0\longrightarrow\underline D^1_{X|R}(\underline E)
\longrightarrow\underline D^1_{X|S}(\underline E)
\longrightarrow
\underline{\operatorname{Hom}}_{\mathcal O_X}
 \bigl((I/I^2)\otimes_R\mathcal O_X,\underline E\bigr)
\longrightarrow0,
\]
where $I=\ker(S\longrightarrow R)$. If, moreover, $R$ is local
with residue field $k$ and $X$ is flat over $R$, then
\[
0\longrightarrow
 \underline D^1_{X_0|k}(\mathcal O_{X_0})
\longrightarrow
 \underline D^1_{X|S}(\mathcal O_{X_0})
\longrightarrow
 D^1(R|S,k)\otimes_k\mathcal O_{X_0}
\longrightarrow0,
\]
where $X_0=X\otimes_Rk$.
\end{enumerate}

\underline{Proof.}
Part (a) follows from the finiteness and low-degree descriptions in
\SGARefLink{sga7.sga7i.evi.n3-12}{3.12}, while (b) follows from \SGARefLink{sga7.sga7i.evi.n3-12}{3.12}(3).

For (c), because $X$ is a relative complete intersection over $R$,
we have $\underline D^i_{X|R}(\underline E)=0$ for every $i\geq2$.
The transitivity sequence \SGARefLink{sga7.sga7i.evi.n3-12}{3.12}(2) therefore gives an exact sequence
\begin{multline*}
0\longrightarrow\underline D^0_{X|R}(\underline E)
\longrightarrow\underline D^0_{X|S}(\underline E)
\longrightarrow\underline D^0_{R|S}(\underline E)
\longrightarrow\underline D^1_{X|R}(\underline E)\\
\longrightarrow\underline D^1_{X|S}(\underline E)
\longrightarrow\underline D^1_{R|S}(\underline E)
\longrightarrow0,
\end{multline*}
where the quasi-coherent sheaves
$\underline D^i_{R|S}(\underline E)$ on $X$ are defined on every
affine open subset $U\subset X$ by
\[
 \bigl[\underline D^i_{R|S}(\underline E)\bigr](U)
 =D^i(R|S,\underline E(U)).
\]
Since $S\longrightarrow R$ is surjective,
$\underline D^0_{R|S}(\underline E)=0$; and the low-degree
description in \SGARefLink{sga7.sga7i.evi.n3-12}{3.12}(8) gives
\[
 \underline D^1_{R|S}(\underline E)
 =\underline{\operatorname{Hom}}_{\mathcal O_X}
  \bigl((I/I^2)\otimes_R\mathcal O_X,\underline E\bigr).
\]
This proves the first exact sequence.

%%%%%%%%%%  scan idx 105  |  folio 94  %%%%%%%%%%

If $X$ is flat over $R$, then, for every quasi-coherent
$\mathcal O_{X_0}$-module $\underline E$, there are canonical
isomorphisms
\[
 \underline D^i_{X_0|k}(\underline E)
 \xrightarrow{\ \sim\ }
 \underline D^i_{X|R}(\underline E)
\]
for every $i$. In particular,
\[
 \underline D^1_{X_0|k}(\mathcal O_{X_0})
 \xrightarrow{\ \sim\ }
 \underline D^1_{X|R}(\mathcal O_{X_0}),
\]
from which the second exact sequence follows.
\hfill$\square$

\underline{Remark \SGARefTarget{sga7.sga7i.evi.n3-16}{3.16}.}
The quasi-coherent sheaves $\underline D^i_{X|R}(\underline E)$
defined above should occur as the terms $E_2^{0,i}$ of a spectral
sequence relating the local theory to the global theory. L.~Illusie
has obtained such a theory in \emph{Lecture Notes in Mathematics}
239. His method uses ``homotopical algebra,'' generalizing the methods
of M.~André and D.~Quillen, rather than ``cohomological algebra.'' It
seems that the approach taken here can be globalized by means of
Grothendieck cohomology on the category of schemes, equipped with an
appropriate topology generalizing the one on the affine category.

% END INCLUDED SGA7I_English_current_buildable_source_20260823/components/37_expose_VI_3_14_through_3_16.tex

% BEGIN INCLUDED SGA7I_English_current_buildable_source_20260823/components/38_expose_VI_section_4_opening_through_4_1.tex
%%%%%%%%%%  scan idx 105  |  folio 94  %%%%%%%%%%

\section*{\SGARefTarget{sga7.sga7i.evi.n4}{4}. Formal moduli of deformations}

Among the natural cofibered categories whose isomorphism classes of
objects are generally not prorepresentable but nevertheless admit a
formal versal object is the cofibered category of deformations. We
return to the situation and notation of section 1. Thus, throughout
this section, $\Lambda$ is a fixed complete noetherian local ring with
residue field $k$, and $\underline C_\Lambda$ is the category of
artinian local $\Lambda$-algebras with residue field $k$.

Let $X_0$ be a fixed scheme over $k$. By an (infinitesimal)
deformation of $X_0$ we mean a pair $(R,X)$, where
$R\in\operatorname{Ob}(\underline C_\Lambda)$ and $X$ is a flat
$R$-scheme, together with a commutative diagram
\[
\begin{tikzcd}[column sep=huge,row sep=huge]
X_0 \arrow[r,"i_X"] \arrow[d] & X \arrow[d,"\mathrm{flat}"] \\
\operatorname{Spec}(k) \arrow[r] & \operatorname{Spec}(R)
\end{tikzcd}
\]
such that the induced morphism
\[
 X_0\longrightarrow
 X\times_{\operatorname{Spec}(R)}\operatorname{Spec}(k)
\]
is an isomorphism. Since $R$ is artinian local,
$i_X:X_0\longrightarrow X$ is necessarily a closed immersion and a
homeomorphism on the underlying topological spaces.

%%%%%%%%%%  scan idx 106  |  folio 95  %%%%%%%%%%

If $(R,X)$ and $(R',X')$ are deformations of $X_0$ over $R$ and
$R'$, respectively, an arrow
\[
 (R',X')\longrightarrow(R,X)
\]
consists of a map $\varphi:R'\longrightarrow R$ in
$\underline C_\Lambda$ and a morphism of schemes
$\Phi:X\longrightarrow X'$ fitting into a commutative diagram
\[
\begin{tikzcd}[column sep=huge,row sep=large]
X \arrow[rr,"\Phi"] \arrow[ddd] && X' \arrow[ddd] \\
& X_0 \arrow[ul,"i_X"'] \arrow[ur,"i_{X'}"] \arrow[d] & \\
& \operatorname{Spec}(k) \arrow[dl] \arrow[dr] & \\
\operatorname{Spec}(R)
 \arrow[rr,"\operatorname{Spec}(\varphi)"']
&& \operatorname{Spec}(R').
\end{tikzcd}
\]
This defines the category $\underline M_{X_0}$ of infinitesimal
deformations of $X_0$. It is equipped with the functor
\[
 p:\underline M_{X_0}\longrightarrow\underline C_\Lambda,
 \qquad (R,X)\longmapsto R.
\]

Let $(R,X)\in\operatorname{Ob}(\underline M_{X_0})$ and let
$\varphi:R\longrightarrow S$ be a map in $\underline C_\Lambda$.
There is a commutative diagram
\[
\begin{tikzcd}[column sep=huge,row sep=huge]
X\times_{\operatorname{Spec}(R)}\operatorname{Spec}(S)
 \arrow[r,"p_1"] \arrow[d,"\mathrm{flat}"']
& X \arrow[d,"\mathrm{flat}"] \\
\operatorname{Spec}(S)
 \arrow[r,"\operatorname{Spec}(\varphi)"']
& \operatorname{Spec}(R),
\end{tikzcd}
\]
and the $\varphi$-arrow
\[
 (\varphi,p_1):(R,X)\longrightarrow
 \bigl(S,X\times_{\operatorname{Spec}(R)}\operatorname{Spec}(S)\bigr)
\]
is cocartesian in $\underline M_{X_0}$. Hence
$\underline M_{X_0}$ is a cofibered category over
$\underline C_\Lambda$. We record the following facts.

%%%%%%%%%%  scan idx 107  |  folio 96  %%%%%%%%%%

\medskip
\noindent\textbf{\SGARefTarget{sga7.sga7i.evi.n4-1.item-a}{\SGARefAliasTarget{sga7.sga7i.evi.n4-1}4.1}(a).}
Let
\[
 (1_R,\Phi):(R,X_1)\longrightarrow(R,X_2)
\]
be an arrow in $\underline M_{X_0}$. Thus
$\Phi:X_2\longrightarrow X_1$ is an $R$-morphism whose reduction
\[
 X_2\otimes_Rk\longrightarrow X_1\otimes_Rk
\]
is an isomorphism. Since $R$ is artinian and the $X_i$ are flat over
$R$, $\Phi$ is an isomorphism. Consequently,
$\underline M_{X_0}$ is a cofibered category in groupoids over
$\underline C_\Lambda$.

\medskip
\noindent\textbf{\SGARefTarget{sga7.sga7i.evi.n4-1.item-b}{4.1}(b).}
Suppose that $X_0$ is affine, say
$X_0=\operatorname{Spec}(A_0)$. Then every
$(R,X)\in\operatorname{Ob}(\underline M_{X_0})$ has $X$ affine as
well. Consequently, for every
$R\in\operatorname{Ob}(\underline C_\Lambda)$, taking global
sections identifies $\underline M_{X_0}(R)$ with the following
category. An object is a flat $R$-algebra $A$ together with an
$R$-algebra map $j_A:A\longrightarrow A_0$ inducing an isomorphism
\[
 A\otimes_Rk\xrightarrow{\ \sim\ }A_0.
\]
A morphism $\Phi:A_1\longrightarrow A_2$ is an $R$-algebra map for
which the diagram
\[
\begin{tikzcd}[column sep=huge,row sep=huge]
A_1 \arrow[r,"\Phi"] \arrow[d,"j_{A_1}"']
& A_2 \arrow[d,"j_{A_2}"] \\
A_0 \arrow[r,equal] & A_0
\end{tikzcd}
\]
commutes.

\medskip
\noindent\textbf{\SGARefTarget{sga7.sga7i.evi.n4-1.item-c}{4.1}(c).}
Consider the cofibered category $\widehat{\underline M}_{X_0}$ in
groupoids over $\widehat{\underline C}_\Lambda$. By definition, for
every $R\in\operatorname{Ob}(\widehat{\underline C}_\Lambda)$, the
category $\widehat{\underline M}_{X_0}(R)$ is canonically equivalent
to the category of formal schemes $\mathfrak X$, adic over
$\operatorname{Spf}(R)$, such that
\[
 \mathfrak X_n
 =\mathfrak X\otimes_R R/\mathfrak m^{n+1}
\]
is flat over $R/\mathfrak m^{n+1}$ for every $n\geq0$. If $X_0$ is
noetherian, this is equivalent to requiring $\mathfrak X$ to be flat
over $\operatorname{Spf}(R)$. Thus, in that case,
$\widehat{\underline M}_{X_0}$ is canonically equivalent over
$\widehat{\underline C}_\Lambda$ to the category of deformations of
$X_0$ over the objects of $\widehat{\underline C}_\Lambda$, in the
sense of formal schemes.

\medskip
\noindent\textbf{\SGARefTarget{sga7.sga7i.evi.n4-1.item-d}{4.1}(d).}
Let $(R,X)\in\operatorname{Ob}(\underline M_{X_0})$, so that there is
a commutative diagram
\[
\begin{tikzcd}[column sep=huge,row sep=huge]
X_0 \arrow[r,"i_X"] \arrow[d] & X \arrow[d,"\mathrm{flat}"] \\
\operatorname{Spec}(k) \arrow[r] & \operatorname{Spec}(R)
\end{tikzcd}
\]
whose induced map
$X_0\longrightarrow X\times_{\operatorname{Spec}(R)}
\operatorname{Spec}(k)$ is an isomorphism.

%%%%%%%%%%  scan idx 108  |  folio 97  %%%%%%%%%%

Let $U_0\subset X_0$ be an open subscheme. Since
$i_X:X_0\longrightarrow X$ is a closed immersion and a homeomorphism
on underlying spaces, restriction gives a commutative diagram
\[
\begin{tikzcd}[column sep=huge,row sep=huge]
U_0 \arrow[r,"i_X|_{U_0}"] \arrow[d]
& X|_{U_0} \arrow[d,"\mathrm{flat}"] \\
\operatorname{Spec}(k) \arrow[r] & \operatorname{Spec}(R).
\end{tikzcd}
\]
The induced map
\[
 U_0\longrightarrow
 (X|_{U_0})\times_{\operatorname{Spec}(R)}\operatorname{Spec}(k)
 =(X\otimes_Rk)|_{U_0}
\]
is an isomorphism, so $(R,X|_{U_0})$ is a deformation of $U_0$ over
$R$. Moreover, every arrow
$(\varphi,\Phi):(R,X)\longrightarrow(R',X')$ in
$\underline M_{X_0}$ restricts to an arrow
\[
 (\varphi,\Phi|_{U_0}):(R,X|_{U_0})
 \longrightarrow(R',X'|_{U_0})
\]
in $\underline M_{U_0}$. We therefore have a restriction functor
\[
 \underline M_{X_0}\longrightarrow\underline M_{U_0}
\]
over $\underline C_\Lambda$.

One may likewise localize and complete at a point. Let
$x_0\in X_0$. If $(R,X)\in\operatorname{Ob}(\underline M_{X_0})$,
then $\operatorname{Spec}(\mathcal O_{X,x_0})$ is a deformation of
$\operatorname{Spec}(\mathcal O_{X_0,x_0})$ over $R$, while
$\operatorname{Spec}(\widehat{\mathcal O}_{X,x_0})$ is a deformation
of $\operatorname{Spec}(\widehat{\mathcal O}_{X_0,x_0})$ over $R$;
the hats denote completion for the topology defined by powers of the
maximal ideal. Thus there are functors
\[
 \underline M_{X_0}
 \longrightarrow
 \underline M_{\operatorname{Spec}(\mathcal O_{X_0,x_0})}
 \longrightarrow
 \underline M_{\operatorname{Spec}(\widehat{\mathcal O}_{X_0,x_0})}.
\]

We first prove that $\underline M_{X_0}$ is a homogeneous cofibered
category over $\underline C_\Lambda$.

% END INCLUDED SGA7I_English_current_buildable_source_20260823/components/38_expose_VI_section_4_opening_through_4_1.tex

% BEGIN INCLUDED SGA7I_English_current_buildable_source_20260823/components/39_expose_VI_lemma_4_2_and_corollary_4_3.tex
%%%%%%%%%%  scan idx 108  |  folio 97  %%%%%%%%%%

\medskip
\noindent\textbf{Lemma \SGARefTarget{sga7.sga7i.evi.n4-2}{4.2}.}
Consider a diagram
\[
\begin{tikzcd}[column sep=huge,row sep=large]
& R' \arrow[d] \\
R \arrow[r,"\mathrm{surjective}"'] & \bar R
\end{tikzcd}
\]
in $\underline C_\Lambda$, and let $E$ and $E'$ be flat modules over
$R$ and $R'$, respectively. Then, for every $R'$-homomorphism
\[
 E'\longrightarrow \bar E=E\otimes_R\bar R,
\]
the fiber product $E\times_{\bar E}E'$ is flat as a module over
$R\times_{\bar R}R'$.

%%%%%%%%%%  scan idx 109  |  folio 98  %%%%%%%%%%

\medskip
\noindent\textit{Proof.}
Let
\[
 0\longrightarrow I\longrightarrow R\longrightarrow\bar R
 \longrightarrow0
\]
be exact, so that
\[
 0\longrightarrow I\longrightarrow R\times_{\bar R}R'
 \longrightarrow R'\longrightarrow0
\]
is exact. Since $E$ is flat over $R$,
\[
 0\longrightarrow I\otimes_RE\longrightarrow E
 \longrightarrow\bar E\longrightarrow0
\]
is exact, and hence we obtain the exact sequence
\begin{equation}\tag{*}
0\longrightarrow I\otimes_RE\longrightarrow
E\times_{\bar E}E'\longrightarrow E'\longrightarrow0.
\end{equation}
Set $S=R\times_{\bar R}R'$. Since $I$ is nilpotent,
$E\times_{\bar E}E'$ is flat as an $S$-module if and only if
$(E\times_{\bar E}E')\otimes_SR'$ is flat as an $R'$-module and
\[
 \operatorname{Tor}_1^S(E\times_{\bar E}E',R')=0.
\]
But
\[
 I\otimes_RE=I\otimes_S(E\times_{\bar E}E'),
 \qquad
 E'=R'\otimes_S(E\times_{\bar E}E'),
\]
and therefore the exact sequence $(*)$ may be rewritten as
\[
0\longrightarrow I\otimes_S(E\times_{\bar E}E')
\longrightarrow E\times_{\bar E}E'
\longrightarrow R'\otimes_S(E\times_{\bar E}E')
\longrightarrow0.
\]
Thus
\[
 \operatorname{Tor}_1^S(R',E\times_{\bar E}E')=0,
 \qquad
 R'\otimes_S(E\times_{\bar E}E')=E',
\]
and $E'$ is flat over $R'$. It follows that
$E\times_{\bar E}E'$ is flat over $S$. \hfill$\square$

\medskip
\noindent\textbf{Corollary \SGARefTarget{sga7.sga7i.evi.n4-3}{4.3}.}
Given a diagram
\[
\begin{tikzcd}[column sep=huge,row sep=large]
(R,X) \arrow[dr,"{(\varphi,\Phi)}"']
&& (R',X') \arrow[dl,"{(\varphi',\Phi')}"] \\
& (\bar R,\bar X) &
\end{tikzcd}
\]
in $\underline M_{X_0}$, the amalgamated sum
$X\amalg_{\bar X}X'$ (in the category of schemes) is a flat scheme
over $R\times_{\bar R}R'$, provided that
$\varphi:R\longrightarrow\bar R$ is surjective. Consequently,
$\underline M_{X_0}$ is a homogeneous cofibered category in groupoids
over $\underline C_\Lambda$.

\medskip
\noindent\textit{Proof.}
For every affine open $U_0\subset X_0$, by definition we have
\[
 (X\amalg_{\bar X}X')|_{U_0}
 \simeq
 \operatorname{Spec}\!\left(
   \Gamma(U_0,X)
   \times_{\Gamma(U_0,\bar X)}
   \Gamma(U_0,X')
 \right).
\]
Consequently, if $\varphi:R\longrightarrow\bar R$ is surjective,
$X\amalg_{\bar X}X'$ is flat over $R\times_{\bar R}R'$ by Lemma
\SGARefLink{sga7.sga7i.evi.n4-2}{4.2}. Moreover, the morphism
\[
 i_X\amalg_{i_{\bar X}}i_{X'}:
 X_0\longrightarrow X\amalg_{\bar X}X'
\]
induces an isomorphism
\[
 X_0\longrightarrow (X\amalg_{\bar X}X')\otimes_Sk,
 \qquad S=R\times_{\bar R}R'.
\]
Therefore
\[
 \bigl(R\times_{\bar R}R',X\amalg_{\bar X}X'\bigr)
 \in\operatorname{Ob}(\underline M_{X_0}),
\]
and, in $\underline M_{X_0}$, we have
\[
 \bigl(R\times_{\bar R}R',X\amalg_{\bar X}X'\bigr)
 =
 (R,X)\times_{(\bar R,\bar X)}(R',X').
\]
Thus $\underline M_{X_0}$ is homogeneous over
$\underline C_\Lambda$. \hfill$\square$

%%%%%%%%%%  scan idx 110  |  folio 99  %%%%%%%%%%

% END INCLUDED SGA7I_English_current_buildable_source_20260823/components/39_expose_VI_lemma_4_2_and_corollary_4_3.tex

% BEGIN INCLUDED SGA7I_English_current_buildable_source_20260823/components/40_expose_VI_lemma_4_4.tex
%%%%%%%%%%  scan idx 110  |  folio 99  %%%%%%%%%%

\medskip
\noindent\textbf{Lemma \SGARefTarget{sga7.sga7i.evi.n4-4}{4.4}.}
For every scheme $X_0$ over $k$, there is a canonical exact sequence
\[
0\longrightarrow H^1(X_0,\underline D^0_{X_0})
\longrightarrow[\underline M_{X_0}(k[\varepsilon])]
\longrightarrow H^0(X_0,\underline D^1_{X_0}),
\]
where the $\underline D^i_{X_0}$ are as defined in section 3.

\medskip
\noindent\textit{Proof.}
First suppose that $X_0$ is affine, and put
$\Gamma(X_0,\mathcal O_{X_0})=A_0$. As noted in \SGARefLink{sga7.sga7i.evi.n4-1.item-b}{4.1}(b), the category
$\underline M_{X_0}(k[\varepsilon])$ is canonically equivalent to a
category $\underline E$ whose objects are flat $k[\varepsilon]$-algebras
$A$ equipped with a $k[\varepsilon]$-algebra map
$j_A:A\longrightarrow A_0$ inducing an isomorphism
\[
k\otimes_{k[\varepsilon]}A\xrightarrow{\ \sim\ }A_0.
\]
Equivalently, an object of $\underline E$ is a $k$-algebra $A$
equipped with an exact sequence
\[
0\longrightarrow A_0\xrightarrow{\ \varepsilon\ }A
\xrightarrow{\ j_A\ }A_0\longrightarrow0
\]
in which $\varepsilon(A_0)^2=0$ as an ideal of $A$.

An arrow $A_1\longrightarrow A_2$ in $\underline E$, in the first
description, is a $k[\varepsilon]$-algebra map
$\Phi:A_1\longrightarrow A_2$ such that
$j_{A_1}=j_{A_2}\circ\Phi$. In the second description it is a
$k$-algebra map $\Phi:A_1\longrightarrow A_2$ yielding a commutative
diagram
\[
\begin{tikzcd}[column sep=large,row sep=large]
0 \arrow[r]
& A_0 \arrow[r,"\varepsilon"] \arrow[d,equal]
& A_1 \arrow[r,"j_{A_1}"] \arrow[d,"\Phi"]
& A_0 \arrow[r] \arrow[d,equal]
& 0 \\
0 \arrow[r]
& A_0 \arrow[r]
& A_2 \arrow[r,"j_{A_2}"]
& A_0 \arrow[r]
& 0.
\end{tikzcd}
\]
It follows that, via $\Gamma(X_0,-)$,
$[\underline M_{X_0}(k[\varepsilon])]$ is canonically isomorphic to
the set of isomorphism classes of $k$-algebra extensions of $A_0$ by
$A_0$. In other words, there is a canonical bijection
\[
[\underline M_{X_0}(k[\varepsilon])]
\xrightarrow{\ \sim\ }H^0(X_0,\underline D^1_{X_0})
\]
which plainly respects the vector-space structures on both sides.

%%%%%%%%%%  scan idx 111  |  folio 100  %%%%%%%%%%

In general, for every affine open $U_0\subset X_0$, restriction and
the affine case give the composite map
\[
[\underline M_{X_0}(k[\varepsilon])]
\longrightarrow[\underline M_{U_0}(k[\varepsilon])]
\xrightarrow{\ \sim\ }H^0(U_0,\underline D^1_{X_0}).
\]
These maps give a canonical map
\[
[\underline M_{X_0}(k[\varepsilon])]
\longrightarrow H^0(X_0,\underline D^1_{X_0}).
\]
Its kernel consists of the isomorphism classes of locally trivial
deformations of $X_0$ over $k[\varepsilon]$, that is, of deformations
locally isomorphic to $X_0\otimes_k k[\varepsilon]$. The sheaf of
automorphisms of $X_0\otimes_k k[\varepsilon]$ inducing the identity
on $X_0$ is
\[
\underline D^0_{X_0}
=\underline{\operatorname{Der}}_k
  (\mathcal O_{X_0},\mathcal O_{X_0}).
\]
The isomorphism classes of such locally trivial deformations are
therefore canonically identified with
$H^1(X_0,\underline D^0_{X_0})$, which gives the asserted exact
sequence. \hfill$\square$

% END INCLUDED SGA7I_English_current_buildable_source_20260823/components/40_expose_VI_lemma_4_4.tex

% BEGIN INCLUDED SGA7I_English_current_buildable_source_20260823/components/41_expose_VI_theorem_4_5.tex
%%%%%%%%%%  scan idx 111  |  folio 100  %%%%%%%%%%

\medskip
\noindent\textbf{Theorem \SGARefTarget{sga7.sga7i.evi.n4-5}{4.5}.}
Let $X_0$ be a scheme of finite type over $k$ for which either

\smallskip
\noindent(i) $X_0$ is proper over $k$, or

\noindent(ii) $X_0$ is affine and
\[
X_0^{\mathrm{sing}}
=\{x\in X_0\mid X_0\text{ is not smooth over }k\text{ at }x\}
\]
is finite.

\smallskip
\noindent\textbf{(a)}
The category $\underline M_{X_0}$ admits a formal versal object over
$\underline C_\Lambda$, called a formal versal deformation of $X_0$.
If $(R,\mathcal X)$ is a formal versal deformation of $X_0$, then for
every $(S,\mathcal Y)\in\operatorname{Ob}(\widehat{\underline M}_{X_0})$
there is a map
\[
(\varphi,\Phi):(R,\mathcal X)\longrightarrow(S,\mathcal Y)
\]
in $\widehat{\underline M}_{X_0}$ such that the following diagram is
cartesian:
\[
\begin{tikzcd}[column sep=huge,row sep=large]
\mathcal Y \arrow[r,"\Phi"] \arrow[d]
& \mathcal X \arrow[d] \\
\operatorname{Spf}(S)
 \arrow[r,"\operatorname{Spf}(\varphi)"']
& \operatorname{Spf}(R).
\end{tikzcd}
\]
Furthermore, a formal versal deformation $(R,\mathcal X)$ of $X_0$
prorepresents the functor $[\underline M_{X_0}]$ if and only if, for
every surjective map $R'\longrightarrow R$ in $\underline C_\Lambda$
and every deformation $(R',X')$ of $X_0$ over $R'$, the homomorphism
\[
\operatorname{Aut}_{R'}(X'\mid X_0)
\longrightarrow
\operatorname{Aut}_R(X'\otimes_{R'}R\mid X_0)
\]
is surjective.

\smallskip
\noindent\textbf{(b)}
If $X_0$ is proper over $k$, then $\underline M_{X_0}$ is smoothly
pseudo-representable over $\underline C_\Lambda$.

%%%%%%%%%%  scan idx 112  |  folio 101  %%%%%%%%%%

\medskip
\noindent\textit{Proof.}
\textbf{(a)} By Corollary 4.3, $\underline M_{X_0}$ is a homogeneous
cofibered category in groupoids over $\underline C_\Lambda$. By \SGARefLink{sga7.sga7i.evi.n1-11}{1.11}
and \SGARefLink{sga7.sga7i.evi.n2-9}{2.9}(b), it therefore suffices to prove that
$[\underline M_{X_0}(k[\varepsilon])]$ is finite-dimensional over
$k$.

Consider the exact sequence
\begin{equation}\tag{*}
0\longrightarrow H^1(X_0,\underline D^0_{X_0})
\longrightarrow[\underline M_{X_0}(k[\varepsilon])]
\longrightarrow H^0(X_0,\underline D^1_{X_0}).
\end{equation}
If $X_0$ is proper over $k$, both
$H^1(X_0,\underline D^0_{X_0})$ and
$H^0(X_0,\underline D^1_{X_0})$ are finite-dimensional, since the
$\underline D^i_{X_0}$ are coherent for all $i\geq0$ by \SGARefLink{sga7.sga7i.evi.n2-13}{2.13}(a).

Now suppose that $X_0$ is affine and $X_0^{\mathrm{sing}}$ is finite.
The sheaf $\underline D^1_{X_0}$ is coherent and
\[
\operatorname{Supp}\underline D^1_{X_0}
\subset X_0^{\mathrm{sing}}
\]
by \SGARefLink{sga7.sga7i.evi.n2-13}{2.13}(a). Hence $H^0(X_0,\underline D^1_{X_0})$ is
finite-dimensional over $k$, whereas
$H^1(X_0,\underline D^0_{X_0})=0$ because $X_0$ is affine. In either
case, the exact sequence $(*)$ shows that
$[\underline M_{X_0}(k[\varepsilon])]$ is finite-dimensional over
$k$. The second assertion of part (a) follows from \SGARefLink{sga7.sga7i.evi.n2-7}{2.7}.

\smallskip
\noindent\textbf{(b)} If $X_0$ is proper, then
\[
\operatorname{Aut}(X_0\otimes_k k[\varepsilon]\mid X_0)
=H^0(X_0,\underline D^0_{X_0})
\]
is finite-dimensional, and the assertion follows from \SGARefLink{sga7.sga7i.evi.n2-13}{2.13}.
\hfill$\square$

% END INCLUDED SGA7I_English_current_buildable_source_20260823/components/41_expose_VI_theorem_4_5.tex

% BEGIN INCLUDED SGA7I_English_current_buildable_source_20260823/components/42_expose_VI_remarks_4_6_and_pair_deformations.tex
%%%%%%%%%%  scan idx 112  |  folio 101  %%%%%%%%%%

\medskip
\noindent\textbf{Remark \SGARefTarget{sga7.sga7i.evi.n4-6.item-a}{\SGARefAliasTarget{sga7.sga7i.evi.n4-6}4.6}(a).}
Let $X_0$ be affine. A formal versal deformation of $X_0$ exists if
and only if the support of $\underline D^1_{X_0}$ is a finite set.
This does not imply, however, that all nonsmooth points of $X_0$ are
isolated: there are affine schemes $X_0$ over $k$ having a
nonisolated nonsmooth point for which
$\operatorname{Supp}\underline D^1_{X_0}$ is finite.

\medskip
\noindent\textbf{Remark \SGARefTarget{sga7.sga7i.evi.n4-6.item-b}{4.6}(b).}
Henceforth, a deformation of $X_0$ means an object of
$\underline M_{X_0}$. Where confusion is possible, such an object
will be called an infinitesimal deformation of $X_0$.

\medskip
\noindent\textbf{Remark \SGARefTarget{sga7.sga7i.evi.n4-6.item-c}{4.6}(c).}
Let $X_0$ be a scheme over $k$. We have fixed a complete noetherian
local ring $\Lambda$ with residue field $k$ and considered
deformations of $X_0$ over $\underline C_\Lambda$.

%%%%%%%%%%  scan idx 113  |  folio 102  %%%%%%%%%%

In practice one usually takes $\Lambda=k$, or $\Lambda=W(k)$ when
$k$ is a perfect field of positive characteristic. The inclusion
$\underline C_k\subset\underline C_\Lambda$ is fully faithful, and
$\underline M_{X_0}|_{\underline C_k}$ is simply the category of
deformations of $X_0$ over the objects of $\underline C_k$.
Consequently, if $(R,\mathfrak X)$ is a formal versal deformation of
$X_0$ over $\underline C_\Lambda$, then
\[
 (k\otimes_\Lambda R,\;k\otimes_\Lambda\mathfrak X)
\]
is a formal versal deformation of $X_0$ over $\underline C_k$. It
follows, for example, that
\[
 \dim_{\underline C_\Lambda}\underline M_{X_0}
 =\dim(k\otimes_\Lambda R)
\]
does not depend on the choice of $\Lambda$, but only on $X_0$.

The preceding theorem admits many generalizations. For example,
consider a pair $Y_0\longrightarrow X_0$, where $X_0$ is a scheme of
finite type over $k$ and $Y_0$ is a fixed closed subscheme of $X_0$.
For $R\in\operatorname{Ob}(\underline C_\Lambda)$, a deformation of
$(X_0,Y_0)$ over $R$ is a commutative diagram
\[
\begin{tikzcd}[row sep=large,column sep=huge]
Y_0 \arrow[rrr] \arrow[dr,hook] \arrow[ddr]
&&& Y \arrow[dl,hook] \arrow[ddl,"\mathrm{flat}"] \\
& X_0 \arrow[r] \arrow[d]
& X \arrow[d,"\mathrm{flat}"'] & \\
& \operatorname{Spec}(k) \arrow[r]
& \operatorname{Spec}(R) &
\end{tikzcd}
\]
such that the induced map
\[
(X_0,Y_0)\longrightarrow
\left(
 X\times_{\operatorname{Spec}(R)}\operatorname{Spec}(k),
 Y\times_{\operatorname{Spec}(R)}\operatorname{Spec}(k)
\right)
\]
is an isomorphism. Defining morphisms
\[
(R';X',Y')\longrightarrow(R;X,Y)
\]
of deformations of $(X_0,Y_0)$ in the evident way gives a cofibered
category $\underline M_{X_0,Y_0}$ in groupoids over
$\underline C_\Lambda$. This construction is frequently used to
rigidify automorphisms in deformation-moduli problems. The argument
of Corollary \SGARefLink{sga7.sga7i.evi.n4-3}{4.3}, based on Lemma \SGARefLink{sga7.sga7i.evi.n4-2}{4.2}, immediately shows that
$\underline M_{X_0,Y_0}$ is a homogeneous cofibered category in
groupoids over $\underline C_\Lambda$.

% END INCLUDED SGA7I_English_current_buildable_source_20260823/components/42_expose_VI_remarks_4_6_and_pair_deformations.tex

% BEGIN INCLUDED SGA7I_English_current_buildable_source_20260823/components/43_expose_VI_lemma_4_7_through_theorem_4_9.tex
%%%%%%%%%%  scan idx 113  |  folio 102  %%%%%%%%%%

\medskip
\noindent\textbf{Lemma \SGARefTarget{sga7.sga7i.evi.n4-7}{4.7}.}
The sequence
\[
0\longrightarrow
H^0\!\left(
Y_0,
\underline{\operatorname{Hom}}_{\mathcal O_{Y_0}}
(\underline I/\underline I^2,\mathcal O_{Y_0})
\right)
\longrightarrow[\underline M_{X_0,Y_0}(k[\varepsilon])]
\longrightarrow[\underline M_{X_0}(k[\varepsilon])]
\]
is exact, where $\underline I$ is the ideal sheaf defining $Y_0$ in
$X_0$; thus
\[
0\longrightarrow\underline I\longrightarrow\mathcal O_{X_0}
\longrightarrow\mathcal O_{Y_0}\longrightarrow0
\]
is exact.

%%%%%%%%%%  scan idx 114  |  folio 103  %%%%%%%%%%

\medskip
\noindent\textit{Proof.}
Let $\underline B$ be the full subcategory of
$\underline M_{X_0,Y_0}(k[\varepsilon])$ whose objects have the form
\[
\bigl(k[\varepsilon];
X_0\times_{\operatorname{Spec}(k)}\operatorname{Spec}(k[\varepsilon]),
Y\bigr).
\]
Then
\[
[\underline B]
=\ker\!\left(
[\underline M_{X_0,Y_0}(k[\varepsilon])]
\longrightarrow[\underline M_{X_0}(k[\varepsilon])]
\right).
\]
The category $\underline B$ is the groupoid of deformations of $Y_0$
over $k[\varepsilon]$ inside
$X_0\times_{\operatorname{Spec}(k)}\operatorname{Spec}(k[\varepsilon])$.
Thus an object is a commutative diagram
\[
\begin{tikzcd}[column sep=huge,row sep=large]
Y \arrow[rr,"i_Y"] \arrow[dr,"\mathrm{flat}"']
&& X_0\times_{\operatorname{Spec}(k)}
   \operatorname{Spec}(k[\varepsilon]) \arrow[dl] \\
& \operatorname{Spec}(k[\varepsilon]) &
\end{tikzcd}
\]
whose special fiber is the fixed closed immersion $i_{Y_0}$. An
arrow $(Y_1,i_{Y_1})\longrightarrow(Y_2,i_{Y_2})$ is a
$k[\varepsilon]$-morphism $\Phi:Y_1\longrightarrow Y_2$ (necessarily
an isomorphism) whose special fiber is $\operatorname{id}_{Y_0}$ and
for which $i_{Y_2}\circ\Phi=i_{Y_1}$.

Equivalently, an object of $\underline B$ is a pair $(Y,j_Y)$, where
$Y$ is a deformation of $Y_0$ over $k[\varepsilon]$ and
$j_Y:Y\longrightarrow X_0$ is a $k$-morphism whose composite
$Y_0\longrightarrow Y\xrightarrow{j_Y}X_0$ is $i_{Y_0}$. An arrow
$(Y_1,j_{Y_1})\longrightarrow(Y_2,j_{Y_2})$ is a morphism
$\Phi:Y_1\longrightarrow Y_2$ whose special fiber is
$\operatorname{id}_{Y_0}$ and for which
$j_{Y_2}\circ\Phi=j_{Y_1}$.

The assignment sending $(Y,j_Y)$ to the diagram
\[
\begin{tikzcd}[column sep=large,row sep=large]
0 \arrow[r]
& \underline I \arrow[r] \arrow[d]
& \mathcal O_{X_0} \arrow[r] \arrow[d,"j_Y"]
& \mathcal O_{Y_0} \arrow[r] \arrow[d,equal]
& 0 \\
0 \arrow[r]
& \mathcal O_{Y_0} \arrow[r,"\varepsilon"]
& \mathcal O_Y \arrow[r]
& \mathcal O_{Y_0} \arrow[r]
& 0
\end{tikzcd}
\]
gives an isomorphism
\[
\begin{aligned}
[\underline B]
&\xrightarrow{\ j_Y|_{\underline I}\ }
\operatorname{Hom}_{\mathcal O_{X_0}}
 (\underline I,\mathcal O_{Y_0}) \\
&=
\operatorname{Hom}_{\mathcal O_{Y_0}}
 (\underline I/\underline I^2,\mathcal O_{Y_0}) \\
&=
H^0\!\left(
Y_0,
\underline{\operatorname{Hom}}_{\mathcal O_{Y_0}}
(\underline I/\underline I^2,\mathcal O_{Y_0})
\right).
\end{aligned}
\]
This proves exactness. \hfill$\square$

%%%%%%%%%%  scan idx 115  |  folio 104  %%%%%%%%%%

\medskip
\noindent\textbf{Remark \SGARefTarget{sga7.sga7i.evi.n4-8}{4.8}.}
For a more general formula implying Lemma \SGARefLink{sga7.sga7i.evi.n4-7}{4.7}, see SGA 3, III, 4.4.
One may also note the following. For every $S$-scheme $X$ and every
quasi-coherent $\mathcal O_X$-module $E$, one can define
quasi-coherent $\mathcal O_X$-modules $\underline D^i_{X|S}(E)$ for
all $i\geq0$, generalizing the case in which $S$ is affine. As usual,
we write
\[
\underline D^i_{X|S}(\mathcal O_X)=\underline D^i_{X|S}.
\]
For the closed immersion $Y_0\hookrightarrow X_0$, we then have
\[
H^0\!\left(
Y_0,
\underline{\operatorname{Hom}}_{\mathcal O_{Y_0}}
(\underline I/\underline I^2,\mathcal O_{Y_0})
\right)
=H^0(Y_0,\underline D^1_{Y_0|X_0}),
\]
where
$0\longrightarrow\underline I\longrightarrow\mathcal O_{X_0}
\longrightarrow\mathcal O_{Y_0}\longrightarrow0$ is exact.

\medskip
\noindent\textbf{Theorem \SGARefTarget{sga7.sga7i.evi.n4-9}{4.9}.}
Let $X_0$ be a $k$-scheme of finite type and $Y_0$ a closed
subscheme of $X_0$. Suppose either

\smallskip
\noindent(i) $X_0$ is proper over $k$, or

\noindent(ii) $X_0$ is affine and smooth over $k$ outside finitely
many points, and $Y_0$ is proper over $k$.

\smallskip
\noindent\textbf{(a)}
The category $\underline M_{X_0,Y_0}$ admits a formal versal
deformation over $\underline C_\Lambda$. It prorepresents the functor
$[\underline M_{X_0,Y_0}]$ if and only if, for every surjective map
$R'\longrightarrow R$ in $\underline C_\Lambda$ and every deformation
$(R';X',Y')$, the map
\[
\operatorname{Aut}_{R'}((X',Y')\mid X_0)
\longrightarrow
\operatorname{Aut}_{R}\!\left(
(X'\otimes_{R'}R,Y'\otimes_{R'}R)\mid X_0
\right)
\]
is surjective.

\smallskip
\noindent\textbf{(b)}
If, in addition,
\[
H^0\!\left(
X_0,
\underline D^0_{Y_0|k}
\times_{\underline D^0_{X_0|k}(\mathcal O_{Y_0})}
\underline D^0_{X_0|k}
\right)
\]
is finite-dimensional, then $\underline M_{X_0,Y_0}$ is smoothly
pseudo-representable.

\medskip
\noindent\textit{Proof.}
This follows immediately from Lemmas \SGARefLink{sga7.sga7i.evi.n4-4}{4.4} and \SGARefLink{sga7.sga7i.evi.n4-7}{4.7} together with \SGARefLink{sga7.sga7i.evi.n1-11}{1.11},
\SGARefLink{sga7.sga7i.evi.n2-7}{2.7}, and \SGARefLink{sga7.sga7i.evi.n2-13}{2.13}. \hfill$\square$

% END INCLUDED SGA7I_English_current_buildable_source_20260823/components/43_expose_VI_lemma_4_7_through_theorem_4_9.tex

% BEGIN INCLUDED SGA7I_English_current_buildable_source_20260823/components/44_expose_VI_localization_and_completion_intro.tex
%%%%%%%%%%  scan idx 115  |  folio 104  %%%%%%%%%%

We now consider localization and completion in the deformation theory
of schemes. We first fix notation. Let $X_0$ be a $k$-scheme of finite
type and let $x_0\in X_0$ be a fixed point. There are natural functors
\[
\underline M_{X_0}
\longrightarrow
\underline M_{\operatorname{Spec}(\mathcal O_{X_0,x_0})}
\longrightarrow
\underline M_{\operatorname{Spec}(\widehat{\mathcal O}_{X_0,x_0})},
\]

%%%%%%%%%%  scan idx 116  |  folio 105  %%%%%%%%%%

and hence functors
\[
\widehat{\underline M}_{X_0}
\longrightarrow
\widehat{\underline M}_{\operatorname{Spec}(\mathcal O_{X_0,x_0})}
\longrightarrow
\widehat{\underline M}_{\operatorname{Spec}
 (\widehat{\mathcal O}_{X_0,x_0})}.
\]
As observed in \SGARefLink{sga7.sga7i.evi.n4-1.item-c}{4.1}(c), $\widehat{\underline M}_{X_0}$ is canonically
equivalent to the category whose objects are pairs $(R,\mathfrak X)$,
where $R\in\operatorname{Ob}(\widehat{\underline C}_\Lambda)$ and
$\mathfrak X$ is a flat $R$-adic formal scheme equipped with a fixed
isomorphism $\mathfrak X_0\xrightarrow{\sim}X_0$.

When there is no danger of confusion, we denote the images of
$(R,\mathfrak X)$ under localization and completion by
$(R,\mathfrak X_{(x_0)})$ and
$(R,\mathfrak X_{(\widehat{x}_0)})$, respectively. Thus, if
\[
\mathfrak X=\varinjlim_\nu\mathfrak X_\nu,
\qquad
\mathfrak X_\nu
=\mathfrak X\times_{\operatorname{Spf}(R)}
 \operatorname{Spec}(R/\mathfrak m^{\nu+1}),
\]
then
\[
\mathfrak X_{(x_0)}
=\varinjlim_\nu
 \operatorname{Spec}(\mathcal O_{\mathfrak X_\nu,x_0})
\]
and
\[
\mathfrak X_{(\widehat{x}_0)}
=\varinjlim_\nu
 \operatorname{Spec}(\widehat{\mathcal O}_{\mathfrak X_\nu,x_0}).
\]

% END INCLUDED SGA7I_English_current_buildable_source_20260823/components/44_expose_VI_localization_and_completion_intro.tex

% BEGIN INCLUDED SGA7I_English_current_buildable_source_20260823/components/45_expose_VI_theorem_4_10.tex
%%%%%%%%%%  scan idx 116  |  folio 105  %%%%%%%%%%

\medskip
\noindent\textbf{Theorem \SGARefTarget{sga7.sga7i.evi.n4-10}{4.10}.}
Let $X_0$ be an affine scheme of finite type over $k$.

\smallskip
\noindent\textbf{(a)}
If $X_0$ is locally a complete intersection, then
$\underline M_{X_0}$ is smooth over $\underline C_\Lambda$.

\smallskip
\noindent\textbf{(b)}
Suppose that, for a fixed closed point $x\in X_0$, the open
$X_0-\{x\}$ is smooth over $k$. Then both functors
\[
\underline M_{X_0}
\longrightarrow
\underline M_{\operatorname{Spec}(\mathcal O_{X_0,x})}
\longrightarrow
\underline M_{\operatorname{Spec}(\widehat{\mathcal O}_{X_0,x})}
\]
are minimally smooth. In particular, if $(R,\mathfrak X)$ is a formal
versal deformation of $X_0$, then $(R,\mathfrak X_{(x)})$ and
$(R,\mathfrak X_{(\widehat x)})$ are formal versal deformations of
$\operatorname{Spec}(\mathcal O_{X_0,x})$ and
$\operatorname{Spec}(\widehat{\mathcal O}_{X_0,x})$, respectively.

\medskip
\noindent\textit{Proof.}
We use the notation for the cohomology of commutative algebras from
section 3, and write $X_0=\operatorname{Spec}(A_0)$.

%%%%%%%%%%  scan idx 117  |  folio 106  %%%%%%%%%%

\smallskip
\noindent\textbf{(a)}
It suffices to show that
$P:[\underline M_{X_0}]\longrightarrow\underline C_\Lambda$ is
smooth. Let $(R,\operatorname{Spec}(A))$ be a deformation of
$\operatorname{Spec}(A_0)$ over
$R\in\operatorname{Ob}(\underline C_\Lambda)$. Given a small
extension $R'\longrightarrow R$ in $\underline C_\Lambda$, so that
\[
0\longrightarrow k\longrightarrow R'\longrightarrow R
\longrightarrow0
\]
is exact, we must lift $(R,\operatorname{Spec}(A))$ to $R'$. The
transitivity sequence for $D^i$ with coefficients in $A_0$, applied
to $R'\longrightarrow R\longrightarrow A$, gives
\[
\cdots\longrightarrow D^1(A|R,A_0)
\longrightarrow D^1(A|R',A_0)
\longrightarrow D^1(R|R',k)\otimes_kA_0
\longrightarrow D^2(A|R,A_0)\longrightarrow\cdots .
\]
Because $R'\longrightarrow R$ is a small extension,
$D^1(R|R',k)$ is one-dimensional over $k$, generated by the class
$[R']$ represented by
$0\longrightarrow k\longrightarrow R'\longrightarrow R
\longrightarrow0$. Since $A$ is flat over $R$,
\[
D^i(A|R,A_0)=D^i(A_0|k,A_0)
\qquad(i>0).
\]
In particular, $D^2(A|R,A_0)=0$ because $A_0$ is locally a complete
intersection, by the low-degree description in \SGARefLink{sga7.sga7i.evi.n3-12}{3.12}(8). We therefore
obtain an exact sequence
\[
\longrightarrow D^1(A_0|k,A_0)
\longrightarrow D^1(A|R',A_0)
\longrightarrow D^1(R|R',k)\otimes_kA_0
\longrightarrow0.
\]
Hence there is an $R'$-algebra extension $A'$ of $A$ by $A_0$ whose
class maps to $[R']\otimes1$. Equivalently,
$(R,\operatorname{Spec}(A))$ lifts to the deformation
$(R',\operatorname{Spec}(A'))$.

\smallskip
\noindent\textbf{(b)}
By \SGARefLink{sga7.sga7i.evi.n3-14}{3.14} there are isomorphisms
\[
D^1(A_0|k,A_0)
\xrightarrow{\ \sim\ }
D^1((A_0)_x|k,(A_0)_x)
\xrightarrow{\ \sim\ }
D^1(\widehat{(A_0)_x}|k,\widehat{(A_0)_x}),
\]
and hence
\[
[\underline M_{X_0}(k[\varepsilon])]
\xrightarrow{\ \sim\ }
[\underline M_{\operatorname{Spec}(\mathcal O_{X_0,x})}
 (k[\varepsilon])]
\xrightarrow{\ \sim\ }
[\underline M_{\operatorname{Spec}(\widehat{\mathcal O}_{X_0,x})}
 (k[\varepsilon])].
\]
By \SGARefLink{sga7.sga7i.evi.n2-7}{2.7}(a), it remains to show that
\[
[\underline M_{X_0}]
\longrightarrow
[\underline M_{\operatorname{Spec}(\mathcal O_{X_0,x})}]
\longrightarrow
[\underline M_{\operatorname{Spec}(\widehat{\mathcal O}_{X_0,x})}]
\]
are smooth functors.

Let $(R,\operatorname{Spec}(A))$ be a deformation of
$X_0=\operatorname{Spec}(A_0)$, and let
$0\longrightarrow k\longrightarrow R'\longrightarrow R
\longrightarrow0$ be a small extension. As above, the homomorphisms
$R'\longrightarrow R\longrightarrow A$ give an exact commutative
diagram

%%%%%%%%%%  scan idx 118  |  folio 107  %%%%%%%%%%

\begin{center}
\scriptsize
\begin{tikzcd}[column sep=tiny,row sep=large]
{} \arrow[r]
& D^1(A_0|k,A_0) \arrow[r] \arrow[d,"\sim"']
& D^1(A|R',A_0) \arrow[r,"p"] \arrow[d]
& D^1(R|R',k)\otimes_kA_0 \arrow[r] \arrow[d]
& D^2(A_0|k,A_0) \arrow[r] \arrow[d,"\sim"]
& \cdots \\
{} \arrow[r]
& D^1(A_0|k,(A_0)_x) \arrow[r]
& D^1(A|R',(A_0)_x) \arrow[r]
& D^1(R|R',k)\otimes_k(A_0)_x \arrow[r]
& D^2(A_0|k,(A_0)_x) \arrow[r]
& \cdots \\
{} \arrow[r]
& D^1((A_0)_x|k,(A_0)_x) \arrow[r] \arrow[u]
& D^1(A_x|R',(A_0)_x) \arrow[r,"p_x"] \arrow[u]
& D^1(R|R',k)\otimes_k(A_0)_x \arrow[r] \arrow[u]
& D^2((A_0)_x|k,(A_0)_x) \arrow[r] \arrow[u]
& \cdots
\end{tikzcd}
\end{center}
The indicated isomorphisms follow from the fact that $A_0$ is smooth
over $k$ away from the single closed point $x$, so that
\[
D^i(A_0|k,(A_0)_x)
=D^i(A_0|k,A_0)_x
=D^i(A_0|k,A_0)
\qquad(i>0).
\]

Suppose we are given an arrow
\[
(R',\operatorname{Spec}(B))
\longrightarrow(R,\operatorname{Spec}(A_x))
\]
in $[\underline M_{\operatorname{Spec}((A_0)_x)}]$. This amounts to
a class $[B]\in D^1(A_x|R',(A_0)_x)$ satisfying
$p_x([B])=[R']$. A diagram chase gives a class
$[A']\in D^1(A|R',A_0)$ such that $p([A'])=[R']$ and $A'_x=B$.
Thus $(R',\operatorname{Spec}(A'))$ is a deformation of
$\operatorname{Spec}(A_0)$ over $R'$ whose localization at $x$ is
$\operatorname{Spec}(B)$. This proves that
\[
\underline M_{X_0}\longrightarrow
\underline M_{\operatorname{Spec}(\mathcal O_{X_0,x})}
\]
is smooth.

The smoothness of
\[
\underline M_{\operatorname{Spec}((A_0)_x)}
\longrightarrow
\underline M_{\operatorname{Spec}(\widehat{(A_0)_x})}
\]
follows from the analogous diagram and the identities
\[
D^i((A_0)_x|k,\widehat{(A_0)_x})
=\widehat{D^i((A_0)_x|k,(A_0)_x)}
=D^i((A_0)_x|k,(A_0)_x)
\qquad(i>0),
\]
since $x$ is the only possible nonsmooth point. \hfill$\square$

% END INCLUDED SGA7I_English_current_buildable_source_20260823/components/45_expose_VI_theorem_4_10.tex

% BEGIN INCLUDED SGA7I_English_current_buildable_source_20260823/components/46_expose_VI_theorem_4_11_and_remark_4_12.tex
%%%%%%%%%%  scan idx 118  |  folio 107  %%%%%%%%%%

\medskip
\noindent\textbf{Theorem \SGARefTarget{sga7.sga7i.evi.n4-11}{4.11}.}
Let $X_0$ be a separated scheme of finite type over $k$, smooth away
from a finite closed subset
\[
X_0^{\mathrm{Sing}}=\{x_1,\ldots,x_r\}.
\]
Let $U_0^1,\ldots,U_0^r$ be affine open neighbourhoods of
$x_1,\ldots,x_r$ such that
$U_0^i\cap X_0^{\mathrm{Sing}}=\{x_i\}$. If
$H^2(X_0,\underline D^0_{X_0})=0$, then the natural functor
\[
\underline M_{X_0}\longrightarrow
\underline M_{U_0^1}\times\cdots\times\underline M_{U_0^r}
\]
is smooth.

\smallskip
\noindent\textit{Proof.}
Complete $U_0^1,\ldots,U_0^r$ to an affine open covering of $X_0$
by choosing $U_0^{r+1},\ldots,U_0^N$ such that
$U_0^j\cap X_0^{\mathrm{Sing}}=\varnothing$ for every $j>r$.
Let $R'\longrightarrow R$ be a small extension in
$\underline C_\Lambda$, and let
$(R',U'^1),\ldots,(R',U'^r)$ be deformations of
$U_0^1,\ldots,U_0^r$ over $R'$ such that
\[
U'^i\otimes_{R'}R=X|_{U_0^i},\qquad 1\leq i\leq r,
\]
for some deformation $(R,X)$ of $X_0$ over $R$. We must find a
deformation $(R',X')$ of $X_0$ over $R'$ such that
\[
X'\otimes_{R'}R=X,
\qquad
X'|_{U_0^i}=U'^i\quad(1\leq i\leq r).
\]

%%%%%%%%%%  scan idx 119  |  folio 108  %%%%%%%%%%

For $j>r$, the scheme $U_0^j$ is smooth over $k$; hence, by \SGARefLink{sga7.sga7i.evi.n4-9}{4.9}(a),
$X|_{U_0^j}$ lifts to a deformation $(R',U'^j)$ of $U_0^j$ over
$R'$. Thus
\[
U'^j\otimes_{R'}R=X|_{U_0^j}\qquad(j>r).
\]
For every pair $i\neq j$, the intersection $U_0^i\cap U_0^j$ is
affine and smooth over $k$. We therefore obtain an isomorphism
\[
\tau_{ij}:U'^i|_{U_0^i\cap U_0^j}
\xrightarrow{\ \sim\ }
U'^j|_{U_0^i\cap U_0^j}
\]
such that
$\tau_{ij}\otimes_{R'}R=\operatorname{id}_{X|_{U_0^i\cap U_0^j}}$.
On triple intersections, the automorphisms
\[
\tau_{ij}\tau_{jk}\tau_{ik}^{-1}
\]
define a cohomology class in
$H^2(X_0,\underline D^0_{X_0})$. This class is the obstruction to
gluing the $U'^i$ to a global deformation. It vanishes by hypothesis,
so the $U'^i$ glue to a deformation $(R',X')$ with
$X'\otimes_{R'}R=X$ and $X'|_{U_0^i}\simeq U'^i$. \hfill$\square$

\medskip
\noindent\textbf{Remark \SGARefTarget{sga7.sga7i.evi.n4-12}{4.12}.}
The preceding proof assumes that $X_0$ is separated, but this
assumption is unnecessary. Indeed, fix a deformation $(R,X)$ of
$X_0$ and deformations $(R',U'^i)$ of the members $U_0^i$ of an open
covering. The deformations $(R',V')$ of open subsets $V$ of $X$ over
$R'$ that are compatible with the already-given deformations
$(R',U'^i)$ form a ``gerbe'' (in Giraud's terminology) whose
structure sheaf is $\underline D^0_{X_0}$. The obstruction to the
existence of a section of this gerbe---that is, to a global
deformation $(R',X')$ of $X$ compatible with the given data
$(R',U'^i)$---lies in $H^2(X_0,\underline D^0_{X_0})$, and hence
vanishes under the hypothesis of the theorem.


% END INCLUDED SGA7I_English_current_buildable_source_20260823/components/46_expose_VI_theorem_4_11_and_remark_4_12.tex

% BEGIN INCLUDED SGA7I_English_current_buildable_source_20260823/components/47_expose_VI_corollary_4_13.tex
\medskip
\noindent\textbf{Corollary \SGARefTarget{sga7.sga7i.evi.n4-13}{4.13}.}
Let $X_0$ be a scheme of finite type over $k$, smooth away from a
finite subset
\[
X_0^{\mathrm{Sing}}=\{x_1,\ldots,x_r\},
\]
and suppose that $H^2(X_0,\underline D^0_{X_0})=0$. Then:

\smallskip
\noindent\textbf{(a)} The functor
\[
\underline M_{X_0}\longrightarrow
\prod_{i=1}^r
\underline M_{\operatorname{Spec}(\widehat{\mathcal O}_{X_0,x_i})}
\]
is smooth. In particular,
\[
\dim\underline M_{X_0}
=\dim_k H^1(X_0,\underline D^0_{X_0})
+\sum_{i=1}^r
\dim\underline M_{\operatorname{Spec}(\widehat{\mathcal O}_{X_0,x_i})}.
\]

\smallskip
\noindent\textbf{(b)} If, in addition, $X_0$ is locally a complete
intersection, then $\underline M_{X_0}$ is smooth.

%%%%%%%%%%  scan idx 120  |  folio 109  %%%%%%%%%%

\smallskip
\noindent\textit{Proof.}
\textbf{(a)} The first assertion follows from \SGARefLink{sga7.sga7i.evi.n4-10}{4.10} and \SGARefLink{sga7.sga7i.evi.n4-9}{4.9}(b).
Let $(R,\mathfrak X)$ and
$(R_1,\mathfrak X_1),\ldots,(R_r,\mathfrak X_r)$ be formal versal
deformations of $X_0$ and of
\[
\operatorname{Spec}(\mathcal O_{X_0,x_1}),\ldots,
\operatorname{Spec}(\mathcal O_{X_0,x_r}),
\]
respectively. We obtain a commutative diagram
\[
\begin{tikzcd}[column sep=large,row sep=large]
\underline h_R
  \arrow[r]
  \arrow[d]
& \underline h_{R_1\widehat\otimes_\Lambda\cdots
  \widehat\otimes_\Lambda R_r}
  \arrow[d] \\
\underline M_{X_0}
  \arrow[r]
& \displaystyle\prod_{i=1}^r
  \underline M_{\operatorname{Spec}(\widehat{\mathcal O}_{X_0,x_i})}.
\end{tikzcd}
\]
The composite functor
\[
\underline h_R\longrightarrow\underline M_{X_0}
\longrightarrow
\prod_{i=1}^r
\underline M_{\operatorname{Spec}(\widehat{\mathcal O}_{X_0,x_i})}
\]
is smooth, whereas
\[
\underline h_{R_1\widehat\otimes_\Lambda\cdots
\widehat\otimes_\Lambda R_r}
\longrightarrow
\prod_{i=1}^r
\underline M_{\operatorname{Spec}(\widehat{\mathcal O}_{X_0,x_i})}
\]
is minimally smooth. It follows from \SGARefLink{sga7.sga7i.evi.n2-7}{2.7}(a) that
\[
\underline h_R\longrightarrow
\underline h_{R_1\widehat\otimes_\Lambda\cdots
\widehat\otimes_\Lambda R_r}
\]
is smooth; equivalently, $R$ is a formal power-series ring over
$R_1\widehat\otimes_\Lambda\cdots\widehat\otimes_\Lambda R_r$.

By \SGARefLink{sga7.sga7i.evi.n4-4}{4.4}, the sequence
\[
0\longrightarrow H^1(X_0,\underline D^0_{X_0})
\longrightarrow[\underline M_{X_0}(k[\varepsilon])]
\longrightarrow
\prod_{i=1}^r
[\underline M_{\operatorname{Spec}(\widehat{\mathcal O}_{X_0,x_i})}
  (k[\varepsilon])]
\]
is exact. Equivalently,
\[
0\longrightarrow H^1(X_0,\underline D^0_{X_0})
\longrightarrow\operatorname{Hom}_{\Lambda\text{-alg}}(R,k[\varepsilon])
\longrightarrow
\operatorname{Hom}_{\Lambda\text{-alg}}
 (R_1\widehat\otimes_\Lambda\cdots\widehat\otimes_\Lambda R_r,
  k[\varepsilon])
\]
is exact. Consequently, $R$ is a formal power-series ring over
$R_1\widehat\otimes_\Lambda\cdots\widehat\otimes_\Lambda R_r$ in
$d$ variables, where
\[
d=\dim_k H^1(X_0,\underline D^0_{X_0}).
\]
Therefore
\begin{align*}
\dim\underline M_{X_0}
 &=\dim(k\otimes_\Lambda R) \\
 &=d+\dim\bigl(k\otimes_\Lambda
   (R_1\widehat\otimes_\Lambda\cdots\widehat\otimes_\Lambda R_r)
   \bigr) \\
 &=d+\dim\bigl((k\otimes_\Lambda R_1)
   \widehat\otimes_k\cdots\widehat\otimes_k
   (k\otimes_\Lambda R_r)\bigr) \\
 &=d+\sum_{i=1}^r\dim(k\otimes_\Lambda R_i) \\
 &=d+\sum_{i=1}^r
   \dim\underline M_{\operatorname{Spec}
   (\widehat{\mathcal O}_{X_0,x_i})}.
\end{align*}

%%%%%%%%%%  scan idx 121  |  folio 110  %%%%%%%%%%

\noindent\textbf{(b)} This follows from \SGARefLink{sga7.sga7i.evi.n4-10}{4.10} and \SGARefLink{sga7.sga7i.evi.n4-9}{4.9}(a).
\hfill$\square$


% END INCLUDED SGA7I_English_current_buildable_source_20260823/components/47_expose_VI_corollary_4_13.tex

% BEGIN INCLUDED SGA7I_English_current_buildable_source_20260823/components/48_expose_VI_4_14_and_4_15.tex
\medskip
\noindent\textbf{\SGARefTarget{sga7.sga7i.evi.n4-14}{4.14}.}
We close this section by constructing explicitly a formal versal
deformation in a simple case. When $X_0$ is affine and a complete
intersection, such a deformation is immediate to write down. Let
\[
A_0=k[x_1,\ldots,x_n]/(f_1,\ldots,f_r)=S/J,
\]
where $S=k[x_1,\ldots,x_n]$, $J=(f_1,\ldots,f_r)$, and
$f_1,\ldots,f_r$ is an $S$-regular sequence. We have an exact
sequence
\[
\cdots\longrightarrow\operatorname{Der}_k(S,A_0)
\longrightarrow\operatorname{Hom}_{A_0}(J/J^2,A_0)
\longrightarrow D^1(A_0|k,A_0)\longrightarrow0.
\]
The $A_0$-module $J/J^2$ is free of rank $r$, generated by the
classes of $f_1,\ldots,f_r$, because these elements form an
$S$-regular sequence.

Assume now that $D^1(A_0|k,A_0)$ is finite-dimensional over $k$;
for example, assume that $A_0$ has only isolated nonsmooth points.
Put
\[
d=\dim_kD^1(A_0|k,A_0),
\]
and choose maps
\[
p_i:J/J^2\longrightarrow A_0\qquad(1\leq i\leq d)
\]
whose classes form a $k$-basis of $D^1(A_0|k,A_0)$. For
$1\leq i\leq d$ and $1\leq j\leq r$, choose
$P_{ij}\in\Lambda[x_1,\ldots,x_n]$ such that
\[
P_{ij}\equiv p_i(f_j)\pmod{\mathfrak m_\Lambda+J},
\]
where $\mathfrak m_\Lambda$ is the maximal ideal of $\Lambda$. Set
\[
R=\Lambda[[t_1,\ldots,t_d]],
\qquad
B=R[x_1,\ldots,x_n]/(G_1,\ldots,G_r),
\]
where
\[
G_j=F_j-\sum_{i=1}^d t_iP_{ij},
\]
and $F_j\in\Lambda[x_1,\ldots,x_n]$ is obtained from $f_j$ by
lifting its coefficients from $k$ to $\Lambda$.

The algebra $R[x_1,\ldots,x_n]$ is flat over $R$, and the images of
$G_1,\ldots,G_r$ modulo $\mathfrak m_R$ form the regular sequence
$f_1,\ldots,f_r$ in
$k\otimes_RR[x_1,\ldots,x_n]$. Hence $B$ is $R$-flat and
$k\otimes_RB=A_0$. By construction, the map
\[
\operatorname{Hom}_{\Lambda\text{-alg}}(R,k[\varepsilon])
\longrightarrow D^1(A_0|k,A_0)
=[\underline M_{\operatorname{Spec}(A_0)}(k[\varepsilon])]
\]
is bijective. Thus, if

%%%%%%%%%%  scan idx 122  |  folio 111  %%%%%%%%%%

\[
\operatorname{Spf}(B)
=\varinjlim_\nu
\operatorname{Spec}\bigl((R/\mathfrak m_R^{\nu+1})\otimes_RB\bigr),
\]
then $(R,\operatorname{Spf}(B))$ is an object of
$\widehat{\underline M}_{\operatorname{Spec}(A_0)}$. The functor
\[
\underline h_R\longrightarrow\underline M_{\operatorname{Spec}(A_0)}
\]
that it induces is smooth by \SGARefLink{sga7.sga7i.evi.n1-6}{1.6}, since $R$ is a formal power-series
ring over $\Lambda$. In other words,
$(R,\operatorname{Spf}(B))$ is a formal versal object of
$\underline M_{\operatorname{Spec}(A_0)}$.

By \SGARefLink{sga7.sga7i.evi.n4-10}{4.10}(b), if instead we put
\[
B=R[[x_1,\ldots,x_n]]/(G_1,\ldots,G_r),
\]
then $(R,\operatorname{Spf}(B))$ is a formal versal deformation of
$\operatorname{Spec}(A_0)$ for
\[
A_0=k[[x_1,\ldots,x_n]]/(f_1,\ldots,f_r).
\]

\medskip
\noindent\textbf{Remark \SGARefTarget{sga7.sga7i.evi.n4-15}{4.15}.}
Let $X_0$ be a scheme of finite type over $k$ with only isolated
nonsmooth points, and let $(R,\mathfrak X)$ be a formal versal
deformation of $X_0$. The preceding construction shows that, when
$X_0$ is affine and locally a complete intersection, $\mathfrak X$
is the formalization of an $R$-scheme of finite type. The same fact
extends readily to the nonaffine case provided that
$H^2(X_0,\underline D^0_{X_0})=0$.


% END INCLUDED SGA7I_English_current_buildable_source_20260823/components/48_expose_VI_4_14_and_4_15.tex

% BEGIN INCLUDED SGA7I_English_current_buildable_source_20260823/components/49_expose_VI_remark_4_16.tex
\medskip
\noindent\textbf{Remark \SGARefTarget{sga7.sga7i.evi.n4-16}{4.16}.}
The global theory of the relative cotangent complex---here reduced to
the truncated complex of length one, which nevertheless suffices for
many purposes in deformation theory---generalizes and simplifies
results such as \SGARefLink{sga7.sga7i.evi.n4-10}{4.10}. Let $S$ be a scheme, and let
$S_0\subset S'\subset S$ be subschemes defined by ideals
$J\supset I$ with $JI=0$. Let $X'$ be flat over $S'$, and write
\[
L_\bullet=L^{X_0/S_0}
\]
for the relative cotangent complex of $X_0$ over $S_0$, regarded as
an object of $D(X_0,\mathcal O_{X_0})$. Let $\underline F$ be the
category of flat $S$-schemes $X$ extending $X'/S'$; thus an object is
equipped with an $S'$-isomorphism
\[
X\times_SS'\xrightarrow{\ \sim\ }X'.
\]
Then:

\smallskip
\noindent\textbf{(a)} For an object $X$ of $\underline F$, its group
of automorphisms inducing the identity on $X'$ is canonically
isomorphic to
\[
\operatorname{Ext}^0(X_0;L_\bullet,I_{X_0}),
\]
where $I_{X_0}$ is the inverse image on $X_0$ of the
$\mathcal O_{S_0}$-module $I$.

%%%%%%%%%%  scan idx 123  |  folio 112  %%%%%%%%%%

\smallskip
\noindent\textbf{(b)} The set of isomorphism classes of objects of
$\underline F$ is either empty or a torsor (a principal homogeneous
space) under
$\operatorname{Ext}^1(X_0;L_\bullet,I_{X_0})$.

\smallskip
\noindent\textbf{(c)} There is a canonical element
\[
c\in\operatorname{Ext}^2(X_0;L_\bullet,I_{X_0})
\]
whose vanishing is necessary and sufficient for $\underline F$ to be
nonempty.

To study these groups, it is often convenient to write
\[
\operatorname{Ext}^i(X_0;L_\bullet,I_{X_0})
=H^i\bigl(X_0,R\underline{\operatorname{Hom}}
  (L_\bullet,I_{X_0})\bigr).
\]
The cohomology sheaves of
$R\underline{\operatorname{Hom}}(L_\bullet,I_{X_0})$ are precisely
the sheaves $\underline D^i_{X_0}(I_{X_0})$ introduced in section 3.
For $i=0,1,2$, they may be viewed respectively as the sheaf of
automorphisms of an object $X$ of $\underline F$ (automorphisms that
induce the identity on $X'$), the sheaf measuring the indeterminacy
of isomorphism classes of local solutions to the flat-deformation
problem, and the sheaf of obstructions to the existence of local
solutions.

The standard spectral sequence is
\[
E_2^{p,q}
=H^p\bigl(X_0,\underline D^q_{X_0}(I_{X_0})\bigr)
\Longrightarrow
\operatorname{Ext}^{p+q}(X_0;L_\bullet,I_{X_0}).
\]
It gives the familiar isomorphisms
\begin{align*}
\operatorname{Aut}(X/X')
&\simeq H^0\bigl(X_0,\underline D^0_{X_0}(I_{X_0})\bigr)\\
&\simeq H^0\bigl(X_0,
  \underline{\operatorname{Hom}}
  (\Omega^1_{X_0/S_0},I_{X_0})\bigr)\\
&=\operatorname{Hom}(\Omega^1_{X_0/S_0},I_{X_0}),
\end{align*}
and the low-degree exact sequence
\[
0\longrightarrow H^1(X_0,\underline D^0_{X_0}(I_{X_0}))
\longrightarrow\operatorname{Ext}^1
\longrightarrow H^0(X_0,\underline D^1_{X_0}(I_{X_0}))
\longrightarrow H^2(X_0,\underline D^0_{X_0}(I_{X_0}))
\longrightarrow\operatorname{Ext}^2.
\]
In particular, if
$H^2(X_0,\underline D^0_{X_0}(I_{X_0}))=0$ and local flat
deformations of $X'/S'$ exist, then compatible local deformations
glue to a global flat deformation. More generally, for a prescribed
section of the sheaf of isomorphism classes of local flat
deformations, the argument of \SGARefLink{sga7.sga7i.evi.n4-10}{4.10} places the obstruction to a
compatible global deformation in
$H^2(X_0,\underline D^0_{X_0}(I_{X_0}))$.

%%%%%%%%%%  scan idx 124  |  folio 113  %%%%%%%%%%

The same spectral sequence shows that $\operatorname{Ext}^2$ vanishes
whenever
\begin{equation}\tag{*}
H^0(X_0,\underline D^2_{X_0})=0,
\qquad
H^1(X_0,\underline D^1_{X_0})=0,
\qquad
H^2(X_0,\underline D^0_{X_0})=0.
\end{equation}
The first equality always holds when $X_0$ is a relative complete
intersection over $S_0$, since then
$\underline D^2_{X_0}(M)=0$ for every $\mathcal O_{X_0}$-module $M$.
The second holds when the support of $\underline D^1_{X_0}$ is
discrete, as happens when $X_0$ is smooth over $S_0$ away from a
discrete subset. If all three equalities in $(*)$ hold, then
$\underline F$ is nonempty by (c), which generalizes \SGARefLink{sga7.sga7i.evi.n4-13}{4.13}(b).

There is also a geometric interpretation, requiring only the
globalized functors $\underline D^i_{X_0}$ rather than the global
theory of the relative cotangent complex. The three successive
obstructions to extending $X'/S'$ flatly over $S$ lie in the three
groups displayed in $(*)$. The local theory defines the first
obstruction
\[
c_1\in H^0(X_0,\underline D^2_{X_0})
\]
as the obstruction to the existence of local solutions. If $c_1=0$,
the sheaf $\underline F$ of isomorphism classes of local solutions is
canonically a torsor under $\underline D^1_{X_0}$. Its obstruction to
having a section is
\[
c_2\in H^1(X_0,\underline D^1_{X_0});
\]
such a section is a coherent system of isomorphism classes of local
solutions. If $c_2=0$, choose a section of $\underline F$. The
obstruction to a global solution realizing the chosen local classes
is then
\[
c_3\in H^2(X_0,\underline D^0_{X_0}).
\]
Changing the section of $\underline F$ by a section of
$\underline D^1_{X_0}$ gives a map
\[
d:H^0(X_0,\underline D^1_{X_0})
\longrightarrow H^2(X_0,\underline D^0_{X_0}).
\]
A solution exists if and only if $c_1=0$, $c_2=0$, and the class of
$c_3$ in $\operatorname{Coker}d$ is zero (each $c_i$ being defined
only after the preceding obstructions vanish).

%%%%%%%%%%  scan idx 125  |  folio 114  %%%%%%%%%%

The three obstruction groups
\[
H^0(X_0,\underline D^2_{X_0}),
\qquad
H^1(X_0,\underline D^1_{X_0}),
\qquad
H^2(X_0,\underline D^0_{X_0})/\operatorname{Im}d
\]
are the initial terms $E_2^{0,2}$, $E_2^{1,1}$, and $E_3^{2,0}$ of
the preceding spectral sequence; up to sign, $d$ is its differential
\[
d_2:E_2^{0,1}\longrightarrow E_2^{2,0}.
\]


% END INCLUDED SGA7I_English_current_buildable_source_20260823/components/49_expose_VI_remark_4_16.tex

% BEGIN INCLUDED SGA7I_English_current_buildable_source_20260823/components/50_expose_VI_section_5_opening_through_5_2.tex
\bigskip
\noindent\textbf{5. Formal Jacobian subschemes}

\medskip
Let $X_0$ be a scheme of finite type over a field $k$, and let
$(R,\mathfrak X)$ be a formal versal deformation of $X_0$. The space
$\mathfrak X$ is not an ordinary $R$-scheme, but an $R$-adic formal
scheme. The purpose of this section is to clarify the meaning of the
``nonsmooth locus of $\mathfrak X$ over $R$.''

We first recall the Fitting ideals of a finite-type module. Let $A$ be
a commutative ring and let $E$ be a finite-type $A$-module. Choose a
presentation
\[
K\longrightarrow F\longrightarrow E\longrightarrow0,
\]
where $F$ is locally free of constant rank $n$. For
$p=0,1,2,\ldots$, define
\[
I_A^p(E)=\operatorname{Im}\left(
 \bigwedge^{n-p}K\otimes_A\bigwedge^{n-p}F^*
 \longrightarrow A\right).
\]
These ideals are independent of the chosen presentation. They are
invariants of the $A$-module $E$, called its \emph{Fitting ideals}.
The following properties follow directly from the definition.

\medskip
\noindent\textbf{\SGARefTarget{sga7.sga7i.evi.n5-1}{5.1}.}
\textbf{(a)} For every ring homomorphism $A\longrightarrow B$, the
canonical map
\[
B\otimes_A A/I_A^p(E)
\xrightarrow{\ \sim\ }
B/I_B^p(B\otimes_AE)
\]
is an isomorphism for every $p$. In particular,
\[
I_{S^{-1}A}^p(S^{-1}E)=S^{-1}I_A^p(E)
\]
for every multiplicative subset $S\subset A$.

%%%%%%%%%%  scan idx 126  |  folio 115  %%%%%%%%%%

\smallskip
\noindent\textbf{(b)} There are inclusions
\[
I_A^0(E)\subset I_A^1(E)\subset I_A^2(E)\subset\cdots.
\]
For every $x\in\operatorname{Spec}(A)$,
\[
I_A^p(E)_x=A_x
\qquad\text{for all }p\geq\dim_{\kappa(x)}E(x),
\]
where $E(x)=\kappa(x)\otimes_AE$.

\smallskip
\noindent\textbf{(c)} A point
$x\in\operatorname{Spec}(A)-\operatorname{Supp}(A/I_A^p(E))$ if and
only if, for any (equivalently, for some) presentation
\[
K\longrightarrow F\longrightarrow E\longrightarrow0,
\qquad F\simeq A^n,
\]
the module $K_x$ contains a copy of $A_x^{n-p}$ that is a direct
summand of $F_x$. In particular,
\[
\operatorname{Supp}(A/I_A^0(E))=\operatorname{Supp}(E).
\]

\smallskip
\noindent\textbf{(d)} For any two finite-type $A$-modules,
\[
I_A^r(E_1\oplus E_2)
=\sum_{p+q=r}I_A^p(E_1)I_A^q(E_2).
\]

\smallskip
\noindent\textbf{(e)} If $A$ is a noetherian adic ring, then
\[
A/I_A^p(E)=\varprojlim_\nu A_\nu/I_{A_\nu}^p(E_\nu),
\]
where $A=\varprojlim_\nu A_\nu$ and $E_\nu=E\otimes_AA_\nu$.

Property \SGARefLink{sga7.sga7i.evi.n5-1}{5.1}(a) globalizes the definition. If $X$ is a scheme and
$\underline E$ is a quasi-coherent $\mathcal O_X$-module of finite
presentation, there are quasi-coherent ideal sheaves
$\underline I_X^p(\underline E)$, $p\geq0$, characterized on every
affine open $U\subset X$ by
\[
\Gamma(U,\underline I_X^p(\underline E))
\xrightarrow{\ \sim\ }
I_{\Gamma(U,\mathcal O_X)}^p(\Gamma(U,\underline E)).
\]

For a morphism of schemes $X\longrightarrow Y$, the relative module
of K\"ahler differentials
$\Omega_{X|Y}=\Omega^1_{X|Y}$ is quasi-coherent. If the morphism is
essentially of finite presentation, then $\Omega_{X|Y}$ is of finite
presentation, so the ideal sheaves
$\underline I_X^p(\Omega_{X|Y})$ are defined.

\medskip
\noindent\textbf{Definition \SGARefTarget{sga7.sga7i.evi.n5-2}{5.2}.}
Let $X\longrightarrow Y$ be a morphism essentially of finite
presentation. For every integer $p\geq0$, define $J_X^p(X|Y)$ to be
the closed subscheme of $X$

%%%%%%%%%%  scan idx 127  |  folio 116  %%%%%%%%%%

\noindent defined by the ideal sheaf
\[
\underline I_X^p(X|Y)=\underline I_X^p(\Omega_{X|Y}).
\]
It is called the $p$th \emph{Jacobian subscheme} of $X$ over $Y$.
For every $x\in X$, put
\[
X_x=\operatorname{Spec}(\mathcal O_{X,x}).
\]


% END INCLUDED SGA7I_English_current_buildable_source_20260823/components/50_expose_VI_section_5_opening_through_5_2.tex

% BEGIN INCLUDED SGA7I_English_current_buildable_source_20260823/components/51_expose_VI_theorem_5_3.tex
\medskip
\noindent\textbf{Theorem \SGARefTarget{sga7.sga7i.evi.n5-3}{5.3}.}
Let $X\longrightarrow Y$ be a morphism essentially of finite
presentation. Then:

\smallskip
\noindent\textbf{(1)} For every base change $Y'\longrightarrow Y$,
the canonical map
\[
J_{X'}^p(X'|Y')\xrightarrow{\ \sim\ }
J_X^p(X|Y)\times_YY'
\]
is an isomorphism for every $p$, where $X'=X\times_YY'$.

\smallskip
\noindent\textbf{(2)} For every $x\in X$, there is a canonical
isomorphism
\[
J_{X_x}^p(X_x|Y)\xrightarrow{\ \sim\ }J_X^p(X|Y)_x.
\]

\smallskip
\noindent\textbf{(3)} There are inclusions
\[
J_X^0(X|Y)\supset J_X^1(X|Y)\supset J_X^2(X|Y)\supset\cdots,
\]
and
\[
x\notin J_X^p(X|Y)
\qquad\text{for every }p\geq
\dim_{\kappa(x)}\Omega_{X|Y}(x).
\]

\smallskip
\noindent\textbf{(4)} A point $x$ does not belong to
$J_X^p(X|Y)$ if and only if it has an open neighbourhood $U$ that
admits a closed immersion $i:U\hookrightarrow U'$ into a $Y$-scheme
$U'$ such that $U'$ is smooth over $Y$ at $x'=i(x)$ and
\[
\dim_{x'}U'_{f(x)}=p.
\]

\smallskip
\noindent\textbf{(5)} For the projections
$\pi_i:X_1\times_YX_2\longrightarrow X_i$,
\[
\underline I_{X_1\times_YX_2}^r(X_1\times_YX_2|Y)
=\sum_{p+q=r}
 \pi_1^*\underline I_{X_1}^p(X_1|Y)\,
 \pi_2^*\underline I_{X_2}^q(X_2|Y).
\]

\smallskip
\noindent\textit{Proof.}
Assertion (3) follows from \SGARefLink{sga7.sga7i.evi.n5-1}{5.1}(b). Assertion (1) follows from \SGARefLink{sga7.sga7i.evi.n5-1}{5.1}(a)
and the canonical isomorphism
\[
P^*\Omega_{X|Y}\xrightarrow{\ \sim\ }\Omega_{X'|Y'},
\qquad P:X\times_YY'\longrightarrow X.
\]
Likewise, (2) follows from \SGARefLink{sga7.sga7i.evi.n5-1}{5.1}(a) and
\[
\mathcal O_{X,x}\otimes_{\mathcal O_X}\Omega_{X|Y}
\xrightarrow{\ \sim\ }\Omega_{X_x|Y}.
\]
Finally, (5) follows from \SGARefLink{sga7.sga7i.evi.n5-1}{5.1}(d) and
\[
\pi_1^*\Omega_{X_1|Y}\oplus\pi_2^*\Omega_{X_2|Y}
\xrightarrow{\ \sim\ }
\Omega_{X_1\times_YX_2|Y}.
\]

For (4), first suppose that there is a diagram
\[
\begin{tikzcd}[column sep=huge,row sep=large]
U \arrow[r,hook,"i"] \arrow[d,"f"']
  & U' \arrow[dl] \\
Y & {}
\end{tikzcd}
\]
where $U'$ is smooth over $Y$ at $x'=i(x)$ and
$\dim_{x'}U'_{f(x)}=p$. Then
\[
\dim_{\kappa(x')}\Omega_{U'|Y}(x')=p,
\]
and the sequence
\[
i^*\Omega_{U'|Y}\longrightarrow\Omega_{U|Y}\longrightarrow0
\]
is exact. Hence
$\dim_{\kappa(x)}\Omega_{U|Y}(x)\leq p$, so
$x\notin J_U^p(U|Y)$ by (3). Since
\[
J_U^p(U|Y)=J_X^p(X|Y)|_U,
\]
we have $x\notin J_X^p(X|Y)$.

%%%%%%%%%%  scan idx 128  |  folio 117  %%%%%%%%%%

Conversely, suppose that $x\notin J_X^p(X|Y)$. The question is
local, so write $X=\operatorname{Spec}(B)$ and
$Y=\operatorname{Spec}(A)$. Choose an $A$-presentation
\[
0\longrightarrow J\longrightarrow P\longrightarrow B
\longrightarrow0,
\qquad P=A[x_1,\ldots,x_n].
\]
There is an exact sequence
\[
J/J^2\xrightarrow{d}B\otimes_P\Omega_{P|A}
\longrightarrow\Omega_{B|A}\longrightarrow0,
\]
and $B\otimes_P\Omega_{P|A}$ is free of rank $n$ over $B$. Thus
\[
I_B^p(\Omega_{B|A})
=\operatorname{Im}\left(
 \bigwedge^{n-p}(J/J^2)\otimes_B
 (B\otimes_P\bigwedge^{n-p}\Omega_{P|A})^*
 \longrightarrow B\right).
\]
Consequently, $x\notin J_X^p(X|Y)$ precisely when there are
$f_1,\ldots,f_{n-p}\in J$ and a choice of $n-p$ among the variables
$x_1,\ldots,x_n$ for which the corresponding Jacobian minor
\[
\det\left(\frac{\partial f_i}{\partial x_j}\right)(x)
\]
is nonzero. Equivalently,
\[
\operatorname{Spec}\bigl(P/(f_1,\ldots,f_{n-p})\bigr)
\longrightarrow\operatorname{Spec}(A)
\]
is smooth at the image $x'$ of $x$. Taking
\[
U'=\operatorname{Spec}\bigl(P/(f_1,\ldots,f_{n-p})\bigr)
\]
gives the required closed immersion locally. \hfill$\square$


% END INCLUDED SGA7I_English_current_buildable_source_20260823/components/51_expose_VI_theorem_5_3.tex

% BEGIN INCLUDED SGA7I_English_current_buildable_source_20260823/components/52_expose_VI_theorem_5_4.tex
\medskip
\noindent\textbf{Theorem \SGARefTarget{sga7.sga7i.evi.n5-4}{5.4}.}
Let $X\xrightarrow{f}Y$ be a morphism essentially of finite
presentation. Suppose that there is a fixed integer $n\geq0$ for
which one of the following conditions holds:

\smallskip
\noindent\textbf{(1)} $f$ is flat and
$\dim_xX_{f(x)}=n$ for every $x\in X$;

\smallskip
\noindent\textbf{(2)} $f$ is a relative complete intersection of
virtual dimension $n$ (SGA 6, VIII, 1.9).

\smallskip
Then $x\notin J_X^n(X|Y)$ if and only if $f$ is smooth at $x$.

\smallskip
\noindent\textit{Proof.}
By \SGARefLink{sga7.sga7i.evi.n5-3}{5.3}(4), $x\notin J_X^n(X|Y)$ if and only if there is a diagram
\[
\begin{tikzcd}[column sep=huge,row sep=large]
U \arrow[r,hook,"i"] \arrow[d,"f"']
  & U' \arrow[dl,"f'"] \\
Y & {}
\end{tikzcd}
\]
such that $f'$ is smooth at $x'$, with
$\dim_{x'}U'_{f(x)}=n$, while $i$ is a closed immersion and
$x'=i(x)$. It therefore suffices to show that under either (1) or
(2), $i$ is an isomorphism at $x$.

Set
\[
A=\mathcal O_{Y,f(x)},\qquad
B'=\mathcal O_{U',x'},\qquad B=\mathcal O_{U,x},
\]
and let
\[
0\longrightarrow I\longrightarrow B'\longrightarrow B
\longrightarrow0
\]
be exact. We must show that $I=0$.

Assume (1). Since $B$ is flat over $A$, the sequence
\[
0\longrightarrow I/\mathfrak mI
\longrightarrow B'/\mathfrak mB'
\longrightarrow B/\mathfrak mB\longrightarrow0
\]
is exact, where $\mathfrak m$ is the maximal ideal of $A$. By
hypothesis,
\[
\dim_xU_{f(x)}=\dim_{x'}U'_{f(x)},
\]
and hence
$\dim(B'/\mathfrak mB')=\dim(B/\mathfrak mB)$. Since
$B'/\mathfrak mB'$ is a regular local ring, it follows that
$I/\mathfrak mI=0$, and therefore $I=0$.

%%%%%%%%%%  scan idx 129  |  folio 118  %%%%%%%%%%

Now assume (2), and consider the exact sequence
\[
I/I^2\xrightarrow{d}B\otimes_{B'}\Omega_{B'|A}
\longrightarrow\Omega_{B|A}\longrightarrow0.
\]
Since $X$ is a relative complete intersection over $Y$ and
$U'\longrightarrow Y$ is smooth at $x'$, the map
\[
d:I/I^2\longrightarrow B\otimes_{B'}\Omega_{B'|A}
\]
is injective at $x$, and both $I/I^2$ and
$B\otimes_{B'}\Omega_{B'|A}$ are free at $x$. By the definition of
the virtual dimension of a relative complete intersection, we must
therefore have
\[
n=\dim_{\kappa(x)}\Omega_{B'|A}(x')
  -\dim_{\kappa(x)}(I/I^2)(x)
 =n-\dim_{\kappa(x)}(I/I^2)(x).
\]
Thus $I/I^2=0$, and hence $I=0$. \hfill$\square$

% END INCLUDED SGA7I_English_current_buildable_source_20260823/components/52_expose_VI_theorem_5_4.tex

% BEGIN INCLUDED SGA7I_English_current_buildable_source_20260823/components/53_expose_VI_5_5_through_5_6.tex
\medskip
\noindent\textbf{Definition \SGARefTarget{sga7.sga7i.evi.n5-5}{5.5}.}
Let $X\xrightarrow{f}Y$ be a morphism essentially of finite
presentation. If there is a fixed integer $n$ satisfying (1) or (2)
of \SGARefLink{sga7.sga7i.evi.n5-4}{5.4}, set
\[
\underline I_X(X|Y)=\underline I_X^n(X|Y),
\qquad
J_X(X|Y)=J_X^n(X|Y).
\]

%%%%%%%%%%  scan idx 130  |  folio 119  %%%%%%%%%%

This notion extends naturally to formal schemes. Let $\mathfrak X$
be a noetherian formal scheme and $\mathcal E$ a coherent
$\mathfrak X$-module. For every affine open $U\subset\mathfrak X$,
define
\[
\underline I_{\mathfrak X}^p(\mathcal E)(U)
=I_{\mathfrak X(U)}^p\bigl(\mathcal E(U)\bigr).
\]
If $\mathfrak X=\varinjlim_\nu X_\nu$, then
$\mathcal E=\varprojlim_\nu\mathcal E_\nu$, where
\[
\mathcal E_\nu
=\mathcal E\otimes_{\mathcal O_{\mathfrak X}}\mathcal O_{X_\nu},
\]
and noetherianness gives
\[
\underline I_{\mathfrak X}^p(\mathcal E)
=\varprojlim_\nu\underline I_{X_\nu}^p(\mathcal E_\nu).
\]

Now let $S$ be a noetherian formal scheme and $\mathfrak X$ an
$S$-adic formal scheme essentially of finite type. Let
$\underline I_{\mathfrak X}$ be the ideal sheaf defining the closed
diagonal immersion
\[
\Delta:\mathfrak X\longrightarrow
\mathfrak X\times_S\mathfrak X
\]
(where the fibre product is taken in the category of formal schemes),
so that
\[
0\longrightarrow\underline I_{\mathfrak X}
\longrightarrow\mathcal O_{\mathfrak X\times_S\mathfrak X}
\longrightarrow\mathcal O_{\mathfrak X}\longrightarrow0
\]
is exact. Define
\[
\Omega_{\mathfrak X|S}
=\underline I_{\mathfrak X}/\underline I_{\mathfrak X}^{,2}.
\]
Writing
\[
\mathfrak X=\varinjlim_\nu X_\nu,
\qquad X_\nu=\mathfrak X\times_SS_\nu,
\]
we have
\[
\Omega_{\mathfrak X|S}
\simeq\varprojlim_\nu\Omega_{X_\nu|S_\nu},
\]
so $\Omega_{\mathfrak X|S}$ is a coherent $\mathfrak X$-module.
For every formal $S$-scheme $S'$ of finite type there is a canonical
isomorphism
\[
\Omega_{\mathfrak X\times_SS'|S'}
\simeq
\mathcal O_{S'}\otimes_{\mathcal O_S}\Omega_{\mathfrak X|S}.
\]

\medskip
\noindent\textbf{Definition \SGARefTarget{sga7.sga7i.evi.n5-6}{5.6}.}
Let $S$ be a noetherian formal scheme and $\mathfrak X$ an $S$-adic
formal scheme essentially of finite type. Define
\[
\underline I_{\mathfrak X}^p(\mathfrak X|S)
=\underline I_{\mathfrak X}^p(\Omega_{\mathfrak X|S}),
\]
and let $J_{\mathfrak X}^p(\mathfrak X|S)$ be the closed formal
subscheme of $\mathfrak X$ defined by this ideal sheaf.

Since
$\Omega_{\mathfrak X|S}=\varprojlim_\nu\Omega_{X_\nu|S_\nu}$,
we have
\[
\underline I_{\mathfrak X}^p(\mathfrak X|S)
=\varprojlim_\nu\underline I_{X_\nu}^p(X_\nu|S_\nu)
\]
and
\[
J_{\mathfrak X}^p(\mathfrak X|S)
=\varinjlim_\nu J_{X_\nu}^p(X_\nu|S_\nu).
\]

%%%%%%%%%%  scan idx 131  |  folio 120  %%%%%%%%%%

The formal scheme $\mathfrak X_{(x)}$ is again $S$-adic and
essentially of finite type, and there is a commutative diagram
\[
\begin{tikzcd}[column sep=huge,row sep=large]
\mathfrak X_{(\widehat{x})}
  \arrow[rr,"i_{(\widehat{x})}"] \arrow[dr]
  && \mathfrak X \\
& \mathfrak X_{(x)} \arrow[ur,"i_{(x)}"'] & {}
\end{tikzcd}
\]
% Direct 1800-dpi source check: the upper-left subscript and its arrow
% label are both \widehat{x}; the lower subscript and label are plain x.
If $x$ is an isolated point of $\mathfrak X$, then
$i_{(x)}$ and $i_{(\widehat{x})}$ are both isomorphisms onto the
connected component of $\mathfrak X$ containing $x$.

In particular, if
\[
|\mathfrak X|=\{x_1,\ldots,x_m\}
\]
is a finite discrete space, then there is a canonical identification
\[
\mathfrak X=\coprod_{1\leq i\leq m}\mathfrak X_{(x_i)}
\qquad\text{(disjoint union)}.
\]

% END INCLUDED SGA7I_English_current_buildable_source_20260823/components/53_expose_VI_5_5_through_5_6.tex

% BEGIN INCLUDED SGA7I_English_current_buildable_source_20260823/components/54_expose_VI_corollary_5_7.tex
\medskip
\noindent\textbf{Corollary \SGARefTarget{sga7.sga7i.evi.n5-7}{5.7}.}
Let $S$ be a noetherian formal scheme and $\mathfrak X$ an $S$-adic
formal scheme essentially of finite type. Then:

\smallskip
\noindent\textbf{(1)} Formation of
$J_{\mathfrak X}^p(\mathfrak X|S)$ commutes with base change. More
precisely, if
\[
\mathfrak X'=\mathfrak X\times_SS',
\]
then
\[
J_{\mathfrak X'}^p(\mathfrak X'|S')
=J_{\mathfrak X}^p(\mathfrak X|S)\times_SS'.
\]

\smallskip
\noindent\textbf{(2)} The connected components of
$J_{\mathfrak X}^p(\mathfrak X|S)$ are in bijective correspondence
with those of $J_{X_0}^p(X_0|S_0)$. For every
\[
x\in\bigl|J_{X_0}^p(X_0|S_0)\bigr|
\]
there is a canonical isomorphism
\[
J_{\mathfrak X}^p(\mathfrak X|S)_{(x)}
\simeq
J_{\mathfrak X_{(x)}}^p(\mathfrak X_{(x)}|S).
\]
In particular, if
\[
\bigl|J_{X_0}^p(X_0|S_0)\bigr|=\{x_1,\ldots,x_m\}
\]
is discrete, then
\[
J_{\mathfrak X}^p(\mathfrak X|S)
=\coprod_{1\leq i\leq m}
 J_{\mathfrak X_{(x_i)}}^p(\mathfrak X_{(x_i)}|S)
\qquad\text{(disjoint union)}.
\]

%%%%%%%%%%  scan idx 132  |  folio 121  %%%%%%%%%%

\smallskip
\noindent\textbf{(3)} Let $S$ be an ordinary noetherian scheme and
$X$ an $S$-scheme essentially of finite type. Given a closed
subscheme $S_0\subset S$, let $\widehat S$ be the formal completion
of $S$ along $S_0$. If
\[
\widehat X=X\times_S\widehat S
\]
is the formal completion of $X$ along
$X_0=X\times_SS_0$, then
\[
J_{\widehat X}^p(\widehat X|\widehat S)
=J_X^p(X|S)\times_S\widehat S.
\]
Equivalently, it is the formal completion of $J_X^p(X|S)$ along
$J_{X_0}^p(X_0|S_0)$.

\smallskip
\noindent\textbf{(4)} Suppose that a fixed integer $n\geq0$
satisfies condition \SGARefLink{sga7.sga7i.evi.n5-4}{5.4}(1) or \SGARefLink{sga7.sga7i.evi.n5-4}{5.4}(2) for every
$X_\nu\longrightarrow S_\nu$. Then
\[
x\notin J_{\mathfrak X}^n(\mathfrak X|S)
\]
if and only if $\mathfrak X$ is formally smooth over $S$ at $x$.

\smallskip
\noindent\textit{Proof.}
Assertion (1) follows from
\[
\Omega_{\mathfrak X'|S'}
\simeq
\mathcal O_{S'}\widehat\otimes_{\mathcal O_S}
\Omega_{\mathfrak X|S},
\]
whereas (2) follows from
\[
\mathcal O_{\mathfrak X_{(x)}}
\widehat\otimes_{\mathcal O_{\mathfrak X}}
\Omega_{\mathfrak X|S}
\simeq\Omega_{\mathfrak X_{(x)}|S}.
\]
For (3),
\[
\begin{aligned}
J_{\widehat X}^p(\widehat X|\widehat S)
&=\varinjlim_\nu
 J_{X\times_SS_\nu}^p(X\times_SS_\nu|S_\nu) \\
&=\varinjlim_\nu\bigl(J_X^p(X|S)\times_SS_\nu\bigr)
 =J_X^p(X|S)\times_S\widehat S.
\end{aligned}
\]
Finally, since $\mathfrak X$ is $S$-adic,
\[
\begin{aligned}
\mathfrak X\to S\text{ is formally smooth at }x
&\Longleftrightarrow
 X_\nu\to S_\nu\text{ is formally smooth at }x
 \text{ for every }\nu \\
&\Longleftrightarrow
 x\notin J_{X_\nu}^n(X_\nu|S_\nu)
 \text{ for every }\nu \\
&\Longleftrightarrow
 x\notin\varinjlim_\nu J_{X_\nu}^n(X_\nu|S_\nu)
 =J_{\mathfrak X}^n(\mathfrak X|S).
\end{aligned}
\]
\hfill$\square$

% END INCLUDED SGA7I_English_current_buildable_source_20260823/components/54_expose_VI_corollary_5_7.tex

% BEGIN INCLUDED SGA7I_English_current_buildable_source_20260823/components/55_expose_VI_remark_5_8.tex
\medskip
\noindent\textbf{Remark \SGARefTarget{sga7.sga7i.evi.n5-8}{5.8}.}
Let $\mathfrak X$ be an $S$-adic formal scheme essentially of finite
type. For a point $x\in\mathfrak X$, consider also the $S$-adic
formal scheme
\[
\mathfrak X_{(\widehat{x})}
=\varinjlim_\nu
 \operatorname{Spec}\bigl(\widehat{\mathcal O_{X_\nu,x}}\bigr).
\]
There are natural morphisms
\[
\mathfrak X_{(\widehat{x})}
\longrightarrow\mathfrak X_{(x)}
\longrightarrow\mathfrak X.
\]
In general, $\widehat{\mathcal O_{X_0,x}}$ is not essentially of
finite type over $S_{(y)}$, where $y$ is the image of $x$ under
$\mathfrak X\longrightarrow S$.

For simplicity, put
\[
A_\nu=\mathcal O_{X_\nu,x},
\qquad R_\nu=\mathcal O_{S_\nu,y}.
\]
The natural map
\[
\Omega_{A_\nu|R_\nu}
\longrightarrow\widehat\Omega_{A_\nu|R_\nu}
\]
induces an isomorphism
\[
\widehat{\Omega_{A_\nu|R_\nu}}
\xrightarrow{\ \sim\ }
\widehat{\widehat\Omega_{A_\nu|R_\nu}},
\]
and therefore
\[
\varprojlim_\nu\widehat{\Omega_{A_\nu|R_\nu}}
\xrightarrow{\ \sim\ }
\varprojlim_\nu
 \widehat{\widehat\Omega_{A_\nu|R_\nu}}.
\]
Here the outer hat denotes the separated completion of a module for
the topology induced by the ring. Consequently, if we define
\[
\widehat\Omega_{\mathfrak X_{(\widehat{x})}|S}
=\varprojlim_\nu
 \widehat{\widehat\Omega_{
   \mathcal O_{X_\nu,x}|\mathcal O_{S_\nu,y}}},
\]
then $\widehat\Omega_{\mathfrak X_{(\widehat{x})}|S}$ is a coherent
$\mathfrak X_{(\widehat{x})}$-module.

%%%%%%%%%%  scan idx 133  |  folio 122  %%%%%%%%%%

Define
$J_{\mathfrak X_{(\widehat{x})}}^p
 (\mathfrak X_{(\widehat{x})}|S)$
to be the closed formal subscheme of
$\mathfrak X_{(\widehat{x})}$ defined by
\[
\underline I_{\mathfrak X_{(\widehat{x})}}^p
 (\mathfrak X_{(\widehat{x})}|S)
=\underline I_{\mathfrak X_{(\widehat{x})}}^p
 (\widehat\Omega_{\mathfrak X_{(\widehat{x})}|S}).
\]
Then
\[
J_{\mathfrak X_{(x)}}^p(\mathfrak X_{(x)}|S)_{(\widehat{x})}
\simeq
J_{\mathfrak X_{(\widehat{x})}}^p
 (\mathfrak X_{(\widehat{x})}|S).
\]
If $x$ is an isolated point of
$J_{\mathfrak X}^p(\mathfrak X|S)$, the natural morphism
\[
J_{\mathfrak X_{(\widehat{x})}}^p
 (\mathfrak X_{(\widehat{x})}|S)
\longrightarrow
J_{\mathfrak X_{(x)}}^p(\mathfrak X_{(x)}|S)
\]
is an isomorphism.

To form the ``image'' of $J_{\mathfrak X}^p(\mathfrak X|S)$ under
$f:\mathfrak X\longrightarrow S$, one could simply take the formal
subscheme of $S$ defined by the ideal sheaf
\[
\ker\left(
 \mathcal O_S\longrightarrow
 \mathcal O_{\mathfrak X}/
 \underline I_{\mathfrak X}^p(\mathfrak X|S)
\right).
\]
For our purposes, however, it is more convenient to use the Fitting
ideal once again.

% END INCLUDED SGA7I_English_current_buildable_source_20260823/components/55_expose_VI_remark_5_8.tex

% BEGIN INCLUDED SGA7I_English_current_buildable_source_20260823/components/56_expose_VI_5_9_through_5_11.tex
\medskip
\noindent\textbf{Definition \SGARefTarget{sga7.sga7i.evi.n5-9}{5.9}.}
Let $S$ be a noetherian formal scheme and
$f:\mathfrak X\longrightarrow S$ an $S$-adic formal scheme
essentially of finite type. Suppose that
$J_{\mathfrak X}^p(\mathfrak X|S)$ is proper over $S$, so that
\[
f_*\mathcal O_{J_{\mathfrak X}^p(\mathfrak X|S)}
\]
is a coherent $\mathcal O_S$-module. Define
\[
\underline I_S^p(\mathfrak X|S)
=\underline I_S^0\!\left(
 f_*\mathcal O_{J_{\mathfrak X}^p(\mathfrak X|S)}
\right),
\]
and let $J_S^p(\mathfrak X|S)$ be the closed formal subscheme of $S$
defined by $\underline I_S^p(\mathfrak X|S)$.

If there is a fixed integer $n$ satisfying condition \SGARefLink{sga7.sga7i.evi.n5-4}{5.4}(1) or
\SGARefLink{sga7.sga7i.evi.n5-4}{5.4}(2) for every $X_\nu\longrightarrow S_\nu$, set
\[
\begin{aligned}
\underline I_{\mathfrak X}(\mathfrak X|S)
 &=\underline I_{\mathfrak X}^n(\mathfrak X|S),
&\qquad
J_{\mathfrak X}(\mathfrak X|S)
 &=J_{\mathfrak X}^n(\mathfrak X|S),\\
\underline I_S(\mathfrak X|S)
 &=\underline I_S^n(\mathfrak X|S),
&
J_S(\mathfrak X|S)
 &=J_S^n(\mathfrak X|S).
\end{aligned}
\]

\medskip
\noindent\textbf{Corollary \SGARefTarget{sga7.sga7i.evi.n5-10}{5.10}.}
Let $S$ be a noetherian formal scheme and
$f:\mathfrak X\longrightarrow S$ an $S$-adic formal scheme
essentially of finite type such that
$J_{\mathfrak X}^p(\mathfrak X|S)$ is proper over $S$. Then:

\smallskip
\noindent\textbf{(1)}
\[
\bigl|J_S^p(\mathfrak X|S)\bigr|
=f\bigl(\,|J_{\mathfrak X}^p(\mathfrak X|S)|\,\bigr).
\]

%%%%%%%%%%  scan idx 134  |  folio 123  %%%%%%%%%%

\smallskip
\noindent\textbf{(2)} Formation of $J_S^p(\mathfrak X|S)$ commutes
with base change: if $\mathfrak X'=\mathfrak X\times_SS'$, then
\[
J_{S'}^p(\mathfrak X'|S')
=J_S^p(\mathfrak X|S)\times_SS'.
\]

\smallskip
\noindent\textbf{(3)} If
\[
J_{\mathfrak X}^p(\mathfrak X|S)
=\coprod_{1\leq i\leq m}\mathfrak Y^{(i)}
\]
is a disjoint union, then
\[
\underline I_S^p(\mathfrak X|S)
=\prod_{1\leq i\leq m}
 \underline I_S^0\bigl(f_*\mathcal O_{\mathfrak Y^{(i)}}\bigr).
\]
In particular, if
\[
\bigl|J_{X_0}^p(X_0|S_0)\bigr|=\{x_1,\ldots,x_m\}
\]
is discrete, then
\[
\underline I_S^p(\mathfrak X|S)
=\prod_{1\leq i\leq m}
 \underline I_S^p(\mathfrak X_{(x_i)}|S)
=\prod_{1\leq i\leq m}
 \underline I_S^p(\mathfrak X_{(\widehat{x_i})}|S).
\]

\smallskip
\noindent\textbf{(4)} Let $S$ be an ordinary noetherian scheme,
$X$ an $S$-scheme essentially of finite type, and
$S_0\hookrightarrow S$ a closed subscheme. Let $\widehat S$ be the
formal completion of $S$ along $S_0$. If
\[
\widehat X=X\times_S\widehat S
\]
is the formal completion of $X$ along
$X_0=X\times_SS_0$, then
\[
J_{\widehat S}^p(\widehat X|\widehat S)
=J_S^p(X|S)\times_S\widehat S.
\]
Thus it is the formal completion of $J_S^p(X|S)$ along
$J_{S_0}^p(X_0|S_0)$.

\smallskip
\noindent\textit{Proof.}
By \SGARefLink{sga7.sga7i.evi.n5-1}{5.1}(c),
\[
\operatorname{Supp}
 f_*\mathcal O_{J_{\mathfrak X}^p(\mathfrak X|S)}
=\bigl|J_S^p(\mathfrak X|S)\bigr|,
\]
which proves (1). Assertions (2) and (4) follow from \SGARefLink{sga7.sga7i.evi.n5-1}{5.1}(a), while
(3) follows from \SGARefLink{sga7.sga7i.evi.n5-1}{5.1}(d) and, for an isolated point
$x\in J_{\mathfrak X}^p(\mathfrak X|S)$, the isomorphism
\[
J_{\mathfrak X_{(\widehat{x})}}^p
 (\mathfrak X_{(\widehat{x})}|S)
\xrightarrow{\ \sim\ }
J_{\mathfrak X_{(x)}}^p(\mathfrak X_{(x)}|S).
\]
\hfill$\square$

\medskip
\noindent\textbf{Remark \SGARefTarget{sga7.sga7i.evi.n5-11}{5.11}.}
Our definition of the Jacobian ideal, adopted from the classical
definition over a base field, commutes with base change and, under the
mild hypothesis of \SGARefLink{sga7.sga7i.evi.n5-4}{5.4}, gives a smoothness criterion. In a general
setting, however, one can probably define in many ways a canonical
subscheme whose support is the nonsmooth locus. We outline one such
definition.

For an $A$-algebra $B$, put
\[
I_{B|A}=\bigcap_E\operatorname{Ann}D^1(B|A,E),
\]
where $E$ ranges over all finite-type $B$-modules. By \SGARefLink{sga7.sga7i.evi.n3-13}{3.13}(c), the
ideal $I_{B|A}$ commutes with localization. Thus a morphism
$f:X\longrightarrow Y$ determines an ideal sheaf
$\underline I_{X|Y}$. If $f$ is essentially of finite type, then by
\SGARefLink{sga7.sga7i.evi.n3-13}{3.13}(a), $f$ is smooth at $x$ if and only if
\[
x\notin\operatorname{Supp}
 \bigl(\mathcal O_X/\underline I_{X|Y}\bigr).
\]

%%%%%%%%%%  scan idx 135  |  folio 124  %%%%%%%%%%

This construction carries over naturally to adic formal schemes
essentially of finite type, and is therefore equally suitable for
deformation theory.

% END INCLUDED SGA7I_English_current_buildable_source_20260823/components/56_expose_VI_5_9_through_5_11.tex

% BEGIN INCLUDED SGA7I_English_current_buildable_source_20260823/components/57_expose_VI_section_6_opening_and_lemma_6_1.tex
\par\bigskip
\noindent\textbf{\SGARefAliasTarget{sga7.sga7i.evi.n6-0}6. Non-degenerate quadratic singularities}

\medskip
The purpose of this section is to study deformations of a $k$-scheme
whose only singularities are non-degenerate quadratic singularities.
In view of \SGARefLink{sga7.sga7i.evi.n4-13}{4.13}, this becomes, in a suitable situation, a purely local
problem.

We return to the notation of \S4. Thus $\Lambda$ is a fixed complete
noetherian local ring with residue field $k$, and
$\underline C_\Lambda$ denotes the category of local artinian
$\Lambda$-algebras having residue field $k$.

Let $X$ be a $k$-scheme of finite type, and let $x$ be a closed point
of $X$. Recall that $x$ is called a \emph{non-degenerate quadratic
singularity rational over $k$} if
\[
\widehat{\mathcal O}_{X,x}
 \simeq k[[x_1,\ldots,x_n]]/(f)
\]
and
\[
\left(\frac{\partial f}{\partial x_1},\ldots,
      \frac{\partial f}{\partial x_n}\right)
=(x_1,\ldots,x_n).
\]
One may always take $f$ to be a quadratic form. Indeed:

\medskip
\noindent\textbf{Lemma \SGARefTarget{sga7.sga7i.evi.n6-1}{6.1}.}
Let $f\in k[[x_1,\ldots,x_n]]$, and put
$M=(x_1,\ldots,x_n)$. If
\[
\left(\frac{\partial f}{\partial x_1},\ldots,
      \frac{\partial f}{\partial x_n}\right)=M,
\]
then there are elements $x_1',\ldots,x_n'$ satisfying
$x_i'\equiv x_i\pmod{M^2}$ such that
\[
f=q(x_1',\ldots,x_n')
\]
for a quadratic form $q$.

\smallskip
\noindent\textit{Proof.}
We use successive approximation. Write
\[
f=f_2+h_3,
\]
where $f_2$ is quadratic and $h_3\in M^3$. Since the partial
derivatives of $f$ generate $M$, there are
$\varepsilon_1^{(1)},\ldots,\varepsilon_n^{(1)}\in M^2$ such that
\[
\sum_i\varepsilon_i^{(1)}\frac{\partial f}{\partial x_i}
 \equiv h_3\pmod{M^4}.
\]
Put $x_i^{(1)}=x_i-\varepsilon_i^{(1)}$. Then
\[
\begin{aligned}
f(x_1^{(1)},\ldots,x_n^{(1)})
&=f(x_1,\ldots,x_n)
 -\sum_i\varepsilon_i^{(1)}\frac{\partial f}{\partial x_i}
 +\cdots\\
&\equiv f_2(x_1,\ldots,x_n)\pmod{M^4}.
\end{aligned}
\]
Thus
\[
f(x_1^{(1)},\ldots,x_n^{(1)})
=f_2(x_1,\ldots,x_n)+h_4,
\qquad h_4\in M^4.
\]
Moreover, the derivatives with respect to the new variables still
generate $M$. Hence one can choose
$\varepsilon_1^{(2)},\ldots,\varepsilon_n^{(2)}\in M^3$ such that
\[
\sum_i\varepsilon_i^{(2)}
 \frac{\partial f}{\partial x_i^{(1)}}
 \equiv h_4\pmod{M^5}.
\]

%%%%%%%%%%  scan idx 136  |  folio 125  %%%%%%%%%%

It follows that
\[
\begin{aligned}
&f(x_1^{(1)}-\varepsilon_1^{(2)},\ldots,
  x_n^{(1)}-\varepsilon_n^{(2)})\\
&\qquad =f(x_1^{(1)},\ldots,x_n^{(1)})
 -\sum_i\varepsilon_i^{(2)}
       \frac{\partial f}{\partial x_i^{(1)}}+\cdots\\
&\qquad \equiv f_2(x_1,\ldots,x_n)\pmod{M^5}.
\end{aligned}
\]
Continuing in this way, one finds elements
$\varepsilon_i^{(\nu)}\in M^{\nu+1}$ such that, on setting
\[
x_i^{(\nu)}
=x_i-\bigl(\varepsilon_i^{(1)}+\cdots+
            \varepsilon_i^{(\nu)}\bigr),
\]
one has
\[
f(x_1^{(\nu)},\ldots,x_n^{(\nu)})
\equiv f_2(x_1,\ldots,x_n)\pmod{M^{\nu+3}}.
\]
Let $\varepsilon_i=\sum_\nu\varepsilon_i^{(\nu)}$. Completeness then
gives
\[
f(x_1-\varepsilon_1,\ldots,x_n-\varepsilon_n)
=f_2(x_1,\ldots,x_n),
\]
or, equivalently,
\[
f(x_1,\ldots,x_n)
=f_2(x_1+\varepsilon_1,\ldots,x_n+\varepsilon_n).
\]
Taking $x_i'=x_i+\varepsilon_i$ proves the assertion.
\hfill$\square$

% END INCLUDED SGA7I_English_current_buildable_source_20260823/components/57_expose_VI_section_6_opening_and_lemma_6_1.tex

% BEGIN INCLUDED SGA7I_English_current_buildable_source_20260823/components/58_expose_VI_lemma_6_2.tex
\par\medskip
\noindent\textbf{Lemma \SGARefTarget{sga7.sga7i.evi.n6-2}{6.2}.}
Let $A_0$ be a local $k$-algebra essentially of finite type such that
\[
\widehat A_0\simeq k[[x_1,\ldots,x_n]]/(f),
\]
where $f$ is a non-degenerate quadratic form, and let
$(R,\mathfrak X)$ be a formal versal deformation of
$X_0=\operatorname{Spec}(A_0)$. Then:

\smallskip
\noindent\textbf{(1)} The canonical maps
\[
J_{\mathfrak X}(\mathfrak X|R)
 \longrightarrow J_R(\mathfrak X|R)
 \longrightarrow \operatorname{Spec}(\Lambda)
\]
are both isomorphisms.

\smallskip
\noindent\textbf{(2)} The ideal
$I_R(\mathfrak X|R)$ is principal and generated by a free formal
generator of $R$ over $\Lambda$; in other words,
\[
I_R(\mathfrak X|R)=(t),
\qquad R=\Lambda[[t]].
\]

\smallskip
\noindent\textbf{(3)} The morphism
$\mathfrak X\longrightarrow\operatorname{Spec}(\Lambda)$ is formally
smooth.

\smallskip
\noindent\textit{Proof.}
Write $\mathfrak X=\operatorname{Spf}(A)$. By \SGARefLink{sga7.sga7i.evi.n4-10}{4.10}(b),
$(R,\operatorname{Spf}(\widehat A))$ is a formal versal deformation
of $\widehat X_0=\operatorname{Spec}(\widehat A_0)$, where
\[
\widehat A=\varprojlim_\nu A\otimes_RR_\nu.
\]
Now
\[
\widehat A_0\simeq k[[x_1,\ldots,x_n]]/(f),
\qquad
\left(\frac{\partial f}{\partial x_1},\ldots,
      \frac{\partial f}{\partial x_n}\right)
=(x_1,\ldots,x_n).
\]
It follows immediately that
\[
\bigl[\underline M_{\widehat X_0}(k[\varepsilon])\bigr]
\simeq k.
\]
Consequently, by \SGARefLink{sga7.sga7i.evi.n4-14}{4.14},
\[
R\simeq\Lambda[[t]],
\qquad
\widehat A\simeq R[[x_1,\ldots,x_n]]/(F-t),
\]
where $F$ is the quadratic polynomial in
$\Lambda[x_1,\ldots,x_n]\subset R[[x_1,\ldots,x_n]]$ obtained from
$f$ by lifting its coefficients to $\Lambda$.

The sequence
\[
0\longrightarrow\widehat A
 \xrightarrow{\ dF\ }
 \widehat A\,dx_1\oplus\cdots\oplus\widehat A\,dx_n
 \longrightarrow\Omega_{\widehat A|R}\longrightarrow0
\]
is exact, where
\[
1\longmapsto
 \frac{\partial F}{\partial x_1}dx_1+\cdots+
 \frac{\partial F}{\partial x_n}dx_n.
\]

%%%%%%%%%%  scan idx 137  |  folio 126  %%%%%%%%%%

Therefore
\[
I_{\widehat A}(\widehat A|R)
=\left(\frac{\partial F}{\partial x_1},\ldots,
       \frac{\partial F}{\partial x_n}\right)
=(x_1,\ldots,x_n),
\]
and hence
\[
\begin{aligned}
\widehat A/I_{\widehat A}(\widehat A|R)
&=\widehat A/(x_1,\ldots,x_n)\\
&=R[[x_1,\ldots,x_n]]/(x_1,\ldots,x_n,F-t)
 \simeq\Lambda.
\end{aligned}
\]
Thus the canonical map
$\Lambda\longrightarrow
\widehat A/I_{\widehat A}(\widehat A|R)$ is an isomorphism. The
completion comparison identifies this quotient with
$A/I_A(A|R)$, and consequently
\[
\Lambda\xrightarrow{\ \sim\ }R/I_R(A|R).
\]
This proves (1) and (2). Finally, (3) follows from the isomorphisms
\[
\begin{aligned}
\widehat A
&\simeq R[[x_1,\ldots,x_n]]/(F-t)\\
&\simeq\Lambda[[t,x_1,\ldots,x_n]]/(F-t)
 \simeq\Lambda[[x_1,\ldots,x_n]]
\end{aligned}
\]
of $\Lambda$-algebras.
\hfill$\square$

% END INCLUDED SGA7I_English_current_buildable_source_20260823/components/58_expose_VI_lemma_6_2.tex

% BEGIN INCLUDED SGA7I_English_current_buildable_source_20260823/components/59_expose_VI_corollary_6_3.tex
\par\medskip
\noindent\textbf{Corollary \SGARefTarget{sga7.sga7i.evi.n6-3}{6.3}.}
Let $X_0$ be a $k$-scheme of finite type, and let $x\in X_0$ be a
non-degenerate quadratic singularity rational over $k$.

\smallskip
\noindent\textbf{(1)} Let $(R,\mathfrak X)$ and
$(R,\mathfrak Y)$ be two deformations of $X_0$, with
$R\in\operatorname{ob}(\widehat{\underline C}_\Lambda)$. Assume that
$k=k^2$. Then
\[
\mathfrak X_{(\widehat x)}\simeq
\mathfrak Y_{(\widehat x)}
\]
as $R$-adic formal schemes if and only if
\[
I_R(\mathfrak X_{(x)}|R)=I_R(\mathfrak Y_{(x)}|R).
\]

\smallskip
\noindent\textbf{(2)} Let
$R\in\operatorname{ob}(\widehat{\underline C}_\Lambda)$ be a regular
local ring, and let $(R,\mathfrak X)$ be a deformation of $X_0$.
Then $\mathcal O_{\mathfrak X,x}$ is regular if and only if
\[
I_R(\mathfrak X_{(x)}|R)\not\subset\mathfrak m_R^2;
\]
equivalently, $I_R(\mathfrak X_{(x)}|R)$ is a regular divisor in
$\operatorname{Spf}(R)$.

\smallskip
\noindent\textit{Proof.}
Let $(R_0,Z)$ be a formal versal deformation of
$\operatorname{Spec}(\widehat{\mathcal O}_{X_0,x})$. Since $x$ is a
non-degenerate quadratic singularity rational over $k$, Lemma \SGARefLink{sga7.sga7i.evi.n6-2}{6.2}
allows us to take
\[
R_0\simeq\Lambda[[t]],
\qquad
Z\simeq\operatorname{Spf}
 \bigl(R_0[[X_1,\ldots,X_n]]/(F-t)\bigr).
\]

For (1), the deformations
$(R,\mathfrak X_{(\widehat x)})$ and
$(R,\mathfrak Y_{(\widehat x)})$ are induced from $(R_0,Z)$ by
continuous local $\Lambda$-algebra homomorphisms
\[
\varphi,\psi:R_0\longrightarrow R.
\]
Put $a=\varphi(t)$ and $b=\psi(t)$. Then
\[
\begin{aligned}
\mathfrak X_{(\widehat x)}
&\simeq\operatorname{Spf}
 \bigl(R[[X_1,\ldots,X_n]]/(F-a)\bigr),\\
\mathfrak Y_{(\widehat x)}
&\simeq\operatorname{Spf}
 \bigl(R[[X_1,\ldots,X_n]]/(F-b)\bigr).
\end{aligned}
\]

%%%%%%%%%%  scan idx 138  |  folio 127  %%%%%%%%%%

Suppose that
$I_R(\mathfrak X_{(x)}|R)=I_R(\mathfrak Y_{(x)}|R)$; equivalently,
$(a)=(b)$ in $R$. Then $a=ub$ for some unit $u\in R$. Since
$k=k^2$ and $R$ is complete, $u=v^2$ for some unit $v\in R$, and
hence $a=v^2b$. Thus
\[
(F-a)\longmapsto v^2(F-b)
\]
under $X_i\mapsto vX_i$, and this substitution induces an
$R$-algebra isomorphism
\[
R[[X_1,\ldots,X_n]]/(F-a)
 \xrightarrow{\ \sim\ }
R[[X_1,\ldots,X_n]]/(F-b),
\]
because $F$ is quadratic. The converse is immediate.

For (2), there is a continuous local $\Lambda$-algebra homomorphism
\[
\varphi:R_0=\Lambda[[t]]\longrightarrow R
\]
such that
\[
\mathfrak X_{(\widehat x)}
\simeq\operatorname{Spf}
 \bigl(R[[X_1,\ldots,X_n]]/(F-\varphi(t))\bigr).
\]
Moreover,
\[
I_R(\mathfrak X_{(x)}|R)
=I_R(\mathfrak X_{(\widehat x)}|R)
=(\varphi(t)),
\]
and
\[
R[[X_1,\ldots,X_n]]/(F-\varphi(t))
=\widehat{\mathcal O}_{\mathfrak X,x}
\]
is regular if and only if $\mathcal O_{\mathfrak X,x}$ is regular.
Since $R$ is regular,
$R[[X_1,\ldots,X_n]]$ is a regular local ring. If $M$ denotes its
maximal ideal, then its quotient by $F-\varphi(t)$ is regular if and
only if
\[
F-\varphi(t)\notin M^2,
\]
which is equivalent to
$\varphi(t)\notin\mathfrak m_R^2$, and hence to
\[
I_R(\mathfrak X_{(x)}|R)\not\subset\mathfrak m_R^2.
\]
\hfill$\square$

% END INCLUDED SGA7I_English_current_buildable_source_20260823/components/59_expose_VI_corollary_6_3.tex

% BEGIN INCLUDED SGA7I_English_current_buildable_source_20260823/components/60_expose_VI_theorem_6_4_through_part_4.tex
\par\medskip
\noindent\textbf{Theorem \SGARefTarget{sga7.sga7i.evi.n6-4}{6.4}.}
Let $X_0$ be a $k$-scheme of finite type whose nonsmooth points are
isolated, and let $(R,\mathfrak X)$ be a formal versal deformation of
$X_0$. Denote the set of nonsmooth points of $X_0$ by
$\{z_1,\ldots,z_m\}$. Assume that
\[
H^2(X_0,\underline D^0_{X_0})=0
\]
and that every $z_i$ is a non-degenerate quadratic singularity
rational over $k$. Then:

\smallskip
\noindent\textbf{(1)} $R$ is formally smooth over $\Lambda$; that is,
$R$ is a formal power-series ring over $\Lambda$.

\smallskip
\noindent\textbf{(2)} There is a disjoint-union decomposition
\[
J_{\mathfrak X}(\mathfrak X|R)
=\coprod_{1\leq i\leq m}
 J_{\mathfrak X_{(z_i)}}(\mathfrak X_{(z_i)}|R),
\]
and the canonical morphism
\[
J_{\mathfrak X_{(z_i)}}(\mathfrak X_{(z_i)}|R)
 \xrightarrow{\ \sim\ }
J_R(\mathfrak X_{(z_i)}|R)
\]
is an isomorphism.

\smallskip
\noindent\textbf{(3)} Each ideal
$I_R(\mathfrak X_{(z_i)}|R)$ is principal, as is
\[
I_R(\mathfrak X|R)
=\prod_{1\leq i\leq m}I_R(\mathfrak X_{(z_i)}|R).
\]
If $I_R(\mathfrak X_{(z_i)}|R)=(t_i)$, then
$t_1,\ldots,t_m$ can be extended to a system of free formal
generators of $R$ over $\Lambda$.

\smallskip
\noindent\textbf{(4)} The morphism
$\mathfrak X\longrightarrow\operatorname{Spec}(\Lambda)$ is formally
smooth; equivalently,
\[
J_{\mathfrak X}(\mathfrak X|\Lambda)=\varnothing.
\]

%%%%%%%%%%  scan idx 139  |  folio 128  %%%%%%%%%%

\smallskip
\noindent\textbf{(5)} If $X_0$ is a complete irreducible curve, then
\[
\dim\underline M_{X_0}=3g-3+a,
\]
where
\[
g=\dim_kH^1(X_0,\mathcal O_{X_0}),
\qquad
a=\dim_kH^0(X_0,\underline D^0_{X_0}).
\]

\smallskip
\noindent\textit{Proof.}
Assertion (1) follows from (2) and (3). The first assertion in (2)
follows from \SGARefLink{sga7.sga7i.evi.n5-7}{5.7}(2). For each $i$, let $(R_i,\mathfrak Y_i)$ be a
formal versal deformation of
$\operatorname{Spec}(\widehat{\mathcal O}_{X_0,z_i})$. Since
$(R,\mathfrak X_{(\widehat z_i)})$ is such a deformation, there is a
continuous local $\Lambda$-algebra homomorphism
\[
\varphi_i:R_i\longrightarrow R
\]
and a cartesian diagram
\[
\begin{tikzcd}[column sep=huge,row sep=large]
\mathfrak X_{(\widehat z_i)}
  \arrow[r] \arrow[d]
& \mathfrak Y_i \arrow[d]\\
\operatorname{Spf}(R)
  \arrow[r,"{\operatorname{Spf}(\varphi_i)}"']
& \operatorname{Spf}(R_i).
\end{tikzcd}
\]
Consequently,
\[
\begin{aligned}
J_{\mathfrak X_{(z_i)}}(\mathfrak X_{(z_i)}|R)
&=J_{\mathfrak X_{(\widehat z_i)}}
   (\mathfrak X_{(\widehat z_i)}|R)\\
&\simeq
J_{\mathfrak Y_i}(\mathfrak Y_i|R_i)
 \times_{\operatorname{Spf}(R_i)}\operatorname{Spf}(R),
\end{aligned}
\]
and
\[
\begin{aligned}
J_R(\mathfrak X_{(z_i)}|R)
&=J_R(\mathfrak X_{(\widehat z_i)}|R)\\
&\simeq
J_{R_i}(\mathfrak Y_i|R_i)
 \times_{\operatorname{Spf}(R_i)}\operatorname{Spf}(R).
\end{aligned}
\]
By Lemma \SGARefLink{sga7.sga7i.evi.n6-2}{6.2}(1),
\[
J_{\mathfrak Y_i}(\mathfrak Y_i|R_i)
 \longrightarrow J_{R_i}(\mathfrak Y_i|R_i)
\]
is an isomorphism, proving the remaining assertion in (2).

By Lemma \SGARefLink{sga7.sga7i.evi.n6-2}{6.2}(2),
\[
I_{R_i}(\mathfrak Y_i|R_i)=(u_i),
\qquad R_i=\Lambda[[u_i]].
\]
It follows that
\[
I_R(\mathfrak X_{(z_i)}|R)
=I_R(\mathfrak X_{(\widehat z_i)}|R)
=(t_i),
\qquad t_i=\varphi_i(u_i),
\]
and therefore
\[
I_R(\mathfrak X|R)
=\prod_{1\leq i\leq m}I_R(\mathfrak X_{(z_i)}|R)
=\left(\prod_{i=1}^m t_i\right).
\]
By \SGARefLink{sga7.sga7i.evi.n4-13}{4.13}, the natural morphism
\[
\underline M_{X_0}\longrightarrow
\underline M_{X_0,(\widehat z_1)}\times\cdots\times
\underline M_{X_0,(\widehat z_m)}
\]
is smooth. Hence
\[
R_1\widehat\otimes_\Lambda\cdots
 \widehat\otimes_\Lambda R_m\longrightarrow R
\]
is formally smooth. Thus $t_1,\ldots,t_m$ extend to a system of free
formal generators of $R$ over $\Lambda$, proving (1) and (3).

%%%%%%%%%%  scan idx 140  |  folio 129  %%%%%%%%%%

For (4), the projection
\[
\mathfrak X_{(\widehat z_i)}
 \xrightarrow{\ \sim\ }
\mathfrak Y_i\times_{\operatorname{Spf}(R_i)}
 \operatorname{Spf}(R)
 \longrightarrow\mathfrak Y_i
\]
is formally smooth because
$\operatorname{Spf}(R)\longrightarrow\operatorname{Spf}(R_i)$ is
formally smooth. Lemma \SGARefLink{sga7.sga7i.evi.n6-2}{6.2}(3) shows that
$\mathfrak Y_i\longrightarrow\operatorname{Spec}(\Lambda)$ is
formally smooth. Hence every
$\mathfrak X_{(\widehat z_i)}\longrightarrow
\operatorname{Spec}(\Lambda)$ is formally smooth. Since the $z_i$
are precisely the nonsmooth points of $X_0$, it follows that
$\mathfrak X\longrightarrow\operatorname{Spec}(\Lambda)$ is formally
smooth.

% END INCLUDED SGA7I_English_current_buildable_source_20260823/components/60_expose_VI_theorem_6_4_through_part_4.tex

% BEGIN INCLUDED SGA7I_English_current_buildable_source_20260823/components/61_expose_VI_theorem_6_4_part_5_and_remark_6_5.tex
For (5), $R$ is formally smooth over $\Lambda$, and by \SGARefLink{sga7.sga7i.evi.n4-4}{4.4} and \SGARefLink{sga7.sga7i.evi.n4-11}{4.11}
the sequence
\[
0\longrightarrow H^1(X_0,\underline D^0_{X_0})
 \longrightarrow[\underline M_{X_0}(k[\varepsilon])]
 \longrightarrow H^0(X_0,\underline D^1_{X_0})
 \longrightarrow0
\]
is exact. Thus
\[
\begin{aligned}
\dim\underline M_{X_0}
&=\dim_kH^1(X_0,\underline D^0_{X_0})
 +\dim_kH^0(X_0,\underline D^1_{X_0})\\
&=\dim_kH^1(X_0,\underline D^0_{X_0})+m.
\end{aligned}
\]
It is therefore enough to show that
\[
\dim_kH^1(X_0,\underline D^0_{X_0})+m
=3g-3+\dim_kH^0(X_0,\underline D^0_{X_0}),
\]
or equivalently
\[
\chi(\underline D^0_{X_0})=3-3g+m,
\]
where, for a coherent $\mathcal O_{X_0}$-module $\mathcal E$, we put
\[
\chi(\mathcal E)
=\dim_kH^0(X_0,\mathcal E)-\dim_kH^1(X_0,\mathcal E).
\]

The Euler characteristic of $\underline D^0_{X_0}$ is unchanged by
extension of the base field to its algebraic closure, so we may assume
that $k$ is algebraically closed. Let $z$ be a nonsmooth point of
$X_0$, put
\[
A=\widehat{\mathcal O}_{X_0,z},
\]
and let $A'$ be the normalization of $A$. By hypothesis,
\[
A\simeq k[[x,y]]/(f),
\]
where $f$ is a non-degenerate quadratic form. Over the algebraically
closed field $k$ we may take $f=xy$, so that
\[
A\simeq k[[x,y]]/(xy).
\]
It follows at once that
\[
\ell(A'/A)=1,
\qquad
\operatorname{Der}_k(A,\mathfrak m_A)
=\operatorname{Der}_k(A,A).
\]
Let $C$ be the conductor of $A$ in $A'$. Since $A$ is a complete
intersection,
\[
\ell(A/C)=\ell(A'/A)=1,
\]
and hence $C=\mathfrak m_A$. Therefore
\[
A'=\{a'\in K\mid a'\mathfrak m_A\subset\mathfrak m_A\},
\]
where $K$ is the total fraction ring of $A$.

We claim that
\[
\operatorname{Der}_k(A,A)
 \subset\operatorname{Der}_k(A',A')
\]
inside $\operatorname{Der}_k(K,K)$. Indeed, let
$D\in\operatorname{Der}_k(A,A)$. For $a'\in A'$ and
$x\in\mathfrak m_A$, one has $xa'\in\mathfrak m_A$, so
$D(xa')\in\mathfrak m_A$ and
$a'D(x)\in\mathfrak m_A$. The identity
\[
D(xa')=a'D(x)+xD(a')
\]
therefore gives $xD(a')\in\mathfrak m_A$ for every
$x\in\mathfrak m_A$, whence $D(a')\in A'$.

%%%%%%%%%%  scan idx 141  |  folio 130  %%%%%%%%%%

Moreover, the exact sequence
\[
0\longrightarrow\operatorname{Der}_k(A,A)
 \longrightarrow\operatorname{Der}_k(A',A')
 \longrightarrow A'/\mathfrak m_AA'\longrightarrow0
\]
may be written
\[
0\longrightarrow\operatorname{Der}_k(A,A)
 \longrightarrow\operatorname{Der}_k(A',A')
 \longrightarrow A'/C\longrightarrow0.
\]
If $X_0'$ is the normalization of $X_0$, these local sequences give
an exact sequence
\[
0\longrightarrow\underline D^0_{X_0}
 \longrightarrow\underline D^0_{X_0'}
 \longrightarrow\mathcal O_{X_0'}/C\longrightarrow0.
\]
Since $X_0$ is locally a complete intersection,
\[
\ell(\mathcal O_{X_0'}/C)
=2\ell(\mathcal O_{X_0'}/\mathcal O_{X_0})=2m,
\]
whereas the Riemann--Roch theorem gives
\[
\chi(\underline D^0_{X_0'})=3-3(g-m).
\]
Consequently,
\[
\begin{aligned}
\chi(\underline D^0_{X_0})
&=\chi(\underline D^0_{X_0'})
 -\ell(\mathcal O_{X_0'}/C)\\
&=3-3(g-m)-2m=3-3g+m.
\end{aligned}
\]
This completes the proof.
\hfill$\square$

\par\medskip
\noindent\textbf{Remark \SGARefTarget{sga7.sga7i.evi.n6-5}{6.5}.}
In the formula of Theorem \SGARefLink{sga7.sga7i.evi.n6-4}{6.4}(5), the number
$a=\dim_kH^0(X_0,\underline D^0_{X_0})$ is easily computed by the
Riemann--Roch theorem:
\[
a=
\begin{cases}
0,&g\geq2,\\
1,&g=1,\\
3,&g=0.
\end{cases}
\]

We finish with a remark on characteristic $2$. If
$\operatorname{char}k=2$, there is no non-degenerate quadratic form
in an odd number of variables. Equivalently, there is no
non-degenerate quadratic singularity of even dimension in
characteristic $2$. In this case one may instead consider a quadratic
form of maximal non-degeneracy.

% END INCLUDED SGA7I_English_current_buildable_source_20260823/components/61_expose_VI_theorem_6_4_part_5_and_remark_6_5.tex

% BEGIN INCLUDED SGA7I_English_current_buildable_source_20260823/components/62_expose_VI_6_6_through_6_8.tex
%%%%%%%%%%  scan idx 142  |  folio 131  %%%%%%%%%%

\par\medskip
\noindent\textbf{Definition \SGARefTarget{sga7.sga7i.evi.n6-6}{6.6}.}
Let $X$ be a $k$-scheme of finite type. A closed point $z\in X$ is
called an \emph{ordinary quadratic singularity rational over $k$} if
the following conditions hold.

\smallskip
\noindent\textbf{(i)} If either $\operatorname{char}k\neq2$, or
$\operatorname{char}k=2$ and $\dim\mathcal O_{X,z}$ is odd, then $z$
is a non-degenerate quadratic singularity rational over $k$.

\smallskip
\noindent\textbf{(ii)} If $\operatorname{char}k=2$ and
$\dim\mathcal O_{X,z}$ is even, then
\[
\overline k\otimes_k\widehat{\mathcal O}_{X,z}
 \simeq
\overline k[[X_0,X_1,\ldots,X_{2n}]]/(f),
\]
where $\overline k$ is an algebraic closure of $k$ and
\[
f=X_0^2+X_1X_2+X_3X_4+\cdots+X_{2n-1}X_{2n}.
\]

\par\medskip
\noindent\textbf{Lemma \SGARefTarget{sga7.sga7i.evi.n6-7}{6.7}.}
Let $k$ be a field of characteristic $2$, and put
\[
A_0=k[[X_0,X_1,\ldots,X_{2n}]]/(f),
\qquad
f=X_0^2+X_1X_2+X_3X_4+\cdots+X_{2n-1}X_{2n}.
\]
Let $(R,\mathfrak X)$ be a formal versal deformation of
$X_0=\operatorname{Spec}(A_0)$. Then
\[
R=\Lambda[[t,t_1]],
\]
where $t,t_1$ are free formal generators of $R$ over $\Lambda$, and
\[
I_R(\mathfrak X|R)=(t^2).
\]
Moreover,
$\mathfrak X\longrightarrow\operatorname{Spec}(\Lambda)$ is formally
smooth.

\smallskip
\noindent\textit{Proof.}
Write $\mathfrak X=\operatorname{Spf}(A)$. Since
\[
\left(f,\frac{\partial f}{\partial X_0},\ldots,
         \frac{\partial f}{\partial X_{2n}}\right)
=(X_0^2,X_1,\ldots,X_{2n}),
\]
we have
\[
D^1(A_0|k,A_0)\simeq k[X_0]/(X_0^2).
\]
Thus $[\underline M_{X_0}(k[\varepsilon])]$ is a two-dimensional
$k$-vector space. By \SGARefLink{sga7.sga7i.evi.n4-14}{4.14},
\[
R\simeq\Lambda[[t,t_1]],
\qquad
A\simeq R[[X_0,X_1,\ldots,X_{2n}]]/(F),
\]
where
\[
F=f+t_1+tX_0.
\]
The exact sequence
\[
0\longrightarrow A\xrightarrow{\ dF\ }
A\,dX_0\oplus\cdots\oplus A\,dX_{2n}
\longrightarrow\widehat\Omega_{A|R}\longrightarrow0,
\]
in which
\[
dF:\ 1\longmapsto
\sum_{0\leq i\leq2n}\frac{\partial f}{\partial X_i}\,dX_i
+t\,dX_0,
\]
gives
\[
I_A(A|R)
=(t,X_1,\ldots,X_{2n})
=(X_0^2+t_1,t,X_1,\ldots,X_{2n}).
\]
% The printed formula omits the coefficient t before dX_0, although
% F=f+t_1+tX_0 and the ensuing ideal both force this derivative term.

%%%%%%%%%%  scan idx 143  |  folio 132  %%%%%%%%%%

Hence
\[
\begin{aligned}
A/I_A(A|R)
&\simeq A/(X_0^2+t_1,t,X_1,\ldots,X_{2n})\\
&\simeq\Lambda[[t_1,X_0]]/(X_0^2+t_1).
\end{aligned}
\]
As an $R$-module this last module is isomorphic to
$R/(t)\oplus R/(t)$. It follows that
\[
I_R(A|R)=(t^2),
\]
which proves the first assertion. On the other hand,
\[
\begin{aligned}
A
&\simeq
\Lambda[[t,t_1,X_0,X_1,\ldots,X_{2n}]]/
 (X_0^2+X_1X_2+\cdots+X_{2n-1}X_{2n}+tX_0+t_1)\\
&\simeq\Lambda[[t,X_0,X_1,\ldots,X_{2n}]],
\end{aligned}
\]
so $A$ is formally smooth over $\Lambda$.
\hfill$\square$

Lemma \SGARefLink{sga7.sga7i.evi.n6-7}{6.7} gives the following analogue of Theorem \SGARefLink{sga7.sga7i.evi.n6-4}{6.4}; we leave its
proof to the reader.

\par\medskip
\noindent\textbf{Corollary \SGARefTarget{sga7.sga7i.evi.n6-8}{6.8}.}
Let $X_0$ be a $k$-scheme of finite type whose nonsmooth points are
isolated, assume that
\[
H^2(X_0,\underline D^0_{X_0})=0,
\]
and let $(R,\mathfrak X)$ be a formal versal deformation of $X_0$.
Denote the set of nonsmooth points by $\{z_1,\ldots,z_m\}$. Suppose
that $\operatorname{char}k=2$, that $\dim X_0$ is even, and that all
the $z_i$ are ordinary quadratic singularities. Then:

\smallskip
\noindent\textbf{(1)} $R$ is formally smooth over $\Lambda$.

\smallskip
\noindent\textbf{(2)} Each ideal
$I_R(\mathfrak X_{(z_i)}|R)$ is principal, as is
\[
I_R(\mathfrak X|R)
=\prod_{1\leq i\leq m}I_R(\mathfrak X_{(z_i)}|R).
\]
More precisely,
\[
I_R(\mathfrak X_{(z_i)}|R)=(t_i^2),
\]
and $t_1,\ldots,t_m$ can be extended to a system of free formal
generators of $R$ over $\Lambda$.

\smallskip
\noindent\textbf{(3)} The morphism
$\mathfrak X\longrightarrow\operatorname{Spec}(\Lambda)$ is formally
smooth.

\smallskip
\noindent\textbf{(4)} If $k$ is algebraically closed and $X_0$ is a
complete irreducible curve, then
\[
\dim\underline M_{X_0}=3g-3+a,
\]
where
\[
g=\dim_kH^1(X_0,\mathcal O_{X_0}),
\qquad
a=\dim_kH^0(X_0,\underline D^0_{X_0}).
\]

% END INCLUDED SGA7I_English_current_buildable_source_20260823/components/62_expose_VI_6_6_through_6_8.tex

\clearpage
\phantomsection\addcontentsline{toc}{section}{Exposé VII}
% BEGIN INCLUDED SGA7I_English_current_buildable_source_20260823/components/63_expose_VII_opening_through_1_1_2_1.tex
% SOURCE AUTHORITY:
% expose_VII_body.tex lines 5--58; scan indices 144--147; source folios 133--136.
% Hard stop before source line 60, paragraph 1.1.3.

%%%%%%%%%%  scan idx 144  |  folio 133  %%%%%%%%%%

\begin{center}
\SGARefAliasTarget{sga7.sga7i.evii.nroot}SGA 7

\vspace{1em}
\underline{Expos\'e VII}

\vspace{1em}
\underline{BIEXTENSIONS OF SHEAVES OF GROUPS}

\vspace{1em}
\underline{by A. Grothendieck}
\end{center}

0. \underline{Introduction}

The notion of a \underline{biextension} was introduced by D.~Mumford
[6], in the context of formal groups, as a convenient way of expressing
the relations between the formal groups associated with an abelian scheme
and the dual abelian scheme. In the present expos\'e we undertake a
systematic study of this notion in the general context of sheaves of
groups on a given site. In particular, we show that it fits remarkably
simply into the classical cohomological formalism of
$\overset{\mathrm L}{\otimes}$ and $\underline{\mathrm{RHom}}$
(\SGARefLink{sga7.sga7i.evii.n3-6-4}{3.6.4} and \SGARefLink{sga7.sga7i.evii.n3-6-5}{3.6.5}). In the following expos\'e we shall make explicit the
most interesting special cases, notably that of biextensions of group
schemes by the multiplicative group $\underline G_m$ (working with such
sites as the Zariski site, the fpqc site, or an intermediate site). We
shall see, in particular, that if $A$ and $B$ are two abelian schemes over
a scheme $S$, biextensions of $(A,B)$ by $\underline G_m$ admit a
remarkable interpretation as, essentially, the classical divisorial
correspondences on $A,B$. This gives a particularly convenient explicit
cohomological meaning to the notion of divisorial correspondence (already
implicit in Tate duality for abelian varieties over $p$-adic fields).
This interpretation of divisorial correspondences is also convenient for
following the fate of such a correspondence when $A$ and $B$ (assumed to
be defined over the fraction field of a discrete valuation ring) are made
to degenerate by means of N\'eron theory.

%%%%%%%%%%  scan idx 145  |  folio 134  %%%%%%%%%%

The results of Expos\'e VIII will show, for example, that one obtains a
biextension of the pair of identity components of the N\'eron models,
hence a biextension of their connected special fibres, which plays the
role of a specialization of the original divisorial correspondence.
These results will be used in Expos\'e IX, in conjunction with the
monodromy theorem (I~3), to prove certain results about the N\'eron models
of abelian varieties, the most important being the ``semistable reduction
theorem'' for the N\'eron model.

Throughout what follows, unless expressly stated otherwise, we work in a
fixed $\underline U$-topos denoted by $\underline T$; that is (SGA~4,
Expos\'e~IV), $\underline T$ is a category equivalent to the category of
sheaves on a suitable $\underline U$-site (which may, if desired, be taken
to be $\underline T$ itself, endowed with its canonical topology). In our
terminology we shall often identify the objects of $\underline T$ with
the sheaves of sets on $\underline T$. All groups, extensions of groups,
and so forth considered below are relative to this topos $\underline T$.

1. \underline{Complements on extensions of groups}

\SGARefTarget{sga7.sga7i.evii.n1-1}{1.1}. Let $P$ and $G$ be two groups of $\underline T$, and let

\begin{equation}\tag{\SGARefTarget{sga7.sga7i.evii.n1-1-1}{1.1.1}}
\begin{tikzcd}[column sep=large]
1 \arrow[r] & G \arrow[r, "i"] & E \arrow[r, "j"] & P \arrow[r] & 1
\end{tikzcd}
\end{equation}

be an extension of $P$ by $G$. Thus $E$ is a group,
$j:E\longrightarrow P$ is an epimorphism of groups, and $i$ induces an
isomorphism of $G$ with $\ker j$. We shall reinterpret the datum of such
an extension in different terms, which will be convenient for the
definition and study of biextensions.

%%%%%%%%%%  scan idx 146  |  folio 135  %%%%%%%%%%

\SGARefTarget{sga7.sga7i.evii.n1-1-2}{1.1.2}. First, the homomorphism $i$ makes $G$ act on $E$ in two ways, on
the left and on the right, by setting
\[
 gxg'=i(g)xi(g')
\]
for $g,g'\in G(S)$ and $x\in E(S)$, where $S$ is any object of
$\underline T$. If, by means of $j$, we regard $E$ as an object ``over
$P$'', that is, as an object of the induced topos $\underline T_{/P}$,
this amounts to saying that the group $G_P=G\times P$ induced by $G$ in
$\underline T_{/P}$ acts on $E$ on the left and on the right. Moreover,
it is well known that each of these actions makes $E$ a
\underline{torsor} (a left torsor and a right torsor, respectively) over
$P$, with group $G_P$; more precisely still, it makes $E$ a
\underline{bitorsor} under $(G_P,G_P)$. Recall [2] that this means that
$E$ is a right torsor $E_d$ under the right action of $G_P$, and that the
left action induces an isomorphism of $G_P$ with the group of
automorphisms of $E_d$ (the sheaf of groups of automorphisms of $E$ that
commute with the right action of $G_P$ on $E$). This condition is
equivalent to the symmetric condition: the left action of $G_P$ on $E$
makes $E$ a left torsor $E_s$, and the right action induces an
isomorphism of the opposite group $G_P^{\mathrm o}$ with the group of
automorphisms of $E_s$.

\SGARefTarget{sga7.sga7i.evii.n1-1-2-1}{1.1.2.1}. When $\underline T$ is the ``point topos'', that is, the
category of sets (equivalent to the category of sheaves on the one-point
topological space), so that $P$ and $G$ are ordinary groups, the datum of
a bitorsor $E$ under $(G_P,G_P)$ is plainly equivalent to the datum, for
each $p\in P$, of a bitorsor $E_p$ under $(G,G)$, namely the fibre of $E$
at $p$. An analogous interpretation is available in the general case.
Indeed, giving a bitorsor $E$ under $(G_P,G_P)$ is equivalent to giving,
for every ``point'' $p\in P(S)$, where $S$ is any object of
$\underline T$, a bitorsor $E_p$ under $(G_S,G_S)$ over $S$, in such a
way that the formation of $E_p$ is ``compatible with base change along
every morphism $S'\longrightarrow S$''. Thus, if
$p':S'\longrightarrow P$ is the composite morphism, one is also given an
isomorphism of bitorsors from $E_{p'}$ to the inverse image of $E_p$ under
$S'\longrightarrow S$, and these isomorphisms satisfy the usual
transitivity condition for a composite
$S''\longrightarrow S'\longrightarrow S$.

%%%%%%%%%%  scan idx 147  |  folio 136  %%%%%%%%%%

When stated in canonical language (SGA~1, Expos\'e~VI), this reduces to
the following tautology: giving a bitorsor under $(G_P,G_P)$ over $P$ is
equivalent to giving a cartesian section, over the category
$\underline T_{/P}$, of the fibred category of bitorsors under
$(G_S,G_S)$ over a variable base $S\in\operatorname{ob}\underline
T_{/P}$. More precisely, the ``value at $P$'' functor induces an
equivalence between the category of the cartesian sections in question
and the category of bitorsors under $(G_P,G_P)$. This follows trivially
from the fact that $P$ is a final object of the base category
$\underline T_{/P}$ of our fibred category: for \underline{every}
category fibred over $\underline T_{/P}$, the ``value at $P$'' functor is
an equivalence between the category of cartesian sections of that fibred
category and its fibre category over $P$. Despite its tautological
nature, however, this interpretation will be useful for expressing
certain compatibilities in language close to the set-theoretic case,
using points with values in an object $S$ in place of ordinary points.


% END INCLUDED SGA7I_English_current_buildable_source_20260823/components/63_expose_VII_opening_through_1_1_2_1.tex

% BEGIN INCLUDED SGA7I_English_current_buildable_source_20260823/components/64_expose_VII_1_1_3_through_1_1_4_2.tex
% SOURCE AUTHORITY:
% expose_VII_body.tex lines 60--137; scan indices 147--151; source folios 136--140.
% Hard stop before source line 139, paragraph 1.1.5.
% The print intermittently writes P_p where its own definitions and the
% associativity diagram require E_p; the coherent E_p notation is retained.

\SGARefTarget{sga7.sga7i.evii.n1-1-3}{1.1.3}. In terms of the extension structure (\SGARefLink{sga7.sga7i.evii.n1-1-1}{1.1.1}), the bitorsor
structure of \SGARefLink{sga7.sga7i.evii.n1-1-2}{1.1.2} on $E$ under $(G_P,G_P)$ used only part of the
multiplicative structure of $E$: namely, multiplication of pairs of
``points'' of which at least one comes from a point of $G$ (via $i$). We
shall now interpret the multiplicative structure of $E$ as an additional
structure on the bitorsor under consideration.

%%%%%%%%%%  scan idx 148  |  folio 137  %%%%%%%%%%

For this purpose, first consider two sections $p,p'$ of $P$. With the
notation of \SGARefLink{sga7.sga7i.evii.n1-1-2-1}{1.1.2.1}, these give two bitorsors $E_p,E_{p'}$ under $(G,G)$,
namely the inverse images in $E$ of the sections
$p,p':e\longrightarrow P$ (where $e$ denotes the final object of
$\underline T$). It is then clear that multiplication in $E$ induces a
morphism of objects of $\underline T$
\[
 \varphi:E_p\times E_{p'}\longrightarrow E_{pp'},
\]
which, in terms of the left and right structures of the three bitorsors
involved (and, as usual, using points with values in an arbitrary object
$S\in\operatorname{ob}\underline T$), satisfies
\begin{equation}\tag{\SGARefTarget{sga7.sga7i.evii.n1-1-3-0}{1.1.3.0}}
\left\{
\begin{aligned}
 \varphi(gx,x')&=g\varphi(x,x'),
 &\qquad \varphi(x,x'g')&=\varphi(x,x')g',\\
 \varphi(xg,x')&=\varphi(x,gx').
\end{aligned}
\right.
\end{equation}
Here $g,g'\in G(S)$, $x\in E_p(S)$, and $x'\in E_{p'}(S)$. In terms of
the contracted product [2]
\[
 E_p\mathbin{\overset{G}{\times}}E_{p'},
\]
the third relation says that $\varphi$ factors as
\[
 E_p\times E_{p'}\longrightarrow
 E_p\mathbin{\overset{G}{\times}}E_{p'}
 \xrightarrow{\ \varphi_{p,p'}\ }E_{pp'},
\]
where the first arrow is the canonical quotient morphism. The first two
relations, on the other hand, say that the morphism thus obtained,
\begin{equation}\tag{\SGARefTarget{sga7.sga7i.evii.n1-1-3-1}{1.1.3.1}}
 \varphi_{p,p'}:
 E_p\mathbin{\overset{G}{\times}}E_{p'}\longrightarrow E_{pp'},
\end{equation}
is a homomorphism of \underline{bitorsors}, where the bitorsor structure
on the source is the one for which the left (respectively, right)
operations of $G$ come from the left (respectively, right) operations of
$G$ on the first factor $E_p$ (respectively, on the second factor
$E_{p'}$).

%%%%%%%%%%  scan idx 149  |  folio 138  %%%%%%%%%%

Notice that, like every homomorphism of bitorsors under a fixed pair of
groups, $\varphi_{p,p'}$ is necessarily an \underline{isomorphism}.

The definition of the preceding isomorphism (\SGARefLink{sga7.sga7i.evii.n1-1-3-1}{1.1.3.1}) immediately extends
to the case in which $p,p'$ are points of $P$ with values in an arbitrary
object $S$ of $\underline T$, rather than sections of $P$: one simply
applies the preceding discussion to the induced topos
$\underline T_{/S}$ and to the extension of groups over $S$ obtained from
(\SGARefLink{sga7.sga7i.evii.n1-1-1}{1.1.1}) by base change. The formation of the isomorphisms
$\varphi_{p,p'}$ is then compatible with every base change
$S'\longrightarrow S$. In other words, it corresponds to a homomorphism
of cartesian sections of the fibred category over the category
$\underline T_{/P\times P}$ of pairs of arrows
$p,p':S\rightrightarrows P$ with target $P$.

Giving a system of $\varphi_{p,p'}$ satisfying this condition is, of
course, equivalent to giving $\varphi_{p,p'}$ in the ``universal'' case
$S=P\times P$, where $p,p'$ are the two projections
$\operatorname{pr}_i$ ($i=1,2$) from $P\times P$ to $P$. It is then an
isomorphism of bitorsors under
$(G_{P\times P},G_{P\times P})$:
\begin{equation}\tag{\SGARefTarget{sga7.sga7i.evii.n1-1-3-2}{1.1.3.2}}
 \varphi:\operatorname{pr}_1^*(E)\,\operatorname{pr}_2^*(E)
 \overset{\sim}{\longrightarrow}\pi^*(E),
\end{equation}
where
\[
 \pi:P\times P\longrightarrow P
\]
is the morphism of objects of $\underline T$ defining the multiplicative
structure of the group $P$. The compatibility properties below are most
conveniently stated, however, in terms of the data (\SGARefLink{sga7.sga7i.evii.n1-1-3-1}{1.1.3.1}), avoiding
diagrams that are less close to intuition.

%%%%%%%%%%  scan idx 150  |  folio 139  %%%%%%%%%%

\SGARefTarget{sga7.sga7i.evii.n1-1-4}{1.1.4}. In writing the relations (\SGARefLink{sga7.sga7i.evii.n1-1-3-0}{1.1.3.0}), which enabled us to express
the given composition law on $E$ in terms of isomorphisms of bitorsors
(\SGARefLink{sga7.sga7i.evii.n1-1-3-1}{1.1.3.1}), we used associativity of this composition law only for a
product of three ``points'' of $E$ of which at least one comes from a
point of $G$. To express full associativity in terms of the
$\varphi_{p,p'}$, one must require, for three points $p,p',p''$ of $P$
with values in an object $S$ of $\underline T$, that the following
diagram be commutative:
\begin{equation}\tag{\SGARefTarget{sga7.sga7i.evii.n1-1-4-1}{1.1.4.1}}
\begin{tikzcd}[column sep=huge, row sep=large]
 & E_{pp'}E_{p''} \arrow[dr, "{\varphi_{pp',p''}}"] & \\
 E_pE_{p'}E_{p''}
   \arrow[ur, "{\varphi_{p,p'}\wedge\id}"]
   \arrow[dr, "{\id\wedge\varphi_{p',p''}}"']
 && E_{pp'p''} \\
 & E_pE_{p'p''} \arrow[ur, "{\varphi_{p,p'p''}}"'] &
\end{tikzcd}
\end{equation}
In this diagram we have written the contracted product of bitorsors
without the signs $\mathbin{\overset{G_S}{\times}}$. It is understood
that the contracted product $Q.R$ of two bitorsors is formed using, in
passing to the quotient in $Q\times R$, the action of the structural
group $G$ \underline{on the right of} $Q$ \underline{and on the left of}
$R$. This contracted-product operation is associative up to a canonical
isomorphism [2], and we make the identification
\[
 (E_pE_{p'})E_{p''}\simeq E_p(E_{p'}E_{p''}),
\]
writing both sides simply as $E_pE_{p'}E_{p''}$. Obvious conventions of
this kind will be used throughout the rest of this expos\'e without being
expressly noted each time.

\SGARefTarget{sga7.sga7i.evii.n1-1-4-2}{1.1.4.2}. Just as the isomorphisms $\varphi_{p,p'}$ of (\SGARefLink{sga7.sga7i.evii.n1-1-3-1}{1.1.3.1}) are
determined by the ``universal'' case $S=P\times P$,
$p=\operatorname{pr}_1$, $p'=\operatorname{pr}_2$ considered in
(\SGARefLink{sga7.sga7i.evii.n1-1-3-2}{1.1.3.2}), so the associativity condition (\SGARefLink{sga7.sga7i.evii.n1-1-4-1}{1.1.4.1}) may likewise be
checked only in the universal case
$S=P\times P\times P$, $p=\operatorname{pr}_1$,
$p'=\operatorname{pr}_2$, $p''=\operatorname{pr}_3$. We leave it to the
reader to write the diagram in whatever notation is preferred.


% END INCLUDED SGA7I_English_current_buildable_source_20260823/components/64_expose_VII_1_1_3_through_1_1_4_2.tex

% BEGIN INCLUDED SGA7I_English_current_buildable_source_20260823/components/65_expose_VII_1_1_5_through_1_1_6.tex
% SOURCE AUTHORITY:
% expose_VII_body.tex lines 139--235; scan indices 151--154; source folios 140--143.
% Hard stop before source line 237, paragraph 1.2.

\SGARefTarget{sga7.sga7i.evii.n1-1-5}{1.1.5}. Conversely, suppose that the groups $P$ and $G$ of
$\underline T$ are given, together with a bitorsor
\begin{equation}\tag{\SGARefTarget{sga7.sga7i.evii.n1-1-5-1}{1.1.5.1}}
 j:E\longrightarrow P
\end{equation}
under $(G_P,G_P)$ (the two actions of $G_P$ on $E$, or equivalently of
$G$ on $E$, being understood in the notation), and a system of
isomorphisms of bitorsors over a variable base $S$,
\begin{equation}\tag{\SGARefTarget{sga7.sga7i.evii.n1-1-5-2}{1.1.5.2}}
 \varphi_{p,p'}:E_pE_{p'}\overset{\sim}{\longrightarrow}E_{pp'}
 \qquad (p,p'\in P(S)),
\end{equation}
compatible with base change $S'\longrightarrow S$; in other words, an
isomorphism of bitorsors (\SGARefLink{sga7.sga7i.evii.n1-1-3-2}{1.1.3.2}). Suppose that the corresponding
diagrams (\SGARefLink{sga7.sga7i.evii.n1-1-4-1}{1.1.4.1}) commute. I claim that \underline{there is a unique
extension structure on $E$ of $P$ by $G$ such that the preceding data
are obtained from it by the constructions of \SGARefLink{sga7.sga7i.evii.n1-1-2}{1.1.2} and \SGARefLink{sga7.sga7i.evii.n1-1-3}{1.1.3}}.

First observe that, on taking $p,p'$ in (\SGARefLink{sga7.sga7i.evii.n1-1-5-2}{1.1.5.2}) to be the identity
section $1$ of $P$, one obtains an isomorphism
\begin{equation}\tag{\SGARefTarget{sga7.sga7i.evii.n1-1-5-3}{1.1.5.3}}
 \varphi_{1,1}:E_1\overset{\sim}{\longleftarrow}
 E_1\mathbin{\overset{G}{\times}}E_1.
\end{equation}
Taking the contracted product, on the left say, of both sides with the
inverse bitorsor $E_1^{-1}$ of $E_1$ [2] therefore gives a canonical
isomorphism
\begin{equation}\tag{\SGARefTarget{sga7.sga7i.evii.n1-1-5-4}{1.1.5.4}}
 E_1\simeq\text{the unit bitorsor }{}_sG_d
\end{equation}
(that is, $G$ acting on itself by left and right translations), in view
of the canonical isomorphism

%%%%%%%%%%  scan idx 152  |  folio 141  %%%%%%%%%%

\[
 E_1^{-1}.E_1\simeq\text{the unit bitorsor }{}_sG_d
\]
of loc.~cit., valid for every bitorsor. Equivalently, one has a section
$1_E$ of $E_1$, hence of $E$,
\begin{equation}\tag{\SGARefTarget{sga7.sga7i.evii.n1-1-5-5}{1.1.5.5}}
 1_E\in\Gamma(E_1)=\operatorname{Hom}(e,E_1)\subset\Gamma(E),
\end{equation}
namely the image under (\SGARefLink{sga7.sga7i.evii.n1-1-5-4}{1.1.5.4}) of the distinguished section of
${}_sG_d$, satisfying
\[
 g.1_E=1_E.g\qquad(g\in G(S)).
\]
By construction this section is characterized by
\begin{equation}\tag{\SGARefTarget{sga7.sga7i.evii.n1-1-5-6}{1.1.5.6}}
 \varphi_{1,1}(1_E,1_E)=1_E.
\end{equation}

On the other hand, giving the morphism of torsors $\varphi$ of
(\SGARefLink{sga7.sga7i.evii.n1-1-3-2}{1.1.3.2}) is plainly equivalent to giving a morphism
\begin{equation}\tag{\SGARefTarget{sga7.sga7i.evii.n1-1-5-7}{1.1.5.7}}
 \pi_E:E\times E\longrightarrow E,
 \qquad\text{also denoted }(x,y)\longmapsto xy,
\end{equation}
satisfying
\begin{equation}\tag{\SGARefTarget{sga7.sga7i.evii.n1-1-5-8}{1.1.5.8}}
 j(xy)=j(x)j(y)\qquad(x,y\in E(S)),
\end{equation}
so that $j$ is compatible with the multiplicative structures $\pi_E$
on $E$ and $\pi$ on $P$, together with the three relations (\SGARefLink{sga7.sga7i.evii.n1-1-3-0}{1.1.3.0})
(preferably beginning with the last of these). Finally define a morphism
of objects of $\underline T$,
\[
 i:G\longrightarrow E,
\]
by the formula
\begin{equation}\tag{\SGARefTarget{sga7.sga7i.evii.n1-1-5-9}{1.1.5.9}}
 i(g)=g.1_E=1_E.g.
\end{equation}

%%%%%%%%%%  scan idx 153  |  folio 142  %%%%%%%%%%

The equality of the last two terms was already observed in constructing
$1_E$. I now claim that \underline{the multiplication law $\pi_E$ on
$E$ makes $E$ a group with identity section $1_E$, and that
$j:E\longrightarrow P$ and $i:G\longrightarrow E$ are homomorphisms of
groups making $E$ an extension of $P$ by $G$}.

Indeed, associativity of the composition law on $E$ is precisely the
commutativity of the diagrams (\SGARefLink{sga7.sga7i.evii.n1-1-4-1}{1.1.4.1}). The construction of $\pi_E$ by
means of the isomorphisms $\varphi_{p,p'}$ of (\SGARefLink{sga7.sga7i.evii.n1-1-3-1}{1.1.3.1}) also shows that,
for every $x\in E(S)$, left and right translation by $x$ in $E(S)$ are
injective. Since $1_E1_E=1_E$ by (\SGARefLink{sga7.sga7i.evii.n1-1-5-6}{1.1.5.6}), it follows at once that
$1_E$ is a two-sided identity of $E$. We leave to the reader the
verification that an $x\in E(S)$ lying over $p\in P(S)$ has an inverse;
one uses the canonical isomorphism deduced from $\varphi_{p,p^{-1}}$,
\begin{equation}\tag{\SGARefTarget{sga7.sga7i.evii.n1-1-5-10}{1.1.5.10}}
 E_{p^{-1}}\simeq(E_p)^{\mathrm o}.
\end{equation}
Thus $E$ is indeed a \underline{group}. We have already observed that
$j$ is compatible with the composition laws (\SGARefLink{sga7.sga7i.evii.n1-1-5-8}{1.1.5.8}), and the same is
true of $i$, since, for $g,g'\in G(S)$,
\begin{align*}
 i(g)i(g')
 &= (g.1_E)(g'.1_E)\\
 &= (g.1_E.g')1_E &&\text{(the last relation of (\SGARefLink{sga7.sga7i.evii.n1-1-3-0}{1.1.3.0}))}\\
 &= g.1_E.g' &&\text{(since $1_E$ is an identity)}\\
 &= g.(1_E.g')=g.(g'.1_E)\\
 &= (gg').1_E &&\text{by (\SGARefLink{sga7.sga7i.evii.n1-1-5-9}{1.1.5.9})}\\
 &= i(gg').
\end{align*}
Since $j$ is the structural projection of a torsor, it is an
epimorphism of objects of $\underline T$. It remains only to see that
$i$ induces an isomorphism of $G$ with $\ker j=E_1$, which is clear from
the definition (\SGARefLink{sga7.sga7i.evii.n1-1-5-9}{1.1.5.9}) and the construction of $1_E$.

%%%%%%%%%%  scan idx 154  |  folio 143  %%%%%%%%%%

\SGARefTarget{sga7.sga7i.evii.n1-1-6}{1.1.6}. We shall accept that the preceding arguments prove the following.
For a fixed topos $\underline T$, the preceding constructions give an
\underline{equivalence between the category of extensions of groups of
$\underline T$} as in (\SGARefLink{sga7.sga7i.evii.n1-1-1}{1.1.1}), where $P,G,E$ are all variable,
\underline{and the category of systems $(P,G,E,\varphi)$ in which $P$
and $G$ are groups of $\underline T$, $E$ is a bitorsor under
$(G_P,G_P)$ over $P$, and $\varphi$ is an isomorphism of bitorsors}
(\SGARefLink{sga7.sga7i.evii.n1-1-3-2}{1.1.3.2}) (or, equivalently, a system of base-change-compatible
isomorphisms of bitorsors (\SGARefLink{sga7.sga7i.evii.n1-1-3-1}{1.1.3.1})), \underline{satisfying the
associativity condition expressed by commutativity of diagram (\SGARefLink{sga7.sga7i.evii.n1-1-4-2}{1.1.4.2})}
(or, equivalently, of the diagrams (\SGARefLink{sga7.sga7i.evii.n1-1-4-1}{1.1.4.1})).

We leave it to the reader to spell out, using the preceding dictionary,
the evident notion of inverse image (respectively, direct image) of an
extension (\SGARefLink{sga7.sga7i.evii.n1-1-1}{1.1.1}) under a homomorphism of groups
\[
 P'\longrightarrow P
 \qquad\bigl(\text{respectively, }G\longrightarrow G'\bigr),
\]
or under both operations simultaneously. We likewise leave to the reader
the compatibility of the preceding equivalence of categories with the
inverse-image functor $f^*$ associated with a morphism of topoi
\[
 f:\underline T'\longrightarrow\underline T.
\]
For this one should use the datum $\varphi$ in the form (\SGARefLink{sga7.sga7i.evii.n1-1-3-2}{1.1.3.2}), in
preference to the $\varphi_{p,p'}$ of (\SGARefLink{sga7.sga7i.evii.n1-1-3-1}{1.1.3.1}); it then defines the
corresponding datum $f^*(\varphi)$ on the topos $\underline T'$.


% END INCLUDED SGA7I_English_current_buildable_source_20260823/components/65_expose_VII_1_1_5_through_1_1_6.tex

% BEGIN INCLUDED SGA7I_English_current_buildable_source_20260823/components/66_expose_VII_section_1_2.tex
% SOURCE AUTHORITY:
% expose_VII_body.tex lines 237--261; scan indices 154--155; source folios 143--144.
% Hard stop before source line 264, section 1.3.

\SGARefTarget{sga7.sga7i.evii.n1-2}{1.2}. \underline{The commutative case}. From \SGARefLink{sga7.sga7i.evii.n1-4}{1.4} onward we shall be
concerned only with \underline{commutative} extensions of groups. We
therefore assume that $P$ and $G$ are commutative, and spell out here,
in terms of the bitorsor description $(P,G,E,\varphi)$ of an extension,
the condition that $E$ be a commutative group.

%%%%%%%%%%  scan idx 155  |  folio 144  %%%%%%%%%%

An obviously necessary first condition---which in fact expresses that
$i(G)$ be a \underline{central} subgroup of $E$---is that the bitorsors
$E_p$, for $p\in P(S)$, be ordinary torsors; that is, the left and right
actions of $G$ on $E_p$ must coincide. Equivalently, $E$ over $P$ is an
ordinary torsor under $G_P$. (Thus there will henceforth be no need to
speak of bitorsors, and we shall consider only ordinary torsors.) Assuming
this preliminary condition, commutativity of $E$ is expressed by
commutativity of the following diagrams, for $p,p'\in P(S)$:
\begin{equation}\tag{\SGARefTarget{sga7.sga7i.evii.n1-2-1}{1.2.1}}
\begin{tikzcd}[column sep=large, row sep=large]
 E_pE_{p'} \arrow[r, "{\varphi_{p,p'}}"] \arrow[d, "\mathrm{sym}"']
   & E_{pp'} \arrow[d, "\id"] \\
 E_{p'}E_p \arrow[r, "{\varphi_{p',p}}"]
   & E_{p'p}.
\end{tikzcd}
\end{equation}
Here the first vertical arrow is the canonical commutativity isomorphism
(defined because $G$ is now commutative); moreover $pp'=p'p$ because
$P$ is commutative. This condition is plainly equivalent to the
corresponding condition in the universal situation
$S=P\times P$, $p=\operatorname{pr}_1$,
$p'=\operatorname{pr}_2$, which we need not write out again.


% END INCLUDED SGA7I_English_current_buildable_source_20260823/components/66_expose_VII_section_1_2.tex

% BEGIN INCLUDED SGA7I_English_current_buildable_source_20260823/components/67_expose_VII_1_3_1_through_1_3_3.tex
% SOURCE AUTHORITY:
% expose_VII_body.tex lines 264--286; scan indices 155--158; source folios 144--147.
% Hard stop before source line 288, paragraph 1.3.4.

\SGARefTarget{sga7.sga7i.evii.n1-3}{1.3}. \underline{Relations with rigidified or birigidified torsors}

\SGARefTarget{sga7.sga7i.evii.n1-3-1-0}{\SGARefAliasTarget{sga7.sga7i.evii.n1-3-1}1.3.1.0}. Let $P$ be an object of $\underline T$ equipped with a section
$e_P$, let $G$ be a group of $\underline T$, and let $E$ be a torsor
under $G_P$. A \underline{rigidification of $E$}, relative to the given
section $e_P$ of $P$, is a trivialization of the torsor
$E_1=e_P^*(E)$ under $G$;

%%%%%%%%%%  scan idx 156  |  folio 145  %%%%%%%%%%

that is, a section of this torsor. This terminology is inspired by the
case where $G$ is commutative and
\begin{equation}\tag{\SGARefTarget{sga7.sga7i.evii.n1-3-1-1}{1.3.1.1}}
 \operatorname{Hom}_{\mathrm{pt}}(P,G)\simeq(e),
\end{equation}
where the left-hand side denotes the set of morphisms carrying $e_P$ to
the identity section $e_G$ of $G$. This relation is equivalent to saying
that the natural homomorphism
\begin{equation}\tag{\SGARefTarget{sga7.sga7i.evii.n1-3-1-2}{1.3.1.2}}
 \Gamma(G)\xrightarrow{\ \sim\ }\operatorname{Hom}(P,G),
\end{equation}
which associates with a section $g$ of $G$ the constant morphism
$P\longrightarrow G$ with value $g$, is an isomorphism. It is immediate
that these equivalent conditions are also equivalent to saying that
every automorphism of the torsor $E$ preserving the given
rigidification is the identity automorphism (such automorphisms
correspond precisely to the elements of the left-hand side of
(\SGARefLink{sga7.sga7i.evii.n1-3-1-1}{1.3.1.1})). In particular, under condition (\SGARefLink{sga7.sga7i.evii.n1-3-1-1}{1.3.1.1}), the category of
torsors under $G_P$ rigidified relative to $e_P$ is ``rigid''.

\SGARefTarget{sga7.sga7i.evii.n1-3-2}{1.3.2}. Let, in addition, $Q$ be an object of $\underline T$ equipped
with a section $e_Q$, and now let $E$ denote a torsor under
$G_{P\times Q}$. Since $e_Q$ defines a section
$e_{P\times Q/P}$ of $P\times Q$ over $P$, one may consider
rigidifications $\alpha$ of $E$ relative to this section, working in the
induced topos $\underline T_{/P}$; explicitly, these are sections of the
torsor induced by $E$ on $P\times e_Q$. Interchanging the roles of $P$
and $Q$, one may likewise consider rigidifications $\beta$ of $E$
relative to the section $e_{P\times Q/Q}$ of $P\times Q$ over $Q$
defined by $e_P$. Two such partial rigidifications are said to be
\underline{compatible} if the sections $\alpha$ and $\beta$ agree on
$e_P\times e_Q$, that is, if they define the same section of the torsor
$E_{1,1}$ under $G$, the inverse image

%%%%%%%%%%  scan idx 157  |  folio 146  %%%%%%%%%%

of $E$ under the section $(e_P,e_Q)$ of $P\times Q$. A
\underline{birigidification} of the torsor $E$ relative to the pair of
sections $(e_P,e_Q)$ is a pair of compatible partial rigidifications.

If $G$ is commutative, then $E$ admits a birigidification
$(\alpha,\beta)$ if and only if it admits partial rigidifications
$\alpha$ and $\beta$; moreover one may find a birigidification in which
$\alpha$ (or $\beta$) is prescribed in advance. Indeed, it suffices to
correct $\beta$ by the morphism $Q\longrightarrow G$ obtained as the
composite
\[
 Q\longrightarrow e\xrightarrow{\lambda}G,
\]
where $\lambda$ is defined by
$\beta e_Q=(\alpha e_P)\lambda$.

Still in the commutative case, if $E,E'$ are two birigidified torsors
under $G_{P\times Q}$, saying that they are isomorphic merely as
$G_{P\times Q}$-torsors is equivalent to saying that they are
isomorphic as birigidified $G_{P\times Q}$-torsors. In other words, if
there is an isomorphism $u:E\longrightarrow E'$ of
$G_{P\times Q}$-torsors, then there is one
$v:E\longrightarrow E'$ that preserves the birigidifications. Indeed,
measuring the discrepancy of $u$ from the given birigidifications gives
morphisms $f:P\longrightarrow G$ and $g:Q\longrightarrow G$ such that
$f(e_P)=g(e_Q)=a\in G(e)$. There is then a morphism
$P\times Q\longrightarrow G$ whose restrictions along
$e_{P\times Q/P}$ and $e_{P\times Q/Q}$ are respectively $f$ and $g$,
namely
\[
 h(p,q)=f(p)g(q)a^{-1},
\]
and one may take $v=uh$.

\SGARefTarget{sga7.sga7i.evii.n1-3-3}{1.3.3}. The terminology of \SGARefLink{sga7.sga7i.evii.n1-3-1}{1.3.1} extends immediately to the case of a
\underline{bitorsor} $E$ under $(G_P,G_P)$. A rigidification of the
bitorsor $E$ is then required to be an isomorphism from $e_P^*(E)$ to
the trivial bitorsor ${}_sG_d$; equivalently, it is defined by a
\underline{central} section $e_E$ of $e_P^*(E)$, one satisfying
\[
 g.e_E=e_E.g\qquad(g\in G(S)).
\]
The terminology of birigidifications extends in the same way to
bitorsors.


% END INCLUDED SGA7I_English_current_buildable_source_20260823/components/67_expose_VII_1_3_1_through_1_3_3.tex

% BEGIN INCLUDED SGA7I_English_current_buildable_source_20260823/components/68_expose_VII_1_3_4.tex
% SOURCE AUTHORITY:
% expose_VII_body.tex lines 288--303; scan indices 158--159; source folios 147--148.
% Hard stop before source line 305, Proposition 1.3.5.

\SGARefTarget{sga7.sga7i.evii.n1-3-4}{1.3.4}. Return to the case in which two groups $P$ and $G$ are given,
and choose the identity section $e_P$ in order to define
rigidifications of bitorsors over $P$ and birigidifications of bitorsors
over $P\times P$. There are natural functors
\begin{equation}\tag{\SGARefTarget{sga7.sga7i.evii.n1-3-4-1}{1.3.4.1}}
 \mathbf{EXT}(P;G)\xrightarrow{\ i\ }
 \mathrm{BITORSRIG}(P,G)\xrightarrow{\ \delta\ }
 \mathbf{BITORSBIRIG}(P,P;G),
\end{equation}
where the three expressions denote, respectively, the category of
extensions of the group $P$ by $G$, the category of bitorsors under
$(G_P,G_P)$ rigidified at $e_P$, and the category of bitorsors under
$(G_{P\times P},G_{P\times P})$ birigidified at $(e_P,e_P)$. The functor
$i$ associates with an extension $E$ of $P$ by $G$ the corresponding
bitorsor over $P$ (\SGARefLink{sga7.sga7i.evii.n1-1-6}{1.1.6}), rigidified by the identity section of $E$.
The functor $\delta$ is defined by
\begin{equation}\tag{\SGARefTarget{sga7.sga7i.evii.n1-3-4-2}{1.3.4.2}}
 \delta(E)=\pi^*(E)\,\pr_1^*(E)^{-1}\,\pr_2^*(E)^{-1},
\end{equation}
where $\pi:P\times P\longrightarrow P$ is the composition law of $P$,
and $\pr_1,\pr_2$ are the two projections. There is a
\underline{canonical birigidification of $\delta(E)$}. Indeed, the
restriction of $\delta(E)$ to $P\times e_P$ is canonically isomorphic
to
\[
 \id_P^*(E)\,\id_P^*(E)^{-1}\,e^*(E)^{-1},
\]
where $e=\pr_2\operatorname{inj}_1$ is the constant morphism
$P\longrightarrow P$ with value $e_P$. The given rigidification of $E$
provides an isomorphism from $e^*(E)$ to the trivial bitorsor, and hence
an isomorphism of the restriction of $\delta(E)$ to $P\times e_P$ with
the trivial bitorsor. Symmetrically one defines a partial
rigidification along $e_P\times P$, and the two partial
rigidifications are immediately seen to be compatible.

%%%%%%%%%%  scan idx 159  |  folio 148  %%%%%%%%%%

\SGARefTarget{sga7.sga7i.evii.n1-3-4-3}{1.3.4.3}. Observe that the composite $\delta i$ of the functors in
(\SGARefLink{sga7.sga7i.evii.n1-3-4-1}{1.3.4.1}) is canonically isomorphic to the trivial functor, that is, to
the constant functor whose value is the trivial birigidified bitorsor
over $P\times P$. Indeed, for every extension $E$ of $P$ by $G$, the
isomorphism $\varphi$ of (\SGARefLink{sga7.sga7i.evii.n1-1-3-2}{1.1.3.2}) plainly defines a trivialization of
$\delta i(E)$ compatible with its canonical birigidification.


% END INCLUDED SGA7I_English_current_buildable_source_20260823/components/68_expose_VII_1_3_4.tex

% BEGIN INCLUDED SGA7I_English_current_buildable_source_20260823/components/69_expose_VII_proposition_1_3_5.tex
% SOURCE AUTHORITY:
% expose_VII_body.tex lines 305--358; scan indices 159--162; source folios 148--151.
% Source-backed corrections: the proposition's opening Q is G throughout its statement;
% E in part (b) is a bitorsor, as required by delta(E) and the ensuing sentence;
% the final associativity specialization follows the three partial products stated in (b).
% Hard stop before source line 360, Corollary 1.3.6.

\underline{Proposition \SGARefTarget{sga7.sga7i.evii.n1-3-5}{1.3.5}. Let $P$ and $G$ be two groups of
$\underline T$.}

a) \underline{Suppose that every morphism of sheaves of sets
$P\times P\longrightarrow G$ which is trivial on $P\times e_P$ and on
$e_P\times P$ is trivial. Then the functor $i$ of (\SGARefLink{sga7.sga7i.evii.n1-3-4-1}{1.3.4.1}) is fully
faithful. In particular, on a bitorsor $E$ under $(G_P,G_P)$ rigidified
relative to $e_P$, that is, equipped with a central section $e_E$ of
$E_1=e_P^*(E)$, there is at most one extension structure of $P$ by $G$
compatible with its bitorsor structure and having $e_E$ as identity
section.}

b) \underline{Suppose in addition that every morphism of sheaves of
sets $P\times P\times P\longrightarrow G$ which is trivial on the
partial products $P\times P\times e_P$, $P\times e_P\times P$, and
$e_P\times P\times P$ is trivial. Let $E$ be a bitorsor under
$(G_P,G_P)$. In order that $E$ admit an extension structure of $P$ by
$G$ compatible with its bitorsor structure, it is necessary and
sufficient that the bitorsor under
$(G_{P\times P},G_{P\times P})$ underlying}
\[
 \delta(E)=\pi^*(E)\,\pr_1^*(E)^{-1}\,\pr_2^*(E)^{-1}
\]
\underline{be trivial. When this is so, for every central section
$e_E$ of $E_1=e_P^*(E)$, that is, every rigidification of $E$ relative
to $e_P$, there is a unique extension structure on $E$ compatible with
its bitorsor structure and having $e_E$ as identity section. In other
words, $E$, regarded as a bitorsor rigidified relative to $e_P$, lies
in the essential image of $i$. Moreover, the birigidified bitorsor
$\delta(E)$ has one and only one trivialization.}

%%%%%%%%%%  scan idx 160  |  folio 149  %%%%%%%%%%

Thus, in terms of diagram (\SGARefLink{sga7.sga7i.evii.n1-3-4-1}{1.3.4.1}), the essential image of the fully
faithful functor $i$ is ``the kernel'' of $\delta$. More picturesquely,
diagram (\SGARefLink{sga7.sga7i.evii.n1-3-4-1}{1.3.4.1}) is an ``exact sequence of pointed categories,'' and
even of categories endowed with pseudo-group structures (in the
terminology introduced in \SGARefLink{sga7.sga7i.evii.n2-5-1}{2.5.1} below).

Let us prove \SGARefLink{sga7.sga7i.evii.n1-3-5}{1.3.5}. For a), let $E,E'$ be two extensions of $P$ by $G$,
and let $f:E\longrightarrow E'$ be a morphism of the underlying
rigidified bitorsors. We show that $f$ is a morphism of extensions, or,
what plainly amounts to the same thing, a homomorphism of groups. We
must prove $f(xy)=f(x)f(y)$. If $j':E'\longrightarrow P$ denotes the
structural morphism of the torsor $E'$ over $P$, then
$j'f(xy)=j'(f(x)f(y))$: the two sides are respectively $j(xy)$ and
$j(x)j(y)$, which are equal because $j$ is multiplicative. Hence
\begin{equation*}
 f(xy)=u(j(x),j(y))f(x)f(y),
\end{equation*}
where
\begin{equation*}
 u:P\times P\longrightarrow G
\end{equation*}
is a uniquely determined morphism. Setting first $x=1$ and then $y=1$
in the preceding relation, and using $f(1)=1$, gives
\begin{equation*}
 u(p,1)=u(1,p)=1.
\end{equation*}
Thus $u$ is trivial on $P\times e_P$ and on $e_P\times P$; by
hypothesis it is trivial, and therefore $f(xy)=f(x)f(y)$, as required.

For b), necessity of the stated condition for $E$ to admit an
extension structure compatible with its bitorsor structure is immediate
from \SGARefLink{sga7.sga7i.evii.n1-3-4-3}{1.3.4.3}. It remains to prove sufficiency for the first assertion.
Pulling a trivialization $\varphi$ of $\delta(E)$ back along the
identity section of $P\times P$ gives a trivialization of
$E_1=e_P^*(E)$ (compare 1.5.1), that is, a rigidification of $E$
relative to $e_P$. Fix this rigidification. A given trivialization
$\varphi$ of $\delta(E)$ need not be compatible with the corresponding
birigidification. Relative to the natural partial rigidifications of
$\delta(E)$, its restrictions to $P\times e_P$ and $e_P\times P$ are
described by central morphisms
\begin{equation*}
 f:P\longrightarrow G,qquad g:P\longrightarrow G,
\end{equation*}
whose values at the identity section of $P$ coincide:
$f(1)=g(1)=a\in\Gamma G$. The central morphism
\begin{equation*}
 h:P\times P\longrightarrow G,qquad h(p,q)=f(p)g(q)a^{-1},
\end{equation*}
has restrictions $f$ and $g$ to $P\times e_P$ and $e_P\times P$.
Correcting $\varphi$ by $h$, that is, replacing it by
$\varphi(p,q)h(p,q)^{-1}$, therefore gives a trivialization of
$\delta(E)$ compatible with its birigidification. Such a trivialization
is unique: two of them differ by a central morphism
$P\times P\longrightarrow G$ trivial on $P\times e_P$ and
$e_P\times P$, hence trivial by hypothesis. This proves the last
assertion of \SGARefLink{sga7.sga7i.evii.n1-3-5}{1.3.5} b).

It remains to prove that $\varphi$, viewed as an isomorphism of the form
(\SGARefLink{sga7.sga7i.evii.n1-1-3-2}{1.1.3.2}), satisfies the associativity condition (\SGARefLink{sga7.sga7i.evii.n1-1-4-1}{1.1.4.1}); by our
choice of $\varphi$, $e_E$ is already a two-sided identity for the
composition law defined by $\varphi$. The two composites in (\SGARefLink{sga7.sga7i.evii.n1-1-4-1}{1.1.4.1})
are isomorphisms of bitorsors under
$(G_{P\times P\times P},G_{P\times P\times P})$, so they differ by a
central morphism
\begin{equation*}
 u:P\times P\times P\longrightarrow G
\end{equation*}
through the formula
\begin{equation*}
 (xy)z=u(jx,jy,jz)\,x(yz).
\end{equation*}

%%%%%%%%%%  scan idx 162  |  folio 151  %%%%%%%%%%

Setting successively $z=1$, $x=1$, and $y=1$, and using the fact that
$1$ is a two-sided identity for the composition law, shows that $u$ is
trivial on the three partial products $P\times P\times e_P$,
$e_P\times P\times P$, and $P\times e_P\times P$. The hypothesis then
implies that $u$ is trivial, so $(xy)z=x(yz)$. This completes the proof.


% END INCLUDED SGA7I_English_current_buildable_source_20260823/components/69_expose_VII_proposition_1_3_5.tex

% BEGIN INCLUDED SGA7I_English_current_buildable_source_20260823/components/70_expose_VII_1_3_6_through_1_3_7.tex
% SOURCE AUTHORITY:
% expose_VII_body.tex lines 360--398; scan indices 162--164; source folios 151--153.
% Source-backed correction: the action in 1.3.6 is defined by the extension E, not G.
% Hard stop before source line 400, paragraph 1.3.8.

\underline{Corollary \SGARefTarget{sga7.sga7i.evii.n1-3-6}{1.3.6}. Suppose that every morphism of sheaves of
sets $P\times G\longrightarrow G$ which is trivial on $P\times e_G$
and on $e_P\times G$ is trivial. Then, if $G$ is commutative, $G$ is
central in every extension of $P$ by $G$. If $P$ is also commutative
and the hypothesis of \SGARefLink{sga7.sga7i.evii.n1-3-5}{1.3.5} a) holds, then every extension $E$ of $P$
by $G$ is commutative.}

For the first assertion, we must prove that the action of $P$ on $G$
defined by the extension $E$ is trivial. This action is expressed by a
morphism $u:P\times G\longrightarrow G$. The restrictions to
$P\times e_G$ and $e_P\times G$ of the corresponding morphism
\[
 (p,g)\longmapsto u(p,g)g^{-1}
\]
are plainly trivial. Hence this last morphism is trivial, so
$u(p,g)=g$, as required. Now, if $E$ is an extension of $P$ by $G$, its
opposite group $E^{\mathrm o}$ is an extension of $P^{\mathrm o}$ by
$G^{\mathrm o}$, and the bitorsor structure underlying
$i(E^{\mathrm o})$ is obtained from $i(E)$ by interchanging the left
and right actions. When $E$ is a central extension, one therefore has
$i(E)=i(E^{\mathrm o})$; if, in addition, $P$ is commutative, so that
$P=P^{\mathrm o}$, the uniqueness assertion \SGARefLink{sga7.sga7i.evii.n1-3-5}{1.3.5} a) gives
$E=E^{\mathrm o}$. Thus $E$ is commutative, proving the corollary.

%%%%%%%%%%  scan idx 163  |  folio 152  %%%%%%%%%%

\SGARefTarget{sga7.sga7i.evii.n1-3-7}{1.3.7}. One interesting special case in which all the hypotheses of
\SGARefLink{sga7.sga7i.evii.n1-3-5}{1.3.5} and \SGARefLink{sga7.sga7i.evii.n1-3-6}{1.3.6} hold is
\begin{equation}\tag{\SGARefTarget{sga7.sga7i.evii.n1-3-7-1}{1.3.7.1}}
 \underline{\operatorname{Hom}}_{\mathrm{pt}}(P,G)=e
 \quad\text{(the final sheaf)},
\end{equation}
where the left-hand side denotes the sheaf of morphisms of pointed
sheaves of sets from $P$ to $G$. Independently of the marked sections,
this condition may also be written in the form, analogous to (\SGARefLink{sga7.sga7i.evii.n1-3-1-2}{1.3.1.2}),
\begin{equation}\tag{\SGARefTarget{sga7.sga7i.evii.n1-3-7-2}{1.3.7.2}}
 G\simeq\underline{\operatorname{Hom}}_{\mathrm{sets}}(P,G),
\end{equation}
where the second member denotes the sheaf of morphisms between the
underlying sheaves of sets of $P$ and $G$. These conditions also mean
that, for every object $Q$ of $\underline T$, the map
$u\longmapsto u\circ\pr_2$
\begin{equation}\tag{\SGARefTarget{sga7.sga7i.evii.n1-3-7-3}{1.3.7.3}}
 \operatorname{Hom}_{\mathrm{sets}}(Q,G)
 \xrightarrow{\ \sim\ }
 \operatorname{Hom}_{\mathrm{sets}}(P\times Q,G)
\end{equation}
is bijective. This property of $P$ (with $G$ fixed) is stable under
Cartesian products; in particular it implies the same property for
$P\times P$, $P\times P\times P$, and so forth. It follows immediately
that it implies the hypotheses made in \SGARefLink{sga7.sga7i.evii.n1-3-5}{1.3.5} and \SGARefLink{sga7.sga7i.evii.n1-3-6}{1.3.6}. It also implies
that all the categories occurring in (\SGARefLink{sga7.sga7i.evii.n1-3-4-1}{1.3.4.1}) are \underline{rigid},
that is, every automorphism of an object of any one of them is the
identity.

Moreover, since condition (\SGARefLink{sga7.sga7i.evii.n1-3-7-1}{1.3.7.1}) on $P,G$ is stable under every base
change $S\longrightarrow e$, descent together with the uniqueness
result \SGARefLink{sga7.sga7i.evii.n1-3-5}{1.3.5} a) shows that a rigidified bitorsor $E$ over $P$ comes from
an extension of $P$ by $G$ as soon as it does so \underline{locally over
$e$}. By \SGARefLink{sga7.sga7i.evii.n1-3-5}{1.3.5} b), this may also be expressed by the vanishing of the
section
\begin{equation*}
 \underline{\delta}(E)\in\Gamma R^1h_*(G_{P\times P})
\end{equation*}
defined by $\delta(E)$, where
\begin{equation*}
 h:P\times P\longrightarrow e
\end{equation*}
is the structural morphism.


% END INCLUDED SGA7I_English_current_buildable_source_20260823/components/70_expose_VII_1_3_6_through_1_3_7.tex

% BEGIN INCLUDED SGA7I_English_current_buildable_source_20260823/components/71_expose_VII_1_3_8.tex
% SOURCE AUTHORITY:
% expose_VII_body.tex lines 400--418; scan indices 164--165; source folios 153--154.
% Source-backed correction: final references 3.5 and 3.6 are 1.3.5 and 1.3.6.
% Hard stop before source line 421, section 1.4.

\SGARefTarget{sga7.sga7i.evii.n1-3-8}{1.3.8}. Let $p:P\longrightarrow S$ be a morphism of schemes which is
\underline{coherent} (that is, quasi-compact and quasi-separated) and
\underline{flat}, and suppose that
\begin{equation}\tag{\SGARefTarget{sga7.sga7i.evii.n1-3-8-1}{1.3.8.1}}
 p_*(\mathcal O_P)\xleftarrow{\ \sim\ }\mathcal O_S.
\end{equation}
Since the formation of $p_*(\mathcal O_P)$ commutes with every flat base
change (by coherence of $p$), the stated conditions on the relative
scheme $P/S$ are stable under Cartesian products and hence hold, in
particular, for the Cartesian powers $P\times P$,
$P\times P\times P$, and so forth. They are also stable under every flat
base change.

Now let $G$ be an \underline{affine} scheme over $S$. Relation
(\SGARefLink{sga7.sga7i.evii.n1-3-7-2}{1.3.7.2}) follows at once from (\SGARefLink{sga7.sga7i.evii.n1-3-8-1}{1.3.8.1}), since
\[
 \operatorname{Hom}_S(P,G)
 =\operatorname{Hom}_{\mathrm{Alg}}(\underline A,p_*(\mathcal O_P))
 \simeq\operatorname{Hom}_{\mathrm{Alg}}(\underline A,\mathcal O_S),
\]
where $\underline A$ is the quasi-coherent algebra over $S$ such that
$G=\operatorname{Spec}(\underline A)$. By the preceding observation,
the same conclusion applies to $P\times P$, $P\times P\times P$, and
so on. Likewise, if $G$ is flat over $S$, the same conclusion applies
to the relative scheme $P_G$ over $G$ with respect to the relative
scheme $G_G$ over $G$.

When $P$ and $G$ are group schemes over $S$, and one works with a
topology on $(\mathrm{Sch})_{/S}$ coarser than the canonical topology,
so that $P$ and $G$ may be identified with sheaves on this site, every
morphism from $P$, $P\times P$, $P\times P\times P$, and so on to $G$
which is compatible with the marked sections is trivial. Similarly, if
$G$ is flat, every $S$-morphism $P\times G\longrightarrow G$ compatible
with the marked sections is trivial, as one sees by interpreting it as
a $G$-morphism from $P_G$ to $G_G$. Thus the hypotheses of \SGARefLink{sga7.sga7i.evii.n1-3-5}{1.3.5} hold,
as do those of \SGARefLink{sga7.sga7i.evii.n1-3-6}{1.3.6} when $G$ is flat. Moreover, the three categories
occurring in (\SGARefLink{sga7.sga7i.evii.n1-3-4-1}{1.3.4.1}) are rigid.

%%%%%%%%%%  scan idx 165  |  folio 154  %%%%%%%%%%

\SGARefTarget{sga7.sga7i.evii.n1-3-8-2}{1.3.8.2}. The most interesting case in which (\SGARefLink{sga7.sga7i.evii.n1-3-8-1}{1.3.8.1}) holds, and
continues to hold after every extension of the base, is that in which
$P$ is an \underline{abelian scheme} over $S$; then the conditions of
\SGARefLink{sga7.sga7i.evii.n1-3-7}{1.3.7} themselves hold. In this case the conclusions of \SGARefLink{sga7.sga7i.evii.n1-3-5}{1.3.5} and \SGARefLink{sga7.sga7i.evii.n1-3-6}{1.3.6}
are therefore valid for every affine group $G$ over $S$. This is the
case originally considered by J.~P. Serre, and the preceding
developments were plainly inspired by it. An interesting case in which
(\SGARefLink{sga7.sga7i.evii.n1-3-8-1}{1.3.8.1}) holds without holding universally is that in which $P$ is the
\underline{N\'eron model} of an abelian scheme over the fraction field
of a discrete valuation ring with spectrum $S$.


% END INCLUDED SGA7I_English_current_buildable_source_20260823/components/71_expose_VII_1_3_8.tex

% BEGIN INCLUDED SGA7I_English_current_buildable_source_20260823/components/72_expose_VII_section_1_4.tex
% SOURCE AUTHORITY:
% expose_VII_body.tex lines 421--460; scan indices 165--167; source folios 154--156.
% Hard stop before source line 465, section 2.

\SGARefTarget{sga7.sga7i.evii.n1-4}{1.4}. \underline{Extensions of a free commutative group}

Let $I$ be an object of $\underline T$, and put
\begin{equation}\tag{\SGARefTarget{sga7.sga7i.evii.n1-4-1}{1.4.1}}
 P=\mathbb Z[I],
\end{equation}
the free $\mathbb Z$-module generated by $I$ (SGA 4, IV, 11). With
every extension $E$ of $P$ by a given $\mathbb Z$-module $G$ (from now
on only commutative group extensions will occur), associate the
corresponding torsor under $G_P$, and then restrict it to $I$ along the
canonical morphism
\[
 I\longrightarrow P=\mathbb Z[I].
\]

%%%%%%%%%%  scan idx 166  |  folio 155  %%%%%%%%%%

This gives a canonical functor
\begin{equation}\tag{\SGARefTarget{sga7.sga7i.evii.n1-4-2}{1.4.2}}
 \operatorname{EXT}(P,G)\longrightarrow
 \operatorname{TORS}(I,G_I).
\end{equation}
I claim that this functor is an \underline{equivalence of categories}.
Equivalently, for every $G$, the maps
\begin{equation}\tag{\SGARefTarget{sga7.sga7i.evii.n1-4-3}{1.4.3}}
 \operatorname{Hom}_{\mathbb Z}(P,G)\longrightarrow
 H^0(I,G_I)=G(I)
\end{equation}
and
\begin{equation}\tag{\SGARefTarget{sga7.sga7i.evii.n1-4-4}{1.4.4}}
 \operatorname{Ext}^1_{\mathbb Z}(P,G)\longrightarrow H^1(I,G_I)
\end{equation}
are isomorphisms. More generally, if $A$ is a ring (not necessarily
commutative) of $\underline T$ and $P=A[I]$ is the free $A$-module
generated by $I$, there are canonical isomorphisms, functorial in the
variable $A$-module $G$,
\begin{equation}\tag{\SGARefTarget{sga7.sga7i.evii.n1-4-5}{1.4.5}}
 \operatorname{Ext}^i_A(A[I],G)\xrightarrow{\ \sim\ }H^i(I,G_I).
\end{equation}
For $A=\mathbb Z$ and $i=0,1$, these coincide with (\SGARefLink{sga7.sga7i.evii.n1-4-3}{1.4.3}) and
(\SGARefLink{sga7.sga7i.evii.n1-4-4}{1.4.4}). To define (\SGARefLink{sga7.sga7i.evii.n1-4-5}{1.4.5}), recall that its left-hand side is the
derived functor of
\[
 G\longmapsto\operatorname{Hom}_A(A[I],G),
\]
which, by the definition of $A[I]$, is isomorphic to
$G\longmapsto G(I)=H^0(I,G_I)$. Since $G\longmapsto G_I$ carries
injective modules to injective modules of $\underline T_{/I}$ (SGA 4,
V), the assertion follows. We leave to the reader the verification of
the stated compatibility with (\SGARefLink{sga7.sga7i.evii.n1-4-3}{1.4.3}) and (\SGARefLink{sga7.sga7i.evii.n1-4-4}{1.4.4}).

Finally, the functor (\SGARefLink{sga7.sga7i.evii.n1-4-2}{1.4.2}) is compatible with inverse images under a
morphism of topoi
\[
 f:\underline T'\longrightarrow\underline T.
\]


% END INCLUDED SGA7I_English_current_buildable_source_20260823/components/72_expose_VII_section_1_4.tex

% BEGIN INCLUDED SGA7I_English_current_buildable_source_20260823/components/73_expose_VII_section_2_0_phi.tex
% SOURCE AUTHORITY:
% expose_VII_body.tex lines 465--540; scan indices 167--169; source folios 156--158.
% Source-backed correction: the printed prose interchanges the P- and Q-direction extension names;
% the displayed phi laws unambiguously multiply the P-coordinate, hence define P_Q by G_Q.
% The missing prime in the upper-left label of (2.0.5) is restored from its nodes and associativity.
% Hard stop before the psi-law continuation on source line 540.

2. \underline{The notion of a biextension of abelian sheaves}

\SGARefTarget{sga7.sga7i.evii.n2-0}{2.0}. Throughout this section, $P,Q,G$ denote commutative groups of the
topos $\underline T$. We shall frequently use the Cartesian diagram
\begin{equation}\tag{\SGARefTarget{sga7.sga7i.evii.n2-0-1}{2.0.1}}
\begin{tikzcd}[column sep=4em,row sep=2.6em]
 P \arrow[d] & P\times Q \arrow[l] \arrow[d] \\
 e & Q \arrow[l]
\end{tikzcd}
\end{equation}
and the groups $G_P$ (respectively $G_Q$, $G_{P\times Q}$) over $P$
(respectively $Q$, $P\times Q$), obtained by inverse image from the
group $G$ over $e$. We identify $G_{P\times Q}$ either with the inverse
image $(G_P)_{P\times Q}$ under $\pr_1$, or with the inverse image
$(G_Q)_{P\times Q}$ under $\pr_2$. In this notation the group
structures of $P,Q,P\times Q$ play no role; indeed, we shall almost
never use the product-group structure on $P\times Q$. We shall instead
use on $P\times Q$ the structure of a group over $P$ obtained by inverse
image from $Q$ over $e$,
\begin{equation}\tag{\SGARefTarget{sga7.sga7i.evii.n2-0-2}{2.0.2}}
 P\times Q=Q_P,
\end{equation}
and, similarly, the structure of a group over $Q$ obtained by inverse
image from $P$ over $e$,
\begin{equation}\tag{\SGARefTarget{sga7.sga7i.evii.n2-0-3}{2.0.3}}
 P\times Q=P_Q.
\end{equation}
Thus over $P$ we have the two groups $G_P,Q_P$, and over $Q$ the two
groups $G_Q,P_Q$.

If $E$ is a torsor over $P\times Q$ under $G_{P\times Q}$, it may be
viewed either as a torsor over $Q_P$ under
$G_{P\times Q}=(G_P)_{Q_P}$, or as a torsor over $P_Q$ under
$G_{P\times Q}=(G_Q)_{P_Q}$. It therefore makes sense to equip $E$ with
a commutative extension structure in either variable. Applied in the
induced topoi $\underline T_{/P}$ and $\underline T_{/Q}$, \SGARefLink{sga7.sga7i.evii.n1-1-6}{1.1.6} gives
a convenient explicit description, using the canonical isomorphisms
\[
 \underline T_{/P\times Q}\simeq
 (\underline T_{/P})_{/P\times Q},\qquad
 \underline T_{/P\times Q}\simeq
 (\underline T_{/Q})_{/P\times Q}.
\]

An extension structure in the $P$-direction---that is, of $P_Q$ by
$G_Q$---is equivalent to giving, for every object $S$ of $\underline T$
and elements $p,p'\in P(S)$, $q\in Q(S)$, an isomorphism of torsors
under $G_S$
\begin{equation}\tag{\SGARefTarget{sga7.sga7i.evii.n2-0-4}{2.0.4}}
 \varphi_{p,p';q}:E_{p,q}E_{p',q}
 \xrightarrow{\ \sim\ }E_{pp',q}.
\end{equation}
As usual (cf. \SGARefLink{sga7.sga7i.evii.n1-1-2-1}{1.1.2.1}), for $(p,q)\in(P\times Q)(S)$, $E_{p,q}$ denotes
the torsor under $G_S$ obtained from $E$ by inverse image along
$(p,q):S\longrightarrow P\times Q$. The data (\SGARefLink{sga7.sga7i.evii.n2-0-4}{2.0.4}) must be
\underline{functorial in $S$}, associative, and commutative. The
associativity condition is the commutativity of
\begin{equation}\tag{\SGARefTarget{sga7.sga7i.evii.n2-0-5}{2.0.5}}
\begin{tikzcd}[column sep=5.1em,row sep=3.4em]
 & E_{pp',q}E_{p'',q}
   \arrow[rd,"{\varphi_{pp',p'';q}}"] & \\
 E_{p,q}E_{p',q}E_{p'',q}
   \arrow[ru,"{\varphi_{p,p';q}\wedge\mathrm{id}}"]
   \arrow[rd,"{\mathrm{id}\wedge\varphi_{p',p'';q}}"']
 && E_{pp'p'',q} \\
 & E_{p,q}E_{p'p'',q}
   \arrow[ru,"{\varphi_{p,p'p'';q}}"'] &
\end{tikzcd}
\end{equation}
for $p,p',p''\in P(S)$ and $q\in Q(S)$. Commutativity is expressed by
\begin{equation}\tag{\SGARefTarget{sga7.sga7i.evii.n2-0-6}{2.0.6}}
\begin{tikzcd}[column sep=5em,row sep=3.1em]
 E_{p,q}E_{p',q} \arrow[r,"{\varphi_{p,p';q}}"]
   \arrow[d,"\mathrm{sym}"']
 & E_{pp',q} \arrow[d,"\mathrm{id}"] \\
 E_{p',q}E_{p,q} \arrow[r,"{\varphi_{p',p;q}}"']
 & E_{p'p,q}.
\end{tikzcd}
\end{equation}
Here $p,p'\in P(S)$ and $q\in Q(S)$.


% END INCLUDED SGA7I_English_current_buildable_source_20260823/components/73_expose_VII_section_2_0_phi.tex

% BEGIN INCLUDED SGA7I_English_current_buildable_source_20260823/components/74_expose_VII_psi_and_definition_2_1.tex
% SOURCE AUTHORITY:
% expose_VII_body.tex lines 540--658; scan indices 169--171; source folios 158--160.
% Source-backed corrections: (2.0.7) targets E_{p,qq'}, not the impossible E_{pp',q};
% the second structure is Q_P by G_P, not the printed reversal G_Q by P_Q.
% Hard stop before source line 660, section 2.2.

Symmetrically, an extension structure in the $Q$-direction---that is,
of $Q_P$ by $G_P$---is described by isomorphisms of torsors under $G_S$
\begin{equation}\tag{\SGARefTarget{sga7.sga7i.evii.n2-0-7}{2.0.7}}
 \psi_{p;q,q'}:E_{p,q}E_{p,q'}
 \xrightarrow{\ \sim\ }E_{p,qq'},
\end{equation}
for $p\in P(S)$ and $q,q'\in Q(S)$. These are functorial in $S$ and
satisfy the symmetric associativity and commutativity conditions
\begin{equation}\tag{\SGARefTarget{sga7.sga7i.evii.n2-0-8}{2.0.8}}
\begin{tikzcd}[column sep=5.1em,row sep=3.4em]
 & E_{p,qq'}E_{p,q''}
   \arrow[rd,"{\psi_{p;qq',q''}}"] & \\
 E_{p,q}E_{p,q'}E_{p,q''}
   \arrow[ru,"{\psi_{p;q,q'}\wedge\mathrm{id}}"]
   \arrow[rd,"{\mathrm{id}\wedge\psi_{p;q',q''}}"']
 && E_{p,qq'q''} \\
 & E_{p,q}E_{p,q'q''}
   \arrow[ru,"{\psi_{p;q,q'q''}}"'] &
\end{tikzcd}
\end{equation}
for $p\in P(S)$ and $q,q',q''\in Q(S)$, and
\begin{equation}\tag{\SGARefTarget{sga7.sga7i.evii.n2-0-9}{2.0.9}}
\begin{tikzcd}[column sep=5em,row sep=3.1em]
 E_{p,q}E_{p,q'} \arrow[r,"{\psi_{p;q,q'}}"]
   \arrow[d,"\mathrm{sym}"']
 & E_{p,qq'} \arrow[d,"\mathrm{id}"] \\
 E_{p,q'}E_{p,q} \arrow[r,"{\psi_{p;q',q}}"']
 & E_{p,q'q}.
\end{tikzcd}
\end{equation}
Here $p\in P(S)$ and $q,q'\in Q(S)$.

%%%%%%%%%%  scan idx 170  |  folio 159  %%%%%%%%%%

\underline{Definition \SGARefTarget{sga7.sga7i.evii.n2-1}{2.1}. Let $P,Q,G$ be abelian groups of
$\underline T$, and let $E$ be a torsor under $G_{P\times Q}$ equipped
with an extension structure in the $P$-direction, defined by the
isomorphisms (\SGARefLink{sga7.sga7i.evii.n2-0-4}{2.0.4}), and an extension structure in the $Q$-direction,
defined by the isomorphisms (\SGARefLink{sga7.sga7i.evii.n2-0-7}{2.0.7}), each functorial in $S$ and subject
to the associativity and commutativity conditions above. The two
structures are called compatible if, for every object $S$ of
$\underline T$ and $p,p'\in P(S)$, $q,q'\in Q(S)$, the following
diagram commutes:}
\begin{equation}\tag{\SGARefTarget{sga7.sga7i.evii.n2-1-1}{2.1.1}}
\begin{tikzcd}[column sep=2.7em,row sep=3.2em,font=\scriptsize]
 && E_{p,qq'}E_{p',qq'}
    \arrow[rrdd,"{\varphi_{p,p';qq'}}"] && \\
 E_{p,q}E_{p,q'}E_{p',q}E_{p',q'}
    \arrow[rru,"{\psi_{p;q,q'}\wedge\psi_{p';q,q'}}"]
    \arrow[d,"{}"'] &&&& \\
 E_{p,q}E_{p',q}E_{p,q'}E_{p',q'}
    \arrow[rrd,"{\varphi_{p,p';q}\wedge\varphi_{p,p';q'}}"']
 &&&& E_{pp',qq'} \\
 && E_{pp',q}E_{pp',q'}
    \arrow[rru,"{\psi_{pp';q,q'}}"'] &&
\end{tikzcd}
\end{equation}
\underline{where the vertical arrow is induced by the symmetry
isomorphism}
\[
 E_{p,q'}E_{p',q}\xrightarrow{\ \sim\ }E_{p',q}E_{p,q'}.
\]
\underline{A \emph{biextension} of $(P,Q)$ by $G$ is a torsor $E$
under $G_{P\times Q}$ equipped with compatible extension structures in
the $P$- and $Q$-directions in the preceding sense.}

%%%%%%%%%%  scan idx 171  |  folio 160  %%%%%%%%%%

\SGARefLink{sga7.sga7i.evii.n2-1-1}{2.1.1}. Equivalently, a biextension structure on the
$G_{P\times Q}$-torsor $E$ is given by two partially defined composition
laws $\varphi(x,y)$ and $\psi(x,y)$ on $E$. The first is defined when
$\pr_2j(x)=\pr_2j(y)$, and the second when
$\pr_1j(x)=\pr_1j(y)$, where $j:E\longrightarrow P\times Q$ is the
structural morphism of the torsor. Besides associativity relations of
the type
\begin{equation}\tag{\SGARefTarget{sga7.sga7i.evii.n2-1-1-1}{2.1.1.1}}
 \varphi(gx,y)=\varphi(x,gy)=g\varphi(x,y)
 \qquad(x,y\in E(S),\ g\in G(S))
\end{equation}
and the analogous relations for $\psi$, the axioms say that these
partial composition laws are associative and commutative and satisfy
\begin{equation}\tag{\SGARefTarget{sga7.sga7i.evii.n2-1-1-2}{2.1.1.2}}
 \varphi(\psi(x,y),\psi(z,t))
 =\psi(\varphi(x,z),\varphi(y,t)),
\end{equation}
where $x,y,z,t$ lie respectively in
$E_{p,q}(S),E_{p,q'}(S),E_{p',q}(S),E_{p',q'}(S)$.


% END INCLUDED SGA7I_English_current_buildable_source_20260823/components/74_expose_VII_psi_and_definition_2_1.tex

% BEGIN INCLUDED SGA7I_English_current_buildable_source_20260823/components/75_expose_VII_section_2_2.tex
% SOURCE AUTHORITY:
% expose_VII_body.tex lines 660--683; scan indices 171--173; source folios 160--162.
% Source-backed corrections: e_{E/P} trivializes over P x 1_Q and e_{E/Q} over 1_P x Q;
% the final restricted torsor is an inverse image in E, not in G.
% Hard stop before source line 688, section 2.3.

\SGARefTarget{sga7.sga7i.evii.n2-2}{2.2}. \underline{Identity sections}. Consider first the identity section
of $E$ regarded as an extension of $Q_P$ by $G_P$,
\begin{equation}\tag{\SGARefTarget{sga7.sga7i.evii.n2-2-1}{2.2.1}}
 e_{E/P}\in\Gamma(E/P)=\operatorname{Hom}_P(P,E),
\end{equation}
and, symmetrically, the identity section of $E$ regarded as an extension
of $P_Q$ by $G_Q$,
\begin{equation}\tag{\SGARefTarget{sga7.sga7i.evii.n2-2-2}{2.2.2}}
 e_{E/Q}\in\Gamma(E/Q)=\operatorname{Hom}_Q(Q,E).
\end{equation}
In particular, \underline{the torsor $E$ over $P\times Q$ under
$G_{P\times Q}$ is canonically trivialized over $P\times1_Q$ by
$e_{E/P}$, and over $1_P\times Q$ by $e_{E/Q}$. These trivializations
coincide over the identity section $1_P\times1_Q$ of $P\times Q$}, and
hence define the same section
\begin{equation}\tag{\SGARefTarget{sga7.sga7i.evii.n2-2-3}{2.2.3}}
 e_E\in\Gamma(E),
\end{equation}
which, by abuse of language, will be called the \underline{identity
section} (or \underline{neutral section}) of $E$.

%%%%%%%%%%  scan idx 172  |  folio 161  %%%%%%%%%%

To prove the asserted equality, let $e,e'$ be the two sections of $E$
deduced respectively from $e_{E/P}$ and $e_{E/Q}$. They are
characterized as the identity sections in the inverse image
$E_{1_P,1_Q}$ of $E$ along the identity section of $P\times Q$, for the
two composition laws $\varphi$ and $\psi$, respectively. Applying
(\SGARefLink{sga7.sga7i.evii.n2-1-1-2}{2.1.1.2}) to the four sections $e,e',e',e$ of $E_{1,1}$ gives $e$ on
the left and $e'$ on the right; hence $e=e'$. It follows moreover that
\underline{the two composition laws $\varphi$ and $\psi$ on $E_{1,1}$
are identical}. Indeed, under transport of structure by the morphisms
$g\mapsto g.e$ and $g\mapsto g.e'$, respectively, both correspond to
the group law of $G$, and $e=e'$.

The extension of $Q_P$ by $G_P$ over the base $P$, restricted over the
identity section $1_P$ of $P$, gives an extension of $Q$ by $G$ on the
inverse image of $1_P\times Q$ in $E$. The homomorphism $e_{E/Q}$ is
immediately seen to be multiplicative for this extension structure, so
the extension is canonically split. Symmetrically, the extension of $P$
by $G$ obtained by restricting the extension of $P_Q$ by $G_Q$ over
$Q$ to the identity section $1_Q$ is canonically split by $e_{E/P}$.


% END INCLUDED SGA7I_English_current_buildable_source_20260823/components/75_expose_VII_section_2_2.tex

% BEGIN INCLUDED SGA7I_English_current_buildable_source_20260823/components/76_expose_VII_section_2_3.tex
% SOURCE AUTHORITY:
% expose_VII_body.tex lines 688--725; scan indices 173--174; source folios 162--163.
% The relative homomorphisms are retained exactly with the two bases made explicit.
% Hard stop before source line 728, section 2.4.

\SGARefTarget{sga7.sga7i.evii.n2-3}{2.3}. \underline{Morphisms of biextensions}. Let $E$ be a biextension
of $(P,Q)$ by $G$, and let $E'$ likewise be a biextension of
$(P',Q')$ by $G'$, where $P',Q',G'$ are again commutative groups of
$\underline T$. A \underline{morphism of biextensions} from $E$ to
$E'$ is a system $(u,v,w,f)$, where
\begin{equation}\tag{\SGARefTarget{sga7.sga7i.evii.n2-3-1}{2.3.1}}
 u:P\longrightarrow P',\qquad
 v:Q\longrightarrow Q',\qquad
 w:G\longrightarrow G'
\end{equation}
are morphisms of groups, and
\begin{equation}\tag{\SGARefTarget{sga7.sga7i.evii.n2-3-2}{2.3.2}}
 f:E\longrightarrow E'
\end{equation}
is a morphism of the sheaves of sets underlying $E$ and $E'$, such
that $f$ is simultaneously a homomorphism of group extensions over
$P$ and $P'$, respectively, associated with the relative group
homomorphisms
\[
 u\times v:Q_P\longrightarrow Q'_{P'},
 \qquad
 u\times w:G_P\longrightarrow G'_{P'},
\]
and a homomorphism of group extensions over $Q$ and $Q'$, respectively,
associated with the relative group homomorphisms
\[
 u\times v:P_Q\longrightarrow P'_{Q'},
 \qquad
 v\times w:G_Q\longrightarrow G'_{Q'}.
\]
Notice moreover that $u,v,w$ are determined by $f$ alone. Indeed,
passing to the quotient recovers $u\times v$ from $f$, and then
recovers $u$ and $v$ by restriction to $P\times1_Q$ and
$1_P\times Q$, respectively; restricting $f$ to $G.e_E$ recovers
$w$. Thus no ambiguity arises from identifying the morphism
$(u,v,w,f)$ with $f$.

After taking the usual precaution of adjoining their sources and
targets to morphisms (the canonical genuflection), one sees that
\underline{the biextensions of the topos $\underline T$ are the
objects of a category}, whose objects and morphisms are those defined
in \SGARefLink{sga7.sga7i.evii.n2-1}{2.1} and \SGARefLink{sga7.sga7i.evii.n2-3}{2.3} and whose composition of morphisms is defined in the
evident way. Moreover, $(P,Q,G,E)\mapsto(P,Q,G)$ gives a canonical
functor
\begin{equation}\tag{\SGARefTarget{sga7.sga7i.evii.n2-3-3}{2.3.3}}
 \operatorname{BIEXT}(\underline T)
 \longrightarrow
 \operatorname{GRAB}(\underline T)\times
 \operatorname{GRAB}(\underline T)\times
 \operatorname{GRAB}(\underline T),
\end{equation}
where $\operatorname{BIEXT}(\underline T)$ denotes the category of
biextensions over $\underline T$, and
$\operatorname{GRAB}(\underline T)$ the category of abelian groups
over $\underline T$. It is useful to decompose this canonical functor
into the two functors $E\mapsto(P,Q)$,
\begin{equation}\tag{\SGARefTarget{sga7.sga7i.evii.n2-3-4}{2.3.4}}
 \operatorname{BIEXT}(\underline T)
 \longrightarrow
 \operatorname{GRAB}(\underline T)\times
 \operatorname{GRAB}(\underline T),
\end{equation}
and $E\mapsto G$,
\begin{equation}\tag{\SGARefTarget{sga7.sga7i.evii.n2-3-5}{2.3.5}}
 \operatorname{BIEXT}(\underline T)
 \longrightarrow
 \operatorname{GRAB}(\underline T).
\end{equation}


% END INCLUDED SGA7I_English_current_buildable_source_20260823/components/76_expose_VII_section_2_3.tex

% BEGIN INCLUDED SGA7I_English_current_buildable_source_20260823/components/77_expose_VII_section_2_4.tex
% SOURCE AUTHORITY:
% expose_VII_body.tex lines 728--773; scan indices 174--176; source folios 163--165.
% Source-backed corrections: the covariant construction begins with a biextension of
% (P,Q), not the printed (P,G); (2.4.2) has BIEXT(P,Q;G x G'), not BIEXT(P,G;...).
% The scanned source confirms the stray <<CONF marker is not mathematical content.
% Hard stop before source line 776, section 2.5.

\SGARefTarget{sga7.sga7i.evii.n2-4}{2.4}. \underline{The category $\operatorname{BIEXT}(P,Q;G)$ for
variable $P,Q,G$}.

I claim that \underline{the canonical functor (\SGARefLink{sga7.sga7i.evii.n2-3-4}{2.3.4}) is fibred}, and
that \underline{the canonical functor (\SGARefLink{sga7.sga7i.evii.n2-3-5}{2.3.5}) is cofibred}. The first
assertion means, in particular, that if $E$ is a biextension of
$(P,Q)$ by $G$, and if
\[
 u:P'\longrightarrow P,\qquad v:Q'\longrightarrow Q
\]
are morphisms of abelian groups, then there is a biextension $E'$ of
$(P',Q')$ by $G$ and a morphism of biextensions
\[
 f:E'\longrightarrow E
\]
compatible with $u,v,\mathrm{id}_G$, having the usual universal
property for morphisms of biextensions into $E$ compatible with
$u,v,\mathrm{id}_G$.

The second assertion is expressed similarly. In particular, if $E$ is
a biextension of $(P,Q)$ by $G$, and if
\[
 w:G\longrightarrow G'
\]
is a morphism of abelian groups, then there is a biextension $E'$ of
$(P,Q)$ by $G'$ and a morphism of biextensions
\[
 g:E\longrightarrow E'
\]
compatible with $\mathrm{id}_P,\mathrm{id}_Q,w$, having the usual
universal property for morphisms of biextensions from $E$ compatible
with $\mathrm{id}_P,\mathrm{id}_Q,w$. We leave it to the reader to
spell out the two constructions whose existence has just been
asserted; they are obtained easily by paraphrasing the familiar
analogous constructions for extensions of ordinary groups.

Thus, for three abelian groups $P,Q,G$ of $\underline T$, write
\begin{equation}\tag{\SGARefTarget{sga7.sga7i.evii.n2-4-1}{2.4.1}}
 \operatorname{BIEXT}(P,Q;G)
\end{equation}
for the category of biextensions of $(P,Q)$ by $G$, defined as the
fibre category of (\SGARefLink{sga7.sga7i.evii.n2-3-3}{2.3.3}) over $(P,Q,G)$ (so its morphisms are the
morphisms of biextensions compatible with the identity morphisms of
$P,Q,G$). Then \underline{this category depends covariantly on $G$ and
contravariantly on $P,Q$}. Moreover, with these variances,
\underline{the category $\operatorname{BIEXT}(P,Q;G)$ depends
additively on each of the three arguments $P,Q,G$}. For example, for
the argument $G$, this means precisely:

\begin{enumerate}[label=\alph*)]
\item $\operatorname{BIEXT}(P,Q;0)$ is equivalent to the one-point
category;
\item for two abelian groups $G,G'$ of $\underline T$, the functor
\begin{equation}\tag{\SGARefTarget{sga7.sga7i.evii.n2-4-2}{2.4.2}}
 \operatorname{BIEXT}(P,Q;G\times G')
 \longrightarrow
 \operatorname{BIEXT}(P,Q;G)\times
 \operatorname{BIEXT}(P,Q;G')
\end{equation}
deduced from the two projections $G\times G'\longrightarrow G$ and
$G\times G'\longrightarrow G'$ is an
\underline{equivalence of categories}.
\end{enumerate}
In the language of [4], one also says that, for fixed $P,Q$, the
category of biextensions of $(P,Q)$ by variable groups $G$ is, via the
functor $E\mapsto G$ induced by (\SGARefLink{sga7.sga7i.evii.n2-3-5}{2.3.5}), an
\underline{additive cofibred category} over
$\operatorname{GRAB}(\underline T)$.

\underline{Remark \SGARefTarget{sga7.sga7i.evii.n2-4-3}{2.4.3}.} In the terminology of \emph{loc.\ cit.}, the
preceding cofibred category is not merely additive but even
\underline{left exact} [4, \SGARefLink{sga7.sga7i.evii.n2-2}{2.2}]. We shall not verify this property here,
since it will reappear below in \SGARefLink{sga7.sga7i.evii.n3-6}{3.6} as a consequence of more precise
facts. We only point out that the general theory of \emph{loc.\ cit.}
would then describe the structure of this cofibred category, up to
equivalence of fibred categories, in terms of a chain complex
$L_\bullet$ in $\operatorname{GRAB}(\underline T)$ depending only on
$P$ and $Q$. We shall also give below a direct construction of this
complex as
\[
 P\overset{\mathbb L}{\otimes}Q,
\]
the tensor product taken in a derived category.


% END INCLUDED SGA7I_English_current_buildable_source_20260823/components/77_expose_VII_section_2_4.tex

% BEGIN INCLUDED SGA7I_English_current_buildable_source_20260823/components/78_expose_VII_section_2_5.tex
% SOURCE AUTHORITY:
% expose_VII_body.tex lines 776--844; scan indices 176--179; source folios 165--168.
% The scan confirms all displayed Biext groups and the derived tensor notation.
% Hard stop before source line 846, section 2.6.

\SGARefTarget{sga7.sga7i.evii.n2-5}{2.5}. \underline{Structure of the category
$\operatorname{BIEXT}(P,Q;G)$}.

The general formal considerations of [4, no.~1] give several useful
consequences of the additive dependence of the categories
$\operatorname{BIEXT}(P,Q;G)$ on $G$.

\SGARefTarget{sga7.sga7i.evii.n2-5-1}{2.5.1}. First, for fixed $G$, the category
$\operatorname{BIEXT}(P,Q;G)$ has a composition law which assigns
functorially to two biextensions $E,E'$ of $(P,Q)$ by $G$ a
\underline{product biextension}, denoted $E.E'$. This construction is
simply the analogue of the familiar product of two extensions of one
commutative group by another. It is compatible with that construction
whether $E,E'$ are viewed as group extensions of $Q_P$ by $G_P$, or
as group extensions of $P_Q$ by $G_Q$. In particular, as a torsor
under $G_{P\times Q}$, $EE'$ is the product of the torsors under
$G_{P\times Q}$ defined by $E$ and $E'$.

The formation of $E.E'$ is moreover \underline{associative and
commutative} up to canonical isomorphisms, and this composition law on
$\operatorname{BIEXT}(P,Q;G)$ has a two-sided neutral object, called
the \underline{trivial biextension of $(P,Q)$ by $G$}. It may be
constructed, as in \emph{loc.\ cit.}, from the unique biextension
$E_0$ of $(P,Q)$ by the zero group $0$, up to isomorphism, by means of
the homomorphism $0\longrightarrow G$. From this description one
immediately obtains an alternative description of the trivial
biextension as the trivial torsor
\[
 E=P\times Q\times G.
\]
Its two group structures over $P$ and over $Q$ are obtained,
respectively, from the product group structures on $Q\times G$ and
$P\times G$ by the base changes $P\longrightarrow e$ and
$Q\longrightarrow e$.

Finally, for every biextension $E$ of $(P,Q)$ by $G$, one can define
the \underline{inverse biextension}, denoted $E^{-1}$, together with
an isomorphism from $EE^{-1}$ to the trivial biextension. This
determines $E^{-1}$ up to unique isomorphism and makes its dependence
on $E$ functorial.

In short, the category $\operatorname{BIEXT}(P,Q;G)$ behaves like a
group object in the category $\mathbf{Cat}$ of categories, except that
associativity, commutativity, unit, and inverse relations hold only up
to specified isomorphisms. One may call it a
\underline{pseudo-group} in the 2-category $\mathbf{Cat}$.

\SGARefTarget{sga7.sga7i.evii.n2-5-2}{2.5.2}. For fixed $P,Q,G$, the preceding properties imply the
following facts. For every biextension $E$ of $(P,Q)$ by $G$, the
monoid of endomorphisms of $E$ is canonically isomorphic to the monoid
of endomorphisms of the trivial biextension $\underline 1$: an
endomorphism $u$ of $\underline 1$ is sent to
$u.\mathrm{id}_E$ on $\underline 1.E\simeq E$. The endomorphism monoid
of $\underline 1$ is in fact a \underline{commutative group}, denoted
\begin{equation}\tag{\SGARefTarget{sga7.sga7i.evii.n2-5-2-1}{2.5.2.1}}
 \operatorname{Biext}^{0}(P,Q;G).
\end{equation}
Its composition law is also given by
$(u,v)\mapsto u.v\in
\operatorname{End}(\underline 1.\underline 1)
\simeq\operatorname{End}(\underline 1)$.

Likewise, the composition law $(E,E')\mapsto EE'$ equips the set
\begin{equation}\tag{\SGARefTarget{sga7.sga7i.evii.n2-5-2-2}{2.5.2.2}}
 \operatorname{Biext}^{1}(P,Q;G)
\end{equation}
of isomorphism classes of biextensions of $(P,Q)$ by $G$ with the
structure of a \underline{commutative group}; the inverse of the class
of $E$ is the class of $E^{-1}$.

\SGARefTarget{sga7.sga7i.evii.n2-5-3}{2.5.3}. Given a group homomorphism $G\longrightarrow G'$, the
corresponding functor
\begin{equation}\tag{\SGARefTarget{sga7.sga7i.evii.n2-5-3-1}{2.5.3.1}}
 \operatorname{BIEXT}(P,Q;G)
 \longrightarrow
 \operatorname{BIEXT}(P,Q;G')
\end{equation}
is compatible, up to canonical isomorphisms, with the operations
$E.E'$ and $E^{-1}$ of \SGARefLink{sga7.sga7i.evii.n2-5-1}{2.5.1}. Consequently, the two evident maps
\begin{equation}\tag{\SGARefTarget{sga7.sga7i.evii.n2-5-3-2}{2.5.3.2}}
 \operatorname{Biext}^{i}(P,Q;G)
 \longrightarrow
 \operatorname{Biext}^{i}(P,Q;G')
 \qquad(i=0,1)
\end{equation}
are \underline{group homomorphisms}. Moreover, if $E$ is a
biextension of $(P,Q)$ by $G$ and $E'$ the corresponding
biextension by $G'$, the natural homomorphism
\begin{equation}\tag{\SGARefTarget{sga7.sga7i.evii.n2-5-3-3}{2.5.3.3}}
 \operatorname{Hom}(E,E)=\operatorname{Aut}(E)
 \longrightarrow
 \operatorname{Hom}(E',E')=\operatorname{Aut}(E')
\end{equation}
is precisely (\SGARefLink{sga7.sga7i.evii.n2-5-3-2}{2.5.3.2}) for $i=0$.

\SGARefTarget{sga7.sga7i.evii.n2-5-4}{2.5.4}. The group $\operatorname{Biext}^{0}(P,Q;G)$ admits a simple
explicit description. If $E$ is a biextension of $(P,Q)$ by $G$, an
automorphism $u$ of $E$ is first of all an automorphism of the
underlying torsor under $G_{P\times Q}$, and hence is given by a
morphism
\[
 u_0:P\times Q\longrightarrow G
\]
through the formula
\[
 u(x)=u_0(j(x))x.
\]
Compatibility of $u$ with the first, respectively the second,
extension structure on $E$ is equivalent to additivity of $u_0$ in
the second, respectively the first, argument. Thus
\underline{$u$ is an automorphism of the biextension structure of $E$
if and only if $u_0$ is bilinear}. This gives a
\underline{canonical bijection}
\begin{equation}\tag{\SGARefTarget{sga7.sga7i.evii.n2-5-4-1}{2.5.4.1}}
 \operatorname{Biext}^{0}(P,Q;G)
 \simeq
 \operatorname{Hom}(P\otimes Q,G).
\end{equation}
It is readily checked that this bijection is independent of $E$,
respects the group structures, and is functorial in $G$ for the
functoriality on the left described in (\SGARefLink{sga7.sga7i.evii.n2-5-3-2}{2.5.3.2}).

Below we shall determine $\operatorname{Biext}^{1}(P,Q;G)$ and find
that, canonically and functorially in $G$, it is isomorphic to
\[
 \operatorname{Ext}^{1}
 \bigl(P\overset{\mathbb L}{\otimes}Q,G\bigr).
\]


% END INCLUDED SGA7I_English_current_buildable_source_20260823/components/78_expose_VII_section_2_5.tex

% BEGIN INCLUDED SGA7I_English_current_buildable_source_20260823/components/79_expose_VII_sections_2_6_and_2_7.tex
% SOURCE AUTHORITY:
% expose_VII_body.tex lines 846--870; scan indices 179--181; source folios 168--170.
% Source-backed corrections: the fibre in 2.6 is BIEXT(P,Q;G), not the printed
% BIEXT(P,P';G); the symmetric biextension has structural target Q x P.
% Hard stop before source line 873, section 2.8.

\SGARefTarget{sga7.sga7i.evii.n2-6}{2.6}. In \SGARefLink{sga7.sga7i.evii.n2-5}{2.5} only $G$ was varied, while $P$ and $Q$ were held fixed.
We now say a few words about variation in $P,Q$. The essential fact,
from which the others follow formally, is that the contravariant
dependence of $\operatorname{BIEXT}(P,Q;G)$ on $P,Q$
\underline{commutes} with its covariant dependence on $G$. Thus the
categories in question may be viewed as the fibre categories of a
\underline{cofibred category}
\begin{equation}\tag{\SGARefTarget{sga7.sga7i.evii.n2-6-1}{2.6.1}}
 \operatorname{BIEXT}'(\underline T)
 \longrightarrow
 \operatorname{GRAB}(\underline T)^{\mathrm{op}}\times
 \operatorname{GRAB}(\underline T)^{\mathrm{op}}\times
 \operatorname{GRAB}(\underline T),
\end{equation}
where the superscript $\mathrm{op}$ denotes the opposite category and
the structural functor is $E\mapsto(P,Q,G)$. This cofibred category is
additive in each of its three arguments, or, as we shall say,
\underline{tri-additive}.

By paraphrasing standard arguments, which we shall not reproduce, it
follows that the composition $E.E'$ in the fibre category
$\operatorname{BIEXT}(P,Q;G)$ may be described using the additive
structure of any one of the three arguments: the three resulting
pseudo-group structures are canonically isomorphic. Consequently, the
abelian groups
\begin{equation}\tag{\SGARefTarget{sga7.sga7i.evii.n2-6-2}{2.6.2}}
 \operatorname{Biext}^{i}(P,Q;G),\qquad i=0,1,
\end{equation}
may be regarded as two trilinear functors in $P,Q,G$, respectively in
$\operatorname{GRAB}(\underline T)^{\mathrm{op}}$,
$\operatorname{GRAB}(\underline T)^{\mathrm{op}}$, and
$\operatorname{GRAB}(\underline T)$. Equivalently, they are
\underline{trilinear functors} of variable $P,Q,G$ in
$\operatorname{GRAB}(\underline T)$, contravariant in $P,Q$ and
covariant in $G$, with values in the category of abelian groups.
For this variance, the isomorphism (\SGARefLink{sga7.sga7i.evii.n2-5-4-1}{2.5.4.1}) is functorial in all of
$P,Q,G$, and not merely in $G$. The analogous assertion will hold
below for the description of $\operatorname{Biext}^{1}$ by an
$\operatorname{Ext}^{1}$.

\SGARefTarget{sga7.sga7i.evii.n2-7}{2.7}. \underline{The symmetric biextension of a given biextension}.

If $E$ is a biextension of $(P,Q)$ by $G$, define its
\underline{symmetric biextension} $\,{}^{s}E$, which is a biextension
of $(Q,P)$ by $G$. The underlying sheaf of sets of $\,{}^{s}E$ is that
of $E$, and the actions of $G$ on $E$ and $\,{}^{s}E$ are the same.
The structural morphism
\[
 {}^{s}E\longrightarrow Q\times P
\]
is the composite of $E\longrightarrow P\times Q$ with the symmetry
$P\times Q\longrightarrow Q\times P$; finally, the left and right
extension structures of $\,{}^{s}E$ are, respectively, the right and
left extension structures of $E$. Evidently $E\mapsto{}^{s}E$ is an
automorphism of the category $\operatorname{BIEXT}(\underline T)$ of
biextensions, and indeed an \underline{involutive automorphism}:
\[
 {}^{s}({}^{s}E)=E.
\]
Every statement about bitorsors $E,\ldots$ therefore gives a symmetric
statement upon being applied to $\,{}^{s}E,\ldots$. We may call it the
statement obtained by \emph{exchanging left and right}. In general we
shall spell out only one of the two symmetric statements, choosing
between left and right, for example, according to our political
preferences.


% END INCLUDED SGA7I_English_current_buildable_source_20260823/components/79_expose_VII_sections_2_6_and_2_7.tex

% BEGIN INCLUDED SGA7I_English_current_buildable_source_20260823/components/80_expose_VII_section_2_8.tex
% SOURCE AUTHORITY:
% expose_VII_body.tex lines 873--899; scan indices 181--182; source folios 170--171.
% The conventional inverse-image notation f^* is used for the scan's typewritten f*.
% Page-header tokens at source lines 889 and 904 are excluded as running matter.
% Hard stop before source line 907, section 2.9.

\SGARefTarget{sga7.sga7i.evii.n2-8}{2.8}. \underline{Change of topos. Localization.}

Let
\[
 f:\underline T'\longrightarrow\underline T
\]
be a morphism of topoi. The inverse-image functor $f^*$ transforms a
biextension of $(P,Q)$ by $G$ into a biextension $f^*(E)$ of
$(f^*P,f^*Q)$ by $f^*G$, as always for a structure expressible in
terms of finite projective limits and arbitrary inductive limits. To
prove compatibility of the two extension structures on $f^*(E)$
induced by those on $E$, one expresses compatibility by commutativity
of (\SGARefLink{sga7.sga7i.evii.n2-1-1}{2.1.1}) in the universal case
\[
 S=P\times P\times Q\times Q,
\]
where $p,p',q,q'$ are the four projections. As expected,
$E\mapsto f^*(E)$ gives a canonical functor
\[
 f^*:\operatorname{BIEXT}(\underline T)
 \longrightarrow
 \operatorname{BIEXT}(\underline T').
\]

In what follows we shall use without further mention the fact that
every construction made with biextensions is compatible with inverse
images under morphisms of topoi. This applies in particular to the
product $EE'$ and inverse $E^{-1}$ considered in \SGARefLink{sga7.sga7i.evii.n2-5}{2.5}.

One may take for $\underline T'$ an \underline{induced topos}
$\underline T_{/S'}$, and for $f$ the canonical morphism, called the
localization morphism (SGA 4 IV 5),
\[
 f=i_{S'}:\underline T_{/S'}\longrightarrow\underline T.
\]
Then $f^*(E)$ is called the \underline{biextension induced} by $E$ on
$S'$, and is also denoted, as usual, by $E|_{S'}$ or $E_{S'}$.

Since restriction has the usual transitivity properties, the
categories $\operatorname{BIEXT}(\underline T_{/S})$, for variable
$S\in\underline T$, form the fibres of a \underline{fibred category}
over $\underline T$ (SGA 1 VI 8). In fact, this category is even a
\underline{stack} over $\underline T$ in the sense of Giraud [2]:
morphisms and objects in its fibres can be glued. For example, to give
a biextension over $S$ up to canonical isomorphism, it suffices to give
a covering family (that is, an epimorphic family) $S_i\longrightarrow
S$, a biextension $E_i$ over each $S_i$ (an object of
$\underline T_{/S_i}$), and, for every pair $(i,j)$, an isomorphism
between $E_i|_{S_j}$ and $E_j|_{S_i}$, subject to the usual cocycle
condition.

We call this the \underline{stack of biextensions over $\underline T$}.
Replacing $\underline T$ in (\SGARefLink{sga7.sga7i.evii.n2-3-3}{2.3.3}), and hence also in (\SGARefLink{sga7.sga7i.evii.n2-3-4}{2.3.4}) and
(\SGARefLink{sga7.sga7i.evii.n2-3-5}{2.3.5}), by $\underline T_{/S}$ for variable $S$ gives morphisms of
stacks over $\underline T$. The same applies to the operations $E.E'$
and $E^{-1}$ of (\SGARefLink{sga7.sga7i.evii.n2-5-1}{2.5.1}) in the categories
$\operatorname{BIEXT}(P_S,Q_S;G_S)$ for variable $S$, where $P,Q,G$
are fixed commutative groups over $\underline T$.


% END INCLUDED SGA7I_English_current_buildable_source_20260823/components/80_expose_VII_section_2_8.tex

% BEGIN INCLUDED SGA7I_English_current_buildable_source_20260823/components/81_expose_VII_section_2_9_opening_and_proposition.tex
% SOURCE AUTHORITY:
% expose_VII_body.tex lines 907--939; scan indices 183--185; source folios 172--174.
% Source-backed corrections: e_p is normalized to e_P; the second extension in
% the proof is of P_Q by G_Q, not the printed G_P.
% Hard stop before source line 941, Corollary 2.9.4.

\SGARefTarget{sga7.sga7i.evii.n2-9}{2.9}. \underline{Relations with birigidified torsors}.

In the terminology of \SGARefLink{sga7.sga7i.evii.n1-3-2}{1.3.2}, we saw in \SGARefLink{sga7.sga7i.evii.n2-2}{2.2} that, for every
biextension $E$ of $(P,Q)$ by $G$, the underlying torsor structure
under $G_{P\times Q}$ is canonically birigidified relative to the pair
of identity sections of $P$ and $Q$. This gives a canonical functor
\begin{equation}\tag{\SGARefTarget{sga7.sga7i.evii.n2-9-1}{2.9.1}}
 \operatorname{BIEXT}(P,Q;G)
 \xrightarrow{\ i\ }
 \operatorname{TORSBIRIG}(P,Q;G),
\end{equation}
where the category on the right consists of birigidified torsors over
$P\times Q$ under $G_{P\times Q}$.

Using the identity sections of $P$ and $Q$, for every product of $n$
objects each isomorphic to $P$ or $Q$ one obtains a canonical
morphism $\operatorname{inj}_i$ from the $i$-th factor into that
product. Thus, for a morphism from such a product to $G$, one may speak
of its restriction to the $i$-th factor. Consider the products
\begin{equation}\tag{\SGARefTarget{sga7.sga7i.evii.n2-9-2}{2.9.2}}
 \left\{
 \begin{array}{l}
 P\times Q,\\[1mm]
 P\times Q\times Q,\quad P\times Q\times Q\times Q,\\[1mm]
 P\times P\times Q,\quad P\times P\times P\times Q,\\[1mm]
 P\times P\times Q\times Q.
 \end{array}
 \right.
\end{equation}

\underline{Proposition \SGARefTarget{sga7.sga7i.evii.n2-9-3}{2.9.3}. Suppose that every morphism from one of
the products (\SGARefLink{sga7.sga7i.evii.n2-9-2}{2.9.2}) to $G$ which is trivial on each factor of that
product (embedded by the identity sections as just explained) is
trivial. Then:}

\medskip
\noindent
a) \underline{The canonical functor (\SGARefLink{sga7.sga7i.evii.n2-9-1}{2.9.1}) is fully faithful, and
both sides of (\SGARefLink{sga7.sga7i.evii.n2-9-1}{2.9.1}) are rigid categories. In particular, if $E$ is a
torsor under $G_{P\times Q}$ birigidified relative to $(e_P,e_Q)$,
there is at most one biextension structure on $E$ compatible with its
birigidified torsor structure.}

\medskip
\noindent
b) \underline{Such a biextension structure exists on $E$ if and only
if $E$ admits an extension structure of $Q_P$ by $G_P$ compatible with
its torsor structure under $G_P$, and an extension structure of $P_Q$
by $G_Q$ compatible with its torsor structure under $G_Q$
(cf.\ \SGARefLink{sga7.sga7i.evii.n1-3-5}{1.3.5}(b)). In that case there is a unique extension structure of
$Q_P$ by $G_P$ compatible with the torsor structure under $G_P$ and
the given rigidification over $P\times e_Q$, and a unique extension
structure of $P_Q$ by $G_Q$ compatible with the torsor structure under
$G_Q$ and the given rigidification over $e_P\times Q$; these two
extension structures are compatible in the sense of \SGARefLink{sga7.sga7i.evii.n2-1}{2.1}.}

\underline{Proof.} a) If $E$ is a birigidified torsor over
$P\times Q$, its automorphism group identifies with the group of
morphisms $P\times Q\longrightarrow G$ trivial on the two partial
factors. By the hypothesis applied to $P\times Q$, this group contains
only the identity automorphism, proving rigidity. Now let $E,E'$ be
biextensions of $(P,Q)$ by $G$. To see that every morphism
$f:E\longrightarrow E'$ of birigidified torsors is a morphism of
biextensions, apply \SGARefLink{sga7.sga7i.evii.n1-3-5}{1.3.5}(a) to the two extension structures defined
by $E,E'$ over $P$ and over $Q$, respectively. This uses the
hypothesis for $P\times Q\times Q$ and $P\times P\times Q$ in
(\SGARefLink{sga7.sga7i.evii.n2-9-2}{2.9.2}).

\noindent
b) Necessity is clear. For sufficiency, \SGARefLink{sga7.sga7i.evii.n1-3-5}{1.3.5}(b) first shows that the
stated condition gives unique extension structures of $Q_P$ by $G_P$
and of $P_Q$ by $G_Q$, compatible with the torsor structure and the
rigidifications. Here the hypothesis is used for
$P\times Q\times Q\times Q$ and $P\times P\times P\times Q$.
It remains to prove that the two extension structures are compatible,
that is, that diagram (\SGARefLink{sga7.sga7i.evii.n2-1-1}{2.1.1}) commutes. Its two composites are
isomorphisms of torsors over $P\times P\times Q\times Q$, and hence
differ by a morphism
\[
 u:P\times P\times Q\times Q\longrightarrow G.
\]
Because the partial trivializations of $E$ are identity sections for
the laws $\varphi$ and $\psi$, respectively, $u$ is trivial on each of
the four factors. The hypothesis for
$P\times P\times Q\times Q$ therefore makes $u$ trivial. This proves
the proposition.


% END INCLUDED SGA7I_English_current_buildable_source_20260823/components/81_expose_VII_section_2_9_opening_and_proposition.tex

% BEGIN INCLUDED SGA7I_English_current_buildable_source_20260823/components/82_expose_VII_corollary_2_9_4_through_remark_2_9_6.tex
% SOURCE AUTHORITY:
% expose_VII_body.tex lines 941--984; scan indices 185--187; source folios 174--176.
% Source-backed corrections: references to "2.9.2 b)" are to Proposition 2.9.3(b);
% Biext in (2.9.6.1) is written Biext^1 to identify the group of isomorphism classes.
% Hard stop before source line 986, section 2.10.

\underline{Corollary \SGARefTarget{sga7.sga7i.evii.n2-9-4}{2.9.4}. Let $P,Q,G$ be commutative groups
satisfying}
\begin{equation}\tag{\SGARefTarget{sga7.sga7i.evii.n2-9-4-1}{2.9.4.1}}
 \underline{\operatorname{Hom}}_{\mathrm{pt}}(P,G)=e,
 \qquad
 \underline{\operatorname{Hom}}_{\mathrm{pt}}(Q,G)=e
\end{equation}
\underline{(cf.\ \SGARefLink{sga7.sga7i.evii.n1-3-7}{1.3.7}). Then the hypothesis of Proposition 2.9.3 is
satisfied, and conclusions (a) and (b) of that proposition hold.
Moreover, if $E$ is a torsor under $G_{P\times Q}$ birigidified
relative to $(e_P,e_Q)$, it admits a biextension structure of $(P,Q)$
by $G$ compatible with its given structures if and only if the two
morphisms of pointed sheaves which it defines,}
\begin{equation}\tag{\SGARefTarget{sga7.sga7i.evii.n2-9-4-2}{2.9.4.2}}
 Q\longrightarrow R^1p_*(G_P),
 \qquad
 P\longrightarrow R^1q_*(G_Q),
\end{equation}
\underline{where $p:P\longrightarrow e$ and $q:Q\longrightarrow e$
are the structural morphisms, are morphisms of groups.}

The first assertion was proved in \SGARefLink{sga7.sga7i.evii.n1-3-7}{1.3.7}. For the second, necessity is
clear. To prove sufficiency, apply Proposition \SGARefLink{sga7.sga7i.evii.n2-9-3}{2.9.3}(b) and \SGARefLink{sga7.sga7i.evii.n1-3-5}{1.3.5}(b).
Multiplicativity of the first morphism in (\SGARefLink{sga7.sga7i.evii.n2-9-4-2}{2.9.4.2}) says that the
rigidified torsor $E$ under $G_P$ satisfies the condition of \SGARefLink{sga7.sga7i.evii.n1-3-5}{1.3.5}(b)
\underline{locally over $P$}. As observed in \SGARefLink{sga7.sga7i.evii.n1-3-7}{1.3.7}, that condition is
then true globally over $P$, so $E$ admits an extension structure of
$Q_P$ by $G_P$ compatible with its rigidified torsor structure.
Symmetrically, $E$ admits an extension structure of $P_Q$ by $G_Q$
compatible with its rigidified torsor structure under $G_Q$. The
conclusion follows from Proposition \SGARefLink{sga7.sga7i.evii.n2-9-3}{2.9.3}(b).

\underline{Example \SGARefTarget{sga7.sga7i.evii.n2-9-5}{2.9.5}.} Let $S$ be a scheme, let $P$ and $Q$ be
abelian schemes over $S$, and let $G=\underline G_{m,S}$ be the
multiplicative group scheme over $S$. We noted in \SGARefLink{sga7.sga7i.evii.n1-3-8}{1.3.8} that
(\SGARefLink{sga7.sga7i.evii.n2-9-4-1}{2.9.4.1}) then holds. A birigidified torsor $E$ over $P\times Q$ is
equivalently a birigidified invertible sheaf over $P\times Q$. Recall
that a \underline{divisorial correspondence on $P,Q$} is the
isomorphism class of such a birigidified invertible sheaf (see, for
example, [1]). In the fppf topology, the morphisms (\SGARefLink{sga7.sga7i.evii.n2-9-4-2}{2.9.4.2}) are
\begin{equation}\tag{\SGARefTarget{sga7.sga7i.evii.n2-9-5-1}{2.9.5.1}}
 Q\longrightarrow\underline{\operatorname{Pic}}_{P/S},
 \qquad
 P\longrightarrow\underline{\operatorname{Pic}}_{Q/S},
\end{equation}
defined by the divisorial correspondence. The theorem of the square
(\emph{loc.\ cit.}) says that these morphisms necessarily respect the
group structures. Thus \SGARefLink{sga7.sga7i.evii.n2-9-4}{2.9.4} shows that
\underline{the category of biextensions of $(P,Q)$ by
$G=\underline G_{m,S}$ is equivalent to the category of birigidified
invertible sheaves on $P\times Q$}; it is therefore a rigid category,
and its set of isomorphism classes identifies with the set of
divisorial correspondences on $(P,Q)$. This identification plainly
respects the natural group laws on the two sets, the first being that
of (\SGARefLink{sga7.sga7i.evii.n2-5-2-2}{2.5.2.2}).

\underline{Remark \SGARefTarget{sga7.sga7i.evii.n2-9-6}{2.9.6}.} Additivity of the first morphism in
(\SGARefLink{sga7.sga7i.evii.n2-9-4-2}{2.9.4.2}) is equivalent to factorization of the second through the
subsheaf $\underline{\operatorname{Ext}}^1(Q,G)$; note that, by \SGARefLink{sga7.sga7i.evii.n1-3-5}{1.3.5},
the canonical morphism
\[
 \underline{\operatorname{Ext}}^1(Q,G)
 \longrightarrow R^1q_*(G_Q)
\]
is injective. Likewise, additivity of the second morphism in
(\SGARefLink{sga7.sga7i.evii.n2-9-4-2}{2.9.4.2}) means that the first factors through
$\underline{\operatorname{Ext}}^1(P,G)$. This gives group
isomorphisms
\begin{equation}\tag{\SGARefTarget{sga7.sga7i.evii.n2-9-6-1}{2.9.6.1}}
 \operatorname{Biext}^{1}(P,Q;G)
 \simeq
 \operatorname{Hom}\bigl(Q,\underline{\operatorname{Ext}}^1(P,G)\bigr)
 \simeq
 \operatorname{Hom}\bigl(P,\underline{\operatorname{Ext}}^1(Q,G)\bigr).
\end{equation}
In the situation of \SGARefLink{sga7.sga7i.evii.n2-9-5}{2.9.5} these reduce to the familiar isomorphisms
\begin{equation}\tag{\SGARefTarget{sga7.sga7i.evii.n2-9-6-2}{2.9.6.2}}
 \operatorname{Corr}(P,Q)
 \simeq \operatorname{Hom}(P,Q')
 \simeq \operatorname{Hom}(Q,P'),
\end{equation}
where
\[
 P'=\underline{\operatorname{Ext}}^1(P,\underline G_m),
 \qquad
 Q'=\underline{\operatorname{Ext}}^1(Q,\underline G_m)
\]
are the dual abelian schemes of $P$ and $Q$, identified respectively
with $\underline{\operatorname{Pic}}^0_{P/S}$ and
$\underline{\operatorname{Pic}}^0_{Q/S}$. In both formulas,
$\operatorname{Hom}$ denotes the group of \underline{group
homomorphisms}.


% END INCLUDED SGA7I_English_current_buildable_source_20260823/components/82_expose_VII_corollary_2_9_4_through_remark_2_9_6.tex

% BEGIN INCLUDED SGA7I_English_current_buildable_source_20260823/components/83_expose_VII_section_2_10.tex
% SOURCE AUTHORITY:
% expose_VII_body.tex lines 986--1007; scan indices 187--190; source folios 176--179.
% Source-backed corrections: in 2.10.2 the multiplicative group is G_{m,S};
% in the A-linear dictionary the torsor is E_p, not the printed P_p.
% Hard stop before source line 1010, section 3.

\SGARefTarget{sga7.sga7i.evii.n2-10}{2.10}. \underline{Various generalizations of the notion of
biextension}. We mention them for completeness, although we shall not
use them.

\SGARefTarget{sga7.sga7i.evii.n2-10-1}{2.10.1}. One can define a biextension of $(P,Q)$ by $G$ without
assuming $P$ or $Q$ commutative, only $G$ being commutative. It is a
torsor under $G_{P\times Q}$ equipped with central extension
structures of $Q_P$ by $G_P$ and of $P_Q$ by $G_Q$, each compatible
with the torsor structure and compatible with one another, meaning
that they make (\SGARefLink{sga7.sga7i.evii.n2-1-1}{2.1.1}) commute. Most of the preceding developments
remain valid, \emph{mutatis mutandis}, for this more general notion.
Even when $P,Q$ are commutative, it is more general than the notion of
\SGARefLink{sga7.sga7i.evii.n2-1}{2.1}, which is the case in which the two underlying relative group
structures on $E$ are \underline{commutative}. In the noncommutative
case, the category of biextensions of $(P,Q)$ by $G$ still depends
additively on $G$, but not on $P$ and $Q$.

\SGARefTarget{sga7.sga7i.evii.n2-10-2}{2.10.2}. Paraphrasing \SGARefLink{sga7.sga7i.evii.n2-1}{2.1}, for a fixed integer $n\geq1$ one defines an
\underline{$n$-extension of a sequence of $n$ commutative groups}
$(P_1,P_2,\ldots,P_n)$ by a commutative group $G$. It consists of a
torsor $E$ under $G_{P_1\times\cdots\times P_n}$, equipped with $n$
partial extension structures over the partial products of $n-1$
factors, with a compatibility relation of type (\SGARefLink{sga7.sga7i.evii.n2-1-1}{2.1.1}) for every pair
of distinct indices between $1$ and $n$.

If $n\geq3$, the $P_i$ are abelian schemes over a scheme $S$, and
$G=\underline G_{m,S}$, every $n$-extension of
$(P_1,\ldots,P_n)$ by $G$ is \underline{trivial}.

\SGARefTarget{sga7.sga7i.evii.n2-10-3}{2.10.3}. The notion of biextension in \SGARefLink{sga7.sga7i.evii.n2-1}{2.1} may be viewed as a variant of
the notion of a commutative extension of commutative groups, that is,
of an extension of $\mathbb Z$-modules. Let $A$ be a ring of
$\underline T$, assumed commutative for simplicity, and let $P,Q,G$
be three $A$-modules. More generally one can define an
\underline{$A$-module biextension} of $(P,Q)$ by $G$, a variant of an
extension of $A$-modules.

Such a biextension is a biextension in the sense of \SGARefLink{sga7.sga7i.evii.n2-1}{2.1} whose
underlying partial extension structures are enriched to extension
structures of $A_P$-modules and $A_Q$-modules, respectively,
compatible with the relative module structures already present on
$Q_P,G_P$ and $P_Q,G_Q$. The two relative module structures on $E$
may be interpreted as two actions of the underlying multiplicative
monoid of $A$ on the sheaf of sets underlying $E$, and these two
actions are required to commute. This condition is added to the
compatibility expressed by (\SGARefLink{sga7.sga7i.evii.n2-1-1}{2.1.1}), and is itself expressed by an
analogous diagram, which we leave to the reader. Finally, the first
action of $A$ on $E$ is required to be distributive with respect not
only to the first composition law on $E$ but also to the second; the
same is required of the second action.

The dictionary of \SGARefLink{sga7.sga7i.evii.n1-1-6}{1.1.6} also has an $A$-linear variant. An extension
of the $A$-module $P$ by the $A$-module $G$ amounts to a family of
torsors $E_p$ under $G_p$, parametrized by points $p\in P(S)$ for
variable $S\in\operatorname{ob}\underline T$, with $E_p$ varying
\underline{$A$-linearly}, rather than merely additively, in $p$.
Besides the isomorphisms $\varphi_{p,p'}$ of (\SGARefLink{sga7.sga7i.evii.n1-1-3-1}{1.1.3.1}), this is
expressed by isomorphisms
\[
 \mu_{a,p}:\,{}^{a}E_p\xrightarrow{\ \sim\ }E_{ap}
 \qquad(p\in P(S),\ a\in A(S)),
\]
where $\,{}^{a}E_p$ denotes the image of the $G_S$-torsor $E_p$ under
extension of the structure group along
\[
 g\longmapsto ag:G_S\longrightarrow G_S.
\]
We leave to the reader the conditions on $\varphi$ and $\mu$ which
make these data correspond to an extension of $A$-modules $E$ of $P$
by $G$ (cf.\ \SGARefLink{sga7.sga7i.evii.n3-5-4}{3.5.4}(b)).


% END INCLUDED SGA7I_English_current_buildable_source_20260823/components/83_expose_VII_section_2_10.tex

% BEGIN INCLUDED SGA7I_English_current_buildable_source_20260823/components/84_expose_VII_section_3_1.tex
% SOURCE AUTHORITY:
% expose_VII_body.tex lines 1010--1110; scan indices 190--194; source folios 179--183.
% Direct scan comparison: physical PDF pages 191--195 at 600 dpi; diagram 3.1.3 checked
% node-by-node and arrow-by-arrow against physical page 192.
% Transparent notation corrections: printed/typewritten subscript letter o is rendered as zero
% throughout L_0, d_0, H_0 and degree 0; the erroneous prose reference 3.1.7.1 is 3.1.8.
% Hard stop before source line 1116, section 3.2.

\section*{\SGARefTarget{sga7.sga7i.evii.n3}{3}. Supplements on homological algebra}

The principal purpose of this section is to define the isomorphism
\[
 \Biext^1(P,Q;G)\simeq
 \Ext^1(P\overset{\mathbb L}{\otimes}Q,G)
\]
announced above. For this we shall need certain developments in
homological algebra which are of independent interest and will occupy
sections \SGARefLink{sga7.sga7i.evii.n3-1}{3.1} through \SGARefLink{sga7.sga7i.evii.n3-4}{3.4}. The developments in \SGARefLink{sga7.sga7i.evii.n3-1}{3.1} and \SGARefLink{sga7.sga7i.evii.n3-3}{3.3} may be
regarded as supplements to [4], especially to section 6.9 therein.

\subsection*{\SGARefTarget{sga7.sga7i.evii.n3-1}{3.1}. Construction of an additive cofibred category
$\underline\Psi_{L_\bullet}$ in terms of a chain complex $L_\bullet$
in an abelian category $\underline C$}

Let $\underline C$ be an abelian category and let $L_\bullet$ be a
chain complex in $\underline C$,
\begin{equation}\tag{\SGARefTarget{sga7.sga7i.evii.n3-1-1}{3.1.1}}
 L_2\longrightarrow L_1\longrightarrow L_0\longrightarrow0,
\end{equation}
written with lower-degree conventions, so that the differential has
degree $-1$. Thus $L_i=0$ for $i<0$. We shall define a cofibred
category
\[
 \underline\Psi=\underline\Psi_{L_\bullet}
\]
over $\underline C$; up to canonical equivalence of cofibred
categories it is the category of the same name in [4, 6.9]. The
objects of $\underline\Psi_{L_\bullet}$ are pairs
\[
 (X,\alpha),
\]
where $X$ is an extension of $L_0$ by an object $A$ of
$\underline C$, depending on $X$,
\begin{equation}\tag{\SGARefTarget{sga7.sga7i.evii.n3-1-2}{3.1.2}}
 0\longrightarrow A\xrightarrow{i}X\xrightarrow{j}L_0
 \longrightarrow0,
\end{equation}
and $\alpha$ is a trivialization of the inverse image $d_0^*(X)$ of
this extension by $d_0:L_1\longrightarrow L_0$. It is required that
the corresponding trivialization of the inverse image
$d_1^*d_0^*(X)$ by $d_1:L_2\longrightarrow L_1$ be the evident one,
obtained from the transitivity isomorphism
\[
 d_1^*d_0^*(X)\simeq(d_0d_1)^*(X)
\]
and the relation $d_0d_1=0$.

Equivalently, $\alpha$ is a lifting of $d_0:L_1\longrightarrow L_0$
to a morphism $\alpha:L_1\longrightarrow X$, subject to
$\alpha d_1=0$:
\begin{equation}\tag{\SGARefTarget{sga7.sga7i.evii.n3-1-3}{3.1.3}}
\begin{tikzcd}[row sep=large,column sep=large]
 &&0\arrow[d]\\
 &&A\arrow[d,"i"]\\
 &&X\arrow[d,"j"]\\
 L_2\arrow[r,"d_1"']&L_1\arrow[r,"d_0"']\arrow[ru,"\alpha"]
   &L_0\arrow[d]\\
 &&0
\end{tikzcd}
\qquad,\qquad \alpha d_1=0.\quad(*)
\end{equation}
Here is a simpler description, pointed out by Deligne:
$\underline\Psi_{L_\bullet}(A)$ is the category of extensions of the
complex $L_\bullet$ by $A$, regarded as a complex concentrated in
degree zero.

In this form one sees that the description of the objects
$(X,\alpha)$ depends only on the chain complex $[L_\bullet]$
\underline{truncated in degree 1}, obtained from $L_\bullet$ by
``killing $H_i(L_\bullet)$ for $i\geq2$'': its terms are zero in
degrees at least $2$, agree with those of $L_\bullet$ in degrees at
most $0$, and its degree-$1$ term is the cokernel of
$d_1:L_2\longrightarrow L_1$.

A homomorphism from $(X,\alpha)$ to $(X',\alpha')$ is a homomorphism
of extensions $u:X\longrightarrow X'$, compatible with the identity
on $L_0$, such that $u\alpha=\alpha'$. Homomorphisms are composed by
composing the corresponding homomorphisms of extensions. This gives
the announced category $\underline\Psi_{L_\bullet}$. Up to canonical
isomorphism, it depends only on the degree-$1$ truncation
$[L_\bullet]$ of $L_\bullet$.

There is an evident functor $(X,\alpha)\longmapsto A$,
\begin{equation}\tag{\SGARefTarget{sga7.sga7i.evii.n3-1-4}{3.1.4}}
 \pi:\underline\Psi_{L_\bullet}\longrightarrow\underline C.
\end{equation}
\underline{This functor is cofibred}. Indeed, let
$u:A\longrightarrow B$ be a morphism of $\underline C$, and let
$(X,\alpha)$ be an object of the fibre
$\underline\Psi_{L_\bullet}(A)$. Thus $X$ is an extension of $L_0$ by
$A$. Form the extension $u_*(X)$ of $L_0$ by $B$, the direct image of
$X$ under $u$, and equip it with the partial trivialization
$\beta=f\alpha$, where $f:X\longrightarrow u_*(X)$ is the canonical
morphism. It follows immediately that
\[
 f:(X,\alpha)\longrightarrow(u_*(X),\beta)
\]
is cocartesian for (\SGARefLink{sga7.sga7i.evii.n3-1-4}{3.1.4}): it has the usual universal property for
$u$-morphisms in $\underline\Psi_{L_\bullet}$ with source
$(X,\alpha)$. This is just the familiar and elementary corresponding
universal property of $f:X\longrightarrow u_*(X)$.

Finally, the cofibred category $\underline\Psi_{L_\bullet}$ is
\underline{additive}; that is, it satisfies
\begin{equation}\tag{\SGARefTarget{sga7.sga7i.evii.n3-1-5}{3.1.5}}
 \underline\Psi_{L_\bullet}(0)
 \quad\text{is equivalent to the one-point category},
\end{equation}
and, for arbitrary $A,B$ in $\underline C$,
\begin{equation}\tag{\SGARefTarget{sga7.sga7i.evii.n3-1-6}{3.1.6}}
 \underline\Psi_{L_\bullet}(A\times B)
 \longrightarrow
 \underline\Psi_{L_\bullet}(A)\times
 \underline\Psi_{L_\bullet}(B)
 \quad\text{is an equivalence}.
\end{equation}
This follows immediately from the corresponding properties of the
cofibred category of extensions of $L_0$ by variable objects of
$\underline C$. We shall freely use the consequences of the
additivity of $\underline\Psi_{L_\bullet}$ established in [4, 1] and
recalled above in \SGARefLink{sga7.sga7i.evii.n2-5}{2.5}.

\SGARefTarget{sga7.sga7i.evii.n3-1-7}{3.1.7}. Among these consequences, for every object $A$ of
$\underline C$ the category $\underline\Psi_{L_\bullet}(A)$ has a
\underline{zero}, or \underline{trivial}, object (\SGARefLink{sga7.sga7i.evii.n2-5-1}{2.5.1}). It is the
trivial extension $X=L_0+A$ of $L_0$ by $A$, equipped with the
partial trivialization $\alpha:L_1\longrightarrow X$ whose components
are $d_0$ and $0$.

Moreover, by \SGARefLink{sga7.sga7i.evii.n2-5-2}{2.5.2}, for every $(X,\alpha)$ in
$\underline\Psi_{L_\bullet}(A)$ the automorphism group of this object
is canonically isomorphic to the automorphism group of the trivial
object. Denote the latter commutative group by
$\Psi^0_{L_\bullet}(A)$. The explicit description above gives the
canonical isomorphisms
\begin{multline}\tag{\SGARefTarget{sga7.sga7i.evii.n3-1-8}{3.1.8}}
 \Psi^0_{L_\bullet}(A)
 =\Ker\bigl(\Hom(L_0,A)\longrightarrow\Hom(L_1,A)\bigr)
 \simeq\Hom(H_0(L_\bullet),A)\\
 \simeq\Ext^0(L_\bullet,A),
\end{multline}
induced by the usual identification of the automorphism group of the
extension $X$ with $\Hom(L_0,A)$. The isomorphism (\SGARefLink{sga7.sga7i.evii.n3-1-8}{3.1.8}) is plainly
functorial in $A$.

\SGARefTarget{sga7.sga7i.evii.n3-1-9}{3.1.9}. Finally, recall from \SGARefLink{sga7.sga7i.evii.n2-5-2}{2.5.2} that the natural multiplicative
structure on the fibre category $\underline\Psi_{L_\bullet}(A)$
endows its set of isomorphism classes with a natural group structure.
We denote this group by $\Psi^1_{L_\bullet}(A)$; it is an additive
functor of $A$.


% END INCLUDED SGA7I_English_current_buildable_source_20260823/components/84_expose_VII_section_3_1.tex

% BEGIN INCLUDED SGA7I_English_current_buildable_source_20260823/components/85_expose_VII_section_3_2.tex
% SOURCE AUTHORITY:
% expose_VII_body.tex lines 1116--1243; scan indices 194--197; source folios 183--186.
% Direct scan comparison: physical PDF pages 195--198 at 600 dpi.  Diagrams (3.2.3),
% (*), and (3.2.5.2) were checked node-by-node, arrow-by-arrow, and label-by-label.
% Transparent notation corrections: printed/typewritten subscript letter o is rendered as zero;
% the source's erroneous ``j : L_0 -> X'' in the injectivity conclusion is rendered as the
% mathematically required ``s : L_0 -> X'', consistently with (3.2.5.1).
% The source's two distinct displays both numbered (3.2.3) are preserved as printed.
% Hard stop before source line 1246, section 3.3.

\subsection*{\SGARefTarget{sga7.sga7i.evii.n3-2}{3.2}. Determination of $\Psi^1_{L_\bullet}(A)$}

Consider an object $(X,\alpha)$ (\SGARefLink{sga7.sga7i.evii.n3-1-3}{3.1.3}) of
$\underline\Psi_{L_\bullet}(A)$. Consider the cochain complex of
length one
\begin{equation}\tag{\SGARefTarget{sga7.sga7i.evii.n3-2-1}{3.2.1}}
 [X\longrightarrow L_0]:\quad
 0\longrightarrow X\longrightarrow L_0\longrightarrow0\longrightarrow
 \cdots
\end{equation}
defined by the morphism $j:X\longrightarrow L_0$, with $X$ as its
degree-zero component. It is a resolution of $A$, and hence may be
identified with $A$ as an object of the derived category
$D(\underline C)$ [5]. The diagram (\SGARefLink{sga7.sga7i.evii.n3-1-3}{3.1.3}) then gives a degree-one
morphism of complexes
\begin{equation}\tag{\SGARefTarget{sga7.sga7i.evii.n3-2-2}{3.2.2}}
 L_\bullet\longrightarrow[X\longrightarrow L_0],
\end{equation}
made explicit by the following diagram with anticommutative squares:
\begin{equation}\tag{\SGARefTarget{sga7.sga7i.evii.n3-2-3}{3.2.3}}
\begin{tikzcd}[column sep=large,row sep=huge]
 \cdots & L_2\arrow[r,"d_1"]
 & L_1\arrow[r,"d_0"]\arrow[d,"u^{-1}=-\alpha"']
 & L_0\arrow[r]\arrow[d,"u^0=\mathrm{id}"'] & 0 & \cdots\\
 \cdots & 0\arrow[r] & X\arrow[r,"j"] & L_0\arrow[r] & 0 & \cdots
\end{tikzcd}
\end{equation}
(We use the usual convention which associates with the chain complex
$L_\bullet$ a complex whose differential has degree $1$ and whose
degree-$i$ component is $L_{-i}$.) We thus obtain a degree-one
morphism in the derived category
\begin{equation}\tag{\SGARefTarget{sga7.sga7i.evii.n3-2-3.occ-2}{3.2.3}}
 c(X,\alpha):L_\bullet\longrightarrow A
 \qquad
 \bigl(\text{a degree-one morphism in }D(\underline C)\bigr),
\end{equation}
which plainly depends only on the isomorphism class of $(X,\alpha)$.
We have thereby defined a canonical map
\begin{equation}\tag{\SGARefTarget{sga7.sga7i.evii.n3-2-4}{3.2.4}}
 \Psi^1_{L_\bullet}(A)\longrightarrow\Ext^1(L_\bullet,A)
 \overset{\mathrm{def}}{=}
 \Hom_{D(\underline C)}(L_\bullet,A[1]).
\end{equation}

\underline{Theorem} \SGARefTarget{sga7.sga7i.evii.n3-2-5}{3.2.5}. \underline{The map} (\SGARefLink{sga7.sga7i.evii.n3-2-4}{3.2.4})
\underline{is an isomorphism of groups, functorial in the object}
$A$ \underline{of} $\underline C$.

We leave to the reader the verification that (\SGARefLink{sga7.sga7i.evii.n3-2-4}{3.2.4}) is additive and
functorial in $A$, and shall confine ourselves to proving that it is
bijective. For injectivity, let us first show that if the morphism of
complexes $u$ in (\SGARefLink{sga7.sga7i.evii.n3-2-3}{3.2.3}) is zero in $D(\underline C)$, then it is
already zero in the homotopy category $K(\underline C)$; that is,
there exists a morphism
\[
 s:L_0\longrightarrow X
\]
such that
\begin{equation}\tag{\SGARefTarget{sga7.sga7i.evii.n3-2-5-1}{3.2.5.1}}
 js=u^0=\mathrm{id}_{L_0},
 \qquad sd_0=-u^{-1}=\alpha.
\end{equation}
Indeed, the hypothesis says that the conclusion becomes true after
the morphism of complexes under consideration is composed with a
suitable quasi-isomorphism $v$
\begin{equation}\tag{*}
\begin{tikzcd}[column sep=large,row sep=huge]
 0\arrow[r] & X\arrow[r,"j"]\arrow[d]
 & L_0\arrow[r]\arrow[d] & 0\\
 0\arrow[r] & C^0\arrow[r,"j'"] & C^1\arrow[r] & C^2\arrow[r] & \cdots
\end{tikzcd}
\end{equation}
(and, after replacing $C^i$ by $0$ for $i\geq2$ and $C^1$ by
$\operatorname{Im}(C^0\longrightarrow C^1)
=\operatorname{Ker}(C^1\longrightarrow C^2)$ if necessary, we may
suppose that $C^i=0$ for $i\geq2$). But the morphism (*) induces the
identity on $A$, and therefore identifies $X$ with the fibre product
of $L_0$ and $C^0$ over $C^1$. A morphism
$s':L_0\longrightarrow C^0$ defining a homotopy $vu\sim0$, and hence
satisfying $j's'=v^1u^0$, consequently factors through a morphism
$s:L_0\longrightarrow X$ satisfying
$js=u^0=\mathrm{id}_{L_0}$. Thus $s$ splits the extension $X$ of
$L_0$ by $A$, so the complex $X\longrightarrow L_0$ is isomorphic
\underline{in} $K(\underline C)$ to the complex $K^\bullet$
concentrated in $A$. For the latter complex it is clear that
\[
 \Hom_{K(\underline C)}(L_\bullet,K^\bullet)
 \longrightarrow
 \Hom_{D(\underline C)}(L_\bullet,K^\bullet)
\]
is an isomorphism, proving the assertion.

Now the choice of a morphism $s:L_0\longrightarrow X$ satisfying
(\SGARefLink{sga7.sga7i.evii.n3-2-5-1}{3.2.5.1}) is precisely equivalent to a trivialization of the
extension $X$ of $L_0$ by $A$ compatible with the partial
trivialization $\alpha$. This proves the injectivity of (\SGARefLink{sga7.sga7i.evii.n3-2-4}{3.2.4}).

For surjectivity, start with an element of the right-hand side. By
definition it may be described by a degree-one morphism of complexes
\begin{equation}\tag{\SGARefTarget{sga7.sga7i.evii.n3-2-5-2}{3.2.5.2}}
\begin{tikzcd}[column sep=large,row sep=huge]
 \ldots & L_2\arrow[r,"d_1"]
 & L_1\arrow[r,"d_0"]\arrow[d,"u^{-1}"']
 & L_0\arrow[r]\arrow[d,"u^0"'] & 0\\
 & 0\arrow[r] & C^0\arrow[r] & C^1\arrow[r] & \ldots
\end{tikzcd}
\end{equation}
with anticommutative squares, where $C^\bullet$ is a resolution of
$A$ and, as above, we may suppose that $C^i=0$ for $i\geq2$. Thus
$C^0$ is an extension of $C^1$ by $A$, and we may form the extension
$X$ obtained by pulling it back along $L_0\longrightarrow C^1$.
Anticommutativity in (\SGARefLink{sga7.sga7i.evii.n3-2-5-2}{3.2.5.2}) shows that $-u^{-1}$ factors through a
morphism $\alpha:L_1\longrightarrow X$ which is a partial splitting,
that is, $j\alpha=d_0$ (with $j$ as in (\SGARefLink{sga7.sga7i.evii.n3-1-2}{3.1.2})). Moreover, the
relation $u^{-1}d_1=0$ in (\SGARefLink{sga7.sga7i.evii.n3-2-5-2}{3.2.5.2}) implies
$\alpha d_1=0$. Hence $(X,\alpha)$ is indeed an object of
$\underline\Psi_{L_\bullet}(A)$. By construction, the morphism
(\SGARefLink{sga7.sga7i.evii.n3-2-5-2}{3.2.5.2}) is the composite of the morphism (\SGARefLink{sga7.sga7i.evii.n3-2-3.occ-2}{3.2.3}) associated with
$(X,\alpha)$ and the morphism
\[
 [X\longrightarrow L_0]\longrightarrow C^\bullet
 =[C^0\longrightarrow C^1]
\]
whose components are $(f,u^0)$, where $f:X\longrightarrow C^0$ is
the canonical morphism obtained from the definition of $X$ as the
pullback of the extension $C^0$.

The preceding morphism is a morphism of resolutions of $A$ and
induces the identity on $A$. It follows, by the definition
$\Ext^1=\Hom^1_{D(\underline C)}$, that the given element of
$\Ext^1(L_\bullet,A)$ described by (\SGARefLink{sga7.sga7i.evii.n3-2-5-2}{3.2.5.2}) is precisely
$c(X,\alpha)$. This proves the surjectivity of (\SGARefLink{sga7.sga7i.evii.n3-2-4}{3.2.4}) and completes
the proof of \SGARefLink{sga7.sga7i.evii.n3-2-5}{3.2.5}.


% END INCLUDED SGA7I_English_current_buildable_source_20260823/components/85_expose_VII_section_3_2.tex

% BEGIN INCLUDED SGA7I_English_current_buildable_source_20260823/components/86_expose_VII_section_3_3.tex
% SOURCE AUTHORITY:
% expose_VII_body.tex lines 1246--1352; scan indices 197--203; source folios 186--192.
% Direct scan comparison: physical PDF pages 198--204 at 600 dpi.
% Transparent source corrections: the two references printed as 3.2.3 in 3.3.2 are 3.3.2.1;
% the quasi-maximal construction is written with the induced map \bar d_0:L'_1->L_0 rather
% than the internally inconsistent printed d'_1/d_1 notation; the canonical map A->X is i,
% not the printed j; and the boundary arrow in (3.3.6.1) is \partial, as the following prose
% and (3.3.6.2) require.  The typo ``auticommutative'' is translated as anticommutative.
% Hard stop before source line 1355, section 3.4.

\subsection*{\SGARefTarget{sga7.sga7i.evii.n3-3}{3.3}. Properties of the cofibred category
$\underline\Psi_{L_\bullet}$}

In this section, which is not essential for understanding what
follows, we freely use the notions of [4], our purpose being to
compare the preceding results with those of that work.

\underline{Theorem} \SGARefTarget{sga7.sga7i.evii.n3-3-1}{3.3.1}. \underline{The additive cofibred category}
$\underline\Psi_{L_\bullet}$ \underline{over} $\underline C$
\underline{is left exact} ([4, \SGARefLink{sga7.sga7i.evii.n2-2}{2.2}]), \underline{admits a
quasi-maximal object} ([4, 4.4]), \underline{and, for every object}
$E=(X,\alpha)$ \underline{of this category, there exists a typical
complex} $L^E_\bullet$ \underline{for} $E$ ([4, \SGARefLink{sga7.sga7i.evii.n3-1}{3.1}]).
\underline{Moreover, the typical pro-complex of}
$\underline\Psi_{L_\bullet}$ ([4, \SGARefLink{sga7.sga7i.evii.n3-7}{3.7} and 4.4])
\underline{is canonically isomorphic in} $D(\underline C)$
\underline{to the truncated complex} $[L_\bullet]$
\underline{considered in} \SGARefLink{sga7.sga7i.evii.n3-1}{3.1}.

We now spell out the meaning of the several assertions in \SGARefLink{sga7.sga7i.evii.n3-3-1}{3.3.1}.

\SGARefTarget{sga7.sga7i.evii.n3-3-2}{3.3.2}. To say that the cofibred category
$\underline\Psi_{L_\bullet}$ is \underline{left exact} means that it
is additive and also satisfies the following condition:

\SGARefTarget{sga7.sga7i.evii.n3-3-2-1}{3.3.2.1}. For every exact sequence
\[
 0\longrightarrow A'\xrightarrow{u}A\xrightarrow{v}A''
 \longrightarrow0
\]
in $\underline C$, the natural functor (in the notation of [4, 1.8])
\[
 \underline\Psi_{L_\bullet}(A)
 \longrightarrow
 \underline\Psi_{L_\bullet}([A\longrightarrow A''])
\]
is an equivalence of categories.

Recall that if $\underline\Psi$ is a cofibred category over
$\underline C$ and $v:A\longrightarrow A''$ is a morphism in
$\underline C$, then $\underline\Psi([A\longrightarrow A''])$ is the
category of pairs $(E,\rho)$ in which $E$ is an object of
$\underline\Psi(A)$ and $\rho$ is a ``trivialization'' of the
corresponding object $v_*(E)$ of $\underline\Psi(A'')$; that is, an
isomorphism from it to the zero object (\SGARefLink{sga7.sga7i.evii.n2-5-1}{2.5.1}) of
$\underline\Psi(A'')$. The definition of the functor
\[
 \underline\Psi(A')\longrightarrow
 \underline\Psi([A\longrightarrow A''])
\]
in \SGARefLink{sga7.sga7i.evii.n3-3-2-1}{3.3.2.1} is evident. When
$\underline\Psi=\underline\Psi_{L_\bullet}$, the objects of
$\underline\Psi_{L_\bullet}([A\longrightarrow A''])$ are explicitly
the triples $(X,\alpha,\rho)$, where $(X,\alpha)$ is as in (\SGARefLink{sga7.sga7i.evii.n3-1-3}{3.1.3})
and $\rho:X\longrightarrow A''$ is a morphism in $\underline C$
satisfying
\begin{equation}\tag{*}
 \rho i=v,\qquad \rho\alpha=0.
\end{equation}
If $v:A\longrightarrow A''$ belongs to a short exact sequence as in
\SGARefLink{sga7.sga7i.evii.n3-3-2-1}{3.3.2.1}, then, as is well known, giving a morphism satisfying the
first relation (*) is equivalent, by the left exactness of the
cofibred category of extensions of $L_0$ by variable objects of
$\underline C$, to giving an extension $X'$ of $L_0$ by
$A'\simeq\Ker v$ which is a subextension of $X$: one takes
$X'=\Ker\rho$. The second relation (*) says that $\alpha$ factors
through a morphism $\alpha':L_1\longrightarrow X'$. Up to the
compatibilities whose verification is immediate, this proves the
left exactness stated in \SGARefLink{sga7.sga7i.evii.n3-3-2-1}{3.3.2.1}.

\SGARefTarget{sga7.sga7i.evii.n3-3-3}{3.3.3}. Recall that an object $E_0$ of a cofibred category
$\underline\Psi$ over $\underline C$ is called \underline{quasi-maximal}
if, for every object $E$ of $\underline\Psi$ lying over an object
$A$ of $\underline C$, there is a morphism $E\longrightarrow E'$ in
$\underline\Psi$ inducing a \underline{monomorphism}
$A\longrightarrow A'$ in $\underline C$, such that $E$ is
\underline{dominated by} $E_0$, meaning that there is a morphism
$E_0\longrightarrow E'$. If $\underline C$ has enough injective
objects, this is equivalent to saying that every object of
$\underline\Psi$ over an injective object of $\underline C$ is
dominated by $E_0$.

For $\underline\Psi=\underline\Psi_{L_\bullet}$ there is the
following canonical quasi-maximal object. Put
\[
 L'_1=\Coker(d_1:L_2\longrightarrow L_1)
\]
and let $\bar d_0:L'_1\longrightarrow L_0$ be the morphism induced
by $d_0:L_1\longrightarrow L_0$. Let $E_0=(X_0,\alpha_0)$ be the
object of $\underline\Psi_{L_\bullet}(L'_1)$ for which
$X_0=L_0\oplus L'_1$ is the trivial extension of $L_0$ by $L'_1$
and $\alpha_0$ is the composite
\[
 L_1\twoheadrightarrow L'_1
 \xrightarrow{(\bar d_0,\,\mathrm{id}_{L'_1})}
 L_0\oplus L'_1.
\]

To see that $E_0$ is quasi-maximal, observe that for every object
$(X,\alpha)$ of $\underline\Psi_{L_\bullet}$ one can find a morphism
from it to an object $(X',\alpha')$, inducing a
\underline{monomorphism} $A\longrightarrow A'$ in $\underline C$,
such that the extension $X'$ of $L_0$ by $A'$ is \underline{split}.
Indeed, it suffices to take $A'=X$ and
$(X',\alpha')=i_*(X,\alpha)$, with the notation of (\SGARefLink{sga7.sga7i.evii.n3-1-2}{3.1.2}). On the
other hand, it is immediate that any object $(X',\alpha')$ of
$\underline\Psi$ for which the extension $X'$ is split is dominated
by $E_0$.

\SGARefTarget{sga7.sga7i.evii.n3-3-4}{3.3.4}. Let $E$ be an object of an additive cofibred category
$\underline\Psi$ over $\underline C$, lying above an object $A$ of
$\underline C$, and consider the following functor of the variable
object $B$ of $\underline C$:
\begin{equation}\tag{\SGARefTarget{sga7.sga7i.evii.n3-3-4-1}{3.3.4.1}}
 B\longmapsto\Hom(E,\theta_B),
\end{equation}
where $\theta_B$ denotes the trivial object of $\underline\Psi(B)$.
We say that $E$ \underline{admits a typical complex} if this functor
is representable by an object $\Omega_E$ of $\underline C$. Since
the functor in question maps canonically to the functor
$B\longmapsto\Hom(A,B)$ represented by $A$, one obtains a canonical
morphism
\begin{equation}\tag{\SGARefTarget{sga7.sga7i.evii.n3-3-4-2}{3.3.4.2}}
 A\longrightarrow\Omega_E,
\end{equation}
and the corresponding length-one \underline{chain complex}
$A\longrightarrow\Omega_E$ in $\underline C$ is called the
\underline{typical complex of} $E$.

When $\underline\Psi=\underline\Psi_{L_\bullet}$ and
$E=(X,\alpha)$, one checks immediately that the functor (\SGARefLink{sga7.sga7i.evii.n3-3-4-1}{3.3.4.1}) is
represented by $\Coker\alpha$, and that the canonical morphism
(\SGARefLink{sga7.sga7i.evii.n3-3-4-2}{3.3.4.2}) is simply the composite
\[
 A\xrightarrow{i}X\xrightarrow{\mathrm{can}}\Coker\alpha.
\]

\SGARefTarget{sga7.sga7i.evii.n3-3-5}{3.3.5}. If an additive cofibred category $\underline\Psi$ over
$\underline C$ admits a quasi-maximal object (\SGARefLink{sga7.sga7i.evii.n3-3-3}{3.3.3}), and the typical
complexes of its objects exist, then the typical complex $T_\bullet$
of a quasi-maximal object $E_0$, regarded as an object of the
\underline{derived category} $D(\underline C)$, is independent, up
to canonical isomorphism, of the choice of $E_0$ ([4, 4.4]). It is
called the \underline{typical complex of the additive cofibred
category} $\underline\Psi$. If, in addition, $\underline\Psi$ is
left exact (\SGARefLink{sga7.sga7i.evii.n3-3-2}{3.3.2}), then $\underline\Psi$ can be reconstructed from
the typical complex $T_\bullet$, up to an equivalence of fibred
categories over $\underline C$ unique up to unique isomorphism
([4, 6.10]). There are also canonical isomorphisms, functorial in
$A\in\underline C$ ([4, 5.4]),
\begin{equation}\tag{\SGARefTarget{sga7.sga7i.evii.n3-3-5-1}{3.3.5.1}}
 \Psi^0(A)\xrightarrow{\sim}\Ext^0(T_\bullet,A),
 \qquad
 \Psi^1(A)\xrightarrow{\sim}\Ext^1(T_\bullet,A).
\end{equation}

To complete the comparison with [4], we shall prove that when
$\underline\Psi=\underline\Psi_{L_\bullet}$ its typical complex is
the truncation $[L_\bullet]$. This gives another definition of the
isomorphisms (\SGARefLink{sga7.sga7i.evii.n3-1-8}{3.1.8}) and (\SGARefLink{sga7.sga7i.evii.n3-2-4}{3.2.4}) in terms of (\SGARefLink{sga7.sga7i.evii.n3-3-5-1}{3.3.5.1}), and one
checks that it is compatible with the definitions in \SGARefLink{sga7.sga7i.evii.n3-1}{3.1} and \SGARefLink{sga7.sga7i.evii.n3-2}{3.2}.

For this assertion we may plainly restrict to the case where
$L_\bullet$ is already truncated in degree $1$, since \SGARefLink{sga7.sga7i.evii.n3-1}{3.1} noted that
truncation does not change $\underline\Psi_{L_\bullet}$ up to an
isomorphism of cofibred categories. The explicit construction in
\SGARefLink{sga7.sga7i.evii.n3-3-4}{3.3.4} then shows immediately that the typical complex of the
quasi-maximal object described in \SGARefLink{sga7.sga7i.evii.n3-3-3}{3.3.3} is canonically isomorphic to
$L_\bullet$.

\underline{Remark} \SGARefTarget{sga7.sga7i.evii.n3-3-6}{3.3.6}. Let
\[
 0\longrightarrow A'\longrightarrow A\longrightarrow A''
 \longrightarrow0
\]
be an exact sequence in $\underline C$. It gives an exact sequence
of $\Ext^i$ groups
\[
 0\longrightarrow\Ext^0(L_\bullet,A')
 \longrightarrow\Ext^0(L_\bullet,A)
 \longrightarrow\cdots\longrightarrow\Ext^1(L_\bullet,A)
 \longrightarrow\Ext^1(L_\bullet,A''),
\]
which could of course be continued to the right by introducing
$\Ext^i$ for $i\geq2$. Using the isomorphisms (\SGARefLink{sga7.sga7i.evii.n3-1-8}{3.1.8}) and (\SGARefLink{sga7.sga7i.evii.n3-2-4}{3.2.4}),
we obtain the exact sequence
\begin{multline}\tag{\SGARefTarget{sga7.sga7i.evii.n3-3-6-1}{3.3.6.1}}
 0\longrightarrow\Psi^0_{L_\bullet}(A')
 \longrightarrow\Psi^0_{L_\bullet}(A)
 \longrightarrow\Psi^0_{L_\bullet}(A'')
 \xrightarrow{\partial}\Psi^1_{L_\bullet}(A')\\
 \longrightarrow\Psi^1_{L_\bullet}(A)
 \longrightarrow\Psi^1_{L_\bullet}(A'').
\end{multline}
All arrows other than $\partial$ are defined simply by functoriality
of $\Psi^0_{L_\bullet}$ and $\Psi^1_{L_\bullet}$. On the other hand,
[4, \SGARefLink{sga7.sga7i.evii.n2-5-2}{2.5.2}] defines an exact sequence (3.3.6.1), for every left-exact
cofibred category $\underline\Psi$ over $\underline C$, through a
morphism
\begin{equation}\tag{\SGARefTarget{sga7.sga7i.evii.n3-3-6-2}{3.3.6.2}}
 \partial:\Psi^0(A'')\longrightarrow\Psi^1(A').
\end{equation}
It associates with an element $u$ of the left-hand side the class,
in the right-hand side, of the object of $\underline\Psi(A')$,
determined up to canonical isomorphism, whose image in
$\underline\Psi([A\longrightarrow A''])$ (notation of 3.3.2) is the
pair $(X_0,u)$, where $X_0$ is the trivial object of
$\underline\Psi(A)$.

For $\underline\Psi=\underline\Psi_{L_\bullet}$, the morphism
(\SGARefLink{sga7.sga7i.evii.n3-3-6-2}{3.3.6.2}) defined by this \underline{``geometric'' construction is
the negative of the boundary morphism} $\partial$ \underline{in}
(\SGARefLink{sga7.sga7i.evii.n3-3-6-1}{3.3.6.1}). Thus, after the identifications above and a change of
sign in $\partial$, the latter exact sequence is precisely that of
[4, \SGARefLink{sga7.sga7i.evii.n2-5-2}{2.5.2}]. This follows from the final assertion of [4, 5.4], in
view of the compatibilities noted in \SGARefLink{sga7.sga7i.evii.n3-3-5}{3.3.5}.

\underline{Remark} \SGARefTarget{sga7.sga7i.evii.n3-3-6-3}{3.3.6.3}. We take this opportunity to clarify a
point in [4, no.~3.11] concerning the fundamental isomorphisms
(3.11.3) and (3.11.4), which play an essential role in the proof of
[4, 5.4]. The diagram on page 35 of [4] is commutative, not
anticommutative, and therefore does not as it stands define a
morphism
\[
 u:[L^X_1\longrightarrow L^X_0]\longrightarrow K^\bullet[1].
\]
It must be made anticommutative by replacing $u^1$ by $-u^1$. One
could instead adopt the opposite convention and replace $u^0$ by
$-u^0$; however, the stated convention, in which $u^1$ becomes
$-u^1$, is the one for which (3.11.3) generalizes the isomorphism
(\SGARefLink{sga7.sga7i.evii.n3-5-3}{3.5.3}) of [4], as claimed there. With this clarification, the claim
made in [4, 5.4] must also be corrected: the two boundary morphisms
considered there \underline{differ by a sign}.

\SGARefTarget{sga7.sga7i.evii.n3-3-7}{3.3.7}. In [4, 6.13] the variance of the cofibred category
$\underline\Psi_{L_\bullet}$ over $\underline C$ as $L_\bullet$
varies was also made explicit, in particular by reconstructing the
category of cocartesian functors from $\underline\Psi_{L_\bullet}$
to $\underline\Psi_{L'_\bullet}$ in terms of the chain complexes
$L_\bullet$ and $L'_\bullet$. The set of isomorphism classes of such
cocartesian functors is canonically in bijection with
\[
 \Hom_{D(\underline C)}(L'_\bullet,L_\bullet).
\]
We shall use only a small part of these results, which we develop in
the next section using the simple description of
$\underline\Psi_{L_\bullet}$ given in \SGARefLink{sga7.sga7i.evii.n3-1}{3.1}.

% END INCLUDED SGA7I_English_current_buildable_source_20260823/components/86_expose_VII_section_3_3.tex

% BEGIN INCLUDED SGA7I_English_current_buildable_source_20260823/components/87_expose_VII_section_3_4_opening_through_3_4_7.tex
% SOURCE AUTHORITY:
% expose_VII_body.tex lines 1355--1434; scan indices 203--205; source folios 192--194.
% Direct scan comparison: physical PDF pages 204--206 at 600 dpi.  Diagrams (3.4.2)
% and (3.4.6) were checked node-by-node, arrow-by-arrow, and label-by-label.
% Transparent corrections: typewritten subscript letter o is rendered as zero; the arrow
% omitted by the working transcription in (3.4.5) is restored from the scan; and the proof's
% unindexed f is rendered as the required f^1.
% Hard stop before source line 1436, 3.4.8.

\subsection*{\SGARefTarget{sga7.sga7i.evii.n3-4}{3.4}. Variance of $\underline\Psi_{L_\bullet}$ with
$L_\bullet$}

Consider a morphism of complexes
\begin{equation}\tag{\SGARefTarget{sga7.sga7i.evii.n3-4-1}{3.4.1}}
 f:L'_\bullet\longrightarrow L_\bullet
\end{equation}
given by a commutative diagram
\begin{equation}\tag{\SGARefTarget{sga7.sga7i.evii.n3-4-2}{3.4.2}}
\begin{tikzcd}[column sep=large,row sep=large]
 {}\arrow[r]&L'_2\arrow[r]\arrow[d,"f_2"']
 &L'_1\arrow[r]\arrow[d,"f_1"']
 &L'_0\arrow[r]\arrow[d,"f_0"']&0\\
 {}\arrow[r]&L_2\arrow[r]&L_1\arrow[r]&L_0\arrow[r]&0.
\end{tikzcd}
\end{equation}
We associate with it a canonical functor
\begin{equation}\tag{\SGARefTarget{sga7.sga7i.evii.n3-4-3}{3.4.3}}
 f^*:\underline\Psi_{L_\bullet}\longrightarrow
 \underline\Psi_{L'_\bullet}.
\end{equation}
It is in fact a cocartesian functor of cofibred categories over
$\underline C$: it is compatible with the structural functors to
$\underline C$ and carries cocartesian morphisms to cocartesian
morphisms; see SGA 1, Exposé VI.

If $(X,\alpha)$ is an object of $\underline\Psi_{L_\bullet}$, define
$f^*(X,\alpha)=(X',\alpha')$ as follows. The extension $X'$ of
$L'_0$ by $A$ is the pullback, along $f_0$, of the given extension
$X$ of $L_0$ by $A$, and $\alpha'$ is the partial trivialization
deduced from $\alpha$ by transitivity in the second square of
(\SGARefLink{sga7.sga7i.evii.n3-4-2}{3.4.2}). The compatibility of $\alpha'$ with the relation
$d'_0d'_1=0$ follows immediately from the analogous compatibility
for $\alpha$. By construction, (\SGARefLink{sga7.sga7i.evii.n3-4-3}{3.4.3}) commutes with the structural
functors to $\underline C$. Its cocartesian character follows at
once from the fact that, in the category of extensions of objects of
$\underline C$, direct images of extensions commute with inverse
images of extensions.

The functor (\SGARefLink{sga7.sga7i.evii.n3-4-3}{3.4.3}) induces morphisms of functors
\begin{equation}\tag{\SGARefTarget{sga7.sga7i.evii.n3-4-4}{3.4.4}}
 f^i:\Psi^i_{L_\bullet}\longrightarrow\Psi^i_{L'_\bullet},
 \qquad i=0,1
\end{equation}
(compare \SGARefLink{sga7.sga7i.evii.n2-6}{2.6}). The morphism (\SGARefLink{sga7.sga7i.evii.n3-4-1}{3.4.1}), on the other hand, induces
morphisms
\begin{equation}\tag{\SGARefTarget{sga7.sga7i.evii.n3-4-5}{3.4.5}}
 f^i:\Ext^i(L_\bullet,-)\longrightarrow\Ext^i(L'_\bullet,-),
 \qquad i\in\mathbb Z.
\end{equation}
For $i=0,1$, the morphisms (\SGARefLink{sga7.sga7i.evii.n3-4-4}{3.4.4}) and (\SGARefLink{sga7.sga7i.evii.n3-4-5}{3.4.5}) are compatible with
the isomorphisms (\SGARefLink{sga7.sga7i.evii.n3-1-8}{3.1.8}) and (\SGARefLink{sga7.sga7i.evii.n3-2-4}{3.2.4}). In other words, for
$A\in\operatorname{Ob}\underline C$ and $i=0,1$, the following
diagram is commutative:
\begin{equation}\tag{\SGARefTarget{sga7.sga7i.evii.n3-4-6}{3.4.6}}
\begin{tikzcd}[column sep=large,row sep=large]
 \Psi^i_{L_\bullet}(A)\arrow[r,"\sim"]\arrow[d]
 &\Ext^i(L_\bullet,A)\arrow[d]\\
 \Psi^i_{L'_\bullet}(A)\arrow[r,"\sim"]
 &\Ext^i(L'_\bullet,A).
\end{tikzcd}
\end{equation}
The verification of this compatibility is essentially immediate and
is left to the reader.

\underline{Proposition} \SGARefTarget{sga7.sga7i.evii.n3-4-7}{3.4.7}. \underline{The functor} (\SGARefLink{sga7.sga7i.evii.n3-4-3}{3.4.3})
\underline{is faithful} (\underline{respectively fully faithful,
respectively an equivalence of categories}) \underline{if and only
if}
\[
 H_0(f):H_0(L'_\bullet)\longrightarrow H_0(L_\bullet)
\]
\underline{is an epimorphism} (\underline{respectively, if}
$H_0(f)$ \underline{is an isomorphism and} $H_1(f)$
\underline{is an epimorphism; respectively, if} $H_0(f)$
\underline{and} $H_1(f)$ \underline{are isomorphisms}).

Indeed, with $f^i$ denoting the functorial morphism (\SGARefLink{sga7.sga7i.evii.n3-4-4}{3.4.4}), the
required condition is that $f^0$ be a monomorphism (respectively,
that $f^0$ be an isomorphism and $f^1$ a monomorphism; respectively,
that $f^0$ and $f^1$ be isomorphisms). By the commutativity of
(\SGARefLink{sga7.sga7i.evii.n3-4-6}{3.4.6}), these conditions may be expressed by the analogous
conditions on the functors $\Ext^i$, $i=0,1$. The proposition then
reduces to a straightforward statement of homological algebra,
which we leave to the reader as an exercise.


% END INCLUDED SGA7I_English_current_buildable_source_20260823/components/87_expose_VII_section_3_4_opening_through_3_4_7.tex

% BEGIN INCLUDED SGA7I_English_current_buildable_source_20260823/components/88_expose_VII_sections_3_4_8_and_3_4_9.tex
% SOURCE AUTHORITY:
% expose_VII_body.tex lines 1436--1542; scan indices 205--209; source folios 194--198.
% Direct scan comparison: physical PDF pages 206--210 at 600 dpi, with the dense
% diagram inspected separately at original detail.  The apparent second arrowhead
% on the X'--L'_0 diagonal belongs to the crossing alpha' arrow: the source has the
% single morphism j':X'->L'_0 required by the extension, not a bidirectional arrow.
% Transparent corrections: 0_{L,A'} is rendered 0_{L,A}; unprimed j in the four
% defining relations is rendered j'; the concluding Psi_L is rendered Psi_{L''};
% and the missing comma in Ext^0(L',A) is restored.
% Hard stop before source line 1545, section 3.5.

\SGARefTarget{sga7.sga7i.evii.n3-4-8}{3.4.8}. Consider now an exact sequence of chain complexes
\begin{equation}\tag{\SGARefTarget{sga7.sga7i.evii.n3-4-8-1}{3.4.8.1}}
 0\longrightarrow L_\bullet\xrightarrow{f}L'_\bullet
 \xrightarrow{g}L''_\bullet\longrightarrow0.
\end{equation}
It gives functors of the form (\SGARefLink{sga7.sga7i.evii.n3-4-3}{3.4.3}),
\[
 \underline\Psi_{L''_\bullet}\xrightarrow{g^*}
 \underline\Psi_{L'_\bullet}\xrightarrow{f^*}
 \underline\Psi_{L_\bullet}.
\]
The transitivity property implicit in functors of the form (\SGARefLink{sga7.sga7i.evii.n3-4-3}{3.4.3})
gives an isomorphism of functors
\[
 f^*g^*\simeq(gf)^*=(0_{L_\bullet\to L''_\bullet})^*.
\]
The last member is evidently canonically isomorphic to the
cocartesian zero functor from $\underline\Psi_{L''_\bullet}$ to
$\underline\Psi_{L_\bullet}$.  The preceding isomorphism therefore
defines a canonical functor
\begin{equation}\tag{\SGARefTarget{sga7.sga7i.evii.n3-4-8-2}{3.4.8.2}}
 \underline\Psi_{L''_\bullet}\longrightarrow
 \underline\Psi_{[L_\bullet\longrightarrow L'_\bullet]}.
\end{equation}
Here $\underline\Psi_{[L_\bullet\longrightarrow L'_\bullet]}$
denotes the category of pairs $(X',t)$ in which $X'$ is an object of
$\underline\Psi_{L'_\bullet}$ and $t$ is a trivialization of its
inverse image $f^*(X')$, that is, an isomorphism
\[
 f^*(X')\simeq0_{L_\bullet,A},
\]
where $0_{L_\bullet,A}$ denotes the zero object of
$\underline\Psi_{L_\bullet}(A)$ and $A$ is the object of
$\underline C$ underlying $X'$.  We claim that \underline{the
functor} (\SGARefLink{sga7.sga7i.evii.n3-4-8-2}{3.4.8.2}) \underline{is an equivalence of categories}.
Together with the left exactness in $A$ established in \SGARefLink{sga7.sga7i.evii.n3-3-1}{3.3.1}, this
completes the exactness property in $L_\bullet$ of the expression
$\underline\Psi_{L_\bullet}(A)$.

Let us sketch the proof.  To give an object of the target of
(\SGARefLink{sga7.sga7i.evii.n3-4-8-2}{3.4.8.2}) is to give an extension
\[
 0\longrightarrow A\xrightarrow{i'}X'\xrightarrow{j'}L'_0
 \longrightarrow0
\]
of $L'_0$ by an object $A$ of $\underline C$, together with a
morphism $\alpha':L'_1\longrightarrow X'$ satisfying
$j'\alpha'=d'_0$ and $\alpha'd'_1=0$, and a morphism
$t:L_0\longrightarrow X'$ satisfying $j't=f_0$ and
$td_0=\alpha'f_1$.  These data are displayed in the following
diagram:
\begin{center}
\begin{tikzpicture}[
  x=1cm,y=1cm,>=Latex,
  every node/.style={font=\small,inner sep=1.2pt}
]
  \node (z00) at (0,0) {$0$};
  \node (l0)  at (1.3,0) {$L_0$};
  \node (lp0) at (2.8,0) {$L'_0$};
  \node (lpp0) at (4.3,0) {$L''_0$};
  \node (z05) at (5.5,0) {$0$};

  \node (z10) at (0,-1.35) {$0$};
  \node (l1)  at (1.3,-1.35) {$L_1$};
  \node (lp1) at (2.8,-1.35) {$L'_1$};
  \node (lpp1) at (4.3,-1.35) {$L''_1$};
  \node (z15) at (5.5,-1.35) {$0$};

  \node (z20) at (0,-2.70) {$0$};
  \node (l2)  at (1.3,-2.70) {$L_2$};
  \node (lp2) at (2.8,-2.70) {$L'_2$};
  \node (lpp2) at (4.3,-2.70) {$L''_2$};
  \node (z25) at (5.5,-2.70) {$0$};

  \node (xp) at (3.25,1.25) {$X'$};
  \node (ap) at (4.0,2.25) {$A$};
  \node (zap) at (4.75,3.25) {$0$};
  \node (xpp) at (5.85,1.25) {$X''$};
  \node (app) at (6.60,2.25) {$A$};
  \node (zapp) at (7.35,3.25) {$0$};

  \draw[->] (z00)--(l0);
  \draw[->] (l0)--node[above]{$f_0$}(lp0);
  \draw[->] (lp0)--node[above]{$g_0$}(lpp0);
  \draw[->] (lpp0)--(z05);
  \draw[->] (z10)--(l1);
  \draw[->] (l1)--node[above]{$f_1$}(lp1);
  \draw[->] (lp1)--node[above]{$g_1$}(lpp1);
  \draw[->] (lpp1)--(z15);
  \draw[->] (z20)--(l2);
  \draw[->] (l2)--(lp2);
  \draw[->] (lp2)--(lpp2);
  \draw[->] (lpp2)--(z25);

  \draw[->] (l1)--node[right]{$d_0$}(l0);
  \draw[->] (lp1)--node[right]{$d'_0$}(lp0);
  \draw[->] (lpp1)--node[right]{$d''_0$}(lpp0);
  \draw[->] (l2)--(l1);
  \draw[->] (lp2)--node[right]{$d'_1$}(lp1);
  \draw[->] (lpp2)--node[right]{$d''_1$}(lpp1);

  \draw[->] (zap)--(ap);
  \draw[->] (ap)--node[left]{$i'$}(xp);
  \draw[->,dashed] (zapp)--(app);
  \draw[->,dashed] (app)--node[left]{$i''$}(xpp);
  \draw[->] (l0) to[bend left=22] node[above]{$t$} (xp);
  \draw[->] (xp)--node[left]{$j'$}(lp0);
  \draw[->] (lp1)--node[pos=.82,right]{$\alpha'$}(xp);
  \draw[->,dashed] (xpp)--node[above]{$j''$}(lpp0);
  \draw[->,dashed] (lpp1) to[bend right=25]
    node[pos=.82,right]{$\alpha''$} (xpp);

  \node[anchor=west] at (7.55,-0.70) {$
    \left\{
    \begin{aligned}
      j'\alpha'&=d'_0,\\
      \alpha'd'_1&=0,\\
      j't&=f_0,\\
      td_0&=\alpha'f_1.
    \end{aligned}
    \right.$};
\end{tikzpicture}
\end{center}

The familiar left exactness, in the first argument, of extensions
of an object $L$ by an object $A$, applied to the first row of the
diagram, shows that a morphism $t$ satisfying $j't=f_0$ allows one
to interpret $X'$ as the inverse image of an extension $X''$ of
$L''_0$ by $A$.  Under this interpretation, $t$ corresponds to the
natural trivialization of the inverse image of that extension along
$g_0f_0=0$.

The morphism $\alpha'$ satisfying $j'\alpha'=d'_0$ and
$\alpha'f_1=td_0$ is interpreted by essentially the same dictionary,
now applied to the second row of the diagram, as arising from a
morphism $\alpha'':L''_1\longrightarrow X''$ satisfying
$j''\alpha''=d''_0$.  Here $\alpha'$ and $\alpha''$ are regarded as
trivializations of the inverse images of $X'$ and $X''$ along
$d'_0$ and $d''_0$, respectively.  Finally, the relation
$\alpha'd'_1=0$ is then equivalent to $\alpha''d''_1=0$.  Thus the
data $(X',\alpha',t)$, an object of
$\underline\Psi_{[L_\bullet\longrightarrow L'_\bullet]}$, are
equivalent to the data $(X'',\alpha'')$, an object of
$\underline\Psi_{L''_\bullet}$.  This proves, in substance, that
(\SGARefLink{sga7.sga7i.evii.n3-4-8-2}{3.4.8.2}) is indeed an equivalence of categories.

\SGARefTarget{sga7.sga7i.evii.n3-4-9}{3.4.9}. Continue to consider the exact sequence (\SGARefLink{sga7.sga7i.evii.n3-4-8-1}{3.4.8.1}), and the
exact sequence of $\Ext^i$ groups that it defines:
\begin{multline}\tag{\SGARefTarget{sga7.sga7i.evii.n3-4-9-1}{3.4.9.1}}
 0\longrightarrow\Ext^0(L''_\bullet,A)
 \xrightarrow{g^0}\Ext^0(L'_\bullet,A)
 \xrightarrow{f^0}\Ext^0(L_\bullet,A)
 \xrightarrow{\partial^0}\Ext^1(L''_\bullet,A)\\
 \xrightarrow{g^1}\Ext^1(L'_\bullet,A)
 \xrightarrow{f^1}\Ext^1(L_\bullet,A)\longrightarrow\cdots.
\end{multline}
The canonical isomorphisms of the forms (\SGARefLink{sga7.sga7i.evii.n3-1-8}{3.1.8}) and (\SGARefLink{sga7.sga7i.evii.n3-2-4}{3.2.4}) give a
geometric interpretation of the first six terms of this exact
sequence in terms of the categories $\underline\Psi_{L_\bullet}$,
$\underline\Psi_{L'_\bullet}$, and
$\underline\Psi_{L''_\bullet}$.  Compatibility (\SGARefLink{sga7.sga7i.evii.n3-4-6}{3.4.6}) likewise
gives a geometric interpretation of the arrows
$f^0,g^0,f^1,g^1$.  We shall now give such an interpretation of
$\partial^0$, using the result just proved that (\SGARefLink{sga7.sga7i.evii.n3-4-8-2}{3.4.8.2}) is an
equivalence of categories.

Let
\[
 u\in\Ext^0(L_\bullet,A),
 \qquad\text{i.e.}\qquad
 u:L_0\longrightarrow A,\quad ud_0=0.
\]
Consider the object $(0,\overline u)$ of
$\underline\Psi_{[L_\bullet\to L'_\bullet]}$.  Here $0$ is the
zero object of $\underline\Psi_{L'_\bullet}$, whose inverse image
along $f_0$ is therefore identified with the zero object of
$\underline\Psi_{L_\bullet}$, and $\overline u$ is the automorphism
of the latter defined by $u$ through the canonical isomorphism
(\SGARefLink{sga7.sga7i.evii.n3-1-8}{3.1.8}).  Since (\SGARefLink{sga7.sga7i.evii.n3-4-8-2}{3.4.8.2}) is an equivalence, $(0,\overline u)$ comes
from an object of $\underline\Psi_{L''_\bullet}(A)$, determined up
to unique isomorphism.  Its isomorphism class is an element of
\[
 \Psi^1_{L''_\bullet}(A)\simeq\Ext^1(L''_\bullet,A).
\]
\underline{We claim that this element is precisely}
$\partial^0(u)$.

We prove this compatibility by reducing it to the analogous one in
[4, Cor.~5.5], where the sign must be corrected as indicated in
3.3.6.3.  The argument also gives another proof of the equivalence
(\SGARefLink{sga7.sga7i.evii.n3-4-8-2}{3.4.8.2}), by reducing it to [4, Prop.~\SGARefLink{sga7.sga7i.evii.n2-7}{2.7}].  Consider the opposite
category $\underline C^{\mathrm o}$ and interpret the construction
of $\underline\Psi_{L_\bullet}(A)$, for fixed $A$, together with
the constructions of \SGARefLink{sga7.sga7i.evii.n3-4-8}{3.4.8} and the present section, in terms of
$\underline C^{\mathrm o}$ and the category $\underline E$ cofibred
over $\underline C^{\mathrm o}$ whose objects are extensions of the
fixed object $A$ of $\underline C^{\mathrm o}$ by a variable object
$L$ of $\underline C^{\mathrm o}$.  This is an additive cofibred
category over $\underline C^{\mathrm o}$, and is even left exact
there (in the sense of [4, \SGARefLink{sga7.sga7i.evii.n2-2}{2.2}]).

A chain complex $L_\bullet$ of $\underline C$ may be interpreted as
a cochain complex of $\underline C^{\mathrm o}$.  Under this
interpretation, the construction of $\underline\Psi_{L_\bullet}(A)$
coincides with the opposite of the category
$\underline{\widehat E}(L_\bullet)$ in the sense of [4, 1.8].  The
functor (\SGARefLink{sga7.sga7i.evii.n3-4-8-2}{3.4.8.2}) is then identified with the opposite of the functor
of the form [4, (2.1.2)], with $E$ replaced by its extension
$\underline{\widehat E}$ to the category of cochain complexes.
Consequently, the assertion that this functor is an equivalence is
the special case [4, \SGARefLink{sga7.sga7i.evii.n2-7}{2.7}, a)$\Rightarrow$d)].

On the other hand, the exact sequence (\SGARefLink{sga7.sga7i.evii.n3-4-9-1}{3.4.9.1}) is identified with
the exact sequence of the groups $\Ext^i(A,L''_\bullet)$,
$\Ext^i(A,L'_\bullet)$, and $\Ext^i(A,L_\bullet)$, apart from a
factor $(-1)^i$ on $\partial^i$ (Verdier dixit).  Similarly, the
isomorphisms (\SGARefLink{sga7.sga7i.evii.n3-1-8}{3.1.8}) and (\SGARefLink{sga7.sga7i.evii.n3-2-4}{3.2.4}) are identified with the analogous
isomorphisms (3.11.5) and (3.11.4).  The asserted compatibility is
therefore the special case of [4, Cor.~5.5] corrected as above.

\underline{Remark} \SGARefTarget{sga7.sga7i.evii.n3-4-9-2}{3.4.9.2}. The editor, being very afraid of signs,
confesses that he does not have complete confidence in the preceding
proof: he fears that signs may have been omitted in passing from
$C$ to $C^{\mathrm o}$.  He recommends that a courageous reader
verify the asserted compatibility directly (or, respectively, its
opposite).  Fortunately, these sign questions will play no role in
the following expos\'es.

% END INCLUDED SGA7I_English_current_buildable_source_20260823/components/88_expose_VII_sections_3_4_8_and_3_4_9.tex

% BEGIN INCLUDED SGA7I_English_current_buildable_source_20260823/components/89_expose_VII_section_3_5.tex
% SOURCE AUTHORITY:
% expose_VII_body.tex lines 1545--1662; scan indices 209--215; source folios 198--204.
% Direct scan comparison: physical PDF pages 210--216 at 600 dpi.
% Transparent corrections: the first overbrace in (3.5.1) is L_2 as printed, not
% the working transcription's L_o; chain subscripts o are rendered as zero; the
% chain-complex relation is written in compositional order d_0d_1=0 throughout;
% and the display used later as (*) receives the tag visible in the scan.
% Hard stop before source line 1668, section 3.6.

\subsection*{\SGARefTarget{sga7.sga7i.evii.n3-5}{3.5}. A canonical partial flat resolution of an
arbitrary $\mathbb Z$-module}

We apply the preceding results when $\underline C$ is the category
of commutative groups of the topos $\underline T$, that is, the
category of $\mathbb Z$-modules of $\underline T$.  Let $P$ be an
object of $\underline C$.  We shall define, functorially in $P$, a
chain complex $L_\bullet(P)$ augmented to $P$:
\begin{equation}\tag{\SGARefTarget{sga7.sga7i.evii.n3-5-1}{3.5.1}}
\begin{tikzcd}[column sep=small,row sep=large]
 0\arrow[r]&\overbrace{\mathbb Z[P\times P]\times
   \mathbb Z[P\times P\times P]}^{L_2}
   \arrow[r,"d_1"]
 &\overbrace{\mathbb Z[P\times P]}^{L_1}
   \arrow[r,"d_0"]
 &\overbrace{\mathbb Z[P]}^{L_0}
   \arrow[r]\arrow[d,"\varepsilon"']&0\\
 &&&P&{}
\end{tikzcd}
\end{equation}
with $L_i(P)=0$ for $i\geq3$.  The components are displayed in
(\SGARefLink{sga7.sga7i.evii.n3-5-1}{3.5.1}), using the notation of \SGARefLink{sga7.sga7i.evii.n1-4}{1.4}.  To describe the differentials
$d_0,d_1$ and the augmentation $\varepsilon$, write $[p]$ for the
section of $\mathbb Z[P](S)$ defined by $p\in P(S)$, and similarly
write $[p,q]$ and $[p,q,r]$ for the corresponding sections of
$\mathbb Z[P\times P](S)$ and
$\mathbb Z[P\times P\times P](S)$.  Then set
\begin{equation}\tag{\SGARefTarget{sga7.sga7i.evii.n3-5-2}{3.5.2}}
\left\{
\begin{aligned}
 \varepsilon[p]&=p,\\
 d_0[p,q]&=[p+q]-[p]-[q],\\
 d_1[p,q]&=[p,q]-[q,p],\\
 d_1[p,q,r]&=[p+q,r]-[p,q+r]+[p,q]-[q,r].
\end{aligned}
\right.
\end{equation}
These homomorphisms of commutative groups do define a complex
$L_\bullet(P)$ augmented to $P$:
\[
 \varepsilon d_0=0,\qquad d_0d_1=0.
\]
The first relation is immediate.  The second splits into two
relations according to the two factors of $L_2(P)$; they express,
respectively, commutativity and associativity of the composition law
on $P$.

The augmented complex depends functorially on $P$.  A morphism
$u:P\longrightarrow P'$ gives a commutative diagram
\[
\begin{tikzcd}[column sep=huge,row sep=large]
 L_\bullet(P)\arrow[r,"L_\bullet(u)"]\arrow[d]
 &L_\bullet(P')\arrow[d]\\
 P\arrow[r,"u"]&P'.
\end{tikzcd}
\]
Its components are flat, as is every free $\mathbb Z$-module.

\underline{Proposition} \SGARefTarget{sga7.sga7i.evii.n3-5-3}{3.5.3}. \underline{One has}
$H_1(L_\bullet(P))=0$, \underline{and the augmentation}
$\varepsilon$ \underline{induces an isomorphism}
\[
 H_0(L_\bullet(P))\xrightarrow{\sim}P.
\]

For the proof one can make the standard reduction to the case in
which $\underline T$ is the category of sets and $P$ is an ordinary
$\mathbb Z$-module, where (\SGARefLink{sga7.sga7i.evii.n3-5-1}{3.5.1}) is likely familiar.  We instead
give a direct proof that does not use this reduction, applying \SGARefLink{sga7.sga7i.evii.n3-4-7}{3.4.7}
to the augmentation
\[
 \varepsilon:L_\bullet(P)\longrightarrow P,
\]
where $P$ is regarded, as usual, as a chain complex concentrated in
degree zero.  It remains to prove that, for every object $G$ of
$\underline C$, the functor
\begin{equation}\tag{*}
 \varepsilon^*:\underline\Psi_P(G)\longrightarrow
 \underline\Psi_{L_\bullet(P)}(G)
\end{equation}
is an equivalence of categories.

The source of (*) is plainly isomorphic to the category of
extensions of the $\mathbb Z$-module $P$ by $G$.  To describe its
target, use \SGARefLink{sga7.sga7i.evii.n1-4}{1.4} to identify the categories of extensions of
$L_0,L_1,L_2$ by $G$, respectively, with the category of torsors
under $G_P$, the category of torsors under $G_{P\times P}$, and the
category of pairs consisting of a torsor under $G_{P\times P}$ and
a torsor under $G_{P\times P\times P}$.

Thus an object of the target of (*) is a pair $(E,\varphi)$, where
$E$ is a torsor under $G_P$ and $\varphi$ is a trivialization of the
torsor $E'$ on $P\times P$ corresponding to the inverse image along
$d_0$ of the extension of $\mathbb Z[P]$ by $G$ defined by $E$.
The trivialization is subject to the compatibility with
$d_0d_1=0$ described below.  Formula (\SGARefLink{sga7.sga7i.evii.n3-5-2}{3.5.2}) immediately gives
\[
 E'=\pi^*E\,\pr_1^*(E)^{-1}\,\pr_2^*(E)^{-1},
\]
and a trivialization of this torsor is the same as an isomorphism,
also denoted by $\varphi$,
\[
 \varphi:\pr_1^*E\,\pr_2^*E\xrightarrow{\sim}\pi^*E.
\]
Expanding the compatibility of $\varphi$ with $d_0d_1=0$ gives the
commutativity condition (\SGARefLink{sga7.sga7i.evii.n1-2-1}{1.2.1}) and the associativity condition
(\SGARefLink{sga7.sga7i.evii.n1-1-4-1}{1.1.4.1}), corresponding to the two factors of $L_2$.  The assertion
that (*) is an equivalence is therefore precisely the dictionary of
\SGARefLink{sga7.sga7i.evii.n1-1-6}{1.1.6}.  This proves \SGARefLink{sga7.sga7i.evii.n3-5-3}{3.5.3}.

\underline{Remark} \SGARefTarget{sga7.sga7i.evii.n3-5-4}{3.5.4}. a) It would be interesting to strengthen
\SGARefLink{sga7.sga7i.evii.n3-5-3}{3.5.3} by constructing a resolution $L_\bullet(P)$ of $P$, functorial
in $P$, whose components are functorially isomorphic to sums of free
$\mathbb Z$-modules generated by cartesian powers of $P$.  In view
of isomorphisms of the form (\SGARefLink{sga7.sga7i.evii.n1-4-5}{1.4.5}), this would give a method for
partially computing $\Ext^i(P,A)$ as the abutment of a spectral
sequence whose initial term is expressed through groups
\[
 H^j(P\times P\times\cdots\times P,A).
\]
Here we obtain only very partial information about $\Ext^1$, and no
information about the higher $\Ext^i$.  There is, in any case, a
natural way to continue the partial resolution $L_\bullet(P)$ by one
degree.\footnote{L. Breen constructed such a resolution (editor's
note).}

b) There is a variant of (\SGARefLink{sga7.sga7i.evii.n3-5-1}{3.5.1}) and \SGARefLink{sga7.sga7i.evii.n3-5-3}{3.5.3} in the category of
$A$-modules, where $A$ is a fixed ring of $\underline T$.  The form
of $L_\bullet(P)$ is more complicated because the description \SGARefLink{sga7.sga7i.evii.n1-1-6}{1.1.6}
must be enlarged (compare \SGARefLink{sga7.sga7i.evii.n2-10-3}{2.10.3}).  One obtains
\begin{equation}\tag{\SGARefTarget{sga7.sga7i.evii.n3-5-5}{3.5.5}}
\left\{
\begin{aligned}
 L_0={}&A[P],\\
 L_1={}&A[P\times P]+A[A\times P],\\
 L_2={}&A[P\times P\times P]+A[P\times P]
       +A[A\times A\times P]+A[A\times P\times P].
\end{aligned}
\right.
\end{equation}
The augmented-complex structure is given by (\SGARefLink{sga7.sga7i.evii.n3-5-2}{3.5.2}), supplemented by
\begin{equation}\tag{\SGARefTarget{sga7.sga7i.evii.n3-5-6}{3.5.6}}
\left\{
\begin{aligned}
 d_0[\lambda,p]={}&[\lambda p]-\lambda[p],\\
 d_1[\lambda,\mu,p]={}&[\lambda,\mu p]-[\lambda\mu,p]
                         +\lambda[\mu,p],\\
 d_1[\lambda,p,q]={}&[\lambda,p+q]-[\lambda,p]-[\lambda,q]
                       +\lambda[p,q]-[\lambda p,\lambda q].
\end{aligned}
\right.
\end{equation}
The relation $d_0d_1=0$ follows from the analogous relation on the
components appearing in (\SGARefLink{sga7.sga7i.evii.n3-5-2}{3.5.2}), which, as we saw, is equivalent to
associativity and commutativity of addition on $P$, together with
the relations
\[
 d_0d_1[\lambda,\mu,p]=0,
 \qquad d_0d_1[\lambda,p,q]=0.
\]
The latter express, respectively, associativity of the given action
of $A$ on $P$,
\[
 \lambda(\mu p)=(\lambda\mu)p,
\]
and its distributivity over addition,
\[
 \lambda(p+q)=\lambda p+\lambda q.
\]

For a torsor $E$ under $G_P$, a trivialization of $d_0^*E$ amounts
to an internal composition law on $E$ (satisfying the compatibility
conditions already described for the internal law on $E$ and its
torsor structure), together with an external composition law of $A$
on $E$, compatible with the action of $A$ on $P$ and on $G$.  The
compatibility of this trivialization with the evident trivialization
of $d_1^*d_0^*E$ is equivalent to the conditions already described
above, which say that the internal law on $E$ is a commutative-group
law extending $P$ by $G$, together with associativity of the external
law and distributivity of the latter with respect to the internal
law.  These conditions say exactly that $E$, with its two
composition laws and its torsor structure, is an $A$-module extension
of the $A$-modules $P$ and $G$.

It follows, by essentially the same method as in \SGARefLink{sga7.sga7i.evii.n3-5-3}{3.5.3}, that
$L_\bullet$ is a partial resolution of the module $P$, that is,
$H_1(L_\bullet)=0$.  It would again be interesting to extend this
partial resolution to an infinite resolution of the same kind, with
components of the form
\[
 A[A\times A\times\cdots\times P\times P\times\cdots].
\]
Compare the analogous question in a) for commutative groups.

% END INCLUDED SGA7I_English_current_buildable_source_20260823/components/89_expose_VII_section_3_5.tex

% BEGIN INCLUDED SGA7I_English_current_buildable_source_20260823/components/90_expose_VII_section_3_6.tex
% SOURCE AUTHORITY:
% expose_VII_body.tex lines 1668--1726; scan indices 215--219; source folios 204--208.
% Direct scan comparison: physical PDF pages 216--220 at 600 dpi.
% Transparent corrections: (3.6.3) and Theorem 3.6.4 use L'_\bullet(P,Q), not
% the impossible L_\bullet(P',Q'); extension notation uses L_0(P,Q); the induced
% maps from P x Q x Q have target P x Q; and the extension of P_Q is by G_Q.
% Typewritten letter-o subscripts are rendered as zero.
% Hard stop before source line 1729, section 3.7.

\subsection*{\SGARefTarget{sga7.sga7i.evii.n3-6}{3.6}. Application to the description of the cofibred
category of biextensions of a given pair $(P,Q)$ of commutative
groups}

Let $P$ and $Q$ be two $\mathbb Z$-modules of the topos
$\underline T$.  Choose a flat resolution $L'_\bullet(P)$ of $P$
that agrees in degrees at most $2$ with the augmented complex
$L_\bullet(P)$ of \SGARefLink{sga7.sga7i.evii.n3-5}{3.5}.  This is possible by \SGARefLink{sga7.sga7i.evii.n3-5-3}{3.5.3} and because every
$\mathbb Z$-module on $\underline T$ is a quotient of a free, hence
flat, $\mathbb Z$-module.  Similarly choose a flat resolution
$L'_\bullet(Q)$ of $Q$ agreeing with $L_\bullet(Q)$ in degrees at
most $2$.  We have canonical homomorphisms
\begin{equation}\tag{\SGARefTarget{sga7.sga7i.evii.n3-6-1}{3.6.1}}
 L_\bullet(P)\longrightarrow L'_\bullet(P),
 \qquad
 L_\bullet(Q)\longrightarrow L'_\bullet(Q),
\end{equation}
and hence a canonical homomorphism
\begin{equation}\tag{\SGARefTarget{sga7.sga7i.evii.n3-6-2}{3.6.2}}
\begin{split}
 L_\bullet(P,Q)
 &\overset{\mathrm{def}}{=}L_\bullet(P)\otimes L_\bullet(Q)
 \longrightarrow L'_\bullet(P,Q)\\
 &\overset{\mathrm{def}}{=}L'_\bullet(P)\otimes L'_\bullet(Q)
 \simeq P\mathbin{\overset{\mathbb L}{\otimes}}Q.
\end{split}
\end{equation}
It is an isomorphism on components in degrees at most $2$.
Consequently, by definition,
\begin{equation}\tag{\SGARefTarget{sga7.sga7i.evii.n3-6-3}{3.6.3}}
 \underline\Psi_{L_\bullet(P,Q)}
 =
 \underline\Psi_{L'_\bullet(P,Q)}.
\end{equation}

\underline{Theorem} \SGARefTarget{sga7.sga7i.evii.n3-6-4}{3.6.4}. \underline{The additive cofibred category
over the category}
\[
 \underline C=\operatorname{GRAB}(\underline T)
\]
\underline{of commutative groups of $\underline T$, consisting of
biextensions of $(P,Q)$ by a variable abelian group $G$} (\SGARefLink{sga7.sga7i.evii.n2-4}{2.4}),
\underline{is canonically equivalent to the additive cofibred
category}
\[
 \underline\Psi_{L_\bullet(P,Q)}
 =
 \underline\Psi_{L'_\bullet(P,Q)}
\]
\underline{defined from}
$L_\bullet(P,Q)=L_\bullet(P)\otimes L_\bullet(Q)$
\underline{as in \SGARefLink{sga7.sga7i.evii.n3-1}{3.1}}.

This and \SGARefLink{sga7.sga7i.evii.n3-2-5}{3.2.5} give the principal result of the present section,
announced above:

\underline{Corollary} \SGARefTarget{sga7.sga7i.evii.n3-6-5}{3.6.5}. \underline{For fixed $P$ and $Q$ there
is a canonical isomorphism, functorial in $G$,}
\begin{equation}\tag{\SGARefTarget{sga7.sga7i.evii.n3-6-5-1}{3.6.5.1}}
 \Biext^1(P,Q;G)
 \xrightarrow{\sim}
 \Ext^1(P\mathbin{\overset{\mathbb L}{\otimes}}Q,G).
\end{equation}

To prove \SGARefLink{sga7.sga7i.evii.n3-6-5}{3.6.5}, first make explicit the components of degrees
$0,1,2$ of $L_\bullet(P,Q)$:
\begin{equation}\tag{\SGARefTarget{sga7.sga7i.evii.n3-6-6}{3.6.6}}
\left\{
\begin{aligned}
 L_0(P,Q)
  ={}&L_0(P)\otimes L_0(Q)
     =\mathbb Z[P\times Q],\\
 L_1(P,Q)
  ={}&L_0(P)\otimes L_1(Q)+L_1(P)\otimes L_0(Q)\\
  ={}&\mathbb Z[P\times Q\times Q]
     +\mathbb Z[P\times P\times Q],\\
 L_2(P,Q)
  ={}&L_0(P)\otimes L_2(Q)+L_1(P)\otimes L_1(Q)
     +L_2(P)\otimes L_0(Q)\\
  ={}&\mathbb Z[P\times Q\times Q]
     +\mathbb Z[P\times Q\times Q\times Q]\\
   &+\mathbb Z[P\times P\times Q\times Q]\\
   &+\mathbb Z[P\times P\times Q]
     +\mathbb Z[P\times P\times P\times Q].
\end{aligned}
\right.
\end{equation}

To describe the cofibred category
$\underline\Psi_{L_\bullet(P,Q)}$ through its fibres
$\underline\Psi_{L_\bullet(P,Q)}(G)$, use \SGARefLink{sga7.sga7i.evii.n1-4}{1.4} to express extensions
of $L_i(P,Q)$, for $i=0,1,2$, by $G$ in terms of torsors.  An
extension of $L_0(P,Q)$ by $G$ is a torsor under $G_{P\times Q}$.
An extension of $L_1(P,Q)$ by $G$ is a pair consisting of a torsor
under $G_{P\times Q\times Q}$ and one under
$G_{P\times P\times Q}$.  An extension of $L_2(P,Q)$ by $G$ is
similarly a system of five torsors, under the groups obtained from
$G$ by base change to
\[
\begin{gathered}
 P\times Q\times Q,\quad
 P\times Q\times Q\times Q,\quad
 P\times P\times Q\times Q,\\
 P\times P\times Q,\quad
 P\times P\times P\times Q.
\end{gathered}
\]

Under this dictionary, a pair $(X,\alpha)$ consisting of an
extension of $L_0(P,Q)$ by $G$ and a trivialization of
$d_0^*X$ on $L_1(P,Q)$ becomes a torsor $E$ on $P\times Q$ equipped
with two trivializations $(\varphi,\psi)$ of the torsors
$(E_1,E_2)$ on $P\times Q\times Q$ and
$P\times P\times Q$, respectively, obtained from $E$ by inverse
image along
\[
 d_0:L_1(P,Q)\longrightarrow L_0(P,Q).
\]
Expanding $d_0$, as in the proof of \SGARefLink{sga7.sga7i.evii.n3-5-3}{3.5.3}, identifies
$\varphi$ and $\psi$ with isomorphisms
\[
 \varphi:q_1^*E\,q_2^*E\xrightarrow{\sim}\pi_Q^*E,
 \qquad
 \psi:p_1^*E\,p_2^*E\xrightarrow{\sim}\pi_P^*E.
\]
Here $\pi_Q,q_1,q_2:P\times Q\times Q\longrightarrow P\times Q$
are induced by the composition law and the first and second
projections $Q\times Q\longrightarrow Q$; the maps
\[
 \pi_P,p_1,p_2:P\times P\times Q\longrightarrow P\times Q
\]
are defined symmetrically.  Thus $\varphi$ and $\psi$ are data of
the forms (\SGARefLink{sga7.sga7i.evii.n2-0-4}{2.0.4}) and (\SGARefLink{sga7.sga7i.evii.n2-0-7}{2.0.7}) appearing in \SGARefLink{sga7.sga7i.evii.n2-1}{2.1}.

Finally, the condition that $(X,\alpha)$ be an object of
$\underline\Psi_{L_\bullet(P,Q)}(G)$, namely compatibility of the
partial trivialization $\alpha$ with $d_0d_1=0$, becomes five
compatibility relations for the structure $(E,\varphi,\psi)$ on
torsors over, respectively,
\[
\begin{gathered}
 P\times Q\times Q,\quad
 P\times Q\times Q\times Q,\quad
 P\times P\times Q\times Q,\\
 P\times P\times Q,\quad
 P\times P\times P\times Q.
\end{gathered}
\]
They correspond to the five terms in the expression (\SGARefLink{sga7.sga7i.evii.n3-6-6}{3.6.6}) for
$L_2(P,Q)$.  As in the proof of \SGARefLink{sga7.sga7i.evii.n3-5-3}{3.5.3}, the first two are precisely
the commutativity relation (\SGARefLink{sga7.sga7i.evii.n2-0-6}{2.0.6}) and the associativity relation
(\SGARefLink{sga7.sga7i.evii.n2-0-5}{2.0.5}); together they say that $\varphi$ gives $E$ the structure
of a commutative extension of $Q_P$ by $G_P$.  Symmetrically, the
last two say that $\psi$ gives $E$ the structure of a commutative
extension of $P_Q$ by $G_Q$.  The middle relation, on torsors over
$P\times P\times Q\times Q$, is exactly the commutativity of
diagram (\SGARefLink{sga7.sga7i.evii.n2-1-1}{2.1.1}).  It says that $\varphi$ and $\psi$ are compatible,
hence, in view of the other conditions, define on $E$ a
biextension structure of $(P,Q)$ by $G$.

The verification of these interpretations is immediate and is left
to the reader.  The preceding discussion indeed defines an
equivalence of categories cofibred over
$\underline C=\operatorname{GRAB}(\underline T)$, as asserted in
\SGARefLink{sga7.sga7i.evii.n3-6-4}{3.6.4}.

\underline{Remark} \SGARefTarget{sga7.sga7i.evii.n3-6-7}{3.6.7}. Each of the three generalizations of
biextensions mentioned in \SGARefLink{sga7.sga7i.evii.n2-10}{2.10} has a corresponding variant of \SGARefLink{sga7.sga7i.evii.n3-6-4}{3.6.4}
and hence of \SGARefLink{sga7.sga7i.evii.n3-6-5}{3.6.5}.  The generalization to any number of factors
(\SGARefLink{sga7.sga7i.evii.n2-10-2}{2.10.2}) is immediate, as is the generalization to biextensions of
$A$-modules (\SGARefLink{sga7.sga7i.evii.n2-10-3}{2.10.3}), once the variant of $L_\bullet(P)$ described
in \SGARefLink{sga7.sga7i.evii.n3-5-4}{3.5.4} b) has been constructed.  From the viewpoint of
homological algebra, the case in which $P$ and $Q$ are not assumed
commutative (\SGARefLink{sga7.sga7i.evii.n2-10-1}{2.10.1}) is more interesting.

Associate with a not necessarily commutative group $P$ of
$\underline T$ a complex $K_\bullet(P)$ truncated in degree $2$.
It is described as $L_\bullet(P)$ was in \SGARefLink{sga7.sga7i.evii.n3-5}{3.5}, except that
\[
 K_2(P)=\mathbb Z[P\times P\times P],
\]
the first factor occurring in $L_2(P)$ being omitted; $d_0$ and
$d_1$ are defined as in (\SGARefLink{sga7.sga7i.evii.n3-5-2}{3.5.2}).  Then
\[
 H_0(K_\bullet(P))=P/[P,P]
 \simeq\underline H_1(P,\mathbb Z),
\]
and similarly
\[
 H_1(K_\bullet(P))\simeq\underline H_2(P,\mathbb Z).
\]
Indeed, $K_\bullet(P)$ is precisely the truncation in degree $2$ of
the chain complex $C_\bullet(P,\mathbb Z)$ of $P$ with coefficients
in $\mathbb Z$.

The argument proving \SGARefLink{sga7.sga7i.evii.n3-5-3}{3.5.3} also shows that the cofibred category
$\underline\Psi_{K_\bullet(P)}$ over
$\underline C=\operatorname{GRAB}(\underline T)$ is canonically
equivalent to the category of central extensions of $P$ by variable
commutative groups $G$.  Likewise, for two groups $P,Q$, the
cofibred category of biextensions of $(P,Q)$ in the sense of \SGARefLink{sga7.sga7i.evii.n2-10-1}{2.10.1}
is equivalent to $\underline\Psi_{K_\bullet(P,Q)}$, where
\[
 K_\bullet(P,Q)
 =
 K_\bullet(P)\mathbin{\overset{\mathbb L}{\otimes}}K_\bullet(Q).
\]

These considerations suggest that, when $P$ is commutative, the
canonical resolution sought in \SGARefLink{sga7.sga7i.evii.n3-5-4}{3.5.4} a) might be obtained by a
suitable modification of $C_\bullet(P,\mathbb Z)$.

% END INCLUDED SGA7I_English_current_buildable_source_20260823/components/90_expose_VII_section_3_6.tex

% BEGIN INCLUDED SGA7I_English_current_buildable_source_20260823/components/91_expose_VII_section_3_7.tex
% SOURCE AUTHORITY:
% expose_VII_body.tex lines 1729--1804; scan indices 219--223; source folios 208--212.
% Direct scan comparison: physical PDF pages 220--224 at 600 dpi.
% Transparent normalization: typewritten letter-o superscripts are rendered as zero;
% the third factor of the base category in Proposition 3.7.6 is covariant, exactly as
% printed; and the chain-complex subscripts use zero.  Hard stop before source line
% 1810, section 3.8.

\subsection*{\SGARefTarget{sga7.sga7i.evii.n3-7}{3.7}. Various compatibilities}

In this section we examine the behaviour, for variable $P$ and $Q$,
of the equivalence of cofibred categories in \SGARefLink{sga7.sga7i.evii.n3-6-4}{3.6.4}.  The variance of
the cofibred category of biextensions of $(P,Q)$ by a variable $G$
with respect to variable $P,Q$ is defined in principle in \SGARefLink{sga7.sga7i.evii.n2-4}{2.4}; it may
also be expressed by saying that one has a cofibred functor (\SGARefLink{sga7.sga7i.evii.n2-6-1}{2.6.1}).
Notice that the variance in $P,Q$ is \underline{contravariant}.  The
variance of the cofibred category
$\underline\Psi_{L_\bullet(P,Q)}$ with respect to variable $P,Q$ is
also clear, since
$L_\bullet(P,Q)=L_\bullet(P)\otimes L_\bullet(Q)$ depends
functorially on the pair $(P,Q)$, while $\underline\Psi_{L_\bullet}$
depends contravariantly on $L_\bullet$ (\SGARefLink{sga7.sga7i.evii.n3-4}{3.4}).  With this understood,
a first compatibility is expressed, for every triple of morphisms
\begin{equation}\tag{\SGARefTarget{sga7.sga7i.evii.n3-7-1}{3.7.1}}
 P'\longrightarrow P,\qquad
 Q'\longrightarrow Q,\qquad
 G\longrightarrow G',
\end{equation}
by the fact that the diagram of functors
\begin{equation}\tag{\SGARefTarget{sga7.sga7i.evii.n3-7-2}{3.7.2}}
\begin{tikzcd}[column sep=huge,row sep=large]
 \underline\Psi_{L_\bullet(P,Q)}(G)
   \arrow[r]\arrow[d]
 &\underline\Psi_{L_\bullet(P',Q')}(G')
   \arrow[d]\\
 \operatorname{BIEXT}(P,Q;G)
   \arrow[r]
 &\operatorname{BIEXT}(P',Q';G')
\end{tikzcd}
\end{equation}
is commutative up to a canonical isomorphism.  We shall admit this
compatibility: making the commutativity isomorphism explicit could
only be as obvious as it is tedious (as, indeed, could the complete
explication it presupposes of the equivalence of cofibred categories
in \SGARefLink{sga7.sga7i.evii.n3-6-4}{3.6.4}).

\SGARefTarget{sga7.sga7i.evii.n3-7-3}{3.7.3}. It follows in particular from the preceding compatibility
that the isomorphism (\SGARefLink{sga7.sga7i.evii.n3-6-5-1}{3.6.5.1}) is functorial not only in $G$, but also
in $P$ and $Q$.  One may state, similarly and more trivially, the
analogous functoriality of the isomorphism defined through (\SGARefLink{sga7.sga7i.evii.n3-1-8}{3.1.8}),
just as (\SGARefLink{sga7.sga7i.evii.n3-6-5-1}{3.6.5.1}) was defined through (\SGARefLink{sga7.sga7i.evii.n3-2-4}{3.2.4}):
\begin{equation}\tag{\SGARefTarget{sga7.sga7i.evii.n3-7-4}{3.7.4}}
 \Biext^0(P,Q;G)
 \xrightarrow{\sim}
 \Ext^0(P\mathbin{\overset{\mathbb L}{\otimes}}Q,G)
 \simeq\Hom(P\otimes Q,G).
\end{equation}
But such an isomorphism was already obtained in (\SGARefLink{sga7.sga7i.evii.n2-5-4-1}{2.5.4.1}), and was
observed in \SGARefLink{sga7.sga7i.evii.n2-6}{2.6} to be functorial in its three arguments.  It is
immediate to verify that the two isomorphisms (\SGARefLink{sga7.sga7i.evii.n3-7-4}{3.7.4}) and (\SGARefLink{sga7.sga7i.evii.n2-5-4-1}{2.5.4.1})
are equal.

\SGARefTarget{sga7.sga7i.evii.n3-7-5}{3.7.5}. The value of the isomorphism in \SGARefLink{sga7.sga7i.evii.n3-6-5}{3.6.5} is that it expresses the
remarkable group of geometric origin $\Biext^1(P,Q;G)$ in terms of
the very simple homological-algebra invariant
$\Ext^1(P\mathbin{\overset{\mathbb L}{\otimes}}Q,G)$.  The latter
belongs to the sequence of invariants
$\Ext^i(P\mathbin{\overset{\mathbb L}{\otimes}}Q,G)$, to which the
usual formalism of homological algebra applies.  In particular, the
sequence of these $\Ext^i$ is a \underline{cohomological functor} in
the sense of [3] (or a \underline{connected sequence of functors} in
the terminology of [5]) in each of the three arguments $P,Q,G$.
Thus an exact sequence
\begin{equation}\tag{\SGARefTarget{sga7.sga7i.evii.n3-7-5-1}{3.7.5.1}}
 0\longrightarrow G'\longrightarrow G\longrightarrow G''
 \longrightarrow0
\end{equation}
(respectively,
\begin{equation}\tag{\SGARefTarget{sga7.sga7i.evii.n3-7-5-2}{3.7.5.2}}
 0\longrightarrow P'\longrightarrow P\longrightarrow P''
 \longrightarrow0,\qquad\text{or}\qquad
 0\longrightarrow Q'\longrightarrow Q\longrightarrow Q''
 \longrightarrow0)
\end{equation}
gives rise to an infinite exact sequence of the $\Ext^i$.  Its first
six terms, which involve only $\Ext^0$ and $\Ext^1$, may be
interpreted in terms of biextensions by \SGARefLink{sga7.sga7i.evii.n3-6-5}{3.6.5} and \SGARefLink{sga7.sga7i.evii.n3-7-4}{3.7.4}.  We have
already indicated in \SGARefLink{sga7.sga7i.evii.n3-3-6}{3.3.6} a geometric interpretation of the exact
sequence deduced from (\SGARefLink{sga7.sga7i.evii.n3-7-5-1}{3.7.5.1}), and especially of the boundary map
\begin{equation}\tag{\SGARefTarget{sga7.sga7i.evii.n3-7-5-3}{3.7.5.3}}
 \partial:\Biext^0(P,Q;G'')\longrightarrow
 \Biext^1(P,Q;G'),
\end{equation}
within the framework of a \underline{left exact cofibred category}.
One may seek an analogous interpretation of the six-term exact
sequence associated with either short exact sequence in (\SGARefLink{sga7.sga7i.evii.n3-7-5-2}{3.7.5.2}),
say the first; equivalently, in view of \SGARefLink{sga7.sga7i.evii.n3-7-3}{3.7.3}, one may seek a
geometric interpretation of the boundary homomorphism
\begin{equation}\tag{\SGARefTarget{sga7.sga7i.evii.n3-7-5-4}{3.7.5.4}}
 \partial:\Biext^0(P',Q;G)\longrightarrow
 \Biext^1(P'',Q;G).
\end{equation}
Indeed, we have the following result.

\underline{Proposition} \SGARefTarget{sga7.sga7i.evii.n3-7-6}{3.7.6}. \underline{The cofibred category over}
\[
 \operatorname{GRAB}(\underline T)^{\mathrm{op}}
 \times\operatorname{GRAB}(\underline T)^{\mathrm{op}}
 \times\operatorname{GRAB}(\underline T)
\]
\underline{of biextensions of $(P,Q)$ by $G$, with $P,Q,G$
variable} (cf. (\SGARefLink{sga7.sga7i.evii.n2-6-1}{2.6.1})), \underline{is left exact not only in the
argument $G$, but also in each of the arguments $P$ and $Q$.
Moreover, the boundary homomorphism (\SGARefLink{sga7.sga7i.evii.n3-7-5-4}{3.7.5.4}) deduced from the exact
sequence of the $\Ext^i$ associated with (\SGARefLink{sga7.sga7i.evii.n3-7-5-2}{3.7.5.2}), by means of the
isomorphisms \SGARefLink{sga7.sga7i.evii.n3-6-5}{3.6.5} and \SGARefLink{sga7.sga7i.evii.n3-7-4}{3.7.4}, is the boundary homomorphism in the
theory of left exact cofibred categories over}
$\underline C=\operatorname{GRAB}(\underline T)^{\mathrm{op}}$
\underline{recalled in \SGARefLink{sga7.sga7i.evii.n3-3-6}{3.3.6}, applied to the cofibred category of
biextensions of $(P,Q)$ by $G$, where $Q$ and $G$ are fixed and $P$
varies} (subject to a possible sign error in \SGARefLink{sga7.sga7i.evii.n3-4-9}{3.4.9}; cf. \SGARefLink{sga7.sga7i.evii.n3-4-9-2}{3.4.9.2}).

We sketch the proof of \SGARefLink{sga7.sga7i.evii.n3-7-6}{3.7.6}.  By \SGARefLink{sga7.sga7i.evii.n3-6-4}{3.6.4} and compatibility (\SGARefLink{sga7.sga7i.evii.n3-7-2}{3.7.2}),
we interpret biextensions of $(P,Q)$ by $G$ as objects of
$\underline\Psi_{L_\bullet(P,Q)}(G)$.  Consider the homomorphism
\[
 L_\bullet(P,Q)
 =L_\bullet(P)\otimes L_\bullet(Q)
 \longrightarrow P\otimes L_\bullet(Q)
\]
defined by the augmentation $L_\bullet(P)\longrightarrow P$.
It follows from \SGARefLink{sga7.sga7i.evii.n3-5-3}{3.5.3} and the flatness of $L_\bullet(Q)$ that this
homomorphism induces isomorphisms on $H_0$ and $H_1$, and hence by
\SGARefLink{sga7.sga7i.evii.n3-4-7}{3.4.7} an equivalence
\[
 \underline\Psi_{P\otimes L_\bullet(Q)}
 \longrightarrow
 \underline\Psi_{L_\bullet(P,Q)}.
\]
We may therefore interpret biextensions of $(P,Q)$ by $G$ as objects
of $\underline\Psi_{P\otimes L_\bullet(Q)}(G)$.  The exact sequence
(\SGARefLink{sga7.sga7i.evii.n3-7-5-2}{3.7.5.2}) now gives an exact sequence of chain complexes
\begin{equation}\tag{\SGARefTarget{sga7.sga7i.evii.n3-7-6-1}{3.7.6.1}}
 0\longrightarrow P'\otimes L_\bullet(Q)
 \longrightarrow P\otimes L_\bullet(Q)
 \longrightarrow P''\otimes L_\bullet(Q)
 \longrightarrow0.
\end{equation}
The asserted left exactness follows from \SGARefLink{sga7.sga7i.evii.n3-4-8}{3.4.8} applied to this exact
sequence of complexes.  (A direct verification would in any event
be immediate.)  Moreover, the homomorphism (\SGARefLink{sga7.sga7i.evii.n3-7-5-4}{3.7.5.4}) is precisely the
boundary homomorphism $\Ext^0\longrightarrow\Ext^1$ deduced from this
exact sequence by taking the $\Ext^i$ with values in $G$.  The last
assertion of \SGARefLink{sga7.sga7i.evii.n3-7-6}{3.7.6} is consequently a special case of \SGARefLink{sga7.sga7i.evii.n3-4-9}{3.4.9}.

% END INCLUDED SGA7I_English_current_buildable_source_20260823/components/91_expose_VII_section_3_7.tex

% BEGIN INCLUDED SGA7I_English_current_buildable_source_20260823/components/92_expose_VII_section_3_8_and_bibliography.tex
% SOURCE AUTHORITY:
% expose_VII_body.tex lines 1810--1977; scan indices 223--228; source folios 212--217.
% Direct scan comparison: physical PDF pages 224--229 at 600 dpi.
% Transparent normalization: typewritten letter-o subscripts are rendered as zero.
% The headless diagonal in (3.8.2) is reproduced as printed, while the prose fixes
% its mathematical role as the trivializing homomorphism.  The derived-category
% arrow in (3.8.10) follows the French authority, from P derived-tensor Q to
% L_\bullet tensor M_\bullet; the inherited abridgment reversed it.  Hard stop at
% the end of the Expose VII bibliography.

\subsection*{\SGARefTarget{sga7.sga7i.evii.n3-8}{3.8}. Biextensions of $(P,Q)$ trivialized over extensions of $P$ and $Q$}

First consider two exact sequences of abelian groups
\begin{equation}\tag{\SGARefTarget{sga7.sga7i.evii.n3-8-1}{3.8.1}}
 0\longrightarrow L_1\xrightarrow{\ i\ }L_0\xrightarrow{\ u\ }P
 \longrightarrow0,\qquad
 0\longrightarrow M_1\xrightarrow{\ j\ }M_0\xrightarrow{\ v\ }Q
 \longrightarrow0.
\end{equation}
We shall describe the category of biextensions $E$ of $(P,Q)$ by an
abelian group $G$ for which the inverse image $(u,v)^*(E)$ is trivial;
that is, those for which there exists a homomorphism $s$ making the
following triangle commutative:
% SGA-SRC-001262 — reconstructed by mathematical inference; scan inconclusive
\begin{equation}\tag{\SGARefTarget{sga7.sga7i.evii.n3-8-2}{3.8.2}}
\begin{tikzcd}[column sep=huge,row sep=huge]
 E\arrow[d]&\\
 P\times Q&L_0\times M_0\arrow[l,"u\times v"']\arrow[ul,"s"]
\end{tikzcd}
\end{equation}
and satisfying the additivity condition required in each of its two
arguments.  The choice of such an $s$ therefore identifies
$(u,v)^*(E)$ with the trivial biextension
$E_0=L_0\times M_0\times G$ of $(L_0,M_0)$ by $G$.  The left
exactness in $P,Q$ of the cofibred category of biextensions of
$(P,Q)$ by fixed $G$ (\SGARefLink{sga7.sga7i.evii.n3-7-6}{3.7.6}) then shows that to give the biextension
$E$ of $(P,Q)$, together with the trivialization $s$ of its inverse
image, is essentially equivalent to giving two partial
trivializations of the trivial biextension $E_0$ of $(L_0,M_0)$ by
$G$, over $(L_1,M_0)$ and $(L_0,M_1)$ respectively, which agree over
$(L_1,M_1)$.  Since the biextensions induced by $E_0$ over
$(L_1,M_0)$, $(L_0,M_1)$, and $(L_1,M_1)$ identify with the trivial
biextensions, a partial trivialization over $(L_1,M_0)$ (respectively
over $(L_0,M_1)$) is equivalent to giving a homomorphism (or pairing)
\begin{equation}\tag{\SGARefTarget{sga7.sga7i.evii.n3-8-3}{3.8.3}}
 f:L_1\otimes M_0\longrightarrow G
 \qquad
 \bigl(\text{respectively }g:L_0\otimes M_1\longrightarrow G\bigr).
\end{equation}
The compatibility in question simply means that the two
homomorphisms agree over $L_1\otimes M_1$, namely
\begin{equation}\tag{\SGARefTarget{sga7.sga7i.evii.n3-8-4}{3.8.4}}
 f(x_1\otimes y_1)=g(x_1\otimes y_1),
\end{equation}
for sections $x_1,y_1$ of $L_1,M_1$ over an object of
$\underline T$.  In terms of the homomorphism $s$ of (\SGARefLink{sga7.sga7i.evii.n3-8-2}{3.8.2}), the
forms $f$ and $g$ are characterized by
\begin{equation}\tag{\SGARefTarget{sga7.sga7i.evii.n3-8-5}{3.8.5}}
 s(x_1,y_0)=f(x_1\otimes y_0)e_{E/Q}(y_0),\qquad
 s(x_0,y_1)=g(x_0\otimes y_1)e_{E/P}(x_0),
\end{equation}
which again makes clear the necessity of condition (\SGARefLink{sga7.sga7i.evii.n3-8-4}{3.8.4}).

Once the trivialization $s$ has been chosen, the other
trivializations are the morphisms
$s':L_0\times M_0\longrightarrow E$ of the form
\begin{equation}\tag{\SGARefTarget{sga7.sga7i.evii.n3-8-6}{3.8.6}}
 s'(x_0,y_0)=s(x_0,y_0)+h(x_0\otimes y_0),
\end{equation}
where
\[
 h:L_0\otimes M_0\longrightarrow G
\]
is an arbitrary form.  It is then clear from (\SGARefLink{sga7.sga7i.evii.n3-8-5}{3.8.5}) that the forms
$f,g$ defined by $s$ are changed into
\begin{equation}\tag{\SGARefTarget{sga7.sga7i.evii.n3-8-7}{3.8.7}}
\begin{aligned}
 f'(x_1\otimes y_0)&=f(x_1\otimes y_0)+h(x_1\otimes y_0),\\
 g'(x_0\otimes y_1)&=g(x_0\otimes y_1)+h(x_0\otimes y_1).
\end{aligned}
\end{equation}
Thus we obtain:

\underline{Proposition} \SGARefTarget{sga7.sga7i.evii.n3-8-8}{3.8.8}.
\underline{The subgroup of $\Biext^1(P,Q;G)$ formed by the classes
of biextensions whose inverse image under $(u,v)$ is a trivial
biextension of $(L_0,M_0)$} (that is, the kernel of
$\Biext^1(P,Q;G)\longrightarrow\Biext^1(L_0,M_0;G)$)
\underline{is canonically isomorphic, by the preceding description,
to the quotient of the group of pairs of forms $(f,g)$ in (\SGARefLink{sga7.sga7i.evii.n3-8-3}{3.8.3})
satisfying compatibility condition (\SGARefLink{sga7.sga7i.evii.n3-8-4}{3.8.4}), by the pairs of forms
which arise by restriction from a form}
$h:L_0\otimes M_0\longrightarrow G$.

We now interpret this isomorphism in terms of homological algebra,
taking into account the canonical isomorphism
\[
 \Biext^1(P,Q;G)\simeq
 \Ext^1(P\mathbin{\overset{\mathbb L}{\otimes}}Q,G)
\]
of \SGARefLink{sga7.sga7i.evii.n3-6-5}{3.6.5}.  To this end, interpret the data (\SGARefLink{sga7.sga7i.evii.n3-8-1}{3.8.1}) as two complexes
resolving $P$ and $Q$, respectively:
\begin{equation}\tag{\SGARefTarget{sga7.sga7i.evii.n3-8-9}{3.8.9}}
 L_\bullet:\ 0\longrightarrow L_1\longrightarrow L_0
 \longrightarrow0,\qquad
 M_\bullet:\ 0\longrightarrow M_1\longrightarrow M_0
 \longrightarrow0,
\end{equation}
and consider the tensor-product complex $L_\bullet\otimes M_\bullet$
\[
 0\longrightarrow L_1\otimes M_1
 \xrightarrow{d_2}
 L_1\otimes M_0+L_0\otimes M_1
 \xrightarrow{d_1}
 L_0\otimes M_0\longrightarrow0.
\]
There is, moreover, a canonical homomorphism in the derived category
\[
 P\mathbin{\overset{\mathbb L}{\otimes}}Q
 \longrightarrow L_\bullet\otimes M_\bullet,
\]
and hence an equally canonical homomorphism
\begin{equation}\tag{\SGARefTarget{sga7.sga7i.evii.n3-8-10}{3.8.10}}
 H^1\!\left(\Hom^{\cdot}(L_\bullet\otimes M_\bullet,G)\right)
 \longrightarrow
 \Ext^1(P\mathbin{\overset{\mathbb L}{\otimes}}Q,G).
\end{equation}
The group of $1$-cocycles of
$\Hom^{\cdot}(L_\bullet\otimes M_\bullet,G)$ identifies with the
group of pairs $(f,g)$ as in (\SGARefLink{sga7.sga7i.evii.n3-8-3}{3.8.3}) which satisfy compatibility
condition (\SGARefLink{sga7.sga7i.evii.n3-8-4}{3.8.4}): to such a pair one associates the homomorphism
$F:L_1\otimes M_0+L_0\otimes M_1\longrightarrow G$ with components
$(f,g)$,
\begin{equation}\tag{\SGARefTarget{sga7.sga7i.evii.n3-8-11}{3.8.11}}
 F(x_1\otimes y_0+x_0\otimes y_1)
 =f(x_1\otimes y_0)+g(x_0\otimes y_1).
\end{equation}
Moreover, the coboundary of the $0$-cochain
$h:L_0\otimes M_0\longrightarrow G$ is precisely the pair of its
restrictions to $L_1\otimes M_0$ and $L_0\otimes M_1$.  Consequently,
the group of pairs $(f,g)$ modulo the pairs $(h,h)$ described in
\SGARefLink{sga7.sga7i.evii.n3-8-8}{3.8.8} is canonically isomorphic to the first term of (\SGARefLink{sga7.sga7i.evii.n3-8-10}{3.8.10}).
With this understood, one has the following compatibility.

\underline{Proposition} \SGARefTarget{sga7.sga7i.evii.n3-8-12}{3.8.12}.
\underline{If}
\[
 \xi\in\Ker\!\left(
 \Biext^1(P,Q;G)\longrightarrow\Biext^1(L_0,M_0;G)\right)
\]
\underline{is described by the pair $(f,g)$ as in \SGARefLink{sga7.sga7i.evii.n3-8-8}{3.8.8}, then the
image of $\xi$ in}
$\Ext^1(P\mathbin{\overset{\mathbb L}{\otimes}}Q,G)$
\underline{under the isomorphism (\SGARefLink{sga7.sga7i.evii.n3-6-5-1}{3.6.5.1}) is the image under
(\SGARefLink{sga7.sga7i.evii.n3-8-10}{3.8.10}) of the element of its first term defined by the cocycle
(\SGARefLink{sga7.sga7i.evii.n3-8-11}{3.8.11})}.

The verification of this compatibility is left to the reader.

\SGARefTarget{sga7.sga7i.evii.n3-8-13}{3.8.13}. We shall apply this especially in the case
\[
 L_0=\mathbb Z[P],\qquad M_0=\mathbb Z[Q],
\]
where $L_1$ and $M_1$ are respectively the kernels of the canonical
homomorphisms $\mathbb Z[P]\longrightarrow P$ and
$\mathbb Z[Q]\longrightarrow Q$.  Notice the canonical isomorphisms
\[
 L_0\mathbin{\overset{\mathbb L}{\otimes}}M_0
 \simeq L_0\otimes M_0
 \simeq\mathbb Z[P\times Q].
\]
Since $L_0$ and $M_0$ are flat, (\SGARefLink{sga7.sga7i.evii.n1-4-5}{1.4.5}) gives canonical isomorphisms
\[
 \Ext^i(L_0\mathbin{\overset{\mathbb L}{\otimes}}M_0,G)
 \simeq H^i(P\times Q,G_{P\times Q}).
\]
It follows from this and \SGARefLink{sga7.sga7i.evii.n3-6-5}{3.6.5}, or directly from the definitions and
\SGARefLink{sga7.sga7i.evii.n1-4}{1.4}, that the functor which assigns to each biextension of
$(L_0,M_0)$ by $G$ the $G_{P\times Q}$-torsor over $P\times Q$ that
it induces is an equivalence of categories.  (This remains valid
without using an additive structure on $P,Q$.)  Consequently, if
$E$ is a biextension of $(P,Q)$ by $G$, giving a trivialization of
its inverse image over $(L_0,M_0)$ is equivalent to giving a section
of $E$ considered only as an object over $P\times Q$, that is, a
trivialization of $E$ considered as a torsor under $G_{P\times Q}$.

Biextensions equipped with this additional structure are therefore
described by the pairs $(f,g)$ of (\SGARefLink{sga7.sga7i.evii.n3-8-3}{3.8.3}) satisfying compatibility
(\SGARefLink{sga7.sga7i.evii.n3-8-4}{3.8.4}), and their classification up to isomorphism is the one given
in \SGARefLink{sga7.sga7i.evii.n3-8-8}{3.8.8}.

With the notation of \SGARefLink{sga7.sga7i.evii.n3-5}{3.5}, one may also regard $L_1$ as the cokernel
of $L_2(P)\longrightarrow L_1(P)$, and similarly for $M_1$.  This
makes it possible to express the data $(f,g)$ in terms of a system
of homomorphisms
\[
\left\{
\begin{aligned}
 f'&:L_1(P)\otimes L_0(Q)
     =\mathbb Z(P\times P\times Q)\longrightarrow G,\\
 g'&:L_0(P)\otimes L_1(Q)
     =\mathbb Z(P\times Q\times Q)\longrightarrow G,
\end{aligned}
\right.
\]
that is, homomorphisms of sheaves of \underline{sets}
\[
 P\times P\times Q\longrightarrow G,\qquad
 P\times Q\times Q\longrightarrow G,
\]
satisfying five compatibility conditions.  These express that $f'$
and $g'$ factor through $f$ and $g$---which gives $2\times2$
equations, in view of the expression of $L_2(P)$ and $L_2(Q)$ as
sums of two terms---together with condition (\SGARefLink{sga7.sga7i.evii.n3-8-4}{3.8.4}).  One then sees
that $f'$ and $g'$ are precisely the partial multiplication data
defining a biextension, and that the five conditions in question are
exactly the five axioms for biextensions considered in \SGARefLink{sga7.sga7i.evii.n2-1}{2.1} (cf. the
proof of \SGARefLink{sga7.sga7i.evii.n3-6-4}{3.6.4}).

\subsection*{Bibliography to Expos\'e VII}

\begin{enumerate}[label={[\arabic*]}]
\item M. Demazure, J. Giraud, and M. Raynaud,
\emph{Schémas abéliens}, Séminaire de la Faculté des Sciences
d'Orsay 1967/68 (to appear).
\item J. Giraud, \emph{Cohomologie non abélienne}, thesis, Paris,
1967 (to appear).
\item A. Grothendieck, \emph{Sur quelques points d'Algèbre
homologique}, Tohoku Mathematical Journal 9 (1956), 119--221.
\item A. Grothendieck, \emph{Catégories cofibrées additives et
Complexe cotangent relatif}, Lecture Notes in Mathematics 79,
Springer, 1968.
\item J.-L. Verdier, \emph{Catégories dérivées des catégories
abéliennes} (to appear).
\item D. Mumford, \emph{Biextensions of formal groups}, Tata
Institute of Fundamental Research, 1968 (to appear).
\end{enumerate}

% END INCLUDED SGA7I_English_current_buildable_source_20260823/components/92_expose_VII_section_3_8_and_bibliography.tex

\clearpage
\phantomsection\addcontentsline{toc}{section}{Exposé VIII}
% BEGIN INCLUDED SGA7I_English_current_buildable_source_20260823/components/93_expose_VIII_opening_and_introduction.tex
% SOURCE AUTHORITY:
% expose_VIII_body.tex lines 5--52; scan indices 229--231; source folios 218--220.
% Direct scan comparison: physical PDF pages 230--232 at 1100 dpi.
% Hard stop before section 1, source line 55.

\begin{center}
\SGARefAliasTarget{sga7.sga7i.eviii.nroot}SGA 7

\vspace{1em}
\underline{Expos\'e VIII}

\vspace{1em}
\underline{COMPLEMENTS ON BIEXTENSIONS. GENERAL PROPERTIES}

\underline{OF BIEXTENSIONS OF GROUP SCHEMES}

\vspace{1em}
\underline{by A. Grothendieck}
\end{center}

\subsection*{Contents}

\begin{enumerate}[label=\arabic*.]
\setcounter{enumi}{-1}
\item Introduction
\item Various special cases over an arbitrary topos
\item Pairings defined by a biextension
\item Biextensions of group schemes $(P,Q)$ by $\underline G_m$: generalities
\item Extensions and biextensions of smooth connected group schemes over a field
\item Extensions and biextensions by constant truncated torsion-free group schemes
\item Extensions and biextensions by the N\'eron model of $\underline G_m$
\item Canonical prolongations of extensions and biextensions by $\underline G_m$
\end{enumerate}

0. \underline{Introduction}

\SGARefTarget{sga7.sga7i.eviii.n0-1}{0.1}. We continue here the general study of biextensions begun in the
preceding expos\'e.  First we give some complements in the case of an
arbitrary base topos, developing in particular, with the detail it
deserves, the formalism of the pairings associated with biextensions
(Section 2), whose importance in the theory of abelian schemes is well
known.  We then begin the study of biextensions of general group
schemes, the case of greatest interest for us being that of
biextensions by the multiplicative group $\underline G_m$.  We shall
nevertheless be obliged, as a technical intermediate step, also to
study biextensions by other types of group schemes (Sections 5 and 6).

The most interesting result for what follows is given in \SGARefLink{sga7.sga7i.eviii.n7-1}{7.1}, which
gives conditions under which, for $P$ and $Q$ smooth over a trait with
generic point $\eta$, a biextension of $(P_\eta,Q_\eta)$ by
$\underline G_{m\eta}$ can be prolonged to a biextension of $(P,Q)$ by
$\underline G_m$, determined up to a unique isomorphism.  This result
will be very convenient in the following expos\'e, devoted to the
theory of the N\'eron model of abelian varieties.  It is likely that
one could manage there at less expense, without the imposing
paperweight of Expos\'es VII and VIII (which will no longer be used in
the remainder of the Seminar); it nevertheless seems likely to us that
the developments given in these expos\'es are destined to be useful
elsewhere as well.

\SGARefTarget{sga7.sga7i.eviii.n0-2}{0.2}. When we work with group schemes over a scheme $S$, it will be
understood that we work with the fppf topos of $S$.  This is the topos
associated with the $\underline U$-site of schemes locally of finite
presentation over $S$ (and belonging to $\underline U$), endowed with
the fppf topology (SGA 3 IV 6.3).  The reader will readily verify that
most of the statements we shall give would also be valid (and
essentially equivalent to the fppf statements) for other topologies on
$S$, whose underlying site contains at least the schemes locally of
finite presentation over $S$, such as the Zariski, \'etale, and fpqc
topologies.  One reason is that, by descent theory, the category of
extensions of a group scheme $P$ locally of finite presentation over
$S$ by the group scheme $\underline G_m$ does not depend, up to
equivalence, on the chosen topology; the same statement holds for the
notion of biextension.

The technical reason for preferring the fppf topology (or a finer
topology, such as fpqc) to coarser topologies such as the Zariski and
\'etale topologies is that, for every integer $n>0$, the homomorphism
$n\cdot\id$ from $\underline G_m$ to itself is an \underline{fppf
epimorphism} (because it is faithfully flat and of finite
presentation), but in general is not a Zariski or an \'etale
epimorphism.  Notice that it is an \'etale epimorphism if and only if
$n$ is prime to the residue characteristics of $S$, and a Zariski
epimorphism if and only if $S$ is empty.

\SGARefTarget{sga7.sga7i.eviii.n0-3}{0.3}. Finally, most of the statements concerning extensions or
biextensions by $\underline G_m$ would remain valid if that group were
replaced by a \underline{torus} $G$.  Since such a $G$ is locally
isomorphic (for the fppf, and even for the \'etale topology) to a group
$\underline G_m^r$, this is an essentially trivial generalization.

% END INCLUDED SGA7I_English_current_buildable_source_20260823/components/93_expose_VIII_opening_and_introduction.tex

% BEGIN INCLUDED SGA7I_English_current_buildable_source_20260823/components/94_expose_VIII_section_1.tex
% SOURCE AUTHORITY:
% expose_VIII_body.tex lines 55--276; scan indices 231--237; source folios 220--226.
% Direct scan comparison: physical PDF pages 232--238 at 1100 dpi.
% Transparent source dispositions:
% - the sentence before (1.1.3) is stated in standard truncation language rather
%   than repeating the printed impossible claim that all Tor_i vanish for i >= 1;
% - the fourth term of (1.1.4) uses Hom(Q,G), as required by the displayed
%   spectral sequence, rather than the printed Hom(Q,P);
% - broken parentheses in (1.3.5), the transition exponent in part (c), and the
%   references to the final member of (1.3.5) are restored from the formulas;
% - in 1.6 the parenthetical Biext argument and the terminal Ext coefficient are
%   corrected to Q and G respectively.
% Hard stop before section 2, source line 279.

\subsection*{\SGARefTarget{sga7.sga7i.eviii.n1}{1}. Various special cases over an arbitrary topos}

\SGARefTarget{sga7.sga7i.eviii.n1-1}{1.1}. In this section and the next, as in Expos\'e VII, $T$ denotes a
topos, and $P,Q,G$ denote three abelian groups of $T$.  Thus there is a
canonical isomorphism (VII \SGARefLink{sga7.sga7i.evii.n3-6-5}{3.6.5})
\begin{equation}\tag{\SGARefTarget{sga7.sga7i.eviii.n1-1-1}{1.1.1}}
 \Biext^1(P,Q;G)
 \simeq
 \Ext^1\!\left(P\mathbin{\overset{\mathbb L}{\otimes}}Q,G\right).
\end{equation}
Recall the canonical isomorphisms (called the Cartan isomorphisms)
\begin{equation}\tag{\SGARefTarget{sga7.sga7i.eviii.n1-1-2}{1.1.2}}
\begin{split}
 \Ext^1\!\left(P\mathbin{\overset{\mathbb L}{\otimes}}Q,G\right)
 &\simeq \Ext^1\!\left(P,R\uHom(Q,G)\right)\\
 &\simeq \Ext^1\!\left(Q,R\uHom(P,G)\right).
\end{split}
\end{equation}
All these isomorphisms are functorial in the three arguments $P,Q,G$.

Taking the standard truncation of the derived tensor product, whose
homology objects are the $\Tor_i(P,Q)$, one obtains from (\SGARefLink{sga7.sga7i.eviii.n1-1-1}{1.1.1}) the
following infinite exact sequence:
\begin{equation}\tag{\SGARefTarget{sga7.sga7i.eviii.n1-1-3}{1.1.3}}
\begin{aligned}
0\longrightarrow{}&
 \Ext^1(P\otimes Q,G)
 \longrightarrow \Biext^1(P,Q;G)
 \longrightarrow \Hom(\Tor_1(P,Q),G)\\
 \longrightarrow{}&
 \Ext^2(P\otimes Q,G)
 \longrightarrow
 \Ext^2\!\left(P\mathbin{\overset{\mathbb L}{\otimes}}Q,G\right)
 \longrightarrow\cdots .
\end{aligned}
\end{equation}
On the other hand, since the cohomology objects of $R\uHom(Q,G)$ are
the $\uExt^i(Q,G)$, the first equality in (\SGARefLink{sga7.sga7i.eviii.n1-1-2}{1.1.2}) gives rise to the
five-term exact sequence
\begin{equation}\tag{\SGARefTarget{sga7.sga7i.eviii.n1-1-4}{1.1.4}}
\begin{aligned}
0\longrightarrow{}&
 \Ext^1(P,\uHom(Q,G))
 \longrightarrow \Biext^1(P,Q;G)
 \longrightarrow \Hom(P,\uExt^1(Q,G))\\
 \longrightarrow{}&
 \Ext^2(P,\uHom(Q,G))
 \longrightarrow \Ext^2(P,R\uHom(Q,G)),
\end{aligned}
\end{equation}
associated with the usual spectral sequence
\[
 \Ext^*(P,R\uHom(Q,G))
 \Longleftarrow
 \Ext^p(P,\uExt^q(Q,G)).
\]
We leave it to the reader, should the need arise, to make the first
arrow of (\SGARefLink{sga7.sga7i.eviii.n1-1-3}{1.1.3}), respectively of (\SGARefLink{sga7.sga7i.eviii.n1-1-4}{1.1.4}), explicit by means of the
\underline{natural fully faithful functor}
\begin{equation}\tag{\SGARefTarget{sga7.sga7i.eviii.n1-1-5}{1.1.5}}
 \operatorname{EXT}(P\otimes Q,G)
 \longrightarrow
 \operatorname{BIEXT}(P,Q;G),
\end{equation}
respectively
\begin{equation}\tag{\SGARefTarget{sga7.sga7i.eviii.n1-1-6}{1.1.6}}
 \operatorname{EXT}(P,\uHom(Q,G))
 \longrightarrow
 \operatorname{BIEXT}(P,Q;G),
\end{equation}
where $\operatorname{EXT}$ and $\operatorname{BIEXT}$ denote the
evident categories of extensions and biextensions.

Of course, the isomorphism
\[
 \Biext^1(P,Q;G)\simeq\Ext^1(Q,R\uHom(P,G))
\]
gives rise to an exact sequence analogous to (\SGARefLink{sga7.sga7i.eviii.n1-1-3}{1.1.3}), with the roles
of $P$ and $Q$ interchanged.  Its first arrow can be made explicit by
means of the fully faithful functor
\begin{equation}\tag{\SGARefTarget{sga7.sga7i.eviii.n1-1-7}{1.1.7}}
 \operatorname{EXT}(Q,\uHom(P,G))
 \longrightarrow
 \operatorname{BIEXT}(P,Q;G)
\end{equation}
analogous to (\SGARefLink{sga7.sga7i.eviii.n1-1-6}{1.1.6}).

\SGARefTarget{sga7.sga7i.eviii.n1-2}{1.2}. Suppose that
\begin{equation}\tag{\SGARefTarget{sga7.sga7i.eviii.n1-2-1}{1.2.1}}
 \Tor_1(P,Q)=0.
\end{equation}
Then (\SGARefLink{sga7.sga7i.eviii.n1-1-3}{1.1.3}) supplies a canonical isomorphism
\begin{equation}\tag{\SGARefTarget{sga7.sga7i.eviii.n1-2-2}{1.2.2}}
 \Biext^1(P,Q;G)\simeq\Ext^1(P\otimes Q,G).
\end{equation}
This holds, for example, if $P$ or $Q$ is a \underline{flat}
$\mathbb Z$-module.  More precisely, under condition (\SGARefLink{sga7.sga7i.eviii.n1-2-1}{1.2.1}), the
functor (\SGARefLink{sga7.sga7i.eviii.n1-1-5}{1.1.5}) is an equivalence of categories.

\SGARefTarget{sga7.sga7i.eviii.n1-3}{1.3}. Suppose that
\begin{equation}\tag{\SGARefTarget{sga7.sga7i.eviii.n1-3-1}{1.3.1}}
 P\otimes Q=0.
\end{equation}
Then (\SGARefLink{sga7.sga7i.eviii.n1-1-3}{1.1.3}) supplies a canonical isomorphism
\begin{equation}\tag{\SGARefTarget{sga7.sga7i.eviii.n1-3-2}{1.3.2}}
 \Biext^1(P,Q;G)\simeq\Hom(\Tor_1(P,Q),G).
\end{equation}
The case of greatest interest to us in which (\SGARefLink{sga7.sga7i.eviii.n1-3-1}{1.3.1}) holds is that in
which there is an integer $n>0$ such that
\begin{equation}\tag{\SGARefTarget{sga7.sga7i.eviii.n1-3-3}{1.3.3}}
 nP=P,\qquad
 Q=\varinjlim_i {}_{n^i}Q;
\end{equation}
that is, $P$ is divisible by $n$, and $Q$ is the limit of subgroups
annihilated by powers of $n$.  Put, as usual,
\[
 T_n(P)=({}_{n^i}P)_{i\geq0}\quad\text{(a projective system)},\qquad
 {}_{n^\infty}P=\varinjlim_i{}_{n^i}P.
\]
From (\SGARefLink{sga7.sga7i.eviii.n1-3-3}{1.3.3}) one obtains the following formulas, more precise than
(\SGARefLink{sga7.sga7i.eviii.n1-3-1}{1.3.1}):
\begin{equation}\tag{\SGARefTarget{sga7.sga7i.eviii.n1-3-4}{1.3.4}}
\begin{split}
 P\otimes Q&=0,\\
 \Tor_1(P,Q)
 &\simeq\varinjlim_i
   \bigl({}_{n^i}P\otimes{}_{n^i}Q\bigr)
 \simeq T_n(P)\otimes Q\\
 &\simeq{}_{n^\infty}P\otimes T_n(Q).
\end{split}
\end{equation}
Thus the isomorphism (\SGARefLink{sga7.sga7i.eviii.n1-3-2}{1.3.2}) here assumes the more concrete form
\begin{equation}\tag{\SGARefTarget{sga7.sga7i.eviii.n1-3-5}{1.3.5}}
\begin{split}
 \Biext^1(P,Q;G)
 &\simeq\Hom(T_n(P),D(Q))\\
 &\simeq\Hom(Q,D(T_n(P)))\\
 &\simeq\Hom(T_n(P)\otimes T_n(Q),T_n(G)).
\end{split}
\end{equation}
Let us specify the meaning of these formulas.

\begin{enumerate}[label=\alph*)]
\item In the inductive system occurring in (\SGARefLink{sga7.sga7i.eviii.n1-3-4}{1.3.4}), the transition
homomorphism
\[
 {}_{n^i}P\otimes{}_{n^i}Q
 \longrightarrow
 {}_{n^j}P\otimes{}_{n^j}Q
 \qquad(i\leq j)
\]
is described by identifying the first member with
$\,{}_{n^j}P\otimes{}_{n^i}Q$ by means of
$n^{j-i}\otimes\id$, and then using the homomorphism
\[
 \id\otimes\operatorname{inj}:
 {}_{n^j}P\otimes{}_{n^i}Q
 \longrightarrow
 {}_{n^j}P\otimes{}_{n^j}Q,
\]
where
$\operatorname{inj}:{}_{n^i}Q\longrightarrow{}_{n^j}Q$
denotes the inclusion.

\item By definition, put
\[
 T_n(P)\otimes Q
 \overset{\mathrm{def}}{=}
 \varinjlim_i\bigl(T_n(P)\otimes{}_{n^i}Q\bigr),
\]
where
\[
 T_n(P)\otimes M
 \overset{\mathrm{def}}{=}
 {}_{n^i}P\otimes M
 \qquad\text{if }n^iM=0.
\]
The meaning of the third term written in (\SGARefLink{sga7.sga7i.eviii.n1-3-4}{1.3.4}) is specified in the
same way when $nQ=Q$ is also assumed, so that $T_n(Q)$, like
$T_n(P)$, is a strict projective system.  The last term of (\SGARefLink{sga7.sga7i.eviii.n1-3-5}{1.3.5}) is
a Hom between projective systems of groups.

\item The letter $D$ denotes $\uHom(-,G)$.  The object $D(Q)$ is
defined as the projective system of groups
\[
 D(Q)=\bigl(D({}_{n^i}Q)\bigr)_{i\geq0}
 =\bigl(\uHom({}_{n^i}Q,G)\bigr)_{i\geq0},
\]
the second member of (\SGARefLink{sga7.sga7i.eviii.n1-3-5}{1.3.5}) being the group of homomorphisms of
projective systems of groups.  The following definition gives a
simpler interpretation of the next member:
\[
 D(T_n(P))=\varinjlim_iD({}_{n^i}P),
\]
where the transition morphisms of this inductive system are obtained
by transposition from the transition morphisms
\[
 n^{j-i}:{}_{n^j}P\longrightarrow{}_{n^i}P
\]
of $T_n(P)$.
\end{enumerate}

We leave it to the reader to make explicit the calculation of
$\Tor_1(P,Q)$ indicated in (\SGARefLink{sga7.sga7i.eviii.n1-3-4}{1.3.4}), and to derive (\SGARefLink{sga7.sga7i.eviii.n1-3-5}{1.3.5}) from it.
We note only that one uses the relations, deduced from the second
hypothesis in (\SGARefLink{sga7.sga7i.eviii.n1-3-3}{1.3.3}),
\[
 \Tor_p(P,Q)\simeq
 \varinjlim_i\Tor_p(P,{}_{n^i}Q).
\]
For $p=0$ this immediately gives $P\otimes Q=0$, and it reduces the
calculation of $\Tor_1(P,Q)$ to that of the
$\Tor_1(P,{}_{n^i}Q)$ and their transition morphisms.  Thus, when
$Q$ is annihilated by an integer $m=n^i$ such that $mP=P$, it remains
to make explicit an isomorphism
\begin{equation}\tag{\SGARefTarget{sga7.sga7i.eviii.n1-3-6}{1.3.6}}
 \Tor_1(P,Q)
 \xleftarrow{\ \sim\ }
 {}_mP\otimes{}_mQ
 \qquad(mP=P,\ mQ=0).
\end{equation}
There are several ways to do this.  One consists in using the exact
Tor sequence associated with
\[
 0\longrightarrow{}_mP\longrightarrow P
 \xrightarrow{\ m\ }P\longrightarrow0.
\]
One may also use the homomorphism (\SGARefLink{sga7.sga7i.eviii.n2-1-1}{2.1.1}) below; the fact that the
transition morphisms are those specified in part (a) is then precisely
\SGARefLink{sga7.sga7i.eviii.n2-1-11}{2.1.11}.  In any case, the choice of (\SGARefLink{sga7.sga7i.eviii.n1-3-6}{1.3.6}) amounts to the choice of a
sign; see \SGARefLink{sga7.sga7i.eviii.n2-1-10}{2.1.10}.

\SGARefTarget{sga7.sga7i.eviii.n1-4}{1.4}. Suppose that
\begin{equation}\tag{\SGARefTarget{sga7.sga7i.eviii.n1-4-1}{1.4.1}}
 \uHom(Q,G)=0.
\end{equation}
Then (\SGARefLink{sga7.sga7i.eviii.n1-1-4}{1.1.4}) supplies a canonical isomorphism
\begin{equation}\tag{\SGARefTarget{sga7.sga7i.eviii.n1-4-2}{1.4.2}}
 \Biext^1(P,Q;G)
 \xrightarrow{\ \sim\ }
 \Hom(P,\uExt^1(Q,G)).
\end{equation}
This case occurs in particular when $Q$ is an abelian scheme and
$G=\underline G_m$; see VII \SGARefLink{sga7.sga7i.evii.n1-3-8}{1.3.8}.  Hence the functor in $P$,
\[
 P\longmapsto\Biext^1(P,Q;G),
\]
is \underline{representable} by $Q^*=\uExt^1(Q,G)$.  In particular,
there is a canonical biextension $E^Q$ of $(Q^*,Q)$ by $G$, determined
up to a unique isomorphism, since its automorphism group is
\[
 \Hom(Q^*,\uHom(Q,G))=0.
\]
This biextension, or its symmetric biextension $\,{}^sE^Q$ of
$(Q,Q^*)$ by $G$, may be called the \underline{Weil biextension} of
$(Q^*,Q)$, respectively $(Q,Q^*)$, by $G$.

\SGARefTarget{sga7.sga7i.eviii.n1-5}{1.5}. Suppose that
\begin{equation}\tag{\SGARefTarget{sga7.sga7i.eviii.n1-5-1}{1.5.1}}
 \Hom(P,\uExt^1(Q,G))=0.
\end{equation}
This holds in particular if $\uExt^1(Q,G)=0$; for example, if $Q$ is
a finite locally free group scheme or a group of multiplicative type
(SGA 3 IX), and $G=\underline G_m$; see \SGARefLink{sga7.sga7i.eviii.n3-4}{3.4} below.  Then (\SGARefLink{sga7.sga7i.eviii.n1-1-4}{1.1.4})
supplies a canonical isomorphism
\begin{equation}\tag{\SGARefTarget{sga7.sga7i.eviii.n1-5-2}{1.5.2}}
 \Biext^1(P,Q;G)
 \xleftarrow{\ \sim\ }
 \Ext^1(P,\uHom(Q,G)),
\end{equation}
where $\uHom(Q,G)$ plays the role of a "dual group" of $Q$ (relative
to $G$).  More precisely, it follows from (\SGARefLink{sga7.sga7i.eviii.n1-5-2}{1.5.2}) that the fully
faithful functor (\SGARefLink{sga7.sga7i.eviii.n1-1-6}{1.1.6}) is an \underline{equivalence of categories}.

\SGARefTarget{sga7.sga7i.eviii.n1-6}{1.6}. Suppose that the two relations
\begin{equation}\tag{\SGARefTarget{sga7.sga7i.eviii.n1-6-1}{1.6.1}}
 \uHom(P,G)=0,\qquad
 \Hom(P,\uExt^1(Q,G))=0
\end{equation}
hold simultaneously.  By \SGARefLink{sga7.sga7i.eviii.n1-4}{1.4} and \SGARefLink{sga7.sga7i.eviii.n1-5}{1.5}, respectively, these two
relations give isomorphisms
\[
 \Biext^1(P,Q;G)
 \xrightarrow{\ \sim\ }
 \Hom(Q,\uExt^1(P,G)),
\]
\[
 \Biext^1(P,Q;G)
 \xleftarrow{\ \sim\ }
 \Ext^1(P,\uHom(Q,G)).
\]
Their composite is the canonical isomorphism
\begin{equation}\tag{\SGARefTarget{sga7.sga7i.eviii.n1-6-2}{1.6.2}}
 \Ext^1(P,\uHom(Q,G))
 \xrightarrow{\ \sim\ }
 \Hom(Q,\uExt^1(P,G))
 \quad\bigl(\simeq\Biext^1(P,Q;G)\bigr),
\end{equation}
which plays the role of a duality formula.

Notice also that the first relation in (\SGARefLink{sga7.sga7i.eviii.n1-6-1}{1.6.1}) plainly implies
$\Hom(P\otimes Q,G)=0$.  Hence the category
$\operatorname{EXT}(P,\uHom(Q,G))$ of commutative extensions of $P$
by $\uHom(Q,G)$ is rigid (every automorphism of an object is the
identity), as is $\operatorname{BIEXT}(P,Q;G)$.  The structure of
these categories is therefore known once the group of isomorphism
classes of objects is known; by (\SGARefLink{sga7.sga7i.eviii.n1-6-2}{1.6.2}), this group is canonically
isomorphic to
\[
 \Hom(Q,\uExt^1(P,G)).
\]

The map (\SGARefLink{sga7.sga7i.eviii.n1-6-2}{1.6.2}) is in fact defined in all cases, without assuming
(\SGARefLink{sga7.sga7i.eviii.n1-6-1}{1.6.1}), and can also be described directly, without passing through
$\Biext^1(P,Q;G)$, in terms of the natural pairing
\[
 Q\times Q'\longrightarrow G,
\qquad Q'=\uHom(Q,G).
\]
Interpreting it as a homomorphism
$Q\longrightarrow\uHom(Q',G)$ gives a composite
\begin{equation}\tag{\SGARefTarget{sga7.sga7i.eviii.n1-6-3}{1.6.3}}
\begin{split}
 \Ext^1(P,Q')\times Q
 &\longrightarrow
 \Ext^1(P,Q')\times\Hom(Q',G)\\
 &\longrightarrow
 \Ext^1(P,G),
\end{split}
\end{equation}
which recovers (\SGARefLink{sga7.sga7i.eviii.n1-6-2}{1.6.2}).

For a useful case in which the considerations of \SGARefLink{sga7.sga7i.eviii.n1-6}{1.6} apply, see \SGARefLink{sga7.sga7i.eviii.n3-7}{3.7}
below.

% END INCLUDED SGA7I_English_current_buildable_source_20260823/components/94_expose_VIII_section_1.tex

% BEGIN INCLUDED SGA7I_English_current_buildable_source_20260823/components/95_expose_VIII_section_2_1_opening_through_2_1_8.tex
% SOURCE AUTHORITY:
% expose_VIII_body.tex lines 279--382; scan indices 237--241; source folios 226--230.
% Direct scan comparison: physical PDF pages 238--242 at 1100 dpi.
% Transparent source dispositions:
% - typewritten letter-o homological indices are normalized to zero;
% - after (2.1.5), the second map is induced by the inclusion {}_nP -> P,
%   not the authority's impossible {}_nP -> Q;
% - the malformed two-argument Tor display in 2.1.8 is set with both arguments.
% Hard stop before 2.1.9, source line 384.

\subsection*{\SGARefTarget{sga7.sga7i.eviii.n2}{2}. Pairings defined by a biextension}

\SGARefTarget{sga7.sga7i.eviii.n2-1}{2.1}. Let $P$ and $Q$ be two abelian groups of $T$, and let $n>0$ be
an integer.  Recall that $\,{}_nP$ denotes the kernel of multiplication
by $n\cdot\id_P$.  We shall define a canonical homomorphism
\begin{equation}\tag{\SGARefTarget{sga7.sga7i.eviii.n2-1-1}{2.1.1}}
 \alpha=\alpha^n_{P,Q}:
 {}_nP\otimes{}_nQ\longrightarrow\uTor_1(P,Q),
\end{equation}
which will be the composite
\begin{equation}\tag{\SGARefTarget{sga7.sga7i.eviii.n2-1-2}{2.1.2}}
 {}_nP\otimes{}_nQ
 \longrightarrow \uTor_1({}_nP,{}_nQ)
 \longrightarrow \uTor_1(P,Q),
\end{equation}
where the first arrow is $\alpha^n_{{}_nP,{}_nQ}$ and the second is
induced by the inclusions $\,{}_nP\longrightarrow P$ and
$\,{}_nQ\longrightarrow Q$.  We are thus reduced to defining (\SGARefLink{sga7.sga7i.eviii.n2-1-1}{2.1.1})
when $nP=0$ and $nQ=0$, that is, when $P$ and $Q$ are
$\mathbb Z/n\mathbb Z$-modules.

Suppose only that $nQ=0$.  Taking a flat resolution $L_\bullet$ of
$P$ (as a $\mathbb Z$-module), one obtains canonical isomorphisms
\[
 P\mathbin{\overset{\mathbb L}{\otimes}_{\mathbb Z}}Q
 \simeq L_\bullet\otimes_{\mathbb Z}Q
 \simeq
 (L_\bullet\otimes_{\mathbb Z}\mathbb Z/n\mathbb Z)
 \otimes_{\mathbb Z/n\mathbb Z}Q,
\]
which may be written as an isomorphism in the derived category
\begin{equation}\tag{\SGARefTarget{sga7.sga7i.eviii.n2-1-3}{2.1.3}}
 P\mathbin{\overset{\mathbb L}{\otimes}_{\mathbb Z}}Q
 \simeq
 \left(P\mathbin{\overset{\mathbb L}{\otimes}_{\mathbb Z}}
       \mathbb Z/n\mathbb Z\right)
 \mathbin{\overset{\mathbb L}{\otimes}_{\mathbb Z/n\mathbb Z}}Q.
\end{equation}
This gives rise to the homological K\"unneth-type spectral sequence
(cf.\ SGA 1 IV 5.2)
\begin{equation}\tag{\SGARefTarget{sga7.sga7i.eviii.n2-1-4}{2.1.4}}
 \uTor^{\mathbb Z}_*(P,Q)
 \Longleftarrow
 E^2_{p,q}
 =
 \uTor^{\mathbb Z/n\mathbb Z}_p
 \bigl(\uTor^{\mathbb Z}_q(P,\mathbb Z/n\mathbb Z),Q\bigr).
\end{equation}
Recall also the canonical isomorphisms
\[
 \uTor^{\mathbb Z}_0(P,\mathbb Z/n\mathbb Z)
 =P_n\mathrel{\mathop:}=P/nP,
 \qquad
 \uTor^{\mathbb Z}_1(P,\mathbb Z/n\mathbb Z)\simeq{}_nP,
\]
and
\[
 \uTor^{\mathbb Z}_i(P,\mathbb Z/n\mathbb Z)=0
 \qquad\text{for }i\geq2.
\]
They follow by computing these objects from the canonical flat
resolution of $(\mathbb Z/n\mathbb Z)_T$:
\[
\begin{tikzcd}[column sep=large,row sep=large]
 0\arrow[r]&
 \mathbb Z_T\arrow[r,"n"]\arrow[d,equal]&
 \mathbb Z_T\arrow[r]\arrow[d,equal]&
 0\\
 &L_1&L_0&
\end{tikzcd}
\]
whose augmentation
$L_0\longrightarrow(\mathbb Z/n\mathbb Z)_T$ sends $1$ to $1$.
Consequently, the spectral sequence (\SGARefLink{sga7.sga7i.eviii.n2-1-4}{2.1.4}) reduces essentially to
the exact sequence
\begin{equation}\tag{\SGARefTarget{sga7.sga7i.eviii.n2-1-5}{2.1.5}}
\begin{split}
0\longrightarrow{}&
 \uTor^{\mathbb Z/n\mathbb Z}_2(P_n,Q)
 \longrightarrow{}_nP\otimes Q
 \xrightarrow{\alpha}\uTor^{\mathbb Z}_1(P,Q)\\
\longrightarrow{}&
 \uTor^{\mathbb Z/n\mathbb Z}_1(P_n,Q)
 \longrightarrow0,
\end{split}
\end{equation}
where $P_n=P\otimes_{\mathbb Z}\mathbb Z/n\mathbb Z=P/nP$.
The desired homomorphism $\alpha$ is the edge homomorphism
$E^2_{0,1}\longrightarrow\uTor^{\mathbb Z}_1$ in this exact sequence.
Since the sequence is plainly functorial in $P$ and $Q$, functoriality
in $P$ shows that the homomorphism in question is indeed the composite
\[
 {}_nP\otimes Q
 \longrightarrow\uTor_1({}_nP,Q)
 \longrightarrow\uTor_1(P,Q),
\]
where the first homomorphism is the $\alpha$ just defined, with
$(P,Q)$ replaced by $({}_nP,Q)$, and the second is induced by the
inclusion $\,{}_nP\longrightarrow P$.

\SGARefTarget{sga7.sga7i.eviii.n2-1-6}{2.1.6}. This therefore defines, for arbitrary $P$ and $Q$, a
homomorphism (\SGARefLink{sga7.sga7i.eviii.n2-1-1}{2.1.1}) satisfying the factorization condition (\SGARefLink{sga7.sga7i.eviii.n2-1-2}{2.1.2}).
Moreover, when $nQ=0$, the exact sequence (\SGARefLink{sga7.sga7i.eviii.n2-1-5}{2.1.5}) shows that $\alpha$
is an isomorphism whenever $P_n$ is flat over
$\mathbb Z/n\mathbb Z$.  This holds in particular if $P_n=0$, that
is, if $nP=P$ and $P$ is $n$-divisible, or if $P$ itself is a flat
$\mathbb Z/n\mathbb Z$-module.  It also holds if $Q$ is a flat
$\mathbb Z/n\mathbb Z$-module.

\SGARefTarget{sga7.sga7i.eviii.n2-1-7}{2.1.7}. Let us make the calculation of the homomorphism $\alpha$ in
(\SGARefLink{sga7.sga7i.eviii.n2-1-5}{2.1.5}) explicit.  To determine it, it suffices, for every section
$q$ of $Q$ (over an object $S$ of $T$; take the final object for
simplicity), to determine the corresponding homomorphism
\[
 {}_nP\longrightarrow\uTor^{\mathbb Z}_1(P,Q).
\]
The section $q$ may be regarded as the image of the section $1$ of
$(\mathbb Z/n\mathbb Z)_T$ under the uniquely determined
homomorphism
\[
 u_q:(\mathbb Z/n\mathbb Z)_T\longrightarrow Q.
\]
Functoriality of $\alpha$ in $Q$ therefore reduces us to the case
$Q=(\mathbb Z/n\mathbb Z)_T$, with $q$ its unit section.  In this
case, identify
\[
 \uTor^{\mathbb Z}_1(P,Q)
 =
 \uTor^{\mathbb Z}_1(P,\mathbb Z/n\mathbb Z)
\]
with $\,{}_nP$ as explained after (\SGARefLink{sga7.sga7i.eviii.n2-1-4}{2.1.4}).  The homomorphism in question
is then the identity.  Equivalently, the edge homomorphism
\[
 E^2_{0,q}
 =
 \uTor^{\mathbb Z}_q(P,\mathbb Z/n\mathbb Z)
 \otimes_{\mathbb Z/n\mathbb Z}\mathbb Z/n\mathbb Z
 \longrightarrow
 \uTor^{\mathbb Z}_q(P,\mathbb Z/n\mathbb Z)
\]
is the identity.  This follows immediately from the definitions.
The analogous assertion remains valid when $\mathbb Z$ and its
quotient $\mathbb Z/n\mathbb Z$ are replaced by any ring $A$ on $T$
and any quotient ring $A/I$.

\SGARefTarget{sga7.sga7i.eviii.n2-1-8}{2.1.8}. The homomorphism $\alpha$ of (\SGARefLink{sga7.sga7i.eviii.n2-1-1}{2.1.1}) may also be determined by
describing, for every pair of sections $p$ of $\,{}_nP$ and $q$ of
$\,{}_nQ$ (over an object $S$ of $T$, again taken to be final to
simplify the notation), the section of
$\uTor^{\mathbb Z}_1(P,Q)$ determined by $(p,q)$.  By functoriality in
$P$ and $Q$, as above, we are reduced to the case
\[
 P=Q=(\mathbb Z/n\mathbb Z)_T,
\]
with $p$ and $q$ the unit section.  Thus we need only determine
\[
 \alpha:
 (\mathbb Z/n\mathbb Z)_T\otimes(\mathbb Z/n\mathbb Z)_T
 \longrightarrow
 \uTor^{\mathbb Z}_1
 \bigl((\mathbb Z/n\mathbb Z)_T,(\mathbb Z/n\mathbb Z)_T\bigr).
\]
The description in \SGARefLink{sga7.sga7i.eviii.n2-1-7}{2.1.7} shows that this is the isomorphism
\begin{equation}\tag{\SGARefTarget{sga7.sga7i.eviii.n2-1-8-1}{2.1.8.1}}
 (\mathbb Z/n\mathbb Z)_T
 \xrightarrow{\sim}
 \uTor^{\mathbb Z}_1
 \bigl((\mathbb Z/n\mathbb Z)_T,(\mathbb Z/n\mathbb Z)_T\bigr)
\end{equation}
deduced from the canonical flat resolution
\begin{equation}\tag{\SGARefTarget{sga7.sga7i.eviii.n2-1-8-2}{2.1.8.2}}
 0\longrightarrow\mathbb Z_T
 \xrightarrow{n}\mathbb Z_T
 \longrightarrow0
\end{equation}
of the \underline{first} factor
$P=(\mathbb Z/n\mathbb Z)_T$ occurring in $\uTor_1$.

% END INCLUDED SGA7I_English_current_buildable_source_20260823/components/95_expose_VIII_section_2_1_opening_through_2_1_8.tex

% BEGIN INCLUDED SGA7I_English_current_buildable_source_20260823/components/96_expose_VIII_section_2_1_9_through_2_1_12.tex
% SOURCE AUTHORITY:
% expose_VIII_body.tex lines 384--547; scan indices 241--244; source folios 230--233.
% Direct scan comparison: physical PDF pages 242--245 at 1100 dpi.
% Transparent source dispositions:
% - typewritten letter-o complex indices are normalized to zero;
% - broken internal references 2.11.1/2.11.2 are restored as 2.1.11.1/2.1.11.2;
% - the upper vertex in the proof of 2.1.11 is set as P/nP, as required by
%   the identification in the prose and the four maps.
% Hard stop before 2.1.13, source line 552.

\SGARefTarget{sga7.sga7i.eviii.n2-1-9}{2.1.9}. The definition given in (\SGARefLink{sga7.sga7i.eviii.n2-1-1}{2.1.1}) is not symmetric.  By
interchanging the roles of $P$ and $Q$ in the preceding construction,
one may define an analogous homomorphism
\begin{equation}\tag{\SGARefTarget{sga7.sga7i.eviii.n2-1-9-1}{2.1.9.1}}
 \beta^n_{P,Q}:{}_nP\otimes{}_nQ\longrightarrow\uTor_1(P,Q),
\end{equation}
first reducing to the case $nP=0$ and then using a flat resolution
\underline{of} $Q$.  Equivalently, $\beta$ may be defined in terms of
$\alpha$ by requiring the following diagram to commute:
\begin{equation}\tag{\SGARefTarget{sga7.sga7i.eviii.n2-1-9-2}{2.1.9.2}}
\begin{tikzcd}[column sep=huge,row sep=huge]
 {}_nP\otimes{}_nQ
   \arrow[r,"\sim"]
   \arrow[d,"{\beta^n_{P,Q}}"']&
 {}_nQ\otimes{}_nP
   \arrow[d,"{\alpha^n_{Q,P}}"]\\
 \uTor_1(P,Q)\arrow[r,"\sim"]&
 \uTor_1(Q,P).
\end{tikzcd}
\end{equation}
Here the horizontal arrows are the symmetry isomorphisms.  With this
convention, the relation between $\alpha$ and $\beta$ is
\begin{equation}\tag{\SGARefTarget{sga7.sga7i.eviii.n2-1-9-3}{2.1.9.3}}
 \beta^n_{P,Q}=-\alpha^n_{P,Q}.
\end{equation}
In other words:

\medskip
\noindent\underline{Lemma \SGARefTarget{sga7.sga7i.eviii.n2-1-10}{2.1.10}.}
\underline{Let $P$ and $Q$ be two commutative groups and let $n>0$ be
an integer.  Then the following diagram is anticommutative, where
$\alpha$ is defined in (\SGARefLink{sga7.sga7i.eviii.n2-1-1}{2.1.1}):}
\begin{equation}\tag{\SGARefTarget{sga7.sga7i.eviii.n2-1-10-1}{2.1.10.1}}
\begin{tikzcd}[column sep=huge,row sep=huge]
 {}_nP\otimes{}_nQ
   \arrow[r,"\sim"]
   \arrow[d,"{\alpha^n_{P,Q}}"']&
 {}_nQ\otimes{}_nP
   \arrow[d,"{\alpha^n_{Q,P}}"]\\
 \uTor_1(P,Q)
   \arrow[r,"\sim"]
   \arrow[ur,phantom,"-" description]&
 \uTor_1(Q,P).
\end{tikzcd}
\end{equation}
\underline{The horizontal arrows are the symmetry isomorphisms.}

\medskip
\noindent\underline{Proof.}
The procedure of \SGARefLink{sga7.sga7i.eviii.n2-1-8}{2.1.8} immediately reduces us to the case
$P=Q=(\mathbb Z/n\mathbb Z)_T$.  We must then prove that the two
isomorphisms (\SGARefLink{sga7.sga7i.eviii.n2-1-8-1}{2.1.8.1}), obtained by resolving respectively the first
and second arguments of $\uTor_1$ by (\SGARefLink{sga7.sga7i.eviii.n2-1-8-2}{2.1.8.2}), differ by a sign
(equivalently, that the symmetry isomorphism on
$\uTor^{\mathbb Z}_1(\mathbb Z/n\mathbb Z,\mathbb Z/n\mathbb Z)$ is
$-\id$).  We may plainly suppose that $T$ is the punctual topos, so
that we work with ordinary commutative groups.  Let $L_\bullet$ be the
complex
\[
 L_\bullet=\bigl(0\longrightarrow\mathbb Z
 \xrightarrow{n}\mathbb Z\longrightarrow0\bigr),
\]
regarded as a flat resolution of
$P=Q=\mathbb Z/n\mathbb Z$.  Use
$L_\bullet\otimes L_\bullet$ to compute
\[
 \Tor_1(P,Q)=\Tor_1(\mathbb Z/n\mathbb Z,\mathbb Z/n\mathbb Z).
\]
The isomorphisms with
$H_1(L_\bullet\otimes\mathbb Z/n\mathbb Z)$, which defines $\alpha$,
and with
$H_1(\mathbb Z/n\mathbb Z\otimes L_\bullet)$, which defines $\beta$,
are induced respectively by the augmentations
$L_\bullet\longrightarrow\mathbb Z/n\mathbb Z$ of the second and
first factors of $L_\bullet\otimes L_\bullet$.  It follows at once
that $\alpha(1,1)$ is represented by the cycle
\[
 -1_0\otimes1_1+1_1\otimes1_0
 \in(L_\bullet\otimes L_\bullet)_1
 =L_0\otimes L_1+L_1\otimes L_0,
\]
where $1_0$ and $1_1$ denote the element $1$ of $L_0$ and $L_1$
respectively, whereas $\beta(1,1)$ is represented by
$1_0\otimes1_1-1_1\otimes1_0$.  This proves \SGARefLink{sga7.sga7i.eviii.n2-1-10}{2.1.10}.

To finish these generalities on the pairings $\alpha$, we must still
specify the relations among the homomorphisms $\alpha^n_{P,Q}$ as
$n$ varies.  The main result is the following.

\medskip
\noindent\underline{Lemma \SGARefTarget{sga7.sga7i.eviii.n2-1-11}{2.1.11}.}
\underline{Let $P,Q$ be two commutative groups of $T$, and let $n,n'$
be positive integers such that $n'$ is a multiple of $n$, say
$n'=mn$.  Then the following diagram is commutative:}
\begin{equation}\tag{\SGARefTarget{sga7.sga7i.eviii.n2-1-11-1}{2.1.11.1}}
\begin{tikzcd}[column sep=large,row sep=large]
 &{}_{n'}P\otimes{}_nQ
   \arrow[ld,"P_m\otimes\id_{{}_nQ}"']
   \arrow[rd,"\id_{{}_{n'}P}\otimes i_m"]&\\
 {}_nP\otimes{}_nQ
   \arrow[rd,"{\alpha^n_{P,Q}}"']&&
 {}_{n'}P\otimes{}_{n'}Q
   \arrow[ld,"{\alpha^{n'}_{P,Q}}"]\\
 &\uTor_1(P,Q)&
\end{tikzcd}
\end{equation}
\underline{Here $P_m:{}_{n'}P\longrightarrow{}_nP$ is the homomorphism
induced by $m\cdot\id_P$, and
$i_m:{}_nQ\longrightarrow{}_{n'}Q$ is the inclusion.  Likewise, the
diagram (\SGARefTarget{sga7.sga7i.eviii.n2-1-11-2}{2.1.11.2}) symmetric to (\SGARefLink{sga7.sga7i.eviii.n2-1-11-1}{2.1.11.1}) is commutative; it has the
same vertices except that the upper vertex
$\,{}_{n'}P\otimes{}_nQ$ is replaced by
$\,{}_nP\otimes{}_{n'}Q$.}

\medskip
It suffices to prove the first assertion; the second follows by
symmetry in view of \SGARefLink{sga7.sga7i.eviii.n2-1-10}{2.1.10}.  By the procedure of \SGARefLink{sga7.sga7i.eviii.n2-1-8}{2.1.8}, we are reduced
to the punctual topos, where $P$ and $Q$ are ordinary commutative
groups and
\[
 P=\mathbb Z/n'\mathbb Z,\qquad Q=\mathbb Z/n'\mathbb Z.
\]
In fact, for $P$ we use only the hypothesis $n'P=0$.  Using the
resolution
\[
 0\longrightarrow\mathbb Z\xrightarrow{n'}\mathbb Z
 \longrightarrow0
\]
of $Q$, identify $\Tor_1(P,Q)$ with $\,{}_{n'}P=P$.  Likewise, using
\[
 0\longrightarrow\mathbb Z\xrightarrow n\mathbb Z
 \longrightarrow0
\]
to resolve
$\,{}_nQ\simeq\mathbb Z/n\mathbb Z$, identify
$\Tor_1({}_nP,{}_nQ)$ with $\,{}_nP$.  The augmentation in the latter
resolution sends $1_0$ to $m\cdot1_Q$, where $1_Q$ denotes the element
$1$ of $Q=\mathbb Z/n'\mathbb Z$.  Under these identifications the
homomorphism
\[
\begin{tikzcd}[column sep=large,row sep=large]
 \Tor_1({}_nP,{}_nQ)\arrow[r]\arrow[d,no head,"\wr"']&
 \Tor_1(P,Q)\arrow[d,no head,"\wr"']\\
 {}_nP&{}_{n'}P=P
\end{tikzcd}
\]
induced by the inclusions
$\,{}_nP\hookrightarrow P$ and $\,{}_nQ\hookrightarrow Q$ is simply
the natural inclusion $\,{}_nP\hookrightarrow{}_{n'}P$.

Consequently, identify the upper vertex of (\SGARefLink{sga7.sga7i.eviii.n2-1-11-1}{2.1.11.1}) with $P/nP$
using the generator $m\cdot1_Q$ of $\,{}_nQ$, which identifies
$\,{}_nQ$ with $\mathbb Z/n\mathbb Z$.  The diagram becomes
\[
\begin{tikzcd}[column sep=large,row sep=large]
 &P/nP
   \arrow[ld,"x\longmapsto mx"']
   \arrow[rd,"x\longmapsto mx"]&\\
 {}_nP\arrow[rd,"x\longmapsto x"']&&
 P\arrow[ld,"x\longmapsto x"]\\
 &{}_{n'}P&
\end{tikzcd}
\]
which is plainly commutative.

\medskip
\noindent\underline{Corollary \SGARefTarget{sga7.sga7i.eviii.n2-1-12}{2.1.12}.}
\underline{With the notation of \SGARefLink{sga7.sga7i.eviii.n2-1-11}{2.1.11}, the following diagram is
commutative:}
\begin{equation}\tag{\SGARefTarget{sga7.sga7i.eviii.n2-1-12-1}{2.1.12.1}}
\begin{tikzcd}[column sep=huge,row sep=huge]
 {}_{n'}P\otimes{}_{n'}Q
   \arrow[r,"m\otimes m"]
   \arrow[d,"{\alpha^{n'}_{P,Q}}"']&
 {}_nP\otimes{}_nQ
   \arrow[d,"{\alpha^n_{P,Q}}"]\\
 {}_{n'}\uTor_1(P,Q)
   \arrow[r,"m"]&
 {}_n\uTor_1(P,Q).
\end{tikzcd}
\end{equation}
\underline{In the description of the horizontal arrows, $m$ denotes
the homomorphisms $\,{}_{n'}E\longrightarrow{}_nE$ induced by
$m\cdot\id_E$, for $E=P,Q$, and $\uTor_1(P,Q)$ respectively.}

\medskip
Indeed, by \SGARefLink{sga7.sga7i.eviii.n2-1-11}{2.1.11}, for
$x\in{}_{n'}P(S)$ and $y\in{}_{n'}Q(S)$ one has
\[
 \alpha^n(mx\otimes my)=\alpha^{n'}(x\otimes my).
\]
The right-hand side equals
$m\alpha^{n'}(x\otimes y)$, since
$x\otimes my=m(x\otimes y)$.  This proves the commutativity of
(\SGARefLink{sga7.sga7i.eviii.n2-1-12-1}{2.1.12.1}).

% END INCLUDED SGA7I_English_current_buildable_source_20260823/components/96_expose_VIII_section_2_1_9_through_2_1_12.tex

% BEGIN INCLUDED SGA7I_English_current_buildable_source_20260823/components/97_expose_VIII_section_2_1_13_and_remark_2_1_14.tex
% SOURCE AUTHORITY:
% expose_VIII_body.tex lines 552--675; scan indices 245--248; source folios 234--237.
% Direct scan comparison: physical PDF pages 246--249 at 1100 dpi.
% Diagram 2.1.14.6: six targeted vector crops at 5000 dpi, covering every
% node, arrow, label, attachment, and terminal punctuation.
% Transparent source dispositions:
% - the authority's self-contradictory phrase "injective isomorphism" in
%   2.1.14 is rendered as "injective homomorphism";
% - in the multiplication-by-pq map following (2.1.14.6), the authority
%   omits pq before the second ideal in the target numerator; it is restored
%   as pq a \cap pq b, as required by multiplication by pq and by the diagram.
% Hard stop before 2.2, source line 676.

\SGARefTarget{sga7.sga7i.eviii.n2-1-13}{2.1.13}. Let $\ell$ be a prime number.  For every commutative group
$E$, denote by $T_\ell(E)$ the inverse system
\begin{equation}\tag{\SGARefTarget{sga7.sga7i.eviii.n2-1-13-1}{2.1.13.1}}
 T_\ell(E)=({}_{\ell^\nu}E)_{\nu\geq0},
\end{equation}
which plainly depends functorially on $E$.  The compatibilities of
\SGARefLink{sga7.sga7i.eviii.n2-1-12}{2.1.12}, for $n,n'$ of the form $\ell^\nu,\ell^{\nu'}$ with
$\nu'\geq\nu$, show that the homomorphisms
$\alpha^{\ell^\nu}_{P,Q}$, as $\nu$ varies, define a homomorphism of
inverse systems, functorial in $P,Q$:
\begin{equation}\tag{\SGARefTarget{sga7.sga7i.eviii.n2-1-13-2}{2.1.13.2}}
 \alpha^{(\ell)}_{P,Q}:
 T_\ell(P)\otimes T_\ell(Q)
 \longrightarrow T_\ell\bigl(\uTor_1(P,Q)\bigr).
\end{equation}
Knowledge of these homomorphisms for every prime $\ell$ is plainly
equivalent to knowledge of the $\alpha^n_{P,Q}$ for every integer
$n$.  The antisymmetry relations of \SGARefLink{sga7.sga7i.eviii.n2-1-10}{2.1.10} then become the analogous
$\ell$-adic antisymmetry relations, for every $\ell$, expressed by
the anticommutativity of the diagrams
\begin{equation}\tag{\SGARefTarget{sga7.sga7i.eviii.n2-1-13-3}{2.1.13.3}}
\begin{tikzcd}[column sep=huge,row sep=huge]
 T_\ell(P)\otimes T_\ell(Q)
   \arrow[r]
   \arrow[d,"{\alpha^{(\ell)}_{P,Q}}"']
   \arrow[rd,phantom,"-" description]&
 T_\ell(Q)\otimes T_\ell(P)
   \arrow[d,"{\alpha^{(\ell)}_{Q,P}}"]\\
 T_\ell\bigl(\uTor_1(P,Q)\bigr)
   \arrow[r]&
 T_\ell\bigl(\uTor_1(Q,P)\bigr).
\end{tikzcd}
\end{equation}
Here the horizontal arrows are the symmetry arrows.

\medskip
\noindent\underline{Remark \SGARefTarget{sga7.sga7i.eviii.n2-1-14}{2.1.14}.}
The construction of the homomorphism (\SGARefLink{sga7.sga7i.eviii.n2-1-1}{2.1.1}), and the results \SGARefLink{sga7.sga7i.eviii.n2-1-10}{2.1.10}
and \SGARefLink{sga7.sga7i.eviii.n2-1-11}{2.1.11} concerning it, may be generalized to the situation in
which one is given a commutative ring $\underline A$ on $T$ and two
ideals $\underline a$ and $\underline b$ of $\underline A$.  In the
case treated above, both ideals are the ideal
$n\mathbb Z_T$ of the constant ring $\mathbb Z_T$ on $T$.  One then
defines a homomorphism, functorial in the $\underline A$-modules
$P,Q$,
\begin{equation}\tag{\SGARefTarget{sga7.sga7i.eviii.n2-1-14-1}{2.1.14.1}}
 {}_{\underline a}P\otimes{}_{\underline b}Q
 \xrightarrow{\ \alpha\ }
 \uHom_{\underline A}
 \left(
   \frac{\underline a\cap\underline b}
        {\underline a\mathbin{\cdot}\underline b},
   \uTor^{\underline A}_1(P,Q)
 \right),
\end{equation}
as the composite of the homomorphisms
\[
\begin{split}
 {}_{\underline a}P\otimes{}_{\underline b}Q
 &\longrightarrow
 \uHom(\underline A/\underline a,P)
 \otimes
 \uHom(\underline A/\underline b,Q)\\
 &\longrightarrow
 \uHom\!\left(
   \uTor_1(\underline A/\underline a,\underline A/\underline b),
   \uTor_1(P,Q)
 \right),
\end{split}
\]
using the well-known canonical isomorphism
\begin{equation}\tag{\SGARefTarget{sga7.sga7i.eviii.n2-1-14-2}{2.1.14.2}}
 \uTor_1(\underline A/\underline a,\underline A/\underline b)
 \simeq
 \frac{\underline a\cap\underline b}
      {\underline a\mathbin{\cdot}\underline b}.
\end{equation}
In fact, the sign in \SGARefLink{sga7.sga7i.eviii.n2-1-10}{2.1.10} is ``hidden'' in the choice of this
isomorphism, according as one works with the resolution
$0\longrightarrow\underline a\longrightarrow\underline A
\longrightarrow0$ of $\underline A/\underline a$, or with the
resolution $0\longrightarrow\underline b\longrightarrow\underline A
\longrightarrow0$ of $\underline A/\underline b$.  Here it becomes
the anticommutativity of the diagram
\begin{equation}\tag{\SGARefTarget{sga7.sga7i.eviii.n2-1-14-3}{2.1.14.3}}
\begin{tikzcd}[column sep=huge,row sep=huge]
 {}_{\underline a}P\otimes{}_{\underline b}Q
   \arrow[r]
   \arrow[d]
   \arrow[rd,phantom,"-1" description]&
 {}_{\underline b}Q\otimes{}_{\underline a}P
   \arrow[d]\\
 \uHom\!\left(
   \dfrac{\underline a\cap\underline b}
         {\underline a\mathbin{\cdot}\underline b},
   \uTor_1(P,Q)
 \right)
   \arrow[r]&
 \uHom\!\left(
   \dfrac{\underline a\cap\underline b}
         {\underline a\mathbin{\cdot}\underline b},
   \uTor_1(Q,P)
 \right).
\end{tikzcd}
\end{equation}
The horizontal arrows are the symmetry arrows.  When
$\underline a=\underline b$, this anticommutativity may also be
expressed by saying that the symmetry automorphism on
$\uTor_1(\underline A/\underline a,\underline A/\underline a)$ is
$-\id$.

In the special case where
$\underline a=\underline b=n\underline A$, with $n$ a section of
$\underline A$, multiplication by $n$ induces an epimorphism
\[
 \underline A/\underline a
 \longrightarrow
 \frac{\underline a\cap\underline a}
      {\underline a\mathbin{\cdot}\underline a}
 =\underline a/\underline a^2,
\]
and hence a functorial injective homomorphism, in the
$\underline A$-module $T$,
\[
 \uHom(\underline a/\underline a^2,T)
 \longrightarrow{}_{\underline a}T.
\]
Thus in this case the homomorphism (\SGARefLink{sga7.sga7i.eviii.n2-1-14-1}{2.1.14.1}) may be interpreted as a
homomorphism
\begin{equation}\tag{\SGARefTarget{sga7.sga7i.eviii.n2-1-14-4}{2.1.14.4}}
 {}_nP\otimes{}_nQ
 \longrightarrow{}_n\uTor_1(P,Q),
\end{equation}
which is closer in form to (\SGARefLink{sga7.sga7i.eviii.n2-1-1}{2.1.1}).

Continue to suppose that
$\underline a=\underline b=n\underline A$, and that $n$ is a regular
section of $\underline A$.  The spectral sequence (\SGARefLink{sga7.sga7i.eviii.n2-1-4}{2.1.4}) still makes
sense in the present setting and yields the following exact sequence
of low-degree terms, generalizing (\SGARefLink{sga7.sga7i.eviii.n2-1-5}{2.1.5}), where $Q$ is now an
$\underline A$-module annihilated by $\underline a$:
\begin{equation}\tag{\SGARefTarget{sga7.sga7i.eviii.n2-1-14-5}{2.1.14.5}}
\begin{split}
 \uTor^{\underline A}_2(P,Q)
 &\longrightarrow
 \uTor^{\underline A/n\underline A}_2(P_n,Q)
 \longrightarrow{}_nP\otimes Q
 \xrightarrow{\ \alpha\ }
 \uTor^{\underline A}_1(P,Q)\\
 &\longrightarrow
 \uTor^{\underline A/n\underline A}_1(P_n,Q)
 \longrightarrow0.
\end{split}
\end{equation}

Return now to arbitrary $\underline a$ and $\underline b$, and let
$p$ and $q$ be two sections of $\underline A$.  We then obtain the
commutative diagram below; (\SGARefLink{sga7.sga7i.eviii.n2-1-11-1}{2.1.11.1}) is essentially the special case
obtained by taking $q=1$:
\begin{equation}\tag{\SGARefTarget{sga7.sga7i.eviii.n2-1-14-6}{2.1.14.6}}
{\footnotesize
\begin{tikzcd}[row sep=large,column sep=tiny]
 &{}_{p\underline a}P\otimes{}_{q\underline b}Q
   \arrow[dl,"p\otimes q"']
   \arrow[dr,"i\otimes j"]&\\
 {}_{\underline a}P\otimes{}_{\underline b}Q
   \arrow[d,"\alpha"']&&
 {}_{pq\underline a}P\otimes{}_{pq\underline b}Q
   \arrow[d,"\alpha"]\\
 \uHom\!\left(
   \dfrac{\underline a\cap\underline b}
         {\underline a\mathbin{\cdot}\underline b},
   \uTor_1(P,Q)
 \right)
   \arrow[dr,"\mathrm{incl.}"']&&
 \uHom\!\left(
   \dfrac{pq\underline a\cap pq\underline b}
         {pq\underline a\mathbin{\cdot}pq\underline b},
   \uTor_1(P,Q)
 \right)
   \arrow[dl,"(pq)^*"]\\
 &\uHom\!\left(
   \dfrac{\underline a\cap\underline b}
         {pq\underline a\mathbin{\cdot}\underline b},
   \uTor_1(P,Q)
 \right)&
\end{tikzcd}}
\end{equation}
Here
\[
 i:{}_{p\underline a}P\longrightarrow{}_{pq\underline a}P,
 \qquad
 j:{}_{q\underline b}Q\longrightarrow{}_{pq\underline b}Q
\]
are the canonical inclusions, and $(pq)^*$ denotes the injective
homomorphism obtained, by passing to $\uHom$, from the surjective
homomorphism
\[
 \frac{\underline a\cap\underline b}
      {pq\underline a\mathbin{\cdot}\underline b}
 \longrightarrow
 \frac{pq\underline a\cap pq\underline b}
      {pq\underline a\mathbin{\cdot}pq\underline b}
\]
induced by multiplication by $pq$.  When
$\underline a=\underline b=n\underline A$, this diagram may be
rewritten in a simpler form, given in (\SGARefLink{sga7.sga7i.eviii.n2-1-11-1}{2.1.11.1}) when $q=1$.

% END INCLUDED SGA7I_English_current_buildable_source_20260823/components/97_expose_VIII_section_2_1_13_and_remark_2_1_14.tex

% BEGIN INCLUDED SGA7I_English_current_buildable_source_20260823/components/98_expose_VIII_section_2_2.tex
% SOURCE AUTHORITY:
% expose_VIII_body.tex lines 676--834; scan indices 248--252; source folios 237--241.
% Direct scan comparison: physical PDF pages 249--253 at 1100 dpi.
% Transparent source dispositions:
% - in 2.2.7, the restriction is to the image of alpha^n_{P,Q} in
%   Tor_1(P,Q), not the authority's impossible phi-map to Tor_1(P,Q;G);
% - the cited pairings are (2.2.5) and (2.2.6), not (2.2.4) and (2.2.6);
% - typewritten letter-o complex subscripts are normalized to zero;
% - the missing closing parenthesis in the lower-left node of (2.2.11.2)
%   is restored.
% Hard stop before 2.3, source line 836.

\SGARefTarget{sga7.sga7i.eviii.n2-2}{2.2}. Consider now, in addition to $P$ and $Q$, a commutative group
$G$.  In (\SGARefLink{sga7.sga7i.eviii.n1-1-3}{1.1.3}) we defined a canonical homomorphism, functorial in
$P,Q$, and $G$,
\begin{equation}\tag{\SGARefTarget{sga7.sga7i.eviii.n2-2-1}{2.2.1}}
 \mathrm{Biext}^1(P,Q;G)
 \longrightarrow
 \operatorname{Hom}\bigl(\operatorname{Tor}_1(P,Q),G\bigr).
\end{equation}
On the other hand, if $n$ is a positive integer (respectively, if
$\ell$ is a prime number), the homomorphism $\alpha^n_{P,Q}$ of
(\SGARefLink{sga7.sga7i.eviii.n2-1-1}{2.1.1}) (respectively, the homomorphism $\alpha^{(\ell)}_{P,Q}$ of
(\SGARefLink{sga7.sga7i.eviii.n2-1-13-2}{2.1.13.2})) gives a homomorphism
\[
 \operatorname{Hom}\bigl(\operatorname{Tor}_1(P,Q),G\bigr)
 \longrightarrow
 \operatorname{Hom}({}_nP\otimes{}_nQ,G)
\]
(respectively, a homomorphism
\[
 \operatorname{Hom}\bigl(\operatorname{Tor}_1(P,Q),G\bigr)
 \longrightarrow
 \operatorname{Hom}\bigl(
   T_\ell(P)\otimes T_\ell(Q),T_\ell(G)
 \bigr)).
\]
These last homomorphisms are also functorial in $P,Q,G$.
Composing them with (\SGARefLink{sga7.sga7i.eviii.n2-2-1}{2.2.1}), we obtain canonical homomorphisms,
functorial in $P,Q,G$,
\begin{equation}\tag{\SGARefTarget{sga7.sga7i.eviii.n2-2-2}{2.2.2}}
 \varphi^n=\varphi^n_{P,Q}:
 \mathrm{Biext}^1(P,Q;G)
 \longrightarrow
 \operatorname{Hom}({}_nP\otimes{}_nQ,G),
\end{equation}
\begin{equation}\tag{\SGARefTarget{sga7.sga7i.eviii.n2-2-3}{2.2.3}}
 \varphi^{(\ell)}=\varphi^{(\ell)}_{P,Q}:
 \mathrm{Biext}^1(P,Q;G)
 \longrightarrow
 \operatorname{Hom}\bigl(
   T_\ell(P)\otimes T_\ell(Q),T_\ell(G)
 \bigr).
\end{equation}
For
\begin{equation}\tag{\SGARefTarget{sga7.sga7i.eviii.n2-2-4}{2.2.4}}
 \xi\in\mathrm{Biext}^1(P,Q;G),
\end{equation}
denote its image under (\SGARefLink{sga7.sga7i.eviii.n2-2-2}{2.2.2}) (respectively, under (\SGARefLink{sga7.sga7i.eviii.n2-2-3}{2.2.3})) by
$\varphi^n_\xi$ (respectively, by $\varphi^{(\ell)}_\xi$):
\begin{equation}\tag{\SGARefTarget{sga7.sga7i.eviii.n2-2-5}{2.2.5}}
 \varphi^n_\xi:{}_nP\otimes{}_nQ
 \longrightarrow G
 \qquad\text{(or, equivalently, }{}_nG\text{)},
\end{equation}
\begin{equation}\tag{\SGARefTarget{sga7.sga7i.eviii.n2-2-6}{2.2.6}}
 \varphi^{(\ell)}_\xi:
 T_\ell(P)\otimes T_\ell(Q)
 \longrightarrow T_\ell(G).
\end{equation}

\SGARefTarget{sga7.sga7i.eviii.n2-2-7}{2.2.7}. Plainly, knowledge of $\varphi^{(\ell)}_\xi$ is equivalent
to knowledge of the $\varphi^n_\xi$ for all $n=\ell^\nu$, with
$\nu\geq0$.  For fixed $n$, knowledge of $\varphi^n_\xi$ is
equivalent to knowledge, on the image of
\[
 \alpha^n_{P,Q}:{}_nP\otimes{}_nQ\longrightarrow\uTor_1(P,Q),
\]
of the homomorphism
\begin{equation}\tag{\SGARefTarget{sga7.sga7i.eviii.n2-2-8}{2.2.8}}
 \uTor_1(P,Q)\longrightarrow G
\end{equation}
which is the image of $\xi$ under (\SGARefLink{sga7.sga7i.eviii.n2-2-1}{2.2.1}).

Observe that if $P=nP$ and $Q=nQ$, then the preceding image is
$\,{}_n\uTor_1(P,Q)$.  This follows readily by considering the exact
sequence (\SGARefLink{sga7.sga7i.eviii.n1-3-5}{1.3.5}) and the analogous exact sequence for $Q$.
Furthermore, if $P$ and $Q$ are divisible, that is, if $nP=P$ and
$nQ=Q$ for every positive integer $n$, one readily checks that
$\uTor_1(P,Q)$ is torsion, that is,
\[
 \uTor_1(P,Q)=\varinjlim_n{}_{n}\uTor_1(P,Q).
\]
Thus, in this case, knowledge of the $\varphi^n_\xi$ for all $n$ is
equivalent to knowledge of (\SGARefLink{sga7.sga7i.eviii.n2-2-8}{2.2.8}), whereas knowledge of
$\varphi^{(\ell)}_\xi$ for a given prime $\ell$ is equivalent to
knowledge of the restriction of (\SGARefLink{sga7.sga7i.eviii.n2-2-8}{2.2.8}) to the $\ell$-primary
component of $\uTor_1(P,Q)$.

We now make explicit the symmetry properties of the pairings (\SGARefLink{sga7.sga7i.eviii.n2-2-5}{2.2.5})
and (\SGARefLink{sga7.sga7i.eviii.n2-2-6}{2.2.6}) associated with biextensions.  First observe that there is
a commutative diagram
\begin{equation}\tag{\SGARefTarget{sga7.sga7i.eviii.n2-2-9}{2.2.9}}
\begin{tikzcd}[column sep=large,row sep=huge]
 \mathrm{Biext}^1(P,Q;G)
   \arrow[r,"\sim"]
   \arrow[d]&
 \mathrm{Biext}^1(Q,P;G)
   \arrow[d]\\
 \operatorname{Hom}\bigl(\uTor_1(P,Q),G\bigr)
   \arrow[r,"\sim"]&
 \operatorname{Hom}\bigl(\uTor_1(Q,P),G\bigr),
\end{tikzcd}
\end{equation}
where the horizontal arrows are the symmetry arrows (the first of
which was defined in VII \SGARefLink{sga7.sga7i.evii.n2-7}{2.7}), and the vertical arrows are (\SGARefLink{sga7.sga7i.eviii.n2-2-1}{2.2.1}).
The commutativity follows from that of the two constituent squares in
the fuller diagram
\begin{equation}\tag{\SGARefTarget{sga7.sga7i.eviii.n2-2-10}{2.2.10}}
\begin{tikzcd}[column sep=large,row sep=large]
 \mathrm{Biext}^1(P,Q;G)
   \arrow[r,"\sim"]
   \arrow[d]&
 \mathrm{Biext}^1(Q,P;G)
   \arrow[d]\\
 \operatorname{Ext}^1\!\left(
   P\mathbin{\overset{\mathbb L}{\otimes}}Q,G
 \right)
   \arrow[r,"\sim"]
   \arrow[d]&
 \operatorname{Ext}^1\!\left(
   Q\mathbin{\overset{\mathbb L}{\otimes}}P,G
 \right)
   \arrow[d]\\
 \operatorname{Hom}\bigl(\uTor_1(P,Q),G\bigr)
   \arrow[r,"\sim"]&
 \operatorname{Hom}\bigl(\uTor_1(Q,P),G\bigr).
\end{tikzcd}
\end{equation}
The first vertical arrow in each column is the canonical isomorphism
of VII \SGARefLink{sga7.sga7i.evii.n3-6-5}{3.6.5}, and the second is the evident homomorphism.  The second
square in (\SGARefLink{sga7.sga7i.eviii.n2-2-10}{2.2.10}) is trivially commutative: by definition, the
symmetry isomorphism
\[
 P\mathbin{\overset{\mathbb L}{\otimes}}Q
 \xrightarrow{\ \sim\ }
 Q\mathbin{\overset{\mathbb L}{\otimes}}P
\]
induces the symmetry isomorphism
$\uTor_1(P,Q)\simeq\uTor_1(Q,P)$, and, for a variable complex
$L_\bullet$, the canonical isomorphism
\[
 \operatorname{Ext}^1(L_\bullet,G)
 \simeq
 \operatorname{Hom}\bigl(H_1(L_\bullet),G\bigr)
\]
is functorial in $L_\bullet$.

As for the commutativity of the first square of (\SGARefLink{sga7.sga7i.eviii.n2-2-10}{2.2.10}), it follows
by making explicit simultaneously the symmetry isomorphisms and the
vertical isomorphisms defined in VII \SGARefLink{sga7.sga7i.evii.n3-6-5}{3.6.5}.  We leave the details of
the verification to the reader.  One uses the fact that the symmetry
isomorphism
\[
 L_\bullet(P)\otimes L_\bullet(Q)
 \simeq
 L_\bullet(Q)\otimes L_\bullet(P)
\]
induces in degree $1$ the isomorphism
\[
 L_0(P)\otimes L_1(Q)+L_1(P)\otimes L_0(Q)
 \simeq
 L_0(Q)\otimes L_1(P)+L_1(Q)\otimes L_0(P),
\]
which, on corresponding summands, is given by the ordinary symmetry
isomorphism of the underlying ungraded modules (and hence introduces
no sign).

In view of the commutativity of (\SGARefLink{sga7.sga7i.eviii.n2-2-9}{2.2.9}), the anticommutativity of
(\SGARefLink{sga7.sga7i.eviii.n2-1-10-1}{2.1.10.1}), in terms of the $\varphi^n$ or
$\varphi^{(\ell)}$, becomes the following relation.

\medskip
\noindent\underline{Proposition \SGARefTarget{sga7.sga7i.eviii.n2-2-11}{2.2.11}.}
\underline{The following diagrams are anticommutative:}
\begin{equation}\tag{\SGARefTarget{sga7.sga7i.eviii.n2-2-11-1}{2.2.11.1}}
\begin{tikzcd}[row sep=large,column sep=huge]
 \mathrm{Biext}^1(P,Q;G)
   \arrow[rr,"\sim"]
   \arrow[dd,"{\varphi^n_{P,Q}}"']&&
 \mathrm{Biext}^1(Q,P;G)
   \arrow[dd,"{\varphi^n_{Q,P}}"]\\
 &- &\\
 \operatorname{Hom}({}_nP\otimes{}_nQ,G)
   \arrow[rr,"\sim"]&&
 \operatorname{Hom}({}_nQ\otimes{}_nP,G),
\end{tikzcd}
\end{equation}
\begin{equation}\tag{\SGARefTarget{sga7.sga7i.eviii.n2-2-11-2}{2.2.11.2}}
\begin{tikzcd}[row sep=large,column sep=huge]
 \mathrm{Biext}^1(P,Q;G)
   \arrow[rr,"\sim"]
   \arrow[dd,"{\varphi^{(\ell)}_{P,Q}}"']&&
 \mathrm{Biext}^1(Q,P;G)
   \arrow[dd,"{\varphi^{(\ell)}_{Q,P}}"]\\
 &- &\\
 \operatorname{Hom}\bigl(
   T_\ell(P)\otimes T_\ell(Q),T_\ell(G)
 \bigr)
   \arrow[rr,"\sim"]&&
 \operatorname{Hom}\bigl(
   T_\ell(Q)\otimes T_\ell(P),T_\ell(G)
 \bigr),
\end{tikzcd}
\end{equation}
\underline{where the horizontal arrows are the symmetry arrows.}

We record the important corollary.

\medskip
\noindent\underline{Corollary \SGARefTarget{sga7.sga7i.eviii.n2-2-12}{2.2.12}.}
\underline{Suppose that $P=Q$, and suppose that $\xi$ is the class of
a symmetric (respectively, antisymmetric) biextension of $(P,P)$ by
$G$; that is, suppose ${}^s\xi=\xi$ (respectively,
${}^s\xi=-\xi$), where ${}^s\xi$ denotes the image of $\xi$ under
symmetry.  Then the bilinear maps $\varphi^n_\xi$, on
${}_nP\times{}_nP$, and $\varphi^{(\ell)}_\xi$, on
$T_\ell(P)\times T_\ell(P)$, are antisymmetric (respectively,
symmetric).}

\medskip
\noindent\underline{Remark \SGARefTarget{sga7.sga7i.eviii.n2-2-13}{2.2.13}.}
Suppose that $\xi$ is symmetric.  If $n$ (respectively, $\ell$) is
odd, then \SGARefLink{sga7.sga7i.eviii.n2-2-12}{2.2.12} shows that $\varphi^n_\xi$ (respectively,
$\varphi^{(\ell)}_\xi$) is in fact an alternating form.  This
conclusion is not valid in general without the restriction that $n$
(respectively, $\ell$) be odd, as is already seen when $P$ is the
ordinary group $\mathbb Z/2\mathbb Z$ and
$G=\mathbb Z/2\mathbb Z$.

Suppose, however, that $2P=P$.  This is the case in particular when
$P$ is an abelian scheme, working as usual in the fppf or fpqc
topology.  The preceding forms are then still alternating.  It is
enough to prove that $\varphi^{2^\nu}_\xi(x,x)=0$.  Since $2P=P$, a
given element $x\in{}_{2^\nu}P(S)$ may be written as $x=2y$, with
$y\in{}_{2^{\nu+1}}P(S)$.  Hence
\[
 \varphi^{2^\nu}_\xi(x,x)
 =2\varphi^{2^{\nu+1}}_\xi(y,y)=0,
\]
because $\varphi^{2^{\nu+1}}_\xi$ is antisymmetric.

% END INCLUDED SGA7I_English_current_buildable_source_20260823/components/98_expose_VIII_section_2_2.tex

% BEGIN INCLUDED SGA7I_English_current_buildable_source_20260823/components/99_expose_VIII_section_2_3_opening_through_proposition_2_3_11.tex
% SOURCE AUTHORITY:
% expose_VIII_body.tex lines 836--923; scan indices 252--254; source folios 241--243.
% Direct scan comparison: physical PDF pages 253--255 at 1100 dpi.
% Transparent source disposition:
% - the right-hand vertex of (2.3.10) is restored as
%   Hom({}_nP,{}_n Ext^1(nQ,G)); the printed/transcribed "nQnG" is a
%   comma-key slip and is incompatible with (2.3.5), (2.3.9), and the prose.
% The authority's numbering jump from (2.3.2) to (2.3.4) is preserved.
% Hard stop before the proof of 2.3.11, source line 924.

\SGARefTarget{sga7.sga7i.eviii.n2-3}{2.3}. \underline{Relation with ``Cartier duality''.}
The description given in \SGARefLink{sga7.sga7i.eviii.n2-2}{2.2} of the pairings defined by a
biextension, through the isomorphism theorem VII \SGARefLink{sga7.sga7i.evii.n3-6-5}{3.6.5} and the
homomorphism (\SGARefLink{sga7.sga7i.eviii.n2-2-1}{2.2.1}), is plainly not the definition formerly used for
abelian schemes, at a time when the very notion of biextension did not
yet exist.  That definition proceeds through Cartier duality, and we
now give its relation to the present one.

Consider again the homomorphism
\begin{equation}\tag{\SGARefTarget{sga7.sga7i.eviii.n2-3-1}{2.3.1}}
 \mathrm{Biext}^1(P,Q;G)
 \longrightarrow
 \operatorname{Hom}\bigl(P,\uExt^1(Q,G)\bigr)
\end{equation}
arising from the exact sequence (\SGARefLink{sga7.sga7i.eviii.n1-4}{1.4}).  It is in fact an isomorphism
when
\begin{equation}\tag{\SGARefTarget{sga7.sga7i.eviii.n2-3-2}{2.3.2}}
 \uHom(Q,G)=0.
\end{equation}
This condition holds in particular when $Q$ is an abelian scheme over
a scheme $S$ and $G$ is the multiplicative group $\underline G_m$
(VII \SGARefLink{sga7.sga7i.evii.n1-3-8}{1.3.8}).

There is also an evident ``restriction to $nQ$'' homomorphism
\[
 \uExt^1(Q,G)\longrightarrow\uExt^1(nQ,G),
\]
and hence a composite homomorphism
\begin{equation}\tag{\SGARefTarget{sga7.sga7i.eviii.n2-3-4}{2.3.4}}
\begin{split}
 \operatorname{Hom}\bigl(P,\uExt^1(Q,G)\bigr)
 &\longrightarrow
 \operatorname{Hom}\bigl(P,\uExt^1(nQ,G)\bigr)\\
 &\longrightarrow
 \operatorname{Hom}\bigl({}_nP,{}_n\uExt^1(nQ,G)\bigr),
\end{split}
\end{equation}
where the second arrow is restriction to $\,{}_nP$.  Composing with
(\SGARefLink{sga7.sga7i.eviii.n2-3-1}{2.3.1}), we obtain a homomorphism
\begin{equation}\tag{\SGARefTarget{sga7.sga7i.eviii.n2-3-5}{2.3.5}}
 \mathrm{Biext}^1(P,Q;G)
 \longrightarrow
 \operatorname{Hom}\bigl({}_nP,{}_n\uExt^1(nQ,G)\bigr).
\end{equation}

On the other hand, the exact sequence
\begin{equation}\tag{\SGARefTarget{sga7.sga7i.eviii.n2-3-6}{2.3.6}}
 0\longrightarrow{}_nQ\longrightarrow Q
 \xrightarrow{\ n\ }nQ\longrightarrow0
\end{equation}
gives rise to a long exact sequence of the $\uExt^i(-,G)$, and hence
to the short exact sequence
\begin{equation}\tag{\SGARefTarget{sga7.sga7i.eviii.n2-3-7}{2.3.7}}
 0\longrightarrow\uHom(Q,G)_n
 \longrightarrow\uHom({}_nQ,G)
 \xrightarrow{\ \partial\ }
 {}_n\uExt^1(nQ,G)
 \longrightarrow0,
\end{equation}
where, as usual, $E_n$ denotes $E/nE$.  We retain from this sequence
the canonical surjective homomorphism
\begin{equation}\tag{\SGARefTarget{sga7.sga7i.eviii.n2-3-8}{2.3.8}}
 \uHom({}_nQ,G)
 \xrightarrow{\hspace{3em}\partial\hspace{3em}}
 {}_n\uExt^1(nQ,G),
\end{equation}
which is an isomorphism if $\uHom(Q,G)_n=0$, and a fortiori if (\SGARefLink{sga7.sga7i.eviii.n2-3-2}{2.3.2})
holds.  From (\SGARefLink{sga7.sga7i.eviii.n2-3-8}{2.3.8}) we deduce a homomorphism
\begin{equation}\tag{\SGARefTarget{sga7.sga7i.eviii.n2-3-9}{2.3.9}}
\begin{split}
 \operatorname{Hom}\bigl({}_nP,\uHom({}_nQ,G)\bigr)
 &\simeq
 \operatorname{Hom}({}_nP\otimes{}_nQ,G)\\
 &\longrightarrow
 \operatorname{Hom}\bigl({}_nP,{}_n\uExt^1(nQ,G)\bigr),
\end{split}
\end{equation}
which is an isomorphism when (\SGARefLink{sga7.sga7i.eviii.n2-3-2}{2.3.2}) holds.

All the homomorphisms under consideration are functorial in $P,Q,G$.
The homomorphisms (\SGARefLink{sga7.sga7i.eviii.n2-3-5}{2.3.5}), (\SGARefLink{sga7.sga7i.eviii.n2-3-9}{2.3.9}), and (\SGARefLink{sga7.sga7i.eviii.n2-2-2}{2.2.2}) form the diagram
\begin{equation}\tag{\SGARefTarget{sga7.sga7i.eviii.n2-3-10}{2.3.10}}
\begin{tikzcd}[column sep=large,row sep=large]
 &\mathrm{Biext}^1(P,Q;G)
   \arrow[ddl,"\varphi^n"']
   \arrow[ddr,"(2.3.5)"]&\\
 &-&\\
 \operatorname{Hom}({}_nP\otimes{}_nQ,G)
   \arrow[rr,"(2.3.9)"]&&
 \operatorname{Hom}\bigl({}_nP,{}_n\uExt^1(nQ,G)\bigr),
\end{tikzcd}
\end{equation}
where the horizontal arrow (\SGARefLink{sga7.sga7i.eviii.n2-3-9}{2.3.9}) is an isomorphism if (\SGARefLink{sga7.sga7i.eviii.n2-3-2}{2.3.2})
holds.

\medskip
\noindent\underline{Proposition \SGARefTarget{sga7.sga7i.eviii.n2-3-11}{2.3.11}.}
\underline{The diagram (\SGARefLink{sga7.sga7i.eviii.n2-3-10}{2.3.10}) is anticommutative.}

When $\uHom(Q,G)=0$, this statement therefore recovers a
characterization of the homomorphism $\varphi^n$ of (\SGARefLink{sga7.sga7i.eviii.n2-2-2}{2.2.2}) in terms
of the homomorphisms (\SGARefLink{sga7.sga7i.eviii.n2-3-5}{2.3.5}) and (\SGARefLink{sga7.sga7i.eviii.n2-3-9}{2.3.9}).  This is the
characterization formerly used in the theory of abelian schemes,
where moreover $nQ=Q$.

% END INCLUDED SGA7I_English_current_buildable_source_20260823/components/99_expose_VIII_section_2_3_opening_through_proposition_2_3_11.tex

% BEGIN INCLUDED SGA7I_English_current_buildable_source_20260823/components/100_expose_VIII_proof_2_3_11_parts_a_and_b.tex
% SOURCE AUTHORITY:
% expose_VIII_body.tex lines 924--1021; scan indices 254--257; source folios 243--246.
% Direct scan comparison: physical PDF pages 255--258 at 1100 dpi.
% Transparent source dispositions:
% - the proof heading's printed 2.5.11 is restored as 2.3.11;
% - typewritten letter-o indices and cochain degrees are normalized to zero;
% - the arrow-key slip inside the defining formula (2.3.14) is restored as
%   the minus sign already printed in (2.3.13) and required by the cocycle.
% Hard stop before part c), source line 1022.

\medskip
\noindent\underline{Proof of \SGARefLink{sga7.sga7i.eviii.n2-3-11}{2.3.11}.}\footnote{For a more elegant
proof, see \SGARefLink{sga7.sga7i.eviii.n2-3-19}{2.3.19} below.}
Let $\xi$ be an element of $\mathrm{Biext}^1(P,Q;G)$, let
\[
 \varphi_\xi:{}_nP\otimes{}_nQ\longrightarrow G
\]
be the associated pairing, and let
\[
 u_\xi:{}_nP\longrightarrow{}_n\uExt^1(nQ,G)
\]
be the image of $\xi$ under (\SGARefLink{sga7.sga7i.eviii.n2-3-5}{2.3.5}).  We must prove that $-u_\xi$ is
obtained from $\varphi_\xi$ by the arrow (\SGARefLink{sga7.sga7i.eviii.n2-3-9}{2.3.9}), that is, that it is
the composite
\begin{equation}\tag{*}
 {}_nP
 \xrightarrow{\ \varphi_\xi\ }
 \uHom({}_nQ,G)
 \xrightarrow{\ (\SGARefLink{sga7.sga7i.eviii.n2-3-8}{2.3.8})\ }
 {}_n\uExt^1(nQ,G).
\end{equation}
To prove this, we may take a section $x$ of $\,{}_nP$ and prove that
its image under $-u_\xi$ equals its image under the composite (*).
If desired, this would allow us to reduce to the case
$P=(\mathbb Z/n\mathbb Z)_T$, but we shall not need that reduction.

\medskip
\noindent a) Choose flat resolutions $L_\bullet$ and $M_\bullet$,
\[
 0\longrightarrow L_1\longrightarrow L_0\longrightarrow0,
 \qquad
 0\longrightarrow M_1\longrightarrow M_0\longrightarrow0,
\]
of $P$ and $Q$, respectively.  A flat resolution of $P$, for example,
is obtained by writing $P$ as a quotient of a flat
$\mathbb Z_T$-module $L_0$ and taking $L_1$ to be the kernel of
$L_0\longrightarrow P$; this kernel is flat because it is
torsion-free.  We obtain the complex $L_\bullet\otimes M_\bullet$
\[
 0\longrightarrow L_1\otimes M_1
 \xrightarrow{d_2}
 L_1\otimes M_0+L_0\otimes M_1
 \xrightarrow{d_1}
 L_0\otimes M_0
 \longrightarrow0,
\]
whose sheaf $\underline H_1$ will serve as $\uTor_1(P,Q)$.

Let $L'_0$ (respectively, $M'_0$) be the subsheaf of $L_0$
(respectively, $M_0$) which is the inverse image of $\,{}_nP$
(respectively, $\,{}_nQ$).  Thus
\[
 0\longrightarrow L_1\longrightarrow L'_0\longrightarrow0
 \qquad
 \left(\text{respectively, }
 0\longrightarrow M_1\longrightarrow M'_0\longrightarrow0\right)
\]
is a flat resolution of $\,{}_nP$ (respectively, of $\,{}_nQ$).  The
homomorphism
\[
 {}_nP\otimes{}_nQ
 \longrightarrow
 \uTor_1(P,Q)=\underline H_1(L_\bullet\otimes M_\bullet)
\]
may then be described, by passage to the quotient, from the
homomorphism
\[
 L'_0\otimes M'_0
 \longrightarrow
 \underline H_1(L_\bullet\otimes M_\bullet)
\]
given by
\[
 x_0\otimes y_0
 \longmapsto
 \text{the class of the $1$-cycle }
 [nx_0]\otimes y_0-x_0\otimes[ny_0].
\]
Here the brackets indicate that the element in question is regarded
as a section of the subsheaf $L_1$ of $L_0$ (respectively, $M_1$ of
$M_0$).  This description follows from the calculation made in the
proof of \SGARefLink{sga7.sga7i.eviii.n2-1-10}{2.1.10}.

\medskip
\noindent b) Take an injective resolution
\[
 C^\bullet=\bigl(0\longrightarrow C^0
 \xrightarrow{d^0}C^1\longrightarrow\cdots\bigr)
\]
of $G$, and identify
\[
 \mathrm{Biext}^1(P,Q;G)
 \simeq
 \operatorname{Ext}^1\!\left(
   P\mathbin{\overset{\mathbb L}{\otimes}}Q,G
 \right)
 =H^1\!\left(
   \operatorname{Hom}^{\boldsymbol\cdot}
   (L_\bullet\otimes M_\bullet,C^{\boldsymbol\cdot})
 \right).
\]
The element $\xi$ may therefore be represented by a pair of
homomorphisms $f_1,f_0$, giving the anticommutative diagram
\begin{equation}\tag{\SGARefTarget{sga7.sga7i.eviii.n2-3-12}{2.3.12}}
\begin{tikzcd}[column sep=huge,row sep=huge]
 L_1\otimes M_0+L_0\otimes M_1
   \arrow[r,"d_1"]
   \arrow[d,"f_1"']
   \arrow[rd,phantom,"-" description]&
 L_0\otimes M_0
   \arrow[d,"f_0"']\\
 C^0\arrow[r,"d^0"]&C^1.
\end{tikzcd}
\end{equation}
By definition, the pairing
\[
 \varphi_\xi:{}_nP\otimes{}_nQ\longrightarrow G
\]
is obtained, by passage to the quotient, from the homomorphism
\[
 \varphi:L'_0\otimes M'_0\longrightarrow G
\]
given by
\begin{equation}\tag{\SGARefTarget{sga7.sga7i.eviii.n2-3-13}{2.3.13}}
 \varphi(x_0\otimes y_0)
 =f_1\bigl([nx_0]\otimes y_0-x_0\otimes[ny_0]\bigr)
 \in G=\ker(C^0\longrightarrow C^1).
\end{equation}
Since assertion \SGARefLink{sga7.sga7i.eviii.n2-3-11}{2.3.11} is local on $T$, we may suppose that the given
section $x$ of $\,{}_nP$ lifts to a section $x_0$ of $L'_0$.
Regarding $y_0$ as variable, the preceding formula then defines a
homomorphism
\begin{equation}\tag{\SGARefTarget{sga7.sga7i.eviii.n2-3-14}{2.3.14}}
 \varphi':M'_0\longrightarrow G,
 \qquad
 y_0\longmapsto
 f_1\bigl([nx_0]\otimes y_0-x_0\otimes[ny_0]\bigr),
\end{equation}
which passes to the quotient to give
\[
 {}_nQ\longrightarrow G,
 \qquad
 y\longmapsto\varphi_\xi(x,y),
\]
the homomorphism associated with $\varphi_\xi$ and $x$.

% END INCLUDED SGA7I_English_current_buildable_source_20260823/components/100_expose_VIII_proof_2_3_11_parts_a_and_b.tex

% BEGIN INCLUDED SGA7I_English_current_buildable_source_20260823/components/101_expose_VIII_proof_2_3_11_parts_c_through_e_and_remark_2_3_16.tex
% SOURCE AUTHORITY:
% expose_VIII_body.tex lines 1022--1221; scan indices 257--260;
% source folios 246--249.
% Direct scan comparison: physical PDF pages 258--261 in overlapping
% 1100-dpi bands.  The ambiguous right-hand component in the first diagram
% of part d was additionally examined at 5000-dpi equivalent magnification.
% Transparent source dispositions:
% - typewritten letter-o indices and cochain degrees are normalized to zero;
% - the source prints f_1 on the right vertical arrow of the degree-one
%   diagram in part d; its codomain is C^1, and the immediately following
%   formula reads g_0=n f_0, so the English diagram restores the type-correct
%   f_0;
% - the uppercase X in Remark 2.3.16 is restored to the section x used in
%   the preceding and following clauses.
% Hard stop before 2.3.17, source line 1222.

\medskip
\noindent c) We must take the image of this homomorphism under the
connecting operator $\partial$ of (\SGARefLink{sga7.sga7i.eviii.n2-3-7}{2.3.7}), associated with the exact
sequence (\SGARefLink{sga7.sga7i.eviii.n2-3-6}{2.3.6}).  Now $nQ\simeq Q/{}_nQ$ is plainly also isomorphic
to $M_0/M'_0$, so that we have a homomorphism of exact sequences
\[
\begin{tikzcd}[column sep=large,row sep=large]
 0\arrow[r]&M'_0\arrow[r]\arrow[d]&M_0\arrow[r]\arrow[d]
   &M_0/M'_0\arrow[r]\arrow[d,"\wr"']&0\\
 0\arrow[r]&{}_nQ\arrow[r]&Q\arrow[r]&nQ\arrow[r]&0&{,}
\end{tikzcd}
\]
where the last vertical arrow is an isomorphism.  This shows that the
desired section of $\uExt^1(nQ,G)$ may also be computed as the image
of $\varphi'$ in (\SGARefLink{sga7.sga7i.eviii.n2-3-14}{2.3.14}) under the connecting operator
\[
 \uHom(M'_0,G)\xrightarrow{\ \partial\ }
 \uExt^1(M_0/M'_0,G),
\]
or, equivalently, as the section of the second term defined by the
element of global $\mathrm{Ext}^1$ which is the image of $\varphi'$
under the connecting operator
\[
 \mathrm{Hom}(M'_0,G)\longrightarrow
 \mathrm{Ext}^1(M_0/M'_0,G).
\]
The usual construction of the connecting operator gives the following
recipe.  Extend the composite (\SGARefLink{sga7.sga7i.eviii.n2-3-13}{2.3.13})
\[
 M'_0\xrightarrow{\ \varphi'\ }G\longrightarrow C^0
\]
to a homomorphism
\[
 \psi:M_0\longrightarrow C^0,
\]
which is possible since $C^0$ is injective, and consider the unique
homomorphism
\begin{equation}\tag{\SGARefTarget{sga7.sga7i.eviii.n2-3-15}{2.3.15}}
 \lambda:M_0/M'_0\longrightarrow C^1
\end{equation}
which makes the diagram
\[
\begin{tikzcd}[column sep=large,row sep=large]
 0\arrow[r]&M'_0\arrow[r]\arrow[d,"\varphi'"']
   &M_0\arrow[r]\arrow[d,"\psi"']
   &M_0/M'_0\arrow[r]\arrow[d,"\lambda"']&0\\
 0\arrow[r]&G\arrow[r]&C^0\arrow[r]&C^1&&{.}
\end{tikzcd}
\]
commutative.  It is clear that (\SGARefLink{sga7.sga7i.eviii.n2-3-15}{2.3.15}) takes values in the
$1$-cocycles---that is, it is a $1$-cocycle---and that this cocycle
represents the desired element of
\[
 \mathrm{Ext}^1(nQ,G)=\mathrm{Ext}^1(M_0/M'_0,G).
\]

\medskip
\noindent d) It remains to describe explicitly the section of
$\uExt^1(nQ,G)$ obtained from $\xi$ and $x$ by means of arrow (\SGARefLink{sga7.sga7i.eviii.n2-3-5}{2.3.5})
in diagram (\SGARefLink{sga7.sga7i.eviii.n2-3-10}{2.3.10}).  For this purpose we may disregard the fact that
the section $x$ of $P$ is annihilated by $n$.  In terms of the chosen
lift $x_0$ of $x$ to a section of $L_0$, the description is as follows.
This lift defines an augmentation homomorphism
\[
 u:M_\bullet\longrightarrow L_\bullet\otimes M_\bullet,
 \qquad
 y_0\longmapsto x_0\otimes y_0,
 \quad
 y_1\longmapsto x_0\otimes y_1.
\]
The section of $\uExt^1(Q,G)$ obtained from (\SGARefLink{sga7.sga7i.eviii.n2-3-1}{2.3.1}) is therefore the
one obtained from the element $\xi$ of
$\mathrm{Ext}^1(L_\bullet\otimes M_\bullet,G)$ by composition with
the preceding homomorphism.  Thus it is represented by the
degree-one homomorphism
\[
\begin{tikzcd}[column sep=large,row sep=huge]
 0\arrow[r]&M_1\arrow[r,"d_1"]
   \arrow[d,"{y_1\mapsto f_1(x_0\otimes y_1)}"']
   \arrow[rd,phantom,"-" description]
   &M_0\arrow[r]\arrow[d,"{y_0\mapsto f_0(x_0\otimes y_0)}"]&0\\
 0\arrow[r]&C^0\arrow[r,"d^0"]&C^1&\cdots&{.}
\end{tikzcd}
\]
To restrict this class to $nQ\hookrightarrow Q$, use the isomorphism
$nQ\simeq M_0/M'_0$, which gives a flat resolution
\[
 0\longrightarrow M'_0\longrightarrow M_0\longrightarrow0
\]
of $nQ$.  The inclusion $nQ\hookrightarrow Q$ lifts to the following
homomorphism of resolutions:
\[
\begin{tikzcd}[column sep=large,row sep=huge]
 0\arrow[r]&M'_0\arrow[r]
   \arrow[d,"{y'_0\mapsto[ny'_0]}"']
   \arrow[rd,phantom,"+" description]
   &M_0\arrow[r]\arrow[d,"{y_0\mapsto ny_0}"]&0\\
 0\arrow[r]&M_1\arrow[r]&M_0\arrow[r]&0&{.}
\end{tikzcd}
\]
By composition, the desired section of $\uExt^1(nQ,G)$ is therefore
associated with the element of $\mathrm{Ext}^1(nQ,G)$ represented by
the degree-one homomorphism of complexes
\[
\begin{tikzcd}[column sep=large,row sep=large]
 0\arrow[r,dash]&M'_0\arrow[r,dash,"d_1"]\arrow[d,"g_1"']
   &M_0\arrow[r,dash]\arrow[d,"g_0"]&0&{}\\
 0\arrow[r,dash]&C^0\arrow[r,dash]&C^1\arrow[r,dash]&\dots&{,}
\end{tikzcd}
\]
where
\[
 g_1(y'_0)=f_1\bigl(x_0\otimes[ny'_0]\bigr),
 \qquad
 g_0(y_0)=n f_0(x_0\otimes y_0).
\]

\medskip
\noindent e) Using now the fact that $x$ is a section of $\,{}_nP$,
we show that the preceding element of $\mathrm{Ext}^1(nQ,G)$ is the
negative of the element defined by (\SGARefLink{sga7.sga7i.eviii.n2-3-15}{2.3.15}).  For any homomorphism
\[
 h:M_0\longrightarrow C^0,
\]
the element of $\mathrm{Ext}^1$ considered in d) is unchanged if
$g_1,g_0$ are replaced, respectively, by
\[
 g'_1=g_1+h d_1,
 \qquad
 g'_0=g_0-d^0h.
\]
We have at our disposal the homomorphism
$\psi:M_0\longrightarrow C^0$ constructed in c), which by definition
satisfies
\[
 \psi(y'_0)=
 f_1\bigl([nx_0]\otimes y'_0-x_0\otimes[ny'_0]\bigr)
\]
for every section $y'_0$ of $M'_0$.  Equivalently,
\[
 f_1\bigl(x_0\otimes[ny'_0]\bigr)
 \bigl(=g_1(y'_0)\bigr)
 =f_1\bigl([nx_0]\otimes y'_0\bigr)-\psi(y'_0).
\]
Define $h:M_0\longrightarrow C^0$ by
\[
 h(y_0)=-f_1\bigl([nx_0]\otimes y_0\bigr)+\psi(y_0).
\]
By construction, then,
\[
 g'_1=0,
 \qquad
 g'_0(y_0)=n f_0(x_0\otimes y_0)
 +d^0f_1\bigl([nx_0]\otimes y_0\bigr)-d^0\psi(y_0).
\]
Since $d^0f_1=-f_0d_1$, we have
\[
 d^0f_1\bigl([nx_0]\otimes y_0\bigr)
 =-f_0(nx_0\otimes y_0),
\]
and hence
\[
 g'_1=0,
 \qquad
 g'_0=-d^0\psi.
\]
Thus the element in question is the one defined by the homomorphism
\[
 M_0/M'_0\simeq nQ\longrightarrow C^1
\]
induced by $-d^0\psi$ upon passage to the quotient.  We therefore
obtain the negative of the element defined by (\SGARefLink{sga7.sga7i.eviii.n2-3-15}{2.3.15}).

\medskip
\noindent\underline{Remarks \SGARefTarget{sga7.sga7i.eviii.n2-3-16}{2.3.16}.}  The preceding proof actually
gives a slightly more precise compatibility than the one stated in
\SGARefLink{sga7.sga7i.eviii.n2-3-11}{2.3.11}, obtained by choosing a section $x$ of $\,{}_nP$.  This section
then defines a homomorphism
\[
 {}_nQ\longrightarrow G,
 \qquad
 y\longmapsto\varphi_\xi(x,y),
\]
and hence, by the connecting homomorphism associated with (\SGARefLink{sga7.sga7i.eviii.n2-3-6}{2.3.6}), an
element of $\mathrm{Ext}^1(nQ,G)$.  On the other hand, if we forget
that $x$ is annihilated by $n$, we consider the homomorphism
$\mathbb Z_T\longrightarrow P$ defined by $x$.  Its inverse image is a
biextension of $(\mathbb Z_T,Q)$ by $G$, whose class may be interpreted
as an element of $\mathrm{Ext}^1(Q,G)$ and therefore induces an element
of $\mathrm{Ext}^1(nQ,G)$.  With these conventions, the two elements
obtained in $\mathrm{Ext}^1(nQ,G)$ are \underline{opposite}.

% END INCLUDED SGA7I_English_current_buildable_source_20260823/components/101_expose_VIII_proof_2_3_11_parts_c_through_e_and_remark_2_3_16.tex

% BEGIN INCLUDED SGA7I_English_current_buildable_source_20260823/components/102_expose_VIII_section_2_3_17.tex
% SOURCE AUTHORITY:
% expose_VIII_body.tex lines 1222--1241; scan indices 260--262;
% source folios 249--251.
% Direct scan comparison: physical PDF pages 261--263 in overlapping
% 1100-dpi bands.
% Transparent source dispositions:
% - typewritten letter-o indices are normalized to zero;
% - the final reference to phi^n is rendered as phi_xi^n, the pairing that
%   is explicitly being identified throughout 2.3.17.
% Hard stop before 2.3.19, source line 1242.

\medskip
\noindent \SGARefTarget{sga7.sga7i.eviii.n2-3-17}{2.3.17}.  Let us use the explicit cocycle description of the
pairing
\[
 \varphi_\xi^n:{}_nP\otimes{}_nQ\longrightarrow G
\]
given in (\SGARefLink{sga7.sga7i.eviii.n2-3-12}{2.3.12}) to obtain another expression for it.  First suppose
that the biextension $E$ of class $\xi$ admits a section $s$ over
$P\times Q$.  As in VII \SGARefLink{sga7.sga7i.evii.n3-8-13}{3.8.13}, take for the flat resolution of $P$
the one with $L_0=\mathbb Z[P]$ and $L_1$ the kernel of the canonical
homomorphism $\mathbb Z[P]\longrightarrow P$, and define
$M_\bullet$ similarly.  There is then a unique biextension
trivialization of the inverse image of $E$ on $(L_0,M_0)$ whose
restriction to $P\times Q$ is given by $s$; denote this trivialization
again by $s$ (cf. VII \SGARefLink{sga7.sga7i.evii.n3-8-2}{3.8.2}).  By VII \SGARefLink{sga7.sga7i.evii.n3-8-12}{3.8.12}, the class of $\xi$ in
\[
 \mathrm{Ext}^1\!\left(
 P\mathbin{\overset{\mathbb L}{\otimes}}Q,G\right)
\]
is represented by the cocycle with components $(f,g)$ characterized by
\[
 s(x_1,y_0)=f(x_1\otimes y_0)e_{E/Q}(y_0),
 \qquad
 s(x_0,y_1)=g(x_0\otimes y_1)e_{E/P}(x_0).
\]
If $x$ and $y$ are sections of $\,{}_nP$ and $\,{}_nQ$, lift them to
the corresponding sections $x_0=[x]$ and $y_0=[y]$ of $L_0$ and
$M_0$.  Applying formula (\SGARefLink{sga7.sga7i.eviii.n2-3-13}{2.3.13}), and putting $z=s(x,y)$, gives
\begin{equation}\tag{\SGARefTarget{sga7.sga7i.eviii.n2-3-18}{2.3.18}}
 \varphi_\xi^n(x,y)
 ={}^n z/e_{E/Q}(y)-z^n/e_{E/P}(x).
\end{equation}
Here the exponent $n$ written on the left, respectively on the right,
of $z$ denotes the $n$th power for the left, respectively right,
multiplicative structure on $E$.  It therefore gives a section of $E$
over $(0,y)$, respectively $(x,0)$; the meaning of the ``quotient'' of
such an element by $e_{E/Q}(y)$, respectively $e_{E/P}(x)$, is clear.

The preceding formula (\SGARefLink{sga7.sga7i.eviii.n2-3-18}{2.3.18}) still makes sense whenever $z$ is a
global section of $E$ lifting the sections $x,y$ of
$\,{}_nP,\,{}_nQ$.  I claim that (\SGARefLink{sga7.sga7i.eviii.n2-3-18}{2.3.18}) is valid for every such
section, without assuming that a section of $E$ over $P\times Q$
exists.  The section $(x,y)$ of $P\times Q$ always lifts locally to a
section of $E$, since $E\longrightarrow P\times Q$ is an epimorphism;
thus the formula computes $\varphi_\xi^n(x,y)$ in every case.
Moreover, if $z$ is replaced by another section $z'$ of $E$ lifting
$(x,y)$, then $z'=z+g$ for a section $g$ of $G$.  The two terms on the
right-hand side of (\SGARefLink{sga7.sga7i.eviii.n2-3-18}{2.3.18}) corresponding to $z'$ are the terms
corresponding to $z$, each increased by $ng$.  Their difference is
therefore independent of $z$, showing that (\SGARefLink{sga7.sga7i.eviii.n2-3-18}{2.3.18}) defines a
well-determined homomorphism
\[
 {}_nP\otimes{}_nQ\longrightarrow G.
\]

It remains to show that this homomorphism is indeed $\varphi_\xi^n$.
For this we reduce at once to the case
\[
 P=Q=(\mathbb Z/n\mathbb Z)_T.
\]
Then $P\times Q$ is of the form $I_T$, where
$I=\mathbb Z/n\mathbb Z\times\mathbb Z/n\mathbb Z$ is a finite set.
In this case the torsor $E$ under $G_{P\times Q}$---which identifies
with a family of torsors $E_i$ under $G$, for $i\in I$---plainly
becomes trivial locally on $T$.  We are thus reduced to the already
treated case in which $E$ has a global section over $P\times Q$.

% END INCLUDED SGA7I_English_current_buildable_source_20260823/components/102_expose_VIII_section_2_3_17.tex

% BEGIN INCLUDED SGA7I_English_current_buildable_source_20260823/components/103_expose_VIII_alternative_proof_2_3_19.tex
% SOURCE AUTHORITY:
% expose_VIII_body.tex lines 1242--1342; scan indices 262--263;
% source folios 251--252.
% Direct scan comparison: physical PDF pages 263--264 in overlapping
% 1100-dpi bands.  The three-square homotopy display and the full 2.3.10
% derived-Hom diagram were additionally inspected in five targeted crops at
% 5000-dpi equivalent magnification.
% Transparent source dispositions:
% - the source's typewritten [o] is normalized to the cohomological [0];
% - the circled source mark "- ?" in the large diagram is retained exactly
%   as an unresolved sign annotation, rather than silently adjudicated.
% Hard stop before section 2.4, source line 1346 / scan index 264.

\medskip
\noindent \SGARefTarget{sga7.sga7i.eviii.n2-3-19}{2.3.19}.  \underline{Another proof of \SGARefLink{sga7.sga7i.eviii.n2-3-11}{2.3.11}} (by
Deligne).

\medskip
\noindent a) Reduction to $P=\mathbb Z/n\mathbb Z$ and to
$T=(\mathrm{Sets})$.

\medskip
\noindent b) Let $\partial$ be the following arrow from $nQ$ to
$\,{}_nQ[1]$:
\[
 nQ\xleftarrow[\text{``$n$''}]{\ \sim\ }
 [\,{}_nQ\longrightarrow Q\,]
 \longrightarrow{}_nQ[1].
\]
The following diagram is anticommutative:
\[
\begin{tikzcd}[row sep=huge,column sep=huge]
 nQ[0]\arrow[r,"1"]\arrow[d,"\partial"']
   &{[Q\xrightarrow{n}Q]}\\
 {}_nQ[1]\arrow[ru,"1"']&{.}
\end{tikzcd}
\]
Indeed, the arrows
\[
\begin{tikzcd}[column sep=small,row sep=large]
 {[{}_nQ}\arrow[r]\arrow[dd,"0"']&{Q]}\arrow[dd,"n"]
   &&{[{}_nQ}\arrow[r]\arrow[dd,"1"']&{Q]}\arrow[dd,"0"']
   &\text{are opposite}
   &{[{}_nQ}\arrow[r]&{Q]}\arrow[ddl,"1"]\\
 &&\text{and}&&&\text{by means of}\quad H\,:&&\\
 {[Q}\arrow[r,"n"]&{Q]}&&{[Q}\arrow[r,"n"]&{Q]}
   &&{[Q}\arrow[r,"n"]&{Q]}
\end{tikzcd}
\]

\medskip
\noindent c) Diagram \SGARefLink{sga7.sga7i.eviii.n2-3-10}{2.3.10} is
\[
\begin{tikzcd}[column sep=small,row sep=small]
 &&&H^1\mathrm{RHom}
   (\mathbb Z\mathbin{\overset{\mathbb L}{\otimes}}Q,G)
   \arrow[ddd]\\
 &&&+\\
 H^1\mathrm{RHom}
   (\mathbb Z/n\mathbb Z\mathbin{\overset{\mathbb L}{\otimes}}Q,G)
   \arrow[rr]\arrow[dd]\arrow[uurrr]
 &&H^1\mathrm{RHom}
   (\mathbb Z/n\mathbb Z,\mathrm{RHom}(Q,G))
   \arrow[dd]&\\
 &&&\mathrm{Ext}^1(Q,G)\arrow[ddl]\\
 \mathrm{Hom}(\mathrm{Tor}_1(\mathbb Z/n\mathbb Z,Q),G)\arrow[d]
 &|[draw,circle,inner sep=1pt]|{-\,?}
 &\mathrm{Hom}(\mathbb Z/n\mathbb Z,\mathrm{Ext}^1(Q,G))
   \arrow[d]&\\
 \mathrm{Hom}({}_nQ,G)\arrow[rr]
 &&\mathrm{Ext}^1(nQ,G)&
\end{tikzcd}
\]
and:

\smallskip
\noindent c1) We have an isomorphism
\[
 \mathbb Z/n\mathbb Z
 \mathbin{\overset{\mathbb L}{\otimes}}Q
 \simeq[Q\xrightarrow{n}Q],
\]
and, under this isomorphism, the first composite vertical arrow is
induced by
\[
 {}_nQ[1]\longrightarrow[Q\xrightarrow{n}Q].
\]

\smallskip
\noindent c2) The second horizontal arrow is the transpose of the
arrow $\partial$ displayed above in b).

\smallskip
\noindent c3) If the diagram is traversed clockwise, it is clear that
one obtains the arrow induced by
\[
 nQ[0]\longrightarrow[Q\xrightarrow{n}Q].
\]
This proves \SGARefLink{sga7.sga7i.eviii.n2-3-11}{2.3.11}.

% END INCLUDED SGA7I_English_current_buildable_source_20260823/components/103_expose_VIII_alternative_proof_2_3_19.tex

% BEGIN INCLUDED SGA7I_English_current_buildable_source_20260823/components/104_expose_VIII_section_2_4.tex
% SOURCE AUTHORITY:
% expose_VIII_body.tex lines 1346--1380; scan index 264; source folio 253.
% Direct scan comparison: physical PDF page 265 in four overlapping
% 1100-dpi bands.
% Hard stop before section 3, source line 1386 / scan index 265.

\medskip
\noindent \SGARefTarget{sga7.sga7i.eviii.n2-4}{2.4}.  Let
\[
 f:\underline T'\longrightarrow\underline T
\]
be a morphism of topoi.  There is then a canonical homomorphism
(VII \SGARefLink{sga7.sga7i.evii.n2-8}{2.8})
\[
 \Biext^1(P,Q;G)\longrightarrow\Biext^1(P',Q';G'),
\]
where primes denote inverse-image groups.  Since, for every
commutative group $E$ of $\underline T$, there is plainly a canonical
isomorphism
\[
 ({}_nE)'\simeq{}_n(E'),
\]
the homomorphisms (\SGARefLink{sga7.sga7i.eviii.n2-2-2}{2.2.2}) over $\underline T$ and $\underline T'$ give
a canonical diagram
\begin{equation}\tag{\SGARefTarget{sga7.sga7i.eviii.n2-4-1}{2.4.1}}
\begin{tikzcd}[column sep=huge,row sep=huge]
 \Biext^1(P,Q;G)\arrow[r]\arrow[d,"{\varphi^n_{P,Q}}"']
   &\Biext^1(P',Q';G')\arrow[d,"{\varphi^n_{P',Q'}}"]\\
 \Hom({}_nP\otimes{}_nQ,G)\arrow[r]
   &\Hom({}_nP'\otimes{}_nQ',G').
\end{tikzcd}
\end{equation}
Similarly, the homomorphisms (\SGARefLink{sga7.sga7i.eviii.n2-2-3}{2.2.3}) give an analogous diagram in
terms of the $T_\ell$.  I claim that these diagrams are
\underline{commutative}.  The verification is essentially immediate
from the definitions and is left to the reader.  More generally,
notice that \underline{all} the canonical homomorphisms considered in
this section and the preceding one are \underline{compatible with
change of topos}, in the evident sense.  We shall henceforth use
compatibilities of this kind without further mention.


% END INCLUDED SGA7I_English_current_buildable_source_20260823/components/104_expose_VIII_section_2_4.tex

% BEGIN INCLUDED SGA7I_English_current_buildable_source_20260823/components/105_expose_VIII_sections_3_1_and_3_2.tex
% SOURCE AUTHORITY:
% expose_VIII_body.tex lines 1386--1427; scan indices 265--266;
% source folios 254--255.
% Direct scan comparison: physical PDF pages 266--267 in four overlapping
% 1100-dpi bands per page.
% Hard stop before section 3.3, source line 1432 / scan index 267.

\subsection*{\SGARefTarget{sga7.sga7i.eviii.n3}{3}. Biextensions of group schemes $(P,Q)$ by
$\underline G_m$: generalities}

Throughout this section we fix a base scheme $S$ and, as announced in
Section 0, work in the fppf topos of $S$.

\medskip
\noindent \SGARefTarget{sga7.sga7i.eviii.n3-1}{3.1}.  Let us first recall the most important example of a
group $\Biext^1(P,Q;\underline G_m)$: the case where $P$ and $Q$ are
two abelian schemes over $S$, and where $\Biext^1$ is interpreted as
the group of divisorial correspondences on $(P,Q)$ (VII \SGARefLink{sga7.sga7i.evii.n2-9-4}{2.9.4}).  The
$\ell$-adic pairings of the preceding section are then the familiar
pairings [1] in the theory of abelian schemes.

\medskip
\noindent \SGARefTarget{sga7.sga7i.eviii.n3-2}{3.2}.  More generally, suppose only that $Q$ is an abelian
scheme.  It is well known (VII \SGARefLink{sga7.sga7i.evii.n1-3-8}{1.3.8}) that
\begin{equation}\tag{\SGARefTarget{sga7.sga7i.eviii.n3-2-1}{3.2.1}}
 \uHom(Q,\underline G_m)=0,
\end{equation}
and hence (\SGARefLink{sga7.sga7i.eviii.n1-4}{1.4}) there is a canonical isomorphism
\begin{equation}\tag{\SGARefTarget{sga7.sga7i.eviii.n3-2-2}{3.2.2}}
 \Biext^1(P,Q;\underline G_m)\simeq\Hom(P,Q^*),
\end{equation}
where
\begin{equation}\tag{\SGARefTarget{sga7.sga7i.eviii.n3-2-3}{3.2.3}}
 Q^*=\uExt^1(Q,\underline G_m)
\end{equation}
is the dual abelian scheme of $Q$.  (For the representability of
$Q^*$, due to M. Raynaud in the general case, see [1].)

When $P$ is also an abelian scheme, symmetrically to (\SGARefLink{sga7.sga7i.eviii.n3-2-2}{3.2.2}), there is
a canonical isomorphism
\begin{equation}\tag{\SGARefTarget{sga7.sga7i.eviii.n3-2-4}{3.2.4}}
 \Biext^1(P,Q;\underline G_m)\simeq\Hom(Q,P^*).
\end{equation}
Thus the datum of a biextension $\xi$ of $(P,Q)$ by $\underline G_m$
may be interpreted equivalently as the datum of a homomorphism
\[
 u:P\longrightarrow Q^*
\]
or of a homomorphism
\[
 v:Q\longrightarrow P^*.
\]

Recall that a homomorphism of abelian schemes $w:A\longrightarrow B$
is an isomorphism if and only if, for every integer $n>0$, the induced
homomorphism ${}_nw:{}_nA\longrightarrow{}_nB$ is an isomorphism.  In
the case of $u$ and $v$ above, ${}_nu:{}_nP\longrightarrow{}_nQ^*$ is
identified with the homomorphism
${}_nP\longrightarrow D({}_nQ)$ defined by the pairing
\[
 \varphi_\xi^n:{}_nP\times{}_nQ\longrightarrow\underline G_m
\]
(see \SGARefLink{sga7.sga7i.eviii.n2-3-11}{2.3.11}), whereas ${}_nv:{}_nQ\longrightarrow{}_nP^*$ is
identified with the homomorphism
${}_nQ\longrightarrow D({}_nP)$ defined by the pairing
\[
 \psi_\xi^n:{}_nP\times{}_nQ\longrightarrow\underline G_m.
\]
Now $\varphi_\xi^n=-\psi_\xi^n$ (\SGARefLink{sga7.sga7i.eviii.n2-1-10}{2.1.10}), so, up to sign,
${}_nu$ and ${}_nv$ are transposes of one another.  Since ${}_nP$ and
${}_nQ$ are finite and locally free over $S$, and therefore satisfy
Cartier duality, it follows that ${}_nu$ is an isomorphism if and only
if ${}_nv$ is an isomorphism.  Consequently $u$ is an isomorphism if
and only if $v$ is an isomorphism.

An equivalent, though slightly less symmetric, formulation is the
following.  Consider the canonical biextension (the Weil
biextension) of $(P,P^*)$ by $\underline G_m$, corresponding to the
identity homomorphism $v:P^*\longrightarrow P^*$.  Then the
corresponding canonical homomorphism
\[
 u:P\longrightarrow P^{**}
\]
is an \underline{isomorphism}.  This is the \underline{biduality
theorem} of Barsotti and Cartier.


% END INCLUDED SGA7I_English_current_buildable_source_20260823/components/105_expose_VIII_sections_3_1_and_3_2.tex

% BEGIN INCLUDED SGA7I_English_current_buildable_source_20260823/components/106_expose_VIII_section_3_3.tex
% SOURCE AUTHORITY:
% expose_VIII_body.tex lines 1432--1481; scan indices 267--270;
% source folios 256--259.
% Direct scan comparison: physical PDF pages 268--271 in four overlapping
% 1100-dpi bands per page.
% Transparent source dispositions:
% - dualizing 0 -> T -> M -> L -> 0 gives the type-correct terminal D(T),
%   not the repeated D(L) in the working transcription;
% - the visibly blank citation after "SGA 3 X 4.9 et" is not guessed; the
%   dangling conjunction is omitted in English.
% Hard stop before Proposition 3.4, source line 1483 / scan index 270.

\medskip
\noindent \SGARefTarget{sga7.sga7i.eviii.n3-3}{3.3}.  We now turn to the case where $Q$ is finite locally
free over $S$, or is a group scheme of multiplicative type over $S$.
The key lemma in this case is the following well-known result.

\medskip
\noindent\underline{Proposition \SGARefTarget{sga7.sga7i.eviii.n3-3-1}{3.3.1}.}  \underline{Let $Q$ be a
commutative group scheme over $S$, and suppose that $Q$ is finite
locally free over $S$, or that $Q$ is of finite type and of
multiplicative type over $S$ (SGA 3 IX).  Then}
\[
 \uExt^1(Q,\underline G_m)=0.
\]

\noindent\underline{Proof.}  We distinguish the two cases.

\medskip
\noindent a) Suppose that $Q$ is finite locally free over $S$.  We
must prove that every extension $E$ of $Q$ by $\underline G_m$ splits
locally for the fppf topology.  We may plainly suppose $S$ affine;
there is then an integer $n>0$ such that ${}_nQ=0$ (SGA 3 IX 6.1).
Moreover, $n\cdot\id$ is an epimorphism of $\underline G_m$ for the
fppf topology, and hence the Kummer exact sequence
\[
 0\longrightarrow\mu_n\longrightarrow\underline G_m
 \xrightarrow{\ n\ }\underline G_m\longrightarrow0
\]
gives an exact sequence of the $\Ext^i(Q,-)$ and, in particular, an
\underline{epimorphism}
\[
 \Ext^1(Q,\mu_n)\longrightarrow
 {}_n\Ext^1(Q,\underline G_m)=\Ext^1(Q,\underline G_m).
\]
The last equality follows because multiplication by $n$, being zero
on $Q$, is also zero on $\Ext^1(Q,\underline G_m)$.  Thus the
extension $E$ comes from an extension $F$ of $Q$ by $\mu_n$.

The group $F$ is represented by a finite locally free scheme over $S$
(fpqc descent theory, SGA 1 VII).  By the known properties of Cartier
duality for finite locally free groups (SGA 3 VII 3),
\[
 D(F)=\uHom(F,\underline G_m)
\]
is an extension of $D(\mu_n)$ by $D(Q)$.  Here one uses the following
fact: if $A\longrightarrow B$ is a monomorphism of commutative finite
locally free groups over $S$, then
$D(B)\longrightarrow D(A)$ is faithfully flat.  Since $D(B)$ and
$D(A)$ are themselves finite locally free over $S$, the fibrewise
criterion for faithful flatness reduces us to the case where $S$ is
the spectrum of a field $k$.  In that case the assertion says that
the homomorphism of affine coordinate rings
\[
 k(D(A))\longrightarrow k(D(B))
\]
is injective (SGA 3 $\mathrm{VI}_B$ 11.14).  This homomorphism is the
transpose of $k(B)\longrightarrow k(A)$, which is surjective because
$A\longrightarrow B$ is a closed immersion (SGA 3 $\mathrm{VI}_B$
1.4.2).  Hence $k(D(A))\longrightarrow k(D(B))$ is indeed injective.

The inclusion $\mu_n\longrightarrow\underline G_m$ identifies with a
section of
\[
 D(\mu_n)=\uHom(\mu_n,\underline G_m),
\]
and it remains only to prove that, locally for the fppf topology on
$S$, this inclusion extends to a homomorphism
$F\longrightarrow\underline G_m$; equivalently, that the section of
$D(\mu_n)$ just considered lifts fppf-locally to a section of $D(F)$.
This is clear, since $D(F)\longrightarrow D(\mu_n)$ is faithfully flat
and of finite presentation.

\medskip
\noindent b) By SGA 3 X 4.5, after replacing $S$ by an \`etale cover,
we may suppose that $Q$ is diagonalizable, say $Q=D(M)$, where $M$ is
a finitely generated $\mathbb Z$-module.  Let $T$ be the torsion
submodule of $M$, and put $L=M/T$, so that $L\simeq\mathbb Z^r$.  The
exact sequence
\[
 0\longrightarrow T\longrightarrow M\longrightarrow L
 \longrightarrow0
\]
then gives an exact sequence (SGA 3 IX 3.1)
\[
 0\longrightarrow D(L)\longrightarrow D(M)\longrightarrow D(T)
 \longrightarrow0.
\]
It is therefore enough to prove separately that the $\uExt^1$ of
$D(L)$ and of $D(T)$ by $\underline G_m$ vanish.  The case of $D(T)$
follows from a), and the case of $D(L)$ reduces to $L=\mathbb Z$.
Thus it remains to check that
\[
 \uExt^1(\underline G_m,\underline G_m)=0,
\]
that is, that every extension $E$ of $\underline G_m$ by
$\underline G_m$ is fppf-locally trivial.

Such an extension is represented by an affine flat group scheme over
$S$ (fpqc descent theory, SGA 1 VII) whose fibres are tori; it is
therefore a torus (SGA 3 X 4.9).  By SGA 3 X 4.5 we may suppose that
$E$ is a split torus
\[
 E=D(R)\simeq\underline G_m^{,2}.
\]
After a Zariski localization, we may also suppose that the
monomorphism $\underline G_m\hookrightarrow E=D(R)$ comes from a
constant morphism $R\longrightarrow\mathbb Z_S$ (SGA 3 VIII 1.4),
which is necessarily an epimorphism (SGA 3 VIII 3.2 b)).  Choosing
$r\in R$ above $1\in\mathbb Z$ gives a homomorphism
$E\longrightarrow\underline G_m$ splitting the extension, and proves
the assertion.

\medskip
\noindent\underline{Corollary \SGARefTarget{sga7.sga7i.eviii.n3-3-2}{3.3.2}.}  \underline{Let $Q$ be as in
\SGARefLink{sga7.sga7i.eviii.n3-3-1}{3.3.1}.  Then, for every $P$, there is a canonical isomorphism}
\begin{equation}\tag{\SGARefTarget{sga7.sga7i.eviii.n3-3-2-1}{3.3.2.1}}
 \Biext^1(P,Q;\underline G_m)
 \xleftarrow{\ \sim\ }\Ext^1(P,D(Q)),
\end{equation}
\underline{where $D(Q)$ denotes the Cartier dual
$\uHom(Q,\underline G_m)$ of $Q$.  More precisely, the canonical
functor (\SGARefLink{sga7.sga7i.eviii.n1-1-6}{1.1.6}) is an equivalence from the category of extensions of
$P$ by $D(Q)$ to the category of biextensions of $(P,Q)$ by
$\underline G_m$.}

This is the special case of \SGARefLink{sga7.sga7i.eviii.n1-5}{1.5} obtained by applying \SGARefLink{sga7.sga7i.eviii.n3-3-1}{3.3.1}.

\medskip
\noindent\underline{Remarks \SGARefTarget{sga7.sga7i.eviii.n3-3-3}{3.3.3}.}  Under the hypotheses of \SGARefLink{sga7.sga7i.eviii.n3-3-1}{3.3.1},
$D(Q)$ is represented by a separated group scheme, locally of finite
presentation and locally quasi-finite over $S$.  More precisely, if
$Q$ is finite locally free over $S$, then so is $D(Q)$, as is well
known (SGA 3 VII 3).  If $Q$ is of multiplicative type, then $D(Q)$
is a \underline{finitely generated twisted constant group} (SGA 3 X
5.1), as was seen in SGA 3 IX 4.5 and X 5.7.  If $P$ is a group
scheme, every extension of $P$ by $D(Q)$ is therefore representable,
by a familiar lemma from fpqc descent theory (SGA 3 X 5.4).

\medskip
\noindent \SGARefTarget{sga7.sga7i.eviii.n3-3-4}{3.3.4}.  It is false that
$\uExt^1(Q,\underline G_m)=0$ for every affine flat group scheme $Q$
of finite presentation over $S$.  This already fails for
\[
 Q=\underline G_{a_S},\qquad
 S=\Spec k[t]/(t^2),
\]
where $k$ is a field of characteristic $p>0$.


% END INCLUDED SGA7I_English_current_buildable_source_20260823/components/106_expose_VIII_section_3_3.tex

% BEGIN INCLUDED SGA7I_English_current_buildable_source_20260823/components/107_expose_VIII_section_3_4.tex
% SOURCE AUTHORITY:
% expose_VIII_body.tex lines 1483--1508; scan indices 270--273;
% source folios 259--262.
% Direct scan comparison: physical PDF pages 271--274 in four overlapping
% 1100-dpi bands per page.  The pages are clear at this resolution and no
% higher-resolution ambiguity pass was required.
% Transparent source disposition:
% - in the proof of Proposition 3.4, the printed phrase "une section f de P
%   sur P" is rendered as the mathematically required section of E over P;
%   the following displayed cocycle condition and Lemma 3.4.2 make the
%   intended object unambiguous.
% Hard stop before Corollary 3.5, source line 1510 / scan index 273.

\medskip
\noindent\underline{Proposition \SGARefTarget{sga7.sga7i.eviii.n3-4}{3.4}.}  \underline{Let $P$ be a smooth
group scheme of finite presentation over $S$, with connected fibres,
and let $Q$ be a torus over $S$.  Then the category of biextensions
of $(P,Q)$ by $\underline G_m$ is equivalent to the one-point
category; that is,}
\begin{equation}\tag{\SGARefTarget{sga7.sga7i.eviii.n3-4-1}{3.4.1}}
 \Biext^0(P,Q;\underline G_m)=
 \Biext^1(P,Q;\underline G_m)=0.
\end{equation}

\noindent\underline{Proof.}  Put $M=D(Q)$.  This is a twisted
constant group over $S$ whose fibres are isomorphic to groups
$\mathbb Z^r$ (\SGARefLink{sga7.sga7i.eviii.n3-3-3}{3.3.3}).  By \SGARefLink{sga7.sga7i.eviii.n3-3-2}{3.3.2}, the category in question is
equivalent to the category of extensions of $P$ by $M$.  This
category is \underline{rigid}: the automorphism group of any object
is the trivial group.  This is the first equality in (\SGARefLink{sga7.sga7i.eviii.n3-4-1}{3.4.1}), which
may also be written
\[
 \Hom(P,M)=0.
\]
It follows at once from the fact that $M$ is etale over $S$, whereas
$P$ has connected fibres.  It remains to prove
\[
 \Ext^1(P,M)=0,
\]
that is, that every extension $E$ of $P$ by $M$ is trivial.

By the uniqueness of a splitting when one exists, it is enough to
prove that $E$ is locally trivial.  We may therefore reduce to the
case $M=\mathbb Z_S^r$, and then to $M=\mathbb Z_S$.  When $S$ is the
spectrum of a field, the conclusion follows from \SGARefLink{sga7.sga7i.eviii.n5-5}{5.5} (i) below.
(Notice that we shall not use \SGARefLink{sga7.sga7i.eviii.n3-4}{3.4} before IX 6, so in particular there
is no vicious circle.)  We now reduce the general case to this
special case.

It is enough to find a section $f$ of $E$ over $P$ such that
$f(0)=0$, for such a section is necessarily multiplicative.  Indeed,
\[
 f(x+y)-f(x)-f(y)=0:
\]
the left-hand side defines a morphism
$P\times_S P\longrightarrow\mathbb Z_S$ which vanishes on the zero
section of $P\times_S P$.  It therefore vanishes identically, since
$P\times_S P$ has connected fibres (the geometric fibres of $P$
being connected).  The required conclusion now follows from the
following lemma.

\medskip
\noindent\underline{Lemma \SGARefTarget{sga7.sga7i.eviii.n3-4-2}{3.4.2}.}  \underline{Let $S$ be a scheme,
$P$ a smooth $S$-scheme of finite presentation with geometrically
connected fibres, $Z$ a group, and $E$ a $Z_P$-torsor (for the fppf,
etale, or fpqc topology; these notions are equivalent here, cf. SGA
3 X 5.4).  Let $e$ be a section of $P$ over $S$, and let $e'$ be a
section of $E$ over $S$ above $e$ (that is, a section of $e^{-1}(E)$
over $S$).}

\underline{The torsor $E$ is trivial if and only if, for every
$s\in S$, the induced torsor $E_s$ over $P_s$ is trivial.  In that
case there is a unique section $f$ of $E$ over $P$ such that
$f\circ e=e'$.}

\noindent\underline{Proof.}  Uniqueness is clear.  Two such sections
$f$ and $g$ ``differ'' by a morphism
\[
 h:P\longrightarrow Z_S,
 \qquad h\circ e=1,
\]
and such a morphism is constantly equal to $1$, since $P$ has
connected fibres.  Moreover, if a section $f$ of $E$ over $P$
exists, it can always be corrected by a morphism
$P\longrightarrow Z_S$ coming from a morphism $z:S\longrightarrow
Z_S$, so as to obtain a section $f'$ satisfying $f'\circ e=e'$.
One need only choose $z$ so that
\[
 f(e(t))=e'(t)z(t).
\]

Applying this observation fibrewise, for every $s\in S$ we obtain a
well-determined section $f_s$ of $E_s$ over $P_s$ taking $e_s$ to
$e'_s$.  It remains to prove that the $f_s$ arise from a section $f$
of $E$ over $P$.  Equivalently (EGA IV 17.9.3), we must prove that
\[
 U=\bigcup_{s\in S}f_s(P_s)\subset E
\]
is open in $E$.

Suppose first that $S$ is noetherian.  Then $U$ is constructible.  To
see this, for each $s\in S$ let $T$ be the closure of $s$, endowed
with its reduced induced structure.  There is an open neighbourhood
$T'$ of $s$ in $T$ and a section of $E_{T'}$ over $P_{T'}$ extending
$f_s$ and compatible with the marked sections $e$ and $e'$ (as
follows directly from EGA IV 8.8.2).  By the criterion of EGA IV
1.10.1, it remains to show that $U$ is stable under generization.
The usual argument using EGA II 7.1.9 reduces this to the case in
which $S$ is the spectrum of a discrete valuation ring.

Since $P$ is smooth over $S$, it is regular and hence normal; and
since it has connected fibres, it is connected.  The given section
of $E$ over the generic fibre of $P$ over $S$ therefore extends
uniquely to a section $f$ of $E$ over $P$ (EGA IV 18.10.19 and
18.10.20), necessarily compatible with $e$ and $e'$.  Thus in this
case $U=f(P)$ is open.  This proves \SGARefLink{sga7.sga7i.eviii.n3-4}{3.4} when $S$ is noetherian.  The
general case reduces readily to this one by the usual limit
arguments of EGA IV, 8 and 9; we leave the details to the reader.

\medskip
\noindent\underline{Remark \SGARefTarget{sga7.sga7i.eviii.n3-4-3}{3.4.3}.}  Lemma 3.4.2 can be ``crushed''
under learned references, namely SGA 4 XV 1.15 (the case $n=1$) and
\SGARefLink{sga7.sga7i.eviii.n4-1}{4.1}.


% END INCLUDED SGA7I_English_current_buildable_source_20260823/components/107_expose_VIII_section_3_4.tex

% BEGIN INCLUDED SGA7I_English_current_buildable_source_20260823/components/108_expose_VIII_sections_3_5_through_3_7.tex
% SOURCE AUTHORITY:
% expose_VIII_body.tex lines 1510--1553; scan indices 273--275;
% source folios 262--264.
% Direct scan comparison: physical PDF pages 274--276, with pages 275--276
% inspected in four overlapping 1100-dpi bands.  The typography is clear at
% that resolution; no 5000--9000-dpi ambiguity escalation was needed.
% Transparent source dispositions:
% - formula (3.6.2) has Biext^i on both sides.  The printed right-hand
%   exponent 1 is incompatible with the following qualifier i=0,1 and with
%   the asserted equivalence of categories;
% - in the proof of 3.7 the second cyclic case is M=Z/nZ, not the printed
%   phrase "Z=Z/nZ".
% Hard stop before section 4, source line 1559 / scan index 276.

\medskip
\noindent\underline{Corollary \SGARefTarget{sga7.sga7i.eviii.n3-5}{3.5}.}  \underline{Let $P$ be a smooth
group scheme of finite presentation over $S$, with connected fibres,
and let}
\[
 0\longrightarrow Q'\longrightarrow Q\longrightarrow Q''
 \longrightarrow0
\]
\underline{be an exact sequence of commutative groups over $S$, with
$Q'$ a torus.  Then the inverse-image functor on biextensions}
\begin{equation}\tag{\SGARefTarget{sga7.sga7i.eviii.n3-5-1}{3.5.1}}
 \operatorname{BIEXT}(P,Q'';\underline G_m)
 \xrightarrow{\ \sim\ }
 \operatorname{BIEXT}(P,Q;\underline G_m)
\end{equation}
\underline{is an equivalence of categories.  In particular, the
following natural homomorphisms are isomorphisms:}
\begin{equation}\tag{\SGARefTarget{sga7.sga7i.eviii.n3-5-2}{3.5.2}}
 \Biext^i(P,Q'';\underline G_m)
 \xrightarrow{\ \sim\ }
 \Biext^i(P,Q;\underline G_m),
 \qquad i=0,1.
\end{equation}

Indeed, up to equivalence, $\operatorname{BIEXT}(P,Q'';\underline
G_m)$ identifies with the category of biextensions $E$ of $(P,Q)$ by
$\underline G_m$ \underline{equipped} with a trivialization of the
induced biextension $E'$ of $(P,Q')$ by $\underline G_m$ (VII
\SGARefLink{sga7.sga7i.evii.n3-7-6}{3.7.6}).  By \SGARefLink{sga7.sga7i.eviii.n3-4}{3.4} there is exactly one such trivialization of $E'$ for
each given biextension $E$.  One can of course also prove directly
that (\SGARefLink{sga7.sga7i.eviii.n3-5-2}{3.5.2}) is an isomorphism for $i=0,1$ - which is equivalent to
the conclusion of \SGARefLink{sga7.sga7i.eviii.n3-5}{3.5} - by using the exact sequence of the
$\Biext^i$ in VII \SGARefLink{sga7.sga7i.evii.n3-7-5}{3.7.5} associated with
\[
 0\longrightarrow Q'\longrightarrow Q\longrightarrow Q''
 \longrightarrow0.
\]

\medskip
\noindent\underline{Remark \SGARefTarget{sga7.sga7i.eviii.n3-6}{3.6}.}  The symmetric statement obtained
from \SGARefLink{sga7.sga7i.eviii.n3-5}{3.5} by interchanging left and right is of course also valid; it
also follows formally from \SGARefLink{sga7.sga7i.eviii.n3-5}{3.5} by passing to symmetric biextensions
(VII \SGARefLink{sga7.sga7i.evii.n2-7}{2.7}).  Suppose, in addition to the data and hypotheses of \SGARefLink{sga7.sga7i.eviii.n3-5}{3.5},
that the symmetric data and hypotheses are given: namely, that $Q$
has connected fibres (equivalently, that $Q''$ has connected fibres,
since $Q'$ has connected fibres), and that there is an exact sequence
\[
 0\longrightarrow P'\longrightarrow P\longrightarrow P''
 \longrightarrow0
\]
with $P'$ a torus.  Successive application of \SGARefLink{sga7.sga7i.eviii.n3-5}{3.5} and of its
symmetric form, applied to this last sequence and to $Q''$, shows
that the functor
\begin{equation}\tag{\SGARefTarget{sga7.sga7i.eviii.n3-6-1}{3.6.1}}
 \operatorname{BIEXT}(P'',Q'';\underline G_m)
 \xrightarrow{\ \sim\ }
 \operatorname{BIEXT}(P,Q;\underline G_m)
\end{equation}
is an equivalence of categories.  In particular, the following
homomorphism is an isomorphism:
\begin{equation}\tag{\SGARefTarget{sga7.sga7i.eviii.n3-6-2}{3.6.2}}
 \Biext^i(P'',Q'';\underline G_m)
 \xrightarrow{\ \sim\ }
 \Biext^i(P,Q;\underline G_m),
 \qquad i=0,1.
\end{equation}

\medskip
\noindent\underline{Proposition \SGARefTarget{sga7.sga7i.eviii.n3-7}{3.7}.}  \underline{Let $A$ be an
abelian scheme over $S$, let $T$ be a group scheme of multiplicative
type (SGA 3 IX) and of finite type over $S$, and let}
\[
 \underline M=\uHom(T,\underline G_m)
\]
\underline{be the Cartier dual of $T$ (a finitely generated twisted
constant group over $S$, by \SGARefLink{sga7.sga7i.eviii.n3-3-3}{3.3.3}).  Then the categories
$\operatorname{EXT}(A,T)$ and
$\operatorname{BIEXT}(A,\underline M;\underline G_m)$, formed
respectively by commutative extensions of $A$ by $T$ and by
biextensions of $(A,\underline M)$ by $\underline G_m$, are rigid and
canonically equivalent.  Their groups of isomorphism classes are
canonically isomorphic to $\Hom(\underline M,A^*)$, where}
\[
 A^*=\uExt^1(A,\underline G_m)
\]
\underline{is the dual abelian scheme of $A$ (\SGARefLink{sga7.sga7i.eviii.n3-2}{3.2}).  In particular,
there is a canonical isomorphism}
\begin{equation}\tag{\SGARefTarget{sga7.sga7i.eviii.n3-7-1}{3.7.1}}
 \Ext^1(A,T)\xrightarrow{\ \sim\ }
 \Hom(\underline M,A^*).
\end{equation}

This is the special case of \SGARefLink{sga7.sga7i.eviii.n1-6}{1.6} obtained by taking $P=A$ and
$Q=\underline M$.  By the biduality theorem (SGA 3 X 5.7), we do
indeed have
\[
 T\simeq\uHom(\underline M,\underline G_m),
\]
and it remains only to verify the conditions (\SGARefLink{sga7.sga7i.eviii.n1-6-1}{1.6.1}), which here read
\[
 \uHom(A,\underline G_m)=0,
 \qquad
 \Hom\bigl(A,\uExt^1(\underline M,\underline G_m)\bigr)=0.
\]
The first equality was recalled in (\SGARefLink{sga7.sga7i.eviii.n3-2-1}{3.2.1}); the second may be
strengthened to
\[
 \uExt^1(\underline M,\underline G_m)=0.
\]
This assertion is local on $S$, so we may suppose that $\underline
M$ is constant, $\underline M=M_S$.  Decomposing $M$ as a direct sum
of cyclic $\mathbb Z$-modules reduces us to $M=\mathbb Z$, which is
trivial, and to $M=\mathbb Z/n\mathbb Z$, which follows from the
$n$-divisibility of $\underline G_m$ (or, more learnedly, from
\SGARefLink{sga7.sga7i.eviii.n3-3-1}{3.3.1}).

When $T=\underline G_m$, and hence $\underline M=\mathbb Z_S$,
(\SGARefLink{sga7.sga7i.eviii.n3-7-1}{3.7.1}) reduces to the identity isomorphism (\SGARefLink{sga7.sga7i.eviii.n3-2-3}{3.2.3}).  Of course, a
direct proof of (\SGARefLink{sga7.sga7i.eviii.n3-7-1}{3.7.1}) would be essentially trivial, without
introducing the group
\[
 \Ext^1\bigl(A\mathbin{\mathop{\otimes}^{\mathbb L}}\underline M,
              \underline G_m\bigr)
 =\Biext^1(A,\underline M;\underline G_m).
\]


% END INCLUDED SGA7I_English_current_buildable_source_20260823/components/108_expose_VIII_sections_3_5_through_3_7.tex

% BEGIN INCLUDED SGA7I_English_current_buildable_source_20260823/components/109_expose_VIII_sections_4_1_through_4_4.tex
% SOURCE AUTHORITY:
% expose_VIII_body.tex lines 1559--1575; scan indices 276--277;
% source folios 265--266.
% Direct scan comparison: physical PDF pages 277--278 in four overlapping
% 1100-dpi bands.  A targeted crop of the formula in Proposition 4.2 confirms
% that the printed argument is P in all three pullbacks, although the sentence
% has just introduced the line bundle L.  The mathematically type-correct L is
% therefore used below and the product is written explicitly as a tensor
% product.  No other high-resolution ambiguity remained.
% Hard stop before Lemma 4.5, source line 1581 / scan index 278.

\subsection*{\SGARefTarget{sga7.sga7i.eviii.n4}{4}. Extensions and biextensions of smooth connected group
schemes over a field}

We recall the following lemma due to Rosenlicht [3 bis, Proposition
3].

\medskip
\noindent\underline{Lemma \SGARefTarget{sga7.sga7i.eviii.n4-1}{4.1}.}  \underline{Let $k$ be a field, and
let $X$ and $Y$ be two separable, geometrically connected $k$-schemes
of finite type, equipped respectively with $k$-rational points $e_X$
and $e_Y$.  Then every morphism}
\[
 X\times Y\longrightarrow\underline G_m
\]
\underline{which is trivial on the two subschemes $X\times e_Y$ and
$e_X\times Y$ is trivial.}

We shall use this result only when $X$ and $Y$ are group schemes over
$k$; in that case the hypotheses imply that $X$ and $Y$ are in fact
smooth over $k$.  By induction, \SGARefLink{sga7.sga7i.eviii.n4-1}{4.1} also implies the analogous result
for any number of factors.  We may therefore apply the results of VII
\SGARefLink{sga7.sga7i.evii.n1-3}{1.3} and VII \SGARefLink{sga7.sga7i.evii.n2-9}{2.9}, and obtain the following two propositions.

\medskip
\noindent\underline{Proposition \SGARefTarget{sga7.sga7i.eviii.n4-2}{4.2}.}  \underline{Let $P$ be a
smooth connected group over the field $k$.  The natural functor from
the category of group extensions of $P$ by $\underline G_m$ to the
category of line bundles on $P$ rigidified at the identity element
$e_P$ is fully faithful.  Its essential image consists of the
rigidified line bundles $\mathbb L$ on $P$ for which}
\[
 \pi^*\mathbb L\simeq
 \pr_1^*\mathbb L\otimes\pr_2^*\mathbb L,
\]
\underline{where $\pi:P\times P\longrightarrow P$ is the group law
of $P$, and $\pr_1,\pr_2$ are the two projections.}

\medskip
\noindent\underline{Proposition \SGARefTarget{sga7.sga7i.eviii.n4-3}{4.3}.}  \underline{Let $P$ and $Q$ be
two smooth connected algebraic groups over the field $k$.  The
natural functor from the category of biextensions of $(P,Q)$ by
$\underline G_m$ to the category of line bundles on $P\times Q$
birigidified with respect to the identity sections $e_P$ and $e_Q$ of
$P$ and $Q$ is fully faithful; moreover, both categories are rigid.
A birigidified line bundle $\mathbb L$ on $P\times Q$ comes from a
biextension if and only if its pullbacks to $P\times P\times Q$ and
$P\times Q\times Q$ by $\pi_P\times\id_Q$ and
$\id_P\times\pi_Q$ satisfy the ``theorem of the cube'' ([1] or
[2]).}

This last condition is simply the translation of the conditions in
VII \SGARefLink{sga7.sga7i.evii.n2-9-3}{2.9.3} b), in the explicit form given in VII \SGARefLink{sga7.sga7i.evii.n3-5}{3.5} b).

\medskip
\noindent \SGARefTarget{sga7.sga7i.eviii.n4-4}{4.4}.  I do not know whether the cube condition in \SGARefLink{sga7.sga7i.eviii.n4-3}{4.3} is
automatically satisfied.  At least it is known [2, IV 2.6] that, for
a given $\mathbb L$, there is always an integer $n\geq1$, a power of
the characteristic exponent of $k$, such that
$\mathbb L^{\otimes n}$ satisfies that condition and hence comes from
a biextension of $(P,Q)$ by $\underline G_m$.  It is also always
possible to find a purely inseparable extension $k'$ of $k$ such that
the required condition becomes true after the base extension
$k\longrightarrow k'$; after this extension, therefore, $\mathbb L$
is a birigidified line bundle arising from a biextension of $(P,Q)$
by $\underline G_m$.

In particular, when $k$ is perfect, the condition in \SGARefLink{sga7.sga7i.eviii.n4-3}{4.3} is always
satisfied.  This last fact, which formally implies the preceding
statements by a trace argument, can also be seen by applying \SGARefLink{sga7.sga7i.eviii.n4-7}{4.7}
below; cf. \SGARefLink{sga7.sga7i.eviii.n4-9}{4.9}.


% END INCLUDED SGA7I_English_current_buildable_source_20260823/components/109_expose_VIII_sections_4_1_through_4_4.tex

% BEGIN INCLUDED SGA7I_English_current_buildable_source_20260823/components/110_expose_VIII_sections_4_5_and_4_6.tex
% SOURCE AUTHORITY:
% expose_VIII_body.tex lines 1581--1605; scan indices 278--280;
% source folios 267--269.
% Direct scan comparison: physical PDF pages 279--281 in four overlapping
% 1100-dpi bands per page.  Prose and formulas are fully legible at that
% resolution.
% Transparent source dispositions in the second case of Proposition 4.6:
% - the descent base is Y=P x B, not the printed Y=P x A;
% - after reducing to the case in which Q is abelian, the rigidifying
%   restriction is P x e_Q, not the printed P x e_A.
% Both readings are forced by the displayed morphism to Pic_{Q/k} and by the
% immediately preceding extension 1 -> M -> Q -> B -> 1.
% Hard stop before Proposition 4.7, source line 1607 / scan index 280.

We recall, in the form of a lemma, the following result from EGA
$\mathrm{Err}_{\mathrm{IV}}$ 53, 21.4.13.

\medskip
\noindent\underline{Lemma \SGARefTarget{sga7.sga7i.eviii.n4-5}{4.5}.}  \underline{Let
$f:X\longrightarrow Y$ be a smooth surjective morphism with integral
fibres, where $Y$ is normal and integral with generic point $y$.  Let
$\mathbb L$ be a line bundle on $X$ whose restriction to $X_y$ is
trivial.  Then there is a line bundle $\mathbb M$ on $Y$ such that}
\[
 \mathbb L\simeq f^*\mathbb M.
\]

\medskip
\noindent\underline{Proposition \SGARefTarget{sga7.sga7i.eviii.n4-6}{4.6}.}  \underline{Let $P,Q$ be two
smooth connected algebraic groups over the field $k$, with $P$
affine.  Suppose in addition either that $P$ admits a composition
series each of whose factors is a torus or is isomorphic to the
additive group $\underline G_a$ (a condition satisfied in particular
when $k$ is perfect [SGA 3, XVII 7.2.1, 4.1.1 (i)$\Longleftrightarrow$
(ii), and 4.1.5]), or that $Q$ is an extension of an abelian variety
$B$ by a smooth connected affine algebraic group satisfying the
preceding condition.  Under either hypothesis, every birigidified
line bundle $\mathbb L$ on $P\times Q$ is trivial; a fortiori, by
\SGARefLink{sga7.sga7i.eviii.n4-3}{4.3}, every biextension of $(P,Q)$ by $\underline G_m$ is trivial.}

\noindent\underline{Proof.}  First observe that every automorphism of
a birigidified line bundle is the identity (\SGARefLink{sga7.sga7i.eviii.n4-3}{4.3}).  Descent theory
therefore shows that it is enough to prove that $\mathbb L$ becomes
trivial after a finite separable extension of $k$.  Passing to the
limit immediately reduces us to the case where $k$ is separably
closed.

Every torus over such a field is split.  Hence, if $P$ admits a
composition series as in \SGARefLink{sga7.sga7i.eviii.n4-6}{4.6}, it also admits one whose quotients are
isomorphic to $\underline G_a$ or $\underline G_m$.  Repeated
application of \SGARefLink{sga7.sga7i.eviii.n4-5}{4.5}, together with
\[
 \Pic(\underline G_{a_K})=\Pic(\underline G_{m_K})=0
\]
for every field $K$, gives
\[
 \Pic(P)=0.
\]
This equality remains true after every extension of the base field
$k$.  Applying \SGARefLink{sga7.sga7i.eviii.n4-5}{4.5} to
\[
 \pr_2:P\times Q\longrightarrow Q
\]
and to $\mathbb L$, we find a line bundle $\mathbb M$ on $Q$ such
that $\mathbb L\simeq\pr_2^*\mathbb M$.  Since the restriction of
$\mathbb L$ to $e_P\times Q$ is trivial, $\mathbb M$ is trivial, and
so is $\mathbb L$.  This proves the first case.

In the second case, let $Q$ be an extension of an abelian scheme $B$
by a group $M$ admitting a composition series with quotients
isomorphic to $\underline G_m$ or $\underline G_a$.  Apply \SGARefLink{sga7.sga7i.eviii.n4-5}{4.5} to
\[
 f:X=P\times Q\longrightarrow Y=P\times B.
\]
The fibre over the generic point $y$ of $P\times B$ is an
$M_y$-torsor.  It is necessarily trivial by the particular structure
of $M$ and the classical equalities
\[
 H^1(K,\underline G_a)=0,
 \qquad
 H^1(K,\underline G_m)=0,
\]
valid for every field $K$, here taken to be $k(y)$.  Since
$\Pic(M)=0$, we have $\Pic(X_y)=0$; thus \SGARefLink{sga7.sga7i.eviii.n4-5}{4.5} applies and shows that
$\mathbb L$ is the pullback of a line bundle $\mathbb N$ on
$Y=P\times B$.

Let $\mathbb N_1$ and $\mathbb N_2$ be the restrictions of $\mathbb
N$ to $P\times e_B$ and $e_P\times B$, respectively.  Replacing
$\mathbb N$ by
\[
 \mathbb N\otimes\pr_1^*(\mathbb N_1)^{-1}
 \otimes\pr_2^*(\mathbb N_2)^{-1},
\]
we may suppose that $\mathbb N$ is birigidified and that $\mathbb L$
is its birigidified pullback.  We are thereby reduced to the case in
which $Q$ itself is an abelian scheme.

Then $\mathbb L$ defines a canonical morphism
\begin{equation}\tag{*}
 P\longrightarrow\underline{\Pic}_{Q/k}
\end{equation}
taking the identity to the identity.  Since $P$ is connected, this
morphism factors through the dual abelian scheme
\[
 Q'=\underline{\Pic}^{,0}_{Q/k}.
\]
It is well known that every morphism of schemes from a smooth
connected affine algebraic group to an abelian scheme is constant.
This is a geometric assertion, so we may suppose $k$ algebraically
closed; a devissage then reduces it immediately to $P=\underline
G_a$ or $P=\underline G_m$, and every rational map from
$\mathbb P^1_k$ to an abelian scheme over $k$ is constant.

The constancy, hence vanishing, of (*) means that $\mathbb L$ is the
pullback of a line bundle $\mathbb M$ on $P$.  Since $\mathbb L$ is
trivial on $P\times e_Q$, $\mathbb M$ is trivial, and therefore so is
$\mathbb L$.  This completes the proof.


% END INCLUDED SGA7I_English_current_buildable_source_20260823/components/110_expose_VIII_sections_4_5_and_4_6.tex

% BEGIN INCLUDED SGA7I_English_current_buildable_source_20260823/components/111_expose_VIII_sections_4_7_through_4_10.tex
% SOURCE AUTHORITY:
% expose_VIII_body.tex lines 1607--1696; scan indices 280--284;
% source folios 269--273.
% Direct scan comparison: physical PDF pages 281--284 in overlapping
% 1100-dpi bands.  The prose and displayed formulas are fully legible at
% that resolution; no 5000--9000-dpi ambiguity escalation was needed.
% Source disposition in the proof of Corollary 4.8: the print displays
% Biext^i(P,M;G_m).  This is retained.  Together with (4.7.1), applied with
% Q=M, those vanishings give the Biext^i(A,M;G_m) vanishings required by
% the exact sequence for 0 -> M -> Q -> B -> 0.
% Hard stop before 4.11, source line 1701 / scan index 284.

\medskip
\noindent\underline{Proposition \SGARefTarget{sga7.sga7i.eviii.n4-7}{4.7}.}  \underline{Let $P$ and $Q$
be two smooth connected algebraic groups over the field $k$.
Suppose that $P$ is an extension of an abelian scheme $A$ by a smooth
connected affine algebraic group $L$ admitting a composition series
each of whose factors is isomorphic to a torus or to $\underline
G_a$ (a condition always satisfied when $k$ is perfect).  Then the
functor considered in \SGARefLink{sga7.sga7i.eviii.n1-3}{1.3} is an equivalence of categories; that is,
every birigidified line bundle $\mathbb L$ on $P\times Q$ comes from
a biextension of $(P,Q)$.  Moreover, the functor between the rigid
categories}
\[
 \operatorname{BIEXT}(A,Q;\underline G_m)
 \longrightarrow
 \operatorname{BIEXT}(P,Q;\underline G_m)
\]
\underline{is an equivalence of categories; in other words, it
induces an isomorphism of groups}
\begin{equation}\tag{\SGARefTarget{sga7.sga7i.eviii.n4-7-1}{4.7.1}}
 \Biext^1(A,Q;\underline G_m)
 \xrightarrow{\ \sim\ }
 \Biext^1(P,Q;\underline G_m).
\end{equation}

For the first assertion, descent and \SGARefLink{sga7.sga7i.eviii.n4-3}{4.3} again reduce us to the case
where $k$ is separably closed.  Proceeding as in the proof of \SGARefLink{sga7.sga7i.eviii.n4-6}{4.6},
we deduce from \SGARefLink{sga7.sga7i.eviii.n4-5}{4.5} and the hypothesis on $L$ that $\mathbb L$ is
isomorphic to the pullback of a birigidified line bundle $\mathbb M$
on $A\times Q$.  For the first assertion of \SGARefLink{sga7.sga7i.eviii.n4-7}{4.7} we are thus reduced
to the case $A=P$.  The line bundle $\mathbb L$ then defines a
morphism of pointed schemes
\[
 f:Q\longrightarrow\underline{\Pic}_{P/k},
\]
which necessarily factors through the dual abelian scheme $P'$ of
$P$.  Let $\mathbb L'$ be the pullback of the Weil sheaf on
$P\times P'$ along
\[
 \id_P\times f:P\times Q\longrightarrow P\times P'.
\]
This is a birigidified line bundle which, by the definition of $f$,
differs from $\mathbb L$ only by the pullback of a line bundle
$\mathbb M$ on $P$.  Since both $\mathbb L'$ and $\mathbb L$ have
trivial restriction to $P\times e_Q$, it follows that $\mathbb M$ is
trivial and hence that $\mathbb L\simeq\mathbb L'$.  We noted in VII
\SGARefLink{sga7.sga7i.evii.n2-9-4}{2.9.4} that divisorial correspondences on a product of two abelian
schemes arise from biextensions.  Consequently $\mathbb L'$, and
therefore $\mathbb L$, comes from a biextension.

Returning to the general case, where $P$ is an extension of $A$ by
$L$, it remains to prove that (\SGARefLink{sga7.sga7i.eviii.n4-7-1}{4.7.1}) is an isomorphism.  This follows
from the exact sequence of the $\Biext^i$, for $i=0,1$, associated
with
\[
 0\longrightarrow L\longrightarrow P\longrightarrow A
 \longrightarrow0,
\]
and from the relations
\[
 \Biext^0(L,Q;\underline G_m)=0,
 \qquad
 \Biext^1(L,Q;\underline G_m)=0.
\]
The first is a special case of the rigidity assertion in \SGARefLink{sga7.sga7i.eviii.n4-3}{4.3}, and
the second is precisely the first case of \SGARefLink{sga7.sga7i.eviii.n4-6}{4.6}.

\medskip
\noindent\underline{Corollary \SGARefTarget{sga7.sga7i.eviii.n4-8}{4.8}.}  \underline{Under the hypotheses
of \SGARefLink{sga7.sga7i.eviii.n4-7}{4.7}, suppose that $Q$ is an extension of an algebraic group $B$
(necessarily smooth and connected) by a smooth connected affine
algebraic group $M$.  Then the functor between the rigid categories}
\[
 \operatorname{BIEXT}(A,B;\underline G_m)
 \longrightarrow
 \operatorname{BIEXT}(P,Q;\underline G_m)
\]
\underline{is an equivalence of categories; in other words, it
induces an isomorphism of groups}
\begin{equation}\tag{\SGARefTarget{sga7.sga7i.eviii.n4-8-1}{4.8.1}}
 \Biext^1(A,B;\underline G_m)
 \xrightarrow{\ \sim\ }
 \Biext^1(P,Q;\underline G_m).
\end{equation}

In view of (\SGARefLink{sga7.sga7i.eviii.n4-7-1}{4.7.1}), it remains to prove that
\[
 \Biext^1(A,B;\underline G_m)
 \xrightarrow{\ \sim\ }
 \Biext^1(A,Q;\underline G_m)
\]
is an isomorphism.  As in the proof of (\SGARefLink{sga7.sga7i.eviii.n4-7-1}{4.7.1}), this follows from the
exact sequence of the $\Biext^i$ associated with
\[
 0\longrightarrow M\longrightarrow Q\longrightarrow B
 \longrightarrow0,
\]
together with the relations
\[
 \Biext^0(P,M;\underline G_m)=0,
 \qquad
 \Biext^1(P,M;\underline G_m)=0.
\]
Indeed, (\SGARefLink{sga7.sga7i.eviii.n4-7-1}{4.7.1}), applied with $M$ in place of $Q$, identifies these
groups with the corresponding groups having $A$ in the first
variable, as required by the preceding exact sequence.  The second
vanishing displayed above is a special case of the second case of
\SGARefLink{sga7.sga7i.eviii.n4-6}{4.6}, after interchanging the two factors.

\medskip
\noindent\underline{Remark \SGARefTarget{sga7.sga7i.eviii.n4-9}{4.9}.}  Let $P$ and $Q$ be two smooth
connected algebraic groups over a \underline{perfect} field $k$.
It is then known [3] that each is an extension of an abelian scheme
by a smooth connected affine group, the latter being the product of
a torus with an iterated extension of copies of $\underline G_a$
(cf. the reference given in \SGARefLink{sga7.sga7i.eviii.n4-6}{4.6}):
\[
 0\longrightarrow L\longrightarrow P\longrightarrow A
 \longrightarrow0,
 \qquad
 0\longrightarrow M\longrightarrow Q\longrightarrow B
 \longrightarrow0.
\]
Thus \SGARefLink{sga7.sga7i.eviii.n4-8}{4.8} applies and shows that the theory of biextensions of $(P,Q)$
by $\underline G_m$ reduces to the analogous (and much more
classical!) theory for the pair $(A,B)$ of abelian varieties.

\medskip
\noindent\underline{Corollary \SGARefTarget{sga7.sga7i.eviii.n4-10}{4.10}.}  \underline{Let $P$ and $Q$ be
two smooth connected algebraic groups over $k$, let $L$ be a smooth
connected affine algebraic subgroup of $P$, and let $E$ be a
biextension of $(P,Q)$ by $\underline G_m$.  By \SGARefLink{sga7.sga7i.eviii.n2-2}{2.2}, $E$ induces a
pairing}
\[
 T_\ell(P)\times T_\ell(Q)\longrightarrow T_\ell(\underline G_m).
\]
\underline{Then the induced pairing}
\[
 T_\ell(L)\times T_\ell(Q)\longrightarrow T_\ell(\underline G_m)
\]
\underline{is zero.}

Indeed, we may plainly suppose that $k$ is algebraically closed.  We
are then under the hypotheses of \SGARefLink{sga7.sga7i.eviii.n4-6}{4.6}, which imply that the
restriction of $E$ to $(L,Q)$ is the trivial biextension.  The
asserted conclusion follows, since the induced pairing in question
is the pairing defined by this restriction.

% END INCLUDED SGA7I_English_current_buildable_source_20260823/components/111_expose_VIII_sections_4_7_through_4_10.tex

% BEGIN INCLUDED SGA7I_English_current_buildable_source_20260823/components/112_expose_VIII_section_4_11.tex
% SOURCE AUTHORITY:
% expose_VIII_body.tex lines 1701--1743; scan indices 284--286;
% source folios 273--275.
% Direct scan comparison: physical PDF pages 285--287 in overlapping
% 1100-dpi bands.  All formulas, primes, kernels, and arrow directions are
% legible at that resolution; no higher-resolution escalation was needed.
% Transparent source disposition: the reference printed as (1.11.1) in the
% middle of the discussion is the pairing (4.11.1), as the local argument
% and every surrounding reference require.
% Hard stop before 4.12, source line 1745 / scan index 286.

\medskip
\noindent \SGARefTarget{sga7.sga7i.eviii.n4-11}{4.11}.  Let $P$ and $Q$ be algebraic groups smooth over
$k$, let $E$ be a biextension of $(P,Q)$ by $\underline G_m$, and
let $T$ be the largest subtorus of $P$.  We shall determine the left
annihilator of the pairing
\begin{equation}\tag{\SGARefTarget{sga7.sga7i.eviii.n4-11-1}{4.11.1}}
 T_\ell(P)\times T_\ell(Q)
 \longrightarrow T_\ell(\underline G_{m,k}),
\end{equation}
where $\ell$ is again a fixed prime number.  By \SGARefLink{sga7.sga7i.eviii.n4-10}{4.10} this annihilator
always contains the subgroup $T_\ell(T)$ of $T_\ell(P)$.

Choose an extension $k'$ of $k$ over which $P$ and $Q$ can each be
written as an extension of an abelian scheme by an affine group
scheme that is an iterated extension of tori and copies of
$\underline G_a$.  If $P$ and $Q$ have unipotent rank zero, we may
take $k'=k$; in general, we may take for $k'$ the perfect closure of
$k$ (or a finite radicial extension).  Let $A'$ and $B'$ be the
abelian parts of $P_{k'}$ and $Q_{k'}$, respectively.  By \SGARefLink{sga7.sga7i.eviii.n4-7}{4.7}, the
biextension $E_{k'}$ of $(P_{k'},Q_{k'})$ comes from a biextension
$F'$ of $(A',B')$, uniquely determined up to unique isomorphism.
Consequently the pairing
\begin{equation}\tag{\SGARefTarget{sga7.sga7i.eviii.n4-11-2}{4.11.2}}
 T_\ell(P_{k'})\times T_\ell(Q_{k'})
 \longrightarrow T_\ell(\underline G_{m,k'}),
\end{equation}
obtained from (\SGARefLink{sga7.sga7i.eviii.n4-11-1}{4.11.1}) by the base-field extension $k\to k'$, and
also defined by the biextension $E_{k'}$, is obtained from the
pairing
\begin{equation}\tag{\SGARefTarget{sga7.sga7i.eviii.n4-11-3}{4.11.3}}
 T_\ell(A')\times T_\ell(B')
 \longrightarrow T_\ell(\underline G_{m,k'})
\end{equation}
defined by $F'$, by composing with the canonical homomorphisms
(which are in fact epimorphisms if $P$ and $Q$ are
$\ell$-divisible)
\begin{equation}\tag{\SGARefTarget{sga7.sga7i.eviii.n4-11-4}{4.11.4}}
 T_\ell(P_{k'})\longrightarrow T_\ell(A'),
 \qquad
 T_\ell(Q_{k'})\longrightarrow T_\ell(B').
\end{equation}

It follows that the left annihilator of (\SGARefLink{sga7.sga7i.eviii.n4-11-2}{4.11.2}), which is obtained
from the left annihilator of (\SGARefLink{sga7.sga7i.eviii.n4-11-1}{4.11.1}) by base change, is the inverse
image under the first homomorphism in (\SGARefLink{sga7.sga7i.eviii.n4-11-4}{4.11.4}) of the left
annihilator of (\SGARefLink{sga7.sga7i.eviii.n4-11-3}{4.11.3}).  Let $C'$ be the largest abelian subscheme
of $A'$ such that the restriction of $F'$ to $(C',B')$ is the
trivial biextension.  Equivalently, $C'$ is the reduced identity
component of the kernel of the homomorphism
\begin{equation}\tag{\SGARefTarget{sga7.sga7i.eviii.n4-11-5}{4.11.5}}
 A'\longrightarrow (B')^*
\end{equation}
defined by $F'$, where $(B')^*$ is the dual abelian scheme of $B'$.
The left annihilator of (\SGARefLink{sga7.sga7i.eviii.n4-11-3}{4.11.3}) identifies with $T_\ell(C')$;
therefore the left annihilator of (\SGARefLink{sga7.sga7i.eviii.n4-11-1}{4.11.1}) is precisely the inverse
image of
\[
 T_\ell(C')\lhook\joinrel\longrightarrow T_\ell(A')
\]
under the first morphism in (\SGARefLink{sga7.sga7i.eviii.n4-11-4}{4.11.4}).

In particular, to say that the left annihilator of (\SGARefLink{sga7.sga7i.eviii.n4-11-1}{4.11.1}) is
exactly $T_\ell(T)$ -- equivalently, that the left annihilator of
(\SGARefLink{sga7.sga7i.eviii.n4-11-2}{4.11.2}) is the kernel $T_\ell(T')$ of the first morphism in
(\SGARefLink{sga7.sga7i.eviii.n4-11-4}{4.11.4}) -- is to say that (\SGARefLink{sga7.sga7i.eviii.n4-11-3}{4.11.3}) is nondegenerate on the left,
that is, that $C'=0$, or equivalently that (\SGARefLink{sga7.sga7i.eviii.n4-11-5}{4.11.5}) has finite
kernel.  In this last form the condition is visibly independent of
$\ell$, and in any of its forms it is plainly independent of the
choice of $k'$.  When it holds, we shall simply say that the given
biextension $E$ of $(P,Q)$ by $\underline G_m$ is
\underline{nondegenerate on the left}.  By duality, this also means
that the homomorphism
\begin{equation}\tag{\SGARefTarget{sga7.sga7i.eviii.n4-11-6}{4.11.6}}
 B'\longrightarrow (A')^*
\end{equation}
defined by $F'$ is an epimorphism, that is, is surjective.

The notion of a biextension \underline{nondegenerate on the right}
is defined symmetrically: (\SGARefLink{sga7.sga7i.eviii.n4-11-6}{4.11.6}) must have finite kernel, or,
equivalently, (\SGARefLink{sga7.sga7i.eviii.n4-11-5}{4.11.5}) must be an epimorphism.  The biextension $E$
of $(P,Q)$ by $\underline G_{m,k}$ is called
\underline{nondegenerate} if it is nondegenerate on both the left
and the right.  This is equivalent to saying that (\SGARefLink{sga7.sga7i.eviii.n4-11-5}{4.11.5}), or
equivalently (\SGARefLink{sga7.sga7i.eviii.n4-11-6}{4.11.6}), is an isogeny.

It is immediate from the definition that these notions are
invariant under every extension of the base field $k$.

% END INCLUDED SGA7I_English_current_buildable_source_20260823/components/112_expose_VIII_section_4_11.tex

% BEGIN INCLUDED SGA7I_English_current_buildable_source_20260823/components/113_expose_VIII_sections_4_12_through_4_14.tex
% SOURCE AUTHORITY:
% expose_VIII_body.tex lines 1745--1799; scan indices 286--289;
% source folios 275--278.
% Direct scan comparison: physical PDF pages 287--290 in overlapping
% 1100-dpi bands.  The formulas and prose are fully legible at that
% resolution; no 5000--9000-dpi ambiguity escalation was needed.
% Transparent source dispositions in 4.12:
% - the canonical quotient is f:P -> A, not the printed f:A -> P; this is
%   forced by the displayed pullback Pic(A) -> Pic(P);
% - the descended relation is L \simeq f^*M, not the printed L^* relation;
% - M is a line bundle on A, not the printed "on M".
% In the last sentence of 4.14, the delta map on Pic(P) is (4.14.2), not
% the printed reference (4.14.1).
% Hard stop before section 5, source line 1800 / scan index 289.

\medskip
\noindent \SGARefTarget{sga7.sga7i.eviii.n4-12}{4.12}.  Let $P$ be a smooth connected group scheme, let
$\mathbb L$ be a line bundle on $P$, and put
\begin{equation}\tag{\SGARefTarget{sga7.sga7i.eviii.n4-12-1}{4.12.1}}
 \delta(\mathbb L)=
 \pi^*\mathbb L\otimes
 \pr_1^*(\mathbb L)^{-1}\otimes
 \pr_2^*(\mathbb L)^{-1},
\end{equation}
with the notation of \SGARefLink{sga7.sga7i.eviii.n4-2}{4.2}.  As already observed in VII \SGARefLink{sga7.sga7i.evii.n1-3-4}{1.3.4}, this is
a birigidified line bundle on $P\times P$.  By \SGARefLink{sga7.sga7i.eviii.n4-3}{4.3}, if
$\delta(\mathbb L)$ comes from a biextension, that biextension is
determined up to unique isomorphism.  In particular, by \SGARefLink{sga7.sga7i.eviii.n4-7}{4.7}, it is
defined when $P$ is an extension of an abelian scheme $A$ by an
iterated extension of tori and copies of $\underline G_a$.  Whenever
it is defined, this biextension is plainly symmetric.

In the preceding favourable case, and assuming that the maximal
torus of $P$ is split, the construction of $\delta(\mathbb L)$ can
be reduced to the case where $P$ is an abelian scheme.  Indeed, the
construction is functorial in $P$, while the homomorphism
\begin{equation}\tag{\SGARefTarget{sga7.sga7i.eviii.n4-12-2}{4.12.2}}
 \Pic(A)\longrightarrow\Pic(P)
\end{equation}
induced by the canonical quotient homomorphism
$f:P\longrightarrow A$ is surjective by \SGARefLink{sga7.sga7i.eviii.n4-5}{4.5} (cf. the beginning of
the proof of \SGARefLink{sga7.sga7i.eviii.n4-6}{4.6}).  We can therefore find an isomorphism
\[
 \mathbb L\simeq f^*\mathbb M,
\]
where $\mathbb M$ is a line bundle on $A$.  Then
$\delta(\mathbb L)$ is the pullback of the biextension
$\delta(\mathbb M)$ of $(A,A)$ by $\underline G_{m,k}$.

\medskip
\noindent\underline{Proposition \SGARefTarget{sga7.sga7i.eviii.n4-13}{4.13}.}  \underline{Under the
hypotheses of \SGARefLink{sga7.sga7i.eviii.n4-12}{4.12}, suppose that $\mathbb L$ is an ample line bundle
on $P$.  If the biextension $\delta(\mathbb L)$ of $(P,P)$ by
$\underline G_{m,k}$ is defined, then it is nondegenerate in the
sense of \SGARefLink{sga7.sga7i.eviii.n4-11}{4.11}.  If $P$ is an extension of an abelian scheme $A$ by a
smooth connected affine group $L$ satisfying the condition of \SGARefLink{sga7.sga7i.eviii.n4-7}{4.7},
then the biextension of $(A,A)$ from which $\delta(\mathbb L)$ comes
by \SGARefLink{sga7.sga7i.eviii.n4-8}{4.8} is of the form $\delta(\mathbb M)$, where $\mathbb M$ is an
ample line bundle on $A$.}

We reduce to the case where $k$ is algebraically closed.  By the
definition, it is then enough, with the notation of \SGARefLink{sga7.sga7i.eviii.n4-12}{4.12}, to prove
that the biextension $\delta(\mathbb M)$ of $(A,A)$ is
nondegenerate.  M. Raynaud proves in his thesis [2, XI 1.11] that if
$\mathbb L=f^*\mathbb M$ is ample, then $\mathbb M$ is ample as
well.  The theory of A. Weil then gives the familiar conclusion that
$\delta(\mathbb M)$ is nondegenerate [1, 2], that is, that the
homomorphism from $A$ to its dual abelian scheme is an isogeny.

\medskip
\noindent\underline{Remark \SGARefTarget{sga7.sga7i.eviii.n4-14}{4.14}.}  Suppose we are in the favourable
case where $P$ is an extension of an abelian scheme $A$ by a group
$L$ which is an iterated extension of tori and copies of
$\underline G_a$, and consider the canonical composite
\begin{equation}\tag{\SGARefTarget{sga7.sga7i.eviii.n4-14-1}{4.14.1}}
 \operatorname{NS}(A)
 \xrightarrow{\ \delta\ }
 \Biext^1(A,A;\underline G_m)
 \xrightarrow{\ \sim\ }
 \Biext^1(P,P;\underline G_m).
\end{equation}
By \SGARefLink{sga7.sga7i.eviii.n4-8}{4.8} the second arrow is bijective.  On the other hand, the theory
of abelian varieties shows that the first arrow identifies the
Neron--Severi group $\operatorname{NS}(A)$ of $A$ with the subgroup
of
\[
 \Biext^1(A,A;\underline G_m)=\operatorname{Corr}(A,A)
\]
consisting of the classes of symmetric biextensions, or equivalently
symmetric divisorial correspondences.  Thus (\SGARefLink{sga7.sga7i.eviii.n4-14-1}{4.14.1}) identifies
$\operatorname{NS}(A)$ with the group of symmetric biextensions of
$(P,P)$ by $\underline G_m$.  By \SGARefLink{sga7.sga7i.eviii.n4-7}{4.7} this group can equally be
described as the group of classes of symmetric divisorial
correspondences on $(P,P)$.

It is natural to call a divisorial correspondence on $(P,P)$
\underline{ample}, and likewise to call the corresponding
biextension of $(P,P)$ by $\underline G_m$ \underline{ample}, when
it comes from an ample element of $\operatorname{NS}(A)$.  This
condition is plainly invariant under extension of the base field;
the notion therefore still makes sense without the preceding
restrictive hypothesis on $P$, assuming only that $P$ is smooth and
connected.

Suppose further that $L$ is ``split solvable over $k$'', in the
sense that it admits a devissage by copies of $\underline G_a$ and
$\underline G_m$.  The preceding observation and the surjectivity of
(\SGARefLink{sga7.sga7i.eviii.n4-12-2}{4.12.2}) show that the image of
\begin{equation}\tag{\SGARefTarget{sga7.sga7i.eviii.n4-14-2}{4.14.2}}
 \Pic(P)\xrightarrow{\ \delta\ }
 \Biext^1(P,P;\underline G_m)
\end{equation}
consists of the symmetric biextensions, and that its kernel is the
subgroup $\Pic^0(P)$ of $\Pic(P)$ which is the image of $\Pic^0(A)$.
Consequently, if we put
\[
 \operatorname{NS}(P)=\Pic(P)/\Pic^0(P)
 \simeq\operatorname{NS}(A),
\]
we may also interpret (\SGARefLink{sga7.sga7i.eviii.n4-14-1}{4.14.1}) as a canonical isomorphism
\begin{equation}\tag{\SGARefTarget{sga7.sga7i.eviii.n4-14-3}{4.14.3}}
 \operatorname{NS}(P)
 \xrightarrow{\ \sim\ }
 \Biext^1(P,P;\underline G_m)^{\mathrm{sym}}.
\end{equation}

Finally, the theorem of Raynaud cited in \SGARefLink{sga7.sga7i.eviii.n4-13}{4.13} says that an element
$\xi$ of $\Pic(P)$ is ample if and only if its image
$\delta(\xi)$ under (\SGARefLink{sga7.sga7i.eviii.n4-14-2}{4.14.2}) is the class of an ample biextension.

% END INCLUDED SGA7I_English_current_buildable_source_20260823/components/113_expose_VIII_sections_4_12_through_4_14.tex

% BEGIN INCLUDED SGA7I_English_current_buildable_source_20260823/components/114_expose_VIII_section_5_opening_through_5_6.tex
% SOURCE AUTHORITY:
% expose_VIII_body.tex lines 1800--1916; scan indices 289--295;
% source folios 278--284.
% Direct scan comparison: physical PDF pages 290--296 in overlapping
% 1100-dpi bands.  The mathematical text and displayed formulas are fully
% legible at that resolution; no higher-resolution escalation was needed.
% Transparent source dispositions:
% - in the proof of 5.1, G_S and the final E are corrected to Z_S and Z;
% - the source repeats the target category in Corollary 5.2; the first
%   category is restored to torsors on P under Z_P;
% - the repeated P_i at the end of that proof is read as P_{i\eta};
% - in 5.5(ii), the second finite-type example is P, and every stray G_k or
%   G in the statement and proof is corrected to the stated coefficient Z;
% - in 5.6, Phi and Psi are groups of connected components (not identity
%   components), and the stray G_k in (5.6.1) is corrected to Z_k.
% Hard stop before 5.7, source line 1917 / scan index 295.

\medskip
\noindent\underline{5.  Extensions and biextensions by torsion-free
truncated constant group schemes.}

In this section and the next we give some auxiliary results that will
be used in section 7 in the case $Z=\mathbb Z$.

\medskip
\noindent\underline{Proposition \SGARefTarget{sga7.sga7i.eviii.n5-1}{5.1}.}  \underline{Let $S$ be an
irreducible geometrically unibranch scheme with generic point $\eta$,
and let $Z$ be a torsion-free group, hence giving a constant group
scheme $Z_S$ over $S$.  Then every torsor $E$ under $Z_S$ is trivial,
and the functor}
\[
 E\longmapsto E(\eta)
\]
\underline{from the category of torsors under $Z_S$ to the category of
torsors under $Z=Z_S(\eta)$ is an equivalence of categories.}

A familiar lemma of descent theory (SGA 3 X 5.4) shows that every
torsor $E$ under $Z_S$ is representable.  In fact, that result applies
to every group scheme over $S$ which is separated, locally of finite
presentation, and locally quasi-finite over $S$.  Since $Z_S$ is
etale over $S$, fpqc descent shows that the same is true of every
torsor $E$ under $Z_S$; such a torsor is therefore locally trivial for
the \underline{etale topology} (and not merely for the fppf topology).
To prove that $E$ is trivial, it is consequently enough to prove
\[
 H^1(S_{\mathrm{\acute{e}t}},Z_S)=e.
\]
Likewise, the full-faithfulness assertion of \SGARefLink{sga7.sga7i.eviii.n5-1}{5.1} reduces to the
relation
\[
 Z\xrightarrow{\ \sim\ }H^0(S_{\mathrm{\acute{e}t}},Z_S).
\]
The latter is immediate.  To prove the former, consider the inclusion
\[
 i:\eta\longrightarrow S.
\]
Because $S$ is geometrically unibranch, we have
\[
 Z_S\xrightarrow{\ \sim\ }i_*(Z_\eta)
\]
(SGA 4 IX 2.14.1).  The noncommutative low-degree exact sequence that
replaces the Leray spectral sequence for $i:\eta\to S$ (SGA 4 XII
3.2) therefore begins
\[
 e\longrightarrow H^1(S,i_*(Z_\eta))
   \longrightarrow H^1(\eta,Z_\eta)\longrightarrow\cdots,
\]
and shows that
\[
 H^1(S,Z_S)\longrightarrow H^1(\eta,Z_\eta)
\]
is injective.  We are thus reduced to the case in which $S$ is the
spectrum of a field $k$.  If $\pi$ is the Galois group of a separable
closure of $k$, then
\[
 H^1(k,Z)=\Hom_{\mathrm{cont}}(\pi,Z)/Z=e,
\]
because $Z$ is torsion-free.  This completes the proof.

\medskip
\noindent\underline{Corollary \SGARefTarget{sga7.sga7i.eviii.n5-2}{5.2}.}  \underline{With the notation of
\SGARefLink{sga7.sga7i.eviii.n5-1}{5.1}, let $P$ be a smooth scheme over $S$.  Then every torsor $E$ under
$Z_P$ is trivial, and the restriction functor}
\[
 E\longmapsto E_\eta
\]
\underline{from the category of torsors on $P$ under $Z_P$ to the
category of torsors on $P_\eta$ under $Z_{P_\eta}$ is an equivalence of
categories.}

Indeed, $P$ is geometrically unibranch (EGA IV 17.5.7).  Moreover,
every irreducible component of $P$ dominates $S$, because $P$ is flat
over $S$, and the set of irreducible components of $P$ is locally
finite.  Thus $P$ is the disjoint sum of irreducible geometrically
unibranch schemes $P_i$.  The assertion now follows immediately from
\SGARefLink{sga7.sga7i.eviii.n5-1}{5.1} applied to the various $P_i$ and their generic fibres $P_{i\eta}$,
noting that the generic point $p_i$ of $P_i$ is also the generic point
of $P_{i\eta}$.

\medskip
\noindent \SGARefTarget{sga7.sga7i.eviii.n5-2-1}{5.2.1}.  The preceding corollary immediately gives the
analogous statement for bitorsors on $P$ under $(Z_P,Z_P)$: restriction
to the generic fibre induces an equivalence from their category to the
category of bitorsors on $P_\eta$ under
$(Z_{P_\eta},Z_{P_\eta})$.  Indeed, it is enough to express the bitorsor
structure on $E$ in terms of its structure as a right torsor under
$Z_P$, together with an isomorphism from $Z_P$ to the commutant of
$Z_P$ for its action on $E$.  Since $E$ is necessarily trivial, that
commutant is itself isomorphic to $Z_P$.  We are thus reduced, on each
irreducible component $X$ of $P$, to the following immediate
observation: if $X$ is a nonempty connected scheme, the functor
$Z\mapsto Z_X$ from sets to $X$-schemes is fully faithful.

\medskip
\noindent\underline{Proposition \SGARefTarget{sga7.sga7i.eviii.n5-3}{5.3}.}  \underline{Let $S$ and $Z$ be
as in \SGARefLink{sga7.sga7i.eviii.n5-1}{5.1}, and let $P$ be a smooth group scheme over $S$.  Consider
extensions $E$ of group schemes of $P$ by $Z_S$.  Then restriction to
the generic fibre}
\[
 E\longmapsto E_\eta
\]
\underline{is an equivalence from the category of these extensions to
the category of extensions of $P_\eta$ by $Z_\eta$.  If $P$ and $Z$
are commutative, it also induces an equivalence between the
corresponding categories of commutative extensions.}

First observe that every extension of sheaves of groups of $P$ by
$Z_S$ (and similarly of $P_\eta$ by $Z_\eta$) gives, by \SGARefLink{sga7.sga7i.eviii.n5-2}{5.2}, a trivial
torsor, hence in particular a representable one.  The categories in
\SGARefLink{sga7.sga7i.eviii.n5-3}{5.3} may therefore equally well be interpreted as categories of
extensions of sheaves of groups.  The description in VII \SGARefLink{sga7.sga7i.evii.n1-1-6}{1.1.6} and
\SGARefLink{sga7.sga7i.evii.n1-2}{1.2} expresses these in terms of bitorsors on
$P$, $P\times P$, and $P\times P\times P$ (respectively on
$P_\eta$, $P_\eta\times P_\eta$, and
$P_\eta\times P_\eta\times P_\eta$), under the pairs of groups
obtained from $(Z,Z)$ by base change.  Applying \SGARefLink{sga7.sga7i.eviii.n5-2}{5.2} in the commutative
case and \SGARefLink{sga7.sga7i.eviii.n5-2-1}{5.2.1} in the general case to the smooth $S$-schemes
$P$, $P\times P$, and $P\times P\times P$ gives the asserted
conclusion.  The same argument gives the following result.

\medskip
\noindent\underline{Corollary \SGARefTarget{sga7.sga7i.eviii.n5-4}{5.4}.}  \underline{Let $S,Z$ be as in
\SGARefLink{sga7.sga7i.eviii.n5-1}{5.1}, and let $P,Q$ be two smooth group schemes over $S$.  Suppose that
$P,Q,$ and $Z$ are commutative, and consider biextensions $E$ of
$(P,Q)$ by $Z_S$.  Then restriction to the generic fibre}
\[
 E\longmapsto E_\eta
\]
\underline{is an equivalence from the category of these biextensions
to the category of biextensions of $(P_\eta,Q_\eta)$ by $Z_\eta$.}

The preceding results reduce the determination of categories of
extensions or biextensions by a torsion-free constant group to the
case in which the base scheme is the spectrum of a field.  In that
case we obtain the following.

\medskip
\noindent\underline{Proposition \SGARefTarget{sga7.sga7i.eviii.n5-5}{5.5}.}  \underline{Let $k$ be a field,
$Z$ a torsion-free group, $P$ a $k$-group scheme locally of finite
type, $P^\circ$ its identity component, and
$\Phi=P/P^\circ$ its group of connected components (SGA 3
$\mathrm{VI}_A$ 4).  Then:}

\smallskip
\noindent (i) \underline{The natural functor from the category of group
schemes which are extensions (respectively commutative extensions) of
$\Phi$ by $Z_k$ to the category of group schemes which are extensions
(respectively commutative extensions) of $P$ by $Z_k$ is an
equivalence.  In particular, if $P$ is connected, every extension of
$P$ by $Z_k$ is trivial, and the category of such extensions is
rigid.}

\smallskip
\noindent (ii) \underline{Suppose that $\Phi$ is a torsion group and
that $P$ and $Z$ are commutative.  Then the category of extensions of
$P$ by $Z_k$ is rigid: every automorphism of an object of this
category is trivial.  If, in addition, there is an integer $n>0$ such
that $n\id_\Phi=0$---for example, if $\Phi$ is finite, as happens when
$P$ is of finite type---then the group of classes of commutative
extensions of $P$ by $Z_k$ is given by the canonical isomorphism}
\begin{equation}\tag{\SGARefTarget{sga7.sga7i.eviii.n5-5-1}{5.5.1}}
 \Ext^1(P,Z_k)
 \xleftarrow{\ \sim\ }
 \Hom\bigl(\Phi,(Z\otimes_{\mathbb Z}\mathbb Q)/Z\bigr),
\end{equation}
\underline{defined by the boundary operator in the exact sequence of
$\Ext$ groups associated with the exact sequence}
\begin{equation}\tag{\SGARefTarget{sga7.sga7i.eviii.n5-5-2}{5.5.2}}
 0\longrightarrow Z\longrightarrow Z\otimes_{\mathbb Z}\mathbb Q
 \longrightarrow (Z\otimes_{\mathbb Z}\mathbb Q)/Z
 \longrightarrow 0.
\end{equation}

The canonical morphism
\[
 P\longrightarrow\Phi
\]
plainly induces an isomorphism on the sets of irreducible components,
since it is flat with irreducible fibres.  By \SGARefLink{sga7.sga7i.eviii.n5-1}{5.1}, the induced functor
from torsors under $Z_\Phi$ to torsors under $Z_P$ is therefore an
equivalence.  (Recall that a group scheme locally of finite type over
$k$ is geometrically unibranch.)  The same conclusion applies to
bitorsors and to the morphisms
\[
 P\times P\longrightarrow\Phi\times\Phi,
 \qquad
 P\times P\times P\longrightarrow\Phi\times\Phi\times\Phi.
\]
The description in VII \SGARefLink{sga7.sga7i.evii.n1-1-6}{1.1.6} of extensions of $P$ by $Z_k$ and of
$\Phi$ by $Z_k$ (and VII \SGARefLink{sga7.sga7i.evii.n1-3}{1.3} in the commutative case) now gives the
assertions in (i).

We restrict attention to commutative extensions.  The automorphism
group of an extension of $P$ by $Z_k$ is
\[
 \Hom(P,Z_k)\simeq\Hom(\Phi,Z_k),
\]
which is plainly zero when $\Phi$ is a torsion group.  Likewise, if
$V=Z\otimes_{\mathbb Z}\mathbb Q$ and $\Phi$ is annihilated by the
integer $n>0$, then
\[
 \Ext^i(\Phi,V_k)=0\qquad\text{for every }i\in\mathbb Z,
\]
because the groups in question are $\mathbb Q$-vector spaces
annihilated by $n$.  The isomorphism (\SGARefLink{sga7.sga7i.eviii.n5-5-1}{5.5.1}) follows at once.

The same method proves, essentially, the following.

\medskip
\noindent\underline{Corollary \SGARefTarget{sga7.sga7i.eviii.n5-6}{5.6}.}  \underline{Let $k$ be a field,
$Z$ a torsion-free commutative group, $P$ and $Q$ two commutative
group schemes locally of finite type over $k$, and let
$\Phi=P/P^\circ$ and $\Psi=Q/Q^\circ$ be their groups of connected
components.  Then:}

\smallskip
\noindent (i) \underline{The natural functor from the category of
biextensions of $(\Phi,\Psi)$ by $Z_k$ to the category of
biextensions of $(P,Q)$ by $Z_k$ is an equivalence.  In particular,
if $P$ or $Q$ is connected, the category of biextensions of $(P,Q)$
by $Z_k$ is equivalent to the terminal category.}

\smallskip
\noindent (ii) \underline{Suppose that $\Phi$ or $\Psi$ is a torsion
group.  Then the category of biextensions of $(P,Q)$ by $Z_k$ is
rigid.  If, in addition, there is an integer $n>0$ such that
$n\id_\Phi=0$ or $n\id_\Psi=0$, then the group of classes of
biextensions of $(P,Q)$ by $Z_k$ is given by the canonical
isomorphism}
\begin{equation}\tag{\SGARefTarget{sga7.sga7i.eviii.n5-6-1}{5.6.1}}
 \Biext^1(P,Q;Z_k)
 \xleftarrow{\ \sim\ }
 \Hom\bigl(\Phi\otimes\Psi,(Z\otimes_{\mathbb Z}\mathbb Q)/Z\bigr),
\end{equation}
\underline{defined by the boundary operator in the exact sequence of
$\Biext^i$ groups $(i=0,1)$ associated with (\SGARefLink{sga7.sga7i.eviii.n5-5-2}{5.5.2}).}


% END INCLUDED SGA7I_English_current_buildable_source_20260823/components/114_expose_VIII_section_5_opening_through_5_6.tex

% BEGIN INCLUDED SGA7I_English_current_buildable_source_20260823/components/115_expose_VIII_sections_5_7_through_5_10.tex
% SOURCE AUTHORITY:
% expose_VIII_body.tex lines 1917--1974; scan indices 295--298;
% source folios 284--287.
% Direct scan comparison: physical PDF pages 296--299 in overlapping
% 1100-dpi bands.  The formulas, subscripts, and arrow directions are fully
% legible at that resolution; no higher-resolution escalation was needed.
% Transparent source dispositions:
% - the reference printed as (2.7.2) in 5.7 is the local formula (5.7.2);
% - stray G_{S,T} occurrences in 5.8 and 5.9 are corrected to the group
%   Z_{S,T} fixed by the surrounding definitions and formulas.
% Hard stop before section 6, source line 1976 / scan index 298.

\medskip
\noindent \SGARefTarget{sga7.sga7i.eviii.n5-7}{5.7}.  Let $S$ be a scheme, $T$ a closed subscheme, and $Z$
a set.  The \underline{constant scheme over $S$, truncated along $T$,
with value $Z$}, denoted $Z_{S,T}$, is the scheme obtained by taking
the sum of copies $S_z$ of $S$, one for each $z\in Z$, and gluing all
these copies along the open subscheme $U=S\setminus T$.  Thus
\begin{equation}\tag{\SGARefTarget{sga7.sga7i.eviii.n5-7-1}{5.7.1}}
 Z_{S,T}\times_S U\simeq U,
 \qquad
 Z_{S,T}\times_S T\simeq Z_T,
\end{equation}
and $Z_{S,T}$ is etale over $S$.  For a variable $S$-scheme $S'$
there is an isomorphism, functorial in $S'$,
\begin{equation}\tag{\SGARefTarget{sga7.sga7i.eviii.n5-7-2}{5.7.2}}
 \Hom_S(S',Z_{S,T})\simeq\Hom_T(S'_T,Z_T).
\end{equation}
The map is restriction to $S'_T$, using the second isomorphism in
(\SGARefLink{sga7.sga7i.eviii.n5-7-1}{5.7.1}); we leave it to the reader to check that the map is bijective.
If
\[
 i:T\longrightarrow S
\]
denotes the inclusion, the functorial isomorphism (\SGARefLink{sga7.sga7i.eviii.n5-7-2}{5.7.2}) may also be
written as an isomorphism of sheaves
\begin{equation}\tag{\SGARefTarget{sga7.sga7i.eviii.n5-7-3}{5.7.3}}
 Z_{S,T}\simeq i_*(Z_T).
\end{equation}
This formula, or the definition directly, shows that as $Z$ varies,
$Z_{S,T}$ depends functorially on $Z$, and that this functor commutes
with finite inverse limits.  In particular, it takes groups to group
schemes and commutative groups to commutative group schemes, as is also
clear from the form (\SGARefLink{sga7.sga7i.eviii.n5-7-2}{5.7.2}) of the functor represented by $Z_{S,T}$.

It is also immediate from the definition, or from (\SGARefLink{sga7.sga7i.eviii.n5-7-3}{5.7.3}), that the
formation of $Z_{S,T}$ commutes with every base change
$f:S'\to S$.  This gives an isomorphism, functorial in $Z$,
\begin{equation}\tag{\SGARefTarget{sga7.sga7i.eviii.n5-7-4}{5.7.4}}
 Z_{S,T}\times_S S'\simeq Z_{S',T'},
 \qquad T'=T\times_S S'=T_{S'}.
\end{equation}

\medskip
\noindent \SGARefTarget{sga7.sga7i.eviii.n5-8}{5.8}.  We shall study torsors $E$ under a group scheme of the
form $Z_{S,T}$, where $Z$ is a group, assumed commutative here for
simplicity.  There is a natural functor
\begin{equation}\tag{\SGARefTarget{sga7.sga7i.eviii.n5-8-0}{5.8.0}}
 E\longmapsto E_T
\end{equation}
from the category of torsors under $Z_{S,T}$ to the category of
torsors under $Z_T$.  We claim that this functor is an
\underline{equivalence of categories}.  Its full faithfulness means
that for two torsors $E,F$ under $Z_{S,T}$, the map
\[
 \Hom(E,F)\longrightarrow\Hom(E_T,F_T)
\]
is bijective.  By fpqc descent this reduces immediately to the case in
which $E$ and $F$ are trivial, hence to $E=F=Z_{S,T}$, where the
assertion is precisely the bijectivity of (\SGARefLink{sga7.sga7i.eviii.n5-7-2}{5.7.2}) for $S'=S$.

It remains to prove that (\SGARefLink{sga7.sga7i.eviii.n5-8-0}{5.8.0}) is essentially surjective.  Let $E_T$
be a torsor under $Z_T$.  We already observed in the proof of \SGARefLink{sga7.sga7i.eviii.n5-1}{5.1} that
$E_T$ is representable and locally trivial for the etale topology.  It
therefore suffices to prove that
\[
 H^1(S_{\mathrm{\acute{e}t}},Z_{S,T})
 \longrightarrow H^1(T_{\mathrm{\acute{e}t}},Z_T)
\]
is bijective.  This follows from (\SGARefLink{sga7.sga7i.eviii.n5-7-3}{5.7.3}) and the relation
\[
 R^1i_{\mathrm{\acute{e}t},*}(Z_T)=0,
\]
which is the special case SGA 4 VIII 5.6.

\medskip
\noindent \SGARefTarget{sga7.sga7i.eviii.n5-8-1}{5.8.1}.  Let now $P$ be a commutative group scheme over $S$.
We wish to determine the category of extensions, as sheaves of
commutative groups, of $P$ by $Z_{S,T}$.  For this we use the
description in VII \SGARefLink{sga7.sga7i.evii.n1-1-6}{1.1.6} and \SGARefLink{sga7.sga7i.evii.n1-3}{1.3} in terms of torsors on
$P$, $P\times P$, and $P\times P\times P$, under the inverse images
of $Z_{S,T}$.  By (\SGARefLink{sga7.sga7i.eviii.n5-7-4}{5.7.4}), these inverse images are again of the same
form as $Z_{S,T}$.  Applying the preceding result to those torsors
shows that they correspond to torsors on
$P_T$, $P_T\times_TP_T$, and
$P_T\times_TP_T\times_TP_T$ under the inverse images of $Z_T$.
We therefore obtain the following.

\medskip
\noindent\underline{Proposition \SGARefTarget{sga7.sga7i.eviii.n5-9}{5.9}.}  \underline{Let $S$ be a scheme,
$T$ a closed subscheme, $Z$ a commutative group, $Z_{S,T}$ the
corresponding truncated constant group scheme over $S$ defined in
\SGARefLink{sga7.sga7i.eviii.n5-7}{5.7}, and $P$ a commutative group scheme over $S$.  Then restriction to
$T$ induces an equivalence between the category of sheaves of
commutative groups which are extensions of $P$ by $Z_{S,T}$ and the
category of sheaves of commutative groups which are extensions of
$P_T$ by $Z_T$.}

An essentially identical argument gives the following.

\medskip
\noindent\underline{Corollary \SGARefTarget{sga7.sga7i.eviii.n5-10}{5.10}.}  \underline{With the notation of
\SGARefLink{sga7.sga7i.eviii.n5-9}{5.9}, let $Q$ be another commutative group scheme over $S$.  Then
restriction to $T$ induces an equivalence between the category of
biextensions of $(P,Q)$ by $Z_{S,T}$ and the category of biextensions
of $(P_T,Q_T)$ by $Z_T$.}


% END INCLUDED SGA7I_English_current_buildable_source_20260823/components/115_expose_VIII_sections_5_7_through_5_10.tex

% BEGIN INCLUDED SGA7I_English_current_buildable_source_20260823/components/116_expose_VIII_section_6_1_and_lemma_6_2.tex
% SOURCE AUTHORITY:
% expose_VIII_body.tex lines 1976--2003; scan indices 298--300;
% source folios 287--289.
% Direct scan comparison: physical PDF pages 299--301 in overlapping
% 1100-dpi bands.  The prose, formulas, and subscripts are fully legible at
% that resolution; no higher-resolution escalation was needed in this unit.
% Transparent source dispositions:
% - the stray \mathbb Z_Z in the paragraph preceding Lemma 6.2 is corrected
%   to \mathbb Z_S, the free group scheme over the stated base S;
% - the corrupt insertion between "over S" and \overline G in Lemma 6.2 is
%   omitted; the authority clearly defines alpha:S->\mathbb Z_{S,T} and then
%   pulls back the extension \overline G along alpha.
% Hard stop before 6.2.1, source line 2004 / scan index 300.

\medskip
\noindent\underline{6.  Extensions and biextensions by the Neron model of
$G_m$.}

\medskip
\noindent \SGARefTarget{sga7.sga7i.eviii.n6-1}{6.1}.  Let $S$ be a scheme, $T$ a closed subscheme, and $Z$ a set,
giving a truncated constant scheme $Z_{S,T}$ over $S$ (\SGARefLink{sga7.sga7i.eviii.n5-7}{5.7}).  There is a
canonical homomorphism, functorial in $Z$, of etale schemes over $S$,
\begin{equation}\tag{\SGARefTarget{sga7.sga7i.eviii.n6-1-1}{6.1.1}}
 Z_S\longrightarrow Z_{S,T},
\end{equation}
which is a \underline{surjective} morphism, as follows at once from (\SGARefLink{sga7.sga7i.eviii.n5-7-1}{5.7.1}).
Its functoriality in $Z$ shows that when $Z$ is a group, (\SGARefLink{sga7.sga7i.eviii.n6-1-1}{6.1.1}) is a
homomorphism of etale group schemes over $S$.  Its kernel is therefore an
etale group scheme over $S$, denoted $Z_{S,U}$, and there is a canonical
exact sequence, functorial in $Z$,
\begin{equation}\tag{\SGARefTarget{sga7.sga7i.eviii.n6-1-2}{6.1.2}}
 0\longrightarrow Z_{S,U}\longrightarrow Z_S\longrightarrow Z_{S,T}
 \longrightarrow 0.
\end{equation}
It follows immediately from (\SGARefLink{sga7.sga7i.eviii.n5-7-1}{5.7.1}) that $Z_{S,U}$ is the sum of a family of
$S$-schemes $S_n$, $n\in Z$, where $S_n=S$ for $n=0$ and $S_n=U$ for
$n\ne0$.  This description of $Z_{S,U}$ readily shows that for every group
$F$ over $S$, restriction to $U$ induces a bijection
\begin{equation}\tag{\SGARefTarget{sga7.sga7i.eviii.n6-1-3}{6.1.3}}
 \mathrm{Hom}_{\mathrm{gr}}(Z_{S,U},F)
 \xrightarrow{\ \sim\ }
 \mathrm{Hom}_{\mathrm{gr}}(Z_U,F|U).
\end{equation}

\medskip
\noindent \SGARefTarget{sga7.sga7i.eviii.n6-1-4}{6.1.4}.  Henceforth we suppose that $Z$ is commutative, and we wish
to examine extensions of commutative groups of $Z_{S,T}$ by a given
commutative group $G$ over $S$.  By the exact sequence (\SGARefLink{sga7.sga7i.eviii.n6-1-2}{6.1.2}), the category
of these extensions is equivalent to the category of extensions $F$ of
$Z_S$ by $G$, furnished with a trivialization of the inverse image of the
extension along $Z_{S,U}\to Z_S$.  By (\SGARefLink{sga7.sga7i.eviii.n6-1-3}{6.1.3}), such a trivialization amounts
to a trivialization of the extension $F|U$ of $Z_U$ by $G|U$.

For what follows we take $Z$ to be the group $\mathbb Z$ of integers.  Thus
$\mathbb Z_S$ is the free group on one generator, and VII \SGARefLink{sga7.sga7i.evii.n1-4}{1.4} gives an
equivalence between the category of extensions of $\mathbb Z_S$ by $G$ and
the category of torsors under $G$.  Applying the same result to
$\mathbb Z_U$ and $G|U$ gives the following.

\medskip
\noindent\underline{Lemma \SGARefTarget{sga7.sga7i.eviii.n6-2}{6.2}.}  \underline{Let $S$ be a scheme, $T$ a
closed subscheme, $U=S\setminus T$, and $G$ a commutative group over $S$.
The category of commutative groups over $S$ which are extensions of
$\mathbb Z_{S,T}$ (defined in (\SGARefLink{sga7.sga7i.eviii.n5-7}{5.7})) by $G$ is equivalent to the category of
pairs $(F,s)$ consisting of a torsor $F$ under $G$ and a section $s$ of
$F|U$.  If $\alpha:S\to\mathbb Z_{S,T}$ is the section of
$\mathbb Z_{S,T}$ which is the image of the section $1$ of $\mathbb Z_S$
over $S$, then the torsor corresponding to an extension $\overline G$ of
$\mathbb Z_{S,T}$ by $G$ is $F=\alpha^*(\overline G)$ under $G$, furnished
with the section of $F|U$ given by restricting the unit section of
$\overline G$ to $U$.}


% END INCLUDED SGA7I_English_current_buildable_source_20260823/components/116_expose_VIII_section_6_1_and_lemma_6_2.tex

% BEGIN INCLUDED SGA7I_English_current_buildable_source_20260823/components/117_expose_VIII_sections_6_2_1_through_6_2_4.tex
% SOURCE AUTHORITY:
% expose_VIII_body.tex lines 2004--2029; scan indices 300--302;
% source folios 289--291.
% Direct scan comparison: physical PDF pages 301--303 in overlapping
% 1100-dpi bands.  The mathematical text is fully legible at that resolution.
% Transparent source disposition: the printed "description 3.2" at the end
% of 6.2.1 is corrected to the immediately preceding description in 6.2.
% Hard stop before 6.3, source line 2031 / scan index 302.

\medskip
\noindent \SGARefTarget{sga7.sga7i.eviii.n6-2-1}{6.2.1}.  If $\overline G$ corresponds to the pair $(F,s)$, then the
inverse image $f_n^*(\overline G)$ in $\overline G$ of the section $f_n$ of
$\mathbb Z_{S,T}$ defined by the constant section of $\mathbb Z_S$ with
value $n\in\mathbb Z$ is canonically identified, as a torsor under $G$, with
the $n$-th power $F^{(n)}$ of the $G$-torsor $F$.  The canonical isomorphism
of $G_U$-torsors
\[
 f_n^*(\overline G)|U\simeq G_U,
\]
defined by the first relation in (\SGARefLink{sga7.sga7i.eviii.n5-7-1}{5.7.1}), is determined by the section
$s^{(n)}$ of $F_U^{(n)}$, the $n$-th power of the given section $s$ of
$F_U$.  This makes the following description of the sections of
$\overline G$ over $S$ explicit.  A section $\overline g$ of $\overline G$
over $S$ is determined by:

\smallskip
\noindent a) a section $\underline n$ of $\mathbb Z_{S,T}$ over $S$, which
may be interpreted as a section of $\mathbb Z_T$, that is, as a locally
constant function on $T$ with values in $\mathbb Z$, or equivalently as a
partition
\[
 T=\coprod_{i\in\mathbb Z}T_i
\]
of $T$ into a disjoint sum of open subsets indexed by $i\in\mathbb Z$;

\smallskip
\noindent b) a family of sections
\[
 \begin{gathered}
  \overline g_i\in\Gamma(F^{(i)}/S_i'),\qquad i\in\mathbb Z,\\
  S_i'=S\setminus T_i',\qquad
  T_i'=\displaystyle\bigcup_{\substack{j\in\mathbb Z\\j\ne i}}T_j,
 \end{gathered}
\]
subject to the following compatibility condition.  For $i,j\in\mathbb Z$,
$\overline g_i$ and $\overline g_j$ agree on $S_i'\cap S_j'=U$ after
$F^{(i)}|U$ and $F^{(j)}|U$ have been identified with $G_U$ by means of the
sections $s^{(i)}$ and $s^{(j)}$, respectively.  In other words, there must
exist a necessarily unique section $h_{ij}\in\Gamma(G/U)$ such that
\[
 \overline g_i=h_{ij}s^{(i)},\qquad
 \overline g_j=h_{ij}s^{(j)}.
\]
Since the description in \SGARefLink{sga7.sga7i.eviii.n6-2}{6.2} of an extension $\overline G$ of
$\mathbb Z_{S,T}$ by $G$ is plainly compatible with every base change
$S'\to S$, the preceding gives a complete description of the sheaf
$\overline G$ as a functor of the argument
$S'\in\operatorname{Ob}(\mathrm{Sch}_{/S})$.

\medskip
\noindent \SGARefTarget{sga7.sga7i.eviii.n6-2-2}{6.2.2}.  If $T$ is a scheme which is a sum of nonempty connected
schemes $T_i$, $i\in I$, the description in \SGARefLink{sga7.sga7i.eviii.n6-2-1}{6.2.1} of the sections of
$\overline G$ over $S$ may be simplified.  Such a section is determined by
a family of integers $(n_i)_{i\in I}$ indexed by $I$, and a family of
sections
\[
 \overline g_i\in\Gamma(F^{(n_i)}/S_i'),\qquad i\in I,
\]
where
\[
 S_i'=S\setminus T_i',
 \qquad
 T_i'=\coprod_{\substack{j\in I\\j\ne i}}T_j.
\]
The $\overline g_i$ satisfy the following compatibility condition: for every
pair $i,j\in I$, the restrictions of $\overline g_i$ and $\overline g_j$ to
$S_i'\cap S_j'=U$ agree under the identifications of $F^{(n_i)}|U$ and
$F^{(n_j)}|U$ deduced from the respective sections $s^{(n_i)}$ and
$s^{(n_j)}$ of these torsors.

\medskip
\noindent \SGARefTarget{sga7.sga7i.eviii.n6-2-3}{6.2.3}.  Finally, if $T$ is assumed nonempty and connected, giving a
section $\overline g$ of $\overline G$ over $S$ amounts to giving an integer
$n$ and a section $\overline g$ of $F^{(n)}$ over $S$.  There is no longer
any compatibility condition to impose, so the section $s$ of $F|U$ does not
occur in the description of the group of sections of $\overline G$ over
$S$.  (Of course, $s$ occurs again as soon as one seeks to describe the
points of $\overline G$ with values in a general $S$-scheme $S'$.)

\medskip
\noindent \SGARefTarget{sga7.sga7i.eviii.n6-2-4}{6.2.4}.  Since $\overline G$ is the union of the ``open
subschemes'' $f_n^*(\overline G)$, $n\in\mathbb Z$, described in \SGARefLink{sga7.sga7i.eviii.n6-2-1}{6.2.1}, it
follows that $\overline G$ is representable if and only if the torsors
$F^{(n)}$ are representable.  For example, it suffices that $G$ be
representable and affine over $S$, by the descent theory of SGA 1 VIII 2.1.

% END INCLUDED SGA7I_English_current_buildable_source_20260823/components/117_expose_VIII_sections_6_2_1_through_6_2_4.tex

% BEGIN INCLUDED SGA7I_English_current_buildable_source_20260823/components/118_expose_VIII_sections_6_3_and_6_4.tex
% SOURCE AUTHORITY:
% expose_VIII_body.tex lines 2031--2063; scan indices 302--304;
% source folios 291--293.
% Direct scan comparison: physical PDF pages 303--305 in overlapping
% 1100-dpi bands.  The mathematical text and displayed formulas are fully
% legible at that resolution.
% Transparent source disposition: formula (6.3.2) prints s_D next to
% O_S(-D), while the immediately following defining sentence correctly
% identifies the canonical rational section as s_{-D}; the English formula
% uses s_{-D}.
% Hard stop before Theorem 6.5, source line 2064 / scan index 304.

\medskip
\noindent \SGARefTarget{sga7.sga7i.eviii.n6-3}{6.3}.  In what follows we suppose that
$G=\underline G_{mS}$.  Every extension $\overline G$ of
$\mathbb Z_{S,T}$ by $G=\underline G_{mS}$ is then representable (\SGARefLink{sga7.sga7i.eviii.n6-2-4}{6.2.4}),
and is of course smooth over $S$, being an extension of the etale group
scheme $\mathbb Z_{S,T}$ by the smooth group scheme $\underline G_{mS}$.
Moreover, giving a torsor $F$ under $\underline G_m$ is equivalent to giving
an invertible Module $\mathbb L$ over $S$ (by associating to the latter the
torsor of invertible sections of $\mathbb L$).  Thus $\overline G$ is
described by a pair $(\mathbb L,s)$, where $\mathbb L$ is an invertible
Module over $S$ and $s$ is a trivialization of $\mathbb L$ over
$U=S\setminus T$, that is, an invertible section of $\mathbb L|U$.

When $U$ is schematically dense in $S$, and $S$ is locally noetherian, or is
reduced with a locally finite set of irreducible components, giving such a
pair is equivalent to giving a \underline{divisor $D$ on $S$ such that}
\begin{equation}\tag{\SGARefTarget{sga7.sga7i.eviii.n6-3-1}{6.3.1}}
 \operatorname{supp}D\subset T,
\end{equation}
by associating to such a divisor the pair
\begin{equation}\tag{\SGARefTarget{sga7.sga7i.eviii.n6-3-2}{6.3.2}}
 \bigl(\underline O_S(-D),s_{-D}|U\bigr)
\end{equation}
formed by the invertible Module $\underline O_S(-D)$ associated with $-D$
(EGA IV 21.2.8.1) and the restriction to $U$ of its canonical rational
section $s_{-D}$ (EGA IV 21.2.11 and 20.2.11(ii)).  In any event, without
any hypothesis on $S$, a divisor $D$ on $S$ satisfying (\SGARefLink{sga7.sga7i.eviii.n6-3-1}{6.3.1}) defines by
(\SGARefLink{sga7.sga7i.eviii.n6-3-2}{6.3.2}) a pair $(\mathbb L,s)$ and hence an extension $\overline G$ of
$\mathbb Z_{S,T}$ by $\underline G_m$.  We call it the
\underline{Neron--Raynaud model of $\underline G_m$} associated with the
divisor $D$ and the closed subscheme $T$ of $S$ satisfying (\SGARefLink{sga7.sga7i.eviii.n6-3-1}{6.3.1}).  If
$T=\operatorname{supp}D$, we omit $T$ from the terminology.

It is immediate that the formation of this extension $\overline G$ commutes
with every base change $f:S'\to S$ for which the inverse-image divisor
$f^*(D)$ is defined (EGA IV 21.4.2), and in particular with every
\underline{flat} base change.

\medskip
\noindent \SGARefTarget{sga7.sga7i.eviii.n6-3-3}{6.3.3}.  The general description in \SGARefLink{sga7.sga7i.eviii.n6-2-1}{6.2.1} of the sections of
$\overline G$ immediately shows that they may be identified with the
sections $f$ of the sheaf $\mathcal M_S^*$ of regular meromorphic sections
on $S$ (EGA IV 20.1.8) whose divisor $\operatorname{div}(f)$ is locally, for
the Zariski topology, an integral multiple of $D$.  If $V$ is an open subset
of $S$ on which the restriction of $\operatorname{div}(f)$ equals that of
$nD$, where $n\in\mathbb Z$, then the section of $\mathbb Z_{S,T}$ defined
by $f$ is simply the image in (\SGARefLink{sga7.sga7i.eviii.n6-1-2}{6.1.2}) of the constant section of
$\mathbb Z_S$ with value $n$.  The preceding description of the group of
sections of $\overline G$ over $S$ remains valid over every open subset of
$S$, and more generally for the group of points of $\overline G$ with
values in any flat $S$-scheme $S'$.  Since $\overline G$ itself is flat
(indeed smooth) over $S$, this description again gives a complete
characterization of the Neron--Raynaud model of $\underline G_m$ defined by
$D$ and $T\supset\operatorname{supp}D$.

\medskip
\noindent \SGARefTarget{sga7.sga7i.eviii.n6-4}{6.4}.  Suppose now that $S$ is \underline{normal}, that its set of
irreducible components is locally finite (so that $S$ is a sum of normal
integral schemes), and that $D$ is a positive divisor with all
multiplicities equal to $1$ (EGA IV 21.6), that $T=\operatorname{supp}D$,
and that $T$ is locally irreducible, that is, its irreducible components are
disjoint.  Then for every rational function $f$ on $S$ which is regular on
$U=S\setminus T$, the divisor $\operatorname{div}(f)$ is locally a multiple
of $D$.  Thus in this case there is a canonical isomorphism of groups
\begin{equation}\tag{\SGARefTarget{sga7.sga7i.eviii.n6-4-1}{6.4.1}}
 \Gamma(S,\overline G)\xrightarrow{\ \sim\ }
 \Gamma(U,\underline G_{mU}).
\end{equation}
Suppose further that $T$ is geometrically unibranch and locally integral, a
hypothesis stable under smooth base change (EGA IV 17.5.7).  Since all the
hypotheses are then stable under base change, (\SGARefLink{sga7.sga7i.eviii.n6-4-1}{6.4.1}) implies that
\underline{for every $S'$ smooth over $S$ there is a canonical group
isomorphism, functorial in $S'$,}
\begin{equation}\tag{\SGARefTarget{sga7.sga7i.eviii.n6-4-2}{6.4.2}}
 \mathrm{Hom}_S(S',\overline G)\xrightarrow{\ \sim\ }
 \mathrm{Hom}_U(S'_U,\underline G_{mU})
 \simeq\Gamma(S'_U,\underline O_{S'_U}^*).
\end{equation}
Since $\overline G$ is itself smooth over $S$, this isomorphism characterizes
the group scheme $\overline G$ over $S$ up to unique isomorphism.  It also
shows that $\overline G$ plays the role of a ``Neron model'' for the group
scheme $\underline G_{mU}$ over the open subset $U$ of $S$.  Together with
the fact that it was first introduced and used by M.~Raynaud, this justifies
the terminology adopted in \SGARefLink{sga7.sga7i.eviii.n6-3}{6.3} for $\overline G$.


% END INCLUDED SGA7I_English_current_buildable_source_20260823/components/118_expose_VIII_sections_6_3_and_6_4.tex

% BEGIN INCLUDED SGA7I_English_current_buildable_source_20260823/components/119_expose_VIII_theorem_6_5_through_corollary_6_7.tex
% SOURCE AUTHORITY:
% expose_VIII_body.tex lines 2064--2102; scan indices 304--307;
% source folios 293--296.
% Direct scan comparison: physical PDF pages 305--308 in overlapping
% 1100-dpi bands.  The theorem statements, formulas, and proof are fully
% legible at that resolution.
% Transparent source disposition: the proof of essential surjectivity in
% Theorem 6.5 prints "local rings of X" once; the stated base is S, and the
% English corrects X to S.
% Hard stop before Remarks 6.8, source line 2103 / scan index 307.

\medskip
\noindent\underline{Theorem \SGARefTarget{sga7.sga7i.eviii.n6-5}{6.5}.}  \underline{Let $S$ be a normal locally
integral scheme and let $D$ be a positive divisor on $S$, without multiple
components, whose support is geometrically unibranch and locally integral.
Let $\overline G$ be the Neron--Raynaud model of $\underline G_m$ over $S$
relative to $D$ (\SGARefLink{sga7.sga7i.eviii.n6-3}{6.3}), and consider the category of torsors $E$ over $S$
under the group scheme $\overline G$.  The restriction functor}
\[
 E\longmapsto E_U=E|U
\]
\underline{from this category to the category of torsors over $U$ under the
group scheme $\underline G_{mU}$ is fully faithful.  If $S$ is regular (more
generally, if its local rings are factorial), this functor is an equivalence
of categories.}

Let us prove that the functor in question is fully faithful, that is, that
for two torsors $E,F$ under $\overline G$, the restriction map
\begin{equation}\tag{*}
 \mathrm{Hom}(E,F)\longrightarrow\mathrm{Hom}(E_U,F_U)
\end{equation}
is bijective, where $\mathrm{Hom}$ denotes the set of homomorphisms of
\underline{torsors}.  First observe that this question is plainly
\underline{local on $S$ for the Zariski topology}.  Indeed, suppose that
$S$ is the union of open subsets $S_i$ such that for every open subset $S'$
of $S$ contained in some $S_i$ (and hence in particular for $S'=S_i$ or
$S_i\cap S_j$), the map
\[
 \mathrm{Hom}(E_{S'},F_{S'})
 \longrightarrow
 \mathrm{Hom}(E_{S'\cap U},F_{S'\cap U})
\]
is bijective.  Then (*) is bijective.  On the other hand, it follows at once
from \SGARefLink{sga7.sga7i.eviii.n5-8}{5.8} applied to $S,T,\mathbb Z$, and from \SGARefLink{sga7.sga7i.eviii.n5-1}{5.1} applied to $T,\mathbb Z$,
that every torsor under $\overline G$ is locally trivial for the Zariski
topology.  These two observations reduce us to proving the bijectivity of
(*) when $E=F=\overline G$.  In that case the map is precisely (\SGARefLink{sga7.sga7i.eviii.n6-4-1}{6.4.1}),
whose bijectivity has already been established.

We now prove that the functor $E\mapsto E_U$ is essentially surjective when
the local rings of $S$ are factorial.  Since $\underline G_{mS}$ is
identified with an open subgroup scheme of $\overline G$, it is enough to
show that every torsor under $\underline G_{mU}$ is isomorphic to the
restriction to $U$ of a torsor under $\underline G_{mS}$.  Equivalently,
every invertible Module over $U$ must be isomorphic to the restriction of an
invertible Module over $S$.  This follows from the hypothesis on $S$
(EGA IV 21.6.11, where only the locally noetherian case is treated; that is
the only case needed below).  This proves \SGARefLink{sga7.sga7i.eviii.n6-5}{6.5}.

\medskip
\noindent \SGARefTarget{sga7.sga7i.eviii.n6-5-1}{6.5.1}.  Suppose now, in addition, that $S$ is locally noetherian
and \underline{regular}.  Since $S$ is then locally factorial, giving $D$
is equivalent to giving a reduced, purely $1$-codimensional closed
subscheme, again denoted $D$, which is geometrically unibranch.  Notice also
that the hypotheses on $S$ and $D$ are stable under every smooth base change
$S'\to S$, so that the conclusion of \SGARefLink{sga7.sga7i.eviii.n6-5}{6.5} applies to $S'$ with the subscheme
$D'=D\times_S S'$.  The description in VII \SGARefLink{sga7.sga7i.evii.n1-1-6}{1.1.6} of extension groups now
immediately gives the following consequence.

\medskip
\noindent\underline{Corollary \SGARefTarget{sga7.sga7i.eviii.n6-6}{6.6}.}  \underline{Let $S$ be a locally
noetherian regular scheme, let $D$ be a reduced, geometrically unibranch,
purely $1$-codimensional closed subscheme (thus defining a divisor on $S$),
let $\overline G$ be the Neron--Raynaud model of $\underline G_m$ relative
to $D$, and let $P$ be a smooth commutative group scheme over $S$.  Consider
the category of commutative extensions, as sheaves of groups, of $P$ by
$\overline G$.  Restriction to $U$, $E\mapsto E|U$, induces an equivalence
between this category and the category of extensions of $P_U$ by
$\underline G_{mU}$.}

\medskip
\noindent \SGARefTarget{sga7.sga7i.eviii.n6-6-1}{6.6.1}.  The same method gives the analogous statement for
extensions which are not necessarily commutative (with $P$ no longer
assumed commutative).  We leave the details of the proof to the reader; they
are slightly more involved because bitorsors must be considered in place of
torsors.

By essentially the same proof as for \SGARefLink{sga7.sga7i.eviii.n6-6}{6.6}, one likewise obtains:

\medskip
\noindent\underline{Corollary \SGARefTarget{sga7.sga7i.eviii.n6-7}{6.7}.}  \underline{Let $S,D,\overline G$ be as
in \SGARefLink{sga7.sga7i.eviii.n6-6}{6.6}, and let $P,Q$ be two smooth commutative group schemes over $S$.
Consider the category of biextensions of $(P,Q)$ by $\overline G$.  Then the
restriction functor to $U$, $E\mapsto E|U$, from this category to the
category of biextensions of $(P_U,Q_U)$ by $\underline G_{mU}$ is an
equivalence of categories.}


% END INCLUDED SGA7I_English_current_buildable_source_20260823/components/119_expose_VIII_theorem_6_5_through_corollary_6_7.tex

% BEGIN INCLUDED SGA7I_English_current_buildable_source_20260823/components/120_expose_VIII_remarks_6_8_and_autocritique_6_9.tex
% SOURCE AUTHORITY:
% expose_VIII_body.tex lines 2103--2170; scan indices 307--310;
% source folios 296--299.
% Direct scan comparison: physical PDF pages 308--311 in overlapping
% 1100-dpi bands.  The prose and formulas are fully legible at that
% resolution.  The hand-corrected diagram on physical page 310 was also
% inspected directly in a 5000-dpi targeted crop.
% Transparent source dispositions:
% - 6.9(a) begins with j:V->S but consistently defines U=S-Y; V is corrected
%   to U, and j_*G_{mS} is corrected to j_*G_{mU};
% - the source diagram retains a faint obsolete right arrowhead on the Y--S
%   stroke beneath a heavy hand-added left arrowhead.  The correction gives
%   Y->S.  The mark at U is an open-immersion hook, and the source's terminal
%   period is retained.
% Hard stop before section 7, source line 2171 / scan index 310.

\medskip
\noindent\underline{Remarks \SGARefTarget{sga7.sga7i.eviii.n6-8}{6.8}.}\quad a) We shall apply \SGARefLink{sga7.sga7i.eviii.n6-6}{6.6} and \SGARefLink{sga7.sga7i.eviii.n6-7}{6.7} when
$S$ is a trait and $D$ is the divisor defined by the inclusion of the closed
point $s$ of $S$ into $S$.  In this case $U$ is the spectrum of the fraction
field $K$ of $S$, and $\overline G$ is the usual Neron model of the
multiplicative group over $K$, introduced by Raynaud.  Notice that even in
this special case, the proof of \SGARefLink{sga7.sga7i.eviii.n6-6}{6.6} and \SGARefLink{sga7.sga7i.eviii.n6-7}{6.7} uses \SGARefLink{sga7.sga7i.eviii.n6-5}{6.5} over bases more general
than traits (such as $P$ and its Cartesian powers).  Nevertheless, in this
case it is enough to know the \underline{existence} of the Neron model
$\overline G$ in the special situation under consideration: \SGARefLink{sga7.sga7i.eviii.n6-5}{6.5} for base
schemes smooth over $S$, and hence \SGARefLink{sga7.sga7i.eviii.n6-6}{6.6} and \SGARefLink{sga7.sga7i.eviii.n6-7}{6.7}, then follow formally.  For
us, the introduction of $\overline G$ will chiefly serve as a convenient
intermediate step in studying, in the next section, the problem of extending
extensions and biextensions by the multiplicative group.

\smallskip
\noindent b) A Neron--Raynaud model $\overline G$ of
$\underline G_{mU}$, satisfying the fundamental condition (\SGARefLink{sga7.sga7i.eviii.n6-4-2}{6.4.2}) for every
$S'$ smooth over $S$, can be defined under much more general conditions than
those considered above, with $U$ merely a schematically dense open subset
of a normal locally noetherian scheme $S$.  In this case $\overline G$ is an
fppf \underline{sheaf}, not necessarily representable.  Let $\underline D$
denote the sheaf, on the small etale site of $S$, of divisors on $S$ whose
support is contained in $T=S\setminus U$, and, by abuse of language, also
denote by the same letter its canonical extension to a sheaf on
$(\mathrm{Sch})_{/S}$.  Its sections over a relative scheme
$f:S'\to S$ are the sections of $f_{\mathrm{\acute{e}t}}^*(\underline D)$
over $S'$.  The desired model is an extension of groups
\begin{equation}\tag{\SGARefTarget{sga7.sga7i.eviii.n6-8-1}{6.8.1}}
 0\longrightarrow\underline G_{mS}\longrightarrow\overline G
 \longrightarrow\underline D\longrightarrow0.
\end{equation}
The construction of a canonical extension (\SGARefLink{sga7.sga7i.eviii.n6-8-1}{6.8.1}), without any normality
hypothesis on $S$, presents no difficulty when one uses the general
description VII \SGARefLink{sga7.sga7i.evii.n1-1-6}{1.1.6}; at the same time one obtains (\SGARefLink{sga7.sga7i.eviii.n6-4-2}{6.4.2}) for $S'$ etale
over $S$.  For the same relation to hold for every $S'$ smooth over $S$, it
is equivalent to verify that the formation of the sheaf $\underline D$
commutes with every smooth base change $f:S'\to S$.  For normal $S$, this
follows from EGA $\mathrm{Err}_{\mathrm{IV}}$ 53, 21.4.9.

\medskip
\noindent \SGARefTarget{sga7.sga7i.eviii.n6-9}{6.9}.\quad ``\underline{Self-criticism}'' (observations of
P.~Deligne).\footnote{The author of the present Expose having admitted
defeat, the reader is invited, if necessary, to rewrite sections 5 and 6
according to the accompanying directions.}

Sections 5 and 6 are really rather uninspiring.  One periodically has the
impression that the formalism of $i_*,i^*$ is being redone ``by hand'',
without seeing what is geometric and what is ``general nonsense''.

I propose extracting the ``geometric'' statements a), b), c):

\smallskip
\noindent a) (cf. \SGARefLink{sga7.sga7i.eviii.n6-8}{6.8} b)) Let $j:U\hookrightarrow S$ be a schematically
dense open subset of a normal noetherian scheme $S$.  The etale sheaf
$\underline G_{mS}$ is then identified with a subsheaf of the etale sheaf
$j_*\underline G_{mU}$.  Let $\underline D$ be the quotient etale sheaf in
the exact sequence
\[
 0\longrightarrow\underline G_{mS}\longrightarrow
 j_*\underline G_{mU}\longrightarrow\underline D\longrightarrow0.
\]
\underline{Then} the formation of $\underline D$ commutes with every smooth
base change (EGA $\mathrm{Err}_{\mathrm{IV}}$ 53, 21.4.9); \underline{moreover},
$R^1j_*\underline G_m=0$ (Theorem 90) if the strictly local rings of $S$ are
factorial.

\smallskip
\noindent b) If $U=S\setminus Y$, where $Y$ is a geometrically unibranch
divisor, and if $i$ is the inclusion of $Y$ into $S$, then
\[
 \underline D=i_*\mathbb Z_Y.
\]

\smallskip
% SGA-SRC-001273
\noindent c) If $Y$ is a closed subscheme of a scheme $X$, with closed immersion
$i:Y\hookrightarrow X$, and if $E$ is a set,
then the etale sheaf $i_*\underline E_Y$ is represented by an etale scheme
over $X$, whose formation commutes with every base change.

\medskip
\noindent\underline{Corollary.}  \underline{Let $S$ be a normal noetherian
scheme, let $Y$ be a geometrically unibranch divisor, and let
$U=S\setminus Y$:}
\begin{center}
\begin{tikzcd}[column sep=huge]
 U \arrow[r, hook, "j"'] & S & Y \arrow[l, "i"']
\end{tikzcd}
\qquad .
\end{center}
\underline{There exists a unique smooth group scheme $\mathcal G_m$ over
$S$ such that after every smooth base change one has
$\mathcal G_m|S_{\mathrm{\acute{e}t}}=j_*(\underline G_m)_U$.}

Indeed, the preceding requirements determine $\mathcal G_m$ as a sheaf, for
the etale topology, on the category of schemes smooth over $S$.  By a) and
b), the sequence
\[
 0\longrightarrow\underline G_m\longrightarrow\mathcal G_m
 \xrightarrow{\ q\ }i_*\mathbb Z_Y\longrightarrow0
\]
is exact.  By c), $i_*\mathbb Z_Y$ is representable.  Finally, the morphism
$q$ is representable.

Use of the preceding ``small smooth site'' should make it possible to deduce
the statements of section 6 formally from a).

% END INCLUDED SGA7I_English_current_buildable_source_20260823/components/120_expose_VIII_remarks_6_8_and_autocritique_6_9.tex

% BEGIN INCLUDED SGA7I_English_current_buildable_source_20260823/components/121_expose_VIII_section_7_opening_and_theorem_7_1_statement.tex
% SOURCE AUTHORITY:
% expose_VIII_body.tex lines 2171--2203; scan indices 310--311;
% source folios 299--300.
% Direct scan comparison: physical PDF pages 311--312 in overlapping
% 1100-dpi bands.  The prose, component-group notation, and obstruction
% formulas are fully legible; no higher-resolution escalation was needed.
% Transparent source dispositions:
% - typewritten letter-o special-fibre subscripts are normalized to zero;
% - in 7.1(b), the source twice drops the eta from E_eta in the extension
%   criterion; the English retains the stated generic-fibre object E_eta;
% - the coefficient is typeset coherently as \underline G_{mS}.
% Hard stop before the proof, source line 2204 / scan index 311.

\medskip
\noindent\underline{7.  Canonical extensions of extensions and
biextensions by $\underline G_m$.}

\medskip
\noindent \SGARefTarget{sga7.sga7i.eviii.n7-0}{7.0}.  We now reap the fruit of our preparations in sections 5 and
6.  First let the base scheme $S$ be a \underline{trait}, that is, the
spectrum of a discrete valuation ring; let $\eta$ be its generic point,
$s$ its closed point, and $k=k(s)$.  Let $P$ be a smooth commutative group
scheme over $S$, and write $P_0=P_s$ for its special fibre.  In the present
section we consider only commutative groups; in particular, all group
extensions under consideration are commutative.

\medskip
\noindent\underline{Theorem \SGARefTarget{sga7.sga7i.eviii.n7-1}{7.1}.}  \underline{The notation is that recalled
in \SGARefLink{sga7.sga7i.eviii.n7-0}{7.0}.}

\smallskip
\noindent a) \underline{Suppose that the group of connected components of
the special fibre of $P$,}
\[
 \Phi=P_0/P_0^{\circ},
\]
\underline{is a torsion group---as is the case, in particular, if $P$ is of
finite type over $S$.  Then the functor $E\mapsto E_\eta$ from the category
of extensions of $P$ by $\underline G_{mS}$ to the category of extensions
of $P_\eta$ by $\underline G_{m\eta}$ is fully faithful.  Suppose that there
is an integer $n>0$ such that $n\Phi=0$.  Then every extension $E_\eta$ of
$P_\eta$ by $\underline G_{m\eta}$ has an associated canonical
homomorphism}
\begin{equation}\tag{\SGARefTarget{sga7.sga7i.eviii.n7-1-1}{7.1.1}}
 d(E_\eta):\Phi\longrightarrow(\mathbb Q/\mathbb Z)_k,
\end{equation}
\underline{whose vanishing is necessary and sufficient for $E_\eta$ to be
isomorphic to the generic fibre of an extension of $P$ by
$\underline G_{mS}$.  In particular, if $P_0$ is connected, the preceding
functor $E\mapsto E_\eta$ is an equivalence of categories.}

\smallskip
\noindent b) \underline{Let $P$ and $Q$ be two smooth group schemes over
$S$, put}
\[
 \Phi=P_0/P_0^{\circ},\qquad \Psi=Q_0/Q_0^{\circ},
\]
\underline{and suppose that either $\Phi$ or $\Psi$ is a torsion group (for
example, that either $P$ or $Q$ is of finite type over $S$).  Then the
functor $E\mapsto E_\eta$ from the category of biextensions of $(P,Q)$ by
$\underline G_{mS}$ to the category of biextensions of
$(P_\eta,Q_\eta)$ by $\underline G_{m\eta}$ is fully faithful.  Suppose that
there is an integer $n>0$ such that $n\Phi=0$ or $n\Psi=0$.  Then every
biextension $E_\eta$ of $(P_\eta,Q_\eta)$ by $\underline G_{m\eta}$ has an
associated canonical pairing}
\begin{equation}\tag{\SGARefTarget{sga7.sga7i.eviii.n7-1-2}{7.1.2}}
 d(E_\eta):\Phi\otimes_{\mathbb Z}\Psi
 \longrightarrow(\mathbb Q/\mathbb Z)_k,
\end{equation}
\underline{whose vanishing is necessary and sufficient for $E_\eta$ to be
isomorphic to the generic fibre of a biextension $E$ of $(P,Q)$ by
$\underline G_{mS}$.  In particular, if either $P_0$ or $Q_0$ is connected,
the preceding functor $E\mapsto E_\eta$ is an equivalence of categories.}


% END INCLUDED SGA7I_English_current_buildable_source_20260823/components/121_expose_VIII_section_7_opening_and_theorem_7_1_statement.tex

% BEGIN INCLUDED SGA7I_English_current_buildable_source_20260823/components/122_expose_VIII_theorem_7_1_proof.tex
% SOURCE AUTHORITY:
% expose_VIII_body.tex lines 2204--2242; scan indices 311--313;
% source folios 300--302.
% Direct scan comparison: physical PDF pages 312--314 in overlapping
% 1100-dpi bands.  Every displayed Ext/Biext formula and tensor product is
% fully legible at that resolution.
% Transparent source dispositions:
% - part (a) prints "biextensions by \overline G" where Theorem 6.6 and the
%   entire argument concern extensions; the English uses extensions;
% - the final target printed Q/Z is the same (Q/Z)_k coefficient used in
%   (7.1.2), 5.6, and the preceding Ext calculation, so the English restores
%   (\mathbb Q/\mathbb Z)_k.
% Hard stop before Remark 7.2, source line 2243 / scan index 313.

\medskip
\noindent\underline{Proof.}\quad a) Let $\overline G$ be the
Neron--Raynaud model of $\underline G_m$.  The functor under consideration
is the composite of the functor ``extension of the kernel
$\underline G_{mS}\to\overline G$'' with the functor ``restriction to
$\eta$'' defined on the category of extensions by $\overline G$.  The latter
functor is an equivalence of categories by \SGARefLink{sga7.sga7i.eviii.n6-6}{6.6}, so it remains to study the
first.  For this purpose we use the exact sequence of group schemes
\begin{equation}\tag{\SGARefTarget{sga7.sga7i.eviii.n7-1-3}{7.1.3}}
 0\longrightarrow\underline G_{mS}\longrightarrow\overline G
 \longrightarrow\mathbb Z_{S,s}\longrightarrow0,
\end{equation}
where the notation $\mathbb Z_{S,s}$ is that of \SGARefLink{sga7.sga7i.eviii.n5-7}{5.7}.  The exact sequence of
$\mathrm{Ext}^i$, $i=0,1$, associated with (\SGARefLink{sga7.sga7i.eviii.n7-1-3}{7.1.3}) shows that the functor in
question is fully faithful if and only if
\begin{equation}\tag{\SGARefTarget{sga7.sga7i.eviii.n7-1-4}{7.1.4}}
 \mathrm{Hom}(P,\mathbb Z_{S,s})=0,
\end{equation}
and that the obstruction for an extension $\overline E$ of $P$ by
$\overline G$ to come from an extension $E$ of $P$ by
$\underline G_{mS}$ lies in
\begin{equation}\tag{\SGARefTarget{sga7.sga7i.eviii.n7-1-5}{7.1.5}}
 \mathrm{Ext}^1(P,\mathbb Z_{S,s}).
\end{equation}
By (\SGARefLink{sga7.sga7i.eviii.n5-7-2}{5.7.2}), the first group in (\SGARefLink{sga7.sga7i.eviii.n7-1-4}{7.1.4}) is isomorphic to
$\mathrm{Hom}(P_0,\mathbb Z_k)$, which in turn is isomorphic to
$\mathrm{Hom}(\Phi,\mathbb Z_k)$.  Thus (\SGARefLink{sga7.sga7i.eviii.n7-1-4}{7.1.4}) is equivalent to
\begin{equation}\tag{7.1.4 bis}
 \mathrm{Hom}(\Phi,\mathbb Z_k)=0.
\end{equation}
This relation certainly holds if $\Phi$ is a torsion group.  Moreover, if
$n\Phi=0$, then \SGARefLink{sga7.sga7i.eviii.n5-9}{5.9} and \SGARefLink{sga7.sga7i.eviii.n5-5}{5.5} give canonical isomorphisms
\begin{equation}\tag{\SGARefTarget{sga7.sga7i.eviii.n7-1-6}{7.1.6}}
 \mathrm{Ext}^1(P,\mathbb Z_{S,s})
 \simeq\mathrm{Ext}^1(P_0,\mathbb Z_k)
 \simeq\mathrm{Hom}\bigl(\Phi,(\mathbb Q/\mathbb Z)_k\bigr),
\end{equation}
which completes the proof of \SGARefLink{sga7.sga7i.eviii.n7-1}{7.1}(a).

\smallskip
\noindent b) The proof is essentially the same.  By \SGARefLink{sga7.sga7i.eviii.n6-7}{6.7}, we reduce to
studying the functor ``extension of the kernel
$\underline G_{mS}\to\overline G$'' from the category of biextensions of
$(P,Q)$ by $\underline G_{mS}$ to the category of biextensions of $(P,Q)$ by
$\overline G$.  We use the exact sequence (\SGARefLink{sga7.sga7i.eviii.n7-1-3}{7.1.3}) and the associated exact
sequence of $\mathrm{Biext}^i$, $i=0,1$ (VII \SGARefLink{sga7.sga7i.evii.n3-7-5}{3.7.5}).  Full faithfulness of
the functor is equivalent to
\begin{equation}\tag{\SGARefTarget{sga7.sga7i.eviii.n7-1-7}{7.1.7}}
 \mathrm{Hom}(P\otimes Q,\mathbb Z_{S,s})=0.
\end{equation}
By (\SGARefLink{sga7.sga7i.eviii.n5-7-2}{5.7.2}), the left side is isomorphic to
$\mathrm{Hom}(P_0\otimes Q_0,\mathbb Z_k)$ and hence to
$\mathrm{Hom}(\Phi\otimes\Psi,\mathbb Z_k)$; this group is indeed zero if
either $\Phi$ or $\Psi$ is a torsion group.  On the other hand, for a given
biextension $\overline E$ of $(P,Q)$ by $\overline G$, the obstruction to
its coming from a biextension $E$ of $(P,Q)$ by $\underline G_{mS}$ lies in
\[
 \mathrm{Biext}^1(P,Q;\mathbb Z_{S,s}).
\]
By \SGARefLink{sga7.sga7i.eviii.n5-10}{5.10} this is isomorphic to
$\mathrm{Biext}^1(P_0,Q_0;\mathbb Z_k)$, which, when $n\Phi=0$ or
$n\Psi=0$, is itself isomorphic by \SGARefLink{sga7.sga7i.eviii.n5-6}{5.6} to
\[
 \mathrm{Hom}\bigl(\Phi\otimes\Psi,(\mathbb Q/\mathbb Z)_k\bigr).
\]
This completes the proof of \SGARefLink{sga7.sga7i.eviii.n7-1}{7.1}.


% END INCLUDED SGA7I_English_current_buildable_source_20260823/components/122_expose_VIII_theorem_7_1_proof.tex

% BEGIN INCLUDED SGA7I_English_current_buildable_source_20260823/components/123_expose_VIII_remark_7_2.tex
% SOURCE AUTHORITY:
% expose_VIII_body.tex lines 2243--2261; scan indices 313--314;
% source folios 302--303.
% Direct scan comparison: physical PDF pages 314--315 in overlapping
% 1100-dpi bands.  The obstruction map, codimension bound, and footnote are
% fully legible; no higher-resolution escalation was needed.
% Transparent source disposition: spacing/typographical slips in "for every"
% and the products P x P x P are normalized without changing content.
% Hard stop before 7.3, source line 2262 / scan index 314.

\medskip
\noindent\underline{Remark \SGARefTarget{sga7.sga7i.eviii.n7-2}{7.2}.}  Given a smooth group scheme $P$ over a
connected regular locally noetherian scheme $S$ with generic point $\eta$,
one may ask to study the functor $E\mapsto E_\eta$ from the category of
extensions of $P$ by $\underline G_{mS}$ to the category of extensions of
$P_\eta$ by $\underline G_{m\eta}$.  This problem is readily reduced to the
one treated in \SGARefLink{sga7.sga7i.eviii.n7-1}{7.1}(a).  For simplicity, suppose that $P$ is of finite type
over $S$.  It follows easily from \SGARefLink{sga7.sga7i.eviii.n7-1}{7.1} that this functor is fully
faithful.\footnote{For this it is enough that $S$ be normal rather than
regular.}  One readily deduces that if $E_\eta$ is an extension of $P_\eta$
by $\underline G_{m\eta}$, then it is isomorphic to the generic fibre of an
extension $E$ of $P$ by $\underline G_{mS}$ if and only if, for every
$s\in S$ such that
\[
 \dim\operatorname{Spec}(\underline O_{S,s})=1,
\]
the analogous extension problem obtained by base change
$\operatorname{Spec}(\underline O_{S,s})\to S$ has a solution.  By \SGARefLink{sga7.sga7i.eviii.n7-1}{7.1}(a),
this is expressed by the vanishing of a certain homomorphism
\begin{equation}\tag{\SGARefTarget{sga7.sga7i.eviii.n7-2-1}{7.2.1}}
 d(E_\eta,s):P_s/P_s^{\circ}
 \longrightarrow(\mathbb Q/\mathbb Z)_{k(s)}.
\end{equation}

The preceding assertion plainly reduces to the following one.  Let $U$ be
an open subset of $S$ whose complement $T$ satisfies
\[
 \operatorname{codim}(T,S)\geq2.
\]
Then the functor $E\mapsto E|U$ from the category of extensions of $P$ by
$\underline G_{mS}$ to the category of extensions of $P_U$ by
$\underline G_{mU}$ is an equivalence of categories.  In terms of the
description VII \SGARefLink{sga7.sga7i.evii.n1-1-6}{1.1.6} this assertion is essentially trivial: for every
smooth $S$-scheme $X$ (such as $P$, $P\times P$,
$P\times P\times P$, and so on), $X$ is locally factorial and
$X_T=X\setminus X_U$ has codimension at least $2$ in $X$.  Consequently,
the functor $F\mapsto F|X_U$ from the category of torsors under
$\underline G_{mX}$ to the category of torsors under
$\underline G_{mX_U}$ is an equivalence of categories.

There are entirely analogous results for the problem of extending
biextensions; we leave their formulation to the reader.


% END INCLUDED SGA7I_English_current_buildable_source_20260823/components/123_expose_VIII_remark_7_2.tex

% BEGIN INCLUDED SGA7I_English_current_buildable_source_20260823/components/124_expose_VIII_section_7_3_1.tex
% SOURCE AUTHORITY:
% expose_VIII_body.tex lines 2262--2314; scan indices 314--315;
% source folios 303--304.
% Direct scan comparison: physical PDF pages 315--316 in overlapping
% 1100-dpi bands. Both triangular diagrams were checked node by node and
% arrow by arrow; their directions, label sides, and terminal punctuation
% are clear at this resolution, so no higher-resolution escalation was needed.
% Transparent source dispositions:
% - typewritten letter-o special-fibre subscripts are normalized to zero;
% - the coefficient of the pulled-back extension in 7.3.1 is restored from
%   the printed bare G_m to the contextually required G_{m\eta};
% - the malformed phrase "si on des morphismes" is rendered "if we have".
% Hard stop before 7.3.2, source line 2316 / scan index 315.

\medskip
\noindent \SGARefTarget{sga7.sga7i.eviii.n7-3}{7.3}.  We shall make explicit some functorial properties of the
obstructions defined in (\SGARefLink{sga7.sga7i.eviii.n7-1-1}{7.1.1}) and (\SGARefLink{sga7.sga7i.eviii.n7-1-2}{7.1.2}), with $P,Q,S$ variable, and
some relations among them.  In what follows, we shall suppose that the
conditions allowing these obstructions to be defined are satisfied.  Since
the verifications present no difficulty, we leave them to the reader.

\medskip
\noindent \SGARefTarget{sga7.sga7i.eviii.n7-3-1}{7.3.1}.  Let
\[
 u:P'\longrightarrow P
\]
be a morphism of smooth group schemes over $S$, and let $E_\eta$ be an
extension of $P_\eta$ by $\underline G_{m\eta}$.  Pullback by
$u_\eta:P'_\eta\to P_\eta$ gives an extension $E'_\eta$ of $P'_\eta$ by
$\underline G_{m\eta}$.  With this notation, the following diagram is
commutative:
\begin{equation}\tag{\SGARefTarget{sga7.sga7i.eviii.n7-3-1-1}{7.3.1.1}}
\begin{tikzcd}[row sep=large,column sep=huge]
 P'_0/{P'_0}^{\circ}
   \arrow[dd,"u_0"']
   \arrow[rd,"d(E'_\eta)"] & \\
 & (\mathbb Q/\mathbb Z)_k \\
 P_0/P_0^{\circ}
   \arrow[ru,"d(E_\eta)"'] &
\end{tikzcd}
\qquad .
\end{equation}

Similarly, if we have morphisms of smooth group schemes over $S$
\[
 u:P'\longrightarrow P,\qquad v:Q'\longrightarrow Q,
\]
and a biextension $E_\eta$ of $(P_\eta,Q_\eta)$ by
$\underline G_{m\eta}$, whose pullback by $(u_\eta,v_\eta)$ is a
biextension $E'_\eta$ of $(P'_\eta,Q'_\eta)$ by
$\underline G_{m\eta}$, then the following diagram is commutative:
\begin{equation}\tag{\SGARefTarget{sga7.sga7i.eviii.n7-3-1-2}{7.3.1.2}}
\begin{tikzcd}[row sep=large,column sep=huge]
 P'_0/{P'_0}^{\circ}\otimes Q'_0/{Q'_0}^{\circ}
   \arrow[dd,"u_0\otimes v_0"']
   \arrow[rd,"d(E'_\eta)"] & \\
 & (\mathbb Q/\mathbb Z)_k \\
 P_0/P_0^{\circ}\otimes Q_0/Q_0^{\circ}
   \arrow[ru,"d(E_\eta)"'] &
\end{tikzcd}
\end{equation}
The essentially mechanical verification of these compatibilities is left
to the reader.

% END INCLUDED SGA7I_English_current_buildable_source_20260823/components/124_expose_VIII_section_7_3_1.tex

% BEGIN INCLUDED SGA7I_English_current_buildable_source_20260823/components/125_expose_VIII_sections_7_3_2_and_7_3_3.tex
% SOURCE AUTHORITY:
% expose_VIII_body.tex lines 2316--2361; scan indices 315--317;
% source folios 304--306.
% Direct scan comparison: physical PDF pages 316--318 in overlapping
% 1100-dpi bands. The kernel formulas and the symmetry diagram were checked
% directly; all nodes, arrows, label sides, and the closing comma are clear.
% Transparent source dispositions:
% - typewritten letter-o special-fibre subscripts are normalized to zero;
% - the biextension denoted once by E rather than E_eta in 7.3.2 is rendered
%   consistently as the stated generic-fibre biextension E_eta;
% - "biextension symetrique" in the construction of {}^sE is rendered
%   "transposed biextension" to distinguish that operation from the later
%   property E_eta \simeq {}^sE_eta of being symmetric.
% Hard stop before 7.3.4, source line 2367 / scan index 317.

\medskip
\noindent \SGARefTarget{sga7.sga7i.eviii.n7-3-2}{7.3.2}.  It follows from the compatibility (\SGARefLink{sga7.sga7i.eviii.n7-3-1-1}{7.3.1.1}) that, if
$E_\eta$ is an extension of $P_\eta$ by $\underline G_{m\eta}$, there is a
largest group subscheme $V$ of $P$ such that the restriction of $E_\eta$ to
$V_\eta$ extends to an extension of $V$ by $\underline G_{mS}$.  It is the
open subgroup $V$ of $P$ containing $P_\eta$ and defined by the condition
\begin{equation}\tag{\SGARefTarget{sga7.sga7i.eviii.n7-3-2-1}{7.3.2.1}}
 V_0/V_0^{\circ}=\ker d(E_\eta).
\end{equation}
If $u:P'\to P$ is then a morphism of smooth group schemes over $S$, the
extension $E'_\eta$, obtained by pulling $E_\eta$ back along $u_\eta$,
extends to an extension of $P'$ by $\underline G_{mS}$ if and only if
\begin{equation}\tag{\SGARefTarget{sga7.sga7i.eviii.n7-3-2-2}{7.3.2.2}}
 u(P')\subset V.
\end{equation}

In the case of a biextension $E_\eta$ of $(P_\eta,Q_\eta)$ by
$\underline G_{m\eta}$, there is likewise a largest group subscheme $V$ of
$P$ such that the restriction of $E_\eta$ to $(V_\eta,Q_\eta)$ extends to
a biextension of $(V,Q)$ by $\underline G_{mS}$.  It is the open subgroup
of $P$ containing the generic fibre $P_\eta$ and characterized by the fact
that $V_0/V_0^{\circ}$ is the left annihilator in $\Phi$ of the pairing
(\SGARefLink{sga7.sga7i.eviii.n7-1-2}{7.1.2}), that is, the kernel of the corresponding homomorphism from $\Phi$
to $\underline{\mathrm{Hom}}(\Psi,(\mathbb Q/\mathbb Z)_k)$:
\begin{equation}\tag{\SGARefTarget{sga7.sga7i.eviii.n7-3-2-3}{7.3.2.3}}
 V_0/V_0^{\circ}
 =\ker\!\left(
   \Phi\xrightarrow{\ d(E_\eta)\ }
   \underline{\mathrm{Hom}}(\Psi,(\mathbb Q/\mathbb Z)_k)
 \right).
\end{equation}
If $u:P'\to P$ is a morphism of smooth group schemes over $S$, the
biextension of $(P'_\eta,Q_\eta)$ by $\underline G_{m\eta}$ obtained by
pulling $E_\eta$ back along $u_\eta$ extends to a biextension of $(P',Q)$
by $\underline G_{mS}$ if and only if (\SGARefLink{sga7.sga7i.eviii.n7-3-2-2}{7.3.2.2}) holds.  Symmetric results
hold upon interchanging the roles of $P$ and $Q$.

Of course, in general there is no largest \underline{pair} of subgroups
$P'\to P$, $Q'\to Q$ for which the restriction $E'_\eta$ of $E_\eta$ to a
biextension of $(P'_\eta,Q'_\eta)$ by $\underline G_{m\eta}$ extends to a
biextension of $(P',Q')$ by $\underline G_{mS}$.  One can only say that,
for a given pair $(P',Q')$, or more generally for a given pair of
morphisms $P'\to P$, $Q'\to Q$ of smooth group schemes, the desired
extension exists if and only if the restriction of the form $d(E_\eta)$
to
\[
 P'_0/{P'_0}^{\circ}\otimes Q'_0/{Q'_0}^{\circ}
\]
is zero.

\medskip
\noindent \SGARefTarget{sga7.sga7i.eviii.n7-3-3}{7.3.3}.  \underline{Symmetry}.  Let $E_\eta$ be a biextension of
$(P_\eta,Q_\eta)$ by $\underline G_{m\eta}$, and consider the transposed
biextension ${}^{s}E_\eta$ (VII \SGARefLink{sga7.sga7i.evii.n2-7}{2.7}), which is a biextension of
$(Q_\eta,P_\eta)$ by $\underline G_{m\eta}$.  The following diagram is
commutative:
\begin{equation}\tag{\SGARefTarget{sga7.sga7i.eviii.n7-3-3-1}{7.3.3.1}}
\begin{tikzcd}[row sep=large,column sep=huge]
 \Phi\otimes\Psi
   \arrow[dd,"\mathrm{sym}"']
   \arrow[rd,"d(E_\eta)"] & \\
 & (\mathbb Q/\mathbb Z)_k \\
 \Psi\otimes\Phi
   \arrow[ru,"d({}^{s}E_\eta)"'] &
\end{tikzcd}
\qquad ,
\end{equation}
where the vertical arrow is the canonical symmetry.  (Compare this
statement with the \underline{antisymmetry} statement of \SGARefLink{sga7.sga7i.eviii.n2-2-10}{2.2.10}.)  It
follows in particular that, if $P=Q$ and $E_\eta$ is a
\underline{symmetric} biextension of $(P_\eta,Q_\eta)$ by
$\underline G_{m\eta}$, that is, if $E_\eta$ is isomorphic to
${}^{s}E_\eta$, then the bilinear form $d(E_\eta)$ on $\Phi\times\Phi$
with values in $(\mathbb Q/\mathbb Z)_k$ is \underline{symmetric} (and not
alternating, as in \SGARefLink{sga7.sga7i.eviii.n2-2-12}{2.2.12}!).  If it is zero, that is, if $E_\eta$ extends
to a biextension $E$ of $(P,P)$ by $\underline G_{mS}$, the latter is
necessarily symmetric as well: by \SGARefLink{sga7.sga7i.eviii.n7-1}{7.1}(a), every isomorphism
${}^{s}E_\eta\to E_\eta$ extends uniquely to an isomorphism
${}^{s}E\to E$.

% END INCLUDED SGA7I_English_current_buildable_source_20260823/components/125_expose_VIII_sections_7_3_2_and_7_3_3.tex

% BEGIN INCLUDED SGA7I_English_current_buildable_source_20260823/components/126_expose_VIII_sections_7_3_4_and_7_3_5.tex
% SOURCE AUTHORITY:
% expose_VIII_body.tex lines 2367--2446; scan indices 317--321;
% source folios 306--310.
% Direct scan comparison: physical PDF pages 318--322 in overlapping
% 1100-dpi bands. Formulas, primes, residue-field subscripts, and the
% ramification index are clear at this resolution; no higher escalation was
% needed.
% Transparent source dispositions:
% - formula (7.3.4.1) is restored from the transcription's impossible E_*
%   to the source scan's q_eta^*(E_eta);
% - the family P_{eta Q_eta} is written explicitly as
%   P_eta x_eta Q_eta;
% - typewritten letter-o special-fibre subscripts are normalized to zero;
% - the malformed phrase "a l'aise du changement de base" is rendered
%   "by base change".
% Hard stop before Proposition 7.4, source line 2449 / scan index 321.

\medskip
\noindent \SGARefTarget{sga7.sga7i.eviii.n7-3-4}{7.3.4}.  \underline{Reduction of obstructions of type (\SGARefLink{sga7.sga7i.eviii.n7-1-2}{7.1.2}) to
those of type (\SGARefLink{sga7.sga7i.eviii.n7-1-1}{7.1.1}).}

Suppose, for definiteness, that there is an integer $n>0$ such that
$n\Phi=0$, and let
\[
 q\in Q(S)
\]
be a section of $Q$.  By an abuse of notation, write
\[
 q_0\in\Psi(k)
\]
for the image in $\Psi(k)$ of the value of $q$ at $s$.  Regarding $E_\eta$
as an extension of $P_\eta\times_\eta Q_\eta$ by
$\underline G_{mQ_\eta}$ over $Q_\eta$, pullback gives an extension
\[
 q_\eta^*(E_\eta),\qquad
 \text{an extension of $P_\eta$ by $\underline G_{m\eta}$},
\]
and hence an obstruction
\begin{equation}\tag{\SGARefTarget{sga7.sga7i.eviii.n7-3-4-1}{7.3.4.1}}
 d_q\overset{\mathrm{dfn}}{=}
 d\bigl(q_\eta^*(E_\eta)\bigr):
 \Phi\longrightarrow(\mathbb Q/\mathbb Z)_k.
\end{equation}
With $d=d(E_\eta)$, one has the relation
\begin{equation}\tag{\SGARefTarget{sga7.sga7i.eviii.n7-3-4-2}{7.3.4.2}}
 d_q(x)=d(x,q_0),
\end{equation}
valid for every point $x$ of $\Phi$ with values in a $k$-scheme $T$.

\smallskip
\noindent \SGARefTarget{sga7.sga7i.eviii.n7-3-4-3}{7.3.4.3}.  For every $q\in Q(S)$, this formula determines $d_q$
in terms of the pairing $d$; conversely, it plainly recovers $d$ from the
invariants $d_q$, at least when every point of $\Psi$ comes from a section
of $Q$.  By Hensel's lemma, this condition holds automatically when $S$ is
strictly local.  The pairing $d$ is known once it is known over the
separable closure of $k$, and its formation commutes with etale localization
on $S$ by \SGARefLink{sga7.sga7i.eviii.n7-3-5-3}{7.3.5.3} below.  Thus, after replacing $S$ by its strict
localization if necessary, the determination of $d$ reduces to that of the
invariants $d_q$.

In particular, when every point of $\Psi$ comes from a section $q$ of $Q$,
the pairing $d$ is zero---that is, $E_\eta$ extends to a biextension $E$ of
$(P,Q)$ by $\underline G_m$---if and only if $d_q$ is zero for every
section $q$ of $Q$, that is, if and only if $q_\eta^*(E_\eta)$ extends to
an extension of $P$ by $\underline G_m$.  With no hypothesis on $\Psi$,
the analogous criterion remains valid provided that the $q$ are allowed to
be points with values in local etale $S$-schemes $S'$.

\smallskip
\noindent \SGARefTarget{sga7.sga7i.eviii.n7-3-4-4}{7.3.4.4}.  Conversely, starting from the situation of \SGARefLink{sga7.sga7i.eviii.n7-1}{7.1}(a), the
obstruction (\SGARefLink{sga7.sga7i.eviii.n7-1-1}{7.1.1}) may be interpreted as an obstruction of type (\SGARefLink{sga7.sga7i.eviii.n7-1-2}{7.1.2}):
take $Q=\mathbb Z_S$ and the biextension of $(P_\eta,\mathbb Z_\eta)$
deduced from $E_\eta$.  Indeed, it suffices to apply the preceding
construction with $q$ equal to the section $1$ of $\mathbb Z_S$.

\medskip
\noindent \SGARefTarget{sga7.sga7i.eviii.n7-3-5}{7.3.5}.  \underline{Effect of a base change $S'\to S$.}  Suppose
given a dominant local morphism of traits
\[
 f:S'\longrightarrow S,
\]
and denote by $\eta'$ the generic point of $S'$, by $s'$ its closed point,
and put $k'=k(s')$.  If $S=\operatorname{Spec}(V)$ and
$S'=\operatorname{Spec}(V')$, then $V'$ is a discrete valuation ring
dominating the discrete valuation ring $V$.  Put
\begin{equation}\tag{\SGARefTarget{sga7.sga7i.eviii.n7-3-5-1}{7.3.5.1}}
 e(S'/S)=e(V'/V)=\operatorname{length}_{V'}(V'/\mathfrak m V'),
\end{equation}
where $\mathfrak m$ is the maximal ideal of $V$.

Let now $P$ be as in \SGARefLink{sga7.sga7i.eviii.n7-1}{7.1}(a), let $E_\eta$ be an extension of $P_\eta$ by
$\underline G_{m\eta}$, put $P'=P\times_S S'$, and let $E'_{\eta'}$ be
the extension of $P'_{\eta'}$ by $\underline G_{m\eta'}$ obtained by
pulling back $E_\eta$.  With the notation of \SGARefLink{sga7.sga7i.eviii.n4-1}{4.1}(a), there is a canonical
isomorphism
\[
 \Phi'\overset{\mathrm{dfn}}{=}P'_0/{P'_0}^{\circ}
 \simeq\Phi\otimes_k k'\simeq\Phi_{k'},
\]
and hence one may consider
\[
 d(E_\eta)_{k'}:\Phi'\longrightarrow(\mathbb Q/\mathbb Z)_{k'}.
\]
Then
\begin{equation}\tag{\SGARefTarget{sga7.sga7i.eviii.n7-3-5-2}{7.3.5.2}}
 d(E'_{\eta'})=e(S'/S)\,d(E_\eta)_{k'}.
\end{equation}

Likewise, under the hypotheses of \SGARefLink{sga7.sga7i.eviii.n7-1}{7.1}(b) and with the evident notation,
(\SGARefLink{sga7.sga7i.eviii.n7-3-5-2}{7.3.5.2}) holds with $d(E_\eta)$ now denoting the pairing (\SGARefLink{sga7.sga7i.eviii.n7-1-2}{7.1.2}),
$d(E_\eta)_{k'}$ the pairing
\[
 \Phi'\otimes\Psi'\longrightarrow(\mathbb Q/\mathbb Z)_{k'}
\]
obtained from it by base change $k\to k'$, and $E'_{\eta'}$ the
biextension of $(P'_{\eta'},Q'_{\eta'})$ by
$\underline G_{m\eta'}$ obtained from the biextension $E_\eta$ of
$(P_\eta,Q_\eta)$ by $\underline G_{m\eta}$ through the base change
$\eta'\to\eta$.

\smallskip
\noindent \SGARefTarget{sga7.sga7i.eviii.n7-3-5-3}{7.3.5.3}.  An interesting case is $e(S'/S)=1$, that is,
\[
 \mathfrak m'=\mathfrak m V',
\]
where $\mathfrak m'$ denotes the maximal ideal of $V'$.  This holds in
particular if $V'$ is etale over $V$, or if $V'$ is the henselization,
strict henselization, or completion of $V$.  In this case (\SGARefLink{sga7.sga7i.eviii.n7-3-5-2}{7.3.5.2}) simply
says that \underline{the formation of obstruction (\SGARefLink{sga7.sga7i.eviii.n7-1-1}{7.1.1}), respectively
(\SGARefLink{sga7.sga7i.eviii.n7-1-2}{7.1.2}), commutes with the base change under consideration $S'\to S$}.

\smallskip
\noindent \SGARefTarget{sga7.sga7i.eviii.n7-3-5-4}{7.3.5.4}.  On the other hand, (\SGARefLink{sga7.sga7i.eviii.n7-3-5-2}{7.3.5.2}) in the general case shows
that, for every extension $E_\eta$ as in \SGARefLink{sga7.sga7i.eviii.n7-1}{7.1}(a) (respectively every
biextension $E_\eta$ as in \SGARefLink{sga7.sga7i.eviii.n7-1}{7.1}(b)), there is a \underline{finite} morphism
of traits $S'\to S$ such that $k(\eta')$ is a separable extension of
$k(\eta)$ and such that the extension (respectively biextension)
$E'_{\eta'}$ obtained from $E_\eta$ by base change $\eta'\to\eta$ extends
to an extension (respectively biextension) $E'$ over all of $S'$.  Indeed,
it is enough to choose $S'/S$ so that $e(S'/S)$ is a multiple of the
integer $n>0$ occurring in \SGARefLink{sga7.sga7i.eviii.n7-1}{7.1}, and such an extension can always be found
with $k(\eta')$ separable over $k(\eta)$.

% END INCLUDED SGA7I_English_current_buildable_source_20260823/components/126_expose_VIII_sections_7_3_4_and_7_3_5.tex

% BEGIN INCLUDED SGA7I_English_current_buildable_source_20260823/components/127_expose_VIII_proposition_7_4_corollary_7_5_and_bibliography.tex
% SOURCE AUTHORITY:
% expose_VIII_body.tex lines 2449--2495; scan indices 321--323;
% source folios 310--312.
% Direct scan comparison: physical PDF pages 322--324 in overlapping
% 1100-dpi bands. The category names, product bases, superscripts, theorem
% underlining, and bibliography were all clear; no higher-resolution
% escalation was needed.
% Transparent source dispositions:
% - source spacing around Z-algebra is normalized without changing meaning;
% - bibliographic accents and punctuation follow the printed authority where
%   the transcription contained obvious keyboard errors.
% This component closes Expose VIII at the end of its bibliography.

\medskip
\noindent\underline{Proposition \SGARefTarget{sga7.sga7i.eviii.n7-4}{7.4}.  Let $S$ be a reduced scheme.}

\smallskip
\noindent a) \underline{Let $P$ be a smooth group scheme over $S$ whose
fibres at the maximal points of $S$ are connected.  Then the natural
functor}
\begin{equation}\tag{\SGARefTarget{sga7.sga7i.eviii.n7-4-1}{7.4.1}}
 \mathrm{EXT}(P,\underline G_{mS})
 \longrightarrow
 \mathrm{TORSRIG}(P,e_P;\underline G_{mS})
\end{equation}
\underline{from the category of extensions of $P$ by $\underline G_{mS}$
to the category of torsors under $\underline G_{mP}$ rigidified along the
identity section of $P$ is fully faithful.  If $S$ is normal and $P$ has
connected fibres, an object $E$ of the target of (\SGARefLink{sga7.sga7i.eviii.n7-4-1}{7.4.1}) belongs to the
essential image of this functor if and only if, for every maximal point
$\eta$ of $S$, the rigidified torsor $E_\eta$ on $P_\eta$ comes from an
extension of $P_\eta$ by $\underline G_{m\eta}$.}

\smallskip
\noindent b) \underline{Let $P,Q$ be two smooth group schemes over $S$
whose fibres at the maximal points of $S$ are connected.  Then the natural
functor}
\begin{equation}\tag{\SGARefTarget{sga7.sga7i.eviii.n7-4-2}{7.4.2}}
 \mathrm{BIEXT}(P,Q;\underline G_{mS})
 \longrightarrow
 \mathrm{TORSBIRIG}(P,Q;\underline G_{mS})
\end{equation}
\underline{from the category of biextensions of $(P,Q)$ by
$\underline G_{mS}$ to the category of torsors under
$\underline G_{m,P\times_S Q}$ birigidified along the pair of identity
sections of $P$ and $Q$ is fully faithful.  If $S$ is normal and either $P$
or $Q$ has connected fibres, an object $E$ of the target of (\SGARefLink{sga7.sga7i.eviii.n7-4-2}{7.4.2}) belongs
to the essential image of the functor if and only if, for every maximal
point $\eta$ of $S$, $E_\eta$ comes from a biextension of
$(P_\eta,Q_\eta)$ by $\underline G_{m\eta}$.}

We prove (a); the essentially identical proof of (b) is left to the reader.
For the full-faithfulness assertion, note that \SGARefLink{sga7.sga7i.eviii.n4-1}{4.1} remains valid if $X$ and
$Y$ are flat schemes over a reduced scheme $S$ whose fibres at the maximal
points satisfy the conditions stated in \SGARefLink{sga7.sga7i.eviii.n4-1}{4.1}.  Indeed, the family of maximal
fibres of $X\times_S Y$ (which is flat over the reduced scheme $S$) is
schematically dense in $X\times_S Y$.  Hence a morphism from
$X\times_S Y$ to the separated $S$-scheme $\underline G_{mS}$ is trivial
as soon as it is trivial on the maximal fibres.  Thus the full faithfulness
of (\SGARefLink{sga7.sga7i.eviii.n7-4-1}{7.4.1}) is proved as in \SGARefLink{sga7.sga7i.eviii.n4-2}{4.2}, and the characterization of its essential
image given there remains valid.

It remains to prove that, if the rigidified torsor $E$ on $P$ is such that
the $E_\eta$ come from extensions, then so does $E$, assuming that $S$ is
normal and $P$ has connected fibres.  The latter condition implies that
$P$ is of finite presentation over $S$ (SGA 3 $\mathrm{VI}_B$ 5.5).
Full faithfulness also shows that the question is local on $S$, so that we
may suppose $S$ affine and local.  The usual limit argument, using
Nagata's theorem on the finiteness of the integral closure of a reduced
finitely generated $\mathbb Z$-algebra (EGA IV 7.8), reduces us further to
the case in which $S$ is noetherian.

Then \SGARefLink{sga7.sga7i.eviii.n7-1}{7.1}(a) shows that there is an open subset $U\subset S$ whose complement
has codimension at least $2$ and such that $E_U$ comes from an extension of
$P_U$ by $\underline G_{mU}$.  In other words, the desired isomorphism
\[
 \pi^*(E)\simeq
 \operatorname{pr}_1^*(E)\operatorname{pr}_2^*(E)
\]
can be constructed over $(P\times_S P)_U$.  This is an open subset of the
normal scheme $P\times_S P$ whose complement has codimension at least $2$.
The isomorphism on $(P\times_S P)_U$ therefore extends over all of
$P\times_S P$, completing the proof.

\medskip
\noindent\underline{Corollary \SGARefTarget{sga7.sga7i.eviii.n7-5}{7.5}.  Let $S$ be a normal scheme and let $P$
and $Q$ be two smooth group schemes over $S$, with connected fibres at the
maximal points, either $P$ or $Q$ having connected fibres.  Suppose in
addition that, for every maximal point $\eta$ of $S$, either $P_\eta$ or
$Q_\eta$ is an extension of an abelian scheme by a torus (or, more
generally, admits a composition series whose factors are abelian schemes,
tori, or groups $\underline G_a$).  Then the functor (\SGARefLink{sga7.sga7i.eviii.n7-4-2}{7.4.2}) is an
equivalence of categories.}

This is an immediate consequence of \SGARefLink{sga7.sga7i.eviii.n7-4}{7.4}(b) and \SGARefLink{sga7.sga7i.eviii.n4-7}{4.7}.

\bigskip
\noindent\underline{BIBLIOGRAPHY}

\begin{enumerate}
\item[{[1]}] M. Demazure, J. Giraud, M. Raynaud,
  \emph{Sch\'emas ab\'eliens}, S\'eminaire Orsay 1967/68, forthcoming.
\item[{[2]}] M. Raynaud, \emph{Faisceaux amples sur les sch\'emas en groupes
  et les espaces homog\`enes}, Paris thesis, 1968 (forthcoming in Lecture
  Notes, Springer).
\item[{[3]}] M. Rosenlicht, \emph{Some basic theorems on algebraic groups},
  Amer. Jour. Math. \textbf{78}, no.~2 (1956), pp.~401--443.
\item[{[3bis]}] M. Rosenlicht, \emph{Some rationality questions on algebraic
  groups}, Annali di Matematica, Series 4, no.~43 (1957), pp.~25--50.
\item[{[4]}] J.-P. Serre, \emph{Corps locaux}, Act. Sci. Ind. 1296 (1962),
  Hermann, Paris.
\end{enumerate}

% END INCLUDED SGA7I_English_current_buildable_source_20260823/components/127_expose_VIII_proposition_7_4_corollary_7_5_and_bibliography.tex

\clearpage
\phantomsection\addcontentsline{toc}{section}{Exposé IX}
% BEGIN INCLUDED SGA7I_English_current_buildable_source_20260823/components/128_expose_IX_title_and_contents.tex
% SOURCE AUTHORITY:
% expose_IX_body.tex lines 5--64; scan indices 324--325; source folios 313--314.
% Direct scan comparison: physical PDF pages 325--326 at 1100 dpi.
% Hard stop before Introduction 0, source line 65.

\begin{center}
\SGARefAliasTarget{sga7.sga7i.eix.nroot}SGA 7

\vspace{1em}
\underline{Expos\'e IX}

\vspace{1em}
\underline{\underline{N\'ERON MODELS AND MONODROMY}}

\vspace{1em}
\underline{by A. Grothendieck}

(with an appendix by M. Raynaud)
\end{center}

\subsection*{Contents}

\begin{enumerate}[label=\arabic*.]
\setcounter{enumi}{-1}
\item Introduction
\item N\'eron model of abelian schemes: notation; the canonical pairing
  $\Phi_o\times\Phi'_o\longrightarrow(\mathbb Q/\mathbb Z)_k$, and the
  canonical biextension $W^o$ of $(A^o,{A'}^o)$
\item Fixed part and toric part of $T_\ell(A_K)$. Criteria for good
  reduction. Orthogonality theorem for $\ell\ne p$
\item The semistable-reduction case. The semistable reduction theorem
\item Application to a conjecture of Serre--Tate and to the conductor
\item The orthogonality theorem for $\ell=p$ (semistable case). Duality
  of the abelian schemes $B_o$ and $B'_o$. Characterization of the fixed
  part. Criteria for good reduction, essentially good reduction, and
  semistable reduction
\item Remarks on the construction of the toric part and the fixed part
  of $T_p(A)$ in the non-semistable-reduction case
\item The Raynaud extension $G^\natural$ over $S$ attached to the
  abelian scheme $A_K$ with semistable reduction. Duality of the abelian
  schemes $B,B'$ over $S$
\item The Raynaud extension $A_K^{\natural\,o}$ in the semistable case,
  and the ind-group $A_K^\natural$
\item Definition of the monodromy pairing
  $\underline M_\ell\otimes\underline M'_\ell\longrightarrow
  \mathbb Z_{\ell S}$
\item Properties of the monodromy pairing. Integrality and positivity
  theorem
\item Connected components of the N\'eron model: duality relation,
  asymptotic behavior
\item N\'eron model of a Jacobian and the Picard--Lefschetz formula
\item Relations with transcendental theory: the complex-analytic case
  \begin{enumerate}[label=\arabic{enumi}.\arabic*.]
  \item Review of analytic groups over $S$
  \item Formal completion along a fiber
  \item The analytic Raynaud extension
  \item Interpretation in terms of Deligne's ``mixed Hodge structures''
  \item The relatively algebraic case
  \end{enumerate}
\item Relations with transcendental theory: the rigid-analytic case.
  The canonical homomorphism
  $\underline M_K\longrightarrow A_K^{\natural\,o}$
\end{enumerate}


% END INCLUDED SGA7I_English_current_buildable_source_20260823/components/128_expose_IX_title_and_contents.tex

% BEGIN INCLUDED SGA7I_English_current_buildable_source_20260823/components/129_expose_IX_introduction_and_leitfaden.tex
% SOURCE AUTHORITY:
% expose_IX_body.tex lines 65--140; scan indices 325--329; source folios 314--318.
% Direct scan comparison: physical PDF pages 326--330 at 1100 dpi.
% Leitfaden checked from four direct-source crops at 5000 dpi; all 18 strokes
% are plain, unlabelled, undirected segments. Hard stop before source line 144.

0. \underline{Introduction}

\SGARefTarget{sga7.sga7i.eix.n0-1}{0.1}. In the present expos\'e we apply the results of the preceding
expos\'es to the study of the \underline{N\'eron model} (\SGARefLink{sga7.sga7i.eix.n1-1}{1.1}) $A$ of an
abelian scheme $A_K$ defined over the function field $K$ of the
Henselian trait $S$. We shall see that the most important properties of
the N\'eron model can be interpreted in terms of the action of the local
monodromy group $I$ on the Tate module
$T_\ell(A_K(\overline K))$ (where $\ell$ is a prime number distinct
from $p$). Since the dual of the latter can be interpreted as the
$\ell$-adic cohomology group $H^1(A_{\overline K},\mathbb Z_\ell)$,
and since, for every smooth projective scheme $X$ over $K$,
$H^1(X_{\overline K},\mathbb Z_\ell)$ is identified with
$H^1(A_{\overline K},\mathbb Z_\ell)$, where $A_{\overline K}$ is the
Albanese variety of $X_{\overline K}$, the present expos\'e may also be
viewed as a detailed study of \underline{local monodromy phenomena for
$\ell$-adic $H^1$} (or, better still, for ``motivic'' $H^1$)
\underline{of smooth projective varieties over} $K$. From this point of
view, it seems clear that the principal results of the present expos\'e
are destined to be subsumed within a ``N\'eron theory'' for motives of
arbitrary weight, that is, for $H^i$ ($\ell$-adic, de Rham, Hodge, and
so forth) with arbitrary $i$, of which at present we are only beginning
to catch a glimpse. (See [9] in this connection, and more particularly
Deligne's Conjectures 9.8 through 9.13 in the cited report.)

\SGARefTarget{sga7.sga7i.eix.n0-2}{0.2}. The most important result of the present expos\'e is the
``\underline{semistable reduction theorem}'' \SGARefLink{sga7.sga7i.eix.n3-6}{3.6}, which asserts that,
after a finite separable extension $K'$ of $K$, the N\'eron model of
$A_{K'}$ has a special fiber whose identity component is an extension
of an abelian scheme by a torus. This result was conjectured by
J.-P. Serre in 1964, and the principle of the proof given here
\footnote{It is developed in a letter from the lecturer to J.-P. Serre
dated 30 October 1964.}
goes back to the same year. It appears in \SGARefLink{sga7.sga7i.eix.n3-5}{3.5} that this property of
``semistable reduction'' is equivalent to the action of the monodromy
group $I$ on $T_\ell(A_K(\overline K))$ being ``essentially unipotent
of step 2'' (where $\ell$ is a fixed prime distinct from the residual
characteristic $p$ of $S$). In this form it was proved, in terms of
$\ell$-adic $H^1$, in Expos\'e III as a special case of the
``\underline{monodromy theorem}'', using resolution of singularities
for excellent schemes of dimension $2$ (due to S. Abhyankar). To prove
the equivalence of the two properties, we shall also need an
``orthogonality theorem'' \SGARefLink{sga7.sga7i.eix.n2-4}{2.4}, due to Igusa [11] in the special case in
which $A_K$ is the Jacobian of a smooth curve $X_K$ over $K$, the
generic fiber of a regular connected projective scheme $X$ over $S$
whose geometric special fiber has no singular points other than
ordinary quadratic singularities. This is indeed a semistable-reduction
case, but it would not suffice for our purposes; it is chiefly in order
to establish the orthogonality theorem in its general form that we use
the convenient tool supplied by the theory of biextensions developed
in the two preceding expos\'es. Let us also point out that D. Mumford
has found a very different proof of the semistable reduction theorem,
one which does not use resolution of singularities and instead uses
his theory of theta functions (unpublished); unfortunately, he must
restrict himself to residual characteristic different from $2$.

\SGARefTarget{sga7.sga7i.eix.n0-3}{0.3}. The semistable reduction theorem is the principal tool used in the
subsequent sections to study N\'eron models, which we undertake here
from the \underline{``algebraic'' point of view}. It also seems
indispensable for studying abelian varieties over $K$ from the
\underline{``rigid-analytic'' point of view}, or from the point of view
of ``theta functions'' (these terms apparently being synonymous when
abelian varieties over complete valued fields are in question), as
developed recently by D. Mumford and M. Raynaud following work of Tate
and Morikawa. (In Sections 13 and 14 we shall say a few words about the
relations between the algebraic and transcendental points of view.)
Finally, the same theorem is a key ingredient in the proof of the
``\underline{semistable reduction theorem for algebraic curves}'' of
Artin--Deligne--Mumford [3],\footnote{P. Deligne gave the proof of this
theorem in the oral Seminar, but it was not considered useful to
reproduce it in the Seminar texts.} which we have already had occasion
to use in the study of the action of monodromy on the fundamental group
(the ``essentially tame action theorem'' of Expos\'e V), and which is
used again in the proof of the integrality and positivity theorem \SGARefLink{sga7.sga7i.eix.n10-4}{10.4}
in the present expos\'e.\footnote{There is another proof of the latter
theorem, by way of rigid-analytic theory, which does not use the
semistable reduction theorem.}

\SGARefTarget{sga7.sga7i.eix.n0-4}{0.4}. The present expos\'e differs from the other expos\'es of the
Seminar in that, in studying the Tate module $T_\ell(A_K)$, we have
systematically taken into account also the case in which $\ell$ equals
the residual characteristic $p$. It turns out that the language of
\underline{Barsotti--Tate groups} [33]\footnote{Or ``$p$-divisible
groups'' in Tate's terminology.} makes it possible to extend to this
case all the most important phenomena established for $\ell\ne p$,
subject on occasion to having available Tate's fundamental theorem
[33, Theorem 4], which at present has been established only when $K$
has characteristic zero. Of course, in the latter case the pro-$p$
Barsotti--Tate group $T_p(A_K)$ may still be interpreted in terms of
the action of the monodromy group $I$ on the $p$-adic module
$T_p(A_K(\overline K))$, but one should bear in mind that most of the
statements suggested by this superficial analogy with the case
$\ell\ne p$ are grossly false. (Thus it is no longer true that the
monodromy action is essentially unipotent when $\ell=p$.) By using more
judicious translations, however (such as \SGARefLink{sga7.sga7i.eix.n5-13}{5.13} below), it is already
clear that the study of pro-$p$ Barsotti--Tate groups defined
geometrically over the function field of a trait of residual
characteristic $p>0$, whether or not the trait has mixed
characteristic, must be regarded as part of the theory of ``local
monodromy'', in which these groups play the role of \underline{$p$-adic
local systems} on $\operatorname{Spec}(K)$. In fact, these are
essentially the ``local systems'' that enter into the \underline{study,
from the ``$p$-adic'' point of view, of $H^1$ of smooth projective
varieties over} $K$. For the still-to-be-developed study of higher
$H^i$ from the $p$-adic point of view, it appears that the category of
pro-$p$ Barsotti--Tate groups over $K$, respectively over $S$, will
have to be suitably enlarged, taking inspiration from the viewpoints
provided by ``crystalline cohomology''.

\SGARefTarget{sga7.sga7i.eix.n0-5}{0.5}. The expos\'e given here follows the oral lectures very closely up
to Section 4, but is markedly more detailed and complete from the
following sections onward, which had been only partially sketched. It
will not be used in the subsequent expos\'es. Some results of the
present expos\'e (\SGARefLink{sga7.sga7i.eix.n10-4}{10.4} and \SGARefLink{sga7.sga7i.eix.n12-5}{12.5}), on the other hand, use the
\underline{Picard--Lefschetz formula}, which will be established only
later.

\begin{center}
\underline{LEITFADEN}

\vspace{0.5em}
\begin{tikzpicture}[line width=0.35pt,every node/.style={fill=white,inner sep=1.2pt}]
  \coordinate (p1)  at ( 0.00, 0.00);
  \coordinate (p2)  at ( 0.00,-0.76);
  \coordinate (p3)  at ( 0.00,-1.53);
  \coordinate (p4)  at (-1.60,-1.92);
  \coordinate (p5)  at ( 0.00,-2.30);
  \coordinate (p6)  at (-1.62,-2.69);
  \coordinate (p7)  at ( 1.32,-2.87);
  \coordinate (p8)  at ( 0.04,-3.35);
  \coordinate (p9)  at ( 3.02,-2.69);
  \coordinate (p10) at ( 3.10,-3.35);
  \coordinate (p11) at ( 2.66,-4.21);
  \coordinate (p12) at ( 3.76,-4.21);
  \coordinate (p13) at ( 0.12,-4.37);
  \coordinate (p14) at ( 1.04,-4.37);
  \draw (p1)--(p2) (p2)--(p3) (p3)--(p4) (p3)--(p5)
        (p5)--(p6) (p5)--(p7) (p5)--(p8) (p5)--(p9)
        (p7)--(p8) (p7)--(p13) (p7)--(p14)
        (p8)--(p13) (p8)--(p14) (p9)--(p10)
        (p10)--(p11) (p10)--(p12) (p10)--(p13) (p10)--(p14);
  \node at (p1)  {1};  \node at (p2)  {2};
  \node at (p3)  {3};  \node at (p4)  {4};
  \node at (p5)  {5};  \node at (p6)  {6};
  \node at (p7)  {7};  \node at (p8)  {8};
  \node at (p9)  {9};  \node at (p10) {10};
  \node at (p11) {11}; \node at (p12) {12};
  \node at (p13) {13}; \node at (p14) {14};
\end{tikzpicture}
\end{center}


% END INCLUDED SGA7I_English_current_buildable_source_20260823/components/129_expose_IX_introduction_and_leitfaden.tex

% BEGIN INCLUDED SGA7I_English_current_buildable_source_20260823/components/130_expose_IX_section_1_0_and_1_1.tex
% Expose IX, 1.0--1.1.
% French transcription lines 144--260; authority scan physical pp. 331--334
% (source folios 319--322).  Ordinary pages were checked in four overlapping
% 1100-dpi bands.  The missing prime on A'_o in the final sentence of 1.1 is
% an ordinary printed/source slip and is restored below.

\medskip
\noindent
\underline{1. Neron models of abelian schemes: notation; the canonical pairing}
\[
 \Phi_o\times\Phi'_o\longrightarrow(\mathbb Q/\mathbb Z)_k,
\]
\underline{and the canonical biextension $W^o$ of $(A^o,{A'}^o)$.}

\medskip
\noindent
\SGARefTarget{sga7.sga7i.eix.n1-0}{1.0}. Throughout this expos\'e we fix a trait $S$, for which we use the
usual notation $s$, $\eta$, $k$, $K$, and so forth, of I~1.  Let $A_K$
be an abelian scheme over $K$, and denote its dual abelian scheme [5] by
$A'_K$.  The duality between $A_K$ and $A'_K$ is given by a divisorial
correspondence $W_K$ on $(A_K,A'_K)$, that is, an invertible sheaf on
$A_K\times A'_K$, birigidified with respect to the identity sections of
$A_K$ and $A'_K$.  It is determined up to unique isomorphism by its
class (cf. VIII~\SGARefLink{sga7.sga7i.eviii.n3-1}{3.1})
\begin{equation}\tag{\SGARefTarget{sga7.sga7i.eix.n1-0-1}{1.0.1}}
 w_K\in\mathrm{Biext}^1(A_K,A'_K;\mathbb G_{mK}).
\end{equation}
An important role will be played by the corresponding pairing
(VIII~\SGARefLink{sga7.sga7i.eviii.n2-2}{2.2})
\begin{equation}\tag{\SGARefTarget{sga7.sga7i.eix.n1-0-2}{1.0.2}}
 \varphi:T_\ell(A_K)\times T_\ell(A'_K)
 \longrightarrow\mathbb Z_\ell(1)_K
 \overset{\mathrm{def}}{=}T_\ell(\mathbb G_{mK}),
\end{equation}
where $\ell$ is an arbitrary prime number.  When
$\ell\ne p=\operatorname{char}k$, giving $\varphi$ is equivalent to
giving a pairing of finite free $\mathbb Z_\ell$-modules
\begin{equation}\tag{\SGARefTarget{sga7.sga7i.eix.n1-0-3}{1.0.3}}
 \varphi:T_\ell(A_K(\overline K))\times
 T_\ell(A'_K(\overline K))\longrightarrow T_\ell(\overline K^*),
\end{equation}
compatible with the action of
\[
 \pi=\operatorname{Gal}(\overline K/K)
\]
on the three modules in question, where $\overline K$ is a separable
closure of $K$.  It is known (VIII~\SGARefLink{sga7.sga7i.eviii.n3-2}{3.2}) that this pairing is a perfect
duality: it identifies each of the modules
$T_\ell(A_K(\overline K))$ and $T_\ell(A'_K(\overline K))$ with the
dual of the other, with values in the invertible module
$\mathbb Z_\ell(1)=T_\ell(\overline K^*)$.

At times we shall suppose that a polarization of $A_K$ has been chosen;
thus we are given an element
\begin{equation}\tag{\SGARefTarget{sga7.sga7i.eix.n1-0-4}{1.0.4}}
 \xi\in\operatorname{NS}(A_K/K)
 =\underline{\operatorname{NS}}_{A_K/K}(K),
\end{equation}
where
\[
 \underline{\operatorname{NS}}_{A_K/K}
 \overset{\mathrm{def}}{=}
 \underline{\operatorname{Pic}}_{A_K/K}/
 \bigl(\underline{\operatorname{Pic}}{}^o_{A_K/K}=A'_K\bigr).
\]
In the usual way it defines a homomorphism
\begin{equation}\tag{\SGARefTarget{sga7.sga7i.eix.n1-0-5}{1.0.5}}
 \psi_\xi:A_K\longrightarrow A'_K,
\end{equation}
hence a homomorphism
\begin{equation}\tag{\SGARefTarget{sga7.sga7i.eix.n1-0-6}{1.0.6}}
 T_\ell(\psi_\xi),\quad\text{or simply}\quad
 \psi_\xi:T_\ell(A_K)\longrightarrow T_\ell(A'_K),
\end{equation}
and a bilinear form $\varphi_\xi$
\begin{equation}\tag{\SGARefTarget{sga7.sga7i.eix.n1-0-7}{1.0.7}}
 \begin{split}
 \varphi_\xi(x,y)&=\varphi(x,\psi_\xi(y)),\\
 \varphi_\xi:T_\ell(A_K)\times T_\ell(A_K)
 &\longrightarrow\mathbb Z_\ell(1)_K,
 \end{split}
\end{equation}
which, for $\ell\ne p$, may also be interpreted as a pairing of
finite free $\mathbb Z_\ell$-modules with operators:
\begin{equation}\tag{\SGARefTarget{sga7.sga7i.eix.n1-0-8}{1.0.8}}
 \varphi_\xi:T_\ell(A_K(\overline K))\times
 T_\ell(A_K(\overline K))\longrightarrow T_\ell(\overline K^*).
\end{equation}
These constructions make sense for every element of the
Neron--Severi group of $A_K$, and yield an \underline{alternating} form
$\varphi_\xi$, in either interpretation (\SGARefLink{sga7.sga7i.eix.n1-0-7}{1.0.7}) or (\SGARefLink{sga7.sga7i.eix.n1-0-8}{1.0.8})
(VIII~\SGARefLink{sga7.sga7i.eviii.n2-2-11}{2.2.11}).  The fact that $\xi$ is a \underline{polarization}
implies that $\psi_\xi$ in (\SGARefLink{sga7.sga7i.eix.n1-0-5}{1.0.5}) is an isogeny; equivalently, since
$\dim A_K=\dim A'_K$, it is an epimorphism of groups.  This is also
equivalent to the injectivity of $T_\ell(\psi_\xi)$ in (\SGARefLink{sga7.sga7i.eix.n1-0-6}{1.0.6}) (or to
its having finite cokernel), and finally to the form $\varphi_\xi$ in
(\SGARefLink{sga7.sga7i.eix.n1-0-7}{1.0.7}) being \underline{separating}.

\medskip
\noindent
\SGARefTarget{sga7.sga7i.eix.n1-1}{1.1}. Recall that there exists a \underline{smooth scheme} $A$ over
$S$, called the \underline{Neron model of $A_K$}, which represents the
functor
\begin{equation}\tag{\SGARefTarget{sga7.sga7i.eix.n1-1-1}{1.1.1}}
 S'\longmapsto\operatorname{Hom}(S'_\eta,A_K)
\end{equation}
on the category of smooth $S$-schemes.  Thus, by definition, there is
an isomorphism, functorial in the smooth relative $S$-scheme $S'$,
\begin{equation}\tag{\SGARefTarget{sga7.sga7i.eix.n1-1-2}{1.1.2}}
 \operatorname{Hom}_S(S',A)\simeq
 \operatorname{Hom}_K(S'_\eta,A_K).
\end{equation}
Neron [19] constructed $A$ when the residue field $k$ of $S$ is
perfect.  Raynaud extended the construction to the general case in an
unwritten seminar held at IHES in 1966--67 (see [21] for the statement
of Raynaud's principal results, which complete those of Neron in
loc. cit.).  We shall admit this existence result here, together with
the fact that $A$ is a quasi-projective \underline{scheme} over $S$.
Applying (\SGARefLink{sga7.sga7i.eix.n1-1-2}{1.1.2}) when $S'$ is smooth \underline{over} $\eta$, one also
finds that the following canonical morphism is an isomorphism:
\begin{equation}\tag{\SGARefTarget{sga7.sga7i.eix.n1-1-3}{1.1.3}}
 A\times_S\operatorname{Spec}(K)\xrightarrow{\ \sim\ }A_K.
\end{equation}
This justifies the paired notation $A,A_K$; when convenient we shall
also write $A_\eta$ in place of $A_K$.  Using (\SGARefLink{sga7.sga7i.eix.n1-1-3}{1.1.3}) as an identifying
isomorphism, the canonical bijection (\SGARefLink{sga7.sga7i.eix.n1-1-2}{1.1.2}) is simply the map given by
\underline{restriction to the generic fibers}.

Since the functor (\SGARefLink{sga7.sga7i.eix.n1-1-1}{1.1.1}) plainly has the structure of a commutative
group, the Neron model $A$ representing it itself has the structure of
a \underline{smooth commutative group scheme over $S$}.  The remainder
of this expos\'e is devoted to the study of this group scheme.

An important invariant arising from the Neron model is, of course, the
smooth commutative group scheme of finite type over $k$
\begin{equation}\tag{\SGARefTarget{sga7.sga7i.eix.n1-1-4}{1.1.4}}
 A_o=A\times_Ss=A_k.
\end{equation}
Following the general notation, we denote its identity component by
$A_o^o$ and also consider the group scheme
\begin{equation}\tag{\SGARefTarget{sga7.sga7i.eix.n1-1-5}{1.1.5}}
 \Phi_o=A_o/A_o^o
\end{equation}
of connected components of $A_o$, which is finite and etale over $k$.
In terms of a separable closure $\overline k$ of $k$, it is therefore
determined by the ordinary finite group
\begin{equation}\tag{\SGARefTarget{sga7.sga7i.eix.n1-1-6}{1.1.6}}
 \Phi_o(\overline k),
\end{equation}
together with the natural action of
\[
 \pi_o=\operatorname{Gal}(\overline k/k)
\]
on this group.

Along with the Neron model $A$ of $A_K$, we shall study
\begin{equation}\tag{\SGARefTarget{sga7.sga7i.eix.n1-1-7}{1.1.7}}
 A'=\text{the Neron model of }A'_K,
\end{equation}
which is likewise a smooth commutative group scheme of finite type over
$S$.  We shall denote by
\begin{equation}\tag{\SGARefTarget{sga7.sga7i.eix.n1-1-8}{1.1.8}}
 \Phi'_o=A'_o/{A'}_o^o
\end{equation}
the group of connected components of its special fiber $A'_o$.

% END INCLUDED SGA7I_English_current_buildable_source_20260823/components/130_expose_IX_section_1_0_and_1_1.tex

% BEGIN INCLUDED SGA7I_English_current_buildable_source_20260823/components/131_expose_IX_sections_1_2_through_1_5.tex
% Expose IX, 1.2--1.5.
% French transcription lines 265--341; authority scan physical pp. 335--338
% (source folios 323--326).  Ordinary pages were checked in four overlapping
% 1100-dpi bands.  Formula (1.4.2) was checked in targeted 5000-dpi crops.

\medskip
\noindent
\SGARefTarget{sga7.sga7i.eix.n1-2}{1.2}. Apply VIII~\SGARefLink{sga7.sga7i.eviii.n7-1}{7.1}(b) to the pair $(A,A')$ and to the biextension
$W_K$ of $(A_K,A'_K)$ by $\mathbb G_{mK}$ in (\SGARefLink{sga7.sga7i.eix.n1-0-1}{1.0.1}).  We thereby
obtain a \underline{canonical pairing}
\begin{equation}\tag{\SGARefTarget{sga7.sga7i.eix.n1-2-1}{1.2.1}}
 \Phi_o\times\Phi'_o\longrightarrow(\mathbb Q/\mathbb Z)_k,
\end{equation}
that is, a pairing of $\pi_o$-modules
\begin{equation}\tag{\SGARefTarget{sga7.sga7i.eix.n1-2-2}{1.2.2}}
 \Phi_o(\overline k)\times\Phi'_o(\overline k)
 \longrightarrow\mathbb Q/\mathbb Z.
\end{equation}
Here the pairing appears as the obstruction to extending $W_K$ to a
biextension $W$ of $(A,A')$ by $\mathbb G_{mS}$; the meaning of this
point of view is explained in VIII~\SGARefLink{sga7.sga7i.eviii.n7-3-2}{7.3.2}.  We record at once the
following conjecture concerning this pairing.

\medskip
\noindent
\underline{Conjecture \SGARefTarget{sga7.sga7i.eix.n1-3}{1.3}.} \underline{The pairing} (\SGARefLink{sga7.sga7i.eix.n1-2-1}{1.2.1})
\underline{is a perfect duality}: \underline{it identifies each of the
component groups} $\Phi_o$ \underline{and} $\Phi'_o$
\underline{with the dual of the other, with values in}
$(\mathbb Q/\mathbb Z)_k$.

\medskip
\noindent
\SGARefTarget{sga7.sga7i.eix.n1-3-1}{1.3.1}. The formation of (\SGARefLink{sga7.sga7i.eix.n1-2-1}{1.2.1}) commutes with formally etale base
change (VIII~\SGARefLink{sga7.sga7i.eviii.n7-3-5-3}{7.3.5.3}), and the formation of Neron models has the
analogous property.  Consequently, to prove Conjecture~\SGARefLink{sga7.sga7i.eix.n1-3}{1.3} it would
suffice to treat the case in which $S$ is complete and has separably
closed residue field.  Artin and Mazur proved the conjecture when
$A_K$ is the Jacobian of a proper smooth curve over $K$, using their
general, unpublished theory of the self-duality of the relative
Jacobian (in the sense of derived categories) of a proper curve scheme
over $S$.  They also proved Conjecture~\SGARefLink{sga7.sga7i.eix.n1-3}{1.3} when the residue field is
\underline{finite}, using Tate duality [32].  We shall likewise see in
\SGARefLink{sga7.sga7i.eix.n11-4}{11.4} and \SGARefLink{sga7.sga7i.eix.n11-5}{11.5} (subject to compatibility checks) that \SGARefLink{sga7.sga7i.eix.n1-3}{1.3} holds on the
$\ell$-primary components for $\ell\ne p$, and that it holds in the
case of ``semistable reduction''.  These results make \SGARefLink{sga7.sga7i.eix.n1-3}{1.3} extremely
plausible.  Moreover, \SGARefLink{sga7.sga7i.eix.n1-3}{1.3} may also be regarded as a by-product of a
general duality theory for group schemes over fraction fields of
complete discrete valuation rings, formally analogous to SGA~5~I.
Such a theory remains heuristic at present; it would also contain
local class field theory in the form given to it by Serre [27].

\medskip
\noindent
\SGARefTarget{sga7.sga7i.eix.n1-3-2}{1.3.2}. For a given abelian scheme $A_K$, the validity of \SGARefLink{sga7.sga7i.eix.n1-3}{1.3} means that
the pairing (\SGARefLink{sga7.sga7i.eix.n1-2-1}{1.2.1}) is \underline{separating} on both the left and the
right.  By VIII~\SGARefLink{sga7.sga7i.eviii.n7-3-2}{7.3.2}, separation on the left says precisely that the
only open subgroup scheme $V$ of $A$ for which $W_K$ can be extended to
a biextension of $(V,A')$ by $\mathbb G_{mS}$ is the identity component
$A^o$ of $A$.  There is a symmetric interpretation of separation on
the right.

\medskip
\noindent
\SGARefTarget{sga7.sga7i.eix.n1-4}{1.4}. More generally, suppose that
\begin{equation}\tag{\SGARefTarget{sga7.sga7i.eix.n1-4-1}{1.4.1}}
 U\subset A\qquad\text{and}\qquad U'\subset A'
\end{equation}
are open subgroup schemes such that $W_K$ extends to a biextension of
$(U,U')$ by $\mathbb G_{mS}$.  By VIII~\SGARefLink{sga7.sga7i.eviii.n7-3-2}{7.3.2}, this says explicitly that
the subgroups
% SGA-SRC-001266
\begin{equation}\tag{\SGARefTarget{sga7.sga7i.eix.n1-4-2}{1.4.2}}
 F_o=U_o/U_o^o\subset\Phi_o,
 \qquad
 F'_o=U'_o/{U'}_o^o\subset\Phi'_o
\end{equation}
defined by $U$ and $U'$ are orthogonal under (\SGARefLink{sga7.sga7i.eix.n1-2-1}{1.2.1}).\footnote{The
printed formula, confirmed in a direct high-resolution check, reuses
$\Phi'_o$ for the first subgroup and writes $\Psi'_o\subset\Psi_o$ for
the second, colliding with the established notation for the component
group of $A'$.  The neutral local names $F_o,F'_o$ above make the
intended two subgroups unambiguous without changing the assertion.}
We denote this extension (unique up to unique isomorphism) of $W_K$ by
$W$.  It is determined up to unique isomorphism by its class
\begin{equation}\tag{\SGARefTarget{sga7.sga7i.eix.n1-4-3}{1.4.3}}
 w\in\mathrm{Biext}^1(U,U';\mathbb G_{mS}),
\end{equation}
which is itself characterized by the condition that its image in
$\mathrm{Biext}^1(A_K,A'_K;\mathbb G_{mK})$ be the class $w_K$ of
(\SGARefLink{sga7.sga7i.eix.n1-0-1}{1.0.1}).

For example, one may take for $U$ and $U'$ the identity components
$A^o$ and $A'{}^o$ of $A$ and $A'$.  One thereby obtains a
\underline{canonical biextension $W^o$, or $W$, of $(A^o,A'{}^o)$ by
$\mathbb G_{mS}$}, giving rise to a class
\begin{equation}\tag{\SGARefTarget{sga7.sga7i.eix.n1-4-4}{1.4.4}}
 w=w^o\in\mathrm{Biext}^1(A^o,A'{}^o;\mathbb G_{mS}),
\end{equation}
with which we shall work chiefly.  Evidently, the restrictions to
$(A^o,A'{}^o)$ of the classes $w$ in (\SGARefLink{sga7.sga7i.eix.n1-4-3}{1.4.3}), for the various pairs
$(U,U')$, are all equal to $w^o$ in (\SGARefLink{sga7.sga7i.eix.n1-4-4}{1.4.4}).

\medskip
\noindent
\underline{Proposition \SGARefTarget{sga7.sga7i.eix.n1-5}{1.5}.} \underline{Consider the special fiber
$W_s$ of the canonical biextension} $W$ \underline{of}
$(A^o,A'{}^o)$ \underline{by} $\mathbb G_{mS}$ \underline{in (\SGARefLink{sga7.sga7i.eix.n1-4-4}{1.4.4}),
so that $W_s$ is a biextension of $(A_o^o,{A'}_o^o)$ by}
$\mathbb G_{mk}$.  \underline{This biextension is separating}
(VIII~\SGARefLink{sga7.sga7i.eviii.n4-11}{4.11}).

Indeed, choose an ample invertible sheaf $\mathbb L$ on $A^o$.  Then
$\mathbb L_\eta$ defines a homomorphism (cf. \SGARefLink{sga7.sga7i.eix.n1-0-5}{1.0.5})
\[
 A_K\longrightarrow A'_K,
\]
which extends to a homomorphism $A\longrightarrow A'$, and hence gives
\[
 \psi:A^o\longrightarrow A'{}^o.
\]
We may consider the biextension $W_{\mathbb L}$ of $(A^o,A^o)$ by
$\mathbb G_{mS}$ obtained by pulling back $W$ along
$(\operatorname{id}_{A^o},\psi)$.  Its restriction to the generic
fibers is precisely the biextension $\delta(\mathbb L_\eta)$ of
VIII~\SGARefLink{sga7.sga7i.eviii.n4-12}{4.12}.  Hence, by VIII~\SGARefLink{sga7.sga7i.eviii.n7-4}{7.4}, $W_{\mathbb L}$ is associated with the
birigidified invertible sheaf $\delta(\mathbb L)$, and consequently
$(W_{\mathbb L})_s$ is defined by $\delta(\mathbb L_s)$.  On the other
hand, it may also be described as the pullback of $W_s$ by the
morphisms
\[
 \operatorname{id}_{A_o^o}:A_o^o\longrightarrow A_o^o,
 \qquad
 \psi_o:A_o^o\longrightarrow {A'}_o^o.
\]
Since $\mathbb L$, and hence $\mathbb L_s$, is ample, VIII~\SGARefLink{sga7.sga7i.eviii.n4-13}{4.13} shows
that $W_{\mathbb L_s}$ is separating, in particular separating on the
left.  A fortiori $W_s$ is therefore separating on the left.  By
symmetry it is also separating on the right, which proves \SGARefLink{sga7.sga7i.eix.n1-5}{1.5}.

We shall substantially sharpen \SGARefLink{sga7.sga7i.eix.n1-5}{1.5} when $A_o^o$ has unipotent rank
zero (\SGARefLink{sga7.sga7i.eix.n5-4}{5.4}).

% END INCLUDED SGA7I_English_current_buildable_source_20260823/components/131_expose_IX_sections_1_2_through_1_5.tex

% BEGIN INCLUDED SGA7I_English_current_buildable_source_20260823/components/132_expose_IX_section_2_1.tex
% Expose IX, section 2 title and 2.1--2.1.14.
% French transcription lines 344--441; authority physical pp. 338--341
% (source folios 326--329).  All pages were checked in four overlapping
% 1100-dpi bands.  Diagram (2.1.6) was checked node-by-node and arrow-by-arrow
% in four overlapping direct-source crops rendered at 5000 dpi.

\medskip
\noindent
\underline{2. Fixed and toric parts of $T_\ell(A_K)$.  Criteria for
good reduction.  Orthogonality theorem for $\ell\ne p$.}

\medskip
\noindent
\SGARefTarget{sga7.sga7i.eix.n2-1}{2.1}. Like every commutative algebraic group over $k$, $A_o$ has a
largest subtorus, whose formation commutes with every extension of the
base field [SGA~3 XVII~7.2.1, XII~1.12]:
\begin{equation}\tag{\SGARefTarget{sga7.sga7i.eix.n2-1-1}{2.1.1}}
 T_o\subset A_o.
\end{equation}
When $k$ is perfect, $T_o$ is characterized by the condition that
$A_o^o/T_o$ be an extension of an abelian scheme $B_o$ over $k$ by a
smooth connected unipotent group [25].  The latter necessarily has a
composition series all of whose quotients are isomorphic to
$\mathbb G_{ak}$ (SGA~3 XVII~4.1.3), and one has an exact sequence
\begin{equation}\tag{\SGARefTarget{sga7.sga7i.eix.n2-1-2}{2.1.2}}
 0\longrightarrow U_o\longrightarrow A_o^o/T_o
 \longrightarrow B_o\longrightarrow0.
\end{equation}
In any case, whether or not $k$ is perfect, \underline{if} $A_o^o$ is
an extension of an abelian scheme $B_o$ by a smooth connected affine
scheme $L_o$,
\begin{equation}\tag{\SGARefTarget{sga7.sga7i.eix.n2-1-3}{2.1.3}}
 0\longrightarrow L_o\longrightarrow A_o^o
 \longrightarrow B_o\longrightarrow0,
\end{equation}
this extension structure is known to be unique [25].  Moreover,
$T_o\subset L_o$, and $T_o$ is the largest subtorus of $L_o$.  Thus
one has the exact sequence (\SGARefLink{sga7.sga7i.eix.n2-1-2}{2.1.2}), with
\begin{equation}\tag{\SGARefTarget{sga7.sga7i.eix.n2-1-4}{2.1.4}}
 U_o=L_o/T_o.
\end{equation}

\medskip
\noindent
\SGARefTarget{sga7.sga7i.eix.n2-1-5}{2.1.5}. Let us note that $A_o^o$ has \underline{unipotent rank zero}
(that is, after passage to an algebraic closure of $k$, it has no
subgroup isomorphic to $\mathbb G_a$) if and only if $A_o^o/T_o$ is an
abelian scheme $B_o$.  In that case the preceding hypotheses apply
with $U_o=0$.

Let $\ell$ be a prime different from the characteristic of $k$, so that
$A_o^o$ is $\ell$-divisible.  The inclusion (\SGARefLink{sga7.sga7i.eix.n2-1-1}{2.1.1}) then gives an exact
sequence of pro-group schemes:
\begin{equation}\tag{\SGARefTarget{sga7.sga7i.eix.n2-1-6}{2.1.6}}
\begin{tikzcd}[row sep=large,column sep=large]
0 \arrow[r] & T_\ell(T_o) \arrow[r]
 & T_\ell(A_o^o) \arrow[r] \arrow[d,"\wr"]
 & T_\ell(A_o^o/T_o) \arrow[r] \arrow[d,"\wr"] & 0\\
 & & T_\ell(A_o) & T_\ell(A_o/T_o) & {,}
\end{tikzcd}
\end{equation}
whose terms may also be interpreted as twisted constant free
$\ell$-adic sheaves.  When an exact sequence (\SGARefLink{sga7.sga7i.eix.n2-1-2}{2.1.2}) is available---for
example, when $k$ is perfect or when $A_o^o$ has unipotent rank
zero---one has
\begin{equation}\tag{\SGARefTarget{sga7.sga7i.eix.n2-1-7}{2.1.7}}
 T_\ell(A_o^o/T_o)\xrightarrow{\ \sim\ }T_\ell(B_o),
\end{equation}
so that the exact sequence (\SGARefLink{sga7.sga7i.eix.n2-1-6}{2.1.6}) becomes
\begin{equation}\tag{\SGARefTarget{sga7.sga7i.eix.n2-1-8}{2.1.8}}
 0\longrightarrow T_\ell(T_o)\longrightarrow T_\ell(A_o^o)
 \longrightarrow T_\ell(B_o)\longrightarrow0.
\end{equation}
Put
\begin{equation}\tag{\SGARefTarget{sga7.sga7i.eix.n2-1-9}{2.1.9}}
 n=\dim A_K=\dim A_o,\qquad
 \mu=\dim T_o,\qquad
 \lambda=\dim U_o,\qquad
 \alpha=\dim B_o.
\end{equation}
Then
\begin{equation}\tag{\SGARefTarget{sga7.sga7i.eix.n2-1-10}{2.1.10}}
 n=\mu+\lambda+\alpha,
\end{equation}
and
\begin{equation}\tag{\SGARefTarget{sga7.sga7i.eix.n2-1-11}{2.1.11}}
 \left\{
 \begin{aligned}
 \operatorname{rank}T_\ell(T_o)&=\mu,&
 \operatorname{rank}T_\ell(A_o^o/T_o)
 &=\operatorname{rank}T_\ell(B_o)=2\alpha,\\
 \operatorname{rank}T_\ell(A_o^o)&=\mu+2\alpha.
 \end{aligned}
 \right.
\end{equation}
The second relation is a well-known theorem of Weil [5, 15].

\medskip
\noindent
\SGARefTarget{sga7.sga7i.eix.n2-1-12}{2.1.12}. The integers $n,\mu,\lambda,\alpha$ are plainly defined, after
passing to an algebraic closure of $k$, even without assuming the
existence of an exact sequence (\SGARefLink{sga7.sga7i.eix.n2-1-2}{2.1.2}).  They are respectively called
the dimension of $A_o$, the \underline{reductive rank of $A_o$}
(equal to the dimension of $T_o$), the \underline{unipotent rank of
$A_o$}, and the \underline{abelian rank of $A_o$}.  They are still
related by (\SGARefLink{sga7.sga7i.eix.n2-1-10}{2.1.10}) and (\SGARefLink{sga7.sga7i.eix.n2-1-11}{2.1.11}), with the term
$\operatorname{rank}T_\ell(B_o)$ omitted from the latter.

\medskip
\noindent
\SGARefTarget{sga7.sga7i.eix.n2-1-13}{2.1.13}. Suppose that $k$ has characteristic $p>0$ and that $A_o^o$ is
$p$-divisible, or equivalently that it has unipotent rank zero
(cf. \SGARefLink{sga7.sga7i.eix.n2-1-5}{2.1.5}).  One should then also consider the pro-group $T_p(A_o^o)$,
which is a \underline{Barsotti--Tate group} (or a
\underline{$p$-divisible group}, in the terminology of [33]).  The
exact sequence (\SGARefLink{sga7.sga7i.eix.n2-1-8}{2.1.8}) remains valid for $\ell=p$ and is then an exact
sequence of Barsotti--Tate groups.  Here $\lambda=0$, and hence
\begin{equation}\tag{\SGARefTarget{sga7.sga7i.eix.n2-1-13-1}{2.1.13.1}}
 n=\mu+\alpha.
\end{equation}
The formulas (\SGARefLink{sga7.sga7i.eix.n2-1-11}{2.1.11}) also remain valid, provided that rank is
interpreted as the rank (or ``height'') of the Barsotti--Tate groups in
question.

\medskip
\noindent
\SGARefTarget{sga7.sga7i.eix.n2-1-14}{2.1.14}. For ${A'}_o^o$ we likewise introduce the invariants
\[
 T'_o\subset {A'}_o^o,\qquad U'_o,\qquad B'_o,
\]
where $T'_o$ is the toric part of ${A'}_o^o$, $B'_o$ its abelian part
(defined in particular when $k$ is perfect or when ${A'}_o^o$ has
unipotent rank zero), and $U'_o$ its unipotent part (defined, for
example, when $k$ is perfect, or trivially when ${A'}_o^o$ has
unipotent rank zero).  There is a corresponding exact sequence of the
form (\SGARefLink{sga7.sga7i.eix.n2-1-6}{2.1.6}), or (\SGARefLink{sga7.sga7i.eix.n2-1-8}{2.1.8}), which we shall not write again.  One could
introduce integers $\mu',\lambda',\alpha'$ for ${A'}_o^o$ as in
(\SGARefLink{sga7.sga7i.eix.n2-1-9}{2.1.9}), but we shall see below (\SGARefLink{sga7.sga7i.eix.n2-2-7}{2.2.7}) that they are the same as those
for $A_o^o$; this spares us from introducing new notation for them.

% END INCLUDED SGA7I_English_current_buildable_source_20260823/components/132_expose_IX_section_2_1.tex

% BEGIN INCLUDED SGA7I_English_current_buildable_source_20260823/components/133_expose_IX_section_2_2.tex
% Expose IX, section 2.2--2.2.9.1.
% French transcription lines 443--599; authority physical pp. 341--348
% (source folios 329--336).  Physical pp. 342--348 were checked in four
% overlapping 1100-dpi bands; the opening on p. 341 was checked with the
% preceding component.  The printed cross-references 1.5.5 and 6.2.5.1 are
% transparently corrected below to 2.1.5 and 2.2.5.1 respectively.

\medskip
\noindent
\SGARefTarget{sga7.sga7i.eix.n2-2}{2.2}. We propose to make more precise the relations between
$T_\ell(A_K)$ and $T_\ell(A_o)=T_\ell(A_o^o)$ by bringing
$T_\ell(A^o)$ into play.  First recall the following.

\medskip
\noindent
\underline{Lemma} \SGARefTarget{sga7.sga7i.eix.n2-2-1}{2.2.1}. \underline{Let $m>0$ be an integer.  Then the
following conditions are equivalent:}

\smallskip
\noindent
(i) $m\,\mathrm{id}_{A^o}$ \underline{is flat}.

\noindent
(ii) $m\,\mathrm{id}_{A^o}$ \underline{is surjective}.

\noindent
(iii) $m\,\mathrm{id}_{A^o}$ \underline{is quasi-finite}.

\noindent
(iv) $m$ \underline{is prime to the characteristic of $k$}, or $A_o^o$
\underline{has unipotent rank zero} (cf. \SGARefLink{sga7.sga7i.eix.n2-1-5}{2.1.5}).

\smallskip
\noindent
\underline{Moreover, if these conditions hold for $m$, they hold for
$m^\nu$ for every $\nu>0$; the groups}
\[
 {}_{m^\nu}A^o\underset{\mathrm{dfn}}{=}
 \ker(m^\nu\,\mathrm{id}_{A^o})
\]
\underline{are quasi-finite, separated, and flat over $S$, and, as
$\nu$ varies, they form an inverse system of group schemes whose
transition morphisms}
\[
 x\longmapsto m^{\nu'-\nu}x:
 {}_{m^{\nu'}}A^o\longrightarrow{}_{m^\nu}A^o
 \qquad(\nu'\geq\nu)
\]
\underline{are faithfully flat}.

Indeed, the fiberwise criterion for flatness (EGA~IV~11.3.5) shows that
(i), like (ii) and (iii), is a condition on the fibers; these three
conditions are equivalent by SGA~3 $\mathrm{VI}_A$~5.6, since the
fibers in question are smooth and connected.  In the form (ii), a
classical theorem of Weil [5, 15] shows that the conditions hold on the
generic fiber, which is an abelian scheme; it is therefore enough to
state them on the special fiber.  Their equivalence with (iv) is then
also well known from the structure theory.  In any of these forms,
stability of the condition under $m\mapsto m^\nu$ is immediate.  It
follows that the ${}_{m^\nu}A^o$ are quasi-finite flat group schemes
over $S$, and they are plainly separated over $S$ because $A^o$ is.
Finally, faithful flatness of the transition morphisms---equivalently,
the fact that they are epimorphisms for the fppf topology---follows
formally from the fact that $m\,\mathrm{id}_{A^o}$ is faithfully flat,
hence an epimorphism for that topology.

\medskip
\noindent
\SGARefTarget{sga7.sga7i.eix.n2-2-2}{2.2.2}. Consequently, we shall take $m$ to be a prime $\ell$ satisfying
the conditions of \SGARefLink{sga7.sga7i.eix.n2-2-1}{2.2.1}.  We then denote by
\begin{equation}\tag{\SGARefTarget{sga7.sga7i.eix.n2-2-2-1}{2.2.2.1}}
 T_\ell(A^o)=({}_{\ell^\nu}A^o)_{\nu>0}
\end{equation}
the inverse system described in \SGARefLink{sga7.sga7i.eix.n2-2-1}{2.2.1}.  Under the two base changes
$\eta\longrightarrow S$ and $s\longrightarrow S$, it recovers the
inverse systems $T_\ell(A_K)$ and
$T_\ell(A_o^o)=T_\ell(A_o)$ considered in \SGARefLink{sga7.sga7i.eix.n1-0}{1.0} and \SGARefLink{sga7.sga7i.eix.n2-1}{2.1} respectively.

\medskip
\noindent
\SGARefTarget{sga7.sga7i.eix.n2-2-3}{2.2.3}. Suppose now that $S$ is henselian.  Every quasi-finite separated
scheme $X$ over $S$ then has a canonical decomposition as a disjoint
union
\begin{equation}\tag{\SGARefTarget{sga7.sga7i.eix.n2-2-3-1}{2.2.3.1}}
 X=X^f\coprod X',
\end{equation}
where $X^f$ is \underline{finite over $S$} and $X'_o=\varnothing$, that
is, $X'$ is confined to its generic fiber $X'_\eta$; thus $X^f$ has the
same special fiber as $X$.  This decomposition is plainly functorial in
$X$.  In particular, if $X$ is a group scheme over $S$, then $X^f$ is
a subgroup scheme.  If $X$ is flat, respectively \'{e}tale, over $S$,
then so is $X^f$.

We apply this to the components ${}_{\ell^\nu}A^o$ of $T_\ell(A^o)$.
This gives a projective subsystem of groups
\begin{equation}\tag{\SGARefTarget{sga7.sga7i.eix.n2-2-3-2}{2.2.3.2}}
 \bigl(({}_{\ell^\nu}A^o)^f\bigr)_{\nu>0}
 \underset{\mathrm{dfn}}{=}T_\ell(A^o)^f
 \subset T_\ell(A^o),
\end{equation}
called the \underline{fixed part} of the Tate module $T_\ell(A^o)$.
By construction, there are canonical isomorphisms, functorial in $A$,
\begin{equation}\tag{\SGARefTarget{sga7.sga7i.eix.n2-2-3-3}{2.2.3.3}}
 T_\ell(A^o)^f\times_S s\xrightarrow{\ \sim\ }
 T_\ell(A^o)\times_Ss\xrightarrow{\ \sim\ }T_\ell(A_o^o).
\end{equation}
On the other hand, the generic fiber of $T_\ell(A^o)^f$ defines a
canonical sub-pro-group of $T_\ell(A_K)$, again called the
``\underline{fixed part}'' of this Tate module:
\begin{equation}\tag{\SGARefTarget{sga7.sga7i.eix.n2-2-3-4}{2.2.3.4}}
 T_\ell(A_K)^f\subset T_\ell(A_K).
\end{equation}

\medskip
\noindent
\SGARefTarget{sga7.sga7i.eix.n2-2-4}{2.2.4}. When $\ell\ne p=\operatorname{char}k$, the fixed part
$T_\ell(A^o)^f$ may be interpreted as a twisted constant $\ell$-adic
sheaf on $S$ (the $({}_{\ell^\nu}A^o)^f$ being finite \'{e}tale over
$S$).  Since the functor
\[
 X\longmapsto X_o=X\times_Ss
\]
from finite \'{e}tale schemes over $S$ to finite \'{e}tale schemes
over $k$ is an equivalence of categories when $S$ is henselian, there
is an analogous equivalence for the categories of twisted constant
$\ell$-adic sheaves over $S$ and over $k$.  Hence (\SGARefLink{sga7.sga7i.eix.n2-2-3-3}{2.2.3.3}), which
gives the special fiber of $T_\ell(A^o)^f$, characterizes this fixed
part up to unique isomorphism.  On the other hand, as an $\ell$-adic
subsheaf of $T_\ell(A^o)$ it is plainly characterized by its generic
fiber (\SGARefLink{sga7.sga7i.eix.n2-2-3-4}{2.2.3.4}), or equivalently by the corresponding
$\mathbb Z_\ell$-submodule, stable under
$\pi=\operatorname{Gal}(\overline K/K)$:
\begin{equation}\tag{\SGARefTarget{sga7.sga7i.eix.n2-2-4-1}{2.2.4.1}}
 T_\ell(A_K)^f(\overline K)
 \underset{\mathrm{dfn}}{=}T_\ell(A(\overline K))^f
 \subset T_\ell(A(\overline K)).
\end{equation}
In fact, one has the following result.

\medskip
\noindent
\underline{Proposition} \SGARefTarget{sga7.sga7i.eix.n2-2-5}{2.2.5}. \underline{Suppose, as above, that
$\ell\ne p$ and that $S$ is henselian.  Then the submodule (\SGARefLink{sga7.sga7i.eix.n2-2-4-1}{2.2.4.1})
of $T_\ell(A(\overline K))=T_\ell(A_K)(\overline K)$ corresponding to
the fixed part $T_\ell(A_K)^f$ in (\SGARefLink{sga7.sga7i.eix.n2-2-3-4}{2.2.3.4}) is the submodule
$T_\ell(A(\overline K))^I$ of invariants under the inertia group
$I\subset\pi$.}

This follows, by passage to the inverse limit, from the analogous
relation valid for every integer $m>0$ prime to
$p=\operatorname{char}k$:
\begin{equation}\tag{\SGARefTarget{sga7.sga7i.eix.n2-2-5-1}{2.2.5.1}}
 ({}_mA)^f(\overline K)=({}_mA(\overline K))^I.
\end{equation}
To prove this relation, introduce the strict localization $\widetilde S$
of $S$, whose fraction field $\widetilde K$ is the maximal unramified
subextension of $\overline K/K$.  Since $({}_mA)^f$ is finite
\'{e}tale over $S$, its points with values in $\overline K$ already
have values in $\widetilde K$.  Thus the left-hand side of (\SGARefLink{sga7.sga7i.eix.n2-2-5-1}{2.2.5.1})
is
$\operatorname{Hom}_S(\widetilde S,({}_mA)^f)$, whereas the right-hand
side is, by definition of $\widetilde K$, equal to ${}_mA(\widetilde K)$.
By the definition of $({}_mA)^f$, and since $\widetilde S$ is an
inverse limit of finite \'{e}tale $S$-schemes $S_i$, the former is also
equal to
\[
 \operatorname{Hom}_S(\widetilde S,{}_mA)
 ={}_m\operatorname{Hom}_S(\widetilde S,A).
\]
By the defining property of the N\'{e}ron model,
$\operatorname{Hom}_S(\widetilde S,A)=A(\widetilde K)$, which proves
(\SGARefLink{sga7.sga7i.eix.n2-2-5-1}{2.2.5.1}).

To deduce \SGARefLink{sga7.sga7i.eix.n2-2-5}{2.2.5}, observe that (\SGARefLink{sga7.sga7i.eix.n2-2-5-1}{2.2.5.1}) gives
\begin{equation}\tag{*}
 T_\ell(A(\overline K))^I
 =\varprojlim_\nu({}_{\ell^\nu}A)^f(\overline K)
 \supset\varprojlim_\nu({}_{\ell^\nu}A^o)^f(\overline K).
\end{equation}
If $(x_\nu)$ belongs to the rightmost term, then for every
$\nu\geq0$, $x_\nu$ is a section of $A$ that is infinitely divisible
by $\ell$.  Its value at $s$ is therefore an element of $A_o(k)$
infinitely divisible by $\ell$, and consequently lies in $A_o^o(k)$.
This proves equality of the last two terms of (*) and completes the
proof of \SGARefLink{sga7.sga7i.eix.n2-2-5}{2.2.5}.

\medskip
\noindent
\underline{Proposition} \SGARefTarget{sga7.sga7i.eix.n2-2-6}{2.2.6}. \underline{Let $C_K$ be a second
abelian scheme over $K$, with N\'{e}ron model $C$, and consider a
homomorphism $u_K:C_K\longrightarrow A_K$, which therefore extends to
a homomorphism $u:C\longrightarrow A$.  Suppose that $u_K$ is an
isogeny.  Then, for every prime $\ell$,
$T_\ell(u_K):T_\ell(C_K)\longrightarrow T_\ell(A_K)$ induces an
isogeny of fixed parts}
\begin{equation}\tag{\SGARefTarget{sga7.sga7i.eix.n2-2-6-1}{2.2.6.1}}
 T_\ell(C_K)^f\longrightarrow T_\ell(A_K)^f,
\end{equation}
\underline{and likewise the map induced by
$u_o:C_o\longrightarrow A_o$}
\begin{equation}\tag{\SGARefTarget{sga7.sga7i.eix.n2-2-6-2}{2.2.6.2}}
 T_\ell(u_o^o):T_\ell(C_o^o)\longrightarrow T_\ell(A_o^o)
\end{equation}
\underline{is an isogeny.  Finally, $u_o:C_o\longrightarrow A_o$
induces an isogeny from the largest subtorus of $C_o$ to the largest
subtorus $T_o$ of $A_o$; and over an extension $k'/k$ on which the
``abelian parts'' of $C_o$ and $A_o$ are defined (for example, over
the perfect closure of $k$), $u_o$ induces an isogeny between the
abelian parts of $C_{ok'}$ and $A_{ok'}$.}

This statement follows formally from the fact that
$T_\ell(A_K)^f$, $T_\ell(A_o^o)$, the largest subtorus of $A_o$, and
the ``abelian part'' of $A_{ok'}$ are additive functors of the varying
$A_K$, and from the fact that $u_K$ is an isogeny precisely in the
sense that there is a morphism $v_K:A_K\longrightarrow C_K$ such that
\[
 u_Kv_K=n\,\mathrm{id},\qquad v_Ku_K=n\,\mathrm{id}
\]
for some integer $n>0$.

\medskip
\noindent
\underline{Corollary} \SGARefTarget{sga7.sga7i.eix.n2-2-7}{2.2.7}. \underline{If $A_K$ and $C_K$ are
isogenous abelian schemes over $K$, then the corresponding invariants
$\mu,\lambda,\alpha$ in (\SGARefLink{sga7.sga7i.eix.n2-1-11}{2.1.11}) are the same for $A_K$ and $C_K$.
In particular, they are the same for $A_K$ and its dual $A'_K$.}

\medskip
\noindent
\underline{Remarks} \SGARefTarget{sga7.sga7i.eix.n2-2-8}{2.2.8}. In the two preceding statements, \SGARefLink{sga7.sga7i.eix.n2-2-6}{2.2.6} and
\SGARefLink{sga7.sga7i.eix.n2-2-7}{2.2.7}, we assumed that $S$ is henselian.  It is immediate, however,
that \SGARefLink{sga7.sga7i.eix.n2-2-7}{2.2.7}, as well as the assertions of \SGARefLink{sga7.sga7i.eix.n2-2-6}{2.2.6} concerning the
properties of the induced homomorphism
$u_o:C_o\longrightarrow A_o$, holds without this hypothesis.  Indeed,
one reduces to the case already treated by the base change
$S^h\longrightarrow S$, where $S^h$ is the henselization of $S$, using
the fact---immediate from the definitions---that formation of N\'{e}ron
models commutes with this base change.

\medskip
\noindent
\underline{Corollary} \SGARefTarget{sga7.sga7i.eix.n2-2-9}{2.2.9}. (\underline{Good reduction.})
\underline{Let $S$ be a trait, let $A_K$ be an abelian scheme over its
fraction field $K$, let $A$ be the N\'{e}ron model of $A_K$, and let
$\ell$ be a prime different from $\operatorname{char}k$.  The following
conditions are equivalent.}

\smallskip
\noindent
(i) \underline{There exists an abelian scheme over $S$ whose generic
fiber is isomorphic to $A_K$.}

\noindent
(ii) $A$ \underline{is an abelian scheme}.

\noindent
(ii$^o$) $A^o$ \underline{is an abelian scheme}.

\noindent
(iii) $A_o$ \underline{is an abelian scheme}.

\noindent
(iii$^o$) $A_o^o$ \underline{is an abelian scheme}, \underline{that is,
with the notation of (\SGARefLink{sga7.sga7i.eix.n2-1-11}{2.1.11}), $\lambda=\mu=0$}.

\noindent
(iv) $A$ \underline{is proper over $S$}.

\noindent
(iv$^o$) $A^o$ \underline{is proper over $S$}.

\noindent
(v) $T_\ell(A_K)$ \underline{is ``unramified over $S$''}, that is,
\underline{it extends to a twisted constant $\ell$-adic sheaf over
$S$}, or equivalently \underline{the inertia group $I$ acts trivially
on $T_\ell(A_K(\overline K))$}.

\noindent
(v bis) \underline{For every integer $\nu\geq0$,
${}_{\ell^\nu}A_K$ is ``unramified over $S$''}, that is,
\underline{it extends to an \'{e}tale covering of $S$}, or equivalently
\underline{the inertia group $I$ acts trivially on
${}_{\ell^\nu}A_K(\overline K)$}.

\smallskip
\noindent
\underline{Finally, an abelian scheme over $S$ as in (i) is a N\'{e}ron
model of $A_K$, and in particular is determined up to unique
isomorphism.}

\medskip
\noindent
\SGARefTarget{sga7.sga7i.eix.n2-2-9-1}{2.2.9.1}. We say that $A_K$ has \underline{good reduction} over $S$ if
it satisfies condition (i) above.  Criterion (v), or (v bis), is known
as the \underline{N\'{e}ron--Ogg--Shafarevich criterion} for good
reduction.

Let us prove \SGARefLink{sga7.sga7i.eix.n2-2-9}{2.2.9}.  It is well known (cf., for example, [5]) that an
abelian scheme over $S$ satisfies the universal property of a
N\'{e}ron model; this reduces to the fact that a rational section of an
abelian scheme over a regular base is defined everywhere.  This proves
the equivalence
$(i)\Longleftrightarrow(ii)\Longleftrightarrow(ii^o)$, as well as the
last assertion of \SGARefLink{sga7.sga7i.eix.n2-2-9}{2.2.9}.  These conditions plainly imply (iii), (iv),
and (iv$^o$), each of which implies (iii$^o$).  Let us prove that
(iii$^o$)$\Longrightarrow$(ii$^o$), which will establish the
equivalence of (i) through (iv$^o$).  This follows from the fact that a
smooth group scheme over $S$ whose fibers are abelian schemes is proper
over $S$ (and hence is an abelian scheme), by EGA~IV~15.7.10.

On the other hand, in view of \SGARefLink{sga7.sga7i.eix.n2-2-5}{2.2.5} and the fact that
\[
 T_\ell(A_K(\overline K))/T_\ell(A_K(\overline K))^I
\]
is plainly a torsion-free $\mathbb Z_\ell$-module, condition (v) says
that $T_\ell(A_K(\overline K))$ and its fixed part have the same rank.
By (\SGARefLink{sga7.sga7i.eix.n2-1-11}{2.1.11}) and (\SGARefLink{sga7.sga7i.eix.n2-2-3-3}{2.2.3.3}), this is the condition
\[
 \mu+2\alpha=2(\lambda+\mu+\alpha),
\]
that is, $2\lambda+\mu=0$, or $\lambda=\mu=0$, which is precisely
(iii$^o$).  Finally, the equivalence of (v) and (v bis) is clear from
the definitions.  This completes the proof of \SGARefLink{sga7.sga7i.eix.n2-2-9}{2.2.9}.

% END INCLUDED SGA7I_English_current_buildable_source_20260823/components/133_expose_IX_section_2_2.tex

% BEGIN INCLUDED SGA7I_English_current_buildable_source_20260823/components/134_expose_IX_sections_2_3_and_2_4.tex
% Expose IX, sections 2.3 and 2.4.
% French transcription lines 601--685; authority physical pp. 348--352
% (source folios 336--340).  Physical pp. 349--352 were checked in four
% overlapping 1100-dpi bands; the opening on p. 348 was checked with the
% preceding component.  The printed reference (1.6.5) is transparently
% corrected to (2.2.5), the missing prime in the second factor of (2.4.5)
% is restored, and the impossible printed reference (1.10) is replaced by
% an unnumbered verbal reference to the displayed inclusion (2.4.7).

\medskip
\noindent
\SGARefTarget{sga7.sga7i.eix.n2-3}{2.3}. Suppose once again that $S$ is henselian and that
$\ell\ne p=\operatorname{char}k$.  We now use the inclusion
\begin{equation}\tag{\SGARefTarget{sga7.sga7i.eix.n2-3-1}{2.3.1}}
 T_\ell(T_o)\hookrightarrow T_\ell(A_o^o)
\end{equation}
of (\SGARefLink{sga7.sga7i.eix.n2-1-6}{2.1.6}), together with the equivalence of categories recalled in
\SGARefLink{sga7.sga7i.eix.n2-2-4}{2.2.4}.  We obtain a canonical twisted constant $\ell$-adic subsheaf
of $T_\ell(A^o)^f$, called the \underline{toric part of}
$T_\ell(A^o)$,
\begin{equation}\tag{\SGARefTarget{sga7.sga7i.eix.n2-3-2}{2.3.2}}
 T_\ell(A^o)^t\hookrightarrow T_\ell(A^o)^f
 \hookrightarrow T_\ell(A^o),
\end{equation}
characterized by the condition that its special fiber is the inclusion
(\SGARefLink{sga7.sga7i.eix.n2-3-1}{2.3.1}):
\begin{equation}\tag{\SGARefTarget{sga7.sga7i.eix.n2-3-3}{2.3.3}}
 T_\ell(A^o)^t\times_Ss=T_\ell(T_o)
 \hookrightarrow T_\ell(A_o^o).
\end{equation}
Later we shall give an interpretation of the toric part of
$T_\ell(A^o)$ as the Tate module $T_\ell(T)$ of a torus $T$ over $S$
lifting the torus $T_o$ over $k$.  That interpretation will allow the
definition of the toric part to be extended to the case $\ell=p$.
For the moment, and to fix ideas, we confine ourselves to the case
$\ell\ne p$.  In this case the toric part of $T_\ell(A^o)$ may be
described by the corresponding $\mathbb Z_\ell$-submodule of
$T_\ell(A(\overline K))$:
\begin{equation}\tag{\SGARefTarget{sga7.sga7i.eix.n2-3-4}{2.3.4}}
 T_\ell(A(\overline K))^t\hookrightarrow
 T_\ell(A(\overline K))^f\hookrightarrow
 T_\ell(A(\overline K)).
\end{equation}
It is stable under the action of $\pi$; its first term is defined as
$T_\ell(A^o)^t(\overline K)$ and is, as one would expect, again called
the \underline{toric submodule} of the Tate module
$T_\ell(A(\overline K))$.  Thus this Tate module carries the canonical
three-step filtration (\SGARefLink{sga7.sga7i.eix.n2-3-4}{2.3.4}).  Applying the same considerations to
the dual abelian scheme $A'_K$, we obtain the analogous three-step
filtration
\begin{equation}\tag{\SGARefTarget{sga7.sga7i.eix.n2-3-5}{2.3.5}}
 T_\ell(A'(\overline K))^t\hookrightarrow
 T_\ell(A'(\overline K))^f\hookrightarrow
 T_\ell(A'(\overline K)).
\end{equation}
The middle terms in (\SGARefLink{sga7.sga7i.eix.n2-3-4}{2.3.4}) and (\SGARefLink{sga7.sga7i.eix.n2-3-5}{2.3.5}), the ``fixed parts,'' are
described in terms of the $\pi$-module structures of the Tate modules
of $A(\overline K)$ and $A'(\overline K)$ as their submodules of
invariants under the inertia group $I$ (\SGARefLink{sga7.sga7i.eix.n2-2-5}{2.2.5}).  The toric parts are
then completely described in terms of the fixed parts and the pairing
$\varphi$ of (\SGARefLink{sga7.sga7i.eix.n1-0-2}{1.0.2}), by the following theorem (a special case of which
already appears in [11]).

\medskip
\noindent
\underline{Theorem} \SGARefTarget{sga7.sga7i.eix.n2-4}{2.4}. (\underline{Orthogonality theorem}.)
\underline{Let $S$ be a henselian trait, let $A_K$ be an abelian scheme
over its fraction field $K$, let $A'_K$ be the dual abelian scheme, and
let $\ell$ be a prime different from the residual characteristic of
$S$.  The Tate modules $T_\ell(A(\overline K))$ and
$T_\ell(A'(\overline K))$ then carry the canonical filtrations (\SGARefLink{sga7.sga7i.eix.n2-3-4}{2.3.4})
and (\SGARefLink{sga7.sga7i.eix.n2-3-5}{2.3.5}), respectively, by the ``fixed part'' and the ``toric
part.''  The toric part $T_\ell(A(\overline K))^t$ is the intersection
of the fixed part $T_\ell(A(\overline K))^f$ with the orthogonal of the
fixed part
$T_\ell(A'(\overline K))^f=T_\ell(A'(\overline K))^I$ with respect to
the canonical pairing $\varphi$ of (\SGARefLink{sga7.sga7i.eix.n1-0-3}{1.0.3}).  Symmetrically, the toric
part $T_\ell(A'(\overline K))^t$ is the intersection of
$T_\ell(A'(\overline K))^f$ with the orthogonal of the fixed part
$T_\ell(A(\overline K))^f=T_\ell(A(\overline K))^I$.  (Here the
superscript $I$ denotes the submodule of invariants under the action of
the inertia group $I$.)}

\smallskip
\noindent
\underline{Proof.}\quad a) Let us prove that
$T_\ell(A(\overline K))^t$ is orthogonal to
$T_\ell(A'(\overline K))^f$.  To do so, we interpret the restriction of
$\varphi$ to the fixed parts of the Tate modules using the canonical
biextension $W^o=W$ of (\SGARefLink{sga7.sga7i.eix.n1-4-4}{1.4.4}).  This is the canonical extension of the
biextension $W_K$ of $(A_K,A'_K)$ by $\mathbb G_{mK}$ to a biextension
of $(A^o,A'{}^o)$ by $\mathbb G_{mS}$.  By VIII~\SGARefLink{sga7.sga7i.eviii.n2-2}{2.2}, $W$ gives rise to
a pairing
\begin{equation}\tag{\SGARefTarget{sga7.sga7i.eix.n2-4-1}{2.4.1}}
 \varphi_W:T_\ell(A^o)\times T_\ell(A'{}^o)
 \longrightarrow T_\ell(\mathbb G_{mS}),
\end{equation}
from which the pairing $\varphi$ of (\SGARefLink{sga7.sga7i.eix.n1-0-3}{1.0.3}) is obtained by passing to
$\overline K$-valued points.  It follows that the restriction of
$\varphi$ to the fixed parts considered in \SGARefLink{sga7.sga7i.eix.n2-4}{2.4} is obtained, by passing
to $\overline K$-valued points, from the pairing
\begin{equation}\tag{\SGARefTarget{sga7.sga7i.eix.n2-4-2}{2.4.2}}
 T_\ell(A^o)^f\times T_\ell(A'{}^o)^f
 \longrightarrow T_\ell(\mathbb G_{mS}),
\end{equation}
the restriction of (\SGARefLink{sga7.sga7i.eix.n2-4-1}{2.4.1}).  Since the two factors on the left of
(\SGARefLink{sga7.sga7i.eix.n2-4-2}{2.4.2}) are now twisted constant $\ell$-adic sheaves, this pairing may
be studied on the special fibers.  The fibers of the two factors are
precisely the Tate modules $T_\ell(A_o^o)$ and
$T_\ell(A'{}_o^o)$ (\SGARefLink{sga7.sga7i.eix.n2-2-3-3}{2.2.3.3}), and the pairing induced on special fibers
by (\SGARefLink{sga7.sga7i.eix.n2-4-2}{2.4.2}) is
\begin{equation}\tag{\SGARefTarget{sga7.sga7i.eix.n2-4-3}{2.4.3}}
 T_\ell(A_o^o)\times T_\ell(A'{}_o^o)
 \longrightarrow T_\ell(\mathbb G_{mk}),
\end{equation}
obtained from the biextension $W_o$ of $(A_o^o,A'{}_o^o)$ by
$\mathbb G_{mk}$, the restriction of $W$ to the special fibers.

Recall from \SGARefLink{sga7.sga7i.eix.n2-3}{2.3} that the toric part of
$T_\ell(A(\overline K))$ is obtained by passing to
$\overline K$-valued points in the twisted constant $\ell$-adic
subsheaf $T_\ell(A^o)^t$ of $T_\ell(A^o)^f$.  Thus the orthogonality
relation to be proved says that the pairing induced by (\SGARefLink{sga7.sga7i.eix.n2-4-2}{2.4.2}),
\begin{equation}\tag{\SGARefTarget{sga7.sga7i.eix.n2-4-4}{2.4.4}}
 T_\ell(A^o)^t\times T_\ell(A'{}^o)^f
 \longrightarrow T_\ell(\mathbb G_{mS}),
\end{equation}
is identically zero.  It is enough to verify this for the corresponding
pairing on the special fibers, which by the definition of the toric
part is
\begin{equation}\tag{\SGARefTarget{sga7.sga7i.eix.n2-4-5}{2.4.5}}
 T_\ell(T_o)\times T_\ell(A'{}_o^o)
 \longrightarrow T_\ell(\mathbb G_{mk}),
\end{equation}
induced by (\SGARefLink{sga7.sga7i.eix.n2-4-3}{2.4.3}).  By VIII~\SGARefLink{sga7.sga7i.eviii.n4-10}{4.10} this pairing is necessarily zero,
which proves the asserted orthogonality.

\smallskip
\noindent
b) Let us prove that the inclusion
\begin{equation}\tag{\SGARefTarget{sga7.sga7i.eix.n2-4-6}{2.4.6}}
 T_\ell(A(\overline K))^t\subset
 T_\ell(A(\overline K))^f\cap
 \bigl(T_\ell(A'(\overline K))^f\bigr)^\perp
\end{equation}
is an equality, where $\perp$ denotes the orthogonal.  Equivalently, in
the pairing (\SGARefLink{sga7.sga7i.eix.n2-4-2}{2.4.2}), the annihilator of the second factor---which, as
we have just seen, contains the toric part $T_\ell(A^o)^t$---is equal
to that toric part.  Since these are again relations between twisted
constant $\ell$-adic sheaves, it is enough to check the assertion on
the special fibers: namely, that the inclusion
\begin{equation}\tag{\SGARefTarget{sga7.sga7i.eix.n2-4-7}{2.4.7}}
 T_\ell(T_o)\subset T_\ell(A'{}_o^o)^\perp
\end{equation}
is an equality.  In other words, we must verify that the biextension
$W_o$ of $(A_o^o,A'{}_o^o)$ is \underline{nondegenerate on the left}
(VIII~\SGARefLink{sga7.sga7i.eviii.n4-11}{4.11}).  But this is precisely what we already observed in \SGARefLink{sga7.sga7i.eix.n1-5}{1.5},
as a consequence of a result from Raynaud's thesis.

% END INCLUDED SGA7I_English_current_buildable_source_20260823/components/134_expose_IX_sections_2_3_and_2_4.tex

% BEGIN INCLUDED SGA7I_English_current_buildable_source_20260823/components/135_expose_IX_section_2_5_and_remarks_2_6.tex
% Expose IX, section 2.5 and Remarks 2.6.
% French transcription lines 687--757; authority physical pp. 352--354
% (source folios 340--342), with the section-3 boundary checked on p. 355.
% Physical pp. 353--354 were checked in four overlapping 1100-dpi bands;
% the opening on p. 352 was checked with the preceding component.  The
% six-arrow filtration diagram (2.5.4) was additionally checked directly
% in five overlapping 5000-dpi crops, including hook sides, label sides,
% arrow directions, and the terminal comma.

\medskip
\noindent
\SGARefTarget{sga7.sga7i.eix.n2-5}{2.5}. The orthogonality theorem has an equivalent formulation using a
polarization $\xi$ of $A_K$ and the alternating bilinear form (\SGARefLink{sga7.sga7i.eix.n1-0-8}{1.0.8}),
compatible with the operations of $\pi$,
\begin{equation}\tag{\SGARefTarget{sga7.sga7i.eix.n2-5-1}{2.5.1}}
 \varphi_\xi:T_\ell(A(\overline K))\times T_\ell(A(\overline K))
 \longrightarrow T_\ell(\overline K^*),
\end{equation}
obtained from the form $\varphi$ of (\SGARefLink{sga7.sga7i.eix.n1-0-3}{1.0.3}) by transporting structure in
the second argument along the isogeny of $\mathbb Z_\ell$-modules
\[
 T_\ell(A(\overline K))\longrightarrow T_\ell(A'(\overline K))
\]
induced by the isogeny $\psi_\xi:A_K\longrightarrow A'_K$ defined by
$\xi$.  The equality in (\SGARefLink{sga7.sga7i.eix.n2-4-6}{2.4.6}) then takes the simple form
\begin{equation}\tag{\SGARefTarget{sga7.sga7i.eix.n2-5-2}{2.5.2}}
 W=V\cap V^\perp,
\end{equation}
where we have set
\begin{equation}\tag{\SGARefTarget{sga7.sga7i.eix.n2-5-3}{2.5.3}}
 \left\{
 \begin{array}{l}
 U=T_\ell(A(\overline K)),\qquad
 V=U^I=\text{the ``fixed part'' of }T_\ell(A(\overline K)),\\[2mm]
 \hspace{4em}W=\text{the toric part of }T_\ell(A(\overline K)),
 \end{array}
 \right.
\end{equation}
and $\perp$ denotes the orthogonal with respect to the polarization form
(\SGARefLink{sga7.sga7i.eix.n2-5-1}{2.5.1}).  These modules give the following diagram of inclusions:
\begin{equation}\tag{\SGARefTarget{sga7.sga7i.eix.n2-5-4}{2.5.4}}
\begin{tikzcd}[column sep=large,row sep=large]
 && V\arrow[dl,hook',"2\lambda"'] && \\
 U & W^\perp=V+V^\perp\arrow[l,hook',"\mu"']
 && W=V\cap V^\perp
    \arrow[ul,hook',"2\alpha"']
    \arrow[dl,hook,"2\lambda"]
 & 0\arrow[l,hook',"\mu"'] \\
 && V^\perp\arrow[ul,hook,"2\alpha"] &&
\end{tikzcd}
\qquad ,
\end{equation}
where the integers $2\alpha$, $2\lambda$, and $\mu$ are the ranks of
the quotients of consecutive modules in the diagram, with the notation
of (\SGARefLink{sga7.sga7i.eix.n2-1-9}{2.1.9}).  (That $W$ has rank $\mu$ is the definition of $\mu$; that
$V/W$ has rank $2\alpha$ follows from the second relation in (\SGARefLink{sga7.sga7i.eix.n2-1-11}{2.1.11})
together with the exact sequence (\SGARefLink{sga7.sga7i.eix.n2-1-6}{2.1.6}).  Hence, by orthogonality,
$\operatorname{rank}(U/W^\perp)=\mu$ and
$\operatorname{rank}(W^\perp/V^\perp)=2\alpha$.  Finally, the formulas
% The print has 2\lambda here.  Diagram (2.5.4), the preceding rank
% computation, and Proposition 3.5 show that the quotient V/W has rank
% 2\alpha; the complementary quotient V^\perp/W has rank 2\lambda.
\begin{equation}\tag{\SGARefTarget{sga7.sga7i.eix.n2-5-5}{2.5.5}}
 \operatorname{rank}\bigl(V/(V\cap V^\perp)\bigr)
 =\operatorname{rank}(V/W)=2\alpha
\end{equation}
and the dual formula follow from the preceding ones, in view of
(\SGARefLink{sga7.sga7i.eix.n2-1-10}{2.1.10}) and $\operatorname{rank}U=n$.)

\medskip
\noindent
\underline{Remarks} \SGARefTarget{sga7.sga7i.eix.n2-6}{2.6}.\quad a) Here we obtained the orthogonality
theorem \SGARefLink{sga7.sga7i.eix.n2-4}{2.4} as an easy consequence of the theory of biextensions
developed in the preceding sections.  The \underline{inclusion}
(\SGARefLink{sga7.sga7i.eix.n2-4-6}{2.4.6}) can be proved without using the theory of biextensions, by an
arithmetic argument in the style of the proof of the ``monodromy
theorem'' given in Expos\'e I.  One reduces to a situation in which $S$
has finite residue field, where the orthogonality result follows from a
simple consideration of \underline{weights} for the Frobenius operation
(using A.~Weil's result on the eigenvalues of that Frobenius).  There is
no need to invoke resolution of singularities here: one may state the
desired result for group schemes $A$ over schemes $S$ under sufficiently
general conditions.  Unfortunately, it does not appear possible to
obtain by this method equality in the preceding inclusion.  That
equality will apparently be used in an essential way in the proof of
the ``semistable reduction theorem'' in the next section.

\smallskip
\noindent
b) The orthogonality theorem \SGARefLink{sga7.sga7i.eix.n2-4}{2.4} allows the most important properties
of the N\'{e}ron model of an abelian variety $A_K$ to be expressed in
terms of the representation of the inertia group $I$ on the Tate module
$T_\ell(A(\overline K))$, where $\ell$ is a fixed prime different from
the residual characteristic.  (The case $\ell=p$ will be considered
separately later, in section 5.)  As a first statement of this kind,
though one that does not use the orthogonality theorem, recall the
good-reduction criterion \SGARefLink{sga7.sga7i.eix.n2-2-9}{2.2.9}.  We shall give many further examples
below, beginning in the next section (\SGARefLink{sga7.sga7i.eix.n3-5}{3.5}).

% END INCLUDED SGA7I_English_current_buildable_source_20260823/components/135_expose_IX_section_2_5_and_remarks_2_6.tex

% BEGIN INCLUDED SGA7I_English_current_buildable_source_20260823/components/136_expose_IX_section_3_1.tex
% Expose IX, section 3.1.
% French transcription lines 763--839; authority physical pp. 355--359
% (source folios 343--347), stopping before Proposition 3.2.
% Physical pp. 355--358 and the opening/boundary on p. 359 were checked
% directly in four overlapping 1100-dpi bands per page.  The text and the
% single hooked-arrow display were unambiguous at that resolution; no
% higher-detail escalation was needed.  The inherited French transcription's
% impossible ``t\in T'' in Lemma 3.1.3 is corrected to the scan-confirmed
% ``t\in S'' in the English text.

\medskip
\noindent
3. \underline{The case of semistable reduction.  The semistable reduction
theorem.}

We first give an auxiliary result in the style of SGA 3.

\medskip
\noindent
\underline{Proposition} \SGARefTarget{sga7.sga7i.eix.n3-1}{3.1}. \underline{Let $S$ be a scheme, let $G$ be a
commutative $S$-group scheme that is flat, separated, and of finite
presentation over $S$, let $H$ be a quasi-separated group scheme locally of
finite type over $S$, let $u:G\longrightarrow H$ be a homomorphism of
$S$-group schemes, and put $K=\ker u$.  For every $s\in S$, let $\bar s$
denote the spectrum of an algebraic closure of $k(s)$, and denote by
$\mu(s)$ (respectively, $\alpha(s)$) the reductive rank (respectively, the
abelian rank) of $K_s$, that is, the dimension of the largest subtorus
(respectively, of the largest abelian-variety quotient) of
$(K_{\bar s})^o_{\mathrm{red}}$.  Let $s\in S$, let $t\in S$ be a
generalization of $s$, and let $p$ be the characteristic exponent of
$k(s)$.}

\smallskip
\noindent
a) \underline{For every integer $n>0$ prime to $p$, ${}_nK$ is etale, of
finite type, and separated over $S$ in a neighborhood of $s$, and one has
the divisibility relation}
\begin{equation}\tag{\SGARefTarget{sga7.sga7i.eix.n3-1-1}{3.1.1}}
 \operatorname{rank}({}_nK_s)\mid \operatorname{rank}({}_nK_t).
\end{equation}

\noindent
b) \underline{One has}
\begin{equation}\tag{\SGARefTarget{sga7.sga7i.eix.n3-1-2}{3.1.2}}
 \mu(s)+2\alpha(s)\leq \mu(t)+2\alpha(t).
\end{equation}

\noindent
c) \underline{If $(K_{\bar t})^o_{\mathrm{red}}$ is unipotent, then so is
$(K_{\bar s})^o_{\mathrm{red}}$.}

\smallskip
\noindent
d) \underline{Suppose that $G_s$ has unipotent rank zero, that $H$ is smooth
over $S$, and that $u_t$ is an isogeny.  Then $u$ is flat and quasi-finite
over an open neighborhood $U$ of $s$; a fortiori, $K|_U$ is flat and
quasi-finite.}

\smallskip
\noindent
e) \underline{Under the hypotheses of d), suppose in addition that $u_t$ is
a monomorphism (and hence an isomorphism from $G_t$ onto an open subgroup of
$H_t$).  Then $u_s$ is an isomorphism from $G_s$ onto an open subgroup of
$H_s$.  If, moreover, $S$ is irreducible with generic point $t$, then there
is an open neighborhood $U$ of $s$ such that $u$ induces an isomorphism from
$G|_U$ onto an open subgroup of $H|_U$.}

\medskip
\noindent
\underline{Proof.}

\smallskip
\noindent
a) We may suppose that $n$ is invertible on $S$.  Then $n\,\mathrm{id}_G$
is etale on every fibre (SGA 3, XV, 1.3); since $G$ is flat and of finite
presentation over $S$, the fibrewise criterion for etaleness (EGA IV,
17.8.2) shows that $n\,\mathrm{id}_G$ is etale.  Its kernel ${}_nG$ is
therefore etale; it is also separated and of finite presentation over $S$,
since $G$ is.  Similarly, ${}_nH$ is locally of finite type,
quasi-separated, and formally unramified over $S$,
\footnote{At least when $H$ is commutative.  In the general case, ${}_nH$
is still unramified over $S$ along its unit section, which is all that is
needed here.}
so its unit section is an open immersion (EGA IV, 17.4.2 b')).  Consequently,
the kernel of
\[
 {}_nu:{}_nG\longrightarrow {}_nH
\]
is an open subscheme of ${}_nG$, hence is etale over $S$; it is
quasi-compact in ${}_nG$ and therefore of finite presentation over $S$, and
it is separated over $S$ because $G$ is.  This kernel is plainly ${}_nK$.

Formula (\SGARefLink{sga7.sga7i.eix.n3-1-1}{3.1.1}) follows from these properties of ${}_nK$.  Indeed, after a
base change we may reduce to the case where $S$ is henselian local with
closed point $s$, and then use the decomposition of ${}_nK$ into
${}_nK^f$ and a scheme with empty special fibre (cf. \SGARefLink{sga7.sga7i.eix.n2-2-3}{2.2.3}).  More precisely,
the same argument gives a monomorphism on geometric fibres
\[
 {}_nK(\bar s)\lhook\joinrel\longrightarrow {}_nK(\bar t).
\]

\smallskip
\noindent
b) Choose an integer $n>0$ prime to the characteristic exponent of $k(s)$
and to the ranks of the quotients of the geometric fibres of $K$ at $s$ and
$t$ by their reduced identity components.  The two sides of (\SGARefLink{sga7.sga7i.eix.n3-1-1}{3.1.1}) are
then obtained by raising $n$, respectively, to the powers whose exponents
are the two sides of (\SGARefLink{sga7.sga7i.eix.n3-1-2}{3.1.2}), as follows from the known structure of
commutative algebraic groups.  Thus b) follows from a).

\smallskip
\noindent
c) The hypothesis that $(K_{\bar t})^o_{\mathrm{red}}$ is unipotent is
equivalent to the equalities $\mu(t)=\alpha(t)=0$, and likewise for
$K_{\bar s}$.  Assertion c) therefore follows formally from b).

\smallskip
\noindent
d) Since $u_t$ is an isogeny, $K_t$ is finite.  We are therefore under the
hypotheses of c), and $(K_{\bar s})^o_{\mathrm{red}}$ is unipotent.  As
$G_s$ has unipotent rank zero, it follows that $K_s$ is finite.  Moreover,
the fact that $u_t$ is an isogeny gives
$\dim G_t=\dim H_t$.  Since $G$ and $H$ are flat and locally of finite
presentation, one has
\[
 \dim G_s=\dim G_t,\qquad \dim H_s=\dim H_t
 \quad\text{(SGA 3, VI$_B$, 4.3)},
\]
and hence $\dim G_s=\dim H_s$.  Since $\dim\ker u_s=0$, it follows that
$\dim H_s=\dim u_s(G_s)$ (SGA 3, VI$_B$, 1.2), which, because $H_s$ is
smooth, is equivalent to $u_s$ being flat (SGA 3, VI$_A$, 5.6).  There is
therefore an open neighborhood $U$ of $s$ such that $u$ is flat and
quasi-finite at the points of the unit section over $U$ (SGA 3, VI$_B$,
2.2(i) and \SGARefLink{sga7.sga7i.eix.n2-5}{2.5}(ii)).  As for $u_s$, it follows that $u_{s'}$ is flat for
every $s'\in U$; hence $u|_{G_U}$ is flat and quasi-finite, since it has
these properties on every fibre.

\smallskip
\noindent
e) For the first assertion, a base change immediately reduces us to the
case where $S$ is integral with generic point $t$, and hence to proving the
second assertion of e).  That assertion follows at once from d) and the
following lemma.

\medskip
\noindent
\underline{Lemma} \SGARefTarget{sga7.sga7i.eix.n3-1-3}{3.1.3}. \underline{Let $S$ be a scheme, let $K$ be a
quasi-finite $S$-scheme of finite presentation that is separated and flat
over $S$, let $s\in S$, and let $t\in S$ be a generalization of $s$.  Then}
\[
 \operatorname{rank}K_s\leq \operatorname{rank}K_t.
\]
\underline{If $\operatorname{rank}K_\eta\leq1$ for every maximal point
$\eta$ of $S$, then $K\longrightarrow S$ is an open immersion (and hence
an isomorphism if it admits a section).}

For the first assertion, strict localization reduces us to the case where
$S$ is strictly local.  Decomposing $K$ as the disjoint union of a finite
scheme $K^f$ and a scheme with empty special fibre, we obtain
\[
 \operatorname{rank}K_s
 =\operatorname{rank}K^f_s
 =\operatorname{rank}K^f_t
 \leq\operatorname{rank}K_t.
\]
The second equality follows because $K^f$ is flat, finite, and of finite
presentation over $S$, and therefore locally free over $S$.

For the second assertion, the preceding argument already shows that
$\operatorname{rank}K_s\leq1$ for every $s\in S$.  Thus $K_s$ is empty or
isomorphic to $\operatorname{Spec}k(s)$.  Consequently,
$K\longrightarrow S$ is a monomorphism (SGA 3, VI$_B$, 2.11), and, being
flat and of finite presentation, it is an open immersion (EGA IV, 17.9.1).
This proves the lemma.

\medskip
\noindent
\underline{Remark} \SGARefTarget{sga7.sga7i.eix.n3-1-4}{3.1.4}.  Part \SGARefLink{sga7.sga7i.eix.n3-1}{3.1} a), and the remainder of Proposition
\SGARefLink{sga7.sga7i.eix.n3-1}{3.1}, give no information about ${}_nK_s$ when $n$ is a power of $p$ (for
$p>1$).  At least when $G$ is smooth (and $S$ is locally noetherian or $H$
is of finite presentation over $S$), essentially the same method as in the
proof of a) gives such information.  Indeed, reducing to the case where
$S$ is complete local noetherian and using SGA 3, XV, 1.6 c), X, 3.2, IX,
\SGARefLink{sga7.sga7i.eix.n3-6}{3.6}, and IX, \SGARefLink{sga7.sga7i.eix.n5-2}{5.2}, one obtains
\begin{equation}\tag{\SGARefTarget{sga7.sga7i.eix.n3-1-4-1}{3.1.4.1}}
 \operatorname{rank}_{\mathrm{mult}}({}_nK_s)
 \mid
 \operatorname{rank}_{\mathrm{mult}}({}_nK_t),
\end{equation}
where $\operatorname{rank}_{\mathrm{mult}}$ denotes the rank of the largest
subgroup of multiplicative type of the commutative algebraic group under
consideration (SGA 3, XVII, 7.2.1).  Notice that a torsion algebraic group
of multiplicative type is finite.  Formula (\SGARefLink{sga7.sga7i.eix.n3-1-4-1}{3.1.4.1}) is valid for every
integer $n\geq1$.  It follows from SGA 3, XVII, 4.6.1(vi) that, in c), if
$K_t$ is unipotent (SGA 3, VII, 1.1), then $K_s$ is unipotent as well.

% END INCLUDED SGA7I_English_current_buildable_source_20260823/components/136_expose_IX_section_3_1.tex

% BEGIN INCLUDED SGA7I_English_current_buildable_source_20260823/components/137_expose_IX_sections_3_2_through_3_4.tex
% Expose IX, sections 3.2--3.4.
% French transcription lines 841--875; authority physical pp. 359--361
% (source folios 347--349), stopping before Proposition 3.5.
% The opening on physical p. 359 was checked with the preceding component;
% physical pp. 360--361 were checked directly in four overlapping 1100-dpi
% bands per page.  The prose and formulas were unambiguous at that resolution,
% so no higher-detail escalation was needed.

\medskip
\noindent
\underline{Proposition} \SGARefTarget{sga7.sga7i.eix.n3-2}{3.2}. \underline{Let $S$ be a regular integral
locally noetherian scheme of dimension $1$, let $K$ be its field of rational
functions, let $A_K$ be an abelian scheme over $K$, and let $A$ be its Neron
model over $S$, which is a smooth separated group scheme over $S$ (and is
quasi-projective over $S$ if $S$ is noetherian [23, VIII, 2, 3\textdegree]).
The following conditions are equivalent:}

\smallskip
\noindent
(i) \underline{For every $s\in S$, $A_s^o$ has unipotent rank zero; that is,
it is an extension of an abelian scheme by a torus (\SGARefLink{sga7.sga7i.eix.n2-1-5}{2.1.5}).}

\smallskip
\noindent
(ii) \underline{There exists a smooth separated group scheme $G$ of finite
type over $S$, extending $A_K$, whose fibres have unipotent rank zero.}

\smallskip
\noindent
\underline{Moreover, if $G$ is as in (ii), then the unique morphism
$G\longrightarrow A$ inducing the given isomorphism on the generic fibres
is an isomorphism from $G$ onto an open subgroup of $A$; it therefore induces
an isomorphism of identity components}
\[
 G^o\xrightarrow{\ \sim\ }A^o.
\]

\medskip
\noindent
\SGARefTarget{sga7.sga7i.eix.n3-2-1}{3.2.1}.  Let us note at once that the definition of the Neron model $A$ of
$A_K$ is the same as in the local case recalled in \SGARefLink{sga7.sga7i.eix.n1-1}{1.1}.  The existence of
this Neron model, and the fact that it is separated and of finite type over
$S$, are well known and in fact reduce immediately to the local case.  Of
course, for every closed point $s\in S$,
\[
 S'=\operatorname{Spec}(\mathcal O_{S,s})
\]
is a trait, and $A\times_S S'$ is precisely the Neron model of $A_K$ over
$S'$ in the sense considered in \SGARefLink{sga7.sga7i.eix.n1-1}{1.1}.

Let us prove \SGARefLink{sga7.sga7i.eix.n3-2}{3.2}.  It is trivial that (i) implies (ii), since it suffices to
take $G=A$.  It remains to prove that if $G$ is as in (ii), then
$G\longrightarrow A$ is an open immersion.  This is indeed a special case
of \SGARefLink{sga7.sga7i.eix.n3-1}{3.1} e).

\medskip
\noindent
\underline{Corollary} \SGARefTarget{sga7.sga7i.eix.n3-3}{3.3}. \underline{With the notation of \SGARefLink{sga7.sga7i.eix.n3-2}{3.2}, suppose
that the identity components of the fibres of the Neron model have
unipotent rank zero.  Let $S'$ be a scheme satisfying the same hypotheses
as $S$, let $K'$ be its field of rational functions, let
$f:S'\longrightarrow S$ be a dominant morphism, let}
\[
 A_{K'}\simeq A_K\otimes_KK'
 \simeq \text{the generic fibre of }A\times_SS',
\]
\underline{let $A'$ be the Neron model of $A_{K'}$, and consider the unique
morphism}
\[
 u:A\times_SS'\longrightarrow A'
\]
\underline{inducing the identity on the generic fibres.  Then $u$ is an
open immersion and therefore induces an isomorphism on identity
components}
\begin{equation}\tag{\SGARefTarget{sga7.sga7i.eix.n3-3-1}{3.3.1}}
 A^o\times_SS'\xrightarrow{\ \sim\ }{A'}^o.
\end{equation}
Indeed, it suffices to apply \SGARefLink{sga7.sga7i.eix.n3-2}{3.2} over $S'$, taking
$G\simeq A\times_SS'$.

\medskip
\noindent
\SGARefTarget{sga7.sga7i.eix.n3-3-2}{3.3.2}.  Thus the equivalent properties of \SGARefLink{sga7.sga7i.eix.n3-2}{3.2} imply a partial
``stability'' property for Neron models.  One should bear in mind that when
$A_K$ does not have good reduction, that is, when $A$ is not an abelian
scheme, it is not true that, for every $S'$ as in \SGARefLink{sga7.sga7i.eix.n3-3}{3.3} (even for $S'$ finite
over $S$ and $K'/K$ separable), the morphism
\[
 A\times_SS'\longrightarrow A'
\]
itself is an isomorphism, or that it induces isomorphisms on the groups of
connected components of the fibres (see \SGARefLink{sga7.sga7i.eix.n11-10}{11.10} below for the behavior of
these groups of connected components under extension of the base).  These
observations, reinforced by \SGARefLink{sga7.sga7i.eix.n3-9}{3.9} below, justify the following terminology,
introduced by Mumford (and which also appears to agree with the notions of
the same name studied in his book [17]).

\medskip
\noindent
\underline{Definition} \SGARefTarget{sga7.sga7i.eix.n3-4}{3.4}. \underline{With the notation of \SGARefLink{sga7.sga7i.eix.n3-2}{3.2}, the
abelian scheme $A_K$ over $K$ is said to have semistable reduction over $S$
if the equivalent conditions (i) and (ii) of \SGARefLink{sga7.sga7i.eix.n3-2}{3.2} are satisfied.}

\medskip
\noindent
\SGARefTarget{sga7.sga7i.eix.n3-4-0}{3.4.0}.  Notice that in the form (ii), this condition does not invoke the
notion of a Neron model.  Moreover, by localization it reduces immediately
to the case where $S$ is a trait, as follows for example from \SGARefLink{sga7.sga7i.eix.n3-2-1}{3.2.1}.
Furthermore, since formation of the Neron model commutes with base change
from $S$ to its henselization or its strict henselization, the notion even
reduces to the case where $S$ is a strictly local trait.

% END INCLUDED SGA7I_English_current_buildable_source_20260823/components/137_expose_IX_sections_3_2_through_3_4.tex

% BEGIN INCLUDED SGA7I_English_current_buildable_source_20260823/components/138_expose_IX_sections_3_5_through_3_9.tex
% Expose IX, Proposition 3.5 through Corollary 3.9.
% French transcription lines 880--960; authority physical pp. 362--366
% (source folios 350--354); hard stop before section 4 on p. 366.
% All five pages were checked directly in four overlapping 1100-dpi bands.
% The empty parenthesis printed after Z_l(1) in 3.5.2 contains no recoverable
% text and is omitted rather than conjecturally supplied.
% Transparent source corrections:
% - Proposition 3.5(ii) prints V \subset V^\perp.  The filtration diagram
%   (2.5.4), condition (iii), and the proof require V^\perp \subset V.
% - In the proof, the print says tautologically that I acts trivially on
%   U^I=V; condition (iv) says that it acts trivially on U/U^I=U/V.
% - The proof of Corollary 3.8 prints (v) => (vi), whereas the argument proves
%   the required converse (vi) => (iv).

\medskip
\noindent
\underline{Proposition} \SGARefTarget{sga7.sga7i.eix.n3-5}{3.5}. \underline{(Galois criterion for semistable
reduction).  Suppose that $S$ is a trait.  The following conditions on the
abelian scheme $A_K$ are equivalent} (the notation being that of (\SGARefLink{sga7.sga7i.eix.n2-5-3}{2.5.3}),
applied to the situation deduced from $(A_K,S)$ by passing to the
henselization of $S$):

\smallskip
\noindent
(i) \underline{$A_K$ has semistable reduction over $S$} (\SGARefLink{sga7.sga7i.eix.n3-4}{3.4}).

\smallskip
\noindent
(ii) \underline{With the notation of (\SGARefLink{sga7.sga7i.eix.n2-5-3}{2.5.3}) and (\SGARefLink{sga7.sga7i.eix.n2-5-4}{2.5.4}), one has}
\[
 V^\perp\subset V.
\]

\noindent
(iii) \underline{For the toric part $W$ of the Tate module one has}
\[
 W=V^\perp.
\]

\noindent
(iv) \underline{The action of the inertia group $I$ on
$U=T_\ell(A(\overline K))$ is unipotent of level $2$} (that is, there is an
$I$-stable submodule $U'$ of $U$ such that $I$ acts trivially on $U'$ and
on $U/U'$).

We may suppose that $S$ is henselian (\SGARefLink{sga7.sga7i.eix.n3-4-0}{3.4.0}).  With the notation of \SGARefLink{sga7.sga7i.eix.n2-1-9}{2.1.9},
(i) means that
\[
 \lambda=0.
\]
By the rank formulas following (\SGARefLink{sga7.sga7i.eix.n2-5-4}{2.5.4}), this is equivalent to
$V^\perp=V\cap V^\perp$, that is, to (ii).  Since
$W=V\cap V^\perp$ by the orthogonality theorem \SGARefLink{sga7.sga7i.eix.n2-4}{2.4}, conditions (ii) and
(iii) are equivalent.  Moreover, (iv) also says that $I$ acts trivially
on $U/U^I=U/V$.  By duality this means that it acts trivially on
$V^\perp$, that is, that $V^\perp\subset U^I=V$, which is (ii).  This
completes the proof.

\medskip
\noindent
\underline{Remarks} \SGARefTarget{sga7.sga7i.eix.n3-5-1}{3.5.1}.  Conditions (ii) and (iv) plainly depend only
on the action of $I$ on the $\mathbb Q_\ell$-vector space
$U\otimes_{\mathbb Z_\ell}\mathbb Q_\ell$, and could have been formulated
in terms of that action.  Let us also note that these conditions are
manifestly stable under every base change by a morphism of traits
$S'\longrightarrow S$.  Finally, if (iv) holds for $A_K$, it plainly holds
for every abelian subscheme and every quotient abelian scheme of $A_K$,
and in particular for every abelian scheme isogenous to $A_K$ (for
example, for $A'_K$).

\medskip
\noindent
\underline{Corollary} \SGARefTarget{sga7.sga7i.eix.n3-5-2}{3.5.2}. \underline{Let $g$ be a topological generator
of $\mathbb Z_\ell(1)=T_\ell(\overline k^{\,*})$.  The equivalent conditions
of \SGARefLink{sga7.sga7i.eix.n3-5}{3.5} are also equivalent to the following condition:}

\smallskip
\noindent
% SGA-SRC-001269 — internal reference omitted in source; target not recoverable
(v) \underline{The action of the inertia group $I$ on
$U=T_\ell(A(\overline K))$ factors through the maximal pro-$\ell$ quotient
$I(\ell)$ of $I$, which is canonically isomorphic to $\mathbb Z_\ell(1)$,
and the action $g_U$ of the generator $g$ on $U$ satisfies}
\begin{equation}\tag{\SGARefTarget{sga7.sga7i.eix.n3-5-3}{3.5.3}}
 (1-g_U)^2=0.
\end{equation}
Indeed, for every continuous unipotent action of a profinite group $I$ on
an $\ell$-adic vector space, one sees at once that the image of $I$ is a
pro-$\ell$ group.  It follows that (v) is equivalent to (iv).

We can now prove the central result of the present expose:

\medskip
\noindent
\underline{Theorem} \SGARefTarget{sga7.sga7i.eix.n3-6}{3.6}. \underline{(Semistable reduction theorem).  Let
$S$ be a connected regular noetherian scheme of dimension $1$, let $K$ be
the field of rational functions of $S$, and let $A_K$ be an abelian scheme
over $K$.  Then there is a finite Galois extension $K'$ of $K$ such that
the abelian scheme $A_{K'}=A_K\otimes_KK'$ has semistable reduction} (\SGARefLink{sga7.sga7i.eix.n3-4}{3.4})
\underline{over the normalization $S'$ of $S$ in $K'$.}

Notice that $S'$ is connected, regular, noetherian, and of dimension $1$,
so that the stated conclusion makes sense.  To prove \SGARefLink{sga7.sga7i.eix.n3-6}{3.6}, choose a
separable closure $\overline K$ of $K$ and write $\overline K$ as the
filtered union of its finite Galois subextensions $K_i$.  Let $S_i$ be the
normalization of $S$ in $K_i$, and let $U$ be the largest open subset of
$S$ over which $A_K$ has good reduction.  Since $S$ is connected,
noetherian, and one-dimensional, and $U\ne\varnothing$, the complement
$T=S-U$ is a finite set of closed points of $S$, and
$\mathcal O_{S,t}$ is a discrete valuation ring for every $t\in T$.  It is
enough to find, for every $t\in T$, an index $i=i(t)$ such that $A_{K_i}$
has semistable reduction at the points of $S_i$ above $t$: one then takes
for $K'$ the compositum of the extensions $K_{i(t)}$, for $t\in T$.
Thus we are reduced to the case where $S$ is a trait.

Moreover, since the Galois group permutes the points of the normalization
lying above the closed point $s$ of $S$, it is enough to find $K'$ such
that $A_{K'}$ has semistable reduction at \emph{one} point of the
normalization $S'$.  Here it is unnecessary at first to require that
$K'/K$ be Galois: it is enough to find a finite etale extension $K'/K$
satisfying the preceding condition, since its Galois closure satisfies the
same condition.  Using criterion \SGARefLink{sga7.sga7i.eix.n3-5}{3.5}(iv), we see that Theorem \SGARefLink{sga7.sga7i.eix.n3-6}{3.6} reduces
to the following assertion.

\medskip
\noindent
\underline{Corollary} \SGARefTarget{sga7.sga7i.eix.n3-7}{3.7}. \underline{Suppose again that $S$ is a trait.
Then there is an open subgroup $\pi'$ of $\pi$ such that the action of
$I'=I\cap\pi'$ on $U=T_\ell(A(\overline K))$ is unipotent of level $2$;
equivalently, there is an open subgroup $I'$ of $I$ whose action on $U$ is
unipotent of level $2$.}

There is a canonical isomorphism, compatible with the actions of $\pi$,
\begin{equation}\tag{*}
 H^1(A_{\overline K},\mathbb Z_\ell)
 \simeq
 \operatorname{Hom}\bigl(T_\ell(A(\overline K)),\mathbb Z_\ell\bigr),
\end{equation}
which shows that assertion \SGARefLink{sga7.sga7i.eix.n3-7}{3.7} is equivalent to the same assertion for
the actions of $I$ on the left-hand side of (*).  But that assertion was
proved in Expose III (using resolution of singularities for excellent
schemes of dimension $2$).

\medskip
\noindent
\underline{Corollary} \SGARefTarget{sga7.sga7i.eix.n3-8}{3.8}. \underline{The equivalent conditions (i)--(iv)
of \SGARefLink{sga7.sga7i.eix.n3-5}{3.5} are also equivalent to the following condition:}

\smallskip
\noindent
(vi) \quad \underline{The action of $I$ on
$U=T_\ell(A(\overline K))$ is unipotent.}

Indeed, to prove that (vi) implies (iv), observe that if a group $I$
acting unipotently on a finite-dimensional vector space $E$ over a field
of characteristic zero (here $\mathbb Q_\ell$) has a subgroup $I'$ of
finite index acting with level $m$ (here $m=2$), then $I$ itself acts with
level $m$.  In fact, since $I'$ has finite index in $I$, the algebraic
group $G'\subset\operatorname{GL}(E)$ that it defines has finite index in
the group $G$ defined by $I$; but $G$ is unipotent and hence connected (in
characteristic zero), so $G=G'$.  Thus $I$ and $I'$ have the same
``level of unipotence.''

We also obtain a converse to \SGARefLink{sga7.sga7i.eix.n3-3}{3.3}:

\medskip
\noindent
\underline{Corollary} \SGARefTarget{sga7.sga7i.eix.n3-9}{3.9}. \underline{The notation is that of \SGARefLink{sga7.sga7i.eix.n3-2}{3.2}.  The
conditions (i) and (ii) of \SGARefLink{sga7.sga7i.eix.n3-2}{3.2} are also equivalent to the following
condition:}

\smallskip
\noindent
(iii) \quad \underline{For every finite etale extension $K'$ of $K$, let
$S'$ denote the normalization of $S$ in $K'$, and let $A'$ denote the
Neron model of
$A_{K'}=A\otimes_SK'=A_K\otimes_KK'$.  Then the canonical morphism}
\[
 A\times_SS'\longrightarrow A'
\]
\underline{induces an isomorphism on identity components.}

The necessity of the condition was proved, in a stronger form, in \SGARefLink{sga7.sga7i.eix.n3-3}{3.3}.
For sufficiency, choose $K'$ as in \SGARefLink{sga7.sga7i.eix.n3-6}{3.6}.  Since the fibres of ${A'}^o$ have
unipotent rank zero, the same is true of the fibres of
$(A\times_SS')^o=A^o\times_SS'$, and hence also of the fibres of $A^o$,
from which the former are obtained by extension of the base field.

% END INCLUDED SGA7I_English_current_buildable_source_20260823/components/138_expose_IX_sections_3_5_through_3_9.tex

% BEGIN INCLUDED SGA7I_English_current_buildable_source_20260823/components/139_expose_IX_sections_4_1_and_4_2.tex
% Expose IX, sections 4.1 and 4.2.
% French transcription lines 963--1093; authority physical pp. 366--370
% (source folios 354--358); hard stop before Theorem 4.3 on p. 370.
% All five pages were checked directly in four overlapping 1100-dpi bands.
% Diagram (4.2.4) was additionally checked in four overlapping 5000-dpi
% crops: all six arrowheads, all three unheaded extension-field segments,
% the two curved segments, node placement, and label sides were compared.
% Transparent source corrections:
% - The footnote to 4.1 prints ell != p, although the main text already
%   assumes ell != p and sections 5--6 treat precisely the case ell = p.
% - The printed reference to the semistable-reduction criterion as (4.5)
%   is corrected to (3.5), where that criterion is stated and proved.
% - After (4.2.11), the source prints the undefined Gamma^o, whereas the
%   exact sequence (4.2.3) makes Gamma the kernel of pi' -> pi_o; hence the
%   displayed character is trivial on Gamma.
% - The parenthetical reference following (4.2.6) is blank after "I" in
%   the source.  It is identified as omitted rather than conjecturally filled.

\medskip
\noindent
\underline{4. Application to a conjecture of Serre--Tate and to the
conductor}

\medskip
\noindent
\SGARefTarget{sga7.sga7i.eix.n4-1}{4.1}. We suppose that $S$ is henselian, and retain the notation of (\SGARefLink{sga7.sga7i.eix.n2-5-3}{2.5.3}),
with $\ell\ne p$.
\footnote{The case $\ell=p$ will be studied in \S\S5 and 6 below.}
To study the action of $\pi=\operatorname{Gal}(\overline K/K)$ and of its
inertia subgroup $I$ on $U=T_\ell(A(\overline K))$, we shall use the fact
that there exists an open subgroup $I'$ of $I$, which we may moreover take
to be normal in $\pi$, such that the action of $I'$ on $U$ is unipotent of
level two; that is,
\[
 U^{I'}\supset (U^{I'})^\perp.
\]
It is then clear that $U^{I'}$ does not depend on the particular choice of
$I'$.  We might call it the \underline{essentially fixed submodule} of
$U$, and shall denote it by $U^{\mathrm{ef}}$.  Its orthogonal will be
called the \underline{essentially toric submodule} of $U$, and will be
denoted by $U^{\mathrm{et}}$.  We thus have a three-step filtration
\begin{equation}\tag{\SGARefTarget{sga7.sga7i.eix.n4-1-1}{4.1.1}}
 U\supset U^{\mathrm{ef}}\supset U^{\mathrm{et}}\supset0,
\end{equation}
which is plainly invariant under the action of $\pi$.

Let us also note that there is a largest subgroup of $I$ acting
unipotently (or, equivalently here, unipotently of level two) on $U$: it
is the subgroup $I'$ consisting of the elements acting unipotently.  It is
also the kernel of the representation of $I$ on the graded module
associated with the filtration (\SGARefLink{sga7.sga7i.eix.n4-1-1}{4.1.1}), or, equivalently in the present
case, the kernel of the induced representation of $I$ on
$U^{\mathrm{ef}}$.  By the Galois criterion for semistable reduction
(\SGARefLink{sga7.sga7i.eix.n3-5}{3.5}), this subgroup $I'$ corresponds to \underline{the smallest Galois
extension} $\widetilde K'$ of $\widetilde K$ (the maximal unramified
extension of $K$) \underline{such that $A_{\widetilde K'}$ has
semistable reduction}.  This description also shows that the subgroup
$I'$ of $I$ is independent of the choice of the prime $\ell\ne p$.

Notice that, once the representation of $\pi$ on $U^{\mathrm{ef}}$, and
hence on $U^{\mathrm{et}}$, is known, the representation on
$U/U^{\mathrm{ef}}$ is known up to isogeny, by means of the isogeny
induced by the polarization pairing (\SGARefLink{sga7.sga7i.eix.n2-5-1}{2.5.1}):
\begin{equation}\tag{\SGARefTarget{sga7.sga7i.eix.n4-1-2}{4.1.2}}
 U/U^{\mathrm{ef}}\longrightarrow (U^{\mathrm{et}})^\vee(1)
 \qquad\text{an isogeny},
\end{equation}
where $(1)$ denotes the Tate twist by
$\mathbb Z_\ell(1)=T_\ell(\overline K^*)$; this homomorphism is plainly
compatible with the operations of $\pi$.  More precisely, there is a
canonical \underline{isomorphism}
\[
 U/U^{\mathrm{ef}}\simeq ((U')^{\mathrm{et}})^\vee(1),
 \qquad U'=T_\ell(A'(\overline K)).
\]
It therefore remains to study more closely the action of $\pi$ on
$U^{\mathrm{ef}}$.

By construction this action is trivial on the normal subgroup $I'$ of
$I$, and hence factors through a representation of the quotient group
\begin{equation}\tag{\SGARefTarget{sga7.sga7i.eix.n4-1-3}{4.1.3}}
 \pi'=\pi/I',
\end{equation}
an action of which we shall now give a description ``independent of
$\ell$.''

\medskip
\noindent
\SGARefTarget{sga7.sga7i.eix.n4-2}{4.2}. Let $K'$ be a Galois extension of $K$ corresponding to an open
normal subgroup $V$ of $\pi$ such that $V\cap I=I'$.  The Galois group
$G$ of $K'/K$ then occurs in an exact sequence
\begin{equation}\tag{\SGARefTarget{sga7.sga7i.eix.n4-2-1}{4.2.1}}
 1\longrightarrow\Gamma\longrightarrow G\longrightarrow G_o
 \longrightarrow1,
\end{equation}
where
\begin{equation}\tag{\SGARefTarget{sga7.sga7i.eix.n4-2-2}{4.2.2}}
 \Gamma=I/I'
\end{equation}
is the inertia subgroup, and $G_o$ is the Galois group of the largest
unramified subextension $K''$ of $K'/K$.  Thus (\SGARefLink{sga7.sga7i.eix.n4-2-1}{4.2.1}) is the image of the
analogous exact sequence
\begin{equation}\tag{\SGARefTarget{sga7.sga7i.eix.n4-2-3}{4.2.3}}
 1\longrightarrow\Gamma\longrightarrow\pi'\longrightarrow\pi_o
 \longrightarrow1,
\end{equation}
where $\pi'$ is the group (\SGARefLink{sga7.sga7i.eix.n4-1-3}{4.1.3}) acting on $U^{\mathrm{ef}}$.  We
therefore have the following diagram of field inclusions:
\begin{equation}\tag{\SGARefTarget{sga7.sga7i.eix.n4-2-4}{4.2.4}}
\begin{tikzpicture}[baseline=(current bounding box.center),
  every node/.style={inner sep=2pt}]
  \node (K) at (0,0) {$K$};
  \node (Ktwo) at (4.25,0) {$K''=\widetilde K\cap K'$};
  \node (Kt) at (8.5,0) {$\widetilde K$};
  \node (Kp) at (4.25,-1.8) {$K'$};
  \node (Ktp) at (8.5,-1.8) {$\widetilde K'$};
  \node (Kb) at (8.5,-3.6) {$\overline K$};
  \draw[->] (K) -- node[above] {$G_o$} (Ktwo);
  \draw[->] (Ktwo) -- (Kt);
  \draw[->] (Ktwo) -- node[right] {$\Gamma$} (Kp);
  \draw[->] (Kp) -- (Ktp);
  \draw[->] (Kt) -- node[right] {$\Gamma$} (Ktp);
  \draw[->] (Ktp) -- node[right] {$I'$} (Kb);
  \draw[-] (K) -- node[below left,pos=.58] {$G$} (Kp);
  \draw[-] (K) to[out=65,in=115,looseness=.85]
    node[above,pos=.78] {$\pi_o$} (Kt);
  \draw[-] (Kt) to[out=12,in=12,looseness=1.25]
    node[right,pos=.52] {$I$} (Kb);
\end{tikzpicture}
\end{equation}

Let $S'$ be the normalization of $S$ in $K'$, and let $A^*$ be the
Neron model of $A_{K'}$.  Then, by the Galois criterion for semistable
reduction (\SGARefLink{sga7.sga7i.eix.n3-5}{3.5}), $A_{K'}$ \underline{has semistable reduction over
$S'$}; that is, $A^{*o}_o$ \underline{is an extension of an abelian
scheme by a torus}.  On the other hand, by transport of structure, $G$
\underline{acts on $A^*$, and hence on $A^{*o}$, by automorphisms in a
manner compatible with the given action of $G$ on the base scheme $S'$
of $A^*$}.  It is then clear that \underline{the action of $\pi'$ on}
\[
 U^{\mathrm{ef}}=T_\ell(A^{*o})^f(\overline K)
\]
\underline{is obtained by transport of structure from the action of
$\pi'$ on $(A^{*o}/K',\widetilde K')$}, where $\pi'$ acts on $A^*$
through $G$.  Taking into account the canonical isomorphism
\begin{equation}\tag{\SGARefTarget{sga7.sga7i.eix.n4-2-6}{4.2.6}}
 U^{\mathrm{ef}}\simeq T_\ell(A^{*o})^f(\overline K)
 \simeq T_\ell(A^{*o}_o)(\overline k)
\end{equation}
% SGA-SRC-001270 — internal reference omitted in source; target not recoverable
(valid for every twisted constant $\ell$-adic sheaf over $S'$; cf.
Expos\'e I [reference omitted in the source]), we obtain the following
description of the action of $\pi'$ on $U^{\mathrm{ef}}$.

Consider the smooth connected group scheme $A^{*o}_o$ over the residue
field $k'\hookrightarrow\overline k$ of $S'$, which is a quasi-Galois
extension of $k$ with Galois group $G_o$.  The quotient group $G$ of
$\pi'$ acts on $A^{*o}_o$ compatibly with its action on $k'$ through
$G_o$.  Moreover, $\pi'$ acts on $\overline k$ through its quotient
$\pi_o$, and hence acts on the situation
$(A^{*o}_o/k',\overline k)$.  Thus, by transport of structure, $\pi'$
acts on
\[
 U^{\mathrm{ef}}=T_\ell(A^{*o}_o)(\overline k),
\]
and this is the desired action of $\pi'$.  In particular, the subgroup
$\Gamma$ of $G$, which acts trivially on $k'$, acts on $A^{*o}_o$ by
$k'$-automorphisms, and the corresponding representation of
$\Gamma=I/I'$ on $T_\ell(A^{*o}_o)(\overline k)$ is the representation
deduced from the natural representation of the inertia group $I$ on
$U^{\mathrm{ef}}$.

On the other hand, the essentially toric part $U^{\mathrm{et}}$ of $U$
corresponds simply to the submodule
\begin{equation}\tag{\SGARefTarget{sga7.sga7i.eix.n4-2-7}{4.2.7}}
 U^{\mathrm{et}}\simeq T_\ell(T^*_o)(\overline k)
 \subset U^{\mathrm{ef}}\simeq T_\ell(A^{*o}_o)(\overline k),
\end{equation}
where $T^*_o$ is the largest subtorus of $A^{*o}_o$.  Likewise, put
\begin{equation}\tag{\SGARefTarget{sga7.sga7i.eix.n4-2-8}{4.2.8}}
 B^*_o=A^{*o}_o/T^*_o
\end{equation}
for the abelian part of $A^{*o}_o$.  The action of $G$ on $A^{*o}_o$
induces separately actions of $G$ on $T^*_o$ and on $B^*_o$, and the
action of $\pi'$ on $U^{\mathrm{et}}$ and on
$U^{\mathrm{ef}}/U^{\mathrm{et}}$ is obtained from these actions by
means of the isomorphism (\SGARefLink{sga7.sga7i.eix.n4-2-7}{4.2.7}) and the corresponding isomorphism
\begin{equation}\tag{\SGARefTarget{sga7.sga7i.eix.n4-2-9}{4.2.9}}
 U^{\mathrm{ef}}/U^{\mathrm{et}}
 \simeq T_\ell(B^*_o)(\overline k).
\end{equation}
It follows from this and the isogeny (\SGARefLink{sga7.sga7i.eix.n4-1-2}{4.1.2}) that the character
$g\longmapsto\operatorname{Tr}g_U$ of the action of $\pi'$ (or of
$\pi$) on $U$ is given by
\begin{equation}\tag{\SGARefTarget{sga7.sga7i.eix.n4-2-10}{4.2.10}}
 \operatorname{Tr}g_{U_\ell}
 =\operatorname{Tr}g_{B^*_o,\ell}
 +\operatorname{Tr}g_{T^*_o,\ell}
 +\chi_\ell(g)\operatorname{Tr}g^{-1}_{T^*_o,\ell}.
\end{equation}
Here $U_\ell$ has been written for $U$ on the left to recall the chosen
prime $\ell\ne p$; $g_{B^*_o,\ell}$ and $g_{T^*_o,\ell}$ denote,
respectively, the actions of $g$ on the Tate modules
$T_\ell(B^*_o)(\overline k)$ and $T_\ell(T^*_o)(\overline k)$ induced by
its actions on $B^*_o$ and $T^*_o$ just described; and
\begin{equation}\tag{\SGARefTarget{sga7.sga7i.eix.n4-2-11}{4.2.11}}
 \chi_\ell:\pi'\longrightarrow\pi_o\longrightarrow\mathbb Z_\ell^*
\end{equation}
is the character of $\pi'$ (or of $\pi_o$) induced by the action on
$\mathbb Z_\ell(1)=T_\ell(\overline k^*)$---a character that is trivial
on $\Gamma$.

We are now in a position to prove a special case of a conjecture of
Serre--Tate [29, Appendix].
\footnote{This is essentially the case $i=1$ of those conjectures for
the $H^i$.}

% END INCLUDED SGA7I_English_current_buildable_source_20260823/components/139_expose_IX_sections_4_1_and_4_2.tex

% BEGIN INCLUDED SGA7I_English_current_buildable_source_20260823/components/140_expose_IX_theorem_4_3_and_corollary_4_4.tex
% Expose IX, Theorem 4.3 and Corollary 4.4.
% French transcription lines 1098--1171; authority physical pp. 371--375
% (source folios 359--363); hard stop before section 4.5 on p. 375.
% All five pages were checked directly in four overlapping 1100-dpi bands.
% The formulas, Frobenius powers, inverses, Tate twists, primes, and bars
% were unambiguous at that resolution; no 5000--9000-dpi escalation was
% warranted for this diagram-free block.
% In applying Lemma 4.3.2, the printed shorthand replacing k by a barred
% k-prime is rendered invariantly as an algebraic closure of k'.

\medskip
\noindent
\underline{Theorem} \SGARefTarget{sga7.sga7i.eix.n4-3}{4.3}. \underline{Let $S$ be a henselian trait with
residue field $k$ and function field $K$, and let $A_K$ be an abelian
scheme over $K$.  Consider the representation of}
$\pi=\operatorname{Gal}(\overline K/K)$ \underline{on the Tate module}
$U_\ell=T_\ell(A(\overline K))$, \underline{where $\ell$ is a prime
distinct from the residual characteristic.}

\smallskip
\noindent
a) \underline{There exists an open subgroup $I'$ of the inertia group
that acts unipotently on $U_\ell$, so that the character of the
representation of $\pi$ on $U_\ell$ is in fact a function on $\pi/I'$.}
\underline{Moreover, on $\Gamma=I/I'$ this function takes integral
values and is independent of $\ell$.  More generally, for every $g\in I$,
the coefficients of the characteristic polynomial of $g_{U_\ell}$ are
rational integers independent of $\ell$.}

\smallskip
\noindent
b) \underline{Suppose that $k$ is a finite field with $q$ elements, so
that}
\[
 \pi_o=\pi/I\simeq\operatorname{Gal}(\overline k/k)
\]
\underline{is identified with the profinite completion $\widehat{\mathbb
Z}$ of $\mathbb Z$, the generator $1$ of $\mathbb Z$ corresponding to
the Frobenius automorphism}
\[
 f_q:x\longmapsto x^q
\]
\underline{of $\overline k/k$.  Then, for every $g\in\pi$ whose image in
$\pi_o$ lies in $\mathbb Z$, the trace (and, more generally, every
coefficient of the characteristic polynomial) of $g_{U_\ell}$ is a
rational number independent of $\ell$.  It is even an integer if the
image of $g$ in $\pi_o$ is of the form $f_q^{,n}$, with $n\geq0$.}

Let us prove, for example, the assertions concerning traces.  The
assertions concerning the coefficients of the characteristic polynomial
could be proved in the same way.  One may also note that they follow
formally from the assertions about traces, by expressing these
coefficients in terms of the traces of the iterates (although those
formulas introduce denominators and therefore do not give the asserted
integrality), together with a lemma of Dwork--Fatou, for which a pleasant
proof is given by Kleiman in [14, 2.8].

\smallskip
\noindent
a) We use formula (\SGARefLink{sga7.sga7i.eix.n4-2-10}{4.2.10}), which here becomes
\begin{equation}\tag{\SGARefTarget{sga7.sga7i.eix.n4-3-0}{4.3.0}}
 \operatorname{Tr}g_{U_\ell}
 =\operatorname{Tr}g_{B^*_o,\ell}
 +\operatorname{Tr}g_{T^*_o,\ell}
 +\operatorname{Tr}g^{-1}_{T^*_o,\ell}.
\end{equation}
Here $g_{B^*_o}$ and $g_{T^*_o}$ are now $k'$-automorphisms of $B^*_o$
and $T^*_o$, respectively.  It therefore suffices to show that the two
corresponding traces are rational integers independent of $\ell$, which
is a general ``geometric'' assertion about algebraic groups.  For
$g_{B^*_o}$ this is a well-known result of Weil in the theory of abelian
varieties [5] [15].  For $g_{T^*_o}$ it is practically immediate: one
describes the representation of $\Gamma$ on $T=T^*_o$ by means of the
corresponding representation on the dual group $M'\simeq\mathbb Z^\mu$.
The required trace is that of the contragredient of the latter
representation, by the canonical isomorphism
\[
 T_\ell(T)\simeq\check M_\ell(1)
 \mathrel{\mathop{=}^{\mathrm{def}}}
 \check M'\otimes_{\mathbb Z}\mathbb Z_\ell(1).
\]

\smallskip
\noindent
b) If the image of $g$ in $\pi_o$ is $f_q^{,n}$, formula (\SGARefLink{sga7.sga7i.eix.n4-2-10}{4.2.10}) now
becomes
\begin{equation}\tag{\SGARefTarget{sga7.sga7i.eix.n4-3-1}{4.3.1}}
 \operatorname{Tr}g_{U,\ell}
 =\operatorname{Tr}g_{B^*_o,\ell}
 +\operatorname{Tr}g_{T^*_o,\ell}
 +q^n\operatorname{Tr}g^{-1}_{T^*_o,\ell}.
\end{equation}
We are thus reduced to proving the following lemma, in which $k$ is to be
replaced by an algebraic closure of $k'$, and $A$ by
$A^{*o}_o\otimes_{k'}\overline k$.

\medskip
\noindent
\underline{Lemma} \SGARefTarget{sga7.sga7i.eix.n4-3-2}{4.3.2}. \underline{Let $k$ be the algebraic closure of
a finite field, let $A$ be a commutative group scheme of finite type over
$k$, and let $g$ be an automorphism of this group scheme compatible with
the automorphism $f$ of $\operatorname{Spec}(k)$ obtained by transport of
structure from an iterated Frobenius automorphism}
\[
 x\longmapsto x^q,
 \qquad q=p^r,\quad r\geq0,
\]
\underline{where $q$ is a power of the characteristic.  Thus the action
of $g$ on functions $\varphi$ satisfies}
\[
 g^*(\lambda\varphi)=f^*(\lambda)g^*(\varphi)
 =\lambda^{q^{-1}}g^*(\varphi)
 =\lambda^{p^{-r}}g^*(\varphi)
 \qquad(\lambda\in k).
\]
\underline{Let $g_\ell$ be the automorphism induced by $g$ on
$T_\ell(A(k))$ by ``transport of structure,'' where $\ell$ is a prime
different from $p$.  Then the trace of $g_\ell$ (and, more generally,
every coefficient of its characteristic polynomial) is a rational
integer independent of $\ell$, and its eigenvalues are algebraic integers
of absolute value $q^{1/2}$ or $q$.  If $A$ is an abelian scheme, their
absolute values are $q^{1/2}$; if $A$ is affine, the eigenvalues are of
the form $q\varepsilon_i$, where the $\varepsilon_i$ are roots of unity.}

Indeed, consider the ``$q$-th power'' endomorphism $f_q^A$ of the
underlying absolute scheme of $A$, defined as the identity on points and
as the $q$-th power map on the structure sheaf.  It is an endomorphism of
the group scheme $A$ compatible with the endomorphism
$f_q^{\operatorname{Spec}(k)}=f^{-1}$ of the base
$\operatorname{Spec}(k)$.  Consequently the composite
\begin{equation}\tag{\SGARefTarget{sga7.sga7i.eix.n4-3-3}{4.3.3}}
 F=f_q^A g
\end{equation}
is an endomorphism of the group scheme $A$ compatible with the identity
on $\operatorname{Spec}(k)$; in other words, it is a $k$-endomorphism of
the group scheme $A$.  I claim that the automorphism $\Gamma_g$ of $A(k)$
defined by $g$ through transport of structure, namely
\begin{equation}\tag{\SGARefTarget{sga7.sga7i.eix.n4-3-4}{4.3.4}}
 \Gamma_g(u)=g\circ u\circ f^{-1}
 \qquad\text{for }u\in A(k),
\end{equation}
may also be described by
\begin{equation}\tag{\SGARefTarget{sga7.sga7i.eix.n4-3-5}{4.3.5}}
 \Gamma_g(u)=F(u)
 \qquad\text{for }u\in A(k).
\end{equation}
To verify this, it is enough to observe that the underlying set maps of
the morphisms $\operatorname{Spec}(k)\longrightarrow A$ represented by
the two sides are the same.  This reduces us to the equality obtained
from (\SGARefLink{sga7.sga7i.eix.n4-3-5}{4.3.5}) by composing its left-hand member on the left with $f_q^A$
and its right-hand member on the right with $f^{-1}$, since $f_q^A$ and
$f^{-1}$ induce the identity on the underlying sets of the source
schemes.  The resulting formula follows immediately from (\SGARefLink{sga7.sga7i.eix.n4-3-3}{4.3.3}) and
(\SGARefLink{sga7.sga7i.eix.n4-3-4}{4.3.4}).

Applying (\SGARefLink{sga7.sga7i.eix.n4-3-5}{4.3.5}) to $u\in{}_{\ell^\nu}A(k)$ and passing to the inverse
limit over $\nu$, we find that $g_\ell$ may also be interpreted as the
endomorphism $T_\ell(F)$ of $T_\ell(A(k))$ associated with the
endomorphism $F$ of $A$.  It is known, on the other hand, that
$T_\ell(F)$ has a characteristic polynomial with integral coefficients
independent of $\ell$.  By devissage, this reduces to the case in which
$A$ is either an abelian scheme or a torus, where the required result was
already used in part a) of the proof.  The same argument also proves the
assertions about the eigenvalues, in view, in the abelian-scheme case, of
the classical results of Weil [15].

This proves \SGARefLink{sga7.sga7i.eix.n4-3-2}{4.3.2} and completes the proof of \SGARefLink{sga7.sga7i.eix.n4-3}{4.3}.  At the same time it
proves the following result.

\medskip
\noindent
\underline{Corollary} \SGARefTarget{sga7.sga7i.eix.n4-4}{4.4}. \underline{Suppose that the hypotheses of
\SGARefLink{sga7.sga7i.eix.n4-3}{4.3} b) hold.  Then, for every $g\in\pi'$ whose image in $\pi_o$ is of the
form $f_q^n$ with $n\geq0$, the eigenvalues of the induced automorphisms
of}
\[
 U^{\mathrm{et}},\qquad
 U^{\mathrm{ef}}/U^{\mathrm{et}},\qquad
 U/U^{\mathrm{ef}}
\]
\underline{are algebraic integers whose absolute values are,
respectively, $q^n$, $q^{n/2}$, and $1$.  For the two extreme pieces one
may say more precisely that the corresponding eigenvalues are,
respectively,}
\[
 (q^n\varepsilon_i)_{1\leq i\leq\mu}
 \qquad\text{and}\qquad
 (\varepsilon_i)_{1\leq i\leq\mu},
\]
\underline{where the $\varepsilon_i$ $(1\leq i\leq\mu)$ are roots of
unity.}

Indeed, formula (\SGARefLink{sga7.sga7i.eix.n4-1-2}{4.1.2}) gives for the last piece the eigenvalues
\[
 \frac{q^n}{q^n\varepsilon_i}=\varepsilon_i^{-1}.
\]
But since $\varepsilon_i$ is a root of unity, one has
$\varepsilon_i^{-1}=\overline{\varepsilon_i}$ (complex conjugate), and
since the family of the $\varepsilon_i$ is stable under complex
conjugation (being the family of roots of an equation with integral
coefficients), the family of the $\varepsilon_i^{-1}$ is also the family
of the $\varepsilon_i$.  This proves the last assertion of \SGARefLink{sga7.sga7i.eix.n4-4}{4.4}.

% END INCLUDED SGA7I_English_current_buildable_source_20260823/components/140_expose_IX_theorem_4_3_and_corollary_4_4.tex

% BEGIN INCLUDED SGA7I_English_current_buildable_source_20260823/components/141_expose_IX_sections_4_5_through_4_7.tex
% Expose IX, sections 4.5--4.7.5.
% French transcription lines 1173--1270; authority physical pp. 375--380
% (source folios 363--368); hard stop before section 5 on physical p. 381.
% Physical pp. 375--380 were checked directly in four overlapping 1100-dpi
% bands. Two one-character/source-logic questions in 4.7.3 were additionally
% inspected in direct 5000-dpi crops.
%
% Transparent source corrections:
% - The source transcription reads I where the printed page has Gamma for
%   the group acting on the dual of L in the proof of Corollary 4.6.
% - In 4.7.3 the print says that E is generated by v but then gives it the
%   finite-order polynomial T^(p^N)-1 belonging to u. Since u=1-nv, the
%   corresponding polynomial in v is (1-nT)^(p^N)-1; that relation is used
%   below. The 5000-dpi crop confirms that the printed letter is v.
% - The print uses n both for the torsion integer and in the displayed root
%   order p^n, but immediately continues with "u != 1, i.e. m >= 1".
%   The root order is therefore rendered p^m, keeping n for the torsion
%   integer throughout.

\medskip
\noindent
\SGARefTarget{sga7.sga7i.eix.n4-5}{4.5}. We shall give an application of \SGARefLink{sga7.sga7i.eix.n4-3}{4.3} a) to the theory of the local
conductor of abelian schemes.  Suppose that $k$ is algebraically
closed.\footnote{It would in fact suffice for $k$ to be perfect, but,
after passing to the strict henselization of $S$, one is immediately
reduced to the case considered here.}  Recall (SGA 5 X 6.2, or [20])
that, for every commutative etale $\ell$-group scheme $F_K$ over $K$
($\ell\neq p$), one defines an integer
\begin{equation}\tag{\SGARefTarget{sga7.sga7i.eix.n4-5-1}{4.5.1}}
 \alpha_s(F_K)=\alpha_s(i_*F_K),
\end{equation}
where $i:\eta\longrightarrow S$ denotes the inclusion, as follows.  One
chooses a finite Galois extension $K'$ of $K$, with group $\Gamma$,
that splits $F_K$, so that $F_K$ is represented by a finite
$\ell$-group $M$ on which $\Gamma$ acts.  On the other hand, one
associates with the extension $K'/K$ a certain projective
$\mathbb Z_\ell[\Gamma]$-module, defined up to isomorphism,
\begin{equation}\tag{\SGARefTarget{sga7.sga7i.eix.n4-5-2}{4.5.2}}
 \operatorname{Sw}=\operatorname{Sw}(K'/K,\mathbb Z_\ell),
\end{equation}
and sets
\begin{equation}\tag{\SGARefTarget{sga7.sga7i.eix.n4-5-3}{4.5.3}}
 \alpha_s(F_K)
 =\operatorname{length}
   \operatorname{Hom}_{\mathbb Z_\ell(\Gamma)}(\operatorname{Sw},M).
\end{equation}
One checks (loc.\ cit.) that this expression is independent of the
choice of the extension $K'/K$.  Moreover, the projectivity of
$\operatorname{Sw}$ as a $\mathbb Z_\ell[\Gamma]$-module implies that
$\alpha_s(F_K)$ is additive in exact sequences.

We apply this definition to the sheaf ${}_\ell A_K$, where $A_K$ is our
abelian scheme over $K$.  In this case one therefore obtains an integer
\begin{equation}\tag{\SGARefTarget{sga7.sga7i.eix.n4-5-4}{4.5.4}}
 \alpha_s(A_K;\ell)
 =\operatorname{length}
   \operatorname{Hom}_{\mathbb Z_\ell(\Gamma)}
   \bigl(\operatorname{Sw},
         U_\ell\otimes_{\mathbb Z_\ell}\mathbb F_\ell\bigr),
 \qquad
 U_\ell=T_\ell(A(\overline K)),
\end{equation}
where $\Gamma$ is chosen as above (a choice that a priori depends on
$\ell$).  The integer (\SGARefLink{sga7.sga7i.eix.n4-5-4}{4.5.4}) is called the \emph{wild conductor
exponent} of the abelian scheme $A_K$ relative to the discrete valuation
ring $V\subset K$.\footnote{One should also define the \emph{tame
conductor exponent} of $A_K$ relative to $V$ as the difference between
the lengths over $\mathbb Z/\ell\mathbb Z$ of the geometric generic and
special fibres of ${}_\ell A^o$.  It is independent of $\ell$ and equals
$2\lambda+\mu$ in the notation of (1.1.11).  The (total) conductor
exponent of $A_K$ is the sum of the tame and wild exponents; it is the
local term occurring in the Euler--Poincare formula for ${}_\ell A^o$ in
SGA 5 X 6.2 and [20].}  After passage to the henselization it is of
course defined without assuming that $V$ is henselian, and one is
immediately reduced to the case in which $S$ is strictly local.  We
intend to prove the following result.

\medskip
\noindent
\underline{Corollary} \SGARefTarget{sga7.sga7i.eix.n4-6}{4.6}. \underline{The wild conductor exponent
(\SGARefLink{sga7.sga7i.eix.n4-5-4}{4.5.4}) is independent of the choice of the prime $\ell\neq p$.}

We may suppose that $S$ is strictly local.  We compute (\SGARefLink{sga7.sga7i.eix.n4-5-4}{4.5.4}), replacing
$\operatorname{Sw}$ by any other projective
$\mathbb Z_\ell(\Gamma)$-module $L$.  Observe that the filtration
(\SGARefLink{sga7.sga7i.eix.n4-1-1}{4.1.1}) of $U_\ell$ consists of $\mathbb Z_\ell$-submodules that are
manifestly direct summands.  It therefore induces a corresponding
filtration of
\[
 U(\ell)=U_\ell\otimes_{\mathbb Z_\ell}\mathbb F_\ell,
\]
whose successive factors are obtained, by tensoring with
$\mathbb F_\ell$, from the factors $V_i$ $(i=1,2,3)$
\[
 U^{\mathrm{et}},\qquad
 U^{\mathrm{ef}}/U^{\mathrm{et}},\qquad
 U/U^{\mathrm{ef}}
\]
appearing in (\SGARefLink{sga7.sga7i.eix.n4-1-1}{4.1.1}).  This filtration is, of course, stable under
$\Gamma$.  Additivity thus decomposes
\[
 \operatorname{length}
 \operatorname{Hom}_{\mathbb Z_\ell(\Gamma)}(L,U(\ell))
\]
as a sum of three terms, the lengths of the groups
\begin{equation}\tag{*}
 \operatorname{Hom}_{\mathbb Z_\ell(\Gamma)}
   (L,V_i\otimes_{\mathbb Z_\ell}\mathbb F_\ell)
 \simeq
 \operatorname{Hom}_{\mathbb Z_\ell(\Gamma)}(L,V_i)
   \otimes_{\mathbb Z_\ell}\mathbb F_\ell.
\end{equation}
The isomorphism in (*) again follows from the assumed projectivity of
$L$; it is only necessary to choose $K'$ large enough for $\Gamma$ to
act on the $V_i$.  Now
$\operatorname{Hom}_{\mathbb Z_\ell(\Gamma)}(L,V_i)$ is plainly a
torsion-free $\mathbb Z_\ell$-module, so its rank is the length of the
right-hand side of (*).  If $\varphi_i$ is the character of the
representation of $\Gamma$ on $V_i$, and $\psi$ the character of the
representation of $\Gamma$ on $\check L$, we obtain
\begin{equation}\tag{\SGARefTarget{sga7.sga7i.eix.n4-6-1}{4.6.1}}
 \operatorname{length}
 \operatorname{Hom}_{\mathbb Z_\ell(\Gamma)}(L,U(\ell))
 =
 \sum_{1\leq i\leq3}N^{-1}
 \sum_{g\in\Gamma}\psi(g)\varphi_i(g)
\end{equation}
\[
 =N^{-1}\sum_{g\in\Gamma}\psi(g)\varphi(g),
 \qquad N=\operatorname{card}\Gamma,
\]
where $\varphi(g)=\sum_i\varphi_i(g)$ is also the function on $\Gamma$
obtained, by passage to the quotient, from the character of the
monodromy representation of the inertia group $I$ on $U$.  By \SGARefLink{sga7.sga7i.eix.n4-3}{4.3} a),
this character is an integral-valued rational function independent of
$\ell$.  Taking $L=\operatorname{Sw}$, the same is true of the character
$\psi(g)$ of $L$, which is
\[
 a(g)-N\delta(g)+1,
\]
where $\delta$ is the Dirac function on $\Gamma$ and $a$ is the Artin
conductor of $K'/K$ [26].  This proves \SGARefLink{sga7.sga7i.eix.n4-6}{4.6} and at the same time gives the
following expression, independent of $\ell$, for the wild conductor
exponent:
\begin{equation}\tag{\SGARefTarget{sga7.sga7i.eix.n4-6-2}{4.6.2}}
 \alpha_s(A_K)
 =N^{-1}\sum_{g\in\Gamma}
   \bigl(a(g)-N\delta(g)+1\bigr)\varphi(g).
\end{equation}
Here $\Gamma$ is the Galois group of a finite Galois extension $K'$ of
$K$ that splits $U^{\mathrm{ef}}$ (for example the smallest such
extension; cf.\ \SGARefLink{sga7.sga7i.eix.n4-1}{4.1}), $a(g)$ is the Artin conductor on $\Gamma$, and
$\varphi(g)$ is the trace function on $\Gamma$ described in (\SGARefLink{sga7.sga7i.eix.n4-3-0}{4.3.0}).
In the semistable-reduction case one may take $K'=K$, hence
$\Gamma=\{e\}$, and obtains
\begin{equation}\tag{\SGARefTarget{sga7.sga7i.eix.n4-6-3}{4.6.3}}
 \alpha_s(A_K)=0
 \qquad\text{(semistable-reduction case).}
\end{equation}

To conclude, let us state explicitly a criterion for semistable
reduction that follows from the general considerations of \SGARefLink{sga7.sga7i.eix.n4-2}{4.2}.

\medskip
\noindent
\underline{Proposition} \SGARefTarget{sga7.sga7i.eix.n4-7}{4.7}. \underline{(M.\ Raynaud's criterion for
semistable reduction.)  Let $S$ be a trait, let $A_K$ be an abelian
scheme over its function field $K$, and let $n>1$ be an integer prime
to the residual characteristic $p$.  Suppose that ${}_nA_K$ is
unramified relative to $S$, that is, that the inertia group acts
trivially on ${}_nA(\overline K)$.  Then either $A_K$ has semistable
reduction over $S$, or $n=2$; in the latter case, if $S$ is strictly
local, $A_K$ acquires semistable reduction over a finite Galois extension
$K'$ of $K$ whose Galois group is isomorphic to
$(\mathbb Z/2\mathbb Z)^r$.}\footnote{As follows from an example
communicated to me by J.\ Tate, \SGARefLink{sga7.sga7i.eix.n4-7}{4.7} is no longer valid when $n$ equals
the residual characteristic $p$, even if $K$ has characteristic zero
and $A_K$ is an elliptic curve with complex multiplication.}

We may plainly suppose that $S$ is strictly local.  With the notation
of \SGARefLink{sga7.sga7i.eix.n4-1}{4.1} and \SGARefLink{sga7.sga7i.eix.n4-2}{4.2}, take $K'$ as small as possible; in other words, suppose
that $\Gamma$ acts faithfully on $A^{*o}_o$.  By hypothesis $\Gamma$
acts trivially on ${}_nA(\overline K)$, and hence a fortiori on
\[
 {}_nA(\overline K)^{\mathrm{ef}}
 \simeq{}_nA^*_o(k)\supset{}_nA^{*o}_o(k).
\]
The conclusion follows from the next two lemmas.

\medskip
\noindent
\underline{Lemma} \SGARefTarget{sga7.sga7i.eix.n4-7-1}{4.7.1}. \underline{(J.-P.\ Serre.)  Let $G$ be a
commutative group scheme over a field $k$, which is an extension of an
abelian scheme by a torus, and let $n\geq2$ be an integer.  Every
finite-order automorphism of $G$ that induces the identity on ${}_nG$ is
the identity if $n\neq2$, and is either the identity or has order two if
$n=2$.}

\medskip
\noindent
\underline{Lemma} \SGARefTarget{sga7.sga7i.eix.n4-7-2}{4.7.2}. \underline{A finite group in which every
element distinct from $1$ has order $2$ is isomorphic to
$(\mathbb Z/2\mathbb Z)^r$ for a suitable $r\geq0$.}

\smallskip
\noindent
\SGARefTarget{sga7.sga7i.eix.n4-7-3}{4.7.3}. Let us recall the well-known proof of \SGARefLink{sga7.sga7i.eix.n4-7-1}{4.7.1}.  Since
$n\,\operatorname{id}_G$ is an epimorphism, the hypothesis that $u$
induces the identity on ${}_nG$, or equivalently that $1-u$ induces zero
on ${}_nG$, is expressed by the condition
\[
 1-u=nv
\]
in the endomorphism ring of $G$.  To prove the asserted conclusion, we
reduce, by decomposing the finite cyclic group generated by $u$ into its
$p$-primary components for the various primes $p$ dividing the order of
$u$, to the case in which $u$ has order $p^N$.

Let $E$ be the subring of $\operatorname{End}(G)$ generated by $v$.
Since $u=1-nv$ and $u^{p^N}=1$, it is a quotient of
\[
 \mathbb Z[T]\big/\bigl((1-nT)^{p^N}-1\bigr);
\]
it is therefore of finite type over $\mathbb Z$.  It is moreover
torsion-free (a formal consequence of the divisibility of $G$), and thus
embeds in
\[
 E_{\mathbb Q}=E\otimes_{\mathbb Z}\mathbb Q.
\]
The latter is a quotient, after the change of variable $u=1-nT$, of
$\mathbb Q[u]/(u^{p^N}-1)$, which is a direct product of finite field
extensions of $\mathbb Q$; the same is therefore true of $E_{\mathbb Q}$.
Considering the images of $E$, $u$, and $v$ in the component fields
$K_i$ of $E_{\mathbb Q}$, it remains to prove the following assertion.
Let $K=\mathbb Q(u)$ be generated by a primitive $p^m$-th root of unity
$u$, and let $n\geq2$ be an integer such that
\[
 u-1\equiv0\pmod n
\]
in the ring of integers of $K$.  Then $u=1$, or $u=-1$ and $n=2$.

This follows at once from the familiar fact that, if $u\neq1$, that is,
$m\geq1$, the ideal $(u-1)$ is prime and $(p)$ is a power of it, with
exponent greater than $1$ except when $p=2$ and $m=1$, that is, when
$u=-1$.  Since $(u-1)$ is prime and divides $p$, an integer $n\geq2$ can
divide $u-1$ only if $n=p$; and $p$ can divide $u-1$ only if
$(u-1)=(p)$, that is, if $u=-1$.  Hence $n=2$.

\smallskip
\noindent
\SGARefTarget{sga7.sga7i.eix.n4-7-4}{4.7.4}. Let us prove \SGARefLink{sga7.sga7i.eix.n4-7-2}{4.7.2}.  It is enough to prove that $\Gamma$ is
commutative.  If $x,y\in\Gamma$, then $(xy)(xy)=1$; since $x^2=1$ and
$y^2=1$, this may also be written
\[
 xy=(xy)^{-1}=y^{-1}x^{-1}=yx,
\]
as required.

\medskip
\noindent
\underline{Remark} \SGARefTarget{sga7.sga7i.eix.n4-7-5}{4.7.5}. The proof given in \SGARefLink{sga7.sga7i.eix.n4-7-3}{4.7.3}, together with
\SGARefLink{sga7.sga7i.eix.n4-7-2}{4.7.2}, shows that if $G$ is a divisible commutative group in a topos,
$n\geq2$ is an integer, and $\Gamma$ is a finite group acting faithfully
on $G$ while inducing the identity on ${}_nG$, then $\Gamma$ is the
trivial group if $n\geq3$, and is isomorphic to a group
$(\mathbb Z/2\mathbb Z)^r$ if $n=2$.

% END INCLUDED SGA7I_English_current_buildable_source_20260823/components/141_expose_IX_sections_4_5_through_4_7.tex

% BEGIN INCLUDED SGA7I_English_current_buildable_source_20260823/components/142_expose_IX_section_5_1.tex
% Expose IX, section 5.1.
% French transcription lines 1276--1348; authority physical pp. 381--383
% (source folios 369--371); hard stop before Theorem 5.2.
% Physical pp. 381--383 were checked directly in four overlapping
% 1100-dpi bands per page.  The print is clear throughout this unit, so
% no higher-resolution ambiguity crop was needed.
%
% Transparent transcription correction:
% - Formula (5.1.7) in the French transcription has Lambda^o, whereas
%   the printed authority and the surrounding definition both have A^o.

\medskip
\noindent
\underline{5. The orthogonality theorem in the case $\ell=p$
(semistable case).  Duality of the abelian schemes $B_o$ and $B'_o$.
Characterization of the fixed part.  Criteria for good reduction,
essentially good reduction, and semistable reduction.}

\medskip
\noindent
\SGARefTarget{sga7.sga7i.eix.n5-1}{5.1}. Let us place ourselves in the situation of \SGARefLink{sga7.sga7i.eix.n2-2-3}{2.2.3}, where $S$ is
henselian and $A^o$ is $\ell$-divisible, so that the \underline{fixed
part}
\[
 T_\ell(A^o)^f\lhook\joinrel\longrightarrow T_\ell(A)
\]
is defined as the pro-$\ell$ Barsotti--Tate group formed by the
$({}_{\ell^\nu}A^o)^f$.  We shall now give a definition of the toric
part (\SGARefLink{sga7.sga7i.eix.n2-3-2}{2.3.2}) of $T_\ell(A^o)$ which is independent of the hypothesis
$\ell\neq p$ made in \SGARefLink{sga7.sga7i.eix.n2-3}{2.3}, but which instead assumes that $S$ is
\underline{complete}.\footnote{This hypothesis is in fact not
essential, as is shown in \SGARefLink{sga7.sga7i.eix.n6-1}{6.1} below.}  For this purpose, for every
integer $n\geq0$ introduce the local Artinian scheme
\begin{equation}\tag{\SGARefTarget{sga7.sga7i.eix.n5-1-1}{5.1.1}}
 S_n=\operatorname{Spec}(V/\mathfrak m^{n+1}),
\end{equation}
the $n$-th infinitesimal neighbourhood of $s$ in $S$, and, for every
$S$-scheme $X$, put $X_n=X\times_S S_n$.  Since
$(A^o)_n=(A_n)^o$ is smooth over $S_n$, it is known (SGA 3 IX 3.6 bis)
that $T_o$ is the special fibre of a unique subtorus
\begin{equation}\tag{\SGARefTarget{sga7.sga7i.eix.n5-1-2}{5.1.2}}
 T_n\subset A_n^o.
\end{equation}
By the uniqueness of the $T_n$, they are obtained from one another by
the base changes $S_{n'}\longrightarrow S_n$.  In other words, they
define a \underline{formal subtorus} $\widehat T$ of the group
$\widehat A^o$ (in the category of formal schemes over $S$), the formal
completion of $A^o$ along the special fibre:
\begin{equation}\tag{\SGARefTarget{sga7.sga7i.eix.n5-1-3}{5.1.3}}
 \widehat T\subset\widehat A^o.
\end{equation}

Since the functor $X\longmapsto\widehat X$ (formal completion along the
special fibre) induces an equivalence between the category of finite
schemes over $S$ and that of finite formal schemes over $S$ (EGA III
4.8), we shall simply identify a finite scheme over $S$ with the
corresponding formal scheme.  With this identification, the definition
of the \emph{fixed parts} immediately gives, for every $\nu\geq0$,
\begin{equation}\tag{\SGARefTarget{sga7.sga7i.eix.n5-1-4}{5.1.4}}
 ({}_{\ell^\nu}A^o)^f={}_{\ell^\nu}(\widehat A^o),
\end{equation}
which may also be written as an identity of projective systems:
\begin{equation}\tag{\SGARefTarget{sga7.sga7i.eix.n5-1-5}{5.1.5}}
 T_\ell(A^o)^f=T_\ell(\widehat A^o).
\end{equation}
We therefore obtain a Barsotti--Tate pro-subgroup of the fixed part of
$T_\ell(A^o)$ by taking
$T_\ell(\widehat T)\subset T_\ell(\widehat A^o)$.  We accordingly
define the \underline{toric part} of $T_\ell(A^o)$ by
\begin{equation}\tag{\SGARefTarget{sga7.sga7i.eix.n5-1-6}{5.1.6}}
 T_\ell(A^o)^t
 \stackrel{\mathrm{def}}{=}
 T_\ell(\widehat T)
 \subset T_\ell(\widehat A^o)
 =T_\ell(A^o)^f.
\end{equation}
Thus, by definition,
\begin{equation}\tag{\SGARefTarget{sga7.sga7i.eix.n5-1-7}{5.1.7}}
 T_\ell(A^o)^t=({}_{\ell^\nu}\widehat T)_{\nu\geq0},
\end{equation}
the finite flat $S$-group ${}_{\ell^\nu}\widehat T$ itself being
described as the inductive limit of the finite flat $S_n$-groups
\begin{equation}\tag{\SGARefTarget{sga7.sga7i.eix.n5-1-8}{5.1.8}}
 {}_{\ell^\nu}(T_n),\qquad n\ \text{variable};
\end{equation}
that is, its affine algebra is described as the projective limit of the
affine algebras of the groups (\SGARefLink{sga7.sga7i.eix.n5-1-8}{5.1.8}).  Of course,
\begin{equation}\tag{\SGARefTarget{sga7.sga7i.eix.n5-1-9}{5.1.9}}
 \bigl(T_\ell(A^o)^t\bigr)_n=T_\ell(T_n).
\end{equation}
For $n=0$ this is precisely formula (\SGARefLink{sga7.sga7i.eix.n2-3-3}{2.3.3}), by which we characterized
the toric part (as a profinite etale group) in the case $\ell\neq p$.
This shows that definition (\SGARefLink{sga7.sga7i.eix.n5-1-6}{5.1.6}) is compatible with that given in
\SGARefLink{sga7.sga7i.eix.n2-3}{2.3} when $\ell\neq p$.

Passing to generic fibres, the fixed and toric parts of $T_\ell(A^o)$
define fixed and toric parts in $T_\ell(A_K)$:
\begin{equation}\tag{\SGARefTarget{sga7.sga7i.eix.n5-1-10}{5.1.10}}
 T_\ell(A_K)^t\subset T_\ell(A_K)^f\subset T_\ell(A_K).
\end{equation}
These are inclusions of pro-$\ell$ Barsotti--Tate groups.  Consider the
analogous filtration corresponding to the dual abelian scheme $A'_K$:
\begin{equation}\tag{\SGARefTarget{sga7.sga7i.eix.n5-1-11}{5.1.11}}
 T_\ell(A'_K)^t\subset T_\ell(A'_K)^f\subset T_\ell(A'_K).
\end{equation}
We then have the following result, which generalizes \SGARefLink{sga7.sga7i.eix.n2-4}{2.4} (at least in
the case where $S$ is complete, to which one could easily have reduced
in \SGARefLink{sga7.sga7i.eix.n2-4}{2.4}).


% END INCLUDED SGA7I_English_current_buildable_source_20260823/components/142_expose_IX_section_5_1.tex

% BEGIN INCLUDED SGA7I_English_current_buildable_source_20260823/components/143_expose_IX_theorem_5_2.tex
% Expose IX, Theorem 5.2 and proof through 5.2.6.
% French transcription lines 1350--1448; authority physical pp. 383--387
% (source folios 371--375); hard stop before section 5.3.
% Physical p. 383 was checked with the preceding unit, and pp. 384--387
% were checked directly in four overlapping 1100-dpi bands per page.
% The formulas and typography are clear throughout; no higher-resolution
% ambiguity crop was needed.
%
% Transparent transcription corrections:
% - The theorem statement twice corrupts the dual abelian scheme A'_K,
%   once as A_{K'} and once as Lambda'_K.  The printed authority has A'_K.

\medskip
\noindent
\underline{Theorem} \SGARefTarget{sga7.sga7i.eix.n5-2}{5.2}. \underline{Suppose that $S$ is a complete
trait, and let $A_K$ be an abelian scheme over $K$, $A'_K$ the dual
abelian scheme, and $\ell$ a prime such that the connected Neron model
$A^o$ is $\ell$-divisible (cf.\ \SGARefLink{sga7.sga7i.eix.n2-2-1}{2.2.1}), so that one then has the
filtrations (\SGARefLink{sga7.sga7i.eix.n5-1-10}{5.1.10}) and (\SGARefLink{sga7.sga7i.eix.n5-1-11}{5.1.11}) of the pro-groups
$T_\ell(A_K)$ and $T_\ell(A'_K)$.  Then the toric parts are
recovered from the fixed parts by the formula}
\begin{equation}\tag{\SGARefTarget{sga7.sga7i.eix.n5-2-1}{5.2.1}}
 T_\ell(A_K)^t
 =
 T_\ell(A_K)^f\cap\bigl(T_\ell(A'_K)^f\bigr)^\perp,
\end{equation}
\underline{and by the analogous formula obtained by interchanging the
roles of $A_K$ and $A'_K$, where $\perp$ denotes the orthogonal
complement with respect to the canonical pairing (\SGARefLink{sga7.sga7i.eix.n1-0-2}{1.0.2}).}

\underline{Finally, when $A_K$ has semistable reduction over $S$
(which follows from the other hypotheses in the case $\ell=p$), one
even has the formula}
\begin{equation}\tag{\SGARefTarget{sga7.sga7i.eix.n5-2-2}{5.2.2}}
 T_\ell(A_K)^t=\bigl(T_\ell(A'_K)^f\bigr)^\perp.
\end{equation}

\smallskip
\noindent
\SGARefTarget{sga7.sga7i.eix.n5-2-3}{5.2.3}. The meaning of formula (\SGARefLink{sga7.sga7i.eix.n5-2-1}{5.2.1}) is quite clear when $\ell$ is
prime to the characteristic of $K$: as in \SGARefLink{sga7.sga7i.eix.n2-4}{2.4}, it can then be
interpreted by passing to the ordinary $\mathbb Z_\ell$-modules, with
an action of
$\pi=\operatorname{Gal}(\overline K/K)$, formed by the
$\overline K$-points of the pro-groups involved.  Let us spell it out
in the general case, using the notion of the dual
\begin{equation}\tag{\SGARefTarget{sga7.sga7i.eix.n5-2-3-1}{5.2.3.1}}
 D(G)=\operatorname{Hom}\bigl(G,T_\ell(\mathbb G_m)\bigr)
\end{equation}
of a pro-$\ell$ Barsotti--Tate group $G=(G(\nu))$.  By definition,
this is the projective system of Cartier duals $D(G(\nu))$
(SGA 3 VII 3.3.1 and VIII 1), whose transition morphisms are defined in
the evident way (for example, by transposition from the inclusion
morphisms among the $G(\nu)$).  The duality theory of abelian
varieties tells us that the pairing mentioned in \SGARefLink{sga7.sga7i.eix.n5-2}{5.2} actually defines
an \underline{isomorphism}
\begin{equation}\tag{\SGARefTarget{sga7.sga7i.eix.n5-2-3-2}{5.2.3.2}}
 T_\ell(A'_K)\simeq D\bigl(T_\ell(A_K)\bigr),
\end{equation}
and hence, by composition, a morphism
\begin{equation}\tag{\SGARefTarget{sga7.sga7i.eix.n5-2-3-3}{5.2.3.3}}
 T_\ell(A_K)^f\longrightarrow D\bigl(T_\ell(A'_K)^f\bigr).
\end{equation}
Formula (\SGARefLink{sga7.sga7i.eix.n5-2-1}{5.2.1}) says that the kernel of the preceding homomorphism
(which a priori is a pro-subgroup of its first term) is precisely the
toric part $T_\ell(A_K)^t$.  Notice that this is an
\underline{isomorphism of pro-objects, and not of projective systems}:
we do not assert (and it would in general be false) that, for every
integer $\nu$, $T_\ell(A_K)^t(\nu)$ equals the kernel of
\[
 T_\ell(A_K)^f(\nu)
 \longrightarrow
 D\bigl(T_\ell(A'_K)^f\bigr)(\nu);
\]
the latter can a priori be larger.  We shall see below, however, that
it is not larger if $A_K$ has semistable reduction; by the hypotheses
made in \SGARefLink{sga7.sga7i.eix.n5-2}{5.2}, this is in particular the case when $\ell=p$.

\smallskip
\noindent
\SGARefTarget{sga7.sga7i.eix.n5-2-4}{5.2.4}. The proof of (\SGARefLink{sga7.sga7i.eix.n5-2-1}{5.2.1}) merely repeats, with practically no change,
the proof of \SGARefLink{sga7.sga7i.eix.n2-4}{2.4}.  Observe that the induced pairing
\begin{equation}\tag{\SGARefTarget{sga7.sga7i.eix.n5-2-4-1}{5.2.4.1}}
 T_\ell(A_K)^f\times T_\ell(A'_K)^f
 \longrightarrow T_\ell(\mathbb G_{mK})
\end{equation}
which gives rise to (\SGARefLink{sga7.sga7i.eix.n5-2-3-3}{5.2.3.3}) is induced by a pairing
\begin{equation}\tag{\SGARefTarget{sga7.sga7i.eix.n5-2-4-2}{5.2.4.2}}
 T_\ell(A^o)^f\times T_\ell({A'}^o)^f
 \longrightarrow T_\ell(\mathbb G_{mS})
\end{equation}
of pro-Barsotti--Tate groups over $S$, itself induced by the pairing
of pro-groups
\begin{equation}\tag{\SGARefTarget{sga7.sga7i.eix.n5-2-4-3}{5.2.4.3}}
 T_\ell(A^o)\times T_\ell({A'}^o)
 \longrightarrow T_\ell(\mathbb G_{mS})
\end{equation}
deduced from the biextension $W$ of $(A^o,{A'}^o)$ by
$\mathbb G_{mS}$.  Consequently, the pairing
\begin{equation}\tag{\SGARefTarget{sga7.sga7i.eix.n5-2-4-4}{5.2.4.4}}
 T_\ell(A_o^o)\times T_\ell({A'}_o^o)
 \longrightarrow T_\ell(\mathbb G_{mk}),
\end{equation}
obtained from (\SGARefLink{sga7.sga7i.eix.n5-2-4-2}{5.2.4.2}) by passing to the special fibres, is precisely
the pairing deduced from the biextension $W_o$ of
$(A_o^o,{A'}_o^o)$ by $\mathbb G_{mk}$.  By VIII \SGARefLink{sga7.sga7i.eviii.n4-11}{4.11} and \SGARefLink{sga7.sga7i.eviii.n1-5}{1.5},
the left annihilator of the pairing (\SGARefLink{sga7.sga7i.eix.n5-2-4-4}{5.2.4.4}) is therefore
$T_\ell(T_o)$, which is precisely the special fibre of the toric part
$T_\ell(A^o)^t$.  Returning to the definition of left annihilators
given in \SGARefLink{sga7.sga7i.eix.n5-2-3}{5.2.3}, we conclude by means of the following lemma.

\medskip
\noindent
\underline{Lemma} \SGARefTarget{sga7.sga7i.eix.n5-2-5}{5.2.5}. \underline{Let $u:G\longrightarrow H$ be a
homomorphism of pro-$\ell$ Barsotti--Tate groups over $S$.  If
$u_o:G_o\longrightarrow H_o$ is zero, then $u$ is zero.  If
$u_o:G_o\longrightarrow H_o$ is an isogeny (respectively an
isomorphism), then $u$ is an isogeny (respectively an isomorphism).}

Indeed, apply the first assertion of \SGARefLink{sga7.sga7i.eix.n5-2-5}{5.2.5} to the composite morphism
\[
 T_\ell(A^o)^t
 \longrightarrow T_\ell(A^o)^f
 \longrightarrow D\bigl(T_\ell({A'}^o)^f\bigr),
\]
where the second morphism is deduced from the pairing (\SGARefLink{sga7.sga7i.eix.n5-2-4-2}{5.2.4.2}).  The
composite is zero, so the pairing under consideration vanishes on the
toric part on the left.  By symmetry it also vanishes on the toric part
on the right, and hence descends to a pairing on the quotients, or
\emph{abelian pieces},
\[
 T_\ell(A^o)^{ab}\times T_\ell({A'}^o)^{ab}
 \longrightarrow T_\ell(\mathbb G_m).
\]
These quotients are again pro-Barsotti--Tate groups, as one checks
readily.  It remains to prove that the morphism corresponding to this
pairing,
\[
 T_\ell(A^o)^{ab}
 \longrightarrow D\bigl(T_\ell({A'}^o)^{ab}\bigr),
\]
is an isogeny, which again follows from \SGARefLink{sga7.sga7i.eix.n5-2-5}{5.2.5}.

We leave to the reader the easy proof of \SGARefLink{sga7.sga7i.eix.n5-2-5}{5.2.5}, a result valid over any
connected Noetherian base.  Moreover, below in \SGARefLink{sga7.sga7i.eix.n7-4-7}{7.4.7}, by a somewhat
more precise analysis of the biextensions $W_n$ of
$(A_n^o,{A'}_n^o)$, we shall find a quite evident direct proof of
the two special cases in which we had to use \SGARefLink{sga7.sga7i.eix.n5-2-5}{5.2.5}.

\smallskip
\noindent
\SGARefTarget{sga7.sga7i.eix.n5-2-6}{5.2.6}. Suppose now that $A_K$ has semistable reduction over $S$,
so that $\ell$ may be any prime whatsoever (equal or unequal to
$p$).  Moreover, by the Galois criterion for semistable reduction
\SGARefLink{sga7.sga7i.eix.n3-5}{3.5}, for $\ell\neq p$ one has
\begin{equation}\tag{\SGARefTarget{sga7.sga7i.eix.n5-2-7}{5.2.7}}
 \bigl(T_\ell(A'_K)^f\bigr)^\perp
 \ \bigl(=T_\ell(A_K)^t\bigr)
 \subset T_\ell(A_K)^f.
\end{equation}
To finish the proof of \SGARefLink{sga7.sga7i.eix.n5-2}{5.2}, it remains to show that this inclusion is
valid without restriction on $\ell$.

Indeed, the left-hand term is canonically isomorphic to
\[
 D\bigl(T_\ell(A'_K)/T_\ell(A'_K)^f\bigr)
\]
and therefore has \emph{rank} (or \emph{height})
$2n-2\alpha-\mu=\mu$, in the notation of (\SGARefLink{sga7.sga7i.eix.n2-1-9}{2.1.9}), since by
hypothesis $\lambda=0$.  This is also the rank of
$T_\ell(A_K)^t$.  Since the inclusion
\begin{equation}\tag{*}
 T_\ell(A_K)^t\subset T_\ell(A_K)
\end{equation}
is \underline{strict}, in the sense that it already induces a
monomorphism on every term of the projective systems under
consideration, it follows readily that the left-hand term of (*) is
maximal among the pro-Barsotti--Tate subgroups of $T_\ell(A_K)$
having the same rank.  Thus equality holds in (\SGARefLink{sga7.sga7i.eix.n5-2-7}{5.2.7}); by the
orthogonality theorem this is equivalent to the inclusion between the
two extreme terms of (\SGARefLink{sga7.sga7i.eix.n5-2-7}{5.2.7}).


% END INCLUDED SGA7I_English_current_buildable_source_20260823/components/143_expose_IX_theorem_5_2.tex

% BEGIN INCLUDED SGA7I_English_current_buildable_source_20260823/components/144_expose_IX_sections_5_3_and_5_4.tex
% Expose IX, sections 5.3--5.4.
% French transcription lines 1450--1499; authority physical pp. 388--390
% (source folios 376--378); hard stop before section 5.5.
% All three pages were checked directly in four overlapping 1100-dpi
% bands.  The print is clear throughout; no higher-resolution ambiguity
% crop was needed.
%
% Transparent source correction:
% - The printed statement of Theorem 5.4 calls A_K an abelian scheme
%   over S.  Its subscript, the surrounding setup, and the reduction
%   language require an abelian scheme over K with semistable reduction
%   over S; the English states this type-correct formulation.

\medskip
\noindent
\SGARefTarget{sga7.sga7i.eix.n5-3}{5.3}. Using the orthogonality correspondence between strict
pro-Barsotti--Tate subgroups (in the sense specified above) of the dual
pro-groups $T_\ell(A_K)$ and $T_\ell(A'_K)$, one formally deduces
from (\SGARefLink{sga7.sga7i.eix.n5-2-2}{5.2.2}) that the pairing (\SGARefLink{sga7.sga7i.eix.n1-0-2}{1.0.2}) between these pro-groups induces a
pairing
\begin{equation}\tag{\SGARefTarget{sga7.sga7i.eix.n5-3-2}{5.3.2}}
 T_\ell(A_K)^{ab}\times T_\ell(A'_K)^{ab}
 \longrightarrow T_\ell(\mathbb G_{mK})
\end{equation}
which is a \underline{perfect duality}, that is, it identifies each of
the two factors on the left with the dual of the other (in the sense of
\SGARefLink{sga7.sga7i.eix.n5-2-3}{5.2.3}).\footnote{The superscripts $ab$ denote the \emph{abelian
parts}, the quotient of the fixed part by the toric part.}  It is also
clear that (\SGARefLink{sga7.sga7i.eix.n5-3-2}{5.3.2}) is induced by the pairing
\begin{equation}\tag{\SGARefTarget{sga7.sga7i.eix.n5-3-3}{5.3.3}}
 T_\ell(A^o)^{ab}\times T_\ell(A^{\prime o})^{ab}
 \longrightarrow T_\ell(\mathbb G_{mS})
\end{equation}
obtained from the pairing (\SGARefLink{sga7.sga7i.eix.n5-2-4-2}{5.2.4.2}) by passing to the quotients.  Since
this is a perfect duality on the generic fibres, that is, since the
corresponding homomorphism
\begin{equation}\tag{\SGARefTarget{sga7.sga7i.eix.n5-3-4}{5.3.4}}
 T_\ell(A^o)^{ab}
 \longrightarrow D\bigl(T_\ell(A^{\prime o})^{ab}\bigr)
\end{equation}
is an isomorphism on the generic fibres, one naturally wishes to deduce
that it is so over the whole of $S$, or equivalently on the special
fibres.  This deduction is trivially valid if $\ell\neq p$.  It is
also valid, at least in unequal characteristics when $\ell=p$, by a
deep theorem of Tate [33, Th.\ 4].  We shall nevertheless establish it
directly below in \SGARefLink{sga7.sga7i.eix.n7-4-6}{7.4.6}, without using Tate's theorem.

To make precise the geometric meaning of the result we wish to prove,
observe that the pairing obtained from (\SGARefLink{sga7.sga7i.eix.n5-3-3}{5.3.3}) by passage to the
special fibres is precisely the pairing
\begin{equation}\tag{\SGARefTarget{sga7.sga7i.eix.n5-3-5}{5.3.5}}
 T_\ell(B_o)\times T_\ell(B'_o)
 \longrightarrow T_\ell(\mathbb G_{mk})
\end{equation}
deduced from the biextension $W_o$ of
$(A_o^o,A_o^{\prime o})$ by passage to the quotient, where
$B_o$ and $B'_o$ denote respectively the abelian parts of
$A_o^o$ and $A_o^{\prime o}$.  By VIII \SGARefLink{sga7.sga7i.eviii.n4-8}{4.8}, the biextension
$W_o$ comes from a biextension $W_o^{ab}$ of
$(B_o,B'_o)$ by $\mathbb G_{mk}$, unique up to a unique
isomorphism.  Thus the pairing (\SGARefLink{sga7.sga7i.eix.n5-3-5}{5.3.5}) may also be interpreted as the
one deduced from the divisorial correspondence $W_o^{ab}$, and the
corresponding homomorphism
\begin{equation}\tag{\SGARefTarget{sga7.sga7i.eix.n5-3-6}{5.3.6}}
 T_\ell(B_o)
 \longrightarrow D\bigl(T_\ell(B'_o)\bigr)
 \simeq T_\ell\bigl(D(B'_o)\bigr)
\end{equation}
(where, in the last term, $D$ denotes the dual abelian scheme) is
precisely the morphism obtained by applying $T_\ell$ functorially to
the homomorphism
\begin{equation}\tag{\SGARefTarget{sga7.sga7i.eix.n5-3-7}{5.3.7}}
 B_o\longrightarrow D(B'_o)
\end{equation}
defined by $W_o^{ab}$.  Consequently, by the standard dictionary
(cf.\ [5]), saying that (5.3.5) is a perfect duality is also saying that
the kernel of (\SGARefLink{sga7.sga7i.eix.n5-3-7}{5.3.7}), which is a finite group over $k$ by \SGARefLink{sga7.sga7i.eix.n1-5}{1.5}, has
zero $\ell$-primary component.  Saying that this holds for every
$\ell$ therefore means that the isogeny (\SGARefLink{sga7.sga7i.eix.n5-3-7}{5.3.7}) is an isomorphism,
that is, that the divisorial correspondence $W_o^{ab}$ is a perfect
duality between $B_o$ and $B'_o$.  The preceding considerations
show only that the kernel of (\SGARefLink{sga7.sga7i.eix.n5-3-7}{5.3.7}) is a $p$-group.  We therefore
state the more complete result.

\medskip
\noindent
\underline{Theorem} \SGARefTarget{sga7.sga7i.eix.n5-4}{5.4}. \underline{Let $S$ be a trait, and let
$A_K$ be an abelian scheme over $K$ having semistable reduction
over $S$.  Let $B_o$ and $B'_o$ denote respectively the
abelian parts of the Neron models of $A_K$ and of the dual abelian
scheme $A'_K$.  Then the canonical divisorial correspondence
$W_o^{ab}$ on $(B_o,B'_o)$, deduced by VIII \SGARefLink{sga7.sga7i.eviii.n4-8}{4.8} from the
canonical biextension $W_o$ of $(A_o^o,A_o^{\prime o})$ (the
special fibre of the biextension $W$ of \SGARefLink{sga7.sga7i.eix.n1-4}{1.4}), is a perfect duality
between $B_o$ and $B'_o$; that is, it defines an isomorphism
of each of these abelian schemes with the dual of the other.}

The proof of \SGARefLink{sga7.sga7i.eix.n5-4}{5.4} in the general case will be given in section 7, as a
consequence of a very natural and important construction due to
Raynaud.  Let us merely observe here that, in proving \SGARefLink{sga7.sga7i.eix.n5-4}{5.4}, one may
assume that $S$ is complete, by \SGARefLink{sga7.sga7i.eix.n3-3}{3.3}.

% END INCLUDED SGA7I_English_current_buildable_source_20260823/components/144_expose_IX_sections_5_3_and_5_4.tex

% BEGIN INCLUDED SGA7I_English_current_buildable_source_20260823/components/145_expose_IX_sections_5_5_and_5_6.tex
\par\medskip
\noindent
\SGARefTarget{sga7.sga7i.eix.n5-5}{5.5}. To express conveniently the extreme factors in the filtrations
(\SGARefLink{sga7.sga7i.eix.n5-1-10}{5.1.10}) and (\SGARefLink{sga7.sga7i.eix.n5-1-11}{5.1.11}), in the case of semistable reduction, let us
introduce the twisted constant group schemes over $S$
\begin{equation}\tag{\SGARefTarget{sga7.sga7i.eix.n5-5-1}{5.5.1}}
 \underline M'_o=D(T_o),\qquad \underline M_o=D(T'_o),
\end{equation}
dual to the maximal tori $T_o,T'_o$ of the special fibres of $A,A'$
(SGA 3, X 5.7). Since $S$ is complete local, these group schemes may
be regarded as the special fibres of twisted constant schemes
\begin{equation}\tag{\SGARefTarget{sga7.sga7i.eix.n5-5-2}{5.5.2}}
 \underline M',\qquad \underline M\quad\hbox{over }S.
\end{equation}
Thus, if the torus $T_o$ is split (for example, if $k$ is separably
closed), that is, if $\underline M'_o$ is the constant group defined
by the ordinary group
\begin{equation}\tag{\SGARefTarget{sga7.sga7i.eix.n5-5-3}{5.5.3}}
 M'=\operatorname{Hom}_{k\text{-}\mathrm{gr}}(T_o,\mathbb G_{mk}),
\end{equation}
then $\underline M'$ is the constant group defined by $M'$:
\begin{equation*}
 \underline M'=M'_S.
\end{equation*}
One may moreover observe that, since $T_o$ and $T'_o$ are isogenous
(\SGARefLink{sga7.sga7i.eix.n2-2-6}{2.2.6}), $T'_o$ is split whenever $T_o$ is. In this case one should
therefore also introduce
\begin{equation}\tag{\SGARefTarget{sga7.sga7i.eix.n5-5-4}{5.5.4}}
 M=\operatorname{Hom}_{k\text{-}\mathrm{gr}}(T'_o,\mathbb G_{mk}),
\end{equation}
and then
\begin{equation*}
 \underline M=M_S.
\end{equation*}
With this notation, for every integer $n$ there are evidently
isomorphisms
\begin{equation}\tag{\SGARefTarget{sga7.sga7i.eix.n5-5-5}{5.5.5}}
 \underline M'_n\simeq D(T_n),\qquad
 \underline M_n\simeq D(T'_n),
\end{equation}
where $T'_n$ is defined in terms of $A'$ as $T_n$ is in terms of $A$
(\SGARefLink{sga7.sga7i.eix.n5-1-2}{5.1.2}). Hence there follow canonical isomorphisms (where
$\mathbb Z_\ell(1)=T_\ell(\mathbb G_{mS})$)
\begin{equation}\tag{\SGARefTarget{sga7.sga7i.eix.n5-5-6}{5.5.6}}
 T_\ell(T_n)\simeq
 \check{\underline M}'_n\otimes_{\mathbb Z}\mathbb Z_\ell(1),
 \qquad
 T_\ell(T'_n)\simeq
 \check{\underline M}_n\otimes_{\mathbb Z}\mathbb Z_\ell(1),
\end{equation}
and, on passing to the limit over $n$, canonical isomorphisms
\begin{equation}\tag{\SGARefTarget{sga7.sga7i.eix.n5-5-7}{5.5.7}}
 \left\{
 \begin{aligned}
  T_\ell(A^o)^t=T_\ell(\widehat T)
   &\simeq \check{\underline M}'\otimes\mathbb Z_\ell(1),\\[1ex]
  T_\ell(A^{\prime o})^t=T_\ell(\widehat T')
   &\simeq \check{\underline M}\otimes\mathbb Z_\ell(1).
 \end{aligned}
 \right.
\end{equation}
In the case of semistable reduction, duality, together with the
orthogonality theorem \SGARefLink{sga7.sga7i.eix.n5-2}{5.2} and relation (\SGARefLink{sga7.sga7i.eix.n5-2-2}{5.2.2}), then gives canonical
isomorphisms
\begin{equation}\tag{\SGARefTarget{sga7.sga7i.eix.n5-5-8}{5.5.8}}
 \left\{
 \begin{aligned}
  T_\ell(A_K)/T_\ell(A_K)^f
   &\simeq \underline M_K\otimes_{\mathbb Z}\mathbb Z_\ell,\\[1ex]
  T_\ell(A'_K)/T_\ell(A'_K)^f
   &\simeq \underline M'_K\otimes_{\mathbb Z}\mathbb Z_\ell.
 \end{aligned}
 \right.
\end{equation}
Likewise, the same formula (\SGARefLink{sga7.sga7i.eix.n5-2-2}{5.2.2}) gives canonical isomorphisms
\begin{equation}\tag{\SGARefTarget{sga7.sga7i.eix.n5-5-9}{5.5.9}}
 \left\{
 \begin{aligned}
  T_\ell(A_K)/T_\ell(A_K)^t
   &\simeq D\bigl(T_\ell(A'_K)^f\bigr),\\[1ex]
  T_\ell(A'_K)/T_\ell(A'_K)^t
   &\simeq D\bigl(T_\ell(A_K)^f\bigr),
 \end{aligned}
 \right.
\end{equation}
where one should recall that the groups $T_\ell(A_K)^f$ and
$T_\ell(A'_K)^f$ occurring on the right are the restrictions to
$\eta$ of the pro-$\ell$ Barsotti--Tate groups $T_\ell(A^o)^f$ and
$T_\ell(A^{\prime o})^f$, respectively, defined over $S$. Here $D$
denotes duality in the sense of pro-Barsotti--Tate groups, namely
$\underline{\operatorname{Hom}}(-,\mathbb Z_\ell(1))$. Thus we have:

\medskip
\noindent
\underline{Proposition} \SGARefTarget{sga7.sga7i.eix.n5-6}{5.6}. \underline{Suppose that $S$ is complete
and that $A_K$ has semistable reduction over $S$. Then, for every
prime number $\ell$ (which may equal the residue characteristic of
$S$), the pro-$\ell$ Barsotti--Tate group $T_\ell(A_K)^t$ is the
generic fibre of a pro-Barsotti--Tate group over $S$ ``of toric
type'', given by (\SGARefLink{sga7.sga7i.eix.n5-5-7}{5.5.7}), whereas
$T_\ell(A_K)/T_\ell(A_K)^f$ is the generic fibre of a
pro-Barsotti--Tate group over $S$ of pro-\'{e}tale type, given by
(\SGARefLink{sga7.sga7i.eix.n5-5-8}{5.5.8}). Finally, the groups $T_\ell(A_K)^f$ and
$T_\ell(A_K)/T_\ell(A_K)^t$ are the generic fibres of
pro-Barsotti--Tate groups defined over $S$, equal respectively to
$T_\ell(A^o)^f$ and to
$D(T_\ell(A^{\prime o})^f)=
\underline{\operatorname{Hom}}(T_\ell(A^{\prime o})^f,
\mathbb Z_\ell(1))$.}

\medskip
\noindent
\SGARefTarget{sga7.sga7i.eix.n5-6-1}{5.6.1}. Observe that, for an \'{e}tale covering $S'$ of $S$, splitting
$T_\ell(A_K)^t$ is equivalent to splitting
$T_\ell(A_K)/T_\ell(A_K)^f$: these conditions say respectively that
$S'$ splits the twisted constant groups $\underline M$ and
$\underline M'$, and these are isogenous because $T_o$ and $T'_o$
are. At the same time one sees that there is a smallest \'{e}tale
extension $S'$ of $S$ which does the job.

We also record the following easy consequence of the definitions in
\SGARefLink{sga7.sga7i.eix.n2-2-3}{2.2.3}.

% END INCLUDED SGA7I_English_current_buildable_source_20260823/components/145_expose_IX_sections_5_5_and_5_6.tex

% BEGIN INCLUDED SGA7I_English_current_buildable_source_20260823/components/146_expose_IX_proposition_5_7_through_corollary_5_10.tex
\par\medskip
\noindent
\underline{Proposition} \SGARefTarget{sga7.sga7i.eix.n5-7}{5.7}. \underline{Suppose that $A^o$ is
$\ell$-divisible and that $S$ is henselian. Let $F_K$ be a
constructible $\ell$-adic sheaf over $K$ which is ``unramified'', that
is, which extends to a twisted constant $\ell$-adic sheaf $F$ over
$S$, and let $u:F_K\longrightarrow T_\ell(A_K)$ be a homomorphism of
pro-groups. Then $u$ factors through a homomorphism from $F$ to
$T_\ell(A_K)^f$.}

Indeed, write
\begin{equation}\tag{\SGARefTarget{sga7.sga7i.eix.n5-7-1}{5.7.1}}
 u=(u(\nu))_{\nu\geqslant0},\qquad
 u(\nu):F_K(\nu)\longrightarrow T_\ell(A_K)(\nu)
 ={}_{\ell^\nu}A_K.
\end{equation}
By the universal property of $A$, since $F(\nu)$ is \'{e}tale and
hence smooth over $S$, the morphism $u(\nu)$ comes from a
homomorphism
\[
 v(\nu):F(\nu)\longrightarrow{}_{\ell^\nu}A.
\]
By definition of the fixed part, this homomorphism factors as
\[
 v(\nu):F(\nu)\longrightarrow({}_{\ell^\nu}A)^f.
\]
Since the image of the induced morphism on special fibres lies in the
infinitely $\ell$-divisible part of $A_o$, it follows that $v(\nu)$
in fact factors through $({}_{\ell^\nu}A^o)^f$. Passing to the limit
over $\nu$ gives a homomorphism
\[
 v:F\longrightarrow T_\ell(A^o)^f,
\]
and the induced morphism on generic fibres plainly gives the required
factorization.

\medskip
\noindent
\underline{Remarks} \SGARefTarget{sga7.sga7i.eix.n5-7-2}{5.7.2}. When $\ell\neq p$, property \SGARefLink{sga7.sga7i.eix.n5-7}{5.7} uniquely
characterizes the sub-pro-group $T_\ell(A_K)^f$ of $T_\ell(A_K)$ as
its largest unramified $\ell$-adic subsheaf. When $\ell=p$, the
statement of \SGARefLink{sga7.sga7i.eix.n5-7}{5.7} ceases to be very satisfactory. We shall give a
considerably stronger property, which can serve as an intrinsic
characterization of the fixed part of $T_p(A_K)$ (in terms of the
pro-Barsotti--Tate group structure of $T_p(A_K)$ alone), at least in
unequal characteristics.

\medskip
\noindent
\underline{Theorem} \SGARefTarget{sga7.sga7i.eix.n5-8}{5.8} (M. Raynaud). \underline{Let $A_K$ be an
abelian scheme with semistable reduction over the function field $K$
of a henselian trait $S$. Let $\ell$ be a prime number (the
interesting case being $\ell=p$), and consider the pro-$\ell$
Barsotti--Tate group $X_\eta=T_\ell(A_K)$ and its ``fixed part''
$X_\eta^f$, which is the generic fibre of the pro-$\ell$
Barsotti--Tate group $X^f=T_\ell(A^o)^f$. If $\ell=p$, suppose that
$\operatorname{char}K=0$ (so that Tate's theorem [33, Th.\ 4] is
available over $S$). Let $F$ be a pro-$\ell$ Barsotti--Tate group over
$S$. Then every homomorphism $F_\eta\longrightarrow X_\eta$ factors
through $X_\eta^f$.}

Passing to the ind-$\ell$ groups associated with the pro-$\ell$
Barsotti--Tate groups under consideration, the preceding statement
becomes equivalent, by the cited theorem of Tate, to the following.

\medskip
\noindent
\underline{Corollary} \SGARefTarget{sga7.sga7i.eix.n5-9}{5.9}. \underline{If $F$ is an ind-$\ell$
Barsotti--Tate group, then every homomorphism
$u_\eta:F_\eta\longrightarrow A_\eta$ extends uniquely to a
homomorphism $u:F\longrightarrow A$.}

In other words, the characteristic Neron property (\SGARefLink{sga7.sga7i.eix.n1-1-2}{1.1.2}) of the
Neron model remains valid for homomorphisms from ind-group schemes of
Barsotti--Tate type, with no smoothness hypothesis on them. (Only the
case of ind-\'{e}tale groups would follow directly from the
definition.)

\medskip
\noindent
\underline{Remark} \SGARefTarget{sga7.sga7i.eix.n5-9-1}{5.9.1}. In fact, M. Raynaud proves \SGARefLink{sga7.sga7i.eix.n5-9}{5.9} without
assuming that $A_K$ has semistable reduction, and even for every Neron
group scheme $A$ over $S$ (whose generic fibre need not be an abelian
scheme). Of course, the assumption that $S$ is henselian is also
unnecessary in the formulation of \SGARefLink{sga7.sga7i.eix.n5-9}{5.9}. Raynaud's proof of this more
general case reduces it to the special case considered in \SGARefLink{sga7.sga7i.eix.n5-8}{5.8}, using
an unpublished property in the construction of Neron models. Here we
shall give only Raynaud's proof of \SGARefLink{sga7.sga7i.eix.n5-9}{5.9}.

\medskip
\noindent
\underline{Proof of \SGARefLink{sga7.sga7i.eix.n5-9}{5.9}.} It is known [33] that $F$ is an extension
of an ind-\'{e}tale Barsotti--Tate group $F'$ by a Barsotti--Tate group
$F''$ with connected special fibre, so that for every integer $n$
there is an exact sequence
\[
 0\longrightarrow F''(n)\longrightarrow F(n)
 \longrightarrow F'(n)\longrightarrow0,
\]
with $F'(n)$ \'{e}tale and $F''(n)_o$ purely infinitesimal.
The given homomorphism $u_\eta$ is defined by a sequence of
homomorphisms
\[
 u_\eta(n):F(n)_\eta\longrightarrow A_\eta,
\]
and it is enough to extend these separately to homomorphisms
$u(n):F(n)\longrightarrow A$. Lemma \SGARefLink{sga7.sga7i.eix.n5-9-2}{5.9.2} below then reduces us to
proving that the homomorphism
$u''_\eta:F''_\eta\longrightarrow A_\eta$ induced by $u_\eta$
extends to $u'':F''\longrightarrow A$. Consequently, \underline{we
may suppose that $F$ has connected special fibre}.

We may regard $u_\eta$ as a homomorphism from $F_\eta$ to the
ind-$\ell$ Barsotti--Tate group $Y_\eta$ associated with
$X_\eta=T_\ell(A_K)$, and then consider the composite
\begin{equation}\tag{*}
 v_\eta:F_\eta\longrightarrow Y_\eta
 \longrightarrow Y_\eta/Y_\eta^f.
\end{equation}
By \SGARefLink{sga7.sga7i.eix.n5-6}{5.6}, the latter group extends to an ind-\'{e}tale
Barsotti--Tate group $Z$ over $S$, namely
\[
 \underline M\otimes_{\mathbb Z}(\mathbb Q_\ell/\mathbb Z_\ell).
\]
By Tate's theorem cited in \SGARefLink{sga7.sga7i.eix.n5-8}{5.8}, the homomorphism (*) then extends to a
homomorphism
\[
 v:F\longrightarrow Z.
\]
Since $F$ has connected special fibre and $Z$ is ind-\'{e}tale, this
morphism is zero. Hence $v_\eta$ is zero as well, and $u_\eta$
therefore factors through $Y_\eta^f$, as required.

It remains to prove the following lemma.

\medskip
\noindent
\underline{Lemma} \SGARefTarget{sga7.sga7i.eix.n5-9-2}{5.9.2}. Let $S$ be a trait, let $A$ be
a Neron group scheme over $S$ (\SGARefLink{sga7.sga7i.eix.n1-1-2}{1.1.2}), let
\[
 1\longrightarrow G'\longrightarrow G\longrightarrow G''
 \longrightarrow1
\]
be an extension of group schemes over $S$, with $G''$ smooth, let
$u:G'\longrightarrow A$ be a group homomorphism, and let
$v_\eta:G_\eta\longrightarrow A_\eta$ be a group homomorphism
extending $u_\eta$. Suppose that the extension of $G''$ by $A$
obtained from the extension $G$ by pushout along $u$ is representable
\footnote{Unpublished results of M. Raynaud imply, moreover, that this
condition is always satisfied.}---as is the case, for example, if
$G\longrightarrow G''$ admits a section locally for the ``finite
locally free'' topology (cf.\ SGA 3, IV 6.3) and if $A$ is
quasi-projective over $S$. Then there exists a unique group
homomorphism $v:G\longrightarrow A$ extending both $v_\eta$ and $u$.

Indeed, the last hypothesis reduces us to the case in which $G'=A$,
$u$ is the identity, and $v_\eta:G_\eta\longrightarrow A_\eta$ is a
projection which we seek to extend to a projection
$v:G\longrightarrow A$. But $G$ is smooth, being an extension of two
smooth group schemes, and the conclusion follows from the definition
of the Neron property of $A$.

\medskip
\noindent
\underline{Corollary} \SGARefTarget{sga7.sga7i.eix.n5-10}{5.10}. \underline{Let $A_K$ be an abelian
scheme over the function field $K$ of a trait $S$, and let $\ell$ be
a prime number. When $\ell=p$, suppose that
$\operatorname{char}K=0$, so that Tate's theorem already invoked in
\SGARefLink{sga7.sga7i.eix.n5-8}{5.8} is available. Then $A_K$ has good reduction over $S$ if and only
if the pro-$\ell$ Barsotti--Tate group $T_\ell(A_K)$ ``has good
reduction over $S$'', that is, extends to a Barsotti--Tate group over
$S$.}

(This is the result conjectured in [8, 4.9].)

Necessity being clear, it remains to prove sufficiency; observe that
the case $\ell\neq p$ has already been treated in \SGARefLink{sga7.sga7i.eix.n2-2-9}{2.2.9}. We may
suppose that $S$ is henselian (\SGARefLink{sga7.sga7i.eix.n2-2-9}{2.2.9}). First suppose that $A_K$ has
semistable reduction over $S$. Applying \SGARefLink{sga7.sga7i.eix.n5-8}{5.8} with $F=X$ gives
\[
 T_\ell(A_K)=T_\ell(A_K)^f.
\]
This also says that the maximal torus of $A_o$ is zero, that is, that
$A_o^o$ is an abelian scheme; the conclusion then follows from
\SGARefLink{sga7.sga7i.eix.n2-2-9}{2.2.9} (iii${}^o$).

We now treat the general case. By \SGARefLink{sga7.sga7i.eix.n2-2-9}{2.2.9} we may plainly assume that
$S$ is strictly local. Let $K'/K$ be a finite Galois extension such
that $A_{K'}$ has semistable reduction over the normalization $S'$ of
$S$ in $K'$ (\SGARefLink{sga7.sga7i.eix.n3-6}{3.6}). If $A^*$ denotes the corresponding Neron model,
$G=\operatorname{Gal}(K'/K)$ acts on the situation $(A^*,S')$, and
hence on $A_o^{*o}$. By \SGARefLink{sga7.sga7i.eix.n4-1}{4.1} it is enough to prove that $G$ acts
trivially on the latter; this indeed implies that $A_K$ has
semistable reduction, reducing us to the case already treated.

Let $X$ be a Barsotti--Tate group over $S$ extending
$X_\eta=T_\ell(A_K)$. By Tate's uniqueness theorem already invoked,
there is a unique isomorphism
\[
 X\times_S S'\simeq T_\ell(A^{*o})
\]
which induces the identity on generic fibres. By transport of
structure, this isomorphism is compatible with the actions of $G$.
Now $G$ plainly acts trivially on the special fibre of
$X\times_S S'$, and hence on $T_\ell(A_o^{*o})$. It follows at once
that $G$ acts trivially on $A_o^{*o}$, completing the proof.

% END INCLUDED SGA7I_English_current_buildable_source_20260823/components/146_expose_IX_proposition_5_7_through_corollary_5_10.tex

% BEGIN INCLUDED SGA7I_English_current_buildable_source_20260823/components/147_expose_IX_remark_5_11_through_corollary_5_15.tex
\medskip
\noindent
\underline{Remark} \SGARefTarget{sga7.sga7i.eix.n5-11}{5.11}. Once Tate's theorem has been proved without any
restrictive hypothesis on $\operatorname{char}K$, the preceding results
\SGARefLink{sga7.sga7i.eix.n5-8}{5.8}, \SGARefLink{sga7.sga7i.eix.n5-9}{5.9}, and \SGARefLink{sga7.sga7i.eix.n5-10}{5.10}, as well as \SGARefLink{sga7.sga7i.eix.n5-13}{5.13} below, will plainly hold without
such a restriction.

We record a definition which will also be useful elsewhere.

\medskip
\noindent
\underline{Definition} \SGARefTarget{sga7.sga7i.eix.n5-12}{5.12}. Let $S$ be a henselian trait, let $\ell$ be
a prime number, and let $X_\eta$ be a pro-$\ell$ Barsotti--Tate group
over the generic point $\eta$ of $S$. We say that $X_\eta$ has
\emph{good reduction over $S$} if it extends to a Barsotti--Tate group
over $S$ (cf. \SGARefLink{sga7.sga7i.eix.n5-10}{5.10}). We say that it has \emph{virtual good reduction
over $S$} if there exists a filtration
\[
 \bigl(\operatorname{Fil}^i(X_\eta)\bigr)_{i\geq0}
\]
by strict Barsotti--Tate subgroups\footnote{That is, subobjects as
projective systems, or equivalently subobjects term by term.} such that
the Barsotti--Tate groups
\[
 \operatorname{Gr}^i(X_\eta)=
 \operatorname{Fil}^i(X_\eta)/\operatorname{Fil}^{i+1}(X_\eta)
\]
have good reduction over $S$. If the filtration can be chosen to have
length at most $n$ (that is, if at most $n$ of the
$\operatorname{Gr}^i$ are nonzero), we say that $X_\eta$ has
\emph{virtual good reduction of level $n$ over $S$}.

Finally, we say that $X_\eta$ has \emph{essential good reduction over
$S$} (respectively, \emph{essential virtual good reduction},
respectively, \emph{essential virtual good reduction of level $n$}) if
there exists a finite extension $K'/K$ such that $X_{K'}$ has good
reduction (respectively, virtual good reduction, respectively, virtual
good reduction of level $n$) over the normalization $S'$ of $S$ in
$K'$.

\medskip
\noindent
\SGARefTarget{sga7.sga7i.eix.n5-12-1}{5.12.1}. When $\ell\neq p$, the categories of pro-$\ell$
Barsotti--Tate groups over $\eta$ and over $S$ may be interpreted in
terms of $\ell$-adic representations of
\[
 \pi=\pi_1(\eta), \qquad \pi_o=\pi_1(S),
\]
respectively. The six notions just introduced can then be expressed in
terms of the natural representation of the inertia group $I$ on
$X_\eta(\overline K)$: the representation is, respectively, trivial,
unipotent, unipotent of level $n$, essentially trivial, essentially
unipotent, or essentially unipotent of level $n$.

\medskip
\noindent
\SGARefTarget{sga7.sga7i.eix.n5-12-2}{5.12.2}. Take $X_\eta=T_\ell(A_K)$, where $A_K$ is an abelian scheme
over $K$. We saw in \SGARefLink{sga7.sga7i.eix.n5-6}{5.6} that, if $A_K$ has semistable reduction over
$S$, then $X_\eta$ has virtual good reduction of level 2 (and we shall
see the converse in a moment). Consequently, by the semistable
reduction theorem \SGARefLink{sga7.sga7i.eix.n3-6}{3.6}, $X_\eta$ has essential virtual good reduction of
level 2. This statement should be regarded as the ``right'' formulation
of the geometric monodromy theorem in Expos\'e III\footnote{The source
leaves the precise internal reference blank.} for $H^1$ of a proper
smooth scheme over $K$, without excluding the case $\ell=p$ (which was
not considered there). The geometric monodromy theorem---or, for the
present, conjecture---for the higher $H^i$ given in Expos\'e III should
likewise be reformulated so as to include ``$p$-adic'' coefficients, as
was observed in \SGARefLink{sga7.sga7i.eix.n0-4}{0.4}.

We now summarize the relations between the reduction properties of
abelian schemes over $K$ and those of the associated Barsotti--Tate
groups.

\medskip
\noindent
\underline{Proposition} \SGARefTarget{sga7.sga7i.eix.n5-13}{5.13}. Let $S$ be a henselian trait, let $A_K$
be an abelian scheme over its fraction field, and let $\ell$ be a prime
number. When $\ell=p$, assume that $K$ has characteristic zero
(cf. \SGARefLink{sga7.sga7i.eix.n5-11}{5.11}).

\smallskip
\noindent
a) The abelian scheme $A_K$ has good reduction over $S$ (\SGARefLink{sga7.sga7i.eix.n2-2-9}{2.2.9}) if and
only if $T_\ell(A_K)$ has good reduction over $S$ (\SGARefLink{sga7.sga7i.eix.n5-12}{5.12}).

\smallskip
\noindent
b) The abelian scheme $A_K$ has essential good reduction over $S$---that
is, there exists a finite extension $K'/K$ such that $A_{K'}$ has good
reduction over the normalization $S'$ of $S$ in $K'$---if and only if
$T_\ell(A_K)$ has essential good reduction over $S$ (\SGARefLink{sga7.sga7i.eix.n5-12}{5.12}). In that
case $K'$ may be chosen to be a finite separable extension of $K$.

\smallskip
\noindent
c) The following conditions are equivalent:
\begin{enumerate}
\item[1)] $A_K$ has semistable reduction over $S$ (3.4);
\item[2)] $T_\ell(A_K)$ has virtual good reduction of level 2 over $S$
  (\SGARefLink{sga7.sga7i.eix.n5-12}{5.12});
\item[3)] $T_\ell(A_K)$ has virtual good reduction over $S$.
\end{enumerate}

\noindent
\underline{Proof}. Part a) is just \SGARefLink{sga7.sga7i.eix.n5-10}{5.10} and has been included for
reference.

For b), the first assertion follows immediately from a) and the
definitions. To prove the final assertion, we may suppose that
$\ell\neq p$. In view of a), it then suffices to take for $K'$ the
finite separable extension of $K$ corresponding to an open subgroup
$\pi'$ of
\[
 \pi=\operatorname{Gal}(\overline K/K)
\]
whose intersection with the inertia group $I$ is contained in the
kernel of the representation of $I$ on
$T_\ell(A_K(\overline K))$; this kernel is open by hypothesis
(cf. \SGARefLink{sga7.sga7i.eix.n5-12-1}{5.12.1}).

For c), we have already noted that 1) implies 2) (\SGARefLink{sga7.sga7i.eix.n5-12-2}{5.12.2}), and it is
immediate that 2) implies 3). It remains to prove that 3) implies 1).
For $\ell\neq p$, this is precisely the Galois criterion for
semistable reduction \SGARefLink{sga7.sga7i.eix.n3-5}{3.5}. For $\ell=p$, argue as follows. Choose a
prime $\ell'\neq p$ and apply the considerations of \SGARefLink{sga7.sga7i.eix.n4-1}{4.1} to it. We may
suppose that $S$ is strictly local (\SGARefLink{sga7.sga7i.eix.n3-4-0}{3.4.0}), and choose a finite Galois
extension $K'/K$, with group $G$, such that $A_{K'}$ has semistable
reduction. We must prove that the representation of $G$ on each of the
three quotients occurring in (\SGARefLink{sga7.sga7i.eix.n4-1-1}{4.1.1}) is trivial. Equivalently, their
direct-sum representation is trivial, or again its trace is the
``trivial'' trace, that is, is constant. We saw in \SGARefLink{sga7.sga7i.eix.n4-3}{4.3}(a) that this
trace function is independent of the choice of $\ell'\neq p$, and in
\SGARefLink{sga7.sga7i.eix.n6-4}{6.4} below we indicate a method of calculating it in terms of
$T_p(A_K)$ as well. That calculation plainly gives a constant function
when $T_p(A_K)$ has virtual good reduction, and this completes the
proof.

\medskip
\noindent
\underline{Remark} \SGARefTarget{sga7.sga7i.eix.n5-13-1}{5.13.1}. It seems plausible that, if $X_\eta$ is a
pro-$\ell$ Barsotti--Tate group over $\eta$ having virtual good
reduction over $S$ and essential virtual good reduction of level $n$,
then it has virtual good reduction of level $n$. (This is true when
$\ell\neq p$.)

The following result, an easy consequence of \SGARefLink{sga7.sga7i.eix.n5-7}{5.7} and of the
semistable reduction theorem \SGARefLink{sga7.sga7i.eix.n3-6}{3.6}, was pointed out to me by
J.-P. Serre.

\medskip
\noindent
\underline{Proposition} \SGARefTarget{sga7.sga7i.eix.n5-14}{5.14} (Criterion for essential good reduction).
Let $S$ be a trait with function field $K$, let $A_K$ be an abelian
scheme over $K$, and let $\ell$ be a prime number distinct from
$\operatorname{char}K$.

\smallskip
\noindent
a) Suppose that $\ell\neq p=\operatorname{char}k$. Then the following
conditions are equivalent:
\begin{enumerate}
\item[(i)] $A_K$ has essential good reduction over $S$; that is, there
  exists a finite separable extension $K'/K$ such that $A_{K'}$ has
  good reduction over the normalization $S'$ of $S$ in $K'$;
\item[(ii)] the action of the inertia group $I$ on
  $U_\ell=T_\ell(A(\overline K))$ is essentially trivial; that is, it
  is trivial on an open subgroup of $I$, or equivalently factors
  through a finite quotient $\Gamma$ of $I$;
\item[(iii)] the action of $I$ on
  $U_\ell\otimes_{\mathbb Z_\ell}\mathbb Q_\ell$ is semisimple.
\end{enumerate}

\noindent
b) Suppose that $\ell=p$. If $n=\dim A_K\neq0$, then the image of $I$
in the automorphism group of $U_\ell$ is not finite; more precisely,
its image in the group $\mathbb Z_\ell^\times$ of automorphisms of
$\bigwedge^{2n}U_\ell$ has finite index. A sufficient (but by no means
necessary) condition for $A_K$ to have essential good reduction is
that the action of $I$ on
$U_\ell\otimes_{\mathbb Z_\ell}\mathbb Q_\ell$ be semisimple.

\noindent
\underline{Proof}. When $\ell\neq p$, the equivalence of (i) and (ii)
follows immediately from the good-reduction criterion \SGARefLink{sga7.sga7i.eix.n2-2-9}{2.2.9}, and (ii)
implies (iii) because $\mathbb Q_\ell$ has characteristic zero. The
homomorphism
\[
 \pi=\operatorname{Gal}(\overline K/K)
 \longrightarrow \mathbb Z_\ell^\times
 \simeq\operatorname{Aut}\!\left(\bigwedge^{2n}U_\ell\right)
\]
is $\chi^n$, where $\chi$ is induced by the character
\[
 \mathbb Z_\ell^\times=\operatorname{Aut}(\mathbb Z_\ell(1))
\]
arising from the action on the $\ell^\nu$-th roots of unity. It is
known that, when $\ell=p$, the image of $I$ under this character has
finite index [33]. The fact that the action of $I$ on
$U_p\otimes\mathbb Q_p$ need not be semisimple even when $A_K$ is an
elliptic curve with good reduction is exhibited in [28, IV A 2.4].

It remains to prove that (iii) implies (i), for every $\ell$. We use
the semistable reduction theorem \SGARefLink{sga7.sga7i.eix.n3-6}{3.6}, which immediately reduces us to
the case where $A_K$ has semistable reduction over $S$: hypothesis
(iii) is preserved upon replacing $K$ by a finite separable extension
$K'$, or equivalently $I$ by an open subgroup $I'$. When $\ell\neq p$,
the action of $I$ on $U_\ell\otimes\mathbb Q_\ell$ is unipotent
(\SGARefLink{sga7.sga7i.eix.n3-5}{3.5}(v)) and semisimple, hence trivial, and the good-reduction
criterion \SGARefLink{sga7.sga7i.eix.n2-2-9}{2.2.9} gives the result. We now give an argument valid with no
restriction on $\ell$, which yields a slightly sharper statement.

\medskip
\noindent
\underline{Corollary} \SGARefTarget{sga7.sga7i.eix.n5-15}{5.15}. With the notation of \SGARefLink{sga7.sga7i.eix.n5-14}{5.14}, suppose that
$A_K$ has semistable reduction over $S$, and that the fixed part
\[
 (U_{\ell\mathbb Q_\ell})^f
 \stackrel{\mathrm{def}}{=}(U_\ell^f)_{\mathbb Q_\ell}
\]
admits an $I$-stable complement in $U_{\ell\mathbb Q_\ell}$. Then
$A_K$ has good reduction.

It is enough to prove under these hypotheses that $U_\ell=U_\ell^f$.
By \SGARefLink{sga7.sga7i.eix.n5-6}{5.6} this says precisely that the toric part $T_o$ of $A_o^o$ is
zero, so that $A_o^o$ is an abelian scheme; hence $A$ is an abelian
scheme by \SGARefLink{sga7.sga7i.eix.n2-2-9}{2.2.9}. We may plainly suppose that $k$ is separably closed.
The intersection of the asserted complement in \SGARefLink{sga7.sga7i.eix.n5-15}{5.15} with $U_\ell$
then defines a pro-$\ell$ subgroup $F_K$ of $T_\ell(A_K)$ which is
isogenous to the cotoric part
\[
 T_\ell(A_K)/T_\ell(A_K)^f.
\]
The latter is pro-\'{e}tale and unramified by \SGARefLink{sga7.sga7i.eix.n5-6}{5.6}, and the same is
therefore true of $F_K$. It follows from \SGARefLink{sga7.sga7i.eix.n5-7}{5.7} that $F_K$ is contained
in $T_\ell(A_K)^f$; thus $F_K=0$, proving that $U_\ell$ is equal to
its fixed part.

% END INCLUDED SGA7I_English_current_buildable_source_20260823/components/147_expose_IX_remark_5_11_through_corollary_5_15.tex

% BEGIN INCLUDED SGA7I_English_current_buildable_source_20260823/components/148_expose_IX_section_6.tex
% Expose IX, section 6.
% French transcription lines 1783--1842; authority physical pp. 404--411
% (scan indices 403--410; source folios 392--399); hard stop before section 7.
% All eight physical pages were checked directly in four overlapping
% 1100-dpi bands per page.  The print, including its small subscripts,
% primes, infinity signs, and trace notation, is clear at that resolution;
% no unresolved feature warranted a higher-resolution ambiguity crop.
%
% Transparent transcription/source corrections:
% - Near the end of 6.1, the working transcription corrupts p^infinity as
%   p^(bar infinity); the print has an ordinary infinity sign.
% - The largest p-divisible subgroup is printed p^N A_o^o (the image of
%   multiplication by p^N), not the left-subscript p^N torsion group in the
%   working transcription.
% - Raynaud's example twice has type-level slips in print: it first calls
%   A_K simple over K' rather than K, and then calls K' a ramified quadratic
%   extension of S rather than of K.  The surrounding argument forces K.
% - Formula (6.3.2) clearly prints A'_{K'} in both orthogonal complements;
%   the working transcription has A'_k.  The following normalization is
%   likewise in K', not K.
% - The forward reference ``6.5'' in 6.4 cannot exist (section 7 follows
%   immediately); the notion of essential good reduction was defined in
%   5.12, which is the type-correct reference used below.
% - The final trace function is clearly on pi-prime; the working
%   transcription's pi-exclamation is a glyph error.

\medskip
\noindent
\underline{6. Remarks on the construction of the toric part and the fixed
part of} $T_p(A)$ \underline{in the case of non-semistable reduction}.

\medskip
\noindent
\SGARefTarget{sga7.sga7i.eix.n6-1}{6.1}. The construction of the toric part of $T_\ell(A_K)$ given in \SGARefLink{sga7.sga7i.eix.n5-1}{5.1}
still makes sense as soon as $S$ is complete, without assuming that
$A^o$ is $\ell$-divisible; that is (when $p=\ell$), without assuming
semistable reduction.  That hypothesis had been used only in order to
formulate the orthogonality theorem \SGARefLink{sga7.sga7i.eix.n5-2}{5.2}.  Moreover, one can construct
$T_p(A_K)^t$ while assuming only that $S$ is henselian: one shows that
the toric part $T_p(A_{\widehat K})^t$ over the completion, constructed
by the method of \SGARefLink{sga7.sga7i.eix.n5-1}{5.1}, comes from a Barsotti--Tate pro-subgroup of
$T_p(A_K)$, necessarily determined uniquely.

More precisely, one can construct a pro-$p$ Barsotti--Tate group
$T_p(A)^t$ \underline{over} $S$ and a homomorphism
\[
 T_p(A)^t\longrightarrow T_p(A)
\]
of projective systems which is a monomorphism term by term and whose
pullback to $\widehat S$ gives the analogous objects constructed by the
method of \SGARefLink{sga7.sga7i.eix.n5-1}{5.1}.  This follows by paraphrasing the construction of \SGARefLink{sga7.sga7i.eix.n5-1}{5.1},
except that, instead of lifting the whole maximal torus $T_o$ of $A_o$,
for each integer $\nu\geq0$ one merely lifts the subgroup
${}_{p^\nu}T_o$ to a multiplicative-type subgroup
$({}_{p^\nu}A^o)^t$ of ${}_{p^\nu}A^o$.

For the existence and uniqueness of such a lifting, note that the
functor
\[
 (\mathrm{Sch})^o_{/S}\longrightarrow\mathrm{Sets},
 \qquad
 S'\longmapsto
 \{\text{finite multiplicative-type subgroups of }A^o_{S'}\}
\]
is represented by an $S$-scheme $X$ of finite type (cf. the proof of
SGA 3 XI 3.12(a)), which is \'etale over $S$ by SGA 3 IX 3.6 bis.
Consequently, there is exactly one section of $X$ over $S$ extending a
given section of $X_s$ over $s$, which proves the assertion.

The same argument shows more generally that, for every
\underline{finite multiplicative-type subgroup} $H_o$ of $A_o$, there
is a unique finite multiplicative-type subgroup $H$ of $A$ extending
it.  Consequently, for every \underline{toric} Barsotti--Tate
$p$-subgroup $H_o$ of $A_o$ (that is, one whose components are of
multiplicative type), there is a unique Barsotti--Tate $p$-subgroup $H$
of $A$ extending it.  This in turn gives a Barsotti--Tate $p$-subgroup
of ${}_{p^\infty}A_K$, or, equivalently, a pro-$p$ Barsotti--Tate
subgroup of $T_p(A_K)$.

In particular, this procedure applies when $H_o$ is the largest toric
Barsotti--Tate $p$-subgroup of $A_o$.  It is obtained by taking the
largest $p$-divisible subgroup $\overline A_o$ of $A_o^o$
($\overline A_o=p^N A_o^o$ for $N$ large) and the ``toric part'' of
${}_{p^\infty}\overline A_o$ (defined, for example, by SGA 3 XVII
7.2.1).  We thus obtain a Barsotti--Tate subgroup of
${}_{p^\infty}A_K$ which might be called the \underline{toric part} of
the latter.

\medskip
\noindent
\SGARefTarget{sga7.sga7i.eix.n6-2}{6.2}. Again let $\overline A_o$ be the largest $p$-divisible subgroup
scheme of $A_o^o$, and consider the Barsotti--Tate $p$-group
\[
 F_o={}_{p^\infty}(\overline A_o).
\]
It is immediate that $F_o$ may be characterized as the largest
Barsotti--Tate $p$-subgroup of $A_o$ (or of $A_o^o$).  It is natural to
ask whether this Barsotti--Tate group can be extended to a
Barsotti--Tate subgroup $F$ of $A$, which would then play the role of
the \underline{fixed part} $T_\ell(A)^f$ of (\SGARefLink{sga7.sga7i.eix.n2-2-3-2}{2.2.3.2}) in the case in
which $A^o$ is $\ell$-divisible.

In the semistable case, where $\overline A_o=A_o^o$, the method of \SGARefLink{sga7.sga7i.eix.n2-2}{2.2}
answers the question affirmatively.  In the general case, whenever the
problem has a solution, it is readily seen to be unique.  The
corresponding pro-subgroups of $T_p(A_K)$ and $T_p(A'_K)$ (which again
deserve to be called the \underline{fixed parts} of the pro-groups in
question) still satisfy the orthogonality theorem \SGARefLink{sga7.sga7i.eix.n5-2}{5.2}, and the fixed
part $T_p(A_K)^f$ may still be characterized as in \SGARefLink{sga7.sga7i.eix.n5-8}{5.8}.

\medskip
\noindent
\SGARefTarget{sga7.sga7i.eix.n6-2-1}{6.2.1}. We shall nevertheless show that in general the question has no
affirmative answer; that is, \underline{the fixed part $T_p(A)^f$
cannot in general be defined} in the case of non-semistable reduction.
Indeed, by a lifting theorem of Serre--Tate [30], when $S$ is complete,
the existence of the fixed part is equivalent to the existence of
infinitesimal liftings $\overline A_n$ of $\overline A_o$ in the $A_n$,
or again to the existence of a smooth group scheme $\overline A$ over
$S$, an extension of an abelian scheme by a torus $T$, having special
fibre $\overline A_o$, together with a formal homomorphism
\[
 \widehat{\overline A}\longrightarrow\widehat A
\]
extending the inclusion on the special fibres.

Thus, when $\overline A_o$ is an abelian scheme, that is, when $T_o=0$,
this is equivalent to the existence of an abelian scheme $\overline A$
extending $\overline A_o$ and a homomorphism
$\overline A\longrightarrow A$ extending the inclusion of the special
fibres.  It will therefore suffice to give an example in which $T_o=0$,
$A_o^o$ is neither unipotent nor proper, and $A_K$ is simple over $K$.

\medskip
\noindent
\SGARefTarget{sga7.sga7i.eix.n6-2-2}{6.2.2}. Here is an example, due to M. Raynaud.  It is based on the
observation that, if $S$ is a trait, $S'$ is the normalization of $S$
in a finite separable extension $K'$ of $K$, and $A'$ is an abelian
scheme over $S'$, then the group scheme
\[
 A=\prod_{S'/S} A'/S'
\]
over $S$ is a N\'eron model: it identifies with the N\'eron model of its
generic fibre, which is indeed an abelian scheme.  This follows
immediately from the definitions and from the fact that $A'$ itself is
a N\'eron model over $S'$.  Then
\[
 A_o=\prod_{S'_s/s}(A'\times_{S'}S'_s)/S'_s
\]
is an extension of an abelian scheme, nonzero if $A'\neq0$, by a
connected unipotent scheme, nonzero if $A'\neq0$ and $S'$ is ramified
over $S$.

It therefore suffices to give an example in which $A'\neq0$, $S'$ is
ramified over $S$ (a complete local trait), and
\[
 A_K=\prod_{K'/K}A'_{K'}/K'
\]
is simple over $K$.  Take $K'/K$ to be a ramified quadratic extension
(such an extension always exists), and take $A'$ to be an elliptic
curve over $S'$.  Then $A$ has relative dimension two, and saying that
$A_K$ is simple over $K$ means that there is no nonzero homomorphism
\[
 B_K\longrightarrow A_K,
\]
where $B_K$ is an abelian scheme of dimension one over $K$.  Giving
such a homomorphism is equivalent to giving a nonzero homomorphism
\[
 B_{K'}\longrightarrow A'_{K'},
\]
which would necessarily be an isogeny.  Thus $A_K$ is simple over $K$
precisely when $A'_{K'}$ is not isogenous to an elliptic curve coming
from $K$.

If the contrary held, the $j$-invariant of $A'$ would be algebraically
dependent (over the prime field) on its conjugate $\overline j$ under
the nontrivial $K$-automorphism of $K'$.  If $j$ is chosen so that
\[
 a=j+\overline j=\operatorname{Tr}_{K'/K}j
\]
is transcendental over the prime field, this is equivalent to saying
that $j$ is algebraically dependent on $a$.  By cardinality
considerations, however, there exists an element $a\in V$
transcendental over the prime subfield $P$ of $K$, and then an element
$b\in V'$ (the normalization of $V$ in $K'$) which is transcendental
over $P(a)$ and satisfies
\[
 \operatorname{Tr}_{V'/V}b=0.
\]
It is enough to take $j=a+b\in V'$.  We may moreover suppose that
$b\in\mathfrak r(V')$, and that the image of $a$ in
$k=V/\mathfrak m$ is a prescribed element of $k$, chosen distinct from
the exceptional values $j_o=0$ and $j_o=12^3$.  This is possible even
when $k$ has only two elements, because then these two contemplated
values of $j_o$ coincide.  Thus the image $j_o$ of $j$ in $k(s')$ is
different from $0$ and $12^3$, and it is known that there then exists
an elliptic curve over $S'$ with $j$-invariant $j$.

\medskip
\noindent
\SGARefTarget{sga7.sga7i.eix.n6-3}{6.3}. Return to the case in which $S$ is merely assumed henselian.  In
analogy with the definitions of \SGARefLink{sga7.sga7i.eix.n4-1}{4.1} for $\ell\neq p$, one can then
define a filtration
\begin{equation}\tag{\SGARefTarget{sga7.sga7i.eix.n6-3-1}{6.3.1}}
 T_p(A_K)\supset T_p(A_K)^{\mathrm{ef}}
 \supset T_p(A_K)^{\mathrm{et}}\supset0
\end{equation}
by strict Barsotti--Tate pro-subgroups, called respectively the
\underline{essentially fixed part} and the
\underline{essentially toric part} of $T_p(A_K)$.

To do so, choose a finite Galois extension $K'$ of $K$ such that
$A_{K'}$ has semistable reduction over the normalization $S'$ of $S$
in $K'$.  The fixed part $T_p(A_{K'})^f$ of
$T_p(A_{K'})=(T_p(A_K))_{K'}$ is then defined (\SGARefLink{sga7.sga7i.eix.n2-2-3}{2.2.3}), as is the toric
part (\SGARefLink{sga7.sga7i.eix.n6-1}{6.1}).  The latter is moreover the orthogonal complement of the
fixed part $T_p(A'_{K'})^f$ in $T_p(A'_{K'})$ (\SGARefLink{sga7.sga7i.eix.n5-2}{5.2}).  Plainly
\[
 T_p(A_{K'})^t\subset T_p(A_{K'})^f
\]
(this follows from \SGARefLink{sga7.sga7i.eix.n5-3}{5.3}, or directly from the construction of the toric
part in \SGARefLink{sga7.sga7i.eix.n6-1}{6.1}), and the Galois group of $K'/K$ preserves these two strict
pro-subgroups of $(T_p(A_K))_{K'}$.  Descent therefore gives the desired
filtration (\SGARefLink{sga7.sga7i.eix.n6-3-1}{6.3.1}).  It is independent of the choice of $K'$, by the
semistable-reduction property \SGARefLink{sga7.sga7i.eix.n3-3}{3.3}.  Descent also plainly gives the
orthogonality property
\begin{equation}\tag{\SGARefTarget{sga7.sga7i.eix.n6-3-2}{6.3.2}}
 \begin{split}
 T_p(A_K)^{\mathrm{et}}
   &=\bigl(T_p(A'_{K'})^{\mathrm{ef}}\bigr)^\perp,\\
 T_p(A_K)^{\mathrm{ef}}
   &=\bigl(T_p(A'_{K'})^{\mathrm{et}}\bigr)^\perp.
 \end{split}
\end{equation}

One could give an intrinsic characterization of the filtration (\SGARefLink{sga7.sga7i.eix.n6-3-1}{6.3.1})
in terms of the Barsotti--Tate pro-group $X=T_p(A_K)$ over the fraction
field $K$ of $S$, in the spirit of \SGARefLink{sga7.sga7i.eix.n5-8}{5.8}---at least, for the moment, in
unequal characteristics.  The essentially fixed part is the largest
Barsotti--Tate pro-subgroup having ``essentially good reduction over
$S$'': after a finite separable extension $K'$ of $K$, it gives a
Barsotti--Tate pro-group $X_{K'}$ having good reduction over the
normalization $S'$ of $S$ in $K'$, that is, extending to a
Barsotti--Tate pro-group over $S'$.  Such an extension is necessarily
unique up to a unique isomorphism, by the already-cited theorem of Tate
[33, Th. 4].  The essentially toric part is then defined in terms of
the essentially fixed part by orthogonality (\SGARefLink{sga7.sga7i.eix.n6-3-2}{6.3.2}).

Let us also record the following immediate consequence of \SGARefLink{sga7.sga7i.eix.n5-8}{5.8} and the
definitions; for $\ell=p$ it reduces to \SGARefLink{sga7.sga7i.eix.n5-12}{5.12}(a).

\medskip
\noindent
\underline{Proposition} \SGARefTarget{sga7.sga7i.eix.n6-3-3}{6.3.3}. \underline{Let $\ell$ be a prime
number.  If $\ell=p$, assume that $\operatorname{char}K=0$.}
\footnote{This hypothesis is probably unnecessary; cf. \SGARefLink{sga7.sga7i.eix.n5-11}{5.11}.}
\underline{Then $A_K$ has essential good reduction over $S$ (\SGARefLink{sga7.sga7i.eix.n5-12}{5.12}(a))
if and only if its essentially fixed part}
$T_\ell(A_K)^{\mathrm{ef}}$
\underline{(\SGARefLink{sga7.sga7i.eix.n4-1-1}{4.1.1} and \SGARefLink{sga7.sga7i.eix.n6-3-1}{6.3.1}) is all of $T_\ell(A_K)$.}

\medskip
\noindent
\SGARefTarget{sga7.sga7i.eix.n6-4}{6.4}. Recall that, when $\ell\neq p$, in \SGARefLink{sga7.sga7i.eix.n4-3}{4.3}(a) we exhibited a
remarkable integer-valued function on the inertia group $I$ of $S$,
independent of the choice of $\ell$.  When the residue field is perfect,
knowledge of this function suffices to reconstruct the conductor
exponent of $A$ (\SGARefLink{sga7.sga7i.eix.n4-5}{4.5}): one takes, on a sufficiently large quotient
$\Gamma=I/I'$ of $I$, the inner product of this function with the Artin
conductor on $\Gamma$ (cf. (\SGARefLink{sga7.sga7i.eix.n4-6-2}{4.6.2})).  It is useful to observe that this
integer-valued function on $I$ can also be recovered from the knowledge
of $T_\ell(A_K)$ for $\ell=p$---at least, for the time being, in unequal
characteristics.

When the residue field $k$ is finite, an analogous observation applies
to the $\mathbb Q$-valued function considered in \SGARefLink{sga7.sga7i.eix.n4-3}{4.3}(b), defined on the
subgroup $\pi'$ of $\pi$ consisting of the elements whose image in
$\pi_o$ has the form $f_q^{\,n}$, $n\in\mathbb Z$.  We indicate the
principle of the verification.  Using the filtration (\SGARefLink{sga7.sga7i.eix.n6-3-1}{6.3.1}), one is
reduced to attaching a suitable function (on $I$, respectively on
$\pi'$) to a pro-$p$ Barsotti--Tate group $Y$ over $K$ having
``essentially good reduction over $S$'' in the sense of \SGARefLink{sga7.sga7i.eix.n5-12}{5.12}.  Here $Y$
will be one of the three quotients arising from the filtration (\SGARefLink{sga7.sga7i.eix.n6-3-1}{6.3.1}).

In the case of \SGARefLink{sga7.sga7i.eix.n4-3}{4.3}(a), a suitable residue-field extension within $S$
(EGA $0_{\mathrm{III}}$ 10.3.1) reduces us to the case of a perfect
residue field, or even an algebraically closed one if desired.  Choose
a finite Galois extension $K'$ of $K$ such that $Y_{K'}$ has good
reduction over $S'$.  We then have a Barsotti--Tate group $Y_{S'}$ over
$S'$ and an action of $G$ on $Y_{S'}$ compatible with its action on
$S'$.  If $k'$ is the residue field of $S'$, then $G$ acts, not
necessarily faithfully, on $k'$ (whose invariant field is $k$), and on
$Y_{k'}$ compatibly with its action on $k'$.

By Dieudonn\'e theory [16 bis], the Barsotti--Tate pro-$p$ group $Y_{k'}$
has a canonically associated finite free module $M$ (of rank equal to
the rank of $Y_{k'}$) over the ring $W(k')$ of Witt vectors of $k'$.
The group $G$ plainly acts on $M$ compatibly with its actions on
$W(k')$, through its actions on $k'$.  Hence the inertia subgroup
$\Gamma$ of $G$ acts $W(k')$-linearly on $M$; since $\Gamma$ is a
quotient of the inertia group $I$, this gives a trace function on
$\Gamma$, and hence on $I$, which is the desired function.

A priori this function takes values in $W(k')$,
\footnote{In fact it is immediate that this function always takes values
in $W(k)$, since it is invariant under $\operatorname{Gal}(k'/k)$.}
but it agrees with the analogous function defined in terms of the
$T_\ell$ for $\ell\neq p$ whenever the action of $G$ on $Y_{S'}$ is
obtained from an action of $G$ on a smooth $p$-divisible group scheme
$C'$ over $S'$ by taking $Y=T_p(C')$ (cf. [25 bis, Th. 7, p. 737]).
It is therefore integer-valued in this case.  The case of \SGARefLink{sga7.sga7i.eix.n4-3}{4.3}(b) is
handled similarly: to define a trace function on $\pi'$, one uses the
method of the proof of \SGARefLink{sga7.sga7i.eix.n4-3-2}{4.3.2} to ``linearize'' on $M$ the action defined
by an element $g\in\pi$ lying above an $f_q^n$, with $n\geq0$.

% END INCLUDED SGA7I_English_current_buildable_source_20260823/components/148_expose_IX_section_6.tex

% BEGIN INCLUDED SGA7I_English_current_buildable_source_20260823/components/149_expose_IX_section_7_1.tex
% Expose IX, section 7 title and 7.1--7.1.6.
% French transcription lines 1845--1887; authority scan physical pp. 411--414
% (scan indices 410--413; source folios 399--402).  All four pages were
% checked directly in four overlapping 1100-dpi bands; no higher-resolution
% escalation was needed.  In 7.1.3 the print itself changes incoherently from
% M on G to L on G_n.  The type-correct reading L on A, L_n on A_n, M_n on
% B_n is fixed by the surrounding construction and by Proposition 7.1.5.

\medskip
\noindent
7. \underline{The Raynaud extension $G^\natural$ over $S$ associated
with the abelian scheme $A_K$ having semistable reduction.  Duality of
the abelian schemes $B,B'$ over $S$.}

\medskip
\noindent
\SGARefTarget{sga7.sga7i.eix.n7-1}{7.1}. Let $S$ be a Noetherian scheme which is the spectrum of a ring
$V$, separated and complete for the topology defined by an ideal
$\underline m$ of $V$.  Let $A$ be a smooth group scheme over $S$, let
$A^o$ be its identity component, and suppose that $A_o/A_o^o$ is
finite (hence finite etale) over $S_o$, and that $A_o^o$ is an
extension of an abelian scheme by a torus.  As usual, we have put
$A_n=A\times_S S_n$, where
$S_n=\operatorname{Spec}(V/\underline m^{n+1})$.  We shall be
interested chiefly in the case where $S$ is a trait, the generic fiber
$A_K$ is an abelian scheme having semistable reduction over $S$, and
$A$ is the Neron model of $A_K$.

Let $T_o$ be the maximal torus of $A_o$, so that
$B_o=A_o^o/T_o$ is an abelian scheme.  The construction of \SGARefLink{sga7.sga7i.eix.n5-1}{5.1} then
applies and gives a formal subtorus $\widehat T$ of the formal
completion $\widehat A$.  For every $n$, the quotient
\begin{equation}\tag{\SGARefTarget{sga7.sga7i.eix.n7-1-1}{7.1.1}}
 B_n=A_n^o/T_n
\end{equation}
is representable (SGA~3 $\mathrm{VI}_A$~3.2); more precisely, it is a
flat scheme of finite type over $S_n$ whose pullback to $S_o$ is the
abelian scheme $B_o$.  Consequently $B_n$ is an \underline{abelian
scheme} over $S_n$.  As $n$ varies, the $B_n$ are obtained from one
another by base change and therefore define a \underline{formal
abelian scheme}
\begin{equation}\tag{\SGARefTarget{sga7.sga7i.eix.n7-1-2}{7.1.2}}
 \widehat B=(B_n)_{n\geqslant0}\,.
\end{equation}
We shall give conditions implying that this formal abelian scheme is
\underline{algebraizable} (EGA~III~5.4.1), more precisely, that it is
the formal completion of a projective abelian scheme over $S$, which
we shall then denote by $B$.  By loc. cit., $B$ is determined up to a
unique isomorphism (as an $S$-scheme; the group structure of
$\widehat B$ then comes uniquely from a group structure on $B$, again
by loc. cit.).  This abelian scheme $B$ will then be called the
\underline{abelian part} of the scheme $A$, and will sometimes be
denoted by $A^{\mathrm{ab}}$.

\medskip
\noindent
\SGARefTarget{sga7.sga7i.eix.n7-1-3}{7.1.3}. By EGA~III~5.4.5, projectivity of $\widehat B$ means that one
can find an invertible sheaf $\widehat{\mathbb M}$ on $\widehat B$,
that is, a sequence of invertible sheaves $\mathbb M_n$ on the $B_n$
which glue as $n$ varies, and which is relatively ample over $S$, that
is, such that $\mathbb M_o$ is ample.  The natural idea would be to
start with an ample invertible sheaf $\mathbb L$ on $A$ and show that,
for every $n$, its restriction $\mathbb L_n$ to $A_n$ is the pullback
of an invertible sheaf $\mathbb M_n$ on $B_n$.\footnote{The printed
text changes incoherently here from an invertible sheaf $\mathbb M$ on
$G$ to the restriction $\mathbb L_n$ of $\mathbb L$ to $G_n$.  The
notation above is the type-correct reading fixed by the surrounding
construction and by Proposition~\SGARefLink{sga7.sga7i.eix.n7-1-5}{7.1.5}.}  I do not know whether such
$\mathbb M_n$ necessarily exist (and strongly doubt it).  We shall
follow a slightly different method, constructing a
\underline{polarization} of $\widehat B$, that is, a projective system
of polarizations $\rho_n$ of the $B_n$ which glue.  It is known [5]
that the doubles of these polarizations canonically define
(symmetric) invertible sheaves $\mathbb M_n$ on the $B_n$.  By their
canonical nature these sheaves necessarily glue and furnish the
desired $\widehat{\mathbb M}$.

We shall interpret an element of $\operatorname{NS}(B_n/S_n)$ as a
symmetric divisorial correspondence on $(B_n,B_n)$ [5], that is, as a
biextension of $(B_n,B_n)$ by $\mathbb G_{mS_n}$.  By VIII~\SGARefLink{sga7.sga7i.eviii.n3-5}{3.5}, these
biextensions correspond bijectively to biextensions of $(A_n,A_n)$ by
$\mathbb G_{mS_n}$.  To obtain a coherent system of such biextensions
$W_n$, it suffices to start with
\begin{equation}\tag{\SGARefTarget{sga7.sga7i.eix.n7-1-4}{7.1.4}}
 W,\quad\text{a symmetric biextension of $(A,A)$ by $\mathbb G_{mS}$},
\end{equation}
and take the $W_n$ obtained from $W$ by the base changes
$S_n\longrightarrow S$.  We must in addition suppose that the element
of $\operatorname{NS}(B_o)$ defined by $W_o$ is a polarization of
$B_o$; we thereby obtain a polarization of $\widehat B$.

\medskip
\noindent
\underline{Proposition \SGARefTarget{sga7.sga7i.eix.n7-1-5}{7.1.5}.} \underline{With the preceding notation
for $S,A,\widehat B$, suppose that $S$ is normal.  Then $\widehat B$ is
algebraizable, and the abelian scheme $B$ over $S$ from which it comes
is projective over $S$.}

We may plainly suppose that $A=A^o$ and that $S$ is integral.  By
Raynaud [23, XI~1.13], the hypothesis that $S$ is normal implies that
$A$ is quasi-projective over $S$.  Choose a relatively ample
invertible sheaf $\mathbb L$ on $A$, and form the birigidified
invertible sheaf
\[
 \delta(\mathbb L)=
 \pi^*(\mathbb L)\operatorname{pr}_1^*(\mathbb L)^{-1}
 \operatorname{pr}_2^*(\mathbb L)^{-1}
\]
on $A\times_S A$, as in VIII~\SGARefLink{sga7.sga7i.eviii.n4-12}{4.12}.  Since the fibers of $A_o$ have
unipotent rank zero, the same is true of the fibers of $A$ (this
follows easily from \SGARefLink{sga7.sga7i.eix.n3-1}{3.1}).  Therefore, by VIII~\SGARefLink{sga7.sga7i.eviii.n7-5}{7.5},
$\delta(\mathbb L)$ defines a biextension $W$ of $(A,A)$, plainly
symmetric, which by VIII~\SGARefLink{sga7.sga7i.eviii.n4-13}{4.13} (that is, essentially by Raynaud
[23, XI~1.11]) satisfies the required ampleness condition at the
points of $S_o$.

\medskip
\noindent
\underline{Remark \SGARefTarget{sga7.sga7i.eix.n7-1-6}{7.1.6}.} Without assuming that $S$ is normal, one
can prove that $\widehat B$ is always algebraizable by reducing to the
normal case as in \SGARefLink{sga7.sga7i.eix.n13-5-3}{13.5.3}(a).

% END INCLUDED SGA7I_English_current_buildable_source_20260823/components/149_expose_IX_section_7_1.tex

% BEGIN INCLUDED SGA7I_English_current_buildable_source_20260823/components/150_expose_IX_lemma_7_2_1.tex
% Expose IX, 7.2 and Lemma 7.2.1 with proof 7.2.1.1--7.2.1.5.
% French transcription lines 1889--1938; authority scan physical pp. 414--418
% (scan indices 413--417; source folios 402--406).  All pages were checked
% directly in four overlapping 1100-dpi bands.  No unresolved feature needed
% higher-resolution escalation.
% Source notation corrected transparently: the adic base is Spf(V), not the
% printed Spf(A), since A is the varying group scheme; likewise Phi_o and
% Phi_n live over S_o and S_n, although 7.2.1.3 prints S.  The parenthetical
% source citation for representability in 7.2.1.3 is blank and is not guessed.

\medskip
\noindent
\SGARefTarget{sga7.sga7i.eix.n7-2}{7.2}. We now record the following lemma.

\medskip
\noindent
\underline{Lemma \SGARefTarget{sga7.sga7i.eix.n7-2-1}{7.2.1}.} Let $S$ be the spectrum of a Noetherian ring
$V$, separated and complete for the topology defined by an ideal
$\underline m$ of $V$.  Let $C$ be the category of group schemes
$\Lambda$ over $S$, smooth and of finite type over $S$, such that
$\Lambda/\Lambda^o$ is a finite etale group scheme over $S$ and
$\Lambda^o$ is an extension of an abelian scheme $B$ by an isotrivial
torus $T$ (SGA~3 IX~1.1).  Let $C'$ be the category of group objects
in the category of formal schemes of finite type (EGA~I~10.3.3) over
$\widehat S=\operatorname{Spf}(V)$, and consider the functor
\[
 A\longmapsto\widehat A
\]
from $C$ to $C'$ which takes the formal completion along
$A_o=A\times_S S_o$, where
$S_o=\operatorname{Spec}(V/\underline m)$.\footnote{At both occurrences
in this statement the print writes $\operatorname{Specf}(A)$, and it
also writes $\operatorname{Spec}(A/\underline m^{n+1})$ below.  Since
$V$ is the adic base ring and $A$ denotes the varying group scheme,
the type-correct base is $\operatorname{Spf}(V)$ and
$S_n=\operatorname{Spec}(V/\underline m^{n+1})$.} Then:

\smallskip
\noindent
\underline{(a) The preceding functor is fully faithful.}

\smallskip
\noindent
\underline{(b) Let $\mathcal A$ be an object of $C'$, smooth over
$\widehat S$}; that is, suppose that the ordinary schemes
\[
 A_n=\mathcal A\times_{\widehat S}S_n,
 \qquad
 S_n=\operatorname{Spec}(V/\underline m^{n+1}),
\]
are smooth over $S_n$.  Suppose that $A_o/A_o^o$ is a finite etale
group scheme over $S_o$, and that $A_o^o$ is an extension of an
abelian scheme $B_o$ by an isotrivial torus $T_o$.  Put
\[
 \mathcal A^o=(A_n^o)_{n\geq0},
 \qquad
 \mathcal A/\mathcal A^o=(A_n/A_n^o)_{n\geq0}.
\]
\underline{Then $\mathcal A/\mathcal A^o$ is finite etale over
$\widehat S$, and $\mathcal A^o$ is, in an essentially unique manner
and functorially in $\mathcal A$, an extension of a formal abelian
scheme $\mathcal B$ over $\widehat S$ by a formal torus $\mathcal T$
over $\widehat S$, reducing respectively to $B_o$ and $T_o$.  The
formal group scheme $\mathcal A$ belongs to the essential image of
$C\longrightarrow C'$---or, as we shall also say, is
``algebraizable''---if and only if $\mathcal B$ is algebraizable.}

\medskip
\noindent
\SGARefTarget{sga7.sga7i.eix.n7-2-1-1}{7.2.1.1}. First observe, over an arbitrary scheme $S$, that if $A$ is
a smooth group scheme such that $A^o$ is an extension of an abelian
scheme by a torus, this extension structure is essentially unique and
is functorial in $A$ in the evident sense.  Indeed, the assertion
reduces at once to the fact that every morphism from a torus $T$ to an
abelian scheme $B$ is zero.  This is well known on geometric fibers
[15], and SGA~3 IX~5.3 then gives the conclusion.  Applying the result
to the compatible systems of smooth group schemes $A_n$ over the
$S_n$ occurring in \SGARefLink{sga7.sga7i.eix.n7-2-1}{7.2.1} gives the analogous result for formal group
schemes over $S$.  Moreover, we already observed in \SGARefLink{sga7.sga7i.eix.n7-1}{7.1} that, if
$A_o^o$ is an extension of an abelian scheme by a torus, the same is
true of $A_n^o$.  This proves the first assertion of (b).

\medskip
\noindent
\SGARefTarget{sga7.sga7i.eix.n7-2-1-2}{7.2.1.2}. Suppose next that the formal group scheme $\mathcal A$ over
$S$ is of the form $\widehat A$, with $A\in\operatorname{ob}C$.  If
$A^o$ is an extension of an abelian scheme $B$ by a torus $T$, then
plainly $\mathcal B\simeq\widehat B$.  Thus algebraizability of
$\mathcal A$ implies algebraizability of $\mathcal B$.

Conversely, suppose that $\mathcal B$ is algebraizable, that is, that
it comes from an abelian scheme $B$, and let us prove that $\mathcal A$
is algebraizable.  We first prove that
$\mathcal A^o=(A_n^o)$ is algebraizable.  By VIII~\SGARefLink{sga7.sga7i.eviii.n3-7}{3.7}, the extensions
$A_n^o$ of $B_n$ by $T_n$ may be interpreted as homomorphisms
\begin{equation}\tag{*}
 \underline M_n\longrightarrow B'_n,
\end{equation}
where $\underline M_n$ is the twisted constant group over $S_n$ dual
to $T_n$ (SGA~3 X~5.7), and $B'_n$ is the abelian scheme dual to
$B_n$.  Restriction to $S_o$, or to any $S_n$, induces an equivalence
from the category of isotrivial tori over $S$ to the corresponding
category over $S_o$ (SGA~3 X~3.2).  Hence the $T_n$ come from a torus
$T$ over $S$, determined up to a unique isomorphism.  Denote its dual
group by $\underline M$; its restriction to $S_n$ is canonically
isomorphic to $\underline M_n$.

The isotriviality of $T$ means that $\underline M$ is a disjoint union
of $S$-schemes $\underline M^i$ which are finite over $S$
(SGA~3 X~5.11).  It follows at once that the homomorphisms $(*)$ come
from a unique homomorphism
\begin{equation}\tag{**}
 \underline M\longrightarrow B',
\end{equation}
where $B'$ is the abelian scheme dual to $B$.  In turn, this
homomorphism is interpreted as an extension of $B$ by $T$, say $A^o$,
which plainly supplies an algebraization of $\mathcal A^o$.

\medskip
\noindent
\SGARefTarget{sga7.sga7i.eix.n7-2-1-3}{7.2.1.3}. On the other hand, by hypothesis $A_o/A_o^o$ is a finite
group scheme $\Phi_o$ over $S_o$, and the $A_n/A_n^o$ are likewise
finite group schemes $\Phi_n$ over $S_n$; they are necessarily etale
over $S_n$.\footnote{The print writes $S$ rather than the fiber bases
$S_o$ and $S_n$ in these sentences, and leaves blank the
parenthetical citation from which representability of $\Phi_n$ is
said to follow.  The bases are restored from the statement of the
lemma; the missing citation is not guessed.}  There is therefore an
etale group scheme $\Phi$ over $S$, determined up to a unique
isomorphism, whose formal completion is the formal group scheme
$\mathcal A/\mathcal A^o$ defined by the $\Phi_n$.

Since $G=A^o$ is divisible, the category of extensions, as fppf
sheaves, of a finite locally constant group over $S$ such as $\Phi$
by the group scheme $A^o$ is canonically equivalent to the category
of torsors under the sheaf
\[
 H=\underline{\operatorname{Hom}}_{\mathrm{gr}}(\Phi,G),
\]
and this equivalence is compatible with every base change.  Applying
the same description to extensions of $\Phi_n$ by $G_n$ in terms of
torsors under
\[
 H_n\simeq
 \underline{\operatorname{Hom}}_{\mathrm{gr}}(\Phi_n,G_n),
\]
the $A_n$ define a coherent system of torsors $P_n$ under the $H_n$.
Because $G$ is an extension of an abelian scheme by a torus, it follows
readily that $H$ is a group scheme \underline{finite and locally free
over $S$}.  Consequently, the functor
\[
 P\longmapsto\widehat P=(P_n)
\]
is an equivalence between the category of torsors over $S$ under $H$
and the category of ``formal'' torsors $\widehat P$ under
$\widehat H$, that is, coherent systems of torsors $P_n$ under the
$H_n$.

Thus the $P_n$ defined by the extensions $A_n$ come from a single
torsor $P$ over $S$ under $H$.  We obtain an extension $A$ of $\Phi$
by $G=A^o$ which algebraizes $\mathcal A$, except that its
representability remains to be checked.  Regard $A$ as a torsor over
$\Phi$ under $A^o_\Phi$.  This torsor becomes trivial after base
change by the finite, surjective, locally free morphism
\[
 P\times_S\Phi\longrightarrow\Phi,
\]
since the base change $P\longrightarrow S$ splits the extension $A$
of $\Phi$ by $A^o$.  One may therefore conclude by one of Raynaud's
many criteria for effectiveness of descent [23, XI~3.1(6)].

\medskip
\noindent
\SGARefTarget{sga7.sga7i.eix.n7-2-1-4}{7.2.1.4}. This completes the proof of (b), and only the full
faithfulness assertion (a) remains.  We leave it to the reader to
verify that this assertion is essentially contained in the method of
proof of (b), or, if preferred, to find an independent direct proof.

\medskip
\noindent
\SGARefTarget{sga7.sga7i.eix.n7-2-1-5}{7.2.1.5}. Notice that the appeal to Raynaud's delicate descent result
becomes unnecessary if $A^o$ is assumed quasi-projective over $S$, or,
equivalently, if $B$ is projective over $S$.  This condition holds in
particular when one starts from the formal group scheme $\widehat A$
over $S$, with $A$ as in \SGARefLink{sga7.sga7i.eix.n7-1-5}{7.1.5}.  This is the only case in which we
shall use \SGARefLink{sga7.sga7i.eix.n7-2-1}{7.2.1}.

% END INCLUDED SGA7I_English_current_buildable_source_20260823/components/150_expose_IX_lemma_7_2_1.tex

% BEGIN INCLUDED SGA7I_English_current_buildable_source_20260823/components/151_expose_IX_sections_7_2_2_through_7_2_6.tex
% Expose IX, 7.2.2--7.2.6.
% French transcription lines 1940--1967; authority scan physical pp. 418--419
% (scan indices 417--418; source folios 406--407).  Both pages were checked
% directly in four overlapping 1100-dpi bands; no escalation was needed.

\medskip
\noindent
\SGARefTarget{sga7.sga7i.eix.n7-2-2}{7.2.2}. We apply \SGARefLink{sga7.sga7i.eix.n7-2-1}{7.2.1} to the formal group scheme $\widehat A$
considered in \SGARefLink{sga7.sga7i.eix.n7-1}{7.1}, in the favorable case treated there in which the
formal abelian scheme $\widehat B$ is algebraizable (cf. \SGARefLink{sga7.sga7i.eix.n7-1-5}{7.1.5}), and
we suppose in addition that the torus $T_o$ over $S_o$ is isotrivial.
We thereby obtain a group scheme $A^\natural$ over $S$, called the
\underline{Raynaud group} of $A$ over $S$, characterized up to
isomorphism by a canonical isomorphism of formal group schemes over
$\operatorname{Spf}(V)=\widehat S$:
\begin{equation}\tag{\SGARefTarget{sga7.sga7i.eix.n7-2-3}{7.2.3}}
 (A^\natural)^\wedge\simeq\widehat A.
\end{equation}
It follows from \SGARefLink{sga7.sga7i.eix.n7-2-1}{7.2.1} that
\begin{equation}\tag{\SGARefTarget{sga7.sga7i.eix.n7-2-4}{7.2.4}}
 \Phi=A^\natural/A^{\natural o}
\end{equation}
is a finite etale group over $S$, and that the identity component
\[
 A^{\natural o}=A^{o\natural}
\]
is endowed with the structure of an extension
\begin{equation}\tag{\SGARefTarget{sga7.sga7i.eix.n7-2-5}{7.2.5}}
 0\longrightarrow T\longrightarrow A^{\natural o}
 \longrightarrow B\longrightarrow0,
\end{equation}
where $B$ is the abelian scheme over $S$ whose formal completion is
the formal abelian scheme $\widehat B$ constructed in \SGARefLink{sga7.sga7i.eix.n7-1}{7.1}, and $T$ is
characterized, up to a unique isomorphism, as the isotrivial torus
over $S$ lifting $T_o$.  On restriction to the $S_n$, the extension
structure (\SGARefLink{sga7.sga7i.eix.n7-2-5}{7.2.5}) gives the extension structures on
$A_n^{\natural o}=A_n^o$ defined by (\SGARefLink{sga7.sga7i.eix.n7-1-1}{7.1.1}).

\medskip
\noindent
\SGARefTarget{sga7.sga7i.eix.n7-2-6}{7.2.6}. It also follows immediately from \SGARefLink{sga7.sga7i.eix.n7-2-1}{7.2.1} that the Raynaud group
$A^\natural$, as well as the exact sequence (\SGARefLink{sga7.sga7i.eix.n7-2-5}{7.2.5}), depends
functorially on $A$ (satisfying the conditions under consideration),
and that its formation commutes with every base change
$S'\longrightarrow S$ associated with a homomorphism of adic
topological rings $V\longrightarrow V'$.

% END INCLUDED SGA7I_English_current_buildable_source_20260823/components/151_expose_IX_sections_7_2_2_through_7_2_6.tex

% BEGIN INCLUDED SGA7I_English_current_buildable_source_20260823/components/152_expose_IX_section_7_3.tex
% Expose IX, section 7.3.
% French transcription lines 1969--2015; authority scan physical pp. 419--421
% (scan indices 418--420; source folios 407--409). All three pages were checked
% directly in four overlapping 1100-dpi bands; no escalation was needed.

\medskip
\noindent
\SGARefTarget{sga7.sga7i.eix.n7-3}{7.3}. Let us now consider, for a fixed prime number $\ell$, the
pro-$\ell$ group $T_{\ell}(A^{o})$. Identifying a finite scheme over
$S$ with the corresponding formal scheme, and defining the
``\underline{fixed part}'' $T_{\ell}(A^{o})^{f}$ as usual, we obtain
an isomorphism as in (\SGARefLink{sga7.sga7i.eix.n5-1-5}{5.1.5}):

\begin{equation}\tag{\SGARefTarget{sga7.sga7i.eix.n7-3-1}{7.3.1}}
T_{\ell}(A^{o})^{f} \simeq T_{\ell}(\widehat{A}^{o}) \quad ,
\end{equation}
and, defining the ``\underline{toric part}'' as in (\SGARefLink{sga7.sga7i.eix.n5-1-6}{5.1.6}),

\begin{equation}\tag{\SGARefTarget{sga7.sga7i.eix.n7-3-2}{7.3.2}}
T_{\ell}(A^{o})^{t} = T_{\ell}(\widehat{T}) \simeq T_{\ell}(T)
 \subset T_{\ell}(A^{o})^{f} \simeq T_{\ell}(\widehat{A}^{o}) \quad ,
\end{equation}
we immediately obtain, for the ``\underline{abelian part},'' an
isomorphism

\begin{equation}\tag{\SGARefTarget{sga7.sga7i.eix.n7-3-3}{7.3.3}}
T_{\ell}(A^{o})^{ab} \stackrel{\mathrm{def}}{=}
T_{\ell}(A^{o})^{f}/T_{\ell}(A^{o})^{t}
 \simeq T_{\ell}(\widehat{B}) \quad .
\end{equation}
We have not yet used here the hypothesis that $\widehat B$ is
algebraizable, the hypothesis that allowed us to construct the
abelian scheme $B$ and the Raynaud extension of $B$ by $T$. Since
there is, of course, a canonical isomorphism (the special case of
(\SGARefLink{sga7.sga7i.eix.n7-3-1}{7.3.1}) applied to the abelian scheme $B$!)
$T_{\ell}(\widehat B)\simeq T_{\ell}(B)$, the isomorphism (\SGARefLink{sga7.sga7i.eix.n7-3-3}{7.3.3})
also gives us

\begin{equation}\tag{\SGARefTarget{sga7.sga7i.eix.n7-3-4}{7.3.4}}
T_{\ell}(A^{o})^{ab} \simeq T_{\ell}(B) \quad .
\end{equation}
Let us also note in passing that (\SGARefLink{sga7.sga7i.eix.n7-2-3}{7.2.3}) and (\SGARefLink{sga7.sga7i.eix.n7-3-1}{7.3.1}) imply a canonical
isomorphism

\begin{equation}\tag{\SGARefTarget{sga7.sga7i.eix.n7-3-5}{7.3.5}}
T_{\ell}(A^{o})^{f} \simeq T_{\ell}(A^{\natural o}) \quad ,
\end{equation}
since the right-hand side is already a profinite group, that is,
identical to its fixed part. The canonical exact sequence

\begin{equation}\tag{\SGARefTarget{sga7.sga7i.eix.n7-3-6}{7.3.6}}
0 \longrightarrow T_{\ell}(A^{o})^{t}
 \longrightarrow T_{\ell}(A^{o})^{f}
 \longrightarrow T_{\ell}(B) \longrightarrow 0 \quad ,
\end{equation}
which can in any case (without the algebraizability hypothesis on
$\widehat B$) be interpreted as the exact sequence obtained by
applying the functor $T_{\ell}$ to the exact sequence

\begin{equation*}
0 \longrightarrow \widehat{T} \longrightarrow \widehat{A}^{o}
 \longrightarrow \widehat{B} \longrightarrow 0 \quad ,
\end{equation*}
can also be regarded as obtained by applying the functor $T_{\ell}$
to the canonical exact sequence of ordinary group schemes

\begin{equation*}
0 \longrightarrow T \longrightarrow A^{\natural o}
 \longrightarrow B \longrightarrow 0 \qquad .
\end{equation*}

% END INCLUDED SGA7I_English_current_buildable_source_20260823/components/152_expose_IX_section_7_3.tex

% BEGIN INCLUDED SGA7I_English_current_buildable_source_20260823/components/153_expose_IX_section_7_4.tex
% Expose IX, section 7.4.
% French transcription lines 2017--2080; authority scan physical pp. 421--424
% (scan indices 420--423; source folios 409--412). All four pages were
% checked directly in four overlapping 1100-dpi bands; no escalation was needed.

\medskip
\noindent
\SGARefTarget{sga7.sga7i.eix.n7-4}{7.4}. Suppose now that a second group scheme $A'$ over $S$ is given,
satisfying the same hypotheses as $A$, and suppose that $A$ and $A'$
have connected fibres. Let $W$ be a biextension of $(A,A')$ by
$\mathbb G_{mS}$. We thereby obtain a pairing (VIII \SGARefLink{sga7.sga7i.eviii.n2-2}{2.2})

\begin{equation}\tag{\SGARefTarget{sga7.sga7i.eix.n7-4-1}{7.4.1}}
T_{\ell}(A) \times T_{\ell}(A') \longrightarrow
T_{\ell}(\mathbb G_{mS}) = \mathbb Z_{\ell}(1)
\end{equation}
and, by restriction to the fixed parts, a pairing

\begin{equation}\tag{\SGARefTarget{sga7.sga7i.eix.n7-4-2}{7.4.2}}
T_{\ell}(A)^{f} \times T_{\ell}(A')^{f}
 = T_{\ell}(\widehat A) \times T_{\ell}(\widehat {A'})
 \longrightarrow \mathbb Z_{\ell}(1) \qquad .
\end{equation}
This can also be interpreted as the pairing defined by the
biextension $\widehat W$ of $(\widehat A,\widehat {A'})$ by
$\widehat{\mathbb G}_{mS}$ obtained from $W$ by completion (here, for
simplicity, defining a biextension of formal group schemes as a
coherent family of biextensions $W_n$ of $(A_n,A'_n)$ by
$\mathbb G_{mS_n}$). Now we know (VIII \SGARefLink{sga7.sga7i.eviii.n3-5}{3.5}) that $W_n$ comes from a
biextension $W_n^{ab}$ of $(B_n,B'_n)$, where $B'_n$ is the abelian
part of $A'_n$, so that (\SGARefLink{sga7.sga7i.eix.n7-4-2}{7.4.2}) factors through

\begin{equation}\tag{\SGARefTarget{sga7.sga7i.eix.n7-4-3}{7.4.3}}
T_{\ell}(A)^{ab} \times T_{\ell}(A')^{ab}
 = T_{\ell}(\widehat B) \times T_{\ell}(\widehat {B'})
 \longrightarrow \mathbb Z_{\ell}(1) \qquad .
\end{equation}
This is also the pairing defined by the biextension

\begin{equation*}
\widehat W^{ab}=(W_n^{ab})
\end{equation*}
of $(\widehat B,\widehat {B'})$ by $\widehat{\mathbb G}_{mS}$. When
$\widehat B$ and $\widehat {B'}$ are algebraizable, so that abelian
schemes $B,B'$ over $S$ are available, EGA IV 5.1.4 shows that
$\widehat W^{ab}$ comes from a biextension, defined up to a unique
isomorphism,

\begin{equation*}
W^{ab}\quad\text{of}\quad(B,B')\quad\text{by}\quad\mathbb G_{mS},
\end{equation*}
and the corresponding pairing

\begin{equation}\tag{\SGARefTarget{sga7.sga7i.eix.n7-4-4}{7.4.4}}
T_{\ell}(B) \times T_{\ell}(B')
 \longrightarrow T_{\ell}(\mathbb G_{mS})
\end{equation}
is canonically isomorphic to (\SGARefLink{sga7.sga7i.eix.n7-4-3}{7.4.3}).

Writing $D(B')$ for the abelian scheme dual to $B'$, $W^{ab}$ defines
a homomorphism

\begin{equation}\tag{\SGARefTarget{sga7.sga7i.eix.n7-4-5}{7.4.5}}
u:B\longrightarrow D(B') \qquad ,
\end{equation}
whose properties reflect those of the pairings (\SGARefLink{sga7.sga7i.eix.n7-4-4}{7.4.4}), that is,
(\SGARefLink{sga7.sga7i.eix.n7-4-3}{7.4.3}). For example, the pairing (\SGARefLink{sga7.sga7i.eix.n7-4-3}{7.4.3}) is a perfect duality at a
point $t\in S$ if and only if
$u_t:B_t\longrightarrow D(B')_t$ is an isogeny whose kernel has rank
prime to $\ell$. This implies that $u$ itself is an isogeny (since
$S$ is local, hence connected), and $\ker u$ is flat and finite over
$S$, so that $\ker u$ has the same rank at every point. Thus the
condition just considered is independent of the chosen point $t$.

\medskip
\noindent
\SGARefTarget{sga7.sga7i.eix.n7-4-6}{7.4.6}. This completes in particular the proof of \SGARefLink{sga7.sga7i.eix.n5-4}{5.4}. Indeed, under
the conditions of that theorem (where $S$ may be assumed complete),
applying the preceding discussion to $A^o$ and ${A'}^o$ shows that,
for every $\ell$, the pairing (\SGARefLink{sga7.sga7i.eix.n7-4-3}{7.4.3}) is a perfect duality at the
generic point $\eta$, and hence also at the point $s$.

\medskip
\noindent
\SGARefTarget{sga7.sga7i.eix.n7-4-7}{7.4.7}. The preceding considerations likewise make the orthogonality
theorem \SGARefLink{sga7.sga7i.eix.n5-2}{5.2} clear, without recourse to Lemma \SGARefLink{sga7.sga7i.eix.n5-2-4}{5.2.4}. Indeed, along the
way in \SGARefLink{sga7.sga7i.eix.n7-4}{7.4} we established that the pairing
$(\SGARefLink{sga7.sga7i.eix.n5-2-3-1}{5.2.3.1})=(\SGARefLink{sga7.sga7i.eix.n7-4-2}{7.4.2})$ vanishes on the toric parts on both sides, and
therefore factors through the pairing (\SGARefLink{sga7.sga7i.eix.n7-4-3}{7.4.3}). It remains only to
establish that this pairing is nondegenerate on the generic fibres of
the pro-groups under consideration, knowing that it is so on the
special fibres. This follows from the fact that if (\SGARefLink{sga7.sga7i.eix.n7-4-5}{7.4.5}) induces an
isogeny on the special fibres, then it is an isogeny.

\medskip
\noindent
\SGARefTarget{sga7.sga7i.eix.n7-4-8}{7.4.8}. Return to the general conditions of \SGARefLink{sga7.sga7i.eix.n7-4}{7.4}, and suppose that an
open subset $U\subset S-S_o$ is given such that $A_U$ and $A'_U$ are
\underline{abelian schemes}, dual to one another through the
biextension $W_U$. Restricting (\SGARefLink{sga7.sga7i.eix.n7-3-6}{7.3.6}), and the analogous exact
sequence for $A'$, to $U$, we obtain canonical filtrations on the
Barsotti--Tate pro-$\ell$ groups $T_\ell(A_U)$ and $T_\ell(A'_U)$:

\begin{equation}\tag{\SGARefTarget{sga7.sga7i.eix.n7-4-9}{7.4.9}}
\left\{
\begin{aligned}
T_\ell(A_U)&\supset T_\ell(A_U)^f\supset T_\ell(A_U)^t\supset0,\\[4ex]
T_\ell(A'_U)&\supset T_\ell(A'_U)^f\supset T_\ell(A'_U)^t\supset0
\qquad .
\end{aligned}
\right.
\end{equation}
I claim that \underline{these two filtrations are dual to one another}
with respect to the pairing (\SGARefLink{sga7.sga7i.eix.n7-4-1}{7.4.1}), which here is a perfect duality
(VIII \SGARefLink{sga7.sga7i.eviii.n3-2}{3.2}); that is, relation (\SGARefLink{sga7.sga7i.eix.n5-2-2}{5.2.2}) and its symmetric counterpart
still hold, with $A_K,A'_K$ replaced by $A_U,A'_U$. We saw above that
this orthogonality assertion is also equivalent to saying that the
biextension $W^{ab}$ of $(B,B')$ by $\mathbb G_{mS}$ is a perfect
duality over $U$ (or, equivalently when $U$ is dense in $S$, over
$S$). It is enough to verify this property at every point $\eta$ of
$U$; the usual argument then reduces us to the already known case in
which $S$ is a trait with generic point $\eta$.

The orthogonality relation also allows us to extend formulas (\SGARefLink{sga7.sga7i.eix.n5-5-7}{5.5.7}),
(\SGARefLink{sga7.sga7i.eix.n5-5-8}{5.5.8}), and (\SGARefLink{sga7.sga7i.eix.n5-5-9}{5.5.9}) to the present case, replacing restriction to the
generic point $\eta$ by restriction to the open subset $U$.
Consequently, \SGARefLink{sga7.sga7i.eix.n5-6}{5.6} likewise extends to the present case. We refrain
from rewriting the preceding formulas and statement here in this
context.

% END INCLUDED SGA7I_English_current_buildable_source_20260823/components/153_expose_IX_section_7_4.tex

% BEGIN INCLUDED SGA7I_English_current_buildable_source_20260823/components/154_expose_IX_proposition_7_5_and_section_7_6.tex
% Expose IX, Proposition 7.5 and section 7.6.
% French transcription lines 2082--2103; authority scan physical pp. 424--425
% (scan indices 423--424; source folios 412--413). Both pages were checked
% directly in four overlapping 1100-dpi bands; no escalation was needed.

Let us summarize, from the point of view of N\'{e}ron models, the
results obtained concerning the abelian parts and their pairings.

\medskip
\noindent
\underline{Proposition} \SGARefTarget{sga7.sga7i.eix.n7-5}{7.5}. \underline{Let $S$ be a complete trait
with function field $K$. To every abelian scheme $A_K$ over $K$
having semistable reduction over $S$, the procedure of \SGARefLink{sga7.sga7i.eix.n7-1}{7.1}, applied
to the identity component $A^o$ of the N\'{e}ron model $A$ of $A_K$,
associates its abelian part $B=A^{ab}$, which is an abelian scheme
over $S$ depending functorially on $A_K$. For every other abelian
scheme $A'_K$ with semistable reduction, one can associate to every
divisorial correspondence $W_K$ on $(A_K,A'_K)$ a unique divisorial
correspondence $W^{ab}$ on $(A^{ab},{A'}^{ab})$ such that, for every
integer $n\geq0$, the pullback of $W_n^{ab}$ along
$A_n^o\longrightarrow B_n=A_n^{oab}$ and
${A'}_n^o\longrightarrow B'_n={A'}_n^{oab}$ is the biextension $W_n$
of $(A_n^o,{A'}_n^o)$ obtained from $W$ by reduction, where $W$ is
the canonical extension (\SGARefLink{sga7.sga7i.eix.n1-4}{1.4}) of $W_K$ to $(A^o,{A'}^o)$.

The formation of $W^{ab}$ is compatible with the formation of
pullbacks of correspondences by morphisms of abelian schemes
$\overline A_K\longrightarrow A_K$ and
$\overline {A'}_K\longrightarrow A'_K$. If $W_K$ is a perfect
duality, then so is $W^{ab}$. If $W_K$ is the divisorial
correspondence on $(A_K,A_K)$ associated with a polarization of
$A_K$, then $W^{ab}$ is associated with a polarization of $A^{ab}$.}

The compatibility of the formation of the $W^{ab}$ with pullbacks is
immediate from the indicated characterization of the $W^{ab}$. The
additional assertion concerning polarizations was established in
\SGARefLink{sga7.sga7i.eix.n7-1-5}{7.1.5} at the same time as the algebraizability of $\widehat B$.

\medskip
\noindent
\SGARefTarget{sga7.sga7i.eix.n7-6}{7.6}. Continue to suppose that $S$ is a complete trait, with $A$ the
N\'{e}ron model of an abelian scheme $A_K$ over $K$. Set

\begin{equation}\tag{\SGARefTarget{sga7.sga7i.eix.n7-6-1}{7.6.1}}
A_K^{\natural}=(A^{\natural})_K \quad ,
\end{equation}
where $A^{\natural}$ is the Raynaud group scheme (\SGARefLink{sga7.sga7i.eix.n7-2-2}{7.2.2}): it is an
extension of the finite etale group $\Phi$ by
$A^{\natural o}$, while $A^{\natural o}$ is an extension of the
abelian scheme $B$ by the torus $T$. Thus $A_K^{\natural}$ is itself
an extension of $\Phi_K$ by
$(A_K^{\natural})^o=(A^{\natural o})_K=(A^{o\natural})_K$, which is an
extension of the abelian scheme $B_K$ by the torus $T_K$. For every
prime number $\ell$, (\SGARefLink{sga7.sga7i.eix.n7-3-5}{7.3.5}) and (\SGARefLink{sga7.sga7i.eix.n7-3-2}{7.3.2}) give

\begin{equation}\tag{\SGARefTarget{sga7.sga7i.eix.n7-6-2}{7.6.2}}
T_{\ell}(A_K^{\natural}) \simeq T_{\ell}(A_K^{\natural o})
 \simeq T_{\ell}(A_K)^f \; , \qquad
T_{\ell}(T_K) \simeq T_{\ell}(A_K)^t \quad .
\end{equation}

\medskip
\noindent
\SGARefTarget{sga7.sga7i.eix.n7-6-3}{7.6.3}. By \SGARefLink{sga7.sga7i.eix.n3-3}{3.3}, the formation of $A_K^{\natural o}$ commutes with
every base change associated with a morphism of complete traits
$S'\longrightarrow S$.

% END INCLUDED SGA7I_English_current_buildable_source_20260823/components/154_expose_IX_proposition_7_5_and_section_7_6.tex

% BEGIN INCLUDED SGA7I_English_current_buildable_source_20260823/components/155_expose_IX_section_8_1.tex
% Expose IX, section 8.1.
% French transcription lines 2109--2154; authority scan physical pp. 426--428
% (scan indices 425--427; source folios 414--416). All three pages were
% checked directly in four overlapping 1100-dpi bands; no escalation was
% needed.

\medskip
\noindent
8. \underline{The Raynaud extension $A_K^{\natural o}$ in the
non-semistable case, and the ind-group $A_K^{\flat}$}

\medskip
\noindent
\SGARefTarget{sga7.sga7i.eix.n8-1}{8.1}. Continue to suppose that the trait $S$ is complete, but no
longer suppose necessarily that $A_K$ has semistable reduction. As
usual, let $K'$ be a finite Galois extension of $K$ such that
$A_{K'}$ has semistable reduction relative to the normalization $S'$
of $S$ in $K'$. Consider the Raynaud extension
$A_{K'}^{\natural}$ associated with $A_{K'}$ (\SGARefLink{sga7.sga7i.eix.n7-6-1}{7.6.1}). It is clear
that the Galois group of $K'/K$ acts by transport of structure on
this group scheme; hence, by Galois descent, it gives a group scheme
over $K$, denoted $A_{K,K'}^{\natural}$, which depends on the choice
of the extension $K'$. Its identity component, however, is
independent of this choice up to canonical isomorphism, by \SGARefLink{sga7.sga7i.eix.n7-6-3}{7.6.3},
and will simply be denoted $A_K^{\natural o}$. Descent shows that
$A_K^{\natural o}$ is an extension

\begin{equation}\tag{\SGARefTarget{sga7.sga7i.eix.n8-1-1}{8.1.1}}
0 \longrightarrow T_K \longrightarrow A_K^{\natural o}
 \longrightarrow B_K \longrightarrow 0 \quad ,
\end{equation}
where $T_K$ is a torus and $B_K$ an abelian scheme over $K$.
Likewise, by the definition of the essentially fixed and essentially
toric parts of $T_{\ell}(A_K)$ (\SGARefLink{sga7.sga7i.eix.n4-1}{4.1} and 6.5), the isomorphisms
(\SGARefLink{sga7.sga7i.eix.n7-6-2}{7.6.2}) give canonical isomorphisms

\begin{equation}\tag{\SGARefTarget{sga7.sga7i.eix.n8-1-2}{8.1.2}}
T_{\ell}(A_K^{\natural}) \simeq T_{\ell}(A_K)^{ef} \quad ,
\qquad T_{\ell}(T_K) \simeq T_{\ell}(A_K)^{et} \quad ,
\end{equation}
and hence, on passing to the quotient,

\begin{equation}\tag{\SGARefTarget{sga7.sga7i.eix.n8-1-3}{8.1.3}}
T_{\ell}(B_K) \simeq
T_{\ell}(A_K)^{ef}/T_{\ell}(A_K)^{et} \quad .
\end{equation}

This therefore determines two of the quotients occurring in the
filtration

\begin{equation}\tag{\SGARefTarget{sga7.sga7i.eix.n8-1-4}{8.1.4}}
T_\ell(A_K) \supset T_\ell(A_K)^{ef}
 \supset T_\ell(A_K)^{et} \supset 0
\end{equation}
Thus the extension of group schemes (\SGARefLink{sga7.sga7i.eix.n8-1-1}{8.1.1}) describes two quotients
in this filtration. For the third quotient
$T_\ell(A_K)/T_\ell(A_K)^{ef}$, (6.5.2) gives a canonical
isomorphism

\begin{equation}\tag{\SGARefTarget{sga7.sga7i.eix.n8-1-5}{8.1.5}}
T_\ell(A_K)/T_\ell(A_K)^{ef}
 \simeq D(T_\ell(A'_K)^{et}) \quad ,
\end{equation}
where $D=\underline{\operatorname{Hom}}(-,T_\ell(\mathbb{G}_m))$
denotes Cartier duality (in the sense of Barsotti--Tate
pro-groups), and $A'_K$ denotes the abelian scheme dual to $A_K$.
Consider then the exact sequence of type (\SGARefLink{sga7.sga7i.eix.n8-1-1}{8.1.1})

\begin{equation}\tag{\SGARefTarget{sga7.sga7i.eix.n8-1-6}{8.1.6}}
0 \longrightarrow T'_K \longrightarrow {A'_K}^{\natural o}
 \longrightarrow B'_K \longrightarrow 0
\end{equation}
associated with the abelian scheme $A'_K$ over $K$, together with
the corresponding isomorphisms of type (\SGARefLink{sga7.sga7i.eix.n8-1-2}{8.1.2}). Formula (\SGARefLink{sga7.sga7i.eix.n8-1-5}{8.1.5})
then gives

\begin{equation}\tag{\SGARefTarget{sga7.sga7i.eix.n8-1-7}{8.1.7}}
T_\ell(A_K)/T_\ell(A_K)^{ef}
 \simeq D(T_\ell(T'_K))
 \simeq \underline{M}_K\otimes_{\mathbb{Z}}\mathbb{Z}_\ell \quad ,
\end{equation}
where, generalizing the notation introduced in (\SGARefLink{sga7.sga7i.eix.n5-5-2}{5.5.2}), we have set

\begin{equation}\tag{\SGARefTarget{sga7.sga7i.eix.n8-1-8}{8.1.8}}
\underline{M}_K=D(T'_K) \quad , \qquad
\underline{M}'_K=D(T_K) \quad ,
\end{equation}
and where $D$ now denotes Cartier duality
$\underline{\operatorname{Hom}}(-,\mathbb{G}_{mK})$ in the sense of
ordinary group schemes. Thus $\underline{M}_K$ and
$\underline{M}'_K$ are twisted constant groups over $K$ which become
constant over $K'$ and isomorphic to groups of the form
$\mathbb{Z}^{r}$, where $r=\dim T=\dim T'$.

% END INCLUDED SGA7I_English_current_buildable_source_20260823/components/155_expose_IX_section_8_1.tex

% BEGIN INCLUDED SGA7I_English_current_buildable_source_20260823/components/156_expose_IX_sections_8_2_and_8_3.tex
% Expose IX, sections 8.2--8.3.
% French transcription lines 2159--2188; authority scan physical pp. 428--429
% (scan indices 427--428; source folios 416--417). Both pages were checked
% directly in four overlapping 1100-dpi bands; physical p. 430 was checked
% for the exact section-9 boundary. The authority's evident A_K -> A'_K
% polarization is restored where the working transcription misplaces the
% prime as A_{K'}; its evident "extensions ... de K'" slip is rendered
% according to the mathematically required extensions of K.

\medskip
\noindent
\SGARefTarget{sga7.sga7i.eix.n8-2}{8.2}. The hypotheses are those of \SGARefLink{sga7.sga7i.eix.n8-1}{8.1}. It is clear that the extension
(\SGARefLink{sga7.sga7i.eix.n8-1-1}{8.1.1}) depends functorially on the abelian scheme $A_K$, and that
its formation commutes with every morphism of traits
$\widetilde S\longrightarrow S$. Choose a polarization of $A_K$,
hence an isogeny $A_K\longrightarrow A'_K$. It induces a
homomorphism from (\SGARefLink{sga7.sga7i.eix.n8-1-1}{8.1.1}) to (\SGARefLink{sga7.sga7i.eix.n8-1-6}{8.1.6}), which is immediately seen, by
looking at the $T_\ell$, to be an isogeny term by term. Likewise,
using \SGARefLink{sga7.sga7i.eix.n7-5}{7.5} and descent, one sees that every biextension by
$\mathbb{G}_{mK}$ of abelian schemes $A_K$ and $\overline A_K$ over
$K$ defines a biextension of $(B_K,\overline B_K)$ by
$\mathbb{G}_{mK}$, where $B_K$ and $\overline B_K$ denote the
abelian quotients of the Raynaud extensions corresponding to $A_K$
and $\overline A_K$, respectively. When $W$ is a perfect duality, so
is the corresponding biextension of $(B_K,\overline B_K)$. In
particular, with the notation of \SGARefLink{sga7.sga7i.eix.n7-6}{7.6}, the abelian schemes $B_K$ and
$B'_K$ are canonically dual to one another.

The correspondence between biextensions on $(A_K,\overline A_K)$
and biextensions on $(B_K,\overline B_K)$ has evident functoriality
properties as $A_K$ and $\overline A_K$ vary. These follow by descent
from the analogous properties stated in \SGARefLink{sga7.sga7i.eix.n7-5}{7.5}, and we leave their
formulation to the reader. We also leave to the reader the
compatibility of the preceding construction with every change of
traits $\widetilde S\longrightarrow S$.

\medskip
\noindent
\SGARefTarget{sga7.sga7i.eix.n8-3}{8.3}. Return to the group $A_{K,K'}^{\natural}$ and study its
variation with $K'$. It is clear that two isomorphic extensions
$K'$ and $K'_1$ of $K$ give canonically isomorphic groups, so it is
enough to consider the subextensions $K'$ of the separable closure
$\overline K$ of $K$. If $K'\hookrightarrow K''$ are two such finite
Galois subextensions of $\overline K$, it follows immediately from
\SGARefLink{sga7.sga7i.eix.n3-3}{3.3} that there is a canonical \underline{open immersion}

\begin{equation}\tag{\SGARefTarget{sga7.sga7i.eix.n8-3-1}{8.3.1}}
A^{\natural}_{K,K'} \hookrightarrow A^{\natural}_{K,K''}
 \qquad ,
\end{equation}
which identifies the first group with an open subgroup of the
second. As $K'$ varies, the $A^{\natural}_{K,K'}$ thus form an
inductive system of smooth groups of finite type over $K$, whose
transition morphisms are open immersions. If we set

\begin{equation}\tag{\SGARefTarget{sga7.sga7i.eix.n8-3-2}{8.3.2}}
\Phi_{K,K'}=A^{\natural}_{K,K'}
 /{A^{\natural}_{K,K'}}^{o} \qquad ,
\end{equation}
the $\Phi_{K,K'}$ themselves form an inductive system of finite
etale groups over $K$, which will be studied in greater detail
below.

We shall set

\begin{equation}\tag{\SGARefTarget{sga7.sga7i.eix.n8-3-3}{8.3.3}}
A^{\flat}_{K}=\varinjlim_{K'} A^{\natural}_{K,K'} \qquad ,
\end{equation}
where the limit is taken in the category of fppf sheaves over $K$.
Since the transition morphisms in this inductive system are open
immersions, it is clear moreover that this limit is represented by
a group scheme over $K$ (in general non-quasi-compact), which we
shall naturally denote by the same symbol $A^{\flat}_{K}$. For this
group scheme one has

\begin{equation}\tag{\SGARefTarget{sga7.sga7i.eix.n8-3-4}{8.3.4}}
{A_{K}^{\flat}}^{o}=A_{K}^{\natural o} \quad , \qquad
A^{\flat}_{K}/{A_{K}^{\flat}}^{o}
 \simeq \varinjlim_{K'}\Phi_{K,K'} \qquad .
\end{equation}
The latter group will be determined in \SGARefLink{sga7.sga7i.eix.n11-10}{11.10}: it is canonically
isomorphic to $\underline{M}_{K}\otimes\mathbb{Q}/\mathbb{Z}$, with
the notation of (\SGARefLink{sga7.sga7i.eix.n8-1-8}{8.1.8}).

% END INCLUDED SGA7I_English_current_buildable_source_20260823/components/156_expose_IX_sections_8_2_and_8_3.tex

% BEGIN INCLUDED SGA7I_English_current_buildable_source_20260823/components/157_expose_IX_section_9_1.tex
% Expose IX, section 9 heading and 9.1.
% French transcription lines 2194--2214; authority scan physical pp. 430--431
% (scan indices 429--430; source folios 418--419). Both pages were checked
% directly in four overlapping 1100-dpi bands; no escalation was needed.

\medskip
\noindent
9. \underline{Definition of the monodromy pairing}
$\underline{M}_{\ell}\otimes\underline{M}'_{\ell}
 \longrightarrow \mathbb{Z}_{\ell S}$.

\medskip
\noindent
\SGARefTarget{sga7.sga7i.eix.n9-1}{9.1}. In this section, suppose that $S$ is henselian and that $A_K$
has semistable reduction over $S$, or, equivalently, that $A'_K$
has semistable reduction over $S$ (\SGARefLink{sga7.sga7i.eix.n3-5-1}{3.5.1}). We then have over $S$ two
twisted constant $\mathbb{Z}$-modules $\underline{M}$ and
$\underline{M}'$ (\SGARefLink{sga7.sga7i.eix.n5-5-2}{5.5.2}). Over a suitable finite etale covering of
$S$ (and hence over $S$ itself if $S$ is strictly local, that is, if
$k$ is separably closed), they split into constant groups with
values $M$ and $M'$, respectively, noncanonically isomorphic to
$\mathbb{Z}^{r}$, where $r$ is the common dimension of the maximal
tori of $A_o$ and $A'_o$. For every prime number $\ell$, we shall
consider the corresponding $\ell$-adic sheaves

\begin{equation}\tag{\SGARefTarget{sga7.sga7i.eix.n9-1-1}{9.1.1}}
\underline{M}_{\ell}=\underline{M}\otimes_{\mathbb{Z}}\mathbb{Z}_{\ell}
 \quad , \qquad
\underline{M}'_{\ell}=\underline{M}'\otimes_{\mathbb{Z}}\mathbb{Z}_{\ell}
 \qquad ,
\end{equation}
whose geometric meaning is made explicit by formulas (\SGARefLink{sga7.sga7i.eix.n5-5-7}{5.5.7}) and
(\SGARefLink{sga7.sga7i.eix.n5-5-8}{5.5.8}). In this section we shall define a canonical pairing, called
the \underline{monodromy pairing},

\begin{equation}\tag{\SGARefTarget{sga7.sga7i.eix.n9-1-2}{9.1.2}}
u_{\ell}:\underline{M}_{\ell}\otimes\underline{M}'_{\ell}
 \longrightarrow\mathbb{Z}_{\ell} \qquad ,
\end{equation}
which will be defined in terms of the Barsotti--Tate pro-$\ell$-groups
$T_{\ell}(A^{o})$ and $T_{\ell}({A'}^{o})$, and even in terms of the
generic fiber $T_{\ell}(A_K)$ of the first of these groups (subject,
when $\ell=p$, to the availability of Tate's theorem [33, Th.~4],
that is, for the time being, provided we are then in unequal
characteristic). When $\ell\neq p$, this is equivalent to saying
that the pairing (\SGARefLink{sga7.sga7i.eix.n9-1-2}{9.1.2}) is described in terms of the Galois
$\pi$-module

\begin{equation}\tag{\SGARefTarget{sga7.sga7i.eix.n9-1-3}{9.1.3}}
U=T_{\ell}(A_K)(\overline K)=T_{\ell}(A_{\overline K})
\end{equation}
considered in (\SGARefLink{sga7.sga7i.eix.n2-5-3}{2.5.3}). More precisely, we shall see that knowing the
pairing is essentially equivalent to knowing the action of the
monodromy group $I$ on $U$.

% END INCLUDED SGA7I_English_current_buildable_source_20260823/components/157_expose_IX_section_9_1.tex

% BEGIN INCLUDED SGA7I_English_current_buildable_source_20260823/components/158_expose_IX_section_9_2.tex
% Expose IX, section 9.2.
% French transcription lines 2216--2276; authority scan physical pp. 431--433
% (scan indices 430--432; source folios 419--421). All three pages were
% checked directly in four overlapping 1100-dpi bands; physical p. 434 was
% checked for the exact 9.3 boundary. The authority controls the evident
% transcription repairs x (not u) in the generator phrase, V^\perp (not
% V^1) in (9.2.4), T'_o in the second line of (9.2.7), and balanced
% parentheses in the displayed formulas.

\medskip
\noindent
\SGARefTarget{sga7.sga7i.eix.n9-2}{9.2}. \underline{The case $\ell\neq p$}. Begin with the simpler case
$\ell\neq p$, using the notation introduced in (\SGARefLink{sga7.sga7i.eix.n2-5-3}{2.5.3}). Note that,
up to torsion, the relation $V=U^I$ also means that $V^{\perp}$ is
the module generated by the elements $gx-x$, where $x\in U$ and
$g\in I$, and in any case one has

\begin{equation}\tag{\SGARefTarget{sga7.sga7i.eix.n9-2-1}{9.2.1}}
gx-x\in V^{\perp}\quad\text{for }x\in U,\ g\in I \qquad .
\end{equation}
Thus, for $g\in I$,

\begin{equation}\tag{\SGARefTarget{sga7.sga7i.eix.n9-2-2}{9.2.2}}
gx=x+u(g)\mathbin{.}x \qquad ,
\end{equation}
where $g\longmapsto u(g)$ is a homomorphism from $U$ to
$V^{\perp}$ which vanishes on $V=U^I$, or equivalently

\begin{equation}\tag{\SGARefTarget{sga7.sga7i.eix.n9-2-3}{9.2.3}}
u\in\operatorname{Hom}(U/V,V^{\perp})
 \simeq (U/V)^{\vee}\otimes V^{\perp} \qquad .
\end{equation}
In the case of semistable reduction under consideration, that is,
when $V^{\perp}\subset V$, formula (\SGARefLink{sga7.sga7i.eix.n9-2-2}{9.2.2}) says that the action of
$g$ on $U$ is a \underline{transvection} parallel to
$V^{\perp}=W$, and composition of these transvections is simply
addition of the corresponding homomorphisms $u$ in (\SGARefLink{sga7.sga7i.eix.n9-2-3}{9.2.3}).
Consequently, the action of the inertia group $I$ on $U$ is
completely described by a group homomorphism

\begin{equation}\tag{\SGARefTarget{sga7.sga7i.eix.n9-2-4}{9.2.4}}
I\longrightarrow (U/V)^{\vee}\otimes V^{\perp} \qquad .
\end{equation}

The orthogonality theorem, applied both to $A_K$ and to $A'_K$,
gives

\begin{equation*}
V^{\perp}=W \quad , \qquad
(U/V)^{\vee}\simeq W'(-1) \qquad ,
\end{equation*}
where $W'$ is the toric part of
$U'=T_{\ell}(A'(\overline K))$, and $(-1)$ denotes the Tate twist.
Moreover, the homomorphism (\SGARefLink{sga7.sga7i.eix.n9-2-4}{9.2.4}) factors through
$I(\ell)\simeq\mathbb{Z}_{\ell}(1)$, so that (\SGARefLink{sga7.sga7i.eix.n9-2-4}{9.2.4}) may be
interpreted as a canonical homomorphism

\begin{equation}\tag{\SGARefTarget{sga7.sga7i.eix.n9-2-5}{9.2.5}}
u:\mathbb{Z}_{\ell}(1)\longrightarrow W'\otimes W(-1)
 \qquad ,
\end{equation}
or, equivalently, as an element

\begin{equation}\tag{\SGARefTarget{sga7.sga7i.eix.n9-2-6}{9.2.6}}
u\in W'\otimes W(-2) \qquad .
\end{equation}
Recall that $W$ and $W'$ are the groups of $\overline K$-valued
points of the two twisted constant $\ell$-adic sheaves
$T_{\ell}(A^o)^t$ and $T_{\ell}({A'}^o)^t$, whose special fibers are
$T_{\ell}(T_o)$ and $T_{\ell}(T'_o)$, respectively. We obtain the
canonical isomorphisms

\begin{equation}\tag{\SGARefTarget{sga7.sga7i.eix.n9-2-7}{9.2.7}}
\left\{
\begin{array}{l}
W\simeq T_{\ell}(T_o)(\overline k)
 \simeq\check M\otimes\mathbb{Z}_{\ell}(1),\\[3ex]
W'\simeq T_{\ell}(T'_o)(\overline k)
 \simeq\check M'\otimes\mathbb{Z}_{\ell}(1),
\end{array}
\right.
\end{equation}
where $M$ and $M'$ denote the character groups over $\overline k$
of $T_o$ and $T'_o$, respectively:

\begin{equation}\tag{\SGARefTarget{sga7.sga7i.eix.n9-2-8}{9.2.8}}
T_o=D(M) \quad , \qquad T'_o=D(M') \qquad .
\end{equation}
It follows that the element $u$ of (\SGARefLink{sga7.sga7i.eix.n9-2-6}{9.2.6}), which describes the
action of the inertia group $I$ on the Tate module $U$ of $A_K$, is
interpreted as a canonical element of
$(\check M'\otimes_{\mathbb{Z}}\check M)
 \otimes_{\mathbb{Z}}\mathbb{Z}_{\ell}$. It is moreover invariant
under $\pi_o=\pi/I=\operatorname{Gal}(\overline k/k)$, since by
definition it is invariant under transport of structure:

\begin{equation}\tag{\SGARefTarget{sga7.sga7i.eix.n9-2-9}{9.2.9}}
u\in
\bigl((\check M'\otimes_{\mathbb{Z}}\check M)
 \otimes_{\mathbb{Z}}\mathbb{Z}_{\ell}\bigr)^I
\simeq
\operatorname{Hom}_{\mathbb{Z}}
 (\underline M\otimes\underline M',\mathbb{Z}_{\ell})
\simeq
\operatorname{Hom}_{\mathbb{Z}_{\ell}}
 (\underline M_{\ell}\otimes\underline M'_{\ell},
  \mathbb{Z}_{\ell}) \quad .
\end{equation}
Up to writing $u$ in place of $u_{\ell}$, this is the element
(\SGARefLink{sga7.sga7i.eix.n9-1-2}{9.1.2}) that we wished to define.

To define (\SGARefLink{sga7.sga7i.eix.n9-1-2}{9.1.2}) in the general case, we shall need a few
preliminaries, which will be developed in the next two sections.

% END INCLUDED SGA7I_English_current_buildable_source_20260823/components/158_expose_IX_section_9_2.tex

% BEGIN INCLUDED SGA7I_English_current_buildable_source_20260823/components/159_expose_IX_section_9_3.tex
% Expose IX, section 9.3.
% French transcription lines 2279--2430; authority scan physical pp. 433--439
% (scan indices 432--438; source folios 421--427). The authority pages were
% checked directly in overlapping 1100-dpi bands. Higher-detail inspection was
% reserved for the functor diagram. Evident print defects are corrected and
% disclosed in reader notes: the P/Q order in the graded pieces; F for E twice
% in the Brauer-pushout paragraph; part (d) for part (c); Ext^2(R,Q) for
% Ext^2(Q,P); omega for omega-prime on the lower functor arrow; and Y for Y'
% in the change-of-base formula.

\medskip
\noindent
\SGARefTarget{sga7.sga7i.eix.n9-3}{9.3}. \underline{Mixed extensions.}

Let $\underline C$ be an abelian category, and let

\begin{equation}\tag{\SGARefTarget{sga7.sga7i.eix.n9-3-1}{9.3.1}}
P,\ R,\ Q
\end{equation}
be three objects of $\underline C$. We are interested in classifying the
objects $X$ of $\underline C$ endowed with a three-step filtration

\begin{equation}\tag{*}
X^0=X\supset X^1\supset X^2\supset X^3=0
\end{equation}
and with given isomorphisms\footnote{The printed source assigns $P,R,Q$,
in that order, to these three graded pieces.  The immediately following exact
sequences and every subsequent construction require $Q,R,P$; the evident
interchange of $P$ and $Q$ has therefore been corrected here.}

\begin{equation*}
X^0/X^1\xrightarrow{\ \sim\ }Q,
\qquad X^1/X^2\xrightarrow{\ \sim\ }R,
\qquad X^2/X^3=X^2\xrightarrow{\ \sim\ }P.
\end{equation*}
Putting

\begin{equation*}
E=X^0/X^2,
\qquad F=X^1/X^3=X^1,
\end{equation*}
we see that $X$ gives rise to two ordinary extensions

\begin{equation}\tag{\SGARefTarget{sga7.sga7i.eix.n9-3-2}{9.3.2}}
0\longrightarrow P\longrightarrow F\longrightarrow R\longrightarrow0,
\qquad
0\longrightarrow R\longrightarrow E\longrightarrow Q\longrightarrow0,
\end{equation}
of $R$ by $P$ and of $Q$ by $R$, respectively. Our point of view here will
be to start, in addition to the objects (\SGARefLink{sga7.sga7i.eix.n9-3-1}{9.3.1}), with the extension
structures (\SGARefLink{sga7.sga7i.eix.n9-3-2}{9.3.2}), which are likewise assumed to be given, and to consider
the corresponding three-step filtered structures $X$.

More precisely, a \underline{mixed extension} of the extension $E$ by the
extension $F$ will mean an object $X$ of $\underline C$, endowed with a
three-step filtration (*) and with given isomorphisms

\begin{equation}\tag{\SGARefTarget{sga7.sga7i.eix.n9-3-3}{9.3.3}}
E\simeq X^0/X^2,
\qquad F\simeq X^1/X^3=X^1,
\end{equation}
compatible with the filtrations: they carry $X^1/X^2$ to
$\operatorname{Im}(R\longrightarrow E)$ and $X^2$ to
$\operatorname{Im}(P\longrightarrow F)$, respectively, and the
isomorphism $R\simeq F/P$ (and hence $R\simeq X^1/X^2$) deduced from
(\SGARefLink{sga7.sga7i.eix.n9-3-3}{9.3.3}) is required to equal the one deduced from the first exact sequence
in (\SGARefLink{sga7.sga7i.eix.n9-3-2}{9.3.2}).

For fixed $P,Q,R$ and fixed extensions $E$ and $F$, the mixed extensions of
$E$ by $F$ form a category in the evident way; it is in fact a groupoid
(every arrow is invertible). The automorphism group of a mixed extension
$X$ of $E$ by $F$ is canonically isomorphic to $\operatorname{Hom}(Q,P)$
through

\begin{equation}\tag{\SGARefTarget{sga7.sga7i.eix.n9-3-4}{9.3.4}}
\operatorname{Hom}(Q,P)
 \xrightarrow{\ \sim\ }
\operatorname{Aut}_{\mathrm{extpan}}(X),
\qquad
u\longmapsto\operatorname{id}_X+\bar u,
\end{equation}
where $\bar u$ is the endomorphism of the object of $\underline C$
underlying $X$ obtained by composing
$X\longrightarrow Q\xrightarrow{u}P\longrightarrow X$; the outer arrows
are respectively the canonical epimorphism and the canonical monomorphism
coming from the given mixed-extension structure on $X$.

One could make the structure of the category
$\operatorname{EXTPAN}(E,F)$ of mixed extensions of $E$ by $F$ precise by
regarding it as a ``torsor'' under the ``group-category'' (VII \SGARefLink{sga7.sga7i.evii.n2-5}{2.5})
$\operatorname{EXT}(Q,P)$ of extensions of $Q$ by $P$, using a natural
functor analogous to the Brauer composition of ordinary extensions:

\begin{equation}\tag{\SGARefTarget{sga7.sga7i.eix.n9-3-5}{9.3.5}}
\begin{aligned}
\omega:
\operatorname{EXT}(Q,P)\times\operatorname{EXTPAN}(E,F)
 &\longrightarrow \operatorname{EXTPAN}(E,F),\\
\omega(G,X)&=G\wedge X.
\end{aligned}
\end{equation}
For every object $X$ of $\operatorname{EXTPAN}(E,F)$, the partial functor
$G\longmapsto G\wedge X$ is an equivalence from
$\operatorname{EXT}(Q,P)$ to $\operatorname{EXTPAN}(E,F)$.

We shall give only a brief description of $\omega$. Regard a mixed
extension $X$ of $E$ by $F$ as defining an ordinary extension, also denoted
$X$, of $Q$ by $F$ (viewed as an object of $\underline C$). For every
extension $G$ of $Q$ by $P$, let $\bar G$ denote the extension of $Q$ by
$F$ induced by the inclusion $P\longrightarrow F$. Set

\begin{equation}\tag{\SGARefTarget{sga7.sga7i.eix.n9-3-6}{9.3.6}}
G\wedge X=\omega(G,X)
 \overset{\mathrm{dfn}}{=}\bar G\wedge X,
\end{equation}
where $\wedge$ denotes Brauer composition in the category of extensions of
$Q$ by $F$. The three-step filtration of the right-hand side $X'$ of
(\SGARefLink{sga7.sga7i.eix.n9-3-6}{9.3.6}) is clear, as is the second isomorphism of type (\SGARefLink{sga7.sga7i.eix.n9-3-3}{9.3.3}) attached to
$X'$.

To define the first isomorphism of type (\SGARefLink{sga7.sga7i.eix.n9-3-3}{9.3.3}) for $X'$, observe that the
functor

\begin{equation*}
Y\longmapsto\widetilde Y,
\qquad
\operatorname{EXT}(Q,F)\longrightarrow\operatorname{EXT}(Q,R),
\end{equation*}
induced by the given morphism $F\longrightarrow R$ in (\SGARefLink{sga7.sga7i.eix.n9-3-2}{9.3.2}) is
compatible with Brauer composition. It therefore sends $X'$ to an
extension $\widetilde X'$ canonically isomorphic to
$\widetilde{\bar G}\wedge\widetilde X$. The first isomorphism (\SGARefLink{sga7.sga7i.eix.n9-3-3}{9.3.3})
gives an isomorphism $\widetilde X\simeq E$, while
$\widetilde{\bar G}$---the extension induced from $G$ by the composite
$P\longrightarrow F\longrightarrow R$---is canonically isomorphic to the
trivial extension, since that composite $P\longrightarrow R$ is zero.
Consequently there is a canonical isomorphism
$\widetilde X'\simeq E$, which specifies on the right-hand side of (\SGARefLink{sga7.sga7i.eix.n9-3-6}{9.3.6})
its mixed-extension structure of $E$ by $F$.
\footnote{At both occurrences in this paragraph the printed source has
$F$ where the objects involved are extensions of $Q$ by $R$ and hence must
be identified with $E$.}

We leave to the reader the definition of $\omega$ on the arrows of
$\operatorname{EXT}(Q,P)\times\operatorname{EXTPAN}(E,F)$, the
verification that (\SGARefLink{sga7.sga7i.eix.n9-3-5}{9.3.5}) has the asserted property---namely, for every
\underline{fixed} mixed extension $X$, the functor
$G\longmapsto G\wedge X$ is an equivalence from
$\operatorname{EXT}(Q,P)$ to $\operatorname{EXTPAN}(E,F)$---and the proof
of the associativity formula

\begin{equation}\tag{\SGARefTarget{sga7.sga7i.eix.n9-3-7}{9.3.7}}
G_1\wedge(G_2\wedge X)
 \simeq (G_1\wedge G_2)\wedge X,
\end{equation}
where $G_1\wedge G_2$ denotes Brauer composition in
$\operatorname{EXT}(Q,P)$.

The preceding discussion immediately gives parts (a) and (b) of the
following result.

\medskip
\noindent
\underline{Proposition \SGARefTarget{sga7.sga7i.eix.n9-3-8}{9.3.8}. Given two extensions $E$ and $F$ as in
(\SGARefLink{sga7.sga7i.eix.n9-3-2}{9.3.2}) in the abelian category $\underline C$, consider the category
$\operatorname{EXTPAN}(E,F)$ of mixed extensions of $E$ by $F$. Then:}

\smallskip
\noindent
a) \underline{$\operatorname{EXTPAN}(E,F)$ is a groupoid. For every one
of its objects the automorphism group is canonically isomorphic to}
$\operatorname{Ext}^0(Q,P)=\operatorname{Hom}(Q,P)$.

\smallskip
\noindent
b) \underline{The set $\operatorname{Extpan}(E,F)$ of isomorphism classes
of objects of $\operatorname{EXTPAN}(E,F)$ is empty, or is naturally a
torsor under the group} $\operatorname{Ext}^1(Q,P)$.

\smallskip
\noindent
c) \underline{One can canonically define a class}

\begin{equation}\tag{\SGARefTarget{sga7.sga7i.eix.n9-3-9}{9.3.9}}
c(E,F)\in\operatorname{Ext}^2(Q,P),
\end{equation}
\underline{whose vanishing is necessary and sufficient for}
$\operatorname{Extpan}(E,F)\neq\varnothing$, \underline{that is, for the
existence of a mixed extension of $E$ by $F$.}

It remains to prove (c).\footnote{The printed source says ``prove (d)'',
although Proposition \SGARefLink{sga7.sga7i.eix.n9-3-8}{9.3.8} has only parts (a), (b), and (c).} A mixed
extension of $E$ by $F$ may also be described as an extension $X$ of $E$
by $P$ together with an isomorphism of extensions of $R$ by $P$

\begin{equation*}
X^1\simeq F,
\end{equation*}
where $X^1$ denotes the image of the extension $X$ under the morphism
$R\longrightarrow E$ in (\SGARefLink{sga7.sga7i.eix.n9-3-2}{9.3.2}). Such a mixed extension therefore exists
if and only if there is an element of $\operatorname{Ext}^1(E,P)$ whose
image in $\operatorname{Ext}^1(R,P)$ is the class $\xi$ of $F$. Consider
the exact sequence of the $\operatorname{Ext}^i(-,P)$ associated with the
second exact sequence in (\SGARefLink{sga7.sga7i.eix.n9-3-2}{9.3.2})\footnote{The printed final term is
$\operatorname{Ext}^2(R,Q)$; applying $\operatorname{Ext}(-,P)$ to
$0\to R\to E\to Q\to0$ gives the displayed $\operatorname{Ext}^2(Q,P)$,
which is also the group specified in (\SGARefLink{sga7.sga7i.eix.n9-3-9}{9.3.9}).}:

\begin{equation*}
\cdots\longrightarrow
\operatorname{Ext}^1(E,P)\longrightarrow
\operatorname{Ext}^1(R,P)
\xrightarrow{\partial}
\operatorname{Ext}^2(Q,P)\longrightarrow\cdots .
\end{equation*}
It is enough to define $c(E,F)$ by

\begin{equation}\tag{\SGARefTarget{sga7.sga7i.eix.n9-3-10}{9.3.10}}
c(E,F)=\partial\xi,
\qquad
\xi=\operatorname{cl}(F)\in\operatorname{Ext}^1(R,P).
\end{equation}

Now suppose that we are given an exact functor

\begin{equation}\tag{\SGARefTarget{sga7.sga7i.eix.n9-3-11}{9.3.11}}
r:\underline C\longrightarrow\underline C',
\end{equation}
where $\underline C'$ is another abelian category. Write
$P',Q',R',E',F'$ for the objects of $\underline C'$ obtained from
$P,Q,R,E,F$ by applying $r$, so that $E'$ is an extension of $Q'$ by
$R'$ and $F'$ is an extension of $R'$ by $P'$. Up to canonical
isomorphism, the following diagram of functors is commutative:

\begin{equation}\tag{\SGARefTarget{sga7.sga7i.eix.n9-3-12}{9.3.12}}
\begin{tikzcd}[column sep=large,row sep=large]
\operatorname{EXT}(Q,P)\times\operatorname{EXTPAN}(E,F)
  \arrow[r,"\omega"]\arrow[d]
& \operatorname{EXTPAN}(E,F)\arrow[d]\\
\operatorname{EXT}(Q',P')\times\operatorname{EXTPAN}(E',F')
  \arrow[r,"\omega'"']
& \operatorname{EXTPAN}(E',F').
\end{tikzcd}
\end{equation}
The vertical arrows are the evident functors induced by the exact functor
$r$, and $\omega'$ is the arrow analogous to $\omega$ in (\SGARefLink{sga7.sga7i.eix.n9-3-5}{9.3.5}), with
$\underline C$ replaced by $\underline C'$. The lower arrow in the
printed diagram is labelled $\omega$ without a prime; the explanatory
sentence identifies it as $\omega'$, which is the notation used here.
We obtain analogous compatibilities, with the functor $r$, of the
structures considered in \SGARefLink{sga7.sga7i.eix.n9-3-8}{9.3.8}; their explicit formulation is left to the
reader.

We shall use the following consequence of this commutativity. Suppose that
$\operatorname{EXTPAN}(E,F)$ is nonempty and choose an object $X$ in it.
Then, by means of $G'\longmapsto G'\wedge X'$, where $X'=r(X)$, the
category $\operatorname{EXTPAN}(E',F')$ is equivalent to
$\operatorname{EXT}(Q',P')$. Thus an object $Y'$ of
$\operatorname{EXTPAN}(E',F')$ is known up to isomorphism when the
isomorphism class of the corresponding object of
$\operatorname{EXT}(Q',P')$ is known, that is, when the corresponding
element

\begin{equation*}
c_X(Y')\in\operatorname{Ext}^1(Q',P')
\end{equation*}
is known.

If $X$ is replaced by another object $X_1=H\wedge X$, where $H$ is an
extension of $Q$ by $P$, then one immediately obtains
\footnote{The printed right-hand side of the following formula has
$c_X(Y)$; the fixed object in this construction is $Y'$.}

\begin{equation*}
c_{X_1}(Y')=c_X(Y')-\operatorname{cl}(H')
\end{equation*}
from the associativity formula (\SGARefLink{sga7.sga7i.eix.n9-3-7}{9.3.7}). Hence, for fixed $Y'$ and varying
$X$, the elements $c_X(Y')$ run through exactly one coset

\begin{equation}\tag{\SGARefTarget{sga7.sga7i.eix.n9-3-13}{9.3.13}}
c(Y')\in
\operatorname{Ext}^1(Q',P')/
\operatorname{Im}\operatorname{Ext}^1(Q,P),
\end{equation}
which is therefore canonically associated with the mixed extension $Y'$ of
$E'$ by $F'$. By construction, its vanishing is necessary and sufficient
for $Y'$ to be isomorphic to the transform under $r$ of a mixed extension
$X$ of $E$ by $F$.

% END INCLUDED SGA7I_English_current_buildable_source_20260823/components/159_expose_IX_section_9_3.tex

% BEGIN INCLUDED SGA7I_English_current_buildable_source_20260823/components/160_expose_IX_section_9_4.tex
% Expose IX, section 9.4.
% French transcription lines 2435--2576; authority scan physical pp. 439--444
% (scan indices 438--443; source folios 427--432). The authority pages were
% checked directly in overlapping 1100-dpi bands. The scan is clear throughout;
% no higher-resolution escalation was required. Evident print/transcription
% defects are corrected and disclosed at the point of use: omitted checks and
% delimiters, Lambda_s misread as a wedge, Q/P reversals in Hom and Ext, and the
% nonexistent reference 9.4.6.2 for 9.4.7.2.

\medskip
\noindent
\SGARefTarget{sga7.sga7i.eix.n9-4}{9.4}. Let $S$ be a henselian trait with generic point $\eta$, let $n$ be an
integer $>0$, put $\Lambda=\mathbb Z/n\mathbb Z$, and let $\underline C$ be
the category of $\Lambda$-modules on $S$ (for the fppf site on $S$; cf.
VIII \SGARefLink{sga7.sga7i.eviii.n0-2}{0.2}) and $\underline C'$ the category of $\Lambda$-modules on $\eta$.
Let $P,R,Q$ be objects of $\underline C$, and let $E,F$ be extensions of
$Q$ by $R$ and of $R$ by $P$, respectively (\SGARefLink{sga7.sga7i.eix.n9-3-2}{9.3.2}). Suppose that $Q$ is a
locally free $\Lambda_S$-module of finite type and that $P$ is represented
by a finite group of multiplicative type over $S$ (SGA 3 IX), so that it is
of the form

\begin{equation*}
P=D(Q^*),
\end{equation*}
where $D=\underline{\operatorname{Hom}}(-,\mathbb G_{mS})$ denotes Cartier
duality. Let $P_\eta,R_\eta,Q_\eta,E_\eta,F_\eta$ denote the restrictions of
$P,R,Q,E,F$ to $\eta$, and suppose that a mixed extension $X_\eta$ of
$E_\eta$ by $F_\eta$ is given. Under these conditions we shall canonically
define a pairing

\begin{equation}\tag{\SGARefTarget{sga7.sga7i.eix.n9-4-1}{9.4.1}}
c(X_\eta):Q\otimes_\Lambda Q^*\longrightarrow\Lambda_S
\qquad(\Lambda=\mathbb Z/n\mathbb Z).
\end{equation}

For this purpose, first \underline{suppose} that $S$ is \underline{strictly local}. Then $Q$
and $Q^*$ are constant, say $Q=Q_{oS}$ and $Q^*=Q^*_{oS}$, and giving
(\SGARefLink{sga7.sga7i.eix.n9-4-1}{9.4.1}) amounts to giving an element

\begin{equation}\tag{\SGARefTarget{sga7.sga7i.eix.n9-4-2}{9.4.2}}
c(X_\eta)\in\check Q_o\otimes\check Q_o^*,
\end{equation}
where the check denotes the ordinary dual $\operatorname{Hom}(-,\Lambda)$.
We apply the considerations of \SGARefLink{sga7.sga7i.eix.n9-3}{9.3} with the restriction functor
$r:\underline C\longrightarrow\underline C'$. First observe that a mixed
extension of $E$ by $F$ exists. Indeed, the obstruction to its existence
lies in $\operatorname{Ext}^2(Q,P)$ (\SGARefLink{sga7.sga7i.eix.n9-3-8}{9.3.8}(c)), and it is enough to prove

\begin{equation*}
\operatorname{Ext}^2(Q,P)=0.
\end{equation*}

This is a special case of the following lemma.

\medskip
\noindent
\underline{Lemma \SGARefTarget{sga7.sga7i.eix.n9-4-3}{9.4.3}. Under the preceding hypotheses, with $S$ strictly
local, one has:}

\smallskip
\noindent
a) \underline{canonical isomorphisms}

\begin{equation*}
\operatorname{Ext}^i_\Lambda(Q,P)\simeq
H^i(S,\check Q\otimes P),
\qquad
\operatorname{Ext}^i_\Lambda(Q_\eta,P_\eta)\simeq
H^i(\eta,\check Q\otimes P);
\end{equation*}

\noindent
b) \underline{a canonical isomorphism}

\begin{equation}\tag{b)}
\check Q\otimes P\simeq D(Q\otimes Q^*);
\end{equation}

\noindent
c) \underline{for every finite constant group $H=H_{oS}$ over $S$,}

\begin{equation}\tag{c)}
H^i(S,D(H))\simeq H^i(\eta,D(H)_\eta)=0
\qquad\underline{\text{for }}i\geq2;
\end{equation}

\noindent
d) \underline{canonical isomorphisms}

\begin{equation}\tag{d)}
\begin{aligned}
H^1(S,D(H))&\simeq\mathbb G_m(S)\otimes\check H_o
              =V^*\otimes\check H_o,\\
H^1(\eta,D(H)_\eta)&\simeq\mathbb G_m(\eta)\otimes\check H_o
              =K^*\otimes\check H_o,
\end{aligned}
\end{equation}
\underline{where the check in (d) denotes $\operatorname{Hom}(-,\mathbb
Q/\mathbb Z)$.}

The formulas in (a) and (b) follow immediately from the fact that $Q$ is a
locally free $\Lambda_S$-module of finite type. To prove (c), we plainly
reduce to the case $H_o=\mathbb Z/m\mathbb Z$, that is,
$D(H)=\mu_m$. The cohomology sequence associated with the Kummer exact
sequence

\begin{equation*}
0\longrightarrow\mu_m\longrightarrow\mathbb G_m
 \xrightarrow{\ m\cdot\operatorname{id}\ }\mathbb G_m
 \longrightarrow0
\end{equation*}
then reduces us to proving

\begin{equation*}
H^i(S,\mathbb G_m)=H^i(\eta,\mathbb G_m)=0
\qquad\text{for }i\geq1.
\end{equation*}
These groups are the groups $H^i_{\mathrm{\acute et}}$ computed on the
étale site [6, Theorem \SGARefLink{sga7.sga7i.eix.n11-7}{11.7}], and therefore vanish for $S$, since $S$ is
strictly local. Their vanishing for $\eta$ is likewise well known
[6, Theorem \SGARefLink{sga7.sga7i.eix.n1-1}{1.1}].

It remains to describe the isomorphisms in (d). Over an arbitrary base $U$
there is a canonical homomorphism

\begin{equation}\tag{\SGARefTarget{sga7.sga7i.eix.n9-4-3-1}{9.4.3.1}}
\mathbb G_m(U)\otimes\check H_o
\longrightarrow H^1(U,D(H_{oU})).
\end{equation}
If $H_o$ is annihilated by an integer $m>0$, this homomorphism is described
using the boundary map
$\mathbb G_m(U)\xrightarrow{\partial}H^1(U,\mu_m)$ associated with the
Kummer sequence: it is the composite

\begin{equation*}
\begin{aligned}
\mathbb G_m(U)\otimes\check H_o
&\longrightarrow H^1(U,\mu_m)\otimes\check H_o\\
&\longrightarrow H^1(U,\mu_m\otimes\check H_o)
 \longrightarrow H^1(U,D(H_{oU})).
\end{aligned}
\end{equation*}
Here the second arrow is deduced from the evident homomorphism
$\check H_o\longrightarrow
\operatorname{Hom}(\mu_m,\mu_m\otimes\check H_o)$, and the third is a
special case of (b), with $m$ in place of $n$. We leave it to the reader to
check that (\SGARefLink{sga7.sga7i.eix.n9-4-3-1}{9.4.3.1}) is independent of the choice of the integer $m$ such
that $mH_o=0$. We also claim that (\SGARefLink{sga7.sga7i.eix.n9-4-3-1}{9.4.3.1}) is always injective and that its
cokernel vanishes if $\operatorname{Pic}(U)=0$. (In fact, the cokernel is
always isomorphic to $\operatorname{Hom}(H_o,\operatorname{Pic}(U))$, but
this is immaterial here.) Indeed, we reduce to the case
$H_o=\mathbb Z/n\mathbb Z$, and the assertion then follows immediately from
the cohomology sequence of the Kummer sequence above. This proves (d), and
immediately yields:

\medskip
\noindent
\underline{Corollary \SGARefTarget{sga7.sga7i.eix.n9-4-4}{9.4.4}. With the notation of \SGARefLink{sga7.sga7i.eix.n9-4-3}{9.4.3}, there is a canonical
isomorphism}

\noindent
e)
\begin{equation*}
H^1(\eta,D(H)_\eta)/\operatorname{Im}H^1(S,D(H))
\simeq\check H_o,
\end{equation*}
\underline{and hence, on combining (a) and (b), a canonical isomorphism}

\begin{equation}\tag{\SGARefTarget{sga7.sga7i.eix.n9-4-5}{9.4.5}}
\operatorname{Ext}^1(Q_\eta,P_\eta)/
\operatorname{Im}\operatorname{Ext}^1(Q,P)
\simeq\check Q_o\otimes\check Q_o^*.
\end{equation}

Returning to the definition of (\SGARefLink{sga7.sga7i.eix.n9-4-2}{9.4.2}) when $S$ is strictly local, the
relation $\operatorname{Ext}^2(Q,P)=0$ makes it possible to define the
canonical element (\SGARefLink{sga7.sga7i.eix.n9-3-13}{9.3.13}), taking for $Y'$ the given mixed extension
$X_\eta$ of $E_\eta$ by $F_\eta$. Formula (\SGARefLink{sga7.sga7i.eix.n9-4-5}{9.4.5}) identifies this element
$c(X_\eta)$ with an element of $\check Q_o\otimes\check Q_o^*$, thereby
completing the definition of (\SGARefLink{sga7.sga7i.eix.n9-4-2}{9.4.2}). The general considerations of \SGARefLink{sga7.sga7i.eix.n9-3}{9.3}
also show that this element is the obstruction to extending $X_\eta$ to a
mixed extension $X$ of $E$ by $F$.

\medskip
\noindent
\SGARefTarget{sga7.sga7i.eix.n9-4-6}{9.4.6}. The definition of (\SGARefLink{sga7.sga7i.eix.n9-4-1}{9.4.1}) when $S$ is henselian now reduces
immediately to the strictly local case by Galois descent from the strict
henselization of $S$.

\medskip
\noindent
\underline{Proposition \SGARefTarget{sga7.sga7i.eix.n9-4-7}{9.4.7}. Let $P,Q,R,E,F$ be as above, except that the
trait $S$ need not be henselian. Consider the functor}

\begin{equation}\tag{\SGARefTarget{sga7.sga7i.eix.n9-4-7-1}{9.4.7.1}}
X\longmapsto X_\eta
\end{equation}
\underline{from the category of mixed extensions of $E$ by $F$ to the
category of mixed extensions of $E_\eta$ by $F_\eta$. For every mixed
extension $X_\eta$ of $E_\eta$ by $F_\eta$, consider the pairing}

\begin{equation}\tag{\SGARefTarget{sga7.sga7i.eix.n9-4-7-2}{9.4.7.2}}
c_o(X_\eta):Q_o\otimes Q_o^*\longrightarrow\Lambda_s
\end{equation}
\underline{induced by the pairing of type (\SGARefLink{sga7.sga7i.eix.n9-4-1}{9.4.1}) defined over the
henselization $S^h$ of $S$. Then the functor (\SGARefLink{sga7.sga7i.eix.n9-4-7-1}{9.4.7.1}) is fully faithful,
and a mixed extension $X_\eta$ of $E_\eta$ by $F_\eta$ belongs to its
essential image if and only if the pairing (\SGARefLink{sga7.sga7i.eix.n9-4-7-2}{9.4.7.2}) is zero.}

The printed target of (\SGARefLink{sga7.sga7i.eix.n9-4-7-2}{9.4.7.2}) is $\Lambda_s$; the French transcription's
$\wedge_s$ is a character-recognition error.

\medskip
\noindent
\SGARefTarget{sga7.sga7i.eix.n9-4-7-3}{9.4.7.3}. This statement still makes sense and remains valid if one assumes
only that $S$ is a connected regular noetherian scheme of dimension $1$.
Instead of the single pairing (\SGARefLink{sga7.sga7i.eix.n9-4-7-2}{9.4.7.2}), one must then consider a pairing
$Q_s\otimes Q_s^*\longrightarrow\Lambda_s$ for each closed point $s$ of
$S$. It will be most convenient to prove \SGARefLink{sga7.sga7i.eix.n9-4-7}{9.4.7} in this more general form.

\smallskip
\noindent
a) For every mixed extension $X$, the map

\begin{equation*}
\operatorname{End}(X)\longrightarrow\operatorname{End}(X_\eta)
\end{equation*}
is bijective. The two sides are respectively the sections over $S$ and
$\eta$ of
$\underline{\operatorname{Hom}}(Q,P)=D(Q\otimes Q^*)$; hence the assertion
follows from the normality of $S$ and the finiteness of
$D(Q\otimes Q^*)$ over $S$.\footnote{The printed proof has
$\underline{\operatorname{Hom}}(P,Q)$ here. Proposition 9.3.8(a), formula
(\SGARefLink{sga7.sga7i.eix.n9-3-4}{9.3.4}), and the displayed identification with $D(Q\otimes Q^*)$ all
require $\underline{\operatorname{Hom}}(Q,P)$.} It follows at once that if $X$
and $Y$ are two mixed extensions already known to be isomorphic, then

\begin{equation*}
\operatorname{Hom}(X,Y)\longrightarrow
\operatorname{Hom}(X_\eta,Y_\eta)
\end{equation*}
is bijective.

\smallskip
\noindent
b) If $X$ and $Y$ are two mixed extensions, every isomorphism
$X_\eta\xrightarrow{\sim}Y_\eta$ comes from an isomorphism
$X\longrightarrow Y$. By (a), the question is local for the étale
topology, so we may reduce to the case in which $S$ is strictly local. In
that case \SGARefLink{sga7.sga7i.eix.n9-4-3}{9.4.3} shows that the map

\begin{equation*}
\operatorname{Ext}^1_\Lambda(Q,P)
\longrightarrow\operatorname{Ext}^1_\Lambda(Q_\eta,P_\eta)
\end{equation*}
identifies with

\begin{equation*}
V^*\otimes\check H_o\longrightarrow K^*\otimes\check H_o,
\qquad H_o=Q_o\otimes Q_o^*.
\end{equation*}
The latter is injective, since the exact sequence
$0\longrightarrow V^*\longrightarrow K^*\longrightarrow\mathbb Z
\longrightarrow0$ is split. By the discussion of \SGARefLink{sga7.sga7i.eix.n9-3}{9.3} this implies that if
$X_\eta$ and $Y_\eta$ are isomorphic, then so are $X$ and $Y$; the
conclusion now follows from (a).\footnote{The printed Ext groups in this
paragraph have the arguments $(P,Q)$. The action of
$\operatorname{Ext}^1(Q,P)$ on mixed extensions in Proposition \SGARefLink{sga7.sga7i.eix.n9-3-8}{9.3.8}(b),
and Lemma \SGARefLink{sga7.sga7i.eix.n9-4-3}{9.4.3} itself, require $(Q,P)$.}

Thus (a) and (b) show that $X\longmapsto X_\eta$ is \underline{fully} faithful.

\smallskip
\noindent
c) Let $X_\eta$ be a mixed extension over $\eta$. We prove that it extends
to $S$ if and only if the pairings (\SGARefLink{sga7.sga7i.eix.n9-4-7-2}{9.4.7.2}) vanish. Here again, the
uniqueness already proved shows that the question is local for the étale
topology, reducing us once more to the strictly local case. The conclusion
then holds by construction, as observed immediately before \SGARefLink{sga7.sga7i.eix.n9-4-6}{9.4.6}.
\footnote{The printed proof cites the nonexistent pairings
``(9.4.6.2)''; these are the pairings (\SGARefLink{sga7.sga7i.eix.n9-4-7-2}{9.4.7.2}).}

This completes the proof of \SGARefLink{sga7.sga7i.eix.n9-4-7}{9.4.7}.

% END INCLUDED SGA7I_English_current_buildable_source_20260823/components/160_expose_IX_section_9_4.tex

% BEGIN INCLUDED SGA7I_English_current_buildable_source_20260823/components/161_expose_IX_section_9_5.tex
% Expose IX, section 9.5.
% French transcription lines 2578--2653; authority scan physical pp. 444--448
% (scan indices 443--447; source folios 432--436). The authority pages were
% checked directly in four overlapping 1100-dpi bands. The scan is clear
% throughout. Four internally forced print slips are corrected and disclosed:
% 9.4.6 for 9.4.7, X_eta for its extension X over S, eta for S as the base of
% the etale group Q(n), and 9.5.6 for the homomorphism (9.5.5).

\medskip
\noindent
\SGARefTarget{sga7.sga7i.eix.n9-5}{9.5}. Fix a prime number $\ell$, pro-$\ell$ Barsotti--Tate groups over $S$

\begin{equation}\tag{\SGARefTarget{sga7.sga7i.eix.n9-5-1}{9.5.1}}
P=(P(n))_{n\geq0},
\qquad R=(R(n))_{n\geq0},
\qquad Q=(Q(n))_{n\geq0},
\end{equation}
and extensions of pro-$\ell$ Barsotti--Tate groups

\begin{equation}\tag{\SGARefTarget{sga7.sga7i.eix.n9-5-2}{9.5.2}}
0\longrightarrow P\longrightarrow F\longrightarrow R\longrightarrow0,
\qquad
0\longrightarrow R\longrightarrow E\longrightarrow Q\longrightarrow0.
\end{equation}
Thus, for every integer $n\geq0$, these define extensions of
$\Lambda_n$-modules, where
$\Lambda_n=\mathbb Z/\ell^{\,n+1}\mathbb Z$:

\begin{equation}\tag{\SGARefTarget{sga7.sga7i.eix.n9-5-3}{9.5.3}}
\begin{gathered}
0\longrightarrow P(n)\longrightarrow F(n)\longrightarrow R(n)
 \longrightarrow0,\\
0\longrightarrow R(n)\longrightarrow E(n)\longrightarrow Q(n)
 \longrightarrow0.
\end{gathered}
\end{equation}
We assume that $Q$ is of étale type and $P$ of multiplicative type.
Finally, suppose that a pro-$\ell$ Barsotti--Tate group $X_\eta$ is given
over $\eta$ which is a \underline{mixed extension} of $E_\eta$ by $F_\eta$. By this
we mean that it is endowed, as a pro-$\ell$ Barsotti--Tate group, with a
three-step filtration whose partial extensions are supplied with
isomorphisms to $E_\eta$ and $F_\eta$, respectively. It follows that, for
every $n$, $X_\eta(n)$ is endowed with the structure of a mixed extension
of $E(n)_\eta$ by $F(n)_\eta$. We may therefore associate with it a pairing
of type (\SGARefLink{sga7.sga7i.eix.n9-4-1}{9.4.1}):

\begin{equation}\tag{\SGARefTarget{sga7.sga7i.eix.n9-5-4}{9.5.4}}
c(X(n)_\eta):Q(n)\otimes Q^*(n)
\longrightarrow\Lambda_n=\mathbb Z/\ell^{\,n+1}\mathbb Z,
\end{equation}
where the $Q^*(n)$ are the finite étale groups over $S$ defined by

\begin{equation*}
P(n)\simeq D(Q^*(n)),
\qquad\text{i.e.}\qquad Q^*(n)=D(P(n)),
\end{equation*}
and $D$ again denotes Cartier duality. The $Q^*(n)$ form a pro-$\ell$ group
in the evident way, denoted by $Q^*$. We leave it to the reader to verify
that, as $n$ varies, the homomorphisms (\SGARefLink{sga7.sga7i.eix.n9-5-4}{9.5.4}) define a homomorphism of
projective systems, that is, a homomorphism of twisted constant
$\ell$-adic sheaves over $S$:

\begin{equation}\tag{\SGARefTarget{sga7.sga7i.eix.n9-5-5}{9.5.5}}
c(X_\eta):Q\otimes Q^*\longrightarrow\mathbb Z_{\ell S}.
\end{equation}

We mention in passing the following result, which will not be used below.

\medskip
\noindent
\underline{Proposition \SGARefTarget{sga7.sga7i.eix.n9-5-6}{9.5.6}. Let $P,Q,R,E,F$ be as above, except that the
trait $S$ need not be henselian. Consider the functor}

\begin{equation}\tag{\SGARefTarget{sga7.sga7i.eix.n9-5-6-1}{9.5.6.1}}
X\longmapsto X_\eta
\end{equation}
\underline{from the category of Barsotti--Tate mixed extensions of $E$ by
$F$ to the category of mixed extensions of $E_\eta$ by $F_\eta$. For every
Barsotti--Tate group $X_\eta$ which is a mixed extension of $E_\eta$ by
$F_\eta$, consider the corresponding pairing}

\begin{equation}\tag{\SGARefTarget{sga7.sga7i.eix.n9-5-6-2}{9.5.6.2}}
Q_o\otimes Q_o^*\longrightarrow\mathbb Z_{\ell s}
\end{equation}
\underline{induced by the pairing of type (\SGARefLink{sga7.sga7i.eix.n9-5-5}{9.5.5}) defined over the
henselization of $S$. Then the functor (\SGARefLink{sga7.sga7i.eix.n9-5-6-1}{9.5.6.1}) is fully faithful, and its
essential image consists of those Barsotti--Tate groups $X_\eta$ which are
mixed extensions of $E_\eta$ by $F_\eta$ and for which the corresponding
pairing (\SGARefLink{sga7.sga7i.eix.n9-5-6-2}{9.5.6.2}) is zero. Moreover, when $\operatorname{char}K=0$ (so that
Tate's theorem [33, Theorem 4] applies), $X_\eta$ belongs to the essential
image if and only if it extends to a Barsotti--Tate group over $S$.}

\medskip
\noindent
\SGARefTarget{sga7.sga7i.eix.n9-5-6-3}{9.5.6.3}. We leave it to the reader to extend this statement to the case in
which $S$ is assumed only to be a connected regular noetherian scheme of
dimension $1$.

The first assertion of \SGARefLink{sga7.sga7i.eix.n9-5-6}{9.5.6} follows immediately from \SGARefLink{sga7.sga7i.eix.n9-4-7}{9.4.7}, applied to the
components of the Barsotti--Tate groups under consideration.
\footnote{The printed proof cites \SGARefLink{sga7.sga7i.eix.n9-4-6}{9.4.6}, which only performs the henselian
descent construction; the full-faithfulness and essential-image statement
being applied componentwise is Proposition \SGARefLink{sga7.sga7i.eix.n9-4-7}{9.4.7}.}
It remains to prove that if a Barsotti--Tate group $X_\eta$, which is a
mixed extension of $E_\eta$ by $F_\eta$, extends to a Barsotti--Tate group
$X$ over $S$, then it also extends as a mixed extension.
\footnote{At this point the printed source gives the extension over $S$ the
name $X_\eta$ as well. We reserve $X_\eta$ for the generic fibre and write
$X$ for its extension over $S$.}
By Tate's theorem [33], the homomorphism
$X_\eta\longrightarrow Q_\eta$ extends uniquely to a homomorphism

\begin{equation*}
X\longrightarrow Q.
\end{equation*}
For every integer $n$, the corresponding homomorphism
$X(n)\longrightarrow Q(n)$ is surjective on the generic fibres and hence is
surjective. Since $Q(n)$ is étale over $S$, the map
$X(n)\longrightarrow Q(n)$ is flat, and consequently it is an epimorphism.
Thus $X\longrightarrow Q$ is an \underline{epimorphism component by component}.
\footnote{The printed sentence says that $Q(n)$ is étale over $\eta$.
The group $Q(n)$ was defined over $S$ and is assumed of étale type there;
the base required for this flatness argument is $S$.}
Let $\overline F$ be the kernel of $X\longrightarrow Q$. Then
$\overline F$ is a Barsotti--Tate group and $X$ is an extension of $Q$ by
$\overline F$. The isomorphism
$F_\eta\xrightarrow{\sim}\overline F_\eta$ extends, by Tate's theorem, to
an isomorphism $F\simeq\overline F$. Hence $X$ is an extension of $Q$ by
$F$. Consider now the quotient $\overline E=X/P$. A further application of
Tate's theorem gives an isomorphism $E\simeq\overline E$. It follows that
$X$ is a mixed extension of $E$ by $F$ extending the mixed extension
$X_\eta$, as required.

\medskip
\noindent
\underline{Remarks \SGARefTarget{sga7.sga7i.eix.n9-5-6-4}{9.5.6.4}.}

\smallskip
\noindent
a) Of course, the last assertion of \SGARefLink{sga7.sga7i.eix.n9-5-6}{9.5.6} should hold without restriction
on $K$, since the same should be true of the invoked theorem of Tate.

\smallskip
\noindent
b) One can give another proof of the good-reduction criterion \SGARefLink{sga7.sga7i.eix.n5-10}{5.10}. As in
\SGARefLink{sga7.sga7i.eix.n5-10}{5.10}, reduce to the case where $A_K$ has semistable reduction, and then
apply \SGARefLink{sga7.sga7i.eix.n9-5-6}{9.5.6} to $T_\ell(A_K)$ regarded as a mixed extension (cf. \SGARefLink{sga7.sga7i.eix.n9-6}{9.6}). This
implies $u_\ell=0$; using \SGARefLink{sga7.sga7i.eix.n11-6-2}{11.6.2}(a) below then gives
$T_\ell(A_K)^t=0$.

\medskip
\noindent
\SGARefTarget{sga7.sga7i.eix.n9-5-8}{9.5.8}. When $\ell\neq p$, the Barsotti--Tate groups under consideration
over $S$ and over $K$, respectively, may simply be interpreted in terms of
twisted constant $\ell$-adic sheaves over $S$, or equivalently in terms of
representations of
$\pi_o=\operatorname{Gal}(\overline k/k)$ and
$\pi=\operatorname{Gal}(\overline K/K)$, respectively, on free
$\mathbb Z_\ell$-modules of finite type. One then sees that the homomorphism
(\SGARefLink{sga7.sga7i.eix.n9-5-5}{9.5.5}) is the canonical homomorphism expressing the action of the inertia
subgroup $I\subset\pi$ on $X_\eta(\overline K)$, as defined by the method
of \SGARefLink{sga7.sga7i.eix.n9-2}{9.2}.
\footnote{The printed source refers here to ``the homomorphism (\SGARefLink{sga7.sga7i.eix.n9-5-6}{9.5.6})'';
the homomorphism constructed in this section is (\SGARefLink{sga7.sga7i.eix.n9-5-5}{9.5.5}), whereas \SGARefLink{sga7.sga7i.eix.n9-5-6}{9.5.6} is
the proposition.}
To verify this, one reduces immediately to the situation of \SGARefLink{sga7.sga7i.eix.n9-4}{9.4}, taking
$n$ there to be a power of $\ell$ and $R$ to be a locally constant
$\Lambda$-module (as are $Q$ and $P$); the same is then true of $E,F$, and
$X_\eta$. Verification that the two definitions agree is left to the
reader.

% END INCLUDED SGA7I_English_current_buildable_source_20260823/components/161_expose_IX_section_9_5.tex

% BEGIN INCLUDED SGA7I_English_current_buildable_source_20260823/components/162_expose_IX_sections_9_6_and_9_7.tex
% Expose IX, sections 9.6--9.7.
% French transcription lines 2655--2711; authority scan physical pp. 448--450
% (scan indices 447--449; source folios 436--438). The authority pages were
% checked directly in four overlapping 1100-dpi bands. The scan is clear.
% Four evident print slips affecting the English are corrected and disclosed:
% F_eta for E_eta in the quotient extension, (9.5.6) for (9.5.5), S for eta
% as the base of X_eta, and a missing closing parenthesis. A missing sentence
% stop in 9.7 is restored.

\medskip
\noindent
\SGARefTarget{sga7.sga7i.eix.n9-6}{9.6}. We are now in a position to define the announced canonical pairing
(\SGARefLink{sga7.sga7i.eix.n9-1-2}{9.1.2}), without any restriction on $\ell$, in such a way that, for
$\ell\neq p$, this definition agrees with the one given in \SGARefLink{sga7.sga7i.eix.n9-2}{9.2}. For this
purpose, we apply the considerations of \SGARefLink{sga7.sga7i.eix.n9-5}{9.5} to the following situation:

\begin{equation}\tag{\SGARefTarget{sga7.sga7i.eix.n9-6-1}{9.6.1}}
\left\{
\begin{aligned}
P&=T_\ell(A^o)^t=\check{\underline M}{}'_\ell(1),
& Q&=\underline M_\ell,
& R&=T_\ell(A^o)^f/T_\ell(A^o)^t,\\[1.25ex]
F&=T_\ell(A^o)^f,
& E&=D(F')=D\bigl(T_\ell({A'}^o)^f\bigr).
\end{aligned}
\right.
\end{equation}
In the last formula, $D$ denotes Cartier duality in the sense of
pro-$\ell$ Barsotti--Tate groups; that is, $D(Y)$ is the
\underline{projective} system of the $D(Y(n))$, with the usual transition
morphisms. The structure of $F$ as an extension of $R$ by $P$ is clear.
Likewise,

\begin{equation*}
F'=T_\ell({A'}^o)^f
\end{equation*}
has the structure of an extension of $R'$ by
$P'=T_\ell({A'}^o)^t=\check{\underline M}_\ell(1)$. Applying Cartier
duality, we conclude that $E=D(F')$ is endowed with the structure of an
extension of $D(P')=\underline M_\ell$ by $D(R')$. Now $D(R')$ is
canonically isomorphic to $R$ by \SGARefLink{sga7.sga7i.eix.n5-5}{5.5} (proved in \SGARefLink{sga7.sga7i.eix.n7-4-6}{7.4.6}); this specifies the
structure of $E$ as an extension of $Q$ by $R$.

As the mixed extension $X_\eta$ of $E_\eta$ by $F_\eta$, we take

\begin{equation}\tag{\SGARefTarget{sga7.sga7i.eix.n9-6-2}{9.6.2}}
X_\eta=T_\ell(A_K)=X^0_\eta
 \supset X^1_\eta=T_\ell(A_K)^f
 \supset X^2_\eta=T_\ell(A_K)^t
 \supset0.
\end{equation}
The isomorphism of extensions of $R_\eta$ by $P_\eta$

\begin{equation*}
X^1_\eta\simeq F_\eta
\end{equation*}
is clear. The isomorphism of extensions of $Q_\eta$ by $R_\eta$

\begin{equation*}
X_\eta/X^2_\eta\simeq E_\eta=D(F'_\eta)
\end{equation*}
is clear from the orthogonality theorem \SGARefLink{sga7.sga7i.eix.n5-2}{5.2}.
\footnote{The printed formula has $F_\eta=D(F'_\eta)$. Since
$E=D(F')$ and the displayed quotient is the extension of $Q_\eta$ by
$R_\eta$, the required term is $E_\eta$.}

We therefore have the data required to define a pairing (\SGARefLink{sga7.sga7i.eix.n9-5-5}{9.5.5}), which
here, since $Q^*=\underline M'_\ell$, takes the form (\SGARefLink{sga7.sga7i.eix.n9-1-2}{9.1.2}).
\footnote{The printed source calls this a pairing ``(\SGARefLink{sga7.sga7i.eix.n9-5-6}{9.5.6})''; the pairing
is the homomorphism (\SGARefLink{sga7.sga7i.eix.n9-5-5}{9.5.5}), whereas \SGARefLink{sga7.sga7i.eix.n9-5-6}{9.5.6} is the subsequent proposition.}
Its compatibility with the definition of \SGARefLink{sga7.sga7i.eix.n9-2}{9.2} follows from \SGARefLink{sga7.sga7i.eix.n9-5-8}{9.5.8}. The
considerations of \SGARefLink{sga7.sga7i.eix.n5-8}{5.8} also show, subject to the validity of Tate's theorem
[33] (hence at least when $\operatorname{char}K=0$), that the pairing in
question can be described in terms of the single pro-$\ell$
Barsotti--Tate group $X_\eta=T_\ell(A_K)$ over $\eta$: its three-step
filtration, and the extensions $F$ and $E$ extending $X^1_\eta$ and
$X_\eta/X^2_\eta$, can be defined canonically in terms of $X_\eta$ alone.
\footnote{The printed sentence says that $X_\eta$ is ``over $S$''; its
notation and definition make it the generic-fibre group over $\eta$. It
also lacks its closing parenthesis.}

\medskip
\noindent
\underline{Remark \SGARefTarget{sga7.sga7i.eix.n9-7}{9.7}.} We no longer necessarily assume that $A_K$ has
semistable reduction over the henselian trait $S$. Let $K'$ be a finite
Galois extension of $K$ such that $A_{K'}$ has semistable reduction over
the normalization $S'$ of $S$ in $K'$ (\SGARefLink{sga7.sga7i.eix.n3-6}{3.6}). With the notation of (\SGARefLink{sga7.sga7i.eix.n8-1-8}{8.1.8}),
Galois descent from the monodromy pairing for $A_{K'}$ then defines a
pairing

\begin{equation}\tag{*}
v_\ell^{K'}:\underline M_K\otimes\underline M'_K
 \longrightarrow\mathbb Z_{\ell K}.
\end{equation}
A priori this pairing depends on the choice of $K'$, but it follows from
(10.4.3) below that the pairing

\begin{equation}\tag{\SGARefTarget{sga7.sga7i.eix.n9-7-1}{9.7.1}}
u_\ell:\underline M_K\otimes\underline M'_K
 \longrightarrow\mathbb Q_{\ell K}
\end{equation}
given by

\begin{equation*}
u_\ell=\frac{1}{e(S'/S)}\,v_\ell^{K'}
\end{equation*}
is independent of it, where

\begin{equation*}
e(S'/S)=\operatorname{length}(V'/\mathfrak mV').
\end{equation*}
We may again call (\SGARefLink{sga7.sga7i.eix.n9-7-1}{9.7.1}) the \underline{monodromy pairing} for the
abelian scheme $A_K$ over $K$ and the prime number $\ell$ under
consideration: when $A_K$ itself has semistable reduction, it plainly
coincides with the restriction to $\eta$ of the monodromy pairing (\SGARefLink{sga7.sga7i.eix.n9-1-2}{9.1.2}).
It follows from \SGARefLink{sga7.sga7i.eix.n10-4}{10.4} below that this pairing in fact always has values in
$\mathbb Q_K$ and is independent of $\ell$.

One may ask whether it has values in $\mathbb Z_K$, which is equivalent to
saying that (\SGARefLink{sga7.sga7i.eix.n9-7-1}{9.7.1}) has values in $\mathbb Z_{\ell K}$ for every prime
number $\ell$. This is so at least when $A_K$ is an elliptic curve. The
only case requiring verification is then the one in which
$\underline M_K\simeq\underline M'_K$ has rank $1$, so that
$\underline M_K\otimes\underline M'_K$ is canonically isomorphic to
$\mathbb Z_K$ and the monodromy pairing (\SGARefLink{sga7.sga7i.eix.n9-7-1}{9.7.1}) can simply be interpreted
as an ordinary rational number. It follows from the results of Néron [19]
that this number is $v(j)$, where $j$ is the absolute invariant of $A_K$
and $v$ is the normalized valuation of $K$.


% END INCLUDED SGA7I_English_current_buildable_source_20260823/components/162_expose_IX_sections_9_6_and_9_7.tex

% BEGIN INCLUDED SGA7I_English_current_buildable_source_20260823/components/163_expose_IX_sections_10_0_through_10_2.tex
% Expose IX, section 10 opening and sections 10.1--10.2.
% French transcription lines 2717--2786; authority scan physical pp. 451--453
% (scan indices 450--452; source folios 439--441). Physical p. 454 was
% checked for the exact 10.3 boundary. All pages were compared in four
% overlapping 1100-dpi bands; the print is clear at that resolution.

\medskip
\noindent
10. \underline{Properties of the monodromy pairing}. \underline{Integrality
and positivity theorem}

In this section we shall first give certain essentially trivial formal
properties of the monodromy pairings, leaving the (somewhat tedious)
details of the verifications to the reader. We shall also state a deeper
theorem (\SGARefLink{sga7.sga7i.eix.n10-4}{10.4}) concerning these pairings and indicate how its proof may be
reduced to a special case of Jacobians, to be treated (by transcendental
methods) in section 12, which is in fact logically independent of section
11.

\medskip
\noindent
\SGARefTarget{sga7.sga7i.eix.n10-1}{10.1}. \underline{Transport of structure}. It is clear from the definition
of the monodromy pairing (\SGARefLink{sga7.sga7i.eix.n9-1-2}{9.1.2}) that this pairing is compatible with
transport of structure in the situation $(S,A_K,A'_K)$.

\medskip
\noindent
\SGARefTarget{sga7.sga7i.eix.n10-2}{10.2}. \underline{Symmetry}. Starting with the dual scheme $A'_K$ of $A_K$
and applying the construction of section 9 to it, we must introduce the
dual scheme $(A'_K)'$ of $A'_K$, which is canonically isomorphic to $A_K$
by the biduality theorem ([5] or VIII 3.2). Denote by
$u_\ell^{A_K}$ the monodromy pairing (\SGARefLink{sga7.sga7i.eix.n9-1-2}{9.1.2}) relative to $A_K$. In the same
way we obtain a pairing

\begin{equation}\tag{\SGARefTarget{sga7.sga7i.eix.n10-2-1}{10.2.1}}
u_\ell^{A'_K}:\underline M'_\ell\otimes\underline M_\ell
 \longrightarrow\mathbb Z_\ell.
\end{equation}
I claim that the latter is obtained from $u_\ell^{A_K}$ by the symmetry

\begin{equation*}
s:\underline M'_\ell\otimes\underline M_\ell
 \xrightarrow{\ \sim\ }
 \underline M_\ell\otimes\underline M'_\ell;
\end{equation*}
that is,

\begin{equation}\tag{\SGARefTarget{sga7.sga7i.eix.n10-2-2}{10.2.2}}
u_\ell^{A'_K}=u_\ell^{A_K}\circ s.
\end{equation}

Taking account of the known anticommutativity property of the pairings
$\varphi_\xi$ ([5] or VIII 2.2.11), the symmetry formula (\SGARefLink{sga7.sga7i.eix.n10-2-2}{10.2.2}) may be
reformulated as follows in terms of a fixed polarization $\xi$ of $A_K$.
It defines a homomorphism

\begin{equation*}
\widetilde\xi:A_K\longrightarrow A'_K
\end{equation*}
and hence a corresponding homomorphism ($t$ standing for ``toric'')

\begin{equation}\tag{\SGARefTarget{sga7.sga7i.eix.n10-2-3}{10.2.3}}
\xi_t:\underline M\longrightarrow\underline M'.
\end{equation}
We may therefore form the pairing

\begin{equation}\tag{\SGARefTarget{sga7.sga7i.eix.n10-2-4}{10.2.4}}
\begin{aligned}
\varphi^t_{\xi,\ell}:\underline M_\ell\otimes\underline M_\ell
 &\longrightarrow\mathbb Z_\ell,\\
\varphi^t_{\xi,\ell}(x,y)&=u_\ell^{A_K}\bigl(x,\xi_t(y)\bigr).
\end{aligned}
\end{equation}
The preceding form $\varphi^t_{\xi,\ell}$ is \underline{symmetric}:

\begin{equation}\tag{\SGARefTarget{sga7.sga7i.eix.n10-2-5}{10.2.5}}
\varphi^t_{\xi,\ell}(x,y)=\varphi^t_{\xi,\ell}(y,x).
\end{equation}
In fact, this assertion is valid for every element $\xi$ of the
Néron--Severi group of $A_K$, not necessarily a polarization. But only
when $\xi$ is nondegenerate---that is, defines an \underline{isogeny}
$A_K\longrightarrow A'_K$, and hence also an isogeny (\SGARefLink{sga7.sga7i.eix.n10-2-3}{10.2.3})---is the
symmetry relation (\SGARefLink{sga7.sga7i.eix.n10-2-5}{10.2.5}) equivalent to (\SGARefLink{sga7.sga7i.eix.n10-2-2}{10.2.2}). When $\xi$ is a
polarization, we shall moreover sharpen (\SGARefLink{sga7.sga7i.eix.n10-2-5}{10.2.5}) considerably in \SGARefLink{sga7.sga7i.eix.n10-4}{10.4}(b)
below.

Let us finally mention a third interpretation of the symmetry relation.
Consider a biextension (that is, a divisorial correspondence) $\xi$ on a
pair $(A_K,\overline A_K)$ of abelian schemes with semistable reduction
over $K$. It defines a homomorphism

\begin{equation*}
A_K\longrightarrow\overline A'_K.
\end{equation*}
Passing to the Néron models, this defines a homomorphism on the maximal
tori of the special fibres, and therefore, in the reverse direction, a
homomorphism on their character groups:

\begin{equation}\tag{\SGARefTarget{sga7.sga7i.eix.n10-2-6}{10.2.6}}
\xi_t:\overline{\underline M}\longrightarrow\underline M',
\end{equation}
where $\overline{\underline M}$ is defined from $\overline A_K$ as
$\underline M$ is defined from $A_K$. We can therefore form the pairing

\begin{equation}\tag{\SGARefTarget{sga7.sga7i.eix.n10-2-7}{10.2.7}}
\begin{aligned}
\varphi^t_\xi:\underline M_\ell\otimes
 \overline{\underline M}_\ell&\longrightarrow\mathbb Z_\ell,\\
\varphi^t_\xi(x,y)&=u_\ell^{A_K}\bigl(x,\xi_t(y)\bigr).
\end{aligned}
\end{equation}
\footnote{In the printed right-hand side the underlying pairing is denoted
only by $\varphi^t$, without a definition. It is the monodromy pairing
$u_\ell^{A_K}$ used in the identical construction (\SGARefLink{sga7.sga7i.eix.n10-2-4}{10.2.4}).}
Notice in passing that, for a variable abelian scheme $A_K$ with semistable
reduction, the twisted constant group $\underline M$ is a covariant
functor of $A_K$. Formation of the pairings (\SGARefLink{sga7.sga7i.eix.n10-2-7}{10.2.7}) associated with
biextensions of abelian schemes with semistable reduction over $(K,S)$ is
then compatible, in the evident sense, with taking inverse images of
biextensions.

When $\xi$ is the canonical (``Weil'') biextension of $(A_K,A'_K)$, the
pairing (\SGARefLink{sga7.sga7i.eix.n10-2-7}{10.2.7}) is of course precisely the monodromy pairing (\SGARefLink{sga7.sga7i.eix.n9-1-2}{9.1.2}).
\footnote{The printed source refers here to (\SGARefLink{sga7.sga7i.eix.n9-2-1}{9.2.1}), which is the inertia
orthogonality formula $gx-x\in V^\perp$, not the monodromy pairing. The
pairing is numbered (\SGARefLink{sga7.sga7i.eix.n9-1-2}{9.1.2}).}
Now consider the transpose ${}^s\xi$ of $\xi$, a biextension of
$(\overline A_K,A_K)$, and hence the pairing

\begin{equation*}
\varphi^t_{{}^s\xi}:\overline{\underline M}_\ell\otimes
 \underline M_\ell\longrightarrow\mathbb Z_\ell.
\end{equation*}
The symmetry relation is

\begin{equation}\tag{\SGARefTarget{sga7.sga7i.eix.n10-2-8}{10.2.8}}
\varphi^t_{{}^s\xi}=\varphi^t_\xi\circ s,
\qquad
s:\overline{\underline M}_\ell\otimes\underline M_\ell
 \xrightarrow{\ \sim\ }
 \underline M_\ell\otimes\overline{\underline M}_\ell,
\end{equation}
where $s$ is again the symmetry isomorphism. Compare this relation with
the antisymmetry relation of VIII \SGARefLink{sga7.sga7i.eviii.n2-2-11}{2.2.11}.


% END INCLUDED SGA7I_English_current_buildable_source_20260823/components/163_expose_IX_sections_10_0_through_10_2.tex

% BEGIN INCLUDED SGA7I_English_current_buildable_source_20260823/components/164_expose_IX_sections_10_3_and_10_4.tex
% Expose IX, sections 10.3--10.4.
% French transcription lines 2792--2883; authority scan physical pp. 454--456
% (scan indices 453--455; source folios 442--444). Physical p. 457 was
% checked for the exact 10.5 boundary. All pages were compared in four
% overlapping 1100-dpi bands; the text, formulas, and simple square are
% clear at that resolution.

\medskip
\noindent
\SGARefTarget{sga7.sga7i.eix.n10-3}{10.3}. \underline{Change of trait}. Let

\begin{equation}\tag{\SGARefTarget{sga7.sga7i.eix.n10-3-1}{10.3.1}}
g:S'\longrightarrow S
\end{equation}
be a dominant morphism of henselian traits, associated with an injective
local homomorphism of discrete valuation rings

\begin{equation*}
V\longrightarrow V'.
\end{equation*}
As usual, put

\begin{equation}\tag{\SGARefTarget{sga7.sga7i.eix.n10-3-2}{10.3.2}}
e(S'/S)=\operatorname{length}_V\bigl(V'/\underline mV'\bigr),
\end{equation}
where $\underline m$ is the maximal ideal of $V$. It is then known that
the homomorphism induced by (\SGARefLink{sga7.sga7i.eix.n10-3-1}{10.3.1}) on the ``tame'' parts of the inertia
groups,

\begin{equation}\tag{\SGARefTarget{sga7.sga7i.eix.n10-3-3}{10.3.3}}
g_*:I'_t\simeq\prod_{\ell\ne p}\mathbb Z_\ell(1)(k')
 \longrightarrow
 I_t=\prod_{\ell\ne p}\mathbb Z_\ell(1)(k),
\end{equation}
is given by

\begin{equation}\tag{\SGARefTarget{sga7.sga7i.eix.n10-3-4}{10.3.4}}
g_*(x)=e(S'/S)x,
\end{equation}
after the two sides of (\SGARefLink{sga7.sga7i.eix.n10-3-3}{10.3.3}) have been identified by means of the
evident isomorphisms
$\mathbb Z_\ell(1)(k)\xrightarrow{\ \sim\ }\mathbb Z_\ell(1)(k')$.
\footnote{The printed source leaves the internal reference following
``Expos\'e I'' blank.}

On the other hand, if $A_K$ is an abelian scheme over $K$ with
semistable reduction over $S$, formation of the twisted lattices
$\underline M$ and $\underline M'$ commutes with every base change
$S'\longrightarrow S$ as above, as follows immediately from the
definitions (especially \SGARefLink{sga7.sga7i.eix.n3-4}{3.4}). We may therefore regard the pairing
$u_\ell^{A_{K'}}$ defined by $A_{K'}$, where $K'$ is the fraction field
of $V'$, as a pairing

\begin{equation*}
u_\ell^{A_{K'}}:
 \underline M_{\ell S'}\otimes\underline M'_{\ell S'}
 \longrightarrow\mathbb Z_{\ell S'},
\end{equation*}
and compare it with the pairing $u_{\ell S'}^{A_K}$ obtained from
$u_\ell^{A_K}$ by base change. When $\ell\ne p$, the construction of the
pairings $u_\ell$ given in \SGARefLink{sga7.sga7i.eix.n9-2}{9.2}, together with formula (\SGARefLink{sga7.sga7i.eix.n10-3-4}{10.3.4}), shows at
once that\footnote{The printed source cites (\SGARefLink{sga7.sga7i.eix.n10-3-3}{10.3.3}), the displayed map
$g_*$; the multiplication formula it uses is (\SGARefLink{sga7.sga7i.eix.n10-3-4}{10.3.4}).}

\begin{equation}\tag{\SGARefTarget{sga7.sga7i.eix.n10-3-5}{10.3.5}}
u_\ell^{A_{K'}}=e(S'/S)u_\ell^{A_K}.
\end{equation}
In fact this formula remains valid for $\ell=p$. To see this, return to
the general construction of \SGARefLink{sga7.sga7i.eix.n9-3}{9.3}--\SGARefLink{sga7.sga7i.eix.n9-6}{9.6} and reduce to the analogous
compatibility assertion for the isomorphism \SGARefLink{sga7.sga7i.eix.n9-4-4}{9.4.4}(e), under a change of
strictly local traits $S'\longrightarrow S$. Notice that the isomorphism
in question is induced by the normalized valuation
$K^*/V^*\simeq\mathbb Z$, and therefore gives the commutative diagram
\footnote{The printed source says that $K^*/V^*\simeq\mathbb Z$ is
deduced from the ``norm.''  The displayed map is the normalized discrete
valuation.}

\begin{center}
\begin{tikzcd}
K^*/V^* \arrow[r, "\sim"] \arrow[d]
  & \mathbb Z \arrow[d, "e(S'/S)"] \\
K^{\prime *}/V^{\prime *} \arrow[r, "\sim"]
  & \mathbb Z .
\end{tikzcd}
\end{center}

Here is a deeper result concerning the monodromy pairings. We shall prove
it by transcendental methods in section 12, relying on results that will
be proved later in the Seminar. (We shall not, in fact, use \SGARefLink{sga7.sga7i.eix.n10-4}{10.4} in the
remainder of the Seminar, any more than we shall use the other results of
the present expos\'e.)

\medskip
\noindent
\underline{Theorem \SGARefTarget{sga7.sga7i.eix.n10-4}{10.4}.} \underline{Let $S$ be a henselian trait, and
let $A_K$ be an abelian scheme over $K$ with semistable reduction over
$S$. Consider the lattices $\underline M$ and $\underline M'$ associated
as in \SGARefLink{sga7.sga7i.eix.n9-1}{9.1} (cf. \SGARefLink{sga7.sga7i.eix.n5-5-2}{5.5.2}). Then the following assertions hold:}

\smallskip
\noindent
a) \underline{There exists a unique pairing}

\begin{equation}\tag{\SGARefTarget{sga7.sga7i.eix.n10-4-1}{10.4.1}}
u:\underline M\otimes\underline M'\longrightarrow\mathbb Z_S
\end{equation}
\underline{such that, for every prime number $\ell$,
$u\otimes_{\mathbb Z}\mathbb Z_\ell$ is the monodromy pairing $u_\ell$
of (\SGARefLink{sga7.sga7i.eix.n9-1-2}{9.1.2}). This pairing is separating.}

\smallskip
\noindent
b) \underline{Let $\xi$ be a polarization of $A_K$, thus defining a
homomorphism (\SGARefLink{sga7.sga7i.eix.n10-2-3}{10.2.3})}

\begin{equation*}
\xi_t:\underline M\longrightarrow\underline M',
\end{equation*}
\underline{and consider the form}

\begin{equation}\tag{\SGARefTarget{sga7.sga7i.eix.n10-4-2}{10.4.2}}
\begin{aligned}
\varphi_\xi^t:\underline M\otimes\underline M
  &\longrightarrow\mathbb Z_S,\\
\varphi_\xi^t(x,y)&=u\bigl(x,\xi_t(y)\bigr).
\end{aligned}
\end{equation}
\underline{This form is symmetric, positive, and nondegenerate.}

Of course, in terms of the geometric fibre $M_o$ of $\underline M$ at a
given geometric point of $S$, the preceding assertion says that the
corresponding form

\begin{equation*}
M_o\otimes M_o\longrightarrow\mathbb Z
\end{equation*}
is symmetric, positive, and nondegenerate. The pairing (\SGARefLink{sga7.sga7i.eix.n10-4-1}{10.4.1}) will
again be called the \underline{monodromy pairing}.

% END INCLUDED SGA7I_English_current_buildable_source_20260823/components/164_expose_IX_sections_10_3_and_10_4.tex

% BEGIN INCLUDED SGA7I_English_current_buildable_source_20260823/components/165_expose_IX_section_10_5.tex
% Expose IX, section 10.5.
% French transcription lines 2889--2955; authority scan physical pp. 457--460
% (scan indices 456--459; source folios 445--448). Physical p. 460 was
% checked through the exact section 11 boundary. All pages were compared in
% four overlapping 1100-dpi bands; the prose and formulas are clear at that
% resolution, so no higher-resolution escalation was necessary.

\medskip
\noindent
\SGARefTarget{sga7.sga7i.eix.n10-5}{10.5}. \underline{Reduction of \SGARefLink{sga7.sga7i.eix.n10-4}{10.4} to a particular case}

\smallskip
\noindent
\SGARefTarget{sga7.sga7i.eix.n10-5-1}{10.5.1}. The uniqueness of a $u$ satisfying the stated condition (even
for a single $\ell$) is clear. On the other hand, existence simply means
that every $u_\ell$ in fact maps $\underline M\otimes\underline M'$ into
$\mathbb Z_S$, and that the homomorphism thereby obtained

\begin{equation*}
\underline M\otimes\underline M'\longrightarrow\mathbb Z_S
\end{equation*}
is independent of $\ell$. In fact it even suffices to prove the analogous
assertion for the $u_\ell\otimes_{\mathbb Z_\ell}\mathbb Q_\ell$ and for
a homomorphism

\begin{equation*}
u_{\mathbb Q}:\underline M_{\mathbb Q}\otimes
 \underline M'_{\mathbb Q}\longrightarrow\mathbb Q_S.
\end{equation*}

Suppose, then, that existence has been established after a base change
$S'\longrightarrow S$ as in \SGARefLink{sga7.sga7i.eix.n9-3}{9.3}. It follows immediately from (\SGARefLink{sga7.sga7i.eix.n10-3-5}{10.3.5})
that there exists a homomorphism

\begin{equation*}
u:\underline M\otimes\underline M'
 \longrightarrow\frac{1}{e(S'/S)}\mathbb Z_S\subset\mathbb Q_S
\end{equation*}
compatible with the $u_\ell\otimes_{\mathbb Z_\ell}\mathbb Q_\ell$;
this proves a). As for b), it is then plain that, if it holds for
$A_{K'}$, it also holds for $A_K$.

Consequently, in proving \SGARefLink{sga7.sga7i.eix.n10-4}{10.4} we may make any dominant change of traits
$S'\longrightarrow S$. In particular, this allows us to suppose that
$S$ is complete with algebraically closed residue field. In that case
$\underline M$ and $\underline M'$ are constant group schemes, with
values $M$ and $M'$, finitely generated free $\mathbb Z$-modules, and a)
and b) may be reformulated more conveniently in terms of the latter.

\smallskip
\noindent
\SGARefTarget{sga7.sga7i.eix.n10-5-2}{10.5.2}. Suppose now that $A_K$ is isomorphic to a direct factor of an
abelian scheme $\overline A_K$ over $K$ with semistable reduction, and
that \SGARefLink{sga7.sga7i.eix.n10-4}{10.4} has been proved for the latter. It follows at once that \SGARefLink{sga7.sga7i.eix.n10-4}{10.4}
also holds for $A_K$, in view of the evident functorial behaviour of the
lattices $\underline M,\underline M'$ and of the monodromy pairings with
respect to direct sums of abelian schemes with semistable reduction.
\footnote{The printed source says that $\overline A_K$ is ``over $k$'';
the notation and the direct-factor argument require ``over $K$.''}

It is equally immediate from \SGARefLink{sga7.sga7i.eix.n10-5-1}{10.5.1} that the validity of \SGARefLink{sga7.sga7i.eix.n10-4}{10.4} depends
only on the isogeny class of $A_K$. Together with the preceding remark
and the ``Poincar\'e reducibility theorem'' [15], this shows that, if
$A_K$ is a quotient of an abelian scheme $\overline A_K$ for which \SGARefLink{sga7.sga7i.eix.n10-4}{10.4}
holds, then \SGARefLink{sga7.sga7i.eix.n10-4}{10.4} also holds for $A_K$.
\footnote{The printed source leaves the pinpoint following citation [15]
blank.}

\smallskip
\noindent
\SGARefTarget{sga7.sga7i.eix.n10-5-3}{10.5.3}. We now use the well-known fact [31, Th.~11 (p.~20)] that, after
making a finite extension of $K$ if necessary (which is permissible by
\SGARefLink{sga7.sga7i.eix.n10-5-1}{10.5.1}), every abelian scheme $A_K$ is a quotient of the Jacobian

\begin{equation*}
\underline{\operatorname{Pic}}^0_{X_K/K}
\end{equation*}
of a proper, smooth, geometrically connected curve $X_K$ over $K$.
Thus, by \SGARefLink{sga7.sga7i.eix.n10-5-2}{10.5.2}, we are reduced to proving \SGARefLink{sga7.sga7i.eix.n10-4}{10.4} for such a Jacobian.
\footnote{The printed source says that the curve $X_K$ is ``over $X$'';
the displayed relative Picard scheme and the argument require ``over
$K$.''}

\smallskip
\noindent
\SGARefTarget{sga7.sga7i.eix.n10-5-4}{10.5.4}. By a theorem of Artin--Mumford [3] (already cited and used in
Expos\'e V), after making a further finite extension $K'$ of $K$ if
necessary, we may suppose that $X_K$ extends to a projective and flat
scheme over $S$ which is regular and whose geometric special fibre
$X_{\bar s}$ has no singularities other than ordinary quadratic
singularities (that is, here, normal crossings of two branches).
\footnote{The printed source leaves the internal reference following
``Expos\'e V'' blank.}
In this case we shall prove \SGARefLink{sga7.sga7i.eix.n10-4}{10.4} a) in \SGARefLink{sga7.sga7i.eix.n12-5}{12.5} below, by describing the
monodromy pairing (\SGARefLink{sga7.sga7i.eix.n10-4-1}{10.4.1}) completely in terms of the geometric data of
the situation; at the same time we shall verify \SGARefLink{sga7.sga7i.eix.n10-4}{10.4} b) for the canonical
polarization of the Jacobian.

\smallskip
\noindent
\SGARefTarget{sga7.sga7i.eix.n10-5-5}{10.5.5}. It remains to see that, if for a given $A_K$ with semistable
reduction assertion \SGARefLink{sga7.sga7i.eix.n10-4}{10.4} a) holds and assertion \SGARefLink{sga7.sga7i.eix.n10-4}{10.4} b) holds for one
given polarization $\xi_o$ of $A_K$, then the same conclusion holds for
every polarization $\xi$ of $A_K$. Indeed, let

\begin{equation*}
\widetilde\xi_o,\widetilde\xi:A_K\longrightarrow A'_K
\end{equation*}
be the homomorphisms defined by $\xi_o$ and $\xi$; these are isogenies.
Then

\begin{equation*}
\widetilde\xi=\widetilde\xi_o\alpha
 ={}^t(\alpha^*)\widetilde\xi_o,
 \qquad
\alpha\in\operatorname{End}(A_K)\otimes_{\mathbb Z}\mathbb Q,
\end{equation*}
where $\alpha\longmapsto\alpha^*$ denotes the canonical involution of
the algebra $\operatorname{End}(A_K)\otimes_{\mathbb Z}\mathbb Q$
defined by the polarization $\xi_o$, and the superscript $t$ in
${}^t(\alpha^*)$ denotes transpose (associating with an endomorphism of
$A_K$ an endomorphism of $A'_K$). On the other hand, it is known
[15, Chap.~V, Th.~3] that

\begin{equation*}
\alpha=\sum_i\beta_i^*\beta_i,
 \qquad
\beta_i\in\operatorname{End}(A_K)\otimes_{\mathbb Z}\mathbb Q,
\end{equation*}
and hence

\begin{equation*}
\widetilde\xi=\sum_i{}^t\beta_i\,\widetilde\xi_o\,\beta_i.
\end{equation*}
Equivalently, in terms of the biextensions (that is, divisorial
correspondences) associated with the polarizations under consideration
(VIII \SGARefLink{sga7.sga7i.eviii.n4-14}{4.14}),

\begin{equation}\tag{\SGARefTarget{sga7.sga7i.eix.n10-5-5-1}{10.5.5.1}}
\delta(\xi)=\sum_i\delta(\xi_o)\circ(\beta_i,\beta_i).
\end{equation}
By what was said after (\SGARefLink{sga7.sga7i.eix.n10-2-7}{10.2.7}), this formula implies (with the notation
of (\SGARefLink{sga7.sga7i.eix.n10-4-2}{10.4.2}))

\begin{equation}\tag{\SGARefTarget{sga7.sga7i.eix.n10-5-5-2}{10.5.5.2}}
\varphi^t_\xi
 =\sum_i\varphi^t_{\xi_o}(\beta_i,\beta_i)
 \in\operatorname{Hom}
 \bigl(M_{\mathbb Q}\otimes_{\mathbb Q}M_{\mathbb Q},\mathbb Q\bigr).
\end{equation}
It indeed follows that, if $\varphi^t_{\xi_o}$ is symmetric and positive,
then so is $\varphi^t_\xi$.

Finally, the nondegeneracy of $\varphi^t_{\xi_o}$ is plainly equivalent
to the assertion that the pairing (\SGARefLink{sga7.sga7i.eix.n10-4-1}{10.4.1}) is separating (since
$\xi_{ot}:M\otimes_{\mathbb Z}\mathbb Q
\longrightarrow M'\otimes_{\mathbb Z}\mathbb Q$ is an isomorphism,
$\widetilde\xi_o$ being an isogeny), and is therefore likewise
independent of the chosen polarization. One may also observe that, in
order to verify that (\SGARefLink{sga7.sga7i.eix.n10-4-1}{10.4.1}) is separating, it suffices to verify the
same assertion for the monodromy pairing $u_\ell$ of (\SGARefLink{sga7.sga7i.eix.n9-1-2}{9.1.2}) for one
prime number $\ell$. Taking $\ell\ne p$, one may use the description of
$u_\ell$ in \SGARefLink{sga7.sga7i.eix.n9-2}{9.2}; the fact that $u_\ell$ is separating on the left (and
hence on both sides, for reasons of rank) is then immediate.

% END INCLUDED SGA7I_English_current_buildable_source_20260823/components/165_expose_IX_section_10_5.tex

% BEGIN INCLUDED SGA7I_English_current_buildable_source_20260823/components/166_expose_IX_sections_11_0_through_11_2.tex
% Expose IX, sections 11.0--11.2.
% French transcription lines 2956--3027; authority scan physical pp. 460--463
% (scan indices 459--462; source folios 448--451). Physical p. 463 was
% checked through the exact 11.3 boundary. All pages were compared in four
% overlapping 1100-dpi bands; the text and formulas are clear at that
% resolution, so no higher-resolution escalation was necessary.

\medskip
\noindent
11. \underline{Connected components of the N\'eron model: duality
relation and asymptotic behaviour}

\smallskip
\noindent
\SGARefTarget{sga7.sga7i.eix.n11-0}{11.0}. The purpose of the present section is to study the finite \'etale
$k$-group

\begin{equation}\tag{\SGARefTarget{sga7.sga7i.eix.n11-0-1}{11.0.1}}
\Phi_o=A_o/A_o^{,o}
\end{equation}
already considered in section 1. For every prime number $\ell$, denote
by

\begin{equation*}
\Phi_o(\ell)
\end{equation*}
its $\ell$-primary component. Thus knowledge of $\Phi_o$ is equivalent
to knowledge of the $\Phi_o(\ell)$, by the formula

\begin{equation}\tag{\SGARefTarget{sga7.sga7i.eix.n11-0-2}{11.0.2}}
\Phi_o=\coprod_\ell\Phi_o(\ell).
\end{equation}

Notice that the groups $\Phi_o$ and $\Phi_o(\ell)$ are determined,
respectively, by the ordinary finite groups $\Phi_o(\bar k)$ and
$\Phi_o(\ell)(\bar k)$ together with the natural action of

\begin{equation*}
\pi_o=\operatorname{Gal}(\bar k/k)
\end{equation*}
on them. On the other hand, it is clear that $\Phi_o(\bar k)$ is also
the group of connected components of the N\'eron model of
$A_{\widetilde K}$, where $\widetilde K$ is the fraction field of the
strict henselization of $S$ relative to the chosen separable closure
$\bar k$ of $k$. Thus, in studying $\Phi_o$, we may reduce to the case
where $S$ is strictly local; the actions of

\begin{equation*}
\pi_o\simeq\operatorname{Gal}(\widetilde K/K)
\end{equation*}
on $\Phi_o(\bar k)$ are then recovered automatically by transport of
structure. We could likewise reduce further to the case where $S$ is
complete, using the base change $\widehat S\longrightarrow S$ and the
fact that formation of the N\'eron model commutes with it
[21, Th.~\SGARefLink{sga7.sga7i.eix.n3-4}{3.4} b)].

\smallskip
\noindent
\SGARefTarget{sga7.sga7i.eix.n11-1}{11.1}. Suppose, then, that $S$ is strictly local, so that $\Phi_o$ may be
identified with the ordinary finite group $\Phi_o(k)$. We obtain

\begin{equation}\tag{\SGARefTarget{sga7.sga7i.eix.n11-1-1}{11.1.1}}
\Phi_o(k)\simeq A_o(k)/A_o^o(k)
 \simeq A(S)/A^o(S)
 \simeq A(K)/A(K)^o,
\end{equation}
where the first equality follows from the fact that $A_o$ is smooth over
$k$, the second from Hensel's lemma since $A$ is smooth over $S$ and
$S$ is henselian, and the third from the relation
$A(S)\xrightarrow{\sim}A(K)$ and from the convention

\begin{equation}\tag{\SGARefTarget{sga7.sga7i.eix.n11-1-2}{11.1.2}}
A(K)^o\mathrel{\mathop:}=A^o(S).
\end{equation}
Thus $A(K)^o$ appears as a natural subgroup of finite index in $A(K)$
(and is not to be confused with $A^o(K)=A(K)$). The isomorphisms
(\SGARefLink{sga7.sga7i.eix.n11-1-1}{11.1.1}) are plainly functorial in $A$ and compatible with transport of
structure in $(S,A_K)$, and in particular with the action of any Galois
groups.

Now let $\ell$ be a prime number different from $p$. Then $A^o$ is
$\ell$-divisible; more precisely, $\ell\operatorname{id}_{A^o}$ is an
\'etale surjective morphism. It consequently induces a surjective
endomorphism of $A^o(S)=A(K)^o$, which is therefore an
\underline{$\ell$-divisible group}. It follows at once that, for every
integer $n\geq 0$ prime to $p=\operatorname{char}(k)$,

\begin{equation}\tag{\SGARefTarget{sga7.sga7i.eix.n11-1-3}{11.1.3}}
{}_n\Phi_o(k)\simeq{}_nA(K)/{}_nA(K)^o,
 \qquad (n,p)=1.
\end{equation}
By the definition (\SGARefLink{sga7.sga7i.eix.n2-2-3}{2.2.3}) of the fixed part of ${}_nA_K$, we have

\begin{equation}\tag{\SGARefTarget{sga7.sga7i.eix.n11-1-4}{11.1.4}}
{}_nA(K)^o=({}_nA_K)^f(\bar K)
 \simeq T_n(A_K)^f(\bar K)
 \otimes_{\mathbb Z_n}\mathbb Z/n\mathbb Z.
\end{equation}
Taking $n$ to be a power of $\ell$, and putting as in (\SGARefLink{sga7.sga7i.eix.n2-5-3}{2.5.3})

\begin{equation}\tag{\SGARefTarget{sga7.sga7i.eix.n11-1-5}{11.1.5}}
U=T_\ell\bigl(A(\bar K)\bigr),
 \qquad
T_\ell\bigl(A(\bar K)\bigr)^f=U^I,
\end{equation}
where $U^I$ denotes the submodule of $U$ consisting of invariants under
the inertia group $I$, we may write (\SGARefLink{sga7.sga7i.eix.n11-1-3}{11.1.3}) in the form

\begin{equation}\tag{\SGARefTarget{sga7.sga7i.eix.n11-1-6}{11.1.6}}
{}_n\Phi_o(k)\simeq (U_n)^I/(U^I)_n,
 \qquad n=\ell^{r+1},
\end{equation}
where the subscript $n$ denotes reduction modulo $n$, that is, tensor
product over $\mathbb Z_\ell$ with
$\mathbb Z_\ell/n\mathbb Z_\ell$. The transition morphisms between the
groups occurring on the two sides of (\SGARefLink{sga7.sga7i.eix.n11-1-6}{11.1.6}) being evident, passage to
the inductive limit over the powers of $\ell$ gives the following
result.

\medskip
\noindent
\underline{Proposition \SGARefTarget{sga7.sga7i.eix.n11-2}{11.2}.} \underline{Suppose that $S$ is strictly
local, and let $\ell$ be a prime number different from $p$. Then there is
a canonical isomorphism}

\begin{equation}\tag{\SGARefTarget{sga7.sga7i.eix.n11-2-1}{11.2.1}}
\Phi_o(\ell)\simeq
 \frac{(U\otimes D_\ell)^I}{U^I\otimes D_\ell},
\end{equation}
\underline{where $U$ is defined in (\SGARefLink{sga7.sga7i.eix.n11-1-5}{11.1.5}) and
$D_\ell=\mathbb Q_\ell/\mathbb Z_\ell$.}

Thus the group $\Phi_o(\ell)$ can be described solely in terms of the
Tate module $U=T_\ell(A(\bar K))$, regarded as a Galois module under the
action of the inertia group $I$.

% END INCLUDED SGA7I_English_current_buildable_source_20260823/components/166_expose_IX_sections_11_0_through_11_2.tex

% BEGIN INCLUDED SGA7I_English_current_buildable_source_20260823/components/167_expose_IX_section_11_3.tex
% Expose IX, section 11.3.
% French transcription lines 3029--3147; authority scan physical pp. 463--466
% (scan indices 462--465; source folios 451--454). Physical p. 466 was
% checked through the exact 11.4 boundary. All pages were compared in four
% overlapping 1100-dpi bands. The text, formulas, arrow directions, labels,
% and two native diagrams are clear at that resolution.

\medskip
\noindent
\SGARefTarget{sga7.sga7i.eix.n11-3}{11.3}. We shall use the explicit expression (\SGARefLink{sga7.sga7i.eix.n11-2-1}{11.2.1}) to sketch the
definition of a canonical pairing

\begin{equation}\tag{\SGARefTarget{sga7.sga7i.eix.n11-3-1}{11.3.1}}
\Phi_o(\ell)\otimes\Phi'_o(\ell)\longrightarrow\mathbb Q/\mathbb Z
\end{equation}
which is a perfect duality; here $\Phi'_o$ denotes, as in section 1, the
group of connected components associated with the N\'eron model of the
\underline{dual} $A'_K$ of $A_K$. There is little doubt (although the
author has dispensed with the verification) that this pairing is none
other than, possibly up to sign, the one defined in (\SGARefLink{sga7.sga7i.eix.n1-2-1}{1.2.1}), and the
reader is invited, as an exercise, to establish the required
compatibility (with the correct sign!).

Put

\begin{equation*}
U'=T_\ell\bigl(A'(\bar K)\bigr),
 \qquad
T=T_\ell(1)(\bar K)=T_\ell(\bar K^*).
\end{equation*}
There is a \underline{perfect pairing} of finitely generated free
$\mathbb Z_\ell$-modules, compatible with the actions of $I$,

\begin{equation}\tag{\SGARefTarget{sga7.sga7i.eix.n11-3-2}{11.3.2}}
\varphi:U\otimes U'\longrightarrow T,
 \qquad (T\simeq\mathbb Z_\ell\text{ noncanonically}),
\end{equation}
and, by the known structure of $I$, an exact sequence
\footnote{The printed source leaves the internal reference following
the first occurrence of $I$ blank.}

\begin{equation}\tag{\SGARefTarget{sga7.sga7i.eix.n11-3-3}{11.3.3}}
1\longrightarrow R\longrightarrow I\longrightarrow T\longrightarrow1,
\end{equation}
where $R$ is a profinite group \underline{prime to $\ell$}. Starting now
from the data (\SGARefLink{sga7.sga7i.eix.n11-3-2}{11.3.2}) and (\SGARefLink{sga7.sga7i.eix.n11-3-3}{11.3.3}), introduce the groups

\begin{equation}\tag{\SGARefTarget{sga7.sga7i.eix.n11-3-4}{11.3.4}}
\begin{aligned}
\Phi&=\frac{(U\otimes D_\ell)^I}{U^I\otimes D_\ell},\qquad
\Phi'=\frac{(U'\otimes D_\ell)^I}{{U'}^I\otimes D_\ell},\qquad
D_\ell=\mathbb Q_\ell/\mathbb Z_\ell.
\end{aligned}
\end{equation}
We shall define a perfect duality

\begin{equation}\tag{\SGARefTarget{sga7.sga7i.eix.n11-3-5}{11.3.5}}
\Phi\otimes\Phi'\longrightarrow\mathbb Q/\mathbb Z,
\end{equation}
which will in particular give the announced pairing (\SGARefLink{sga7.sga7i.eix.n11-3-1}{11.3.1}). Consider
the exact sequence of $I$-modules

\begin{equation}\tag{\SGARefTarget{sga7.sga7i.eix.n11-3-6}{11.3.6}}
0\longrightarrow U\xrightarrow{\ i\ }
 U\otimes_{\mathbb Z_\ell}\mathbb Q_\ell
 \xrightarrow{\ j\ }U\otimes D_\ell\longrightarrow0
\end{equation}
and the associated exact cohomology sequence, where the cohomology of
$I$ is defined by \underline{continuous cochains} with coefficients in
the topological groups under consideration:

\begin{equation}\tag{\SGARefTarget{sga7.sga7i.eix.n11-3-7}{11.3.7}}
\resizebox{\textwidth}{!}{%
\begin{tikzcd}[ampersand replacement=\&, column sep=small, row sep=large]
0 \arrow[r]
 \& H^o(I,U) \arrow[r, "i^o"] \arrow[d, no head, "\wr"]
 \& H^o(I,U\otimes\mathbb Q_\ell) \arrow[r, "j^o"]
      \arrow[d, no head, "\wr"]
 \& H^o(I,U\otimes D_\ell) \arrow[r, "\partial"]
      \arrow[d, no head, "\wr"]
 \& H^1(I,U) \arrow[r, "i^1"]
 \& H^1(I,U\otimes\mathbb Q_\ell) \arrow[d, no head, "\wr"] \\
{}
 \& U^I
 \& (U\otimes\mathbb Q_\ell)^I\simeq U^I\otimes_{\mathbb Z}\mathbb Q
 \& (U\otimes D_\ell)^I
 \& {}
 \& H^1(I,U)\otimes_{\mathbb Z}\mathbb Q .
\end{tikzcd}}
\end{equation}
It is immediate that the group $\Phi$ defined by (\SGARefLink{sga7.sga7i.eix.n11-3-4}{11.3.4}) is precisely
$\operatorname{coker}j^o$, and is therefore canonically isomorphic to
$\ker i^1$. Thus we obtain a canonical exact sequence

\begin{equation}\tag{\SGARefTarget{sga7.sga7i.eix.n11-3-8}{11.3.8}}
0\longrightarrow\Phi\longrightarrow H^1(I,U)
 \longrightarrow H^1(I,U)\otimes_{\mathbb Z}\mathbb Q,
\end{equation}
which shows that $\Phi$ \underline{is canonically isomorphic to the
torsion subgroup of $H^1(I,U)$}.

On the other hand, from the exact sequence (\SGARefLink{sga7.sga7i.eix.n11-3-3}{11.3.3}) we deduce, for every
$n=\ell^r$, an isomorphism

\begin{equation}\tag{\SGARefTarget{sga7.sga7i.eix.n11-3-9}{11.3.9}}
H^1(I,T_n)\simeq\operatorname{Hom}(I,T_n)
 \xleftarrow{\sim}\operatorname{Hom}(T,T_n)
 \xleftarrow{\sim}\mathbb Z/n\mathbb Z.
\end{equation}
(Notice that $I$ acts trivially on $T$, and hence on $T_n=T/nT$.) It
follows that, for every finite $I$-$\mathbb Z_\ell$-module $M$, if

\begin{equation}\tag{\SGARefTarget{sga7.sga7i.eix.n11-3-10}{11.3.10}}
M'=\operatorname{Hom}(M,T\otimes D_\ell)
\end{equation}
denotes its ``Cartier dual,'' there is a canonical pairing

\begin{equation}\tag{\SGARefTarget{sga7.sga7i.eix.n11-3-11}{11.3.11}}
H^i(I,M)\otimes H^{1-i}(I,M')
 \longrightarrow H^1(I,T\otimes D_\ell)\simeq D_\ell,
\end{equation}
where the last isomorphism is deduced from the isomorphisms (\SGARefLink{sga7.sga7i.eix.n11-3-9}{11.3.9}) by
passage to the inductive limit over the powers $n$ of $\ell$. It is well
known that these pairings are \underline{perfect dualities} (see, for
example, SGA 5 I 5). Taking $M$ in these pairings to be of the form
$U_n$, so that $M'=U'_n$, and passing to the limit, we obtain a perfect
pairing

\begin{equation}\tag{\SGARefTarget{sga7.sga7i.eix.n11-3-12}{11.3.12}}
H^o(I,U\otimes D_\ell)\times H^1(I,U')\longrightarrow D_\ell,
\end{equation}
where the second factor on the left is a finitely generated
$\mathbb Z_\ell$-module. A further passage to the limit in the pairings
(\SGARefLink{sga7.sga7i.eix.n11-3-12}{11.3.12}), for transition morphisms of the form
$\ell^r\operatorname{id}$, gives a second perfect pairing

\begin{equation}\tag{\SGARefTarget{sga7.sga7i.eix.n11-3-13}{11.3.13}}
H^o(I,U\otimes\mathbb Q_\ell)
 \times H^1(I,U'\otimes\mathbb Q_\ell)
 \longrightarrow D_\ell
\end{equation}
of finite-dimensional vector spaces over $\mathbb Q_\ell$. Let $i'$ and
$j'$ denote the morphisms defined as in (\SGARefLink{sga7.sga7i.eix.n11-3-6}{11.3.6}), with $U$ replaced by
$U'$. It is then easy to see that the following homomorphisms

\begin{center}
\begin{tikzcd}[column sep=6em, row sep=5em]
H^o(I,U\otimes\mathbb Q_\ell) \arrow[r, "j^o"]
  & H^o(I,U\otimes D_\ell) \\
H^1(I,U'\otimes\mathbb Q_\ell)
  & H^1(I,U') \arrow[l, "{i'}^1"']
\end{tikzcd}
\end{center}
are \underline{transposes} of one another under the pairings (\SGARefLink{sga7.sga7i.eix.n11-3-12}{11.3.12})
and (\SGARefLink{sga7.sga7i.eix.n11-3-13}{11.3.13}). It follows immediately that
$\operatorname{coker}j^o$ and $\ker {i'}^1$ are dual to one another by a
canonical pairing with values in $D_\ell$; this is the announced pairing
(\SGARefLink{sga7.sga7i.eix.n11-3-5}{11.3.5}).

% END INCLUDED SGA7I_English_current_buildable_source_20260823/components/167_expose_IX_section_11_3.tex

% BEGIN INCLUDED SGA7I_English_current_buildable_source_20260823/components/168_expose_IX_sections_11_4_and_11_5.tex
% Expose IX, sections 11.4 and 11.5.
% French transcription lines 3148--3191; authority scan physical pp. 466--468
% (scan indices 465--467; source folios 454--456). Physical p. 468 was
% checked through the exact 11.6 boundary. All pages were compared in four
% overlapping 1100-dpi bands. The text, formulas, primes, checks, underlines,
% and displayed arrows are clear at that resolution; no escalation was needed.

\medskip
\noindent
\SGARefTarget{sga7.sga7i.eix.n11-4}{11.4}. The preceding description of a pairing (\SGARefLink{sga7.sga7i.eix.n11-3-1}{11.3.1}), by means of the
expression (\SGARefLink{sga7.sga7i.eix.n11-2-1}{11.2.1}) for the $\ell$-primary component of the group of
connected components, is valid only when $\ell\ne p$. Suppose now that
$A_K$ has \underline{semistable reduction} over $S$. We shall again
define pairings that are perfect dualities

\begin{equation}\tag{\SGARefTarget{sga7.sga7i.eix.n11-4-1}{11.4.1}}
\Phi_o(\ell)\otimes\Phi'_o(\ell)\longrightarrow\mathbb Q/\mathbb Z
\end{equation}
without any restriction on $\ell$. As in \SGARefLink{sga7.sga7i.eix.n11-3}{11.3},\footnote{The printed
source reads ``as in \SGARefLink{sga7.sga7i.eix.n11-4}{11.4}'' here, within section \SGARefLink{sga7.sga7i.eix.n11-4}{11.4} itself. The
reference is corrected to the immediately preceding construction in
\SGARefLink{sga7.sga7i.eix.n11-3}{11.3}.} one should establish that this recovers (up to a sign, perhaps,
which ought to be specified) the pairing defined in (\SGARefLink{sga7.sga7i.eix.n1-2-1}{1.2.1}). This would
therefore establish Conjecture \SGARefLink{sga7.sga7i.eix.n1-3}{1.3} in the case of semistable reduction,
just as it does when one restricts to the $\ell$-primary components for
$\ell\ne p$.

To define (\SGARefLink{sga7.sga7i.eix.n11-4-1}{11.4.1}), we shall give, essentially, an expression for
$\Phi_o(\ell)$ in terms of the Barsotti--Tate group $T_\ell(A_K)$ over
$K$; in the present case this will play the same role as \SGARefLink{sga7.sga7i.eix.n11-2}{11.2} did in
the case treated above.

\medskip
\noindent
\underline{Theorem} \SGARefTarget{sga7.sga7i.eix.n11-5}{11.5}. \underline{Suppose that $S$ is henselian and
that $A_K$ has semistable reduction over $S$. Let $\ell$ be a prime and
consider the monodromy pairing} (\SGARefLink{sga7.sga7i.eix.n9-1-2}{9.1.2}), \underline{and hence the
canonical homomorphism}

\begin{equation*}
\widetilde u_\ell:\underline M_\ell
 \longrightarrow\underline{\check M}'_\ell.
\end{equation*}
\underline{Put $\Phi(\ell)=\operatorname{coker}\widetilde u_\ell$.
Then there is a canonical isomorphism}

\begin{equation}\tag{\SGARefTarget{sga7.sga7i.eix.n11-5-1}{11.5.1}}
\Phi_o(\ell)\xrightarrow{\ \sim\ }\Phi(\ell)\times_S s.
\end{equation}

This statement immediately gives the definition of a perfect pairing
(\SGARefLink{sga7.sga7i.eix.n11-4-1}{11.4.1}). Indeed, applying \SGARefLink{sga7.sga7i.eix.n11-5}{11.5} to $A'_K$ in place of $A_K$ and using
(\SGARefLink{sga7.sga7i.eix.n10-2-2}{10.2.2}), we also obtain a canonical isomorphism

\begin{equation*}
\Phi'_o(\ell)\xrightarrow{\ \sim\ }\Phi'(\ell)\times_S s,
\qquad
\Phi'(\ell)=\operatorname{coker}\bigl(
 {}^t\!\widetilde u_\ell:\underline M'_\ell
 \longrightarrow\underline{\check M}_\ell\bigr).
\end{equation*}

\medskip
\noindent
\underline{Remarks} \SGARefTarget{sga7.sga7i.eix.n11-5-2}{11.5.2}. a) Theorem 11.5 implies in particular that
$\widetilde u_\ell$ is an isogeny, that is, that $u_\ell$ is
separating, including in the case $\ell=p$, where this assertion is not
immediate from the definition in \SGARefLink{sga7.sga7i.eix.n9-6}{9.6}. Notice that we do not use here
the (considerably deeper) Theorem \SGARefLink{sga7.sga7i.eix.n10-4}{10.4}, which will be proved in the next
section.

\smallskip
\noindent
b) Using \SGARefLink{sga7.sga7i.eix.n10-4}{10.4}, which supplies us with a monodromy pairing

\begin{equation*}
u:\underline M\otimes\underline M'\longrightarrow\mathbb Z_S,
\end{equation*}
and denoting by $\Phi$ the cokernel of the corresponding homomorphism

\begin{equation*}
\widetilde u:\underline M\longrightarrow\underline{\check M}',
\end{equation*}
we obtain an equivalent formulation of \SGARefLink{sga7.sga7i.eix.n11-5}{11.5}, independent of the choice
of $\ell$, in the canonical isomorphism

\begin{equation}\tag{\SGARefTarget{sga7.sga7i.eix.n11-5-3}{11.5.3}}
\Phi_o\xrightarrow{\ \sim\ }\Phi\times_S s.
\end{equation}

% END INCLUDED SGA7I_English_current_buildable_source_20260823/components/168_expose_IX_sections_11_4_and_11_5.tex

% BEGIN INCLUDED SGA7I_English_current_buildable_source_20260823/components/169_expose_IX_section_11_6_opening_through_lemma_11_6_13.tex
% Expose IX, section 11.6 opening through Lemma 11.6.13 (statement).
% French transcription lines 3192--3299; authority physical pp. 468--471
% (scan indices 467--470; source folios 456--459). Physical p. 471 was
% checked through the exact proof boundary. All pages were compared in four
% overlapping 1100-dpi bands. Checks, primes, subscripts, hooks, underlines,
% quotient bars, and formula punctuation are clear at that resolution.

\medskip
\noindent
\SGARefTarget{sga7.sga7i.eix.n11-6}{11.6}. To define the isomorphism (\SGARefLink{sga7.sga7i.eix.n11-5-1}{11.5.1}), we may restrict ourselves to
the case in which $S$ is strictly local (\SGARefLink{sga7.sga7i.eix.n11-0}{11.0}), so that $\underline M$
and $\underline M'$ are constant, with values

\begin{equation}\tag{\SGARefTarget{sga7.sga7i.eix.n11-6-1}{11.6.1}}
M=D(T'_o),\qquad M'=D(T_o),
\end{equation}
where $T_o$ and $T'_o$ are respectively the maximal tori of $A_o$ and
$A'_o$. Likewise identifying $\Phi_o$ with an ordinary finite group,
we must define a canonical isomorphism between its component
$\Phi_o(\ell)$ and

\begin{equation*}
\Phi(\ell)=\operatorname{coker}\bigl(
 \widetilde u_\ell:M_\ell\longrightarrow\check M'_\ell\bigr).
\end{equation*}
In fact, we shall directly define a canonical isomorphism

\begin{equation}\tag{\SGARefTarget{sga7.sga7i.eix.n11-6-2}{11.6.2}}
\Phi_o(\ell)\xrightarrow{\ \sim\ }
 \ker\bigl(\widetilde u_\ell\otimes D_\ell:
 M_\ell\otimes D_\ell\longrightarrow
 \check M'_\ell\otimes D_\ell\bigr).
\end{equation}
Since $\Phi_o(\ell)$ is finite and $M_\ell$ and $\check M'_\ell$ have
the same rank, it will first follow that
$\widetilde u_\ell:M_\ell\to\check M'_\ell$ is an isogeny. It then
follows that the right-hand side of (\SGARefLink{sga7.sga7i.eix.n11-6-2}{11.6.2}) is canonically isomorphic
to $\Phi(\ell)$:

\begin{equation}\tag{\SGARefTarget{sga7.sga7i.eix.n11-6-3}{11.6.3}}
\Phi(\ell)=\operatorname{coker}\widetilde u_\ell
 \xrightarrow{\ \sim\ }
 \ker(\widetilde u_\ell\otimes D_\ell).
\end{equation}
Indeed, identifying $M_\ell$ by means of $\widetilde u_\ell$ with a
finite-index submodule of $\check M'_\ell$, we may write the source and
target of $\widetilde u_\ell\otimes D_\ell$ as\footnote{At this point
the printed source labels $V/M_\ell$ as $\Phi(\ell)$ and
$V/\check M'_\ell$ as $\Phi'(\ell)$. Those quotients are instead the
source and target of the map whose kernel is the finite cokernel
$\check M'_\ell/M_\ell$; the contradictory labels are not reproduced.}

\begin{equation*}
M_\ell\otimes D_\ell=V/M_\ell,\qquad
\check M'_\ell\otimes D_\ell=V/\check M'_\ell,
\end{equation*}
where

\begin{equation*}
V=M_\ell\otimes\mathbb Q_\ell
 \xrightarrow{\ \sim\ }\check M'_\ell\otimes\mathbb Q_\ell.
\end{equation*}
Composing (\SGARefLink{sga7.sga7i.eix.n11-6-2}{11.6.2}) with (\SGARefLink{sga7.sga7i.eix.n11-6-3}{11.6.3}) now immediately gives the isomorphism
(\SGARefLink{sga7.sga7i.eix.n11-5-1}{11.5.1}) to be defined, so everything reduces to defining the canonical
isomorphism (\SGARefLink{sga7.sga7i.eix.n11-6-2}{11.6.2}).

Writing $M_\ell\otimes D_\ell$ (respectively,
$\check M'_\ell\otimes D_\ell$) as the inductive limit of the $M_n$
(respectively, the $\check M'_n$), where $n$ ranges over the powers of
$\ell$, it is equivalent to define canonical isomorphisms

\begin{equation}\tag{\SGARefTarget{sga7.sga7i.eix.n11-6-4}{11.6.4}}
{}_n\Phi_o\xrightarrow{\ \sim\ }
 \ker\bigl(\widetilde u_n:M_n\longrightarrow\check M'_n\bigr),
 \qquad n=\ell^{r+1},
\end{equation}
compatible with the evident transition morphisms as $n$ varies. Here,
by an abuse of notation, $\widetilde u_n$ now denotes the homomorphism
obtained by reducing
$\widetilde u_\ell:M_\ell\to\check M'_\ell$ modulo $n$.

For this purpose, consider the group scheme ${}_nA$ over $S$. By the
semistable-reduction hypothesis it is flat and quasi-finite over $S$
(\SGARefLink{sga7.sga7i.eix.n2-2-1}{2.2.1}). There is a canonical filtration

\begin{equation}\tag{\SGARefTarget{sga7.sga7i.eix.n11-6-5}{11.6.5}}
{}_nA\supset({}_nA^o)^f\supset({}_nA^o)^t\supset0,
\end{equation}
where the superscript $f$ (respectively, $t$) denotes the ``fixed part''
(\SGARefLink{sga7.sga7i.eix.n2-2-3}{2.2.3}) (respectively, the ``toric part'' (\SGARefLink{sga7.sga7i.eix.n5-1}{5.1})) of the group scheme to
which it is applied. In the definition of the toric part in \SGARefLink{sga7.sga7i.eix.n5-1}{5.1}, we had
assumed $S$ complete; this is permissible here (\SGARefLink{sga7.sga7i.eix.n11-0}{11.0}). (We saw moreover
in \SGARefLink{sga7.sga7i.eix.n6-1}{6.1} that the construction generalizes to the henselian case, if one
insists.) Put

\begin{equation}\tag{\SGARefTarget{sga7.sga7i.eix.n11-6-6}{11.6.6}}
{}_n\Psi={}_nA/({}_nA^o)^f.
\end{equation}
This is a flat, separated, quasi-finite group scheme over $S$
(SGA 3 $\mathrm{VI}_A$ 3.2 and $\mathrm{VI}_B$ \SGARefLink{sga7.sga7i.eix.n9-2}{9.2}), with generic fibre

\begin{equation*}
{}_n\Psi_K=({}_nA_K)/({}_nA_K)^f,
\end{equation*}
where, by definition,
$({}_nA_K)^f=({}_nA^o)^f_K$ denotes the $r$-th constituent
$(n=\ell^{r+1})$ of the Barsotti--Tate pro-group
$T_\ell(A_K)^f$. The first of relations (\SGARefLink{sga7.sga7i.eix.n5-5-8}{5.5.8}) therefore gives a
canonical isomorphism

\begin{equation}\tag{\SGARefTarget{sga7.sga7i.eix.n11-6-7}{11.6.7}}
{}_n\Psi_K\xrightarrow{\ \sim\ }\underline M_{nK},
\qquad
\underline M_n\mathrel{\mathop:}=\underline M
 \otimes_{\mathbb Z}(\mathbb Z/n\mathbb Z).
\end{equation}
In fact, we have:

\medskip
\noindent
\underline{Lemma} \SGARefTarget{sga7.sga7i.eix.n11-6-8}{11.6.8}. \underline{The group scheme ${}_n\Psi$ over
$S$ defined in (\SGARefLink{sga7.sga7i.eix.n11-6-6}{11.6.6}) is \emph{\'etale} over $S$, and the isomorphism
(\SGARefLink{sga7.sga7i.eix.n11-6-7}{11.6.7}) extends uniquely to a morphism}

\begin{equation}\tag{\SGARefTarget{sga7.sga7i.eix.n11-6-8-1}{11.6.8.1}}
{}_n\Psi\lhook\joinrel\longrightarrow\underline M_n
\end{equation}
\underline{which is an open immersion.}

To see that ${}_n\Psi$ is \'etale, in view of (\SGARefLink{sga7.sga7i.eix.n11-6-7}{11.6.7}) it remains only
to verify this at the points of its special fibre, or equivalently to
verify that this fibre is \'etale. This is immediate, since the fibre in
question is

\begin{equation*}
{}_nA_o/{}_nA_o^o\subset A_o/A_o^o=\Phi_o.
\end{equation*}
It follows that ${}_n\Psi$ is normal; and since $\underline M_n$ is
finite over $S$, the existence of the extension (\SGARefLink{sga7.sga7i.eix.n11-6-8-1}{11.6.8.1}) follows at
once. This is a morphism of \'etale group schemes over $S$, so its
kernel is \'etale over $S$. Since ${}_n\Psi$, and hence this kernel, is
separated, while its generic fibre is the unit group by (\SGARefLink{sga7.sga7i.eix.n11-6-7}{11.6.7}), the
kernel is the unit group (\SGARefLink{sga7.sga7i.eix.n3-1-3}{3.1.3}). Thus (\SGARefLink{sga7.sga7i.eix.n11-6-8-1}{11.6.8.1}) is an \'etale
monomorphism, hence an open immersion (EGA IV 17.9.1), proving the lemma.

Now put

\begin{equation}\tag{\SGARefTarget{sga7.sga7i.eix.n11-6-9}{11.6.9}}
{}_n\Phi=({}_nA)^f/({}_nA^o)^f\subset{}_n\Psi.
\end{equation}
Thus ${}_n\Phi$ is a finite open subscheme of ${}_n\Psi$, since
$({}_nA)^f$ is a finite open subgroup scheme of ${}_nA$. More
precisely, it is immediate that

\begin{equation}\tag{\SGARefTarget{sga7.sga7i.eix.n11-6-10}{11.6.10}}
{}_n\Phi\xrightarrow{\ \sim\ }({}_n\Psi)^f.
\end{equation}
Composing (\SGARefLink{sga7.sga7i.eix.n11-6-9}{11.6.9}) with (\SGARefLink{sga7.sga7i.eix.n11-6-8-1}{11.6.8.1}), we obtain an open immersion

\begin{equation}\tag{\SGARefTarget{sga7.sga7i.eix.n11-6-11}{11.6.11}}
{}_n\Phi\lhook\joinrel\longrightarrow\underline M_n.
\end{equation}
On the other hand,

\begin{equation*}
{}_n\Phi\times_Ss={}_n\Psi\times_Ss={}_nA_o/{}_nA_o^o,
\end{equation*}
and, since $A_o^o$ is $n$-divisible (\SGARefLink{sga7.sga7i.eix.n2-2-1}{2.2.1}), the last member is
immediately identified with ${}_n\Phi_o$. We therefore obtain an
isomorphism

\begin{equation}\tag{\SGARefTarget{sga7.sga7i.eix.n11-6-12}{11.6.12}}
{}_n\Phi\times_Ss\xrightarrow{\ \sim\ }{}_n\Phi_o,
\end{equation}
which shows that ${}_n\Phi$ is characterized, up to unique isomorphism,
as the finite \'etale scheme over $S$ whose special fibre is
${}_n\Phi_o$. Thus, to define the desired isomorphism (\SGARefLink{sga7.sga7i.eix.n11-6-4}{11.6.4}) by means
of the open immersion (\SGARefLink{sga7.sga7i.eix.n11-6-11}{11.6.11}), it is enough to prove the following.

\medskip
\noindent
\underline{Lemma} \SGARefTarget{sga7.sga7i.eix.n11-6-13}{11.6.13}. \underline{The image of the canonical
homomorphism (\SGARefLink{sga7.sga7i.eix.n11-6-11}{11.6.11}) is equal to the kernel of}

\begin{equation*}
\widetilde u_n:\underline M_n\longrightarrow\underline{\check M}'_n.
\end{equation*}

The required compatibility of the homomorphisms (\SGARefLink{sga7.sga7i.eix.n11-6-11}{11.6.11}) with the
transition homomorphisms as $n$ varies is immediate from the
definitions. Consequently Lemma \SGARefLink{sga7.sga7i.eix.n11-6-13}{11.6.13} will complete the proof of
Theorem \SGARefLink{sga7.sga7i.eix.n11-5}{11.5}.

% END INCLUDED SGA7I_English_current_buildable_source_20260823/components/169_expose_IX_section_11_6_opening_through_lemma_11_6_13.tex

% BEGIN INCLUDED SGA7I_English_current_buildable_source_20260823/components/170_expose_IX_lemma_11_6_13_proof_and_remarks_11_7.tex
% Expose IX, proof of Lemma 11.6.13 and Remarks 11.7.
% French transcription lines 3304--3391; authority physical pp. 472--474
% (scan indices 471--473; source folios 460--462). Physical p. 474 was
% checked through the exact 11.8 boundary. All pages were compared in four
% overlapping 1100-dpi bands. The two compact diagrams, their nodes, ten plus
% five arrows, equality/isomorphism marks, labels, bars, and punctuation are
% all decisive at that resolution; no higher-resolution escalation was needed.

To prove the lemma, note that

\begin{equation*}
K_n=\ker\widetilde u_n
\end{equation*}
is also the left annihilator of the canonical pairing

\begin{equation*}
u_n:\underline M_n\times\underline M'_n
 \longrightarrow(\mathbb Z/n\mathbb Z)_S,
\end{equation*}
whose explicit definition is given in \SGARefLink{sga7.sga7i.eix.n9-4}{9.4} and \SGARefLink{sga7.sga7i.eix.n9-6}{9.6} in terms of the
extensions

\begin{equation}\tag{\SGARefTarget{sga7.sga7i.eix.n11-6-14}{11.6.14}}
\resizebox{\textwidth}{!}{%
\begin{tikzcd}[ampersand replacement=\&, column sep=small, row sep=large]
0 \arrow[r]
 \& {\ensuremath{\check{\underline M}'_n(1)}} \arrow[r]
 \& F_n \arrow[r] \arrow[d, equal]
 \& R_n \arrow[r]
 \& 0
 \& {,}
 \& 0 \arrow[r]
 \& {\ensuremath{\check{\underline M}_n(1)}} \arrow[r]
 \& F'_n \arrow[r] \arrow[d, equal]
 \& R'_n \arrow[r]
 \& 0 \\
{}
 \& {}
 \& ({}_nA^o)^f
 \& {}
 \& {}
 \& {}
 \& {}
 \& {}
 \& ({}_nA'^o)^f
 \& {}
 \& {}
\end{tikzcd}}
\end{equation}
(where $R_n$ and $R'_n$ are Cartier dual to one another), and in terms
of the mixed extension $X_{n\eta}$ of $E_{n\eta}$ by $F_{n\eta}$,
which is the generic fibre of (\SGARefLink{sga7.sga7i.eix.n11-6-5}{11.6.5}),\footnote{The printed source
refers here to nonexistent formula (11.7.5); the filtration whose
generic fibre is being used is (\SGARefLink{sga7.sga7i.eix.n11-6-5}{11.6.5}).} where we put

\begin{equation}\tag{\SGARefTarget{sga7.sga7i.eix.n11-6-15}{11.6.15}}
\begin{tikzcd}[column sep=large, row sep=large]
E_n=D(F_n),
 & 0 \arrow[r]
 & R_n \arrow[r] \arrow[d, no head, "\wr"']
 & E_n \arrow[r]
 & \underline M_n \arrow[r]
 & 0
 & {.} \\
{}
 & {}
 & D(R'_n)
 & {}
 & {}
 & {}
 & {}
\end{tikzcd}
\end{equation}
It follows immediately from Proposition \SGARefLink{sga7.sga7i.eix.n9-4-7}{9.4.7}\footnote{The printed
source cites \SGARefLink{sga7.sga7i.eix.n9-4-6}{9.4.6}, but the required extension criterion is Proposition
\SGARefLink{sga7.sga7i.eix.n9-4-7}{9.4.7}.} that $K_n$ is the largest among the finite \'etale subgroups
$K'_n$ of $\underline M_n$ for which the inverse image
$\overline X_{n\eta}$ of $K'_n$ in the mixed extension $X_{n\eta}$,
regarded as a mixed extension of $\overline E_{n\eta}$ by $F_{n\eta}$,
extends to a mixed extension of $\overline E_n$ by $F_n$. Here
$\overline E_n$ is the evident subgroup extension of $K'_n$ by
$R_n$.\footnote{The printed source says ``of $K_n$'' in this general
criterion, although the subgroup being varied at this point is $K'_n$;
the specialization $K'_n=K_n$ is made only below.}

This condition is plainly satisfied for the subgroup ${}_n\Phi$ of
$\underline M_n$ in (\SGARefLink{sga7.sga7i.eix.n11-6-11}{11.6.11}), since the corresponding mixed extension
$({}_nA)^f$ over $S$ is available. Hence ${}_n\Phi\subset K_n$. To
prove the opposite inclusion, consider the mixed extension
$\overline X_n$ of $\overline E_n$ by $F_n$ extending
$\overline X_{n\eta}$, defined as above by taking $K'_n=K_n$. Thus it
is an extension

\begin{equation*}
0\longrightarrow F_n\longrightarrow\overline X_n
 \longrightarrow K_n\longrightarrow0,
\end{equation*}
and there is a canonical homomorphism

\begin{equation}\tag{*}
F_n=({}_nA^o)^f\lhook\joinrel\longrightarrow A,
\end{equation}
which on the generic fibre is induced by the canonical homomorphism

\begin{equation}\tag{**}
\overline X_{n\eta}\lhook\joinrel\longrightarrow X_{n\eta}
 ={}_nA_K\lhook\joinrel\longrightarrow A_K.
\end{equation}
It follows immediately from (\SGARefLink{sga7.sga7i.eix.n5-9-2}{5.9.2}) that the homomorphism $(**)$ extends
to a unique homomorphism

\begin{equation}\tag{***}
\overline X_n\longrightarrow A
\end{equation}
inducing $(*)$. Since this last homomorphism maps the generic fibre into
${}_nA$, which is a closed subscheme of $A$, it factors through
$\overline X_n\to{}_nA$; and since $\overline X_n$ is finite, it factors
through $({}_nA)^f$. Looking at the induced homomorphism on generic
fibres, we conclude that

\begin{equation*}
\overline X_{n\eta}\subset({}_nA)^f_\eta,
\quad\text{i.e.}\quad
K_{n\eta}\subset{}_n\Phi_\eta,
\quad\text{i.e.}\quad
K_n\subset{}_n\Phi.
\end{equation*}
This proves the lemma.

\medskip
\noindent
\underline{Remarks} \SGARefTarget{sga7.sga7i.eix.n11-7}{11.7}. The construction made in \SGARefLink{sga7.sga7i.eix.n11-6}{11.6}\footnote{The
printed source says ``the construction made in \SGARefLink{sga7.sga7i.eix.n11-7}{11.7}'' at the opening of
\SGARefLink{sga7.sga7i.eix.n11-7}{11.7}; the reference is to the construction just completed in \SGARefLink{sga7.sga7i.eix.n11-6}{11.6}.}
shows that, from the extensions (\SGARefLink{sga7.sga7i.eix.n11-6-14}{11.6.14}) over $S$ and the mixed
extension $X_{n\eta}={}_nA_K$ of $E_n$ by $F_n$, one can reconstruct,
up to unique isomorphism, the mixed extension (\SGARefLink{sga7.sga7i.eix.n11-6-5}{11.6.5}), and hence also
the subgroup

\begin{equation*}
({}_nA)^f\supset({}_nA^o)^f
\end{equation*}
of ${}_nA$. Indeed, by \SGARefLink{sga7.sga7i.eix.n9-4}{9.4} these data determine the pairing

\begin{equation*}
u_n:\underline M_n\times\underline M'_n
 \longrightarrow(\mathbb Z/n\mathbb Z)_S,
\end{equation*}
and the preceding construction gives $({}_nA)^f$ as the mixed extension
$\overline X_n$ of $\overline E_n$ by $F_n$, extending the mixed
extension $\overline X_{n\eta}$ induced by $X_{n\eta}$. Here
$\overline E_n$ is the inverse image of $\ker\widetilde u_n$ under
$E_n\to\underline M_n$. By (9.4.8), this mixed extension
$\overline X_n$ is thus determined \underline{up to unique
isomorphism}. We then recover ${}_nA$ by gluing $\overline X_n$ and
$X_{n\eta}$ along their common open subscheme $\overline X_{n\eta}$.

Thus, if the structural data (\SGARefLink{sga7.sga7i.eix.n9-6-1}{9.6.1}) and the mixed extension

\begin{equation*}
X_\eta=T_\ell(A_K)
\end{equation*}
of $E_\eta$ by $F_\eta$ are known---equivalently, if the extensions
(\SGARefLink{sga7.sga7i.eix.n11-6-14}{11.6.14}) and $X_{n\eta}$ are known for every $n=\ell^{r+1}$,
$r\geq0$---then one can reconstruct the inductive (or, if preferred,
projective) system of mixed extensions (\SGARefLink{sga7.sga7i.eix.n11-6-5}{11.6.5}).

% END INCLUDED SGA7I_English_current_buildable_source_20260823/components/170_expose_IX_lemma_11_6_13_proof_and_remarks_11_7.tex

% BEGIN INCLUDED SGA7I_English_current_buildable_source_20260823/components/171_expose_IX_sections_11_8_through_11_11.tex
% Expose IX, sections 11.8--11.11.
% French transcription lines 3393--3446; authority physical pp. 474--476
% (scan indices 473--475; source folios 462--464). Physical p. 477 was
% checked through the exact section-12 boundary. All pages were compared in
% four overlapping 1100-dpi bands; every formula, accent, subscript, arrow,
% underline, and punctuation mark was decisive without higher escalation.
% The page images restore two transcription slips: the quotient sign in
% (11.8.1), and Phi rather than Psi on the left of (11.8.3).

\SGARefTarget{sga7.sga7i.eix.n11-8}{11.8}. As in \SGARefLink{sga7.sga7i.eix.n8-3}{8.3}, regard the separable closure $\overline K$ of $K$ as the
inductive limit of its finite subextensions $K'$. For every such $K'$, let
$A[K']$ denote the N\'{e}ron model of

\begin{equation*}
A_{K'}=A_K\otimes_K K'.
\end{equation*}
It gives a finite \'{e}tale group scheme

\begin{equation}\tag{\SGARefTarget{sga7.sga7i.eix.n11-8-1}{11.8.1}}
\Phi_o[K']=A[K']_o/A[K']_o^o
\end{equation}
of connected components, defined over the residue field of the normalization
$S[K']$ of $S$ in $K'$. For simplicity, suppose that the residue field of
$S$ is separably closed, so that the same is true of the residue fields of
the $S[K']$. The groups (\SGARefLink{sga7.sga7i.eix.n11-8-1}{11.8.1}) may then be identified with ordinary finite
groups $\Phi[K']$, and it is immediate that, as $K'$ varies, these form an
inductive system, which we now propose to study. We shall be interested in
particular in the group

\begin{equation}\tag{\SGARefTarget{sga7.sga7i.eix.n11-8-2}{11.8.2}}
\Psi=\varinjlim_{K'}\Phi[K'].
\end{equation}
We know from \SGARefLink{sga7.sga7i.eix.n3-6}{3.6} that, for $K'$ sufficiently large, say
$K'\supset K'_o$, $A[K']$ has semistable reduction. Hence for
$K'\supset K'_o$ the transition morphisms are injective (\SGARefLink{sga7.sga7i.eix.n3-3}{3.3}). Moreover,
with the notation of \SGARefLink{sga7.sga7i.eix.n5-5}{5.5}, retained in the present section, there is for
$K'\supset K'_o$ a canonical isomorphism

\begin{equation}\tag{\SGARefTarget{sga7.sga7i.eix.n11-8-3}{11.8.3}}
\Phi[K'](\ell)\simeq
\ker\bigl(\widetilde u[K']_\ell\otimes D_\ell:
 M_\ell\otimes D_\ell\longrightarrow
 \check M'_\ell\otimes D_\ell\bigr),
\end{equation}
coming from (\SGARefLink{sga7.sga7i.eix.n11-6-2}{11.6.2}) applied to $A[K']$. One also verifies immediately that
the corresponding immersions

\begin{equation*}
\Phi[K'](\ell)\lhook\joinrel\longrightarrow M_\ell\otimes D_\ell
\end{equation*}
are compatible with the transition morphisms as $K'$ varies. Passing to the
limit therefore gives a canonical monomorphism

\begin{equation}\tag{\SGARefTarget{sga7.sga7i.eix.n11-8-4}{11.8.4}}
\Psi(\ell)\mathrel{\mathop{=}^{\mathrm{def}}}
 \varinjlim_{K'}\Phi[K'](\ell)
 \lhook\joinrel\longrightarrow M_\ell\otimes D_\ell
\end{equation}
from the $\ell$-primary component of $\Psi$; summing over all $\ell$ gives

\begin{equation}\tag{\SGARefTarget{sga7.sga7i.eix.n11-8-5}{11.8.5}}
\Psi\mathrel{\mathop{=}^{\mathrm{def}}}
 \varinjlim_{K'}\Phi[K']
 \longrightarrow M\otimes\mathbb Q/\mathbb Z.
\end{equation}

\medskip
\noindent
\underline{Theorem} \SGARefTarget{sga7.sga7i.eix.n11-9}{11.9}. \underline{The canonical homomorphism (\SGARefLink{sga7.sga7i.eix.n11-8-5}{11.8.5})
is an isomorphism}.

We already know that it is a monomorphism. To prove that it is an
isomorphism, we must use the characterization (\SGARefLink{sga7.sga7i.eix.n11-8-3}{11.8.3}) of the image of
$\Phi[K'](\ell)$ and the behaviour of $\widetilde u[K']$ as $K'$ varies.
That behaviour is given by (\SGARefLink{sga7.sga7i.eix.n10-3-4}{10.3.4}).

% SGA-SRC-001281 — internal reference omitted in source; target not recoverable
We may reduce to the case in which $S$ is strictly local. To prove \SGARefLink{sga7.sga7i.eix.n11-9}{11.9} we
are then reduced to the following familiar assertion, applied to $S[K'_o]$:
if $S$ is a strictly local trait with fraction field $K$, then for every
integer $n\geq1$ there is a separable extension $K'$ of $K$ whose degree is
divisible by $n$.
\footnote{The printed source appends ``cf. I'' followed by a blank internal
reference. No target is supplied, so none is invented here.}

\medskip
\noindent
\underline{Corollary} \SGARefTarget{sga7.sga7i.eix.n11-10}{11.10}. \underline{Under the hypotheses ($S$
complete) and with the notation of \SGARefLink{sga7.sga7i.eix.n8-3}{8.3}, the ind-finite group scheme over
$K$ of connected components of the group scheme $A_K^{\flat}$,
namely $A_K^{\flat}/{A_K^{\flat}}^o$, is canonically isomorphic to}
$\underline M_K\otimes\mathbb Q/\mathbb Z$, \underline{where
$\underline M_K$ is the twisted constant group over $K$ defined in
(\SGARefLink{sga7.sga7i.eix.n8-1-8}{8.1.8})}.

This follows immediately from (8.2.4) and \SGARefLink{sga7.sga7i.eix.n11-9}{11.9}, in view of the definition
of $\underline M_K$.

\medskip
\noindent
\underline{Remarks} \SGARefTarget{sga7.sga7i.eix.n11-11}{11.11}. a) If one restricts to the $\ell$-primary
components for $\ell\neq p$, Theorem \SGARefLink{sga7.sga7i.eix.n11-9}{11.9} is immediate in terms of \SGARefLink{sga7.sga7i.eix.n11-2}{11.2} and
does not require the more delicate considerations of \SGARefLink{sga7.sga7i.eix.n11-6}{11.6} (via \SGARefLink{sga7.sga7i.eix.n9-4}{9.4}).
Likewise, if $K$ is assumed to be \underline{of characteristic zero}, \SGARefLink{sga7.sga7i.eix.n11-9}{11.9}
can be proved without using (\SGARefLink{sga7.sga7i.eix.n11-8-3}{11.8.3}): it is immediate that the image of
${}_n\Phi[K'](\ell)$ in

\begin{equation*}
M_n={}_n(M_\ell\otimes D_\ell)
\end{equation*}
contains the image of the group ${}_nA(K')$ of $K'$-valued points of
${}_nA$. For $K'$ sufficiently large this group is equal to
${}_nA(\overline K)$, and therefore maps onto $M_n$. If
$\operatorname{char}K=p$, on the other hand, it appears that there is no
``trivial'' proof---that is, no proof avoiding \SGARefLink{sga7.sga7i.eix.n9-4}{9.4}---of \SGARefLink{sga7.sga7i.eix.n11-10}{11.10} for the
$p$-primary component.

b) The existence of an isomorphism

\begin{equation*}
\Psi\simeq(\mathbb Q/\mathbb Z)^r,
\end{equation*}
together with the correct value of $r$, was conjectured in 1964 by
J.-P. Serre.

% END INCLUDED SGA7I_English_current_buildable_source_20260823/components/171_expose_IX_sections_11_8_through_11_11.tex

% BEGIN INCLUDED SGA7I_English_current_buildable_source_20260823/components/172_expose_IX_section_12_1_through_12_2.tex
\medskip
\noindent
12. \underline{N\'eron model of a Jacobian and the Picard--Lefschetz
formula}

Let us first recall the following general results of M. Raynaud
[22, Ths.~5 and 7; 22 bis].

\medskip
\noindent
\underline{Theorem} \SGARefTarget{sga7.sga7i.eix.n12-1}{12.1} (M. Raynaud). \underline{Let $S$ be a strictly
local trait, and let $f:X\to S$ be a scheme proper and flat over $S$, with
one-dimensional fibres, such that}

\begin{equation}\tag{\SGARefTarget{sga7.sga7i.eix.n12-1-1}{12.1.1}}
 f_*\underline{\mathcal O}_X\simeq\underline{\mathcal O}_S,
\end{equation}

\noindent
\underline{and such that the local rings of $X$ are factorial. For every
maximal point $x_i$ of $X_o$, put}

\begin{equation}\tag{\SGARefTarget{sga7.sga7i.eix.n12-1-2}{12.1.2}}
 d_i=\operatorname{length}\underline{\mathcal O}_{X_o,x_i},
 \qquad d_i^t=\mu_i d_i,
\end{equation}

\noindent
\underline{where $\mu_i$ is the radicial multiplicity of $k(x_i)$ over
$k$} (EGA IV, 4.7.4) \underline{(thus a power of the characteristic
exponent of $k$). Let $d$ (respectively $d^t$) be the gcd of the integers
$d_i$ (respectively $d_i^t$); hence $d=d^t$ if $k$ is perfect. Put}

\begin{equation}\tag{\SGARefTarget{sga7.sga7i.eix.n12-1-3}{12.1.3}}
 P=\underline{\operatorname{Pic}}_{X/S}:
   (\operatorname{Sch})_{/S}^{,o}\longrightarrow(\operatorname{Ab}).
\end{equation}

\noindent
a) \underline{Consider the following conditions}:

\smallskip
\noindent
1) $d^t=1$.

\smallskip
\noindent
2) $X$ \underline{has a zero-cycle of degree $1$}.

\smallskip
\noindent
3) \underline{There exists an invertible sheaf $L$ on $X\times P^o$ such
that the corresponding homomorphism}
$P^o_\eta\to\underline{\operatorname{Pic}}_{X_\eta/\eta}=P_\eta$
\underline{is the canonical inclusion}. (It is known, by J.-P. Murre,
that $P_\eta$ is representable by a group scheme locally of finite type
over $\eta$; see, for example, SGA 6, XII, 1.5.)

\smallskip
\noindent
4) $X$ \underline{is cohomologically flat over $S$} (EGA III, 7.8.1)
(equivalently, cohomologically flat in dimension zero, or equivalently
$k\xrightarrow{\sim}H^o(X_o,\underline{\mathcal O}_{X_o})$),
\underline{and $d=1$}.

\smallskip
\noindent
5) $P$ \underline{is representable} (equivalently, representable by a
group scheme smooth over $S$).

\smallskip
\noindent
6) $P^o$ \underline{is representable} (equivalently, representable by a
group scheme smooth over $S$). (For the definition of $P^o$, see SGA 3,
$\mathrm{VI}_B$, \SGARefLink{sga7.sga7i.eix.n3-1}{3.1}.)

\smallskip
\noindent
7) $P^o$ \underline{is separated over $S$}, that is, it satisfies the
valuative criterion of separatedness of EGA II, 7.2.3 (equivalently,
every section $s$ of $P^o$ over $S$ for which $s(\eta)=0$ is zero).

\smallskip
\noindent
8) $d=1$.

\smallskip
\noindent
\underline{Then one has the implications}

\begin{equation}\tag{\SGARefTarget{sga7.sga7i.eix.n12-1-4}{12.1.4}}
 1)\Longleftrightarrow2)\Longrightarrow3)\Longrightarrow4)
 \Longleftrightarrow5)\Longleftrightarrow6)\Longleftrightarrow7)
 \Longrightarrow8).
\end{equation}

\noindent
\underline{In particular, if $d=d^t$, for example if $k$ is perfect, the
conditions 1) through 8) are equivalent}.

\medskip
\noindent
b) \underline{Suppose that $d^t=1$, so that $P$ is representable by a
group scheme smooth over $S$. Then the schematic closure of the identity
section of $P$ is a group subscheme $E$ which is \`etale over $S$. If $I$
is the set of reduced irreducible components $X_o^i$ of $X_o$, one obtains
an isomorphism}

\begin{equation}\tag{\SGARefTarget{sga7.sga7i.eix.n12-1-5}{12.1.5}}
 \mathbb Z^I\simeq\Gamma(S,E)
\end{equation}

\noindent
\underline{by associating to every $i\in I$ the section $S_i$ of $P$ over
$S$ (zero over $\eta$) defined by the invertible sheaf
$\underline{\mathcal O}_X(X_o^i)$ on $X$}\footnote{The printed source says
``on $S$'' here and later writes $X_i^o$ for the same divisor. The sheaf is
the invertible sheaf on $X$ associated with the component $X_o^i$; the
English text keeps that notation consistently.}. (Because the local rings of $X$
are factorial by hypothesis, the cycles $X_o^i$ are indeed divisors; thus
$\underline{\mathcal O}_X(X_o^i)$ has its stated meaning as the invertible
sheaf associated with that divisor; see EGA IV, 21.2.8.1.) \underline{One
has}

\begin{equation}\tag{\SGARefTarget{sga7.sga7i.eix.n12-1-6}{12.1.6}}
 P^o\cap E=\underline{\text{the zero section of }P}.
\end{equation}

\noindent
c) \underline{Under the conditions of b), the quotient}

\begin{equation}\tag{\SGARefTarget{sga7.sga7i.eix.n12-1-7}{12.1.7}}
 Q=P/E
\end{equation}

\noindent
\underline{is representable by a smooth separated group scheme over $S$,
and this group scheme is N\'eron, that is, it satisfies the universal
property (\SGARefLink{sga7.sga7i.eix.n1-1-2}{1.1.2}) in terms of its special fibre}.

\medskip
\noindent
d) \underline{Under the conditions of c), suppose in addition that the
geometric generic fibre of $X$ is integral. Then there is a canonical
exact sequence}

\begin{equation}\tag{\SGARefTarget{sga7.sga7i.eix.n12-1-8}{12.1.8}}
 0\longrightarrow A\longrightarrow Q
   \xrightarrow{\ \deg\ }\mathbb Z_S\longrightarrow0,
\end{equation}

\noindent
\underline{where $\deg$ is induced by the ``degree'' homomorphism
$P\to\mathbb Z_S$, and where $A$ is a smooth separated group scheme of
finite type over $S$, which is likewise N\'eron, that is, a N\'eron model
of}

\begin{equation}\tag{\SGARefTarget{sga7.sga7i.eix.n12-1-9}{12.1.9}}
 A_\eta=P_\eta^o=
 \underline{\operatorname{Pic}}^o_{X_\eta/\eta}.
\end{equation}

\noindent
\underline{Moreover, the composite homomorphism $P^o\to P\to Q$ induces
an isomorphism}

\begin{equation}\tag{\SGARefTarget{sga7.sga7i.eix.n12-1-10}{12.1.10}}
 P^o\xrightarrow{\ \sim\ }A^o.
\end{equation}

\medskip
\noindent
\SGARefTarget{sga7.sga7i.eix.n12-1-11}{12.1.11}. The case of interest to us is that in which $X_\eta$ is smooth
over $\eta$, so that, in view of (\SGARefLink{sga7.sga7i.eix.n12-1-1}{12.1.1}), the geometric generic fibre of
$X$ is smooth and connected (hence integral), and in which $X_\eta$ has a
zero-cycle of degree one. All the conditions considered in \SGARefLink{sga7.sga7i.eix.n12-1}{12.1} (and in
particular conditions 1) through 8) of a)) then hold. Moreover,
$A_\eta=P_\eta^o$ is then an abelian scheme, and d) gives us a very
explicit construction of its N\'eron model $A$. In particular, (\SGARefLink{sga7.sga7i.eix.n12-1-10}{12.1.10})
gives an isomorphism

\begin{equation}\tag{\SGARefTarget{sga7.sga7i.eix.n12-1-12}{12.1.12}}
 A_o^{,o}\simeq\underline{\operatorname{Pic}}^o_{X_o/k},
\end{equation}

\noindent
which shows in particular that $A$ \underline{has semistable reduction
over $S$ if and only if}
$\underline{\operatorname{Pic}}^o_{X_o/k}$ \underline{has unipotent rank
zero}. When $X$ is regular and $X_o$ is separable over $k$, this is known
to hold if and only if the geometric special fibre $\overline X_o$ has no
singular points other than ordinary double points (that is, points whose
completed local ring is the same as that of the union of the two
coordinate axes of affine $2$-space at the origin). Verification of this
fact, together with the ``well-known'' description of the Picard scheme
of a proper separable curve over a field $k$, using EGA IV, 21.8.5, is left
to the reader as an exercise. Notice also that the implication
``$\overline X_o$ has ordinary double points'' $\Longrightarrow$
``$A_\eta=\underline{\operatorname{Pic}}^o_{X_\eta/\eta}$ has semistable
reduction'' follows immediately from the ``geometric monodromy theorem''
III\footnote{The printed source leaves the internal locator following
``III'' blank; no target is invented here.}, in view of the Galois
criterion for semistable reduction in \SGARefLink{sga7.sga7i.eix.n3-5}{3.5}. This argument, however, does
not give the important isomorphism (\SGARefLink{sga7.sga7i.eix.n12-1-10}{12.1.10}) and its consequence
(\SGARefLink{sga7.sga7i.eix.n12-1-12}{12.1.12}).

\medskip
\noindent
\SGARefTarget{sga7.sga7i.eix.n12-2}{12.2}. For the proof of \SGARefLink{sga7.sga7i.eix.n12-1}{12.1}, we refer to [22 bis]; in fact we shall need
only the case where $X_\eta$ is smooth over $\eta$. Let us also note that
[22 bis] proves relations between Picard schemes and N\'eron models under
conditions substantially more general than those of \SGARefLink{sga7.sga7i.eix.n12-1}{12.1} c) and d). As
their statements are more technical, we have chosen not to include them
in this review. We mention only that b), c), and d) are easy consequences
of a), and that in a) the only delicate points are the implications
$3)\Longrightarrow4)\Longrightarrow5)$ and $7)\Longrightarrow4)$, the
first being the most difficult. This point is trivial, however, if one
assumes that $X_o$ is separable over $k$ (which implies that $X$ has a
section over $S$, since $S$ is strictly local, and hence that condition
1), $d^t=1$, holds). We advise the reader to give the proof in this
particularly simple case, assuming in addition, if necessary, that
$X_\eta$ is smooth, so that $P_\eta^o$ is an abelian scheme. One may then
use its N\'eron model $A$, and $A^o$ represents
$\underline{\operatorname{Pic}}^o_{X/S}$. We shall use \SGARefLink{sga7.sga7i.eix.n12-1}{12.1} only in this
particular case, and under the further assumption that $\overline X_o$
has no singularities other than ordinary double points, that is, by
(\SGARefLink{sga7.sga7i.eix.n12-1-11}{12.1.11}), that $A_\eta$ has semistable reduction. This is the case
considered in 12.5.4, to which we restrict ourselves from now on.

% END INCLUDED SGA7I_English_current_buildable_source_20260823/components/172_expose_IX_section_12_1_through_12_2.tex

% BEGIN INCLUDED SGA7I_English_current_buildable_source_20260823/components/173_expose_IX_section_12_3.tex
\medskip
\noindent
\SGARefTarget{sga7.sga7i.eix.n12-3}{12.3}. Let $C$ be a proper separable curve over a field $k$, which for
simplicity we suppose separably closed (or even algebraically closed, if
the reader prefers). If $\widetilde C$ is the normalization of $C$, it is
known that $\underline{\operatorname{Pic}}^o_{\widetilde C/k}$ has
``reductive rank'' zero (that is, its maximal torus is zero). Consequently
the maximal torus of $\underline{\operatorname{Pic}}^o_{C/k}$ is contained
in

\begin{equation*}
 N=\ker\bigl(\underline{\operatorname{Pic}}_{C/k}
       \longrightarrow\underline{\operatorname{Pic}}_{\widetilde C/k}\bigr),
\end{equation*}

\noindent
and is therefore canonically isomorphic to the maximal torus of this last
group. The structure of $N$ is easily made explicit by EGA IV, 21.8.5:
$N$ is canonically an extension of a torus $T$, which we shall describe,
by a connected unipotent group $U$ (which itself has, over $k$, a
composition series whose quotients are isomorphic to the additive group
$\mathbb G_a$). Moreover, if $T'$ is the maximal torus of $N$, then
$T'\cap U$ is both unipotent and of multiplicative type, and is therefore
zero (SGA 3, XVII, 2.4), while the image of $T'$ in $T$ is all of $T$
(SGA 3, XII, 7.1 e)). Thus $T'\to T$ is an isomorphism. Hence $T$ is
canonically isomorphic to the maximal torus of $N$, and therefore to that
of $\underline{\operatorname{Pic}}^o_{C/k}$.
\footnote{The printed source has $\underline{\operatorname{Pic}}^o_{X/k}$
at the end of this argument. The curve throughout \SGARefLink{sga7.sga7i.eix.n12-3}{12.3} is $C$, and the
preceding and following occurrences both read
$\underline{\operatorname{Pic}}^o_{C/k}$; the English text restores $C$.}

To make the structure of $N$ explicit, it is convenient to suppose that
$k$ is separably closed. The irreducible components $C_i$ of $C$ are then
geometrically irreducible, and the nonnormal points of $C$ are radicial
over $k$. Let $J$ be the set of irreducible components of $\widetilde C$,
and for every $x\in C$ let $J(x)$ be the set of ``branches'' of $C$ through
$x$ (that is, the points of $\widetilde C$ above $x$). Finally put

\begin{equation}\tag{\SGARefTarget{sga7.sga7i.eix.n12-3-1}{12.3.1}}
 R(x)=\mathbb Z^{J(x)}/\operatorname{Im}\mathbb Z,
\end{equation}

\noindent
where $\mathbb Z\to\mathbb Z^{J(x)}$ is the diagonal map, and

\begin{equation}\tag{\SGARefTarget{sga7.sga7i.eix.n12-3-2}{12.3.2}}
 R=\left(\bigoplus_x R(x)\right)/\operatorname{Im}\mathbb Z^J.
\end{equation}

\noindent
Clearly it is enough to take the sum over the finite set of points of $C$
which belong to at least two branches of $C$. Thus $R$ is a finitely
generated $\mathbb Z$-module. There is then a canonical isomorphism

\begin{equation}\tag{\SGARefTarget{sga7.sga7i.eix.n12-3-3}{12.3.3}}
 T\ \bigl(\simeq\text{ the maximal torus of }
       \underline{\operatorname{Pic}}^o_{C/k}\bigr)
 \simeq R\otimes\mathbb G_{m,k}.
\end{equation}

\noindent
In other words, the character group $M$ of $T$ is

\begin{equation}\tag{\SGARefTarget{sga7.sga7i.eix.n12-3-4}{12.3.4}}
 M\mathrel{\mathop{=}^{\mathrm{def}}}D(T)=\check R
 \simeq\ker\left(\bigoplus_x\check R(x)\longrightarrow\mathbb Z^J\right),
\end{equation}

\noindent
where $\check R(x)$ may also be interpreted as the set of
integer-valued functions on $J(x)$ whose values have sum zero.
\footnote{The printed source numbers these two displays (12.2.3) and
(12.2.4). They belong to section \SGARefLink{sga7.sga7i.eix.n12-3}{12.3} and are cited below as (\SGARefLink{sga7.sga7i.eix.n12-3-3}{12.3.3}) and
(\SGARefLink{sga7.sga7i.eix.n12-3-4}{12.3.4}); those coherent numbers are used here.}

\medskip
\noindent
\SGARefTarget{sga7.sga7i.eix.n12-3-5}{12.3.5}. We shall need only the case in which at most two branches of $C$
pass through any point $x$ (as happens, in particular, for ordinary double
points). If $\operatorname{card}J(x)=2$, then $R(x)$ is isomorphic to
$\mathbb Z$, but the isomorphism depends on an ordering of the two branches
$(C',C'')$ through $x$, or equivalently on the choice of one branch $C'$:
to $C'$ one associates the element of $R(x)$ which is the canonical image
of $0_{C'}+1_{C''}$. One then obtains an isomorphism

\begin{equation}\tag{\SGARefTarget{sga7.sga7i.eix.n12-3-6}{12.3.6}}
 R=\mathbb Z^I/\operatorname{Im}\mathbb Z^J,
 \qquad
 T\simeq\mathbb G_m^I/\operatorname{Im}\mathbb G_m^J,
 \qquad
 M=\ker\bigl(\mathbb Z^I\longrightarrow\mathbb Z^J\bigr),
\end{equation}

\noindent
where $I$ is the set of points $x\in C$ which belong to two branches of
$C$, after a choice of signs as just explained.

\medskip
\noindent
\underline{Remarks} \SGARefTarget{sga7.sga7i.eix.n12-3-7}{12.3.7}. A sometimes convenient way to describe the
character group $M$ of the maximal torus $T$ of
$\underline{\operatorname{Pic}}^o_{C/k}$ is the following. Consider the
``unoriented one-dimensional graph'' $\Gamma(C)$ (N. Bourbaki,
\emph{Groupes et alg\`ebres de Lie}, Chapter IV, Appendix), whose vertices
are the irreducible components of $\widetilde C$. Two vertices
$\widetilde C_i$ and $\widetilde C_j$ are joined by a set of edges in
bijective correspondence with the pairs of \underline{distinct} branches
of $C$, contained respectively in $\widetilde C_i$ and
$\widetilde C_j$, which pass through the same point of $C$. This set is
$\widetilde C_i\times_C\widetilde C_j$ when $i\neq j$; when $i=j$, it is

\begin{equation*}
 \bigl(\widetilde C_i\times_C\widetilde C_i-
       \Delta(\widetilde C_i)\bigr)/(\text{symmetry}),
\end{equation*}

\noindent
where $\Delta:\widetilde C\to\widetilde C\times_C\widetilde C$ is the
diagonal. (It is therefore empty if $C_i$ is geometrically unibranch.) One
then has canonical isomorphisms

\begin{equation}\tag{\SGARefTarget{sga7.sga7i.eix.n12-3-7.occ-2}{12.3.7}}
 M=D(T)\xleftarrow{\ \sim\ }H_1(\Gamma(C),\mathbb Z),
 \qquad
 \check M=R\simeq H^1(\Gamma(C),\mathbb Z).
\end{equation}

We shall not need this combinatorial interpretation of $M$, and shall
confine ourselves to defining the canonical homomorphism

\begin{equation}\tag{\SGARefTarget{sga7.sga7i.eix.n12-3-8}{12.3.8}}
 H_1(\Gamma(C),\mathbb Z)\longrightarrow M,
\end{equation}

\noindent
leaving it to the reader to prove that this is an isomorphism. Let $I$ be
the set of points of $C$ which belong to at least two branches of $C$,
let $J$, as above, be the set of irreducible components of $\widetilde C$
(or of $C$, which amounts to the same thing), and let $K$ be the set of
``branches'' of $C$ at the points of $I$, that is, the inverse image of
$I$ in $\widetilde C$. There is a canonical map

\begin{equation}\tag{\SGARefTarget{sga7.sga7i.eix.n12-3-9}{12.3.9}}
 h:K\longrightarrow I\times J,
 \qquad h=(f,g),
\end{equation}

\noindent
and the description (\SGARefLink{sga7.sga7i.eix.n12-3-4}{12.3.4}) of $M$ is expressed in terms of $h$ as

\begin{equation}\tag{\SGARefTarget{sga7.sga7i.eix.n12-3-10}{12.3.10}}
 M\simeq\ker\left(
  \mathbb Z^{(h)}:\mathbb Z^{(K)}\longrightarrow
  \mathbb Z^{(I\times J)}\right).
\end{equation}

Indeed, the expression (\SGARefLink{sga7.sga7i.eix.n12-3-3}{12.3.3}) for $R=\check M$ may plainly be
interpreted as the cokernel of the homomorphism
$\mathbb Z^{I\times J}\to\mathbb Z^K$ induced by $h$, which is the
transpose of $\mathbb Z^{(h)}$. The graph $\Gamma(C)$ is likewise described
in terms of $h$ by

\begin{equation*}
 \Gamma_0=J,
 \qquad
 \Gamma_1=(K\times_I K-\Delta(K))/(\text{symmetry}).
\end{equation*}

Choose an orientation on each edge of $\Gamma(C)$, and denote by
$\Gamma_1$ the set of these oriented edges $(x,C',C'')$. Then

\begin{equation}\tag{\SGARefTarget{sga7.sga7i.eix.n12-3-11}{12.3.11}}
 H_1(\Gamma(C),\mathbb Z)=
 \ker\left(d:\mathbb Z^{(\Gamma_1)}\longrightarrow\mathbb Z^{(J)}\right),
\end{equation}

\noindent
where the differential is

\begin{equation}\tag{\SGARefTarget{sga7.sga7i.eix.n12-3-12}{12.3.12}}
 d(x,C',C'')=g(C'')-g(C').
\end{equation}

The homomorphism (\SGARefLink{sga7.sga7i.eix.n12-3-8}{12.3.8}) is induced by the following homomorphism $v$ of
chain complexes:\footnote{The printed display defining (\SGARefLink{sga7.sga7i.eix.n12-3-8}{12.3.8}) carries
no number even though the following text cites it, and (\SGARefLink{sga7.sga7i.eix.n12-3-13}{12.3.13}) writes
$\mathbb Z^{(\Gamma)}$ where the edge chain group is
$\mathbb Z^{(\Gamma_1)}$. The English text supplies the cited label and
the required subscript. The working transcription's
$\widetilde C_i\times_C\widetilde C_1$ in the loop-edge description is
also restored from the scan as
$\widetilde C_i\times_C\widetilde C_i$.}

\begin{equation}\tag{\SGARefTarget{sga7.sga7i.eix.n12-3-13}{12.3.13}}
 \begin{split}
 v=(v_0,v_1):C_\bullet(\Gamma(C))
   &=\left[\mathbb Z^{(\Gamma_1)}\longrightarrow\mathbb Z^{(J)}\right]\\
   &\longrightarrow
     \left[\mathbb Z^{(K)}\longrightarrow\mathbb Z^{(I\times J)}\right],
 \end{split}
\end{equation}

\noindent
given by

\begin{equation}\tag{\SGARefTarget{sga7.sga7i.eix.n12-3-14}{12.3.14}}
 v_1(x,C',C'')=(x,C'')-(x,C'),
 \qquad
 v_0(\widetilde C_j)=(0,\widetilde C_j).
\end{equation}

% END INCLUDED SGA7I_English_current_buildable_source_20260823/components/173_expose_IX_section_12_3.tex

% BEGIN INCLUDED SGA7I_English_current_buildable_source_20260823/components/174_expose_IX_section_12_4_and_theorem_12_5.tex
\medskip
\noindent
\SGARefTarget{sga7.sga7i.eix.n12-4}{12.4}. We now place ourselves in the case considered at the end of \SGARefLink{sga7.sga7i.eix.n12-2}{12.2}.
By (\SGARefLink{sga7.sga7i.eix.n12-1-12}{12.1.12}) and (\SGARefLink{sga7.sga7i.eix.n12-3-6}{12.3.6}), applied with $C=X_o$, the character group
$M$ of the maximal torus $T_o$ of the special fibre $A_o$ of the
N\'eron model $A$ of
$A=\underline{\operatorname{Pic}}^o_{X_\eta/\eta}$ is given by

\begin{equation}\tag{\SGARefTarget{sga7.sga7i.eix.n12-4-1}{12.4.1}}
 M=\ker\left(\mathbb Z^I\xrightarrow{\ d\ }\mathbb Z^J\right),
\end{equation}

\noindent
where $I$ is the set of singular points of $X_o$, and $J$ is the set of
its irreducible components, the homomorphism
$\mathbb Z^I\to\mathbb Z^J$ being given by

\begin{equation}\tag{\SGARefTarget{sga7.sga7i.eix.n12-4-2}{12.4.2}}
 d(x_i)=C'_i-C''_i.
\end{equation}

Here $C'_i$ and $C''_i$ are the two irreducible components of
$\widetilde C$ corresponding to the branches through the double point
$x_i$. The operator $d$ depends on the choice of an order
$(C'_i,C''_i)$ on these two branches, and the isomorphism (\SGARefLink{sga7.sga7i.eix.n12-4-1}{12.4.1})
plainly depends on this system of choices. Nevertheless, if, for
$i\in I$, we denote by

\begin{equation}\tag{\SGARefTarget{sga7.sga7i.eix.n12-4-3}{12.4.3}}
 \varphi_{x_i}:M\longrightarrow\mathbb Z
\end{equation}

\noindent
the composite

\begin{equation*}
 M\longrightarrow\mathbb Z^I
   \xrightarrow{\ \operatorname{pr}_i\ }\mathbb Z,
\end{equation*}

\noindent
where the first homomorphism comes from (\SGARefLink{sga7.sga7i.eix.n12-4-1}{12.4.1}), then, on returning to
the expression (\SGARefLink{sga7.sga7i.eix.n12-3-4}{12.3.4}), one sees at once that $\varphi_{x_i}$ is
well-defined up to sign. More precisely, it depends only on the order
chosen between the two branches $C'_i,C''_i$ through $x_i$, and changes
sign when that order is reversed. It follows in particular that the
bilinear form

\begin{equation}\tag{\SGARefTarget{sga7.sga7i.eix.n12-4-4}{12.4.4}}
 \varphi_{x_i}\otimes\varphi_{x_i}
 \in \check M\otimes_{\mathbb Z}\check M
 \simeq\operatorname{Hom}(M\otimes M,\mathbb Z)
\end{equation}

\noindent
is determined canonically by the singular point $x_i$, without any sign
ambiguity. We may therefore put

\begin{equation}\tag{\SGARefTarget{sga7.sga7i.eix.n12-4-5}{12.4.5}}
 u=\sum_i\varphi_{x_i}\otimes\varphi_{x_i}:
 M\otimes_{\mathbb Z}M\longrightarrow\mathbb Z.
\end{equation}

This gives a well-defined pairing, independent of every choice of
signs. It is plainly \underline{positive definite}, since it is the
restriction to $M\hookrightarrow\mathbb Z^I$ of the standard quadratic
form $\sum_i X_i^2$ on $\mathbb Z^I$.

The notation $M$ for the character group of $T_o$ is not, a priori,
consistent with that of the preceding sections, where this group was
written $M'$, while $M$ denoted the analogous group associated with the
dual abelian scheme $A'_K$ of $A_K$. But since $A_K$ is a Jacobian, it
has a canonical polarization $\xi$ [5]\footnote{The printed citation
following the canonical polarization is ``[5]'' followed by a blank
pinpoint. The English text preserves the reference without inventing a
missing locator.}, which induces a canonical
isomorphism

\begin{equation}\tag{\SGARefTarget{sga7.sga7i.eix.n12-4-6}{12.4.6}}
 \psi_\xi:A_K\xrightarrow{\ \sim\ }A'_K.
\end{equation}

\noindent
By means of this isomorphism we shall identify $A_K$ and $A'_K$, so that
there is no longer any need to distinguish $M$ from $M'$. In particular,
for every prime number $\ell$, the monodromy pairing (\SGARefLink{sga7.sga7i.eix.n9-1-2}{9.1.2}) may be
regarded as a pairing

\begin{equation}\tag{\SGARefTarget{sga7.sga7i.eix.n12-4-7}{12.4.7}}
 u_\ell:M_\ell\otimes_{\mathbb Z_\ell}M_\ell
 \longrightarrow\mathbb Z_\ell.
\end{equation}

With this understood, we can state the principal result of the present
section:

\medskip
\noindent
\underline{Theorem} \SGARefTarget{sga7.sga7i.eix.n12-5}{12.5} (\underline{Picard--Lefschetz formula}).
\underline{Let $S$ be a strictly local trait, and let $X$ be a regular
scheme, projective and flat over $S$, with one-dimensional fibres, whose
geometric generic fibre is smooth and connected and whose geometric
special fibre has no singularities other than ordinary quadratic
singularities. With the notation introduced in \SGARefLink{sga7.sga7i.eix.n12-4}{12.4}, for every prime
number $\ell$, the monodromy pairing (\SGARefLink{sga7.sga7i.eix.n12-4-7}{12.4.7}) is equal to the pairing
obtained from (\SGARefLink{sga7.sga7i.eix.n12-4-5}{12.4.5}) by extension of scalars from $\mathbb Z$ to
$\mathbb Z_\ell$.}

% END INCLUDED SGA7I_English_current_buildable_source_20260823/components/174_expose_IX_section_12_4_and_theorem_12_5.tex

% BEGIN INCLUDED SGA7I_English_current_buildable_source_20260823/components/175_expose_IX_sections_12_6_and_12_7.tex
\medskip
\noindent
\SGARefTarget{sga7.sga7i.eix.n12-6}{12.6}. Notice that this statement proves \SGARefLink{sga7.sga7i.eix.n10-4}{10.4} in the particular case
under consideration when, in \SGARefLink{sga7.sga7i.eix.n10-4}{10.4} b), one takes for $\xi$ the canonical
polarization considered there. As we showed in \SGARefLink{sga7.sga7i.eix.n10-5}{10.5}, this suffices to
imply \SGARefLink{sga7.sga7i.eix.n10-4}{10.4} in the general case.

\medskip
\noindent
\SGARefTarget{sga7.sga7i.eix.n12-7}{12.7}. \underline{Proof of \SGARefLink{sga7.sga7i.eix.n12-5}{12.5}. The case $\ell\ne p$.} It relies
essentially on the \underline{cohomological Picard--Lefschetz formula}
in the situation considered in \SGARefLink{sga7.sga7i.eix.n12-5}{12.5}, a formula giving the action of the
monodromy group $I$ on the following group:\footnote{The working
transcription adds a dual superscript to $T_\ell(A_{\bar\eta})$ in
(\SGARefLink{sga7.sga7i.eix.n12-7-1}{12.7.1}). A direct 5,000-dpi inspection of the printed formula shows no
such superscript, in agreement with the standard identification
$H^1(A_{\bar\eta},\mathbb Z_\ell)=
\operatorname{Hom}(T_\ell(A_{\bar\eta}),\mathbb Z_\ell)$.}

\begin{equation}\tag{\SGARefTarget{sga7.sga7i.eix.n12-7-1}{12.7.1}}
 H^1(X_{\bar\eta},\mathbb Z_\ell)
 \simeq H^1(A_{\bar\eta},\mathbb Z_\ell)
 \simeq\operatorname{Hom}
   \bigl(T_\ell(A_{\bar\eta}),\mathbb Z_\ell\bigr)
\end{equation}

\noindent
when $\ell\ne p$. This formula, which will be proved in a later expos\'e
(Expos\'e XV), and indeed in a substantially more general case, is
essentially identical to assertion \SGARefLink{sga7.sga7i.eix.n12-5}{12.5} for the prime $\ell$ in
question, in view of the definition in \SGARefLink{sga7.sga7i.eix.n9-2}{9.2} of $u_\ell$ for $\ell\ne p$.
It therefore remains to treat the case $\ell=p$. We shall do this by a
method which reduces it to characteristic zero (and, moreover, to the
case where $X_o$ has only one singular point), which we now explain.

\medskip
\noindent
\underline{Remarks \SGARefTarget{sga7.sga7i.eix.n12-7-2}{12.7.2}.} Notice also that, once the announced
reduction has been made, the usual arguments would allow us to reduce to
the case in which the base field is the field $\mathbb C$ of complex
numbers, and then to invoke the Picard--Lefschetz formula in its classical
form---except that there does not seem to be a satisfactory reference
for that formula. The proof of the Picard--Lefschetz formula given in our
Seminar will proceed similarly, by reduction to the transcendental case,
which is amenable to properly transcendental methods. One may therefore
say that, in any event, the proof of \SGARefLink{sga7.sga7i.eix.n12-5}{12.5} (and hence of \SGARefLink{sga7.sga7i.eix.n10-4}{10.4}) proposed
here is transcendental in nature. Let us also point out that the only
other known proof of \SGARefLink{sga7.sga7i.eix.n10-4}{10.4} comes from rigid-analytic theory [24], and must
for that reason likewise be regarded as essentially ``transcendental.''
It would be interesting to find purely algebraic proofs of these results
(although it is by no means evident that this is possible!).

% END INCLUDED SGA7I_English_current_buildable_source_20260823/components/175_expose_IX_sections_12_6_and_12_7.tex

% BEGIN INCLUDED SGA7I_English_current_buildable_source_20260823/components/176_expose_IX_section_12_8_opening_through_before_12_8_9.tex
\medskip
\noindent
\SGARefTarget{sga7.sga7i.eix.n12-8}{12.8}. \underline{Proof of \SGARefLink{sga7.sga7i.eix.n12-5}{12.5}: reduction to the case in which $S$ has
characteristic zero.} We shall use an argument based on the same principle
as the proof of the ``theorem on the essentially tame action of monodromy
on $\pi_1$'' in Expos\'e V.\footnote{The printed source leaves the internal
pinpoint following ``V'' blank; no section target is supplied.} Notice that
we may assume $S$ complete, by \SGARefLink{sga7.sga7i.eix.n10-3}{10.3}. Let $W$ be a Cohen $p$-ring with
residue field $k$, and choose a homomorphism $W\longrightarrow V$ inducing
the identity on residue fields (EGA $0_{\mathrm{IV}}$, 19.8.6). Let
$(x_\lambda)_{\lambda\in L}$ be the set of singular points of $X_o$, and
consider the formal moduli scheme $\underline S$ of $X_o$ over $W$, together
with the ``universal curve''
\[
 \underline f:\underline X\longrightarrow\underline S
\]
whose special fibre is $X_o$ (Expos\'e VI).\footnote{The printed source
likewise leaves the internal pinpoint following ``VI'' blank.} We saw in
the cited place that $\underline S$ is isomorphic to the formal spectrum of
a power-series ring $W[[T_1,\ldots,T_N]]$, and that in $\underline S$ one
can find a normal-crossings divisor

\begin{equation}\tag{\SGARefTarget{sga7.sga7i.eix.n12-8-1}{12.8.1}}
 D=\sum_{\lambda\in L}D_\lambda,
\end{equation}

\noindent
whose irreducible components $D_\lambda$ are indexed by $L$, each
$D_\lambda$ being described by an equation of the form $T_{i_\lambda}=0$.
(If desired, one may suppose $i_\lambda=\lambda$ after indexing $L$ by a
set of integers $[1,r]$.) The component $D_\lambda$ is the image under
$\underline f:\underline X\longrightarrow\underline S$ of the component
$D'_\lambda$ through the singular point $x_\lambda$ of $X_o$ in the
non-smoothness locus
\[
 D'=\bigcup_{\lambda\in L}D'_\lambda
\]
of $\underline f$. By the construction itself,
$f:X\longrightarrow S$ is isomorphic to the pullback of
$\underline f:\underline X\longrightarrow\underline S$ by a suitable
$W$-morphism $g:S\longrightarrow\underline S$:

\begin{equation}\tag{\SGARefTarget{sga7.sga7i.eix.n12-8-2}{12.8.2}}
\begin{tikzcd}
 \underline X \arrow[d,"\underline f"'] &
 X \arrow[l] \arrow[d,"f"] \\
 \underline S & S \arrow[l,"g"']
\end{tikzcd}
\qquad .
\end{equation}

\noindent
Set

\begin{equation}\tag{\SGARefTarget{sga7.sga7i.eix.n12-8-3}{12.8.3}}
 P=\underline{\operatorname{Pic}}^{o}_{\underline X/\underline S}.
\end{equation}

\noindent
Since the special fibre of $\underline f$ is geometrically reduced of
dimension $1$, it follows that $\underline f$ is cohomologically flat in
dimension zero,\footnote{The working transcription reads $\underline F$
here. Direct inspection of the printed page shows $\underline f$, as also
required by the sentence's reference to the special fibre of the universal
curve.} by EGA III, 7.8.1, and that $P$ is formally smooth over
$\underline S$ (the latter fact follows from
$H^2(X_o,\underline{\mathcal O}_{X_o})=0$, since $X_o$ has dimension $1$).
For simplicity, we shall assume that $P$ \underline{is representable by a
group scheme over $\underline S$}; it will necessarily be a smooth group
scheme with connected fibres. This is a special case of an unpublished
result of Mumford cited in [22] (and considerably generalized by Raynaud in
the note cited there, when the base is a trait). Since no proof of either
Mumford's or Raynaud's result has been published, the reader may prefer to
use the general results of M. Artin [1]. They give the representability of
$P$---under the sole assumption that $\underline f$ is proper, flat, and
cohomologically flat in dimension zero---by an ``algebraic space'' in
Artin's sense, which would suffice for our purpose. Put

\begin{equation}\tag{\SGARefTarget{sga7.sga7i.eix.n12-8-4}{12.8.4}}
 U=\underline S-\operatorname{supp}D.
\end{equation}

\noindent
Thus $U$ is the largest open subset of $\underline S$ over which
$\underline X$ is smooth, and hence
\[
 P|_U=\underline{\operatorname{Pic}}^{o}_{\underline X_U/U}
\]
is an abelian scheme over $U$. Its pullback to $\eta$ by
$g:\eta\longrightarrow U$ is the abelian scheme $A_\eta$ over $\eta$ that
we wish to study. More precisely, there is a canonical isomorphism over
$S$ (\SGARefLink{sga7.sga7i.eix.n12-1-10}{12.1.10}):

\begin{equation}\tag{\SGARefTarget{sga7.sga7i.eix.n12-8-5}{12.8.5}}
 P\times_{\underline S}S\simeq A^{o}.
\end{equation}

\noindent
Since the special fibre $P_o=A_o^{o}$ is $\ell$-divisible for every
$\ell$, the same holds for $P$ itself; all its fibres are therefore
extensions of an abelian variety by a torus. Consider the pro-$\ell$ group
$T_\ell(P)$ over $\underline S$, whose pullback by $g$ is precisely
$T_\ell(A^{o})$. Applying the constructions of \SGARefLink{sga7.sga7i.eix.n7-3}{7.3}, one defines a
filtration of this pro-group by a ``fixed part'' and a ``toric part'':

\begin{equation}\tag{\SGARefTarget{sga7.sga7i.eix.n12-8-6}{12.8.6}}
 T_\ell(P)\supset T_\ell(P)^f=T_\ell(\widehat P)
 \supset T_\ell(P)^t=T_\ell(T)\supset0.
\end{equation}

\noindent
Its pullback by $g$ is the analogous filtration of $T_\ell(A^{o})$
studied in sections 2 and 5.

Consider the restriction of this filtration over the open subset $U$ of
$\underline S$. There we have the canonical self-duality of $P_U$, as the
Jacobian of a smooth curve scheme over $U$, and hence a corresponding
pairing

\begin{equation}\tag{\SGARefTarget{sga7.sga7i.eix.n12-8-7}{12.8.7}}
 \varphi:T_\ell(P_U)\otimes T_\ell(P_U)
 \longrightarrow T_\ell(\mathbb G_{mU}).
\end{equation}

\noindent
This is a \underline{self-duality}, in the sense of Barsotti--Tate groups,
of $T_\ell(P_U)$.\footnote{The printed source reads $T_\ell(B_U)$ here,
although no $B_U$ has been introduced and the pairing and every adjacent
term involve $P_U$. The English text restores the internally forced
$P_U$.} We saw in \SGARefLink{sga7.sga7i.eix.n7-4-8}{7.4.8} that, for this pairing,
$T_\ell(P_U)^f$ \underline{and} $T_\ell(P_U)^t$ \underline{are each
other's orthogonal complements}. It follows that
\[
 E_U=T_\ell(P_U)/T_\ell(P_U)^t
\]
is canonically isomorphic to the Barsotti--Tate pro-group dual to
$T_\ell(P_U)^f=F_U$. It therefore extends canonically to a Barsotti--Tate
group $E$ over $\underline S$, namely

\begin{equation}\tag{\SGARefTarget{sga7.sga7i.eix.n12-8-8}{12.8.8}}
 E=D(F),\qquad\text{where}\qquad F=T_\ell(P)^f.
\end{equation}

\noindent
We are now in a typical situation of mixed extensions (\SGARefLink{sga7.sga7i.eix.n9-3}{9.3}), and can repeat
the constructions of \SGARefLink{sga7.sga7i.eix.n9-4}{9.4}--\SGARefLink{sga7.sga7i.eix.n9-6}{9.6} (for whose having been written in the stale
and piss-poor setting of a base scheme reduced to a trait we apologize!):
$T_\ell(P_U)$ is a mixed extension of $E_U$ by $F_U$, where $E$ and $F$
are Barsotti--Tate groups over $\underline S$ which are extensions, dual to
one another:

% END INCLUDED SGA7I_English_current_buildable_source_20260823/components/176_expose_IX_section_12_8_opening_through_before_12_8_9.tex

% BEGIN INCLUDED SGA7I_English_current_buildable_source_20260823/components/177_expose_IX_section_12_8_9_through_12_10.tex
\begin{equation}\tag{\SGARefTarget{sga7.sga7i.eix.n12-8-9}{12.8.9}}
\begin{tikzcd}[row sep=small,column sep=small]
0 \arrow[r] & P \arrow[rr] \arrow[d,dash,"\wr"]
 & & F \arrow[r] \arrow[d,dash,"\wr"] & R \arrow[r] & 0 & \\
& T_\ell(P)^t\simeq
  \underline{\operatorname{Hom}}
    \bigl(M,T_\ell(\mathbb G_{m\underline S})\bigr)
& & T_\ell(P)^f & & & \\
& & & & & & \\
0 \arrow[r] & R \arrow[rr] \arrow[d,dash,"\wr"]
& & E \arrow[r] & Q \arrow[r] \arrow[d,dash,"\wr"] & 0 & \\
& D(R) & & & M\otimes\mathbb Z_{\ell\,\underline S} & & .
\end{tikzcd}
\end{equation}

\noindent
The printed upper identification in (\SGARefLink{sga7.sga7i.eix.n12-8-9}{12.8.9}) omits the outer closing
parenthesis after
$\underline{\operatorname{Hom}}(M,T_\ell(\mathbb G_{m\underline S}))$;
the balanced expression has been supplied here. Notice that $R$ is
canonically self-dual, as follows from the self-duality (\SGARefLink{sga7.sga7i.eix.n12-8-7}{12.8.7}) and the
orthogonality statement noted above. The only difference from the cited
construction is that the pairing obtained on $M_\ell\times M_\ell$ takes
values in the group
\[
 \bigl(\mathbb G_m(U)/\mathbb G_m(\underline S)\bigr)
 \otimes_{\mathbb Z}\mathbb Z_\ell,
\]
which is no longer canonically isomorphic to $\mathbb Z_\ell$, but rather
to $\mathbb Z_\ell^{\,L}$, by means of the isomorphism

\begin{equation}\tag{\SGARefTarget{sga7.sga7i.eix.n12-8-10}{12.8.10}}
 \mathbb Z^L\simeq
 \mathbb G_m(U)/\mathbb G_m(\underline S).
\end{equation}

\noindent
It maps the element indexed by $\lambda$ in the canonical basis of
$\mathbb Z^L$ to the class of $T_{i_\lambda}$ in the right-hand side,
where $T_{i_\lambda}$ is an equation for the divisor $D_\lambda$. We thus
obtain a canonical pairing

\begin{equation}\tag{\SGARefTarget{sga7.sga7i.eix.n12-8-11}{12.8.11}}
 \underline{u}_\ell:M_\ell\otimes M_\ell
 \longrightarrow\mathbb Z_\ell^{\,L}\qquad .
\end{equation}

\noindent
It is essentially immediate that the monodromy pairing $u_\ell$ (\SGARefLink{sga7.sga7i.eix.n12-4-7}{12.4.7})
over $S$ which we propose to study is the composite of
$\underline{u}_\ell$ with the homomorphism
$\mathbb Z_\ell^{\,L}\longrightarrow\mathbb Z_\ell$ obtained, by
tensoring with $\mathbb Z_\ell$ over $\mathbb Z$, from the homomorphism

\begin{center}
\begin{tikzcd}
 \mathbb G_m(U)/\mathbb G_m(\underline S)
   \arrow[r] \arrow[d,dash,"\wr"'] &
 \mathbb G_m(\eta)/\mathbb G_m(S)
   \arrow[d,dash,"\wr"'] \\
 \mathbb Z^L & \mathbb Z
\end{tikzcd}
\end{center}

\noindent
induced by $g:S\longrightarrow\underline S$ (which sends $\eta$ into
$U$). This homomorphism is given by an element

\begin{equation}\tag{\SGARefTarget{sga7.sga7i.eix.n12-8-12}{12.8.12}}
 (d_\lambda)_{\lambda\in L}\in\mathbb Z^L,\qquad
 d_\lambda\geqslant1\quad\text{for every }\lambda\in L,
\end{equation}

\noindent
where, for every $\lambda\in L$, $d_\lambda$ is the multiplicity of the
inverse image in $S$ of the divisor $D_\lambda$ on $\underline S$. It in
fact depends only on the completion
$\widehat{\underline{\mathcal O}}_{X,x_\lambda}$ as a $V$-algebra
(Expos\'e VI).\footnote{Both printed citations to Expos\'e VI in this
sentence leave their internal section pinpoints blank.} Moreover, it is
known (Expos\'e VI) that, for a given $\lambda$, one has $d_\lambda=1$ if
and only if $X$ is \underline{regular} at $x_\lambda$.\footnote{The
printed source has the lowercase letter $x$ as the subject of ``is
regular.'' The following sentence and the mathematical statement require
the scheme $X$.} Since in \SGARefLink{sga7.sga7i.eix.n12-5}{12.5} we assumed $X$ regular, that is, all the
$d_\lambda$ equal to $1$, assertion \SGARefLink{sga7.sga7i.eix.n12-5}{12.5} will therefore follow from the
following lemma, in which $X,S,g$ no longer occur.

\medskip
\noindent
\underline{Lemma} \SGARefTarget{sga7.sga7i.eix.n12-8-13}{12.8.13}. \underline{If $\lambda\in L$, the composite of
(\SGARefLink{sga7.sga7i.eix.n12-8-11}{12.8.11}) with $\operatorname{pr}_\lambda$ is equal to the form
$\varphi_{x_\lambda}\otimes\varphi_{x_\lambda}$ of (\SGARefLink{sga7.sga7i.eix.n12-4-4}{12.4.4}).}

\medskip
To prove this last assertion, consider the completed strict localization
$\underline S^\lambda$ of $\underline S$ at the generic point of
$D_\lambda$,\footnote{In this reduction the printed source drops the
subscript $\lambda$ from $D_\lambda$ and prints an un-underlined $X$ as
the right-hand factor of the following fibre product. Both marks are
restored: the localization is at the component $D_\lambda$, and the
fibre product is formed from the universal curve $\underline X$ over
$\underline S$.} and put
\[
 \underline X^\lambda=
 \underline X\times_{\underline S}\underline S^\lambda.
\]
The special fibre $\underline X_o^\lambda$ then has only one singular
point $y_\lambda$, which is again an ordinary quadratic singularity. In
the corresponding monodromy pairing

\begin{equation}\tag{\SGARefTarget{sga7.sga7i.eix.n12-8-14}{12.8.14}}
 u_\ell^\lambda:M_\ell^\lambda\otimes M_\ell^\lambda
 \longrightarrow\mathbb Z_\ell\qquad ,
\end{equation}

\noindent
$M^\lambda$ consequently has rank $1$ or $0$. We leave to the reader the
verification of the following easy facts, which make explicit, in the
special case in which we need them, the behaviour under localization of
the preceding constructions.

\medskip
\noindent
a) There is a unique subtorus $T_\lambda$ of
$P\times_{\underline S}D_\lambda$ whose restriction to the generic point
$s_\lambda$ of $D_\lambda$ is the maximal torus of $P_{s_\lambda}$ (which
has dimension $0$ or $1$).

\medskip
\noindent
b) Consider the transpose homomorphism $M\longrightarrow M^\lambda$ on
character groups of the inclusion
$T_{\lambda o}\hookrightarrow T_o$, where $T_o$ is the maximal torus of
$P_o$, and the form
$\varphi_{y_\lambda}:M^\lambda\longrightarrow\mathbb Z$, defined up to
sign as in \SGARefLink{sga7.sga7i.eix.n12-4-3}{12.4.3} (with $X$ replaced by $\underline X^\lambda$). Then,
up to sign, the following diagram is commutative:

\begin{equation}\tag{\SGARefTarget{sga7.sga7i.eix.n12-8-15}{12.8.15}}
\begin{tikzcd}[column sep=huge,row sep=large]
 M \arrow[rr] \arrow[dr,"\varphi_{x_\lambda}"']
 & & M^\lambda \arrow[dl,"\varphi_{y_\lambda}"] \\
 & \mathbb Z & \qquad .
\end{tikzcd}
\end{equation}

\noindent
The working transcription misreads the right diagonal label in (\SGARefLink{sga7.sga7i.eix.n12-8-15}{12.8.15})
as $\varphi_{\lambda_\lambda}$. The printed diagram, the immediately
preceding definition, and the formula below all give
$\varphi_{y_\lambda}$. (In fact, choosing an ordering of the two branches
of $X_o$ through $x_\lambda$ determines an analogous choice at the point
$y_\lambda$ of $\underline X_o^\lambda$, and for choices corresponding in
this way the preceding diagram is commutative.)

\medskip
\noindent
c) The following diagram is commutative:

\begin{center}
\begin{tikzcd}[column sep=huge,row sep=large]
 M_\ell\otimes M_\ell
   \arrow[rr] \arrow[dr,"\operatorname{pr}_\lambda\underline{u}_\ell"']
 & & M_\ell^\lambda\otimes M_\ell^\lambda
   \arrow[dl,"u_\ell^\lambda"] \\
 & \mathbb Z_\ell & \qquad ,
\end{tikzcd}
\end{center}

\noindent
where the notation is that of \SGARefLink{sga7.sga7i.eix.n12-8-13}{12.8.13} and (\SGARefLink{sga7.sga7i.eix.n12-8-14}{12.8.14}). These relations show
at once that, to prove \SGARefLink{sga7.sga7i.eix.n12-8-13}{12.8.13}, it is enough to prove the analogous
relation for $\underline X^\lambda$:
\[
 u_\ell^\lambda=
 (\varphi_{y_\lambda}\otimes\varphi_{y_\lambda})
 \otimes_{\mathbb Z}\mathbb Z_\ell,
\]
which is nothing other than \SGARefLink{sga7.sga7i.eix.n12-5}{12.5} for $\underline X^\lambda$ over
$\underline S^\lambda$. Now $\underline S^\lambda$ is a trait of
\underline{characteristic zero}.\footnote{At this final step the printed
source writes $\underline S$ rather than the localized trait
$\underline S^\lambda$ used throughout the reduction. The English text
restores the mathematically required superscript.} This completes the
announced reduction.

\medskip
\noindent
\underline{Remarks \SGARefTarget{sga7.sga7i.eix.n12-9}{12.9}.} Notice that this reduction simultaneously
reduces us to the case in which $X_o$ has only one singular point, the case
classically considered in the Picard--Lefschetz formula. Notice also that,
in this case, $X_o$ has one or two irreducible components, corresponding
respectively to $\operatorname{rank}M=1$ and $\operatorname{rank}M=0$ (the
latter being the case in which the Jacobian of the generic fibre has good
reduction).

\medskip
\noindent
\SGARefTarget{sga7.sga7i.eix.n12-10}{12.10}. The proof given in \SGARefLink{sga7.sga7i.eix.n12-8}{12.8} yields a result more general than \SGARefLink{sga7.sga7i.eix.n12-5}{12.5},
valid without assuming that $X$ is regular, by introducing the integers
$d_\lambda\geqslant1$ of (\SGARefLink{sga7.sga7i.eix.n12-8-12}{12.8.12}): all the monodromy pairings are then
deduced from the pairing

\begin{equation}\tag{\SGARefTarget{sga7.sga7i.eix.n12-10-1}{12.10.1}}
 u=\sum_{\lambda\in L}d_\lambda\,
   \varphi_{x_\lambda}\otimes\varphi_{x_\lambda}:
   M\otimes M\longrightarrow\mathbb Z\qquad .
\end{equation}

% END INCLUDED SGA7I_English_current_buildable_source_20260823/components/177_expose_IX_section_12_8_9_through_12_10.tex

% BEGIN INCLUDED SGA7I_English_current_buildable_source_20260823/components/178_expose_IX_sections_13_1_and_13_2.tex
\medskip
\noindent
13. \underline{Links with transcendental theory: the complex analytic
case}.

For the general notions concerning analytic spaces, we refer to [2].

\medskip
\noindent
\SGARefTarget{sga7.sga7i.eix.n13-1}{13.1}. \underline{Review of analytic groups over $S$.} Let $S$ be an
analytic space, and let $A$ be an analytic space over $S$, endowed with
the structure of a commutative group over $S$ (that is, $A$ is a
commutative group in the category of analytic spaces over $S$). We assume
that $A$ is \underline{smooth} over $S$, and denote by

\begin{equation}\tag{\SGARefTarget{sga7.sga7i.eix.n13-1-1}{13.1.1}}
 \underline E=\underline{\operatorname{Lie}}(A/S)
\end{equation}

\noindent
the locally free module on $S$ obtained by pulling the relative tangent
sheaf of $A$ over $S$ back along the zero section. We shall admit here the
definition of the exponential homomorphism

\begin{equation}\tag{\SGARefTarget{sga7.sga7i.eix.n13-1-2}{13.1.2}}
 \exp:E\longrightarrow A\qquad ,
\end{equation}

\noindent
where $E$ is the analytic vector bundle on $S$ whose module of
holomorphic sections is $\underline E$. In a neighbourhood of a point
$s\in S$, this homomorphism is first defined near the zero section of $E$
by the familiar methods of the theory of formal groups in characteristic
zero, and is then extended in the evident way, by multiplicativity, to
all of $E$. It is a homomorphism of relative analytic groups over $S$;
it is \underline{\'etale}, and its image is the open subspace $A^o$ of
$A$ which is the union of the identity components of the fibres of $A$
over $S$. If $R$ is the analytic subspace which is the kernel of $\exp$,
we therefore have an exact sequence of relative analytic groups

\begin{equation}\tag{\SGARefTarget{sga7.sga7i.eix.n13-1-3}{13.1.3}}
 0\longrightarrow R\longrightarrow E\longrightarrow A^o
 \longrightarrow0\qquad .
\end{equation}

\noindent
Here the morphism $E\longrightarrow A^o$ is surjective and \'etale, so
$R$ is a relative analytic group \underline{\'etale} over $S$. It is an
analytic subspace of $E$, and is closed in $E$ if and only if $A^o$ is
separated over $S$, that is, if and only if the zero section of $A^o$ is
closed in $A^o$.

\medskip
\noindent
\SGARefTarget{sga7.sga7i.eix.n13-1-4}{13.1.4}. Of course, the exact sequence (\SGARefLink{sga7.sga7i.eix.n13-1-3}{13.1.3}) is
\underline{functorial in $A$}, and commutes with every base change
$S'\longrightarrow S$. In this way one obtains an equivalence of
categories, compatible with base changes, between the category of
relative analytic groups over $S$ which are smooth and have connected
fibres, and the category of pairs $(E,R)$ consisting of a locally free
vector bundle $E$ over $S$ and an analytic subgroup $R$ of $E$ which is
\'etale over $S$.

\medskip
\noindent
\SGARefTarget{sga7.sga7i.eix.n13-1-5}{13.1.5}. When $A$ is an ``\underline{abeloid}'' group over $S$, by which
we mean that it is smooth, proper, and separated over $S$, with connected
fibres---so that in particular $A=A^o$---then $R$ is a locally constant
group over $S$, whose rank at each point $s$ (as a $\mathbb Z$-module) is
equal to $2\operatorname{rank}_s(E)$, and conversely. In this way one
obtains the familiar interpretation of an abeloid space over $S$ in
terms of a ``\underline{local system}'' $R$ over $S$ and a
\underline{complex} analytic vector-bundle structure on the real
analytic bundle $R\otimes_{\mathbb Z}\mathbb R$ defined by $R$.

\medskip
\noindent
\SGARefTarget{sga7.sga7i.eix.n13-2}{13.2}. \underline{Formal completion along a fibre.} We now fix a point
$s\in S$, and shall work locally in a neighbourhood of $s$. For every
integer $n\geqslant0$, let $S_n$ denote the $n$-th infinitesimal
neighbourhood of $s$ in $S$, and, for every analytic space $X$ over $S$,
put likewise $X_n=X\times_S S_n$. As $n$ varies, the system of the $X_n$
then forms an inductive system of analytic spaces over $S$, which we
denote by $\widehat X$ and call the \underline{formal completion of $X$
along $X_s$}. It clearly depends functorially on $X$, and its formation
commutes with finite projective limits; hence $X\longmapsto\widehat X$
takes groups to groups. Notice that $\widehat X$ in fact depends only on
the restriction of $X$ to an arbitrary open neighbourhood $U$ of $s$ in
$S$.

\medskip
\noindent
\underline{Lemma \SGARefTarget{sga7.sga7i.eix.n13-2-1}{13.2.1}.} \underline{The data are those of \SGARefLink{sga7.sga7i.eix.n13-1}{13.1}, and
suppose that the following conditions hold:} a) $A=A^o$;
b) $R$ \underline{is constant in a neighbourhood of $s$}; and
c) $R_s$ \underline{spans the complex vector space $E_s$}.
\underline{Let $A'$ be a second smooth analytic group over $S$, and
consider the map}

\begin{equation*}
 u\longmapsto\widehat u:
 \operatorname{Hom}_{S\text{-}\mathrm{gr}}(A,A')
 \longrightarrow
 \operatorname{Hom}_{\widehat S\text{-}\mathrm{gr}}
   (\widehat A,\widehat {A}')\quad ,
\end{equation*}

\noindent
\underline{where the notation is that of \SGARefLink{sga7.sga7i.eix.n13-2}{13.2}, and the analogous map
obtained by restricting over a variable open neighbourhood $U$ of $s$
and passing to the inductive limit:}

\begin{equation}\tag{\SGARefTarget{sga7.sga7i.eix.n13-2-1-1}{13.2.1.1}}
 u\longmapsto\widehat u:
 \operatorname{Hom}_{(S,s)\text{-}\mathrm{gr}}(A,A')
 \longrightarrow
 \operatorname{Hom}_{\widehat S\text{-}\mathrm{gr}}
   (\widehat A,\widehat {A}')\quad .
\end{equation}

\noindent
(\underline{Here $(S,s)$ denotes the germ of the analytic space $S$ at
$s$.}) \underline{Then the latter map is bijective.}

Indeed, consider the exact sequence

\begin{equation*}
 0\longrightarrow R'\longrightarrow E'\longrightarrow {A'}^o
 \longrightarrow C
\end{equation*}

\noindent
analogous to (\SGARefLink{sga7.sga7i.eix.n13-1-3}{13.1.3}) and defined by $A'$. Since the homomorphisms
$A\longrightarrow A'$ necessarily factor through ${A'}^o$, and the same
holds after the base changes $U\longrightarrow S$ and
$S_n\longrightarrow S$, we may assume that $A'={A'}^o$. By the
dictionary of \SGARefLink{sga7.sga7i.eix.n13-1-4}{13.1.4}, it remains to prove that giving a germ of a vector
bundle homomorphism $E\longrightarrow E'$ which takes $R$ into $R'$ is
equivalent to giving a coherent system of homomorphisms
$E_n\longrightarrow E'_n$ which take $R_n$ into $R'_n$.

If $\underline{\mathcal O}_s$ denotes the local ring of $S$ at $s$,
giving a coherent system of homomorphisms
$E_n\longrightarrow E'_n$ is equivalent to giving a homomorphism of
$\widehat{\underline{\mathcal O}}_s$-modules

\begin{equation*}
 \widehat F:\widehat{\underline E}_s
 \longrightarrow\widehat{\underline E}'_s,
\end{equation*}

\noindent
where $\underline E_s$ and $\underline E'_s$ are the free
$\underline{\mathcal O}_s$-modules obtained at $s$ from the
$\underline{\mathcal O}_s$-modules $\underline E$ and $\underline E'$,
corresponding to $E$ and $E'$. On the other hand, a germ of a vector
bundle homomorphism $E\longrightarrow E'$ is equivalent to giving a
homomorphism of $\underline{\mathcal O}_s$-modules

\begin{equation*}
 F:\underline E_s\longrightarrow\underline E'_s\quad .
\end{equation*}

\noindent
This already shows that (\SGARefLink{sga7.sga7i.eix.n13-2-1-1}{13.2.1.1}) is injective. For surjectivity, it is
enough to prove that if $\widehat F$ is a homomorphism as above such that,
for every $n$, the induced homomorphism
$F_n:E_n\longrightarrow E'_n$ takes $R_n$ into $R'_n$, then
$\widehat F$ comes from an $F$---that is, it takes $\underline E_s$ into
$\underline E'_s$---and the corresponding germ of homomorphisms
$E\longrightarrow E'$ takes $R$ into $R'$.

By hypothesis c), $R_s=R_o$ spans $E_s$; it follows from Nakayama's
lemma that the group $R_o$ (identified either with the group of sections
of $\widehat R$ over $\widehat S$, or with the group of germs of sections
of $R$ over $S$) generates $\widehat E_s$. Thus, to prove that
$\widehat F(\underline E_s)\subset\underline E'_s$, it is enough to prove
that $\widehat F(R_o)\subset\underline E'_s$. But
$\widehat F(R_o)\subset R'_o$, where $R'_o$ is itself interpreted as the
group of germs of sections of $R'\subset E'$ in a neighbourhood of $s$.
This gives the asserted inclusion, and hence the existence of $F$.
Moreover, we see that $F(R_o)\subset R'_o$, where $R_o$ and $R'_o$ are
regarded as groups of germs of sections of $E$ and $E'$, respectively.
By b), $R_o$ generates $R$ in a neighbourhood of $s$; it follows that
$F(R)\subset R'$ in a neighbourhood of $s$, which completes the proof.

\medskip
\noindent
\underline{Corollary \SGARefTarget{sga7.sga7i.eix.n13-2-2}{13.2.2}.} \underline{If $A$ satisfies conditions
a), b), and c) of \SGARefLink{sga7.sga7i.eix.n13-2-1}{13.2.1}, it is determined up to unique isomorphism in a
neighbourhood of $s$ by the knowledge of the $\widehat S$-group
$\widehat A$.}\footnote{The printed source refers here to the conditions
of \SGARefLink{sga7.sga7i.eix.n13-2-2}{13.2.2} itself; the three conditions are introduced in Lemma \SGARefLink{sga7.sga7i.eix.n13-2-1}{13.2.1}.}

% END INCLUDED SGA7I_English_current_buildable_source_20260823/components/178_expose_IX_sections_13_1_and_13_2.tex

% BEGIN INCLUDED SGA7I_English_current_buildable_source_20260823/components/179_expose_IX_section_13_3.tex
\medskip
\noindent
\SGARefTarget{sga7.sga7i.eix.n13-3}{13.3}. \underline{The analytic Raynaud extension}.

We are now in a position to paraphrase, in analytic geometry, the
construction of the Raynaud extension given in 7. Suppose that $A$ is
as in \SGARefLink{sga7.sga7i.eix.n13-1}{13.1}, with $A=A^o$,
\footnote{The printed source has $A=A_o$ here. The required hypothesis,
and the one used below, is $A=A^o$; compare \SGARefLink{sga7.sga7i.eix.n13-1}{13.1} and Lemma \SGARefLink{sga7.sga7i.eix.n13-2-1}{13.2.1}.}
and consider $R_o=R_s$, which is a free $\mathbb Z$-module of finite
type. After restricting $S$ to a sufficiently small open neighbourhood
of $s$, and interpreting $R_o$ as the group of germs of sections of $R$
near $s$, one obtains a homomorphism

\begin{equation}\tag{\SGARefTarget{sga7.sga7i.eix.n13-3-1}{13.3.1}}
 R_{oS}\longrightarrow R\subset E,
\end{equation}

\noindent
where $R_{oS}$ denotes the constant analytic group over $S$ with value
$R_o$. On choosing the neighbourhood of $s$ small enough that every
section of $\underline{\mathcal O}_S$ whose germ at $s$ is zero is itself
zero, one sees that the preceding homomorphism is injective. Its image is
therefore a constant analytic subgroup of $R$; it is plainly the largest
constant subgroup of $R$, and is called the \underline{fixed part} of
$R$ and denoted by $R^f$:

\begin{equation}\tag{\SGARefTarget{sga7.sga7i.eix.n13-3-2}{13.3.2}}
 R_{oS}\simeq R^f\hookrightarrow R,
 \qquad\text{inducing an isomorphism }R_s^f\simeq R_s.
\end{equation}

\noindent
We may then consider the quotient group

\begin{equation}\tag{\SGARefTarget{sga7.sga7i.eix.n13-3-3}{13.3.3}}
 A^\natural=E/R^f,
\end{equation}

\noindent
and, in view of (\SGARefLink{sga7.sga7i.eix.n13-1-3}{13.1.3}), obtain a canonical exact sequence

\begin{equation}\tag{\SGARefTarget{sga7.sga7i.eix.n13-3-4}{13.3.4}}
 0\longrightarrow R/R^f\stackrel{i}{\longrightarrow}A^\natural
 \longrightarrow A\longrightarrow0.
\end{equation}

\noindent
Notice that, after restricting $S$ to a neighbourhood of $s$, $R^f$ is
always closed in $E$, and hence $A^\natural$ is separated over $S$.
Notice also that the homomorphism $A^\natural\longrightarrow A$ is
\'etale, and therefore induces an isomorphism

\begin{equation}\tag{\SGARefTarget{sga7.sga7i.eix.n13-3-5}{13.3.5}}
 \underline{\operatorname{Lie}}(A^\natural)
 \stackrel{\sim}{\longrightarrow}
 \underline{\operatorname{Lie}}(A)=\underline E.
\end{equation}

\noindent
Moreover, for every $n\geqslant0$, the canonical homomorphism
$A^\natural\longrightarrow A$ induces an isomorphism
$A_n^\natural\simeq A_n$ (since $R_n^f\stackrel{\sim}{\longrightarrow}
R_n$), and consequently an isomorphism

\begin{equation}\tag{\SGARefTarget{sga7.sga7i.eix.n13-3-6}{13.3.6}}
 \widehat{A^\natural}\stackrel{\sim}{\longrightarrow}\widehat A.
\end{equation}

\noindent
This last fact shows that $A^\natural$ does indeed play the role of the
group, likewise denoted by $A^\natural$, that we considered in 7. Here
it has been constructed without any hypothesis on the fibre
$A_s=A_o$ of the form ``an extension of an abeloid group by a torus''.
Such a hypothesis will nevertheless be necessary if we wish to recover
on $A^\natural$ the important extension structure (\SGARefLink{sga7.sga7i.eix.n7-2-3}{7.2.3}). By contrast,
the homomorphism $A^\natural\longrightarrow A$ in (\SGARefLink{sga7.sga7i.eix.n13-3-4}{13.3.4}) is essentially
transcendental in nature and has no counterpart in the purely algebraic
theory of \S~7. In \S~14 we shall return to the question of a purely
algebraic interpretation of the homomorphism
$i:R/R^f\longrightarrow A^\natural$ in (\SGARefLink{sga7.sga7i.eix.n13-3-4}{13.3.4}).

Suppose first that we are given a complex subtorus

\begin{equation}\tag{\SGARefTarget{sga7.sga7i.eix.n13-3-7}{13.3.7}}
 T_o\subset A_o,
\end{equation}

\noindent
where by a ``\underline{complex torus}'' we mean an analytic group
isomorphic to a Cartesian power of $\mathbb C^*$. Let

\begin{equation}\tag{\SGARefTarget{sga7.sga7i.eix.n13-3-8}{13.3.8}}
 T=T_o\times S
\end{equation}

\noindent
be the constant torus over $S$ with value $T_o$. Proceeding as in \SGARefLink{sga7.sga7i.eix.n5-1}{5.1},
one obtains a unique homomorphism (indeed, an immersion) of inductive
systems of analytic groups

\begin{equation*}
 \widehat T\longrightarrow\widehat{A^\natural},
\end{equation*}

\noindent
which over $S_o$ reduces to (\SGARefLink{sga7.sga7i.eix.n13-3-7}{13.3.7}). By \SGARefLink{sga7.sga7i.eix.n13-2-2}{13.2.2}, this homomorphism comes
from a unique germ of homomorphisms

\begin{equation}\tag{\SGARefTarget{sga7.sga7i.eix.n13-3-9}{13.3.9}}
 T\longrightarrow A^\natural,
\end{equation}

\noindent
and, after taking $S$ small enough around $s$, we may suppose that it is
defined over all of $S$. The corresponding homomorphism

\begin{equation*}
 \underline{\operatorname{Lie}}(T)\longrightarrow
 \underline{\operatorname{Lie}}(A^\natural)\simeq\underline E
\end{equation*}

\noindent
is injective over each $S_n$, hence over $\widehat S$, and therefore over
$\underline{\mathcal O}_s$ by flat descent. Thus there is a neighbourhood
of $s$ in $S$, which we may take to be all of $S$, on which
$\underline{\operatorname{Lie}}(T)$ identifies with a locally direct
summand $\underline E^t$ of $\underline E$; denote the corresponding
vector subbundle of $E$ by $E^t$:

\begin{equation}\tag{\SGARefTarget{sga7.sga7i.eix.n13-3-10}{13.3.10}}
 \underline{\operatorname{Lie}}(T)\xrightarrow{\ \sim\ }
 \underline E^t\subset\underline E\simeq
 \underline{\operatorname{Lie}}(A^\natural)\simeq
 \underline{\operatorname{Lie}}(A),
 \qquad E^t\subset E.
\end{equation}

\noindent
Likewise, let

\begin{equation}\tag{\SGARefTarget{sga7.sga7i.eix.n13-3-11}{13.3.11}}
 R_o^t\subset R_o,
 \qquad R^t=R_{oS}^t\subset R_{oS}\simeq R^f
\end{equation}

\noindent
denote respectively the sublattice of $R_o$ defined by the subtorus
$T_o$ of $A_o$, and the corresponding constant analytic group over $S$.
Then the homomorphism (\SGARefLink{sga7.sga7i.eix.n13-3-9}{13.3.9}) corresponds simply to the homomorphism

\begin{equation*}
 E^t/R_{oS}^t=E^t/R^t\longrightarrow E/R_{oS}=E/R^f
\end{equation*}

\noindent
induced by the inclusions

\begin{equation*}
 E^t\hookrightarrow E,
 \qquad R^t\hookrightarrow R^f.
\end{equation*}

\noindent
Using the fact that (\SGARefLink{sga7.sga7i.eix.n13-3-9}{13.3.9}) induces an immersion on the fibre over $s$,
we conclude that the homomorphism

\begin{equation*}
 R_o/R_o^t\longrightarrow E_o/E_o^t=(E/E^t)_o
\end{equation*}

\noindent
is injective. It follows at once that

\begin{equation*}
 R_o/R_o^t\longrightarrow(E/E^t)_{s'}
\end{equation*}

\noindent
is still injective for $s'$ near $s$. Thus, after shrinking $S$ around
$s$, we may suppose that (\SGARefLink{sga7.sga7i.eix.n13-3-9}{13.3.9}) \underline{is a closed immersion}.

We are therefore led to introduce

\begin{equation}\tag{\SGARefTarget{sga7.sga7i.eix.n13-3-12}{13.3.12}}
 B=A^\natural/T\simeq(E/E^t)/(R^f/R^t),
\end{equation}

\noindent
so that $A^\natural$ acquires an extension structure

\begin{equation}\tag{\SGARefTarget{sga7.sga7i.eix.n13-3-13}{13.3.13}}
 0\longrightarrow T\longrightarrow A^\natural
 \longrightarrow B\longrightarrow0,
\end{equation}

\noindent
where $B$ is a relative analytic group over $S$, smooth and separated
over $S$, with connected fibres. By construction we therefore have

\begin{equation}\tag{\SGARefTarget{sga7.sga7i.eix.n13-3-14}{13.3.14}}
 \underline{\operatorname{Lie}}(B)\simeq
 \underline E/\underline E^t,
\end{equation}

\noindent
and the exponential exact sequence for $B$ is

\begin{equation}\tag{\SGARefTarget{sga7.sga7i.eix.n13-3-15}{13.3.15}}
 0\longrightarrow R^f/R^t\longrightarrow E/E^t
 \longrightarrow B\longrightarrow0.
\end{equation}

\noindent
Notice that

\begin{equation}\tag{\SGARefTarget{sga7.sga7i.eix.n13-3-16}{13.3.16}}
 R^f/R^t=(R_o/R_o^t)_S
\end{equation}

\noindent
is a constant analytic group over $S$. It follows from \SGARefLink{sga7.sga7i.eix.n13-1-5}{13.1.5} that $B$
is an abeloid group over $S$ in a neighbourhood of $s$ if and only if
its fibre at $s$,

\begin{equation*}
 B_o=A_o/T_o,
\end{equation*}

\noindent
is an abeloid analytic group; in that case $A_o$ appears as an extension
of such a group $B_o$ by a torus $T_o$. \underline{Every way of writing
$A_o$ as an extension of an abeloid group by a complex torus $(*)$
therefore gives rise on $A$, in a neighbourhood of $s$, to a unique
extension structure} (\SGARefLink{sga7.sga7i.eix.n13-3-13}{13.3.13}) \underline{of an abeloid group $B$ by a
torus $T$, inducing the given extension structure at $s$}.

% END INCLUDED SGA7I_English_current_buildable_source_20260823/components/179_expose_IX_section_13_3.tex

% BEGIN INCLUDED SGA7I_English_current_buildable_source_20260823/components/180_expose_IX_section_13_4.tex
\medskip
\noindent
\SGARefTarget{sga7.sga7i.eix.n13-4}{13.4}. \underline{Interpretation in terms of Deligne's ``mixed Hodge
structures''.}

Let us first return to the dictionary of \SGARefLink{sga7.sga7i.eix.n13-1-4}{13.1.4}, which we shall
reformulate in a slightly different way, better suited to introducing
the ``Hodge structures'' relevant to the present situation.  For this
purpose, let $\underline R$ be the sheaf of sections of $R$ over $S$, and
introduce the $\underline{\mathcal O}_S$-module

\begin{equation}\tag{\SGARefTarget{sga7.sga7i.eix.n13-4-1}{13.4.1}}
 \mathcal E(A)=\mathcal E=
 \underline R\otimes_{\mathbb Z}\underline{\mathcal O}_S\qquad .
\end{equation}

\noindent
It is locally free of finite type when $R$ is locally constant.  (Note
that $\mathcal E$ need not be coherent when $R$ is not locally
constant.)  The homomorphism (\SGARefLink{sga7.sga7i.eix.n13-1-3}{13.1.3})
$\underline R\longrightarrow\underline E$ then defines a canonical
homomorphism

\begin{equation}\tag{\SGARefTarget{sga7.sga7i.eix.n13-4-2}{13.4.2}}
 \mathcal E=\underline R\otimes_{\mathbb Z}
 \underline{\mathcal O}_S\longrightarrow\underline E\qquad .
\end{equation}

\noindent
We denote its kernel by $\underline F(A)$, or simply by $\underline F$;
thus there is an exact sequence

\begin{equation}\tag{\SGARefTarget{sga7.sga7i.eix.n13-4-3}{13.4.3}}
 0\longrightarrow\underline F\longrightarrow\mathcal E
 \longrightarrow\underline E\qquad .
\end{equation}

\noindent
The most interesting case is that in which (\SGARefLink{sga7.sga7i.eix.n13-4-2}{13.4.2}) is an epimorphism,
that is, in which $R_s$ spans the complex vector space $E_s$ for every
$s\in S$.  We then have an exact sequence

\begin{equation}\tag{\SGARefTarget{sga7.sga7i.eix.n13-4-4}{13.4.4}}
 0\longrightarrow\underline F\longrightarrow\mathcal E
 \longrightarrow\underline E\longrightarrow0\qquad .
\end{equation}

\noindent
If $R$ is locally constant, this is a short exact sequence of locally
free modules over $S$.  Knowledge of the pair $(E,R)$ is then equivalent
to knowledge of the pair $(R,\underline F)$, where $R$ is a locally
constant analytic group over $S$ whose fibres are free $\mathbb Z$-modules
of finite type, and $\underline F$ is a locally direct-summand submodule
of $\mathcal E=\underline R\otimes_{\mathbb Z}
\underline{\mathcal O}_S$, subject in addition to

\begin{equation}\tag{\SGARefTarget{sga7.sga7i.eix.n13-4-5}{13.4.5}}
 \underline F\cap\underline R=0\qquad .
\end{equation}

\noindent
From the standpoint of Hodge structures, $\mathcal E$ should be regarded
as filtered by the ``\underline{Hodge filtration}''

\begin{equation}\tag{\SGARefTarget{sga7.sga7i.eix.n13-4-6}{13.4.6}}
 \operatorname{Fil}^{1}(\mathcal E)=\mathcal E,
 \qquad \operatorname{Fil}^{0}(\mathcal E)=\underline F,
 \qquad \operatorname{Fil}^{-1}(\mathcal E)=0\qquad .
\end{equation}

\noindent
Notice that the exact sequence (\SGARefLink{sga7.sga7i.eix.n13-4-3}{13.4.3}) still depends functorially on the
analytic group $A$ over $S$, and commutes with arbitrary base change.

Returning now to the conditions of \SGARefLink{sga7.sga7i.eix.n13-3}{13.3}, put

\begin{equation}\tag{\SGARefTarget{sga7.sga7i.eix.n13-4-7}{13.4.7}}
 \mathcal E^f=\underline R^f\otimes_{\mathbb Z}
 \underline{\mathcal O}_S,
 \qquad
 \mathcal E^t=\underline R^t\otimes_{\mathbb Z}
 \underline{\mathcal O}_S\qquad .
\end{equation}

\noindent
These modules are defined on a suitable neighbourhood of $s$, which we
may suppose to be all of $S$, and define a filtration of $\mathcal E$,
called the \underline{weight filtration},

\begin{equation}\tag{\SGARefTarget{sga7.sga7i.eix.n13-4-8}{13.4.8}}
 \mathcal E\supset\mathcal E^f\supset\mathcal E^t\supset0\qquad .
\end{equation}

\noindent
Here $\mathcal E^t$, $\mathcal E^f$, and $\mathcal E$ are to be regarded
as having weights $-2$, $-1$, and $0$, respectively; thus this is an
increasing filtration.  Like $R^t$, the module $\mathcal E^t$ depends on
the choice of a subtorus $T_o$ of $A_o=A_s$.  Since $R^f$ and $R^t$ are
locally constant, \underline{$\mathcal E^f$ and $\mathcal E^t$ are
locally free}; they may also be interpreted as

\begin{equation}\tag{\SGARefTarget{sga7.sga7i.eix.n13-4-9}{13.4.9}}
 \mathcal E^f=\mathcal E(A^\natural),
 \qquad \mathcal E^t=\mathcal E(T)\qquad .
\end{equation}

\noindent
Applying the functoriality of (\SGARefLink{sga7.sga7i.eix.n13-4-3}{13.4.3}), as $A$ varies, to the
homomorphisms

\begin{equation*}
 T\longrightarrow A^\natural\longrightarrow A
\end{equation*}

\noindent
shows at once that $\underline F(A^\natural)$ and $\underline F(T)$
identify respectively with

\begin{equation}\tag{\SGARefTarget{sga7.sga7i.eix.n13-4-10}{13.4.10}}
 \underline F(A^\natural)\simeq
 \underline F^f\stackrel{\mathrm{def}}{=}
 \underline F\cap\mathcal E^f,
 \qquad
 \underline F(T)=\underline F^t=
 \underline F\cap\mathcal E^t=0\qquad .
\end{equation}

\noindent
The last equality follows immediately from the fact that
$\underline F(T)=0$ for a torus $T$: in this case the exact sequence of
type (\SGARefLink{sga7.sga7i.eix.n13-4-3}{13.4.3}) is
$0\longrightarrow0\longrightarrow\mathcal E(T)
\longrightarrow\underline{\operatorname{Lie}}(T)\longrightarrow0$.
At the same time we see that the composite homomorphism

\begin{equation*}
 \mathcal E^t\longrightarrow\mathcal E\longrightarrow\underline E
\end{equation*}

\noindent
induces an isomorphism

\begin{equation}\tag{\SGARefTarget{sga7.sga7i.eix.n13-4-11}{13.4.11}}
 \mathcal E^t\stackrel{\sim}{\longrightarrow}\underline E^t
 \simeq\underline{\operatorname{Lie}}(T)\qquad .
\end{equation}

\noindent
\underline{Suppose now that $R_s$ spans the complex vector space $E_s$},
as happens in particular in the favourable case \SGARefLink{sga7.sga7i.eix.n13-3-16}{13.3.16}, where $A_s$ is
an extension of an abeloid group by a complex torus $T_o$.  After
restricting $S$ to a suitable neighbourhood of $s$, the canonical
homomorphism of type (\SGARefLink{sga7.sga7i.eix.n13-4-2}{13.4.2}), with $A$ replaced by $A^\natural$,

\begin{equation*}
 \mathcal E^f=\mathcal E(A^\natural)
 \longrightarrow\underline{\operatorname{Lie}}(A^\natural)
 \simeq\underline E
\end{equation*}

\noindent
is an epimorphism; a fortiori, so is (\SGARefLink{sga7.sga7i.eix.n13-4-2}{13.4.2}).  We therefore have the
exact sequences (\SGARefLink{sga7.sga7i.eix.n13-4-4}{13.4.4}) and

\begin{equation}\tag{\SGARefTarget{sga7.sga7i.eix.n13-4-12}{13.4.12}}
 0\longrightarrow\underline F^f\longrightarrow\mathcal E^f
 \longrightarrow\underline E\longrightarrow0\qquad .
\end{equation}

\noindent
In particular,

\begin{equation}\tag{\SGARefTarget{sga7.sga7i.eix.n13-4-13}{13.4.13}}
 \underline F+\mathcal E^f=\mathcal E\qquad ,
\end{equation}

\noindent
since $\mathcal E/\underline F\simeq\underline E$.  Moreover, since
$B=A^\natural/T$ is a quotient of $A^\natural$,
\footnote{The printed source says here that $B$ is a quotient of $A$;
the displayed definition $B=A^\natural/T$ gives the intended group.}
it satisfies the analogous surjectivity condition, and hence gives an
exact sequence of type (\SGARefLink{sga7.sga7i.eix.n13-4-4}{13.4.4}):

\begin{equation}\tag{\SGARefTarget{sga7.sga7i.eix.n13-4-14}{13.4.14}}
 0\longrightarrow\underline F(B)\longrightarrow\mathcal E(B)
 \longrightarrow\underline{\operatorname{Lie}}(B)
 \longrightarrow0\qquad .
\end{equation}

\noindent
Its interpretation in terms of the modules already introduced is
immediate.  Indeed, by (\SGARefLink{sga7.sga7i.eix.n13-3-12}{13.3.12}),

\begin{equation}\tag{\SGARefTarget{sga7.sga7i.eix.n13-4-15}{13.4.15}}
 \left\{
 \begin{array}{l}
  \mathcal E(B)\simeq
  (\underline R^f/\underline R^t)\otimes_{\mathbb Z}
  \underline{\mathcal O}_S
  \simeq\mathcal E^f/\mathcal E^t,\\[2ex]
  \underline{\operatorname{Lie}}(B)
  \simeq\underline E/\underline E^t.
 \end{array}
 \right.
\end{equation}

\noindent
It follows from (\SGARefLink{sga7.sga7i.eix.n13-4-11}{13.4.11}) and the relation
$\underline F\cap\mathcal E^t=0$ in (\SGARefLink{sga7.sga7i.eix.n13-4-10}{13.4.10}) that
$\underline F(B)=\ker(\mathcal E^f/\mathcal E^t
\longrightarrow\underline E/\underline E^t)$ is isomorphic to
$\underline F^f$ through the homomorphism

\begin{equation}\tag{\SGARefTarget{sga7.sga7i.eix.n13-4-16}{13.4.16}}
 \underline F^f\simeq\underline F(A^\natural)
 \stackrel{\sim}{\longrightarrow}\underline F(B)\qquad .
\end{equation}

\noindent
Let us summarize the properties just obtained for the weight filtration
(\SGARefLink{sga7.sga7i.eix.n13-4-8}{13.4.8}), and its relation with the Hodge filtration (\SGARefLink{sga7.sga7i.eix.n13-4-6}{13.4.6}) of
$\mathcal E$ defined by the submodule $\underline F$, in the favourable
case \SGARefLink{sga7.sga7i.eix.n13-3-16}{13.3.16} where $A_s$ is an extension of an abeloid analytic group
$B_o$ by a torus $T_o$.  We use this torus to define $R^t$, so that $B$
is an abeloid analytic group over $S$.

\medskip
\noindent
\underline{Proposition \SGARefTarget{sga7.sga7i.eix.n13-4-17}{13.4.17}.} \underline{Under the preceding
conditions, the Hodge and weight filtrations on}
$\mathcal E=\underline R\otimes_{\mathbb Z}\underline{\mathcal O}_S$,

\begin{equation}\tag{\SGARefTarget{sga7.sga7i.eix.n13-4-17-1}{13.4.17.1}}
 \mathcal E\supset\underline F,
 \qquad
 \mathcal E\supset\mathcal E^f\supset\mathcal E^t\supset0,
\end{equation}

\noindent
\underline{satisfy the following conditions:}

\smallskip
\noindent
a) \underline{The weight filtration is obtained, by tensoring with
$\underline{\mathcal O}_S$ over $\mathbb Z$, from the analogous
filtration on $\underline R$}

\begin{equation*}
 \underline R\supset\underline R^f\supset\underline R^t\supset0,
\end{equation*}

\noindent
\underline{where $\underline R^f$ and $\underline R^t$ are locally
constant analytic groups over $S$, with values free $\mathbb Z$-modules
of finite type}, namely
$R_o^f=R_s=\pi_1(A_s)$ and $R_o^t=\pi_1(T_o)$.

\smallskip
\noindent
b) \underline{One has the relations}

\begin{equation}\tag{*}
 \underline F\cap\mathcal E^t=0,
 \qquad \underline F+\mathcal E^f=\mathcal E\qquad .
\end{equation}

\noindent
\underline{In other words, the filtration on $\mathcal E^t$ induced by
the Hodge filtration}, with the degree conventions of (\SGARefLink{sga7.sga7i.eix.n13-4-6}{13.4.6}),
\underline{is concentrated in degree $-1$}, whereas the filtration on
$\mathcal E/\mathcal E^f$ induced by the Hodge filtration is
\underline{concentrated in degree $0$}.

\smallskip
\noindent
c) \underline{Consider on}

\begin{equation*}
 \mathcal E^{\mathrm{ab}}
 \stackrel{\mathrm{def}}{=}\mathcal E^f/\mathcal E^t
 \simeq\underline R^{\mathrm{ab}}\otimes_{\mathbb Z}
 \underline{\mathcal O}_S,
 \qquad
 \underline R^{\mathrm{ab}}=\underline R^f/\underline R^t
 \quad\text{(cf. a))},
\end{equation*}

\noindent
\underline{the filtration induced by the Hodge filtration on
$\mathcal E$}.  It is therefore defined by the submodule
$\underline F^{\mathrm{ab}}$ of $\mathcal E^{\mathrm{ab}}$ which is the
image of

\begin{equation*}
 \underline F^f\stackrel{\mathrm{def}}{=}
 \underline F\cap\mathcal E^f\longrightarrow\mathcal E^{\mathrm{ab}}.
\end{equation*}

\noindent
The arrow is a monomorphism by the first relation $(*)$.
\underline{Then the pair}

\begin{equation*}
 \bigl(\underline R^{\mathrm{ab}},
 \underline F^{\mathrm{ab}}\subset
 \underline R^{\mathrm{ab}}\otimes_{\mathbb Z}
 \underline{\mathcal O}_S=\mathcal E^{\mathrm{ab}}\bigr)
\end{equation*}

\noindent
\underline{is isomorphic to the pair}
$(\underline R(B),\underline F(B))$
\underline{associated with an abeloid group over $S$}, as explained
after (\SGARefLink{sga7.sga7i.eix.n13-4-4}{13.4.4}).  Thus

\begin{equation*}
 B\simeq\mathcal E^{\mathrm{ab}}/
 (F^{\mathrm{ab}}+R^{\mathrm{ab}}),
\end{equation*}

\noindent
where the symbols without underlining denote the corresponding analytic
groups over $S$.

\medskip
\noindent
\underline{Remark \SGARefTarget{sga7.sga7i.eix.n13-4-18}{13.4.18}.} In Deligne's terminology [4], the conclusion
of the preceding proposition, after restriction to an open subset $U$
of $S$ on which $R$ is locally constant, says that $R$, endowed with the
two filtrations (\SGARefLink{sga7.sga7i.eix.n13-4-17-1}{13.4.17.1}) on
$\mathcal E=\underline R\otimes_{\mathbb Z}\underline{\mathcal O}_S$,
is an ``analytic family of mixed Hodge structures over $U$''.  The
situation just described served as the model for Deligne's conjectures
on N\'eron theory in higher dimensions, stated in [9, 9.8--9.13].

There is, however, an important feature of the algebraic theory studied
in \S~7 which is missing from the description in \SGARefLink{sga7.sga7i.eix.n13-4-17}{13.4.17}.  Suppose that
$A$ is an abeloid group over an open subset $U$ complementary to a
nowhere-dense analytic subset $Z$ of $S$---say $Z=\{s\}$ when
$\dim S=1$---so that, by (\SGARefLink{sga7.sga7i.eix.n13-1-5}{13.1.5}), $R|U$ is a local system over $U$.
Then the local system $R/R^t$ over $U$ ought to be \underline{constant};
more precisely, it ought to be dual to the local system ${R'}^f$, where
$R'$ is the analytic group, \`etale over $S$, which would correspond
(as $R$ corresponds to $A$) to a suitable analytic group $A'$ over $S$
such that $A'|U$ is the abeloid group dual to $A|U$.  In the notation of
\S~5, the algebraic analogue is the duality between
$T_\ell(A_K)/T_\ell(A_K)^t$ and $T_\ell(A'_K)^f$ in (\SGARefLink{sga7.sga7i.eix.n5-5-9}{5.5.9}), which
follows from the duality theorem \SGARefLink{sga7.sga7i.eix.n5-2}{5.2}.

I do not know whether such a statement holds under only the hypotheses
of \SGARefLink{sga7.sga7i.eix.n13-4-17}{13.4.17}, or even whether the action of $\pi_1(U,t)$ on $R_t$ is
unipotent for $t\in U$.  Nor do the naive methods used in the present
paragraph seem capable of proving such a result, even under the
additional hypothesis that $A|U$ is endowed with a ``polarization'',
that is, with an invertible sheaf ample on every fibre.  This is,
however, something that the ``semi-simple group people'' apparently know
how to prove, at least when $S$ is nonsingular of dimension $1$; see the
report [9].  In the next subsection we shall show that everything works
under relative algebraicity hypotheses on $A$, and at the same time
relate the result to the construction of \S~7.

% END INCLUDED SGA7I_English_current_buildable_source_20260823/components/180_expose_IX_section_13_4.tex

% BEGIN INCLUDED SGA7I_English_current_buildable_source_20260823/components/181_expose_IX_section_13_5.tex
\medskip
\noindent
\SGARefTarget{sga7.sga7i.eix.n13-5}{13.5}. \underline{The case of a relatively algebraic situation.}
We now suppose that $A$ comes from a relative group scheme over $S$
which is smooth and separated, and whose fibres are extensions of
abelian schemes by tori.  We denote this relative scheme by the same
letter $A$.  (We allow ourselves this abuse of notation because the
functor $G\longmapsto G^{\mathrm{an}}$, from the category of relative
group schemes of the type under consideration to the category of
analytic groups over $S$, is fully faithful, as would be easy to prove
using results of relative GAGA type [10].)  For everything concerning
the notion of a relative scheme and the relations between relative
schemes over $S$ and analytic spaces over $S$, we refer to [10] (see
also [2, Exp. 15]).

If $\underline O$ denotes the local ring of $S$ at $s$, then the
relative scheme $A$ is determined near $s$ by the ordinary scheme
$A_{\underline O}$ over $\operatorname{Spec}(\underline O)$ obtained
from $A$ by localization at $s$.  It is then clear that the formal
completion of $A_{\underline O}$ along its special fibre (isomorphic to
$A_s=A_o$) is isomorphic to the formal completion of the analytic group
$A$ along $A_s$ (\SGARefLink{sga7.sga7i.eix.n13-2}{13.2}), where, by the abuse of notation introduced
above, the $A_n=A\times_S S_n$ are identified with algebraic schemes
over $S$.  In the favourable case in which $\hat A$, viewed as a formal
group scheme over $\hat{\underline O}$, is algebraizable (\SGARefLink{sga7.sga7i.eix.n7-2-1}{7.2.1}), as is
the case in particular when $S$ is normal at $s$ and $A$ is
quasi-projective over $S$ (\SGARefLink{sga7.sga7i.eix.n7-1-5}{7.1.5}), we therefore obtain a group scheme
$A^{\natural}_{\hat{\underline O}}$ over
$\operatorname{Spec}(\hat{\underline O})$, which is an extension of an
abelian scheme $B_{\hat{\underline O}}$ by a torus.  For simplicity, we
suppose that $B_{\hat{\underline O}}$ is projective over
$\hat{\underline O}$ (as is indeed the case, we know, if
$\underline O$, and hence $\hat{\underline O}$, is normal).  We shall
then prove:

\medskip
\noindent
\underline{Proposition \SGARefTarget{sga7.sga7i.eix.n13-5-1}{13.5.1}.} \underline{Under the preceding
conditions, the analytic group $A^{\natural}$ of (\SGARefLink{sga7.sga7i.eix.n13-3-3}{13.3.3}) comes, in a
neighbourhood of $s$, from a relative group scheme determined up to a
unique isomorphism, which is an extension of a projective relative
abelian scheme over $S$ by a torus.  If
$A^{\natural}_{\underline O}$ denotes the group scheme it defines over
$\underline O=\underline O_{S,s}$, then the group scheme}

\begin{equation*}
 A^{\natural}_{\underline O}\otimes_{\underline O}
 \hat{\underline O}
\end{equation*}

\noindent
\underline{is canonically isomorphic, with its extension structure, to
the Raynaud extension $A^{\natural}_{\hat{\underline O}}$ introduced
above.}

The latter assertion follows immediately from the first, in view of
\SGARefLink{sga7.sga7i.eix.n7-2-1}{7.2.1}, since, by (\SGARefLink{sga7.sga7i.eix.n13-3-6}{13.3.6}), the formal completion of
$A^{\natural}_{\underline O}\otimes_{\underline O}\hat{\underline O}$
along its special fibre is $\hat A_{\underline O}$.  To prove the
algebraicity of $A^{\natural}$ near $s$, proceeding as in \SGARefLink{sga7.sga7i.eix.n7-2-1}{7.2.1}, we are
immediately reduced to proving that of the abelian quotient $B$ of
$A^{\natural}$, which follows from the next lemma.

\medskip
\noindent
\underline{Lemma \SGARefTarget{sga7.sga7i.eix.n13-5-2}{13.5.2}.} \underline{Let $B$ be an abeloid analytic
group over the analytic space $S$, and let $s$ be a point of $S$ such
that the completion $\hat B$ along the fibre $B_s$ is algebraizable;
that is, the $B_n=B\times_S S_n$ are algebraizable and $\hat B$, now
viewed as a formal scheme over $\hat{\underline O}_{S,s}$, is
algebraizable, or equivalently comes from an abelian scheme
$B_{\hat{\underline O}}$ over $\hat{\underline O}$.  Suppose moreover
that $B_{\hat{\underline O}}$ is projective over
$\hat{\underline O}$.  Then $B$ is projective over $S$ in a
neighbourhood of $s$.}

The hypothesis means that there is a polarization of the formal
analytic space $\hat B$, that is, a sequence of polarizations $p_n$ of
the $B_n$ which ``glue together''.  The conclusion means that there is
a polarization $p$ of $B$ over a neighbourhood of $s$.  Let $P$ be the
analytic space over $S$ of relative polarizations of $B$ over $S$,
which represents the functor

\begin{equation*}
 S'\longmapsto
 \{\text{polarizations of $B_{S'}$ over $S'$}\}
\end{equation*}

\noindent
on the category of analytic spaces $S'$ over $S$.  This is an open
subfunctor of the N\'eron--Severi functor $\underline{\mathrm{NS}}_{B/S}$
over $S$.  The simplest standard arguments, which we omit here, show
that $\underline{\mathrm{NS}}_{B/S}$, and hence also $P$, is represented
by an \underline{unramified} analytic space over $S$.  The hypothesis
then says that there is a \underline{formal} section $p$ of $P$ over
$S$ at $s$.  Since $P$ is unramified and therefore locally quasi-finite
over $S$, it follows from [2, Exp. 18, no.~1] that every formal section
of $P$ over $S$ at $s$ comes from an analytic section near $s$, that is,
from a polarization of $B$ over a neighbourhood of $s$.  This proves
the assertion.

\medskip
\noindent
\underline{Remarks \SGARefTarget{sga7.sga7i.eix.n13-5-3}{13.5.3}.}

\smallskip
\noindent
a) If in \SGARefLink{sga7.sga7i.eix.n13-5-2}{13.5.2} we omit the hypothesis that
$B_{\hat{\underline O}}$ is projective over $\hat{\underline O}$,
\footnote{The printed source says ``projective over $S$'' here; the
hypothesis of Lemma \SGARefLink{sga7.sga7i.eix.n13-5-2}{13.5.2} and the ensuing normality argument require
``projective over $\hat{\underline O}$''.}
we can still conclude that $B$ is algebraizable near $s$.  For this one
uses an unpublished result (due independently to M.~Artin and
P.~Deligne, using the theory of the sheaves $\underline T^i$ of VI 2)
which asserts that, if $X$ is proper over $S$ and $X_{\mathrm{red}}$ is
algebraizable, then $X$ is algebraizable.  This reduces us to the case
where $S$ is reduced.  On the other hand, when $S$ is normal, so that
$\hat{\underline O}$ is normal, every abelian scheme over
$\hat{\underline O}$ is known to be projective over
$\hat{\underline O}$ (see, for example, [23]), and 13.5.2 applies.  To
obtain the case of an arbitrary reduced $S$, one uses non-flat descent
from the normalization $S'$ of $S$, by means of the following
(delicate) result, which we shall not prove here: a proper morphism of
noetherian schemes $f:S'\longrightarrow S$ which is an epimorphism
(that is, such that
$\underline O_S\longrightarrow f_*(\underline O_{S'})$ is injective)
is an \underline{effective descent morphism} for the fibred category of
relative abelian schemes.

\smallskip
\noindent
b) It appears that the proof of \SGARefLink{sga7.sga7i.eix.n13-5-2}{13.5.2} should also prove the analogous
result for every analytic space $B$ proper over $S$.  Probably more
delicate is the question whether an analytic space $X$ proper over $S$
whose formal completion $\hat X$ along $X_s$ is algebraizable is itself
algebraizable over $S$ near $s$.  We merely point out that here the
notion of relative algebraizability of $X$ over $S$ should be taken in
a broader sense than in [10]: near $s$, $X$ should be required to come
from an \underline{algebraic space} over $\underline O_{S,s}$ in the
sense of M.~Artin [1].

\medskip
\noindent
\SGARefTarget{sga7.sga7i.eix.n13-5-4}{13.5.4}. Let us place ourselves under the conditions of \SGARefLink{sga7.sga7i.eix.n13-5-1}{13.5.1}.  For
every integer $n>0$, choose also the standard isomorphism of groups over
$\mathbb C$

\begin{equation*}
 \mu_{n\mathbb C}\simeq(\mathbb Z/n\mathbb Z)_{\mathbb C},
 \qquad m\longmapsto\exp(2\pi i m/n)\qquad .
\end{equation*}

\noindent
By passing to the inverse limit over the powers of a fixed prime
$\ell$, this gives a canonical isomorphism of $\ell$-adic sheaves over
$\mathbb C$,

\begin{equation*}
 \mathbb Z_\ell(1)_{\mathbb C}\simeq\mathbb Z_{\ell\mathbb C}
 \qquad ,
\end{equation*}

\noindent
which we shall use to identify the two sides.

Applying the exact sequence of
$\underline{\operatorname{Ext}}^i_{\mathbb Z}
(\mathbb Z/n\mathbb Z,-)$ to the exact sequence (\SGARefLink{sga7.sga7i.eix.n13-1-3}{13.1.3}), we obtain a
canonical isomorphism

\begin{equation}\tag{\SGARefTarget{sga7.sga7i.eix.n13-5-5}{13.5.5}}
 {}_nA^{\mathrm o}\simeq
 \underline R\otimes_{\mathbb Z}\mathbb Z/n\mathbb Z
 \qquad ,
\end{equation}

\noindent
which, as $n$ runs through the powers of a given prime $\ell$, gives an
isomorphism of inverse systems of groups over $S$,

\begin{equation}\tag{\SGARefTarget{sga7.sga7i.eix.n13-5-6}{13.5.6}}
 T_\ell(A^{\mathrm o})\simeq
 \underline R\otimes_{\mathbb Z}\mathbb Z_\ell
 \qquad ,
\end{equation}

\noindent
where $\mathbb Z_\ell$ is regarded as the inverse system of the
$\mathbb Z/\ell^{n+1}\mathbb Z$.  In the situation considered here, in
which $A$ comes from a separated group scheme over $S$, it is then
immediately clear that the
$\underline R\otimes\mathbb Z/n\mathbb Z$ may be regarded as relative
schemes.  This gives a meaning to the symbol
$(\underline R\otimes_{\mathbb Z}\mathbb Z_\ell)_{\underline O}$
(whereas $R_{\underline O}$ would in general have none!), and yields a
canonical isomorphism

\begin{equation}\tag{\SGARefTarget{sga7.sga7i.eix.n13-5-7}{13.5.7}}
 T_\ell(A_{\underline O})\simeq
 (\underline R\otimes_{\mathbb Z}\mathbb Z_\ell)_{\underline O}
 \qquad .
\end{equation}

\noindent
It follows immediately from the constructions of paragraph 2 and \SGARefLink{sga7.sga7i.eix.n13-3}{13.3}
that the preceding isomorphism carries the fixed part to the fixed part
and the toric part to the toric part:

\begin{equation}\tag{\SGARefTarget{sga7.sga7i.eix.n13-5-8}{13.5.8}}
 T_\ell(A_{\underline O})^f\simeq
 (\underline R^f\otimes_{\mathbb Z}\mathbb Z_\ell)_{\underline O},
 \qquad
 T_\ell(A_{\underline O})^t\simeq
 (\underline R^t\otimes_{\mathbb Z}\mathbb Z_\ell)_{\underline O}
 \qquad .
\end{equation}

\noindent
If $A'$ is a relative analytic group over $S$ satisfying the same
hypotheses as $A$, and if $W_{\underline O}$ is a biextension of
$(A_{\underline O},A'_{\underline O})$ by
$\mathbb G_{m\underline O}$, hence defines over a neighbourhood of $s$
a biextension of $(A,A')$ by $\mathbb G_{mS}$ (now using the notion in
the complex-analytic sense, obtained by working in the site of analytic
spaces over $S$ for a suitable topology), then, by VIII \SGARefLink{sga7.sga7i.eviii.n2-2}{2.2}, $W$
defines pairings

\begin{equation}\tag{\SGARefTarget{sga7.sga7i.eix.n13-5-9}{13.5.9}}
 (\underline R\otimes\mathbb Z_\ell)\times
 (\underline R'\otimes_{\mathbb Z}\mathbb Z_\ell)
 \simeq T_\ell(A)\times T_\ell(A')
 \longrightarrow\mathbb Z_\ell(1)_S
 \simeq\mathbb Z_{\ell S}\qquad ,
\end{equation}

\noindent
and it is again clear that the corresponding pairing

\begin{equation}\tag{\SGARefTarget{sga7.sga7i.eix.n13-5-10}{13.5.10}}
 T_\ell(A_{\underline O})\times T_\ell(A'_{\underline O})
 \longrightarrow\mathbb Z_\ell(1)_S
\end{equation}

\noindent
is precisely the one associated, by loc. cit., with the biextension
$W_{\underline O}$ of $(A_{\underline O},A'_{\underline O})$.  It is
also not difficult to see that, after restricting to a neighbourhood
of $s$, the pairings (\SGARefLink{sga7.sga7i.eix.n13-5-9}{13.5.9}) arise from a unique pairing (no longer
depending on $\ell$)

\begin{equation}\tag{\SGARefTarget{sga7.sga7i.eix.n13-5-11}{13.5.11}}
 \underline R\otimes\underline R'\longrightarrow\mathbb Z
 \qquad .
\end{equation}

\noindent
With these preliminaries, the algebraic results concerning the
$T_\ell(A_{\underline O})$ can be used to draw conclusions about the
sheaf $\underline R$.  Thus the pairing (\SGARefLink{sga7.sga7i.eix.n13-5-11}{13.5.11}) is zero on
$\underline R^t\times\underline R'{}^f$ and on
$\underline R^f\times\underline R'{}^t$ (\SGARefLink{sga7.sga7i.eix.n7-4-8}{7.4.8}).  Moreover, if $S$ is
nonsingular of dimension $1$ and, over $U=S-\{s\}$, $A$ and $A'$ are
relative abelian schemes dual to one another through $W$, then
$\underline R^t_U$ and $\underline R'{}^f_U$, as well as
$\underline R^f_U$ and $\underline R'{}^t_U$, are exact annihilators of
one another over a sufficiently small neighbourhood of $s$ (\SGARefLink{sga7.sga7i.eix.n7-4-8}{7.4.8}).
Since (\SGARefLink{sga7.sga7i.eix.n13-5-11}{13.5.11}) is a perfect duality over $U$, it follows that there is
a duality over $U$ between $\underline R/\underline R^t$ and
$\underline R'{}^f$.  In particular,
$\underline R/\underline R^t$ is constant, and a fortiori the action of
$\pi_1(U,t)$ on $\underline R_t$ is unipotent of level $2$.  Similarly,
over $U$ there is a perfect duality between
$\underline R/\underline R^f$ and $\underline R'{}^t$.  Since the
latter is also isomorphic to the dual torus of the lattice $M_S$ over
$S$, where $M$ is the group defined in (\SGARefLink{sga7.sga7i.eix.n5-5-4}{5.5.4}), we obtain a canonical
isomorphism

\begin{equation}\tag{\SGARefTarget{sga7.sga7i.eix.n13-5-12}{13.5.12}}
 \underline R_U/\underline R^f_U\simeq M_U
 \qquad ,
\end{equation}

\noindent
and the homomorphism $i$ of (\SGARefLink{sga7.sga7i.eix.n13-3-4}{13.3.4}) may be written as a homomorphism

\begin{equation}\tag{\SGARefTarget{sga7.sga7i.eix.n13-5-13}{13.5.13}}
 i:M_U\longrightarrow A^\natural
 \qquad ,
\end{equation}

\noindent
which we shall encounter again in a slightly different context in the
next paragraph.

% END INCLUDED SGA7I_English_current_buildable_source_20260823/components/181_expose_IX_section_13_5.tex

% BEGIN INCLUDED SGA7I_English_current_buildable_source_20260823/components/182_expose_IX_section_14_1.tex
\medskip
\noindent
14. \underline{Links with transcendental theory: the rigid-analytic
case.  The canonical homomorphism}
$\underline M_K\longrightarrow A_K^{\natural\mathrm o}$.

\medskip
\noindent
\SGARefTarget{sga7.sga7i.eix.n14-1}{14.1}. We assume here the foundations of the theory of rigid-analytic
spaces, for which one may consult [12], [13], and [34].

Let $S$ be a complete trait, and let $A_K$ be an abelian scheme over its
field of fractions.  In \SGARefLink{sga7.sga7i.eix.n8-1}{8.1}, by descent from the Raynaud extension in
the semistable case, we constructed an extension
$G_K=A_K^{\natural\mathrm o}$ (\SGARefLink{sga7.sga7i.eix.n8-1-1}{8.1.1})

\begin{equation}\tag{\SGARefTarget{sga7.sga7i.eix.n14-1-1}{14.1.1}}
 0\longrightarrow T_K\longrightarrow G_K\longrightarrow B_K
 \longrightarrow0
\end{equation}

\noindent
of an abelian scheme $B_K$ by a torus $T_K$.  Considering the dual
abelian scheme $A'_K$, we likewise have an extension
$G'_K=A'{}^{\natural\mathrm o}_K$

\begin{equation}\tag{\SGARefTarget{sga7.sga7i.eix.n14-1-2}{14.1.2}}
 0\longrightarrow T'_K\longrightarrow G'_K\longrightarrow B'_K
 \longrightarrow0\qquad ,
\end{equation}

\noindent
where $B'_K$ is moreover identified with the abelian scheme dual to
$B_K$ (\SGARefLink{sga7.sga7i.eix.n8-2}{8.2}).  As in (\SGARefLink{sga7.sga7i.eix.n8-1-8}{8.1.8}), let $\underline M_K$ and
$\underline M'_K$ denote the lattices over $K$ dual to $T'_K$ and
$T_K$, respectively:

\begin{equation}\tag{\SGARefTarget{sga7.sga7i.eix.n14-1-3}{14.1.3}}
 \underline M_K=D(T'_K),\qquad \underline M'_K=D(T_K)\qquad .
\end{equation}

\noindent
Thus the extension structures (\SGARefLink{sga7.sga7i.eix.n14-1-2}{14.1.2}) and (\SGARefLink{sga7.sga7i.eix.n14-1-1}{14.1.1}) may respectively be
interpreted (VIII \SGARefLink{sga7.sga7i.eviii.n3-7}{3.7}) as equivalent to the data of homomorphisms

\begin{equation}\tag{\SGARefTarget{sga7.sga7i.eix.n14-1-4}{14.1.4}}
 \underline M_K\xrightarrow{\ p_K\ }B_K,\qquad
 \underline M'_K\xrightarrow{\ p'_K\ }B'_K\qquad .
\end{equation}

Let $W_K^{\natural}$ be the biextension of $(B_K,B'_K)$ by
$\mathbb G_{mK}$ which defines the canonical duality between these
abelian schemes (\SGARefLink{sga7.sga7i.eix.n8-2}{8.2}), and let $E_K$ be the biextension of
$(\underline M_K,\underline M'_K)$ by $\mathbb G_{mK}$ obtained by
pulling $W_K^{\natural}$ back along the homomorphisms (\SGARefLink{sga7.sga7i.eix.n14-1-4}{14.1.4}).  One
checks immediately that trivializations of $E_K$ correspond
canonically and bijectively to homomorphisms

\begin{equation}\tag{\SGARefTarget{sga7.sga7i.eix.n14-1-5}{14.1.5}}
 i:\underline M_K\longrightarrow G_K
\end{equation}

\noindent
lifting $p_K:\underline M_K\longrightarrow B_K$, or, symmetrically, to
homomorphisms

\begin{equation}\tag{\SGARefTarget{sga7.sga7i.eix.n14-1-6}{14.1.6}}
 i':\underline M'_K\longrightarrow G'_K
\end{equation}

\noindent
lifting $p'_K:\underline M'_K\longrightarrow B'_K$.

With this understood, write a superscript ${}^{\mathrm{an}}$ for the
rigid-analytic groups over $K$ associated with the algebraic groups
under consideration [13].  In [24] one defines a canonical exact
sequence

\begin{equation}\tag{\SGARefTarget{sga7.sga7i.eix.n14-1-7}{14.1.7}}
 0\longrightarrow\underline M_K^{\mathrm{an}}
 \xrightarrow{\ i^{\mathrm{an}}\ }G_K^{\mathrm{an}}
 \xrightarrow{\ e\ }A_K^{\mathrm{an}}\longrightarrow0\qquad .
\end{equation}

\noindent
Here $i^{\mathrm{an}}$, which may also be interpreted as a homomorphism
$i:\underline M_K\longrightarrow G_K$ of ordinary group schemes,
\underline{lifts the homomorphism} $p$ in (\SGARefLink{sga7.sga7i.eix.n14-1-4}{14.1.4}), that is, it is of
the form considered in (\SGARefLink{sga7.sga7i.eix.n14-1-5}{14.1.5}).  This exact sequence is the
rigid-analytic analogue of the exact sequence (\SGARefLink{sga7.sga7i.eix.n13-3-4}{13.3.4}).  It depends
functorially on $A_K$, and its formation is compatible with every
change of traits $S'\longrightarrow S$.

The latter fact also shows that, in order to define and study (\SGARefLink{sga7.sga7i.eix.n14-1-7}{14.1.7}),
finite descent reduces us to the case where $A_K$ has semistable
reduction.  In this case the group schemes
$T_K,G_K,B_K,\underline M_K,\ldots$ are the generic fibres of group
schemes $T,G,B,M,\ldots$ over $S$, with $G$ an extension of the abelian
scheme $B$ by the torus $T$ (paragraph 7).  We may then consider the
formal completion $\hat G$ of $G$ along its special fibre, and by
construction there is a canonical isomorphism (\SGARefLink{sga7.sga7i.eix.n7-2-6}{7.2.6})

\begin{equation}\tag{\SGARefTarget{sga7.sga7i.eix.n14-1-8}{14.1.8}}
 \hat G\xrightarrow{\ \sim\ }\widehat{A^{\mathrm o}}
 \qquad .
\end{equation}

\noindent
Passing to the ``generic fibres'', which are rigid-analytic groups over
$K$, gives a canonical isomorphism

\begin{equation}\tag{\SGARefTarget{sga7.sga7i.eix.n14-1-9}{14.1.9}}
 \hat G^{\mathrm{an}}\xrightarrow{\ \sim\ }
 \widehat{A^{\mathrm o}}\,{}^{\mathrm{an}}
 \qquad .
\end{equation}

\noindent
On the other hand, there are canonical homomorphisms (rigid-analytic
open immersions, but not in general isomorphisms)

\begin{equation}\tag{\SGARefTarget{sga7.sga7i.eix.n14-1-10}{14.1.10}}
 \hat G^{\mathrm{an}}\lhook\joinrel\longrightarrow G_K^{\mathrm{an}},
 \qquad
 \widehat{A^{\mathrm o}}\,{}^{\mathrm{an}}
 \lhook\joinrel\longrightarrow A_K^{\mathrm{an}}
 \qquad .
\end{equation}

\noindent
The homomorphism $e$ in (\SGARefLink{sga7.sga7i.eix.n14-1-7}{14.1.7}) is uniquely characterized by the
commutativity of the diagram

\begin{equation}\tag{\SGARefTarget{sga7.sga7i.eix.n14-1-11}{14.1.11}}
\begin{tikzcd}
 \hat G^{\mathrm{an}} \arrow[r] \arrow[d]
   & \widehat{A^{\mathrm o}}\,{}^{\mathrm{an}} \arrow[d] \\
 G_K^{\mathrm{an}} \arrow[r] & A_K^{\mathrm{an}}
\end{tikzcd}
\end{equation}

\noindent
where the unnamed arrows are those of (\SGARefLink{sga7.sga7i.eix.n14-1-9}{14.1.9}) and (\SGARefLink{sga7.sga7i.eix.n14-1-10}{14.1.10}).

\medskip
\noindent
\underline{Remarks \SGARefTarget{sga7.sga7i.eix.n14-1-12}{14.1.12}.}

\smallskip
\noindent
a) Notice that the uniqueness of $e$ making the preceding diagram
commute is immediate, since the rigid-analytic space
$G_K^{\mathrm{an}}$ is connected.  The existence of $e$, on the other
hand, is by no means trivial, and on this point we merely refer to the
cited note of M.~Raynaud.  Moreover, the definition of
$i^{\mathrm{an}}$ given in loc. cit. is likewise transcendental in
nature, although it may be interpreted as an object of an essentially
algebraic nature (\SGARefLink{sga7.sga7i.eix.n14-1-5}{14.1.5}), once the construction of $G$ itself, and
hence of $G_K$, has been obtained through the constructions of
\underline{formal geometry} in paragraph 7.  I know of no purely
algebro-geometric description of it, but below (\SGARefLink{sga7.sga7i.eix.n14-4}{14.4}) I shall propose
an ``arithmetic'' characterization, obtained by reducing to a situation
of finite type over the ring $\mathbb Z$.  Thus, from the standpoint of
algebraic geometry, the homomorphism (\SGARefLink{sga7.sga7i.eix.n14-1-5}{14.1.5}) seems to be a more hidden
piece of structure than those encountered in the preceding paragraphs.

\smallskip
\noindent
b) As explained above, the homomorphism $i$ is canonically associated
with a well-determined trivialization of the biextension $E_K$ of
$(\underline M_K,\underline M'_K)$, or equivalently with a
well-determined homomorphism (\SGARefLink{sga7.sga7i.eix.n14-1-6}{14.1.6}) lifting $p'_K$.  One then proves
that the latter is precisely the homomorphism occurring in the exact
sequence

\begin{equation}\tag{\SGARefTarget{sga7.sga7i.eix.n14-1-13}{14.1.13}}
 0\longrightarrow\underline M'{}^{\mathrm{an}}_K
 \xrightarrow{\ i'{}^{\mathrm{an}}\ }G'{}^{\mathrm{an}}_K
 \longrightarrow A'{}^{\mathrm{an}}_K\longrightarrow0
\end{equation}

\noindent
analogous to (\SGARefLink{sga7.sga7i.eix.n14-1-7}{14.1.7}), associated with the abelian scheme $A'_K$ dual
to $A_K$ (unless it is the opposite???).

% END INCLUDED SGA7I_English_current_buildable_source_20260823/components/182_expose_IX_section_14_1.tex

% BEGIN INCLUDED SGA7I_English_current_buildable_source_20260823/components/183_expose_IX_section_14_2.tex
\par\medskip
\noindent
\SGARefTarget{sga7.sga7i.eix.n14-2}{14.2}. \underline{Relations of the homomorphism}
$i:\underline M_K\longrightarrow G_K$ \underline{with monodromy.}

Let $n$ be an integer $>0$.  From the exact sequence (\SGARefLink{sga7.sga7i.eix.n14-1-7}{14.1.7}), by
means of the exact sequence of the
$\underline{\operatorname{Ext}}^i(\mathbb Z/n\mathbb Z,-)$ on the
site of rigid-analytic spaces over $K$, one then obtains a canonical
exact sequence (taking account of the fact that
$\underline M_K^{\mathrm{an}}$ is torsion-free and
$G_K^{\mathrm{an}}$ is $n$-divisible):

\begin{equation*}
 0\longrightarrow{}_nG_K^{\mathrm{an}}
 \longrightarrow{}_nA_K^{\mathrm{an}}
 \longrightarrow\underline M_{K,n}^{\mathrm{an}}
 \longrightarrow0\qquad ,
\end{equation*}

\noindent
where
$\underline M_{K,n}\overset{\mathrm{dfn}}{=}
\underline M_K/n\underline M_K$.  Since all the terms occurring in
this exact sequence come from finite group schemes over $K$, it may be
rewritten as arising from an exact sequence

\begin{equation}\tag{\SGARefTarget{sga7.sga7i.eix.n14-2-1}{14.2.1}}
 0\longrightarrow{}_nG_K\longrightarrow{}_nA_K
 \longrightarrow\underline M_{K,n}\longrightarrow0\qquad .
\end{equation}

\noindent
As $n$ varies, these exact sequences behave in the expected way, and,
on passing to the inverse limit over the powers of a prime number
$\ell$, one obtains the analogous exact sequence

\begin{equation}\tag{\SGARefTarget{sga7.sga7i.eix.n14-2-2}{14.2.2}}
 0\longrightarrow T_\ell(G_K)\longrightarrow T_\ell(A_K)
 \longrightarrow\underline M_K\otimes_{\mathbb Z}\mathbb Z_\ell
 \longrightarrow0\qquad .
\end{equation}

\noindent
But one already has such an exact sequence by purely algebraic means,
thanks to the first formula (\SGARefLink{sga7.sga7i.eix.n8-1-2}{8.1.2}) and to (\SGARefLink{sga7.sga7i.eix.n8-1-7}{8.1.7}), which give

\begin{equation}\tag{\SGARefTarget{sga7.sga7i.eix.n14-2-3}{14.2.3}}
 T_\ell(G_K)\xrightarrow{\ \sim\ }T_\ell(A_K)^{\mathrm{ef}}
 \subset T_\ell(A_K),\qquad
 T_\ell(A_K)/T_\ell(A_K)^{\mathrm{ef}}
 \simeq\underline M_K\otimes_{\mathbb Z}\mathbb Z_\ell\qquad .
\end{equation}

\noindent
With this understood, I claim that \underline{the homomorphisms}
(\SGARefLink{sga7.sga7i.eix.n14-2-3}{14.2.3}) \underline{arising from the algebraic theory are precisely
those deduced from the exact sequence} (\SGARefLink{sga7.sga7i.eix.n14-2-2}{14.2.2}) \underline{arising
from the transcendental theory}.  For the first homomorphism, which
corresponds to the first homomorphism of (\SGARefLink{sga7.sga7i.eix.n14-2-2}{14.2.2}), our assertion follows
at once, after reduction to the semistable situation, from the
commutativity of (\SGARefLink{sga7.sga7i.eix.n14-1-11}{14.1.11}), since the vertical arrows there induce
isomorphisms on the $T_\ell$, while applying $T_\ell$ to the horizontal
arrows gives, respectively, the arrow under consideration in (\SGARefLink{sga7.sga7i.eix.n14-2-3}{14.2.3})
and that in (\SGARefLink{sga7.sga7i.eix.n14-2-2}{14.2.2}).  We shall admit the asserted compatibility for
the two isomorphisms obtained between
$T_\ell(A_K)/T_\ell(A_K)^{\mathrm{ef}}$ and
$\underline M_K\otimes_{\mathbb Z}\mathbb Z_\ell$.  Notice, moreover,
that \underline{this compatibility} (for a given $\ell$)
\underline{already suffices to characterize uniquely, by
transcendental means, the homomorphism} (\SGARefLink{sga7.sga7i.eix.n14-1-5}{14.1.5}) \underline{occurring
in} (\SGARefLink{sga7.sga7i.eix.n14-1-7}{14.1.7}), interpreted as an isomorphism from
$\underline M_K^{\mathrm{an}}$ onto $\operatorname{Ker}e$.

One may also say that the homomorphism $e$ of (\SGARefLink{sga7.sga7i.eix.n14-1-7}{14.1.7}) identifies the
$T_\ell(A_K)/T_\ell(A_K)^{\mathrm{ef}}$ with the $\ell$-adic
completions $\underline R_K\otimes\mathbb Z_\ell$ of a single lattice
$\underline R_K=\operatorname{Ker}e$ over $K$, and that this
``integral structure'' on the preceding system of pro-$\ell$ groups is
compatible, in the evident sense, with the integral structure arising
from formulas (\SGARefLink{sga7.sga7i.eix.n8-1-7}{8.1.7}) (orthogonality theorem \SGARefLink{sga7.sga7i.eix.n5-2}{5.2}).  This then defines
canonically an isomorphism

\begin{equation*}
 i:\underline M_K\xrightarrow{\ \sim\ }\operatorname{Ker}e
 \qquad .
\end{equation*}

Since the homomorphism $i$ (\SGARefLink{sga7.sga7i.eix.n14-1-5}{14.1.5}) enables one to reconstruct the
usual extension structure on $T_\ell(A_K)$, namely (\SGARefLink{sga7.sga7i.eix.n14-2-2}{14.2.2}), and hence
also its structure as a \underline{blended extension}, described in
\SGARefLink{sga7.sga7i.eix.n9-6}{9.6}, by means of the exact sequence deduced from (\SGARefLink{sga7.sga7i.eix.n14-1-1}{14.1.1}),

\begin{equation*}
\begin{array}{ccccccccc}
0&\longrightarrow&T_\ell(T_K)&\longrightarrow&T_\ell(G_K)
 &\longrightarrow&T_\ell(B_K)&\longrightarrow&0\\[5ex]
 &&\underline M'_K\otimes\mathbb Z_\ell(1)&&&&&&
\end{array}
\end{equation*}

\noindent
it also enables one, in the semistable case, to reconstruct the
\underline{monodromy pairings} (\SGARefLink{sga7.sga7i.eix.n9-1-2}{9.1.2})

\begin{equation}\tag{\SGARefTarget{sga7.sga7i.eix.n14-2-4}{14.2.4}}
 u_\ell:\underline M\otimes\underline M'
 \longrightarrow\mathbb Z_\ell\qquad ,
\end{equation}

\noindent
these having been defined precisely in terms of blended extensions in
\SGARefLink{sga7.sga7i.eix.n9-6}{9.6}.  Let us merely indicate the result of translating the construction
of \SGARefLink{sga7.sga7i.eix.n9-6}{9.6}.  For this purpose, observe that in the semistable case $p_K$
and $p'_K$ arise from homomorphisms

\begin{equation*}
 p:\underline M\longrightarrow B,\qquad
 p':\underline M'\longrightarrow B'\qquad ,
\end{equation*}

\noindent
where $B,B'$ are abelian schemes over $S$, dual to one another.  Pulling
back by $(p,p')$ therefore gives a biextension $E$ of
$(\underline M,\underline M')$ by $\mathbb G_{mS}$ whose generic fibre
is the biextension $E_K$ already considered in \SGARefLink{sga7.sga7i.eix.n4-1}{4.1}.

With this said, suppose for simplicity that $\underline M$ and
$\underline M'$ are constant, with values $M$ and $M'$.  For every
pair $(m,m')\in M\times M'$, $E_{m,m'}$ is a torsor under
$\mathbb G_{mS}$, and hence corresponds to an invertible module
$L_{m,m'}$ over the ring $V$ of $S$.  It therefore admits a basis
$u_0$, determined up to multiplication by a unit of $V$.  Consequently
one obtains a canonical ``valuation'' map

\begin{equation*}
 \underline v:E_{m,m'}(K)\longrightarrow\mathbb Z
\end{equation*}

\noindent
which associates with every basis $u$ of
$L_{m,m'}\otimes_V K$, written $u=\lambda u_0$ ($\lambda\in K^*$),
the element

\begin{equation*}
 \underline v(u)=\underline v(\lambda)\in\mathbb Z\qquad ,
\end{equation*}

\noindent
where $\underline v$ on the right denotes the valuation associated
with $V$ (normalized in the usual way, so that its value group is
$\mathbb Z$); the resulting integer is plainly independent of the
choice of $u_0$.  If $\varphi$ is now a trivialization of $E$, one may
consider the corresponding map

\begin{equation}\tag{\SGARefTarget{sga7.sga7i.eix.n14-2-5}{14.2.5}}
 (m,m')\longmapsto\underline v\bigl(\varphi(m,m')\bigr):
 M\times M'\longrightarrow\mathbb Z\qquad .
\end{equation}

\noindent
With this understood, one finds that \underline{the monodromy
homomorphism} (\SGARefLink{sga7.sga7i.eix.n14-2-4}{14.2.4}) \underline{takes its values in} $\mathbb Z$
\underline{and is precisely the homomorphism} (\SGARefLink{sga7.sga7i.eix.n14-2-5}{14.2.5}), \underline{in
which} $\varphi$ \underline{is the splitting of} $E_K$
\underline{associated with the homomorphism} $i$ (\SGARefLink{sga7.sga7i.eix.n14-1-5}{14.1.5}).

\medskip
\noindent
\underline{Remarks \SGARefTarget{sga7.sga7i.eix.n14-2-6}{14.2.6}.}  In particular, one recovers (by means of
the transcendental rigid-analytic theory) the fact that the monodromy
pairing is integral and independent of $\ell$ (\SGARefLink{sga7.sga7i.eix.n10-4}{10.4} a)), proved by a
different (and likewise transcendental) method in \SGARefLink{sga7.sga7i.eix.n12-5}{12.5}.  According to
M.~Raynaud, rigid-analytic theory also recovers the positivity assertion
\SGARefLink{sga7.sga7i.eix.n10-4}{10.4} b) (which is likewise contained in \SGARefLink{sga7.sga7i.eix.n12-5}{12.5}).

\medskip
\noindent
\underline{Remarks \SGARefTarget{sga7.sga7i.eix.n14-2-7}{14.2.7}.}  Regard
$\underline M_K\otimes_{\mathbb Z}\mathbb Q$ as an extension of
$\underline M_K\otimes_{\mathbb Z}\mathbb Q/\mathbb Z$ by
$\underline M_K$.  Using the homomorphism (\SGARefLink{sga7.sga7i.eix.n14-1-5}{14.1.5})
$i:\underline M_K\longrightarrow G_K$, one deduces an extension
$A_K^b$

\begin{equation}\tag{\SGARefTarget{sga7.sga7i.eix.n14-2-8}{14.2.8}}
 0\longrightarrow G_K\longrightarrow A_K^b
 \longrightarrow\underline M_K\otimes\mathbb Q/\mathbb Z
 \longrightarrow0\qquad .
\end{equation}

\noindent
Observe moreover that, since $G_K$ is a divisible group, the datum of
such an extension is equivalent to that of a family of extensions of
the form (\SGARefLink{sga7.sga7i.eix.n14-2-2}{14.2.2}), namely those deduced from (\SGARefLink{sga7.sga7i.eix.n14-2-8}{14.2.8}) by applying the
$T_\ell$.  Since $i$ is a monomorphism, the extensions in question
which define (\SGARefLink{sga7.sga7i.eix.n14-2-8}{14.2.8}) are precisely the extensions (\SGARefLink{sga7.sga7i.eix.n14-2-2}{14.2.2}) arising
from the exact sequence (\SGARefLink{sga7.sga7i.eix.n14-2-1}{14.2.1}).  As these latter extensions are
defined by purely algebraic means, the same is true of the extension
(\SGARefLink{sga7.sga7i.eix.n14-2-8}{14.2.8}).  In fact, the characterization just given shows at once that
this extension is canonically isomorphic to the group, likewise denoted
$A_K^b$, introduced in (8.2.3), in view of the isomorphism \SGARefLink{sga7.sga7i.eix.n11-11}{11.11} from
the group of connected components of the latter to
$\underline M_K\otimes\mathbb Q/\mathbb Z$.

% END INCLUDED SGA7I_English_current_buildable_source_20260823/components/183_expose_IX_section_14_2.tex

% BEGIN INCLUDED SGA7I_English_current_buildable_source_20260823/components/184_expose_IX_section_14_3.tex
\par\medskip
\noindent
\SGARefTarget{sga7.sga7i.eix.n14-3}{14.3}. \underline{Application to an arithmetic characterization of}
$i:\underline M_K\longrightarrow G_K$.

The characterization we have in mind rests on the observation that,
for every prime number $\ell$, the extension (\SGARefLink{sga7.sga7i.eix.n14-2-2}{14.2.2}) of
$\underline M_K\otimes\mathbb Z_\ell$ by $T_\ell(G_K)$ deduced from
$i$ is known by purely algebraic means.  In particular, the image of
$i$ under the canonical homomorphism

\begin{equation}\tag{\SGARefTarget{sga7.sga7i.eix.n14-3-1}{14.3.1}}
 \operatorname{Hom}(\underline M_K,G_K)\longrightarrow
 \prod_\ell\operatorname{Ext}^1
 \bigl(\underline M_K\otimes\mathbb Z_\ell,T_\ell(G_K)\bigr)
 \qquad .
\end{equation}

\noindent
is known.  In the case (to which one may always reduce by descent) in
which $\underline M_K$ is constant, with value $M$, the preceding
homomorphism may be interpreted as

\begin{equation}\tag{\SGARefTarget{sga7.sga7i.eix.n14-3-2}{14.3.2}}
 \operatorname{Hom}\bigl(M,G_K(K)\bigr)\longrightarrow
 \prod_\ell\operatorname{Hom}
 \bigl(M,H^1(K,T_\ell(G_K))\bigr),
\end{equation}

\noindent
deduced from the canonical homomorphism

\begin{equation}\tag{\SGARefTarget{sga7.sga7i.eix.n14-3-3}{14.3.3}}
 G_K(K)\longrightarrow\prod_\ell H^1(K,T_\ell(G_K))
 \qquad ,
\end{equation}

\noindent
which comes from the connecting homomorphisms in the Kummer exact
sequences

\begin{equation*}
 0\longrightarrow{}_nG_K\longrightarrow G_K
 \xrightarrow{\ n\ }G_K\longrightarrow0\qquad .
\end{equation*}

(The homomorphism (\SGARefLink{sga7.sga7i.eix.n14-3-1}{14.3.1}) could in fact always be interpreted as the
connecting homomorphism analogous to (\SGARefLink{sga7.sga7i.eix.n14-3-3}{14.3.3}) for the twisted group
$\check{\underline M}_K\otimes G_K=
\underline{\operatorname{Hom}}(\underline M_K,G_K)$.)  The kernel of
(\SGARefLink{sga7.sga7i.eix.n14-3-3}{14.3.3}) plainly consists of the elements of $G_K(K)$ which are
\underline{divisible by} $n$ \underline{for every integer}
$n\geq1$, or, as we shall say, which are \underline{infinitely
divisible}.  Thus this already determines $i$ modulo an element of
$\operatorname{Hom}(M,g)$, where $g$ is the subgroup of $G_K(K)$
formed by its infinitely divisible elements:

\begin{equation*}
 g=\bigcap_{n\geq1}nG_K(K)\qquad .
\end{equation*}

We now have the following result.

\medskip
\noindent
\underline{Lemma \SGARefTarget{sga7.sga7i.eix.n14-3-4}{14.3.4}.}  \underline{Let} $V$ \underline{be a
discrete valuation ring whose residue field} $k$ \underline{has
characteristic} $p>0$ \underline{and is of finite type over its prime
field, and let} $G_K$ \underline{be a group scheme of finite type over}
$K$, \underline{of unipotent rank zero if} $\operatorname{char}K=0$.
\underline{Then the group} $g$ \underline{of infinitely divisible
elements of} $G_K(K)$ \underline{is zero; that is,} (\underline{if}
$G_K$ \underline{is divisible}) \underline{the homomorphism} (\SGARefLink{sga7.sga7i.eix.n14-3-3}{14.3.3})
\underline{is injective}.

After making a finite extension $K'/K$ and replacing $V$ by its
normalization in $K'$, we may suppose that $G_K$ admits a composition
series whose successive quotients are of one of the following types:
a constant finite group, an abelian variety, $\mathbb G_{mK}$, or
$\mathbb G_{aK}$.  We are therefore immediately reduced to the case
where $G_K$ itself is of one of these types.  The case of a constant
finite group, or of $\mathbb G_{aK}$, is trivial.  In the abelian case,
$G_K$ may be regarded as the generic fibre of its N\'{e}ron model $G$,
so we are reduced to showing that, if $G$ is a scheme of finite type
over $S$, the subgroup of $G(S)$ formed by the infinitely divisible
points is zero.

This is true of the group $G(k)$, as follows at once from the usual
d\'{e}vissage of $G_o$ and from the \underline{Mordell--Weil theorem}
[16] (which says that $G_o(k)$ is a finitely generated $\mathbb Z$-module
if $G_o$ is an abelian scheme over the field of finite type $k$).
Consequently $g$ is contained in the kernel $G(S)^{(1)}$ of
$G(S)\longrightarrow G(k)$.  The kernel of
$G(S)^{(1)}\longrightarrow G(V/\underline m^{,n+1})$ has a filtration
whose quotients are isomorphic to vector groups over $k$ (after
supposing, as we may, that $S$ is complete), and these are killed by
$p$.  It follows that $g$ is contained in the kernel of the preceding
homomorphism for every $n$, and hence that it is zero.  The case
$G_K=\mathbb G_{mK}$ reduces to the preceding one, since
$G_K(K)=K^*$ is an extension of $\mathbb Z$ (which has no nonzero
infinitely divisible elements) by
$V^*=\mathbb G_m(S)$, and the latter satisfies the same condition by
what we have just proved.

\medskip
\noindent
\underline{Remarks \SGARefTarget{sga7.sga7i.eix.n14-3-4-1}{14.3.4.1}.}  Notice that in \SGARefLink{sga7.sga7i.eix.n14-3-4}{14.3.4} it is essential
to require divisibility by \underline{all} integers $\geq1$, and not
only by the powers of a fixed prime $\ell$; equivalently, one must
consider the kernel of (\SGARefLink{sga7.sga7i.eix.n14-3-3}{14.3.3}), rather than restricting the right-hand
side to a single factor $H^1(K,T_\ell(G_K))$.  Indeed, if $G_K$ comes
from a smooth group scheme $G$ over $S$ and $S$ is henselian, every
element of the kernel of $G(S)\longrightarrow G(k)$ is infinitely
$\ell$-divisible for every $\ell\neq p$!  The argument proving \SGARefLink{sga7.sga7i.eix.n14-3-4}{14.3.4}
shows that the subgroup of $G_K(K)$ formed by the infinitely
$p$-divisible elements (that is, the kernel of
$G_K(K)\longrightarrow H^1(K,T_p(G_K))$ when $G_K$ is $p$-divisible)
is a \underline{finite} group.

\medskip
\noindent
\SGARefTarget{sga7.sga7i.eix.n14-3-5}{14.3.5}. The preceding considerations show that, when the residue
field $k$ has characteristic $p>0$ and is of finite type in the
absolute sense, the homomorphism $i$ (\SGARefLink{sga7.sga7i.eix.n14-1-5}{14.1.5}) is completely determined
by its image under (\SGARefLink{sga7.sga7i.eix.n14-3-1}{14.3.1}), that is, by the extensions (\SGARefLink{sga7.sga7i.eix.n14-2-2}{14.2.2}).  To
deduce from this a principle characterizing $i$ in the general case,
with no restriction on $k$, observe that in paragraph 7 we were able to
construct the scheme $G$ over $S$ under conditions considerably more
general than those in which $S$ is a complete trait and one has an
abelian variety with semistable reduction over its field of fractions.

More generally, let $S$ denote the spectrum of a normal noetherian ring
$V$, separated and complete for an $\underline m$-adic topology, and
consider a smooth quasi-projective group scheme $A$ over $S$, with
connected fibres.  Suppose that
$A_o=A\times_S S_o$, where
$S_o=\operatorname{Spec}(V/\underline m)$, is an extension of an
abelian scheme $B_o$ by an isotrivial torus $T_o$, and that $A_U$ is an
abelian scheme, where $U\subset S-S_o$ is a nonempty open subset.  In
7.22 we then constructed a scheme $G$ over $S$, an extension of an
abelian scheme $B$ by a torus $T$, characterized up to unique
isomorphism by the isomorphism (\SGARefLink{sga7.sga7i.eix.n14-1-8}{14.1.8})

\begin{equation}\tag{\SGARefTarget{sga7.sga7i.eix.n14-3-5-1}{14.3.5.1}}
 \hat G\xrightarrow{\ \sim\ }\hat A
\end{equation}

\noindent
on the formal completions along $G_o$ and $A_o$.  For every prime
number $\ell$ one deduces an isomorphism (\SGARefLink{sga7.sga7i.eix.n7-3-5}{7.3.5})

\begin{equation}\tag{\SGARefTarget{sga7.sga7i.eix.n14-3-5-2}{14.3.5.2}}
 T_\ell(G)\xrightarrow{\ \sim\ }T_\ell(A)^f\qquad ,
\end{equation}

\noindent
and hence, on restricting to $U$, the structure of an extension of
pro-$\ell$ Barsotti--Tate groups

\begin{equation}\tag{\SGARefTarget{sga7.sga7i.eix.n14-3-5-3}{14.3.5.3}}
\begin{tikzcd}[row sep=large,column sep=large]
0 \arrow[r] & T_\ell(G_U) \arrow[r] \arrow[d,equal]
 & T_\ell(A_U) \arrow[r]
 & T_\ell(A_U)/T_\ell(A_U)^f \arrow[r] & 0\quad .\\
& T_\ell(A_U)^f
\end{tikzcd}
\end{equation}

Suppose now that a second group scheme $A'$ over $S$ is given,
satisfying the same hypotheses as $A$, together with a biextension $W$
of $(A,A')$ by $\mathbb G_m$ which over $U$ induces a duality between
the abelian schemes $A_U$ and $A'_U$.  (In fact, by VIII \SGARefLink{sga7.sga7i.eviii.n7-1}{7.1}, giving
$W$ is equivalent to giving such a duality between $A_U$ and $A'_U$.)
Then $W$ defines pairings

\begin{equation}\tag{\SGARefTarget{sga7.sga7i.eix.n14-3-5-4}{14.3.5.4}}
 T_\ell(A_U)\times T_\ell(A'_U)\longrightarrow\mathbb Z_\ell(1)
 \qquad ,
\end{equation}

\noindent
which give rise to the orthogonality theorem (\SGARefLink{sga7.sga7i.eix.n7-4-8}{7.4.8}):
$T_\ell(A_U)^f$ is the orthogonal of
$T_\ell(A'_U)^t=T_\ell(T'_U)$, where $T'$ is the toric part of the
Raynaud extension $G'={A'}^{\natural}$ associated with $A'$.  Thus one
again obtains an isomorphism

\begin{equation}\tag{\SGARefTarget{sga7.sga7i.eix.n14-3-5-5}{14.3.5.5}}
 T_\ell(A_U)/T_\ell(A_U)^f
 \simeq\underline M_U\otimes\mathbb Z_\ell\qquad ,
\end{equation}

\noindent
where $\underline M=D(T')$ is the lattice over $S$ dual to the torus
$T'$.  Consequently (\SGARefLink{sga7.sga7i.eix.n14-3-5-3}{14.3.5.3}) and (\SGARefLink{sga7.sga7i.eix.n14-3-5-5}{14.3.5.5}) give an extension

\begin{equation}\tag{\SGARefTarget{sga7.sga7i.eix.n14-3-5-6}{14.3.5.6}}
 0\longrightarrow T_\ell(G_U)\longrightarrow T_\ell(A_U)
 \longrightarrow\underline M_U\otimes\mathbb Z_\ell
 \longrightarrow0\qquad .
\end{equation}

With this understood, we may state the following conjecture.

% END INCLUDED SGA7I_English_current_buildable_source_20260823/components/184_expose_IX_section_14_3.tex

% BEGIN INCLUDED SGA7I_English_current_buildable_source_20260823/components/185_expose_IX_sections_14_4_and_14_5.tex
\par\medskip
\noindent
\underline{Conjecture \SGARefTarget{sga7.sga7i.eix.n14-4}{14.4}.}  The hypotheses being those just set out in
\SGARefLink{sga7.sga7i.eix.n14-3-5}{14.3.5}, there exists a homomorphism

\begin{equation}\tag{\SGARefTarget{sga7.sga7i.eix.n14-4-1}{14.4.1}}
 i:\underline M_U\longrightarrow G_U
\end{equation}

\noindent
such that, for every prime number $\ell$, the extension (\SGARefLink{sga7.sga7i.eix.n14-3-5-6}{14.3.5.6}) is
isomorphic to the extension deduced from $i$ as explained for (\SGARefLink{sga7.sga7i.eix.n14-2-2}{14.2.2}).
More precisely, there is one and only one way of associating with every
situation of the type considered in \SGARefLink{sga7.sga7i.eix.n14-3-5}{14.3.5} a homomorphism (\SGARefLink{sga7.sga7i.eix.n14-4-1}{14.4.1})
satisfying the preceding condition, in such a way that the formation of
$i$ is compatible with every base change $S'\longrightarrow S$
associated with a continuous homomorphism $V\longrightarrow V'$ of
noetherian adic rings.

\medskip
\noindent
\SGARefTarget{sga7.sga7i.eix.n14-4-2}{14.4.2}. It seems that the existence of $i$, compatibly with base change,
follows from M. Raynaud's rigid-analytic theory, in view of the fact that
he is developing it in a framework more general than that stated in
\SGARefLink{sga7.sga7i.eix.n14-1}{14.1}, working, as here, over general noetherian adic rings.  Moreover,
to prove the existence and uniqueness asserted in \SGARefLink{sga7.sga7i.eix.n14-4}{14.4}, one is
immediately reduced to considering only those $S$ for which $S_o$ is a
scheme of finite type over $\mathbb Z$.  Indeed, $V$ is the filtered
inductive limit of subrings $V_i$ of finite type over $\mathbb Z$; these
may be supposed normal by Nagata's theorem (EGA IV 7.8), and, for $i$
sufficiently large, the situation $(A,A',W)$ comes from an analogous
situation over $V_i$, or alternatively over the separated completion of
$V_i$ for the topology defined by $m_i=\underline m V_i$.  This gives
the asserted reduction.

\medskip
\noindent
\SGARefTarget{sga7.sga7i.eix.n14-4-3}{14.4.3}. Observe now that $S-S_o$ is a Jacobson scheme (EGA IV 10.5.8).
Thus, to see that two solutions $i,i'$ of the problem, for a given
situation $(A,A',W)$, are equal, it is enough to prove that these
homomorphisms are equal at every \underline{closed} point $\eta$ of
$U$.  Let $S'$ be the normalization of the closure of $\eta$ in $S$.
Then $S'$ is a complete trait, and the uniqueness result of \SGARefLink{sga7.sga7i.eix.n4-3}{4.3} may be
applied to the pullback to $S'$ of the situation over $S$, since the
residue field of $S'$ is \underline{finite}.  Once uniqueness is thus
established, it implies compatibility of the formation of $i$ with base
change, subject to existence.

\medskip
\noindent
\underline{Remarks \SGARefTarget{sga7.sga7i.eix.n14-5}{14.5}.}  It would remain to formulate with the desired
generality, and to establish, a compatibility between the homomorphism
$i$ introduced in \SGARefLink{sga7.sga7i.eix.n14-4}{14.4} by a rigid-analytic method and the analogous
homomorphism (\SGARefLink{sga7.sga7i.eix.n13-5-13}{13.5.13}) of complex-analytic theory---or, better still, to
find a common generalization of them!


% END INCLUDED SGA7I_English_current_buildable_source_20260823/components/185_expose_IX_sections_14_4_and_14_5.tex

% BEGIN INCLUDED SGA7I_English_current_buildable_source_20260823/components/186_expose_IX_bibliography.tex
\bigskip
\noindent\underline{BIBLIOGRAPHY}

\begin{enumerate}
\item[{[1]}] M. Artin, \emph{Algebraization of formal moduli I}
  (forthcoming).
\item[{[2]}] H. Cartan, \emph{Familles d'espaces complexes et fondements
  de la g\'{e}om\'{e}trie analytique}, S\'{e}minaire ENS 1960/61.
\item[{[3]}] P. Deligne, D. Mumford, \emph{The irreducibility of the space
  of curves of given genus}, Publ. Math. no.~36 (1969).
\item[{[4]}] P. Deligne, work in preparation on Hodge theory (forthcoming
  in Publ. Math., 1970).
\item[{[5]}] M. Demazure, J. Giraud, M. Raynaud, \emph{Sch\'{e}mas
  ab\'{e}liens}, S\'{e}minaire Orsay 1967/68, forthcoming.
\item[{[6]}] A. Grothendieck, \emph{Le Groupe de Brauer III}, in
  \emph{Dix Expos\'{e}s sur la cohomologie des sch\'{e}mas}, North-Holland
  Publishing Co.
\item[{[7]}] A. Grothendieck, \emph{Crystals and the De Rham cohomology of
  schemes} (notes by I. Coates and O. Jussila), in \emph{Dix expos\'{e}s
  sur la cohomologie des sch\'{e}mas}, North-Holland Publishing Co.
\item[{[8]}] A. Grothendieck, \emph{Un th\'{e}or\`{e}me sur les
  homomorphismes de sch\'{e}mas ab\'{e}liens}, Invent. Math. \textbf{2}
  (1966), pp.~59--78.
\item[{[9]}] P. A. Griffiths, \emph{Report on variation of Hodge
  structures} (forthcoming).
\item[{[10]}] M. Hakim, \emph{Topos annel\'{e}s et sch\'{e}mas relatifs},
  forthcoming in Grundlehren der Mathematik, Springer (1970).
\item[{[11]}] J. I. Igusa, \emph{Abstract vanishing cycle theory},
  Proceedings of the Japan Academy \textbf{34}, no.~9 (1958),
  pp.~589--593.
\item[{[12]}] R. Kiehl, \emph{Der Endlichkeitssatz f\"{u}r eigentliche
  Abbildungen}, Invent. Math. \textbf{2} (1967), pp.~191--214.
\item[{[13]}] R. Kiehl.
\item[{[14]}] S. L. Kleiman, \emph{Algebraic cycles and the Weil
  conjectures}, in \emph{Dix expos\'{e}s sur la cohomologie des sch\'{e}mas},
  North-Holland Publishing Co.
\item[{[15]}] S. Lang, \emph{Abelian varieties}, Interscience Publishers,
  no.~7 (1959).
\item[{[16]}] S. Lang and A. N\'{e}ron, \emph{Rational points of abelian
  varieties over function fields}, Amer. J. Math. \textbf{81}, no.~1
  (1959), pp.~95--118.
\item[{[16bis]}] U. I. Manin, \emph{Th\'{e}orie des groupes formels
  commutatifs}, Uspekhi Mat. Nauk \textbf{XVIII} (1963), pp.~4--90.
\item[{[17]}] D. Mumford, \emph{Geometric invariant theory}, Ergebnisse
  der Mathematik 34, Springer (1965).
\item[{[18]}] D. Mumford, \emph{Abstract Theta functions}, Advanced
  Science Seminar in Algebraic Geometry, Bowdoin College, 1967 (notes by
  H. Pittie).
\item[{[19]}] A. N\'{e}ron, \emph{Mod\`{e}les minimaux des vari\'{e}t\'{e}s
  ab\'{e}liennes sur les corps locaux et globaux}, Publ. Math. 21 (1964),
  pp.~5--128.
\item[{[20]}] M. Raynaud, \emph{Caract\'{e}ristique d'Euler--Poincar\'{e}
  d'un faisceau et cohomologie des vari\'{e}t\'{e}s ab\'{e}liennes}, in
  \emph{Dix expos\'{e}s sur la cohomologie des sch\'{e}mas}, North-Holland
  Publishing Co.
\item[{[21]}] M. Raynaud, \emph{Mod\`{e}les de N\'{e}ron}, C. R. Acad.
  Sci. Paris \textbf{262} (February 1966), pp.~413--416.
\item[{[22]}] M. Raynaud, \emph{Sp\'{e}cialisation du foncteur de Picard},
  C. R. Acad. Sci. \textbf{264}, pp.~941--943 and 1001--1004.
\item[{[22bis]}] M. Raynaud, \emph{Sp\'{e}cialisation du foncteur de
  Picard}, forthcoming in Publ. Math.
\item[{[23]}] M. Raynaud, \emph{Faisceaux amples sur les sch\'{e}mas en
  groupes et les espaces homog\`{e}nes}, Paris thesis, 1968 (forthcoming
  in Lecture Notes, Springer, 1970).
\item[{[24]}] M. Raynaud.
\item[{[25]}] M. Rosenlicht, \emph{Some basic theorems on algebraic
  groups}, Amer. J. Math. \textbf{78}, no.~2 (1956), pp.~401--443.
\item[{[25bis]}] J.-P. Serre, \emph{Quelques propri\'{e}t\'{e}s des
  vari\'{e}t\'{e}s ab\'{e}liennes en caract\'{e}ristique $p$}, Amer. J.
  Math. \textbf{80} (1958), pp.~715--739.
\item[{[26]}] J.-P. Serre, \emph{Corps locaux}, Act. Sci. Ind. 1296
  (1962), Hermann, Paris.
\item[{[27]}] J.-P. Serre, \emph{Sur les corps locaux \`{a} corps
  r\'{e}siduel alg\'{e}briquement clos}, Bull. Soc. Math. France
  \textbf{89} (1961), pp.~105--154.
\item[{[28]}] J.-P. Serre, \emph{Abelian $\ell$-adic representations and
  elliptic curves}, Benjamin Inc. (1968).
\item[{[29]}] J.-P. Serre and J. Tate, \emph{Good reduction of abelian
  varieties}, Ann. of Math. \textbf{88}, no.~3 (1968), pp.~492--517.
\item[{[30]}] J.-P. Serre and J. Tate, \emph{Elliptic curves and formal
  groups} (notes of the Summer Institute on Algebraic Geometry, Woods
  Hole, 1964).
\item[{[31]}] J.-P. Serre, \emph{Morphismes universels et vari\'{e}t\'{e}
  d'Albanese}, S\'{e}minaire Chevalley 1958/59, Expos\'{e} 10.
\item[{[32]}] J. Tate, \emph{WC-groups over $p$-adic fields}, S\'{e}minaire
  Bourbaki no.~156 (1957).
\item[{[33]}] J. Tate, \emph{$p$-divisible groups}, Proceedings of a
  Conference on Local Fields, Driebergen (1966), Springer.
\item[{[34]}] J. Tate, \emph{Rigid analytic spaces} (mimeographed by
  IHES).
\end{enumerate}


% END INCLUDED SGA7I_English_current_buildable_source_20260823/components/186_expose_IX_bibliography.tex

% BEGIN INCLUDED SGA7I_English_current_buildable_source_20260823/components/187_expose_IX_terminal_publisher_matter.tex
\clearpage
\begingroup
\small
\sloppy
Lecture Notes in Mathematics


Comprehensive leaflet on request

Vol. 111: K. H. Mayer, Relationen zwischen charakteristischen Zahlen. III, 99 Seiten. 1969. DM 16,-

Vol. 112: Colloquium on Methods of Optimization. Edited by N. N. Moiseev. IV, 293 pages. 1970. DM 18,-

Vol. 113: R. Wille, Kongruenzklassengeometrien. III, 99 Seiten. 1970. DM 16,-

Vol. 114: H. Jacquet and R. P. Langlands, Automorphic Forms on GL (2). VII, 548 pages. 1970. DM 24,-

Vol. 115: K. H. Roggenkamp and V. Huber-Dyson, Lattices over Orders I. XIX, 290 pages. 1970. DM 18,-

Vol. 116: Sèminaire Pierre Lelong (Analyse) Année 1969. IV, 195 pages. 1970. DM 16,-

Vol. 117 : Y. Meyer, Nombres de Pisot, Nombres de Salem et Analyse Harmonique. 63 pages. 1970. DM 16,-

Vol. 118: Proceedings of the 15th Scandinavian Congress, Oslo 1968. Edited by K. E. Aubert and W. Ljunggren. IV, 162 pages. 1970. DM 16,-

Vol. 119: M. Raynaud, Faisceaux amples sur les schémas en groupes et les espaces homogènes. III, 219 pages. 1970. DM 16,-

Vol. 120: D. Siefkes, Büchi's Monadic Second Order Successor Arithmetic. XII, 130 Seiten. 1970. DM 16,-

Vol. 121: H. S. Bear, Lectures on Gleason Parts. III, 47 pages. 1970. DM 16,-

Vol. 122: H. Zieschang, E. Vogt und H.-D. Coldewey, Flächen und ebene diskontinuierliche Gruppen. VIII, 203 Seiten. 1970. DM 16,-

Vol. 123: A. V. Jategaonkar, Left Principal Ideal Rings.VI, 145 pages. 1970. DM 16,-

Vol. 124: Séminare de Probabilités IV. Edited by P. A. Meyer. IV, 282 pages. 1970. DM 20,-

Vol. 125: Symposium on Automatic Demonstration. V, 310 pages.1970. DM 20,-

Vol. 126: P. Schapira, Théorie des Hyperfonctions. XI,157 pages.1970. DM 16,-

Vol. 127: I. Stewart, Lie Algebras. IV, 97 pages. 1970. DM 16,-

Vol. 128: M. Takesaki, Tomita's Theory of Modular Hilbert Algebras and its Applications. II, 123 pages. 1970. DM 16,-

Vol. 129: K. H. Hofmann, The Duality of Compact Semigroups and C$^*$- Bigebras. XII, 142 pages. 1970. DM 16,-

Vol. 130: F. Lorenz, Quadratische Formen über Körpern. II, 77 Seiten. 1970. DM 16,-

Vol. 131: A Borel et al., Seminar on Algebraic Groups and Related Finite Groups. VII, 321 pages. 1970. DM 22,-

Vol. 132: Symposium on Optimization. III, 348 pages. 1970. DM 22,-

Vol. 133: F. Topsøe, Topology and Measure. XIV, 79 pages. 1970. DM 16,-

Vol. 134: L. Smith, Lectures on the Eilenberg-Moore Spectral Sequence. VII, 142 pages. 1970. DM 16,-

Vol. 135: W. Stoll, Value Distribution of Holomorphic Maps into Compact Complex Manifolds. II, 267 pages. 1970. DM 18,-

Vol. 136 : M. Karoubi et al., Séminaire Heidelberg-Saarbrücken-Strasbuorg sur la K-Théorie. IV, 264 pages. 1970. DM 18,-

Vol. 137 : Reports of the Midwest Category Seminar IV. Edited by S. MacLane. III, 139 pages. 1970. DM 16,-

Vol. 138: D. Foata et M. Schützenberger, Théorie Géométrique des Polynômes Eulériens. V, 94 pages. 1970. DM 16,-

Vol. 139: A. Badrikian, Séminaire sur les Fonctions Aléatoires Linéaires et les Mesures Cylindriques. VII, 221 pages. 1970. DM 18,-

Vol. 140: Lectures in Modern Analysis and Applications II. Edited by C. T. Taam. VI, 119 pages. 1970. DM 16,-

Vol. 141: G. Jameson, Ordered Linear Spaces. XV, 194 pages. 1970. DM 16,-

Vol. 142: K. W. Roggenkamp, Lattices over Orders II. V, 388 pages. 1970. DM 22,-

Vol. 143: K. W. Gruenberg, Cohomological Topics in Group Theory. XIV, 275 pages. 1970. DM 20,-

Vol. 144: Seminar on Differential Equations and Dynamical Systems, II. Edited by J. A. Yorke. VIII, 268 pages. 1970. DM 20,-

Vol. 145: E. J. Dubuc, Kan Extensions in Enriched Category Theory. XVI, 173 pages. 1970. DM 16,-

Vol. 146: A. B. Altman and S. Kleiman, Introduction to Grothendieck Duality Theory. II, 192 pages. 1970. DM 18,-

Vol. 147: D. E. Dobbs, Cech Cohomological Dimensions for Commutative Rings. VI, 176 pages. 1970. DM 16,-

Vol. 148: R. Azencott, Espaces de Poisson des Groupes Localement Compacts. IX, 141 pages. 1970. DM 16,-

Vol. 149: R. G. Swan and E. G. Evans, K-Theory of Finite Groups and Orders. IV, 237 pages. 1970. DM 20,-

Vol. 150: Heyer, Dualität lokalkompakter Gruppen. XIII, 372 Seiten. 1970. DM 20,-

Vol. 151: M. Demazure et A. Grothendieck, Schémas en Groupes I. (SGA 3). XV, 562 pages. 1970. DM 24,-

Vol. 152: M. Demazure et A. Grothendieck, Schémas en Groupes II. (SGA 3). IX, 654 pages. 1970. DM 24,-

Vol. 153: M. Demazure et A. Grothendieck, Schémas en Groupes III. (SGA 3). VIII, 529 pages. 1970. DM 24,-

Vol. 154: A. Lascoux et M. Berger, Variétés Kähleriennes Compactes. VII, 83 pages. 1970. DM 16,-

Vol. 155: Several Complex Variables I, Maryland 1970. Edited by J. Horváth. IV, 214 pages. 1970. DM 18,-

Vol. 156: R. Hartshorne, Ample Subvarieties of Algebraic Varieties. XIV, 256 pages. 1970. DM 20,-

Vol. 157: T. tom Dieck, K. H. Kamps und D. Puppe, Homotopietheorie. VI, 265 Seiten. 1970. DM 20,-

Vol. 158: T. G. Ostrom, Finite Translation Planes. IV. 112 pages. 1970. DM 16,-

Vol. 159: R. Ansorge und R. Hass. Konvergenz von Differenzenverfahren für lineare und nichtlineare Anfangswertaufgaben. VIII, 145 Seiten. 1970. DM 16,-

Vol. 160: L. Sucheston, Constributions to Ergodic Theory and Probability. VII, 277 pages. 1970. DM 20,-

Vol. 161: J. Stasheff, H-Spaces from a Homotopy Point of View. VI, 95 pages. 1970. DM 16,-

Vol. 162: Harish-Chandra and van Dijk, Harmonic Analysis on Reductive p-adic Groups. IV, 125 pages. 1970. DM 16,-

Vol. 163: P. Deligne, Equations Différentielles à Points Singuliers Reguliers. III, 133 pages. 1970. DM 16,-

Vol. 164: J. P. Ferrier, Seminaire sur les Algebres Complètes. II, 69 pages. 1970. DM 16,-

Vol. 165: J. M. Cohen, Stable Homotopy. V, 194 pages. 1970. DM 16,-

Vol. 166: A. J. Silberger, PGL$_2$ over the p-adics: its Representations, Spherical Functions, and Fourier Analysis. VII, 202 pages. 1970. DM 18,-

Vol. 167: Lavrentiev, Romanov and Vasiliev, Multidimensional Inverse Problems for Differential Equations. V, 59 pages. 1970. DM 16,-

Vol. 168: F. P. Peterson, The Steenrod Algebra and its Applications: A conference to Celebrate N. E. Steenrod's Sixtieth Birthday. VII, 317 pages. 1970. DM 22,-

Vol. 169: M. Raynaud, Anneaux Locaux Henséliens. V, 129 pages. 1970. DM 16,-

Vol. 170: Lectures in Modern Analysis and Applications III. Edited by C. T. Taam. VI, 213 pages. 1970. DM 18,-

Vol. 171: Set-Valued Mappings, Selections and Topological Properties of $2^X$. Edited by W. M. Fleischman. X, 110 pages. 1970. DM 16,-

Vol. 172: Y.-T. Siu and G. Trautmann, Gap-Sheaves and Extension of Coherent Analytic Subsheaves. V, 172 pages. 1971. DM 16,-

Vol. 173: J. N. Mordeson and B. Vinograde, Structure of Arbitrary Purely Inseparable Extension Fields. IV, 138 pages. 1970. DM 16,-

Vol. 174: B. Iversen, Linear Determinants with Applications to the Picard Scheme of a Family of Algebraic Curves. VI, 69 pages. 1970. DM 16,-

Vol. 175: M. Brelot, On Topologies and Boundaries in Potential Theory. VI, 176 pages. 1971. DM 18,-

Vol. 176: H. Popp, Fundamentalgruppen algebraischer Mannigfaltigkeiten. IV, 154 Seiten. 1970. DM 16,-

Vol. 177: J. Lambek, Torsion Theories, Additive Semantics and Rings of Quotients. VI, 94 pages. 1971. DM 16,-

Please turn over


%%%%%%%%%%  scan idx 537  |  folio 526  %%%%%%%%%%

Vol. 178: Th. Bröcker und T. tom Dieck, Kobordismentheorie. XVI, 191 Seiten. 1970. DM 18,-

Vol. 179: Seminaire Bourbaki -- vol. 1968/69. Exposés 347-363. IV. 295 pages. 1971. DM 22,-

Vol. 180: Séminaire Bourbaki -- vol. 1969/70. Exposés 364-381. IV, 310 pages. 1971. DM 22,-

Vol. 181: F. DeMeyer and E. Ingraham, Separable Algebras over Commutative Rings. V, 157 pages. 1971. DM 16.-

Vol. 182: L. D. Baumert. Cyclic Difference Sets. VI, 166 pages. 1971. DM 16,-

Vol. 183: Analytic Theory of Differential Equations. Edited by P. F. Hsieh and A. W. J. Stoddart. VI, 225 pages. 1971. DM 20,-

Vol. 184: Symposium on Several Complex Variables, Park City, Utah, 1970. Edited by R. M. Brooks. V, 234 pages. 1971. DM 20,-

Vol. 185: Several Complex Variables II, Maryland 1970. Edited by J. Horváth. III, 287 pages. 1971. DM 24,-

Vol. 186: Recent Trends in Graph Theory. Edited by M. Capobianco/J. B. Frechen/M. Krolik. VI, 219 pages. 1971. DM 18.-

Vol. 187: H. S. Shapiro, Topics in Approximation Theory. VIII, 275 pages. 1971. DM 22,-

Vol. 188: Symposium on Semantics of Algorithmic Languages. Edited by E. Engeler. VI, 372 pages. 1971. DM 26,-

Vol. 189: A. Weil, Dirichlet Series and Automorphic Forms. V, 164 pages. 1971. DM 16,-

Vol. 190: Martingales. A Report on a Meeting at Oberwolfach, May 17-23, 1970. Edited by H. Dinges. V, 75 pages. 1971. DM 16,-

Vol. 191: Séminaire de Probabilités V. Edited by P. A. Meyer. IV, 372 pages. 1971. DM 26,-

Vol. 192: Proceedings of Liverpool Singularities -- Symposium I. Edited by C. T. C. Wall. V, 319 pages. 1971. DM 24,-

Vol. 193: Symposium on the Theory of Numerical Analysis. Edited by J. Ll. Morris. VI, 152 pages. 1971. DM 16,-

Vol. 194: M. Berger, P. Gauduchon et E. Mazet. Le Spectre d'une Variété Riemannienne. VII, 251 pages. 1971. DM 22,-

Vol. 195: Reports of the Midwest Category Seminar V. Edited by J.W. Gray and S. Mac Lane.III, 255 pages. 1971. DM 22,-

Vol. 196: H-spaces -- Neuchâtel (Suisse)- Août 1970. Edited by F. Sigrist, V, 156 pages. 1971. DM 16,-

Vol. 197: Manifolds -- Amsterdam 1970. Edited by N. H. Kuiper. V, 231 pages. 1971. DM 20,-

Vol. 198: M. Hervé, Analytic and Plurisubharmonic Functions in Finite and Infinite Dimensional Spaces. VI, 90 pages. 1971. DM 16.-

Vol. 199: Ch. J. Mozzochi, On the Pointwise Convergence of Fourier Series. VII, 87 pages. 1971. DM 16,-

Vol. 2 00: U. Neri, Singular Integrals. VII, 272 pages. 1971. DM 22,-

Vol. 201 : J. H. van Lint, Coding Theory. VII, 136 pages. 1971. DM 16,-

Vol. 202: J. Benedetto, Harmonic Analysis on Totally Disconnected Sets. VIII, 261 pages. 1971. DM 22,-

Vol. 203: D. Knutson, Algebraic Spaces. VI, 261 pages. 1971. DM 22,-

Vol. 204: A. Zygmund, Intégrales Singulières. IV, 53 pages. 1971. DM 16,-

Vol. 205: Séminaire Pierre Lelong (Analyse) Année 1970. VI, 243 pages. 1971. DM 20,-

Vol. 206: Symposium on Differential Equations and Dynamical Systems. Edited by D. Chillingworth. XI, 173 pages. 1971. DM 16,-

Vol. 207: L. Bernstein, The Jacobi-Perron Algorithm -- Its Theory and Application. IV, 161 pages. 1971. DM 16,-

Vol. 208: A. Grothendieck and J. P. Murre, The Tame Fundamental Group of a Formal Neighbourhood of a Divisor with Normal Crossings on a Scheme. VIII, 133 pages. 1971. DM 16,-

Vol. 209: Proceedings of Liverpool Singularities \quad Symposium II. Edited by C. T. C. Wall. V, 280 pages. 1971. DM 22,-

Vol. 210: M. Eichler, Projective Varieties and Modular Forms. III, 118 pages. 1971. DM 16,-

Vol. 211: Théorie des Matroïdes. Edité par C. P. Bruter. III, 108 pages. 1971. DM 16,-

Vol. 212: B. Scarpellini, Proof Theory and Intuitionistic Systems. VII, 291 pages. 1971. DM 24,-

Vol. 213: H. Hogbe-Nlend, Théorie des Bornologies et Applications. V, 168 pages. 1971. DM 18,-

Vol. 214: M. Smorodinsky, Ergodic Theory, Entropy. V, 64 pages. 1971. DM 16,-

Vol. 215: P. Antonelli, D. Burghelea and P. J. Kahn, The Concordance-Homotopy Groups of Geometric Automorphism Groups. X, 140 pages. 1971. DM 16,-

Vol. 216: H. Maaß, Siegel's Modular Forms and Dirichlet Series. VII, 328 pages. 1971. DM 20,-

Vol. 217: T. J. Jech, Lectures in Set Theory with Particular Emphasis on the Method of Forcing. V, 137 pages. 1971. DM 16,-

Vol. 218: C. P. Schnorr, Zufälligkeit und Wahrscheinlichkeit. IV, 212.Seiten 1971. DM 20,-

Vol. 219: N. L. Alling and N. Greenleaf, Foundations of the Theory of Klein Surfaces. IX, 117 pages. 1971. DM 16,-

Vol. 220: W. A. Coppel, Disconjugacy. V, 148 pages. 1971. DM 16,-

Vol. 221: P. Gabriel und F. Ulmer, Lokal präsentierbare Kategorien. V, 200 Seiten. 1971. DM 18,-

Vol. 222: C. Meghea, Compactification des Espaces Harmoniques. III, 108 pages. 1971. DM 16,-

Vol. 223: U. Felgner, Models of ZF-Set Theory. VI, 173 pages. 1971. DM 16,-

Vol. 224: Revètements Etales et Groupe Fondamental. (SGA 1). Dirigé par A. Grothendieck XXII, 447 pages. 1971. DM 30,-

Vol. 225: Théorie des Intersections et Théorème de Riemann-Roch. (SGA 6). Dirigé par P. Berthelot, A. Grothendieck et L. Illusie. XII, 700 pages. 1971. DM 40,-

Vol. 226: Seminar on Potential Theory, II. Edited by H. Bauer. IV, 170 pages. 1971. DM 18,-

Vol. 227: H. L. Montgomery, Topics in Multiplicative Number Theory. IX, 178 pages. 1971. DM 18,-

Vol. 228: Conference on Applications of Numerical Analysis. Edited by J. Ll. Morris. X, 358 pages. 1971. DM 26,-

Vol. 229: J. Väisälä, Lectures on n-Dimensional Quasiconformal Mappings. XIV, 144 pages. 1971. DM 16,-

Vol. 230: L. Waelbroeck, Topological Vector Spaces and Algebras. VII, 158 pages. 1971. DM 16,-

Vol. 231: H. Reiter, L$^1$-Algebras and Segal Algebras. XI, 113 pages. 1971. DM 16,-

Vol. 232: T. H. Ganelius, Tauberian Remainder Theorems. VI, 75 pages. 1971. DM 16,-

Vol. 233: C. P. Tsokos and W. J. Padgett. Random Integral Equations with Applications to Stochastic Systems. VII, 174 pages. 1971. DM 18,-

Vol. 234: A. Andreotti and W. Stoll. Analytic and Algebraic Dependence of Meromorphic Functions. III, 390 pages. 1971. DM 26,-

Vol. 235: Global Differentiable Dynamics. Edited by O. Hájek, A. J. Lohwater, and R. McCann. X, 140 pages. 1971. DM 16,-

Vol. 236 : M. Barr, P. A. Grillet, and D. H. van Osdol. Exact Categories and Categories of Sheaves. VII, 239 pages. 1971, DM 20,-

Vol. 237: B. Stenström. Rings and Modules of Quotients. VII, 136 pages. 1971. DM 16,-

Vol. 238: Der kanonische Modul eines Cohen-Macaulay-Rings. Herausgegeben von Jürgen Herzog und Ernst Kunz. VI, 103 Seiten. 1971. DM 16,-

Vol. 239: L. Illusie, Complexe Cotangent et Déformations I. XV, 355 pages. 1971. DM 26,-

Vol. 240: A. Kerber, Representations of Permutation Groups I. VII, 192 pages. 1971. DM 18,-

Vol. 241: S. Kaneyuki, Homogeneous Bounded Domains and Siegel Domains. V, 89 pages. 1971. DM 16,-

Vol. 242: R. R. Coifman et G. Weiss, Analyse Harmonique Non-Commutative sur Certains Espaces. V, 160 pages. 1971. DM 16,-

Vol. 243: Japan-United States Seminar on Ordinary Differential and Functional Equations. Edited by M. Urabe. VIII, 332 pages. 1971. DM 26,-

Vol. 244: Séminaire Bourbaki -- vol. 1970/71. Exposés 382-399. IV, 356 pages. 1971. DM 26,-

Vol. 245: D. E. Cohen, Groups of Cohomological Dimension One. V, 99 pages. 1972. DM 16,-


%%%%%%%%%%  scan idx 538  |  folio 527  %%%%%%%%%%

\noindent\rule{\linewidth}{0.4pt}

Vol. 246: Lectures on Rings and Modules. Tulane University Ring and Operator Theory Year, 1970--1971. Volume I. X, 661 pages. 1972. DM 40,--

Vol. 247: Lectures on Operator Algebras. Tulane University Ring and Operator Theory Year, 1970--1971. Volume II. XI, 786 pages. 1972. DM 40,--

Vol. 248: Lectures on the Applications of Sheaves to Ring Theory. Tulane University Ring and Operator Theory Year, 1970--1971. Volume III. VIII, 315 pages. 1971. DM 26,--

Vol. 249: Symposium on Algebraic Topology. Edited by P. J. Hilton. VII, 111 pages. 1971. DM 16,--

Vol. 250: B. Jónsson, Topics in Universal Algebra. VI, 220 pages. 1972. DM 20,--

Vol. 251: The Theory of Arithmetic Functions. Edited by A. A. Gioia and D. L. Goldsmith VI, 287 pages. 1972. DM 24,--

Vol. 252: D. A. Stone, Stratified Polyhedra. IX, 193 pages. 1972. DM 18,--

Vol. 253: V. Komkov, Optimal Control Theory for the Damping of Vibrations of Simple Elastic Systems. V, 240 pages. 1972. DM 20,--

Vol. 254: C. U. Jensen, Les Foncteurs Dérivés de $\varprojlim$ et leurs Applications en Théorie des Modules. V, 103 pages. 1972. DM 16,--

Vol. 255: Conference in Mathematical Logic -- London '70. Edited by W. Hodges. VIII, 351 pages. 1972. DM 26,--

Vol. 256: C. A. Berenstein and M. A. Dostal, Analytically Uniform Spaces and their Applications to Convolution Equations. VII, 130 pages. 1972. DM 16,--

Vol. 257: R. B. Holmes, A Course on Optimization and Best Approximation. VIII, 233 pages. 1972. DM 20,--

Vol. 258: Séminaire de Probabilités VI. Edited by P. A. Meyer. VI, 253 pages. 1972. DM 22,--

Vol. 259: N. Moulis, Structures de Fredholm sur les Variétés Hilbertiennes. V, 123 pages. 1972. DM 16,--

Vol. 260: R. Godement and H. Jacquet, Zeta Functions of Simple Algebras. IX, 188 pages. 1972. DM 18,--

Vol. 261: A. Guichardet, Symmetric Hilbert Spaces and Related Topics. V, 197 pages. 1972. DM 18,--

Vol. 262: H. G. Zimmer, Computational Problems, Methods, and Results in Algebraic Number Theory. V, 103 pages. 1972. DM 16,--

Vol. 263: T. Parthasarathy, Selection Theorems and their Applications. VII, 101 pages. 1972. DM 16,--

Vol. 264: W. Messing, The Crystals Associated to Barsotti-Tate Groups: with Applications to Abelian Schemes. III, 190 pages. 1972. DM 18,--

Vol. 265: N. Saavedra Rivano, Catégories Tannakiennes. II, 418 pages. 1972. DM 26,--

Vol. 266: Conference on Harmonic Analysis. Edited by D. Gulick and R. L. Lipsman. VI, 323 pages. 1972. DM 24,--

Vol. 267: Numerische Lösung nichtlinearer partieller Differential- und Integro-Differentialgleichungen. Herausgegeben von R. Ansorge und W. Törnig, VI, 339 Seiten. 1972. DM 26,--

Vol. 268: C. G. Simader, On Dirichlet's Boundary Value Problem. IV, 238 pages. 1972. DM 20,--

Vol. 269: Théorie des Topos et Cohomologie Etale des Schémas. (SGA 4). Dirigé par M. Artin, A. Grothendieck et J. L. Verdier. XIX, 525 pages. 1972. DM 50,--

Vol. 270: Théorie des Topos et Cohomologie Etale des Schémas. Tome 2. (SGA 4). Dirigé par M. Artin, A. Grothendieck et J. L. Verdier. V, 314 pages. 1972. DM 50,--

Vol. 271: J. P. May, The Geometry of Iterated Loop Spaces. IX, 175 pages. 1972. DM 18,--

Vol. 272: K. R. Parthasarathy and K. Schmidt, Positive Definite Kernels, Continuous Tensor Products, and Central Limit Theorems of Probability Theory. VI, 107 pages. 1972. DM 16,--

Vol. 273: U. Seip, Kompakt erzeugte Vektorräume und Analysis. IX, 119 Seiten. 1972. DM 16,--

Vol. 274: Toposes, Algebraic Geometry and Logic. Edited by. F. W. Lawvere. VI, 189 pages. 1972. DM 18,--

Vol. 275: Séminaire Pierre Lelong (Analyse) Année 1970--1971. VI, 181 pages. 1972. DM 18,--

Vol. 276: A. Borel, Représentations de Groupes Localement Compacts. V, 98 pages. 1972. DM 16,--

Vol. 277: Séminaire Banach. Edité par C. Houzel. VII, 229 pages. 1972. DM 20,--

Vol. 278: H. Jacquet, Automorphic Forms on GL(2). Part II. XIII, 142 pages. 1972. DM 16,--

Vol. 279: R. Bott, S. Gitler and I. M. James, Lectures on Algebraic and Differential Topology. V, 174 pages. 1972. DM 18,--

Vol. 280: Conference on the Theory of Ordinary and Partial Differential Equations. Edited by W. N. Everitt and B. D. Sleeman. XV, 367 pages. 1972. DM 26,--

Vol. 281: Coherence in Categories. Edited by S. Mac Lane. VII, 235 pages. 1972. DM 20,--

Vol. 282: W. Klingenberg und P. Flaschel, Riemannsche Hilbertmannigfaltigkeiten. Periodische Geodätische. VII, 211 Seiten. 1972. DM 20,--

Vol. 283: L Illusie, Complexe Cotangent et Déformations II. VII, 304 pages. 1972. DM 24,--

Vol. 284: P. A. Meyer, Martingales and Stochastic Integrals I. VI, 89 pages. 1972. DM 16,--

Vol. 285: P. de la Harpe, Classical Banach-Lie Algebras and Banach-Lie Groups of Operators in Hilbert Space. III, 160 pages. 1972. DM 16,--

Vol. 286: S. Murakami, On Automorphisms of Siegel Domains. V, 95 pages. 1972. DM 16,--

Vol. 288: Groupes de Monodromie en Géométrie Algébrique. (SGA 7 I). Dirigé par A. Grothendieck. IX, 523 pages. 1972. DM 50,--

\noindent\rule{\linewidth}{0.4pt}


%%%%%%%%%%  scan idx 539  |  folio 528  %%%%%%%%%%

ISBN 3-540-05987-3

ISBN 0-387-05987-3
\endgroup


% END INCLUDED SGA7I_English_current_buildable_source_20260823/components/187_expose_IX_terminal_publisher_matter.tex

\end{document}
